\documentclass[11pt]{article}

\usepackage[margin=1in]{geometry}
\usepackage{amsmath,amssymb,amsthm,mathtools,mathrsfs,bm}
\usepackage{enumitem}
\usepackage{booktabs,longtable,array,tabularx}
\usepackage{xcolor}
\usepackage{microtype}
\usepackage[colorlinks=true,linkcolor=blue,citecolor=blue,urlcolor=blue]{hyperref}
\usepackage[nameinlink,capitalise]{cleveref}
\setlength{\emergencystretch}{4em}
\raggedbottom

\newtheorem{theorem}{Theorem}[section]
\newtheorem{proposition}[theorem]{Proposition}
\newtheorem{lemma}[theorem]{Lemma}
\newtheorem{corollary}[theorem]{Corollary}
\newtheorem{countertheorem}[theorem]{Countertheorem}
\newtheorem{openproblem}[theorem]{Open Problem}
\theoremstyle{definition}
\newtheorem{definition}[theorem]{Definition}
\newtheorem{construction}[theorem]{Construction}
\newtheorem{obstruction}[theorem]{Obstruction}
\newtheorem{acquisition}[theorem]{Acquisition Card}
\theoremstyle{remark}
\newtheorem{remark}[theorem]{Remark}
\newtheorem{assessment}[theorem]{Assessment}
\newtheorem{firewall}[theorem]{Firewall}

\newcommand{\C}{\mathbb C}
\newcommand{\R}{\mathbb R}
\newcommand{\N}{\mathbb N}
\newcommand{\one}{\mathbf 1}
\newcommand{\dd}{\,\mathrm d}
\newcommand{\Ran}{\operatorname{Ran}}
\newcommand{\Ker}{\operatorname{Ker}}
\newcommand{\rank}{\operatorname{rank}}
\newcommand{\supp}{\operatorname{supp}}
\newcommand{\spec}{\operatorname{spec}}
\newcommand{\Tr}{\operatorname{Tr}}
\newcommand{\dist}{\operatorname{dist}}
\newcommand{\diag}{\operatorname{diag}}
\newcommand{\Span}{\operatorname{span}}
\newcommand{\norm}[1]{\left\lVert #1\right\rVert}
\newcommand{\abs}[1]{\left\lvert #1\right\rvert}
\newcommand{\ip}[2]{\left\langle #1,#2\right\rangle}
\newcommand{\adj}{^{*}}
\newcommand{\Psrc}{P_{B}}
\newcommand{\Rdyn}{R_{\mathrm{dyn}}}
\newcommand{\RDuh}{R^{\mathrm{Duh}}}
\newcommand{\ThetaD}{\Theta^{\mathrm{Duh}}}
\newcommand{\PhiLev}{\Phi^{\mathrm{lev}}}
\newcommand{\Gres}{\mathsf G^{\mathrm{res}}}
\newcommand{\Gaux}{\mathsf G_{\mathrm{aux}}}
\newcommand{\Caux}{\mathsf C_{\mathrm{aux}}}
\newcommand{\Msrc}{\mathsf M}
\newcommand{\Ntar}{\mathsf N}
\newcommand{\Wphys}{\mathscr W^{\Msrc}}
\newcommand{\What}{\widehat{\mathscr W}}
\newcommand{\Rphys}{\mathscr R^{\Msrc}}
\newcommand{\Rhat}{\widehat{\mathscr R}}

\definecolor{deepgreen}{RGB}{24,112,68}
\definecolor{deepred}{RGB}{150,42,42}
\definecolor{deepblue}{RGB}{35,74,145}
\newcommand{\passstatus}{\textcolor{deepgreen}{\textbf{PASS}}}
\newcommand{\obsstatus}{\textcolor{deepred}{\textbf{TYPED OBSTRUCTION}}}
\newcommand{\stopstatus}{\textcolor{deepblue}{\textbf{STOP}}}

\title{\textbf{Renewal Einstein--Standard--Model Physical Source Metric,\\
Duhamel Ancestry and Quantum Continuum Programme}\\[0.45em]
\large Version V1: Paired Reserve--Duhamel Cards, an Exact Two-Tail Squeeze,\\
and the First Terminating Continuum Handoff}
\author{A.~P\'elissier}
\date{2 September 2026}

\begin{document}
\maketitle

\begin{abstract}
This is the first note in a fresh consolidated Einstein--Standard--Model
continuum series.  It does not restart the finite field-shell, determinant-line,
reflection, domain, clock, source-ancestry, or compact-horizon constructions
already proved in the attached programmes.  Instead it imports two completed
theorem layers.  The sector-patched line/clock series reduces the internal
Duhamel target, on every left-energy fibre, to five finite same-card matrices
and a positive omitted-left-energy tail.  The main continuum series proves that
positive leverage is physically meaningful only after the source-action metric
has been frozen; it then separates unit-action reserve, polar-normalized
source-orbit geometry, the small-positive leverage tail, and the exact
zero-leverage residual.

The first new result is a single-card compression.  One positive static Gram of
the full primitive bank and the Duhamel target gives the complete dynamic
residual by one supported Schur complement.  The five-matrix formula is the
orthogonal nested factorization of that same complement and remains available
when attribution between right-clock variance and prior-primitive absorption
is required.  The source-action metric, the static Gram, and the raw delayed
source/source--target kernels then reconstruct both the physical-metric and the
polar-normalized finite-horizon compilers without an independently chosen
whitening or source orientation.

The central theorem is an exact lower--upper sandwich.  A finite left-energy
Duhamel head is a positive lower bound for the exact ancestry residual, while a
physical-metric finite-horizon Tikhonov card is a positive upper bound.  Their
difference is exactly the sum of two positive quantities: the omitted
left-energy Duhamel tail and the small-positive-leverage tail.  Consequently one
same-card cross-gap controls both tails at once.  A positive Duhamel head proves
a dynamic birth without completing the horizon tail; a vanishing horizon upper
card proves followerhood without resolving every left-energy fibre; and a
cofinal collapse of the cross-gap squeezes the exact residual to a unique
transported limit.  Failure returns a typed source-metric obstruction, reserve
collapse, Duhamel-incidence defect, support transition, clock/route defect, or a
normalized extremizing sequence living in one of the two positive tails.

The attached corpus does not contain the actual routed source-action matrix,
cardwise full primitive/target Grams, or complete delayed row-summed weak kernel
on one selected ESM card.  Version~V1 therefore proves the combined compiler
and its stopping theorem but does not invent a numerical residual.  Current,
stress, determinant/Pfaffian orientation, interacting continuum, Lorentzian,
Einstein, and pure-colour conclusions remain closed until the paired physical
card is populated and transported.
\end{abstract}

\tableofcontents

% =============================================================================
\section{Context, imported theorem layer, and append-only rule}
\label{sec:context}
% =============================================================================

\subsection{Why this fresh programme is narrower}

The mature programme architecture has reached a finite-physical-value boundary.
The relevant chain is
\begin{equation}
 \boxed{
 \begin{gathered}
 \text{physical source metric and unit-action reserve}
 \longrightarrow
 \text{polar-normalized compact-horizon geometry}\\
 \longrightarrow
 \text{Duhamel fibre residual and cofinal ancestry packet}\\
 \longrightarrow
 \text{current/stress incidence}
 \longrightarrow
 \text{Einstein handoff}.
 \end{gathered}}
 \label{eq:master-chain}
\end{equation}
The first three arrows are the sole active target here.  The last two are kept
visible only to state precisely what the ancestry packet must export.  No new
general-purpose Hilbert-space compiler is introduced.

\subsection{Imported results which are not reproved}

The following facts are treated as theorem inputs.

\begin{enumerate}[label=\textup{(I\arabic*)},leftmargin=3.2em]
\item The sector-patched fermion-line and physical-clock programme constructs a
      target-independent line/reflection/domain/$H_Q$/OS--KMS card and proves
      that the surviving internal target satisfies, on each left-energy fibre,
      a calibrated Duhamel relation
      \[
       Y_\lambda=\mathfrak m_\lambda(H^R_\lambda)A_\lambda.
      \]
\item On that card the full primitive fibre residual is positive and has the
      five-matrix representation
      \[
       R_\lambda^{\mathrm{Duh}}
       =U_\lambda-T_\lambda^*S_\lambda^\dagger T_\lambda
        -C_\lambda^*G_\lambda^\dagger C_\lambda.
      \]
      Its positive norm-convergent sum is the complete dynamic residual, with a
      positive omitted-left-energy tail for every finite head.
\item The main quantum-state/dynamic-source/continuum programme freezes a
      physical source-action form $\Msrc$, proves the null-cost condition
      $\Ker\Msrc\subseteq\Ker B$, constructs the unit-action covariance
      $B\Msrc^\dagger B^*$ and its supported reserve, and then constructs the
      range-normalized horizon Gramian
      \[
       \What_{T,B}=\int_0^T e^{-tH}\Psrc e^{-tH}\dd t.
      \]
      Its zero atom is the same $\Rdyn$.
\item Under a two-sided reserve window, physical-metric and range-normalized
      Tikhonov tails are equivalent after rescaling the regularization
      parameter.  Source-map, source-metric, source-range, and clock transport
      remain separate primitive rows.
\item The Grand-Tensor terminal order keeps line orientation, locality,
      represented current/stress, charge, regional occurrence, joint spectrum,
      and Einstein closure as later gates.  A dynamic-source result alone does
      not populate any of them.
\end{enumerate}

The role of the present note is to prove the missing \emph{compatibility and
squeeze theorem} between items (I2) and (I3), reduce their common physical
acquisition, and state an exact terminating card.

\begin{firewall}[No repetition of the imported programmes]
\label{fire:no-repetition}
The proofs of the line construction, reflection positivity, common graph
core, KMS identity, Duhamel divided difference, source-action variational
formula, polar reserve factorization, finite-horizon support theorem,
block-Jacobi depth theorem, and primitive cutoff transport are not repeated.
Whenever one of those statements is used, it is named as an imported theorem
layer.  New proofs begin only where the two layers are assembled on one card.
\end{firewall}

\subsection{Append-only convention}

This file begins the series
\begin{center}
\path{Renewal_Einstein_Standard_Model_Physical_Source_Metric_Duhamel_Ancestry_and_Quantum_Continuum_Programme_V1.tex}.
\end{center}
Every later version must preserve the complete preceding source verbatim,
remove only its sole final \verb|\end{document}|, append the new material, and
restore one final \verb|\end{document}|.  The version number in the filename is
incremented.  Previously proved statements remain in place; later notes add
new theorems, corrections, or typed obstructions at the end rather than
silently rewriting the record.

\begin{definition}[Three terminal programme outcomes]
\label{def:three-outcomes}
Every iteration must end in exactly one of the following classes:
\begin{equation}
 \boxed{
 \mathsf{PASS}
 \quad\big|\quad
 \mathsf{TYPED\ OBSTRUCTION}
 \quad\big|\quad
 \mathsf{STOP}.}
 \label{eq:three-outcomes}
\end{equation}
A pass proves the declared finite or cofinal branch.  A typed obstruction
returns an explicit vector, positive residual, support transition, reserve
collapse, tail survivor, or route defect.  A stop means that the theorem layer
has terminated and the next missing object is a named physical matrix or
kernel.  ``Develop more compiler machinery'' is not a fourth outcome.
\end{definition}

% =============================================================================
\section{The compatible paired source--Duhamel card}
\label{sec:paired-card}
% =============================================================================

All Hilbert spaces are complex.  At one fixed cutoff, coefficient spaces are
finite dimensional; the physical carrier may be finite dimensional or a
Hilbert space on which the displayed finite-rank syntheses and semigroups are
well defined.

Let $D$ be the common-action/deformation coefficient space, let $E_{\rm aux}$
carry every independently occurring primitive which precedes the dynamic
short, and put
\begin{equation}
 E=D\oplus E_{\rm aux}.
 \label{eq:E-split}
\end{equation}
For each retained left energy $\lambda$, let
\begin{equation}
 A_\lambda:D\to\mathcal K_\lambda,
 \qquad
 F_\lambda:E_{\rm aux}\to\mathcal K_\lambda,
 \qquad
 B_\lambda=(A_\lambda,F_\lambda):E\to\mathcal K_\lambda.
 \label{eq:fibre-sources}
\end{equation}
The Duhamel target is
\begin{equation}
 Y_\lambda=\mathfrak m_\lambda(H^R_\lambda)A_\lambda:D\to\mathcal K_\lambda,
 \label{eq:Duhamel-relation}
\end{equation}
where $H^R_\lambda$ is the inherited right-energy operator and
$\mathfrak m_\lambda$ is the calibrated real multiplier.  The symbol
$\mathfrak m_\lambda$ is deliberately distinct from the physical source-action
metric $\Msrc$ below.

Let
\begin{equation}
 \mathcal K=\bigoplus_\lambda\mathcal K_\lambda,
 \qquad
 B=\bigoplus_\lambda B_\lambda,
 \qquad
 Y=\bigoplus_\lambda Y_\lambda,
 \label{eq:global-BY}
\end{equation}
and let $H=H^*\succeq0$ be the completed reflected Hamiltonian used by the
finite-horizon ancestry compiler.  Let
\begin{equation}
 \Msrc=\Msrc^*\succeq0\quad\text{on }E
 \label{eq:source-metric}
\end{equation}
be the physical source-action form after the deterministic profile, Higgs/
electroweak router, Ward quotient, and primitive-role assembly have all been
performed on the same unquenched history.

\begin{definition}[Compatible paired card]
\label{def:compatible-card}
A \emph{compatible paired source--Duhamel card} consists of:
\begin{enumerate}[label=\textup{(P\arabic*)},leftmargin=3.2em]
\item one target-independent occurring history carrying $B$, $Y$, $H$, and
      $\Msrc$;
\item the left-energy projections and the literal relations
      \eqref{eq:Duhamel-relation};
\item the full static positive Grams of $(A_\lambda,F_\lambda,Y_\lambda)$;
\item the delayed raw blocks
      \begin{equation}
       K_{BB}(t)=B^*e^{-tH}B,
       \qquad
       K_{BY}(t)=B^*e^{-tH}Y,
       \qquad 0\le t\le2T;
       \label{eq:raw-delayed-blocks}
      \end{equation}
\item the inherited same-card identification
      \begin{equation}
       \boxed{
       Y^*(I-P_H^B)Y
       =\sum_\lambda R_\lambda^{\mathrm{Duh}}.}
       \label{eq:compatibility-identification}
      \end{equation}
\end{enumerate}
Here $P_H^B$ is the projection onto the completed $H$-cyclic source space.
Item (P5) is not asserted for an arbitrary abstract pair $(B,H)$; it is the
specific compatibility already proved for the inherited ESM line/clock and
main ancestry cards.
\end{definition}

\begin{firewall}[Same history before every short]
\label{fire:same-history}
The metric $\Msrc$, primitive synthesis $B$, target $Y$, static mixed Gram, and
delayed blocks are assembled before any source whitening, role short, polar
normalization, left-energy truncation, or Tikhonov inversion.  Separate
averaging of source orientations or fitting a mixed block after inspecting the
residual does not produce a compatible paired card.
\end{firewall}

% =============================================================================
\section{One static Gram, one supported Schur complement}
\label{sec:static-Schur}
% =============================================================================

The V15 five-matrix packet is indispensable when the amount of right-clock
variance absorbed by prior primitives is itself an output.  For the main
continuum ancestry verdict, however, the complete residual can be obtained
from one full primitive block and one supported pseudoinverse.  This section
proves that compression.

\begin{lemma}[Support incidence of a positive block]
\label{lem:positive-block-support}
Let
\begin{equation}
 \Gamma=\begin{pmatrix}Q&L\\L^*&U\end{pmatrix}\succeq0
 \label{eq:positive-two-block}
\end{equation}
with $Q=Q^*\succeq0$.  Then
\begin{equation}
 \boxed{\Ker Q\subseteq\Ker L^*,
 \qquad \Ran L\subseteq\Ran Q.}
 \label{eq:block-support-incidence}
\end{equation}
Consequently the supported Schur complement
\begin{equation}
 U-L^*Q^\dagger L
 \label{eq:supported-Schur}
\end{equation}
is well defined and positive.
\end{lemma}

\begin{proof}
Let $x\in\Ker Q$ and let $y$ be arbitrary.  Positivity of \eqref{eq:positive-two-block}
implies, for every real $t$,
\[
 0\le
 \left\langle\binom{x}{ty},
 \Gamma\binom{x}{ty}\right\rangle
 =2t\operatorname{Re}\langle x,Ly\rangle+t^2\langle y,Uy\rangle.
\]
The linear coefficient must vanish.  Replacing $y$ by $\mathrm i y$ also kills
the imaginary part, so $L^*x=0$.  In finite dimension,
$\Ker Q\subseteq\Ker L^*$ is equivalent to
$\Ran L\subseteq(\Ker Q)^\perp=\Ran Q$.

Write $L=Q^{1/2}C$ on $\supp Q$.  Positivity of the block is equivalent, after
completion of the square on that support, to $U-C^*C\succeq0$.  Since
$C=Q^{\dagger/2}L$, this is precisely
$U-L^*Q^\dagger L\succeq0$.
\end{proof}

\begin{theorem}[Direct full-primitive residual]
\label{thm:direct-residual}
Let $B:E\to\mathcal K$ and $Y:D\to\mathcal K$ be finite-rank syntheses and set
\begin{equation}
 Q=B^*B,
 \qquad
 L=B^*Y,
 \qquad
 U=Y^*Y.
 \label{eq:QLU}
\end{equation}
Then
\begin{equation}
 P_B:=BQ^\dagger B^*
 \label{eq:PB-Q}
\end{equation}
is the orthogonal projection onto $\Ran B$, and
\begin{equation}
 \boxed{
 R^{\mathrm{full}}
 :=U-L^*Q^\dagger L
 =Y^*(I-P_B)Y\succeq0.}
 \label{eq:direct-full-residual}
\end{equation}
Moreover,
\begin{equation}
 \boxed{
 \Ker R^{\mathrm{full}}
 =\{d\in D:Yd\in\Ran B\},
 \qquad
 \rank R^{\mathrm{full}}
 =\dim\Ran((I-P_B)Y).}
 \label{eq:direct-rank-kernel}
\end{equation}
Thus the full primitive ancestry verdict is determined by the one positive
Gram
\begin{equation}
 \begin{pmatrix}Q&L\\L^*&U\end{pmatrix}
 =[B,Y]^*[B,Y].
 \label{eq:one-static-Gram}
\end{equation}
\end{theorem}

\begin{proof}
The Moore--Penrose range formula gives
$BQ^\dagger B^*=B(B^*B)^\dagger B^*=P_{\Ran B}$.  Therefore
\[
 L^*Q^\dagger L
 =Y^*BQ^\dagger B^*Y
 =Y^*P_BY,
\]
which proves \eqref{eq:direct-full-residual}.  If
$Z=(I-P_B)Y$, then $R^{\mathrm{full}}=Z^*Z$.  Hence
$R^{\mathrm{full}}d=0$ exactly when $Zd=0$, equivalently
$Yd\in\Ran B$, and the rank of the Gram $Z^*Z$ equals the rank of $Z$.
\end{proof}

\begin{theorem}[Nested Duhamel attribution from the same Gram]
\label{thm:nested-attribution}
Let $B=(A,F)$ and write the complete static Gram as
\begin{equation}
 \Gamma_{AFY}
 =\begin{pmatrix}
 S&X&T\\
 X^*&V&Z\\
 T^*&Z^*&U
 \end{pmatrix}
 :=[A,F,Y]^*[A,F,Y].
 \label{eq:AFY-Gram}
\end{equation}
Define
\begin{equation}
 \Gaux:=V-X^*S^\dagger X,
 \qquad
 \Caux:=Z-X^*S^\dagger T.
 \label{eq:Gaux-Caux}
\end{equation}
Then
\begin{align}
 \mathbb V^{\mathrm{Duh}}
 &:=U-T^*S^\dagger T\succeq0,
 \label{eq:VDuh-new}\\
 \mathbb A^{\mathrm{old}}
 &:=\Caux^*\Gaux^\dagger\Caux\succeq0,
 \label{eq:Aold-new}\\
 R^{\mathrm{full}}
 &=\mathbb V^{\mathrm{Duh}}-\mathbb A^{\mathrm{old}}
 \nonumber\\
 &=U-T^*S^\dagger T-\Caux^*\Gaux^\dagger\Caux.
 \label{eq:nested-full-residual}
\end{align}
In particular, the five V15 matrices
$(S,T,U,\Gaux,\Caux)$ are deterministic supported Schur transforms of the one
static Gram \eqref{eq:AFY-Gram}; they do not constitute five independently
orientable data sets.
\end{theorem}

\begin{proof}
Let $P_A=A(A^*A)^\dagger A^*$ and set $N=(I-P_A)F$.  Direct calculation gives
\[
 N^*N=F^*(I-P_A)F=V-X^*S^\dagger X=\Gaux
\]
and
\[
 N^*Y=F^*(I-P_A)Y=Z-X^*S^\dagger T=\Caux.
\]
The range decomposition
\[
 \Ran(A,F)=\Ran A\oplus\Ran N
\]
is orthogonal.  Hence
\[
 P_{(A,F)}=P_A+P_N,
 \qquad
 P_N=N(N^*N)^\dagger N^*.
\]
Using \Cref{thm:direct-residual},
\begin{align*}
 R^{\mathrm{full}}
 &=Y^*(I-P_A-P_N)Y\\
 &=\bigl(U-T^*S^\dagger T\bigr)
   -\Caux^*\Gaux^\dagger\Caux,
\end{align*}
which proves the identities and positivity.  Every displayed matrix is a
corner or a supported Schur transform of \eqref{eq:AFY-Gram}.
\end{proof}

\begin{corollary}[Direct verdict versus provenance attribution]
\label{cor:direct-versus-nested}
For the main ancestry verdict, the direct pair $(Q,L,U)$ with
$Q=(A,F)^*(A,F)$ is sufficient.  Separate source-support floors for $S$ and
$\Gaux$ are needed only when one also wants the attribution
\begin{equation}
 \boxed{
 \text{right-clock variance}
 =\text{prior-primitive absorption}
  +\text{new dynamic residual}.}
 \label{eq:attribution-split}
\end{equation}
Thus exact residual transport can be performed with one full primitive support
chart, while the line/clock side programme retains the finer five-matrix split
as an explanatory export.
\end{corollary}

\begin{proposition}[Support-stable perturbation bound for the direct residual]
\label{prop:direct-stability}
For $j=0,1$, let
\[
 \Gamma_j=\begin{pmatrix}Q_j&L_j\\L_j^*&U_j\end{pmatrix}\succeq0
\]
act on fixed coefficient spaces.  Suppose both $Q_j$ have the same support
projection $P$ and
\begin{equation}
 q_*P\preceq Q_j\preceq q^*P,
 \qquad q_*>0.
 \label{eq:Q-floor}
\end{equation}
Put
\begin{equation}
 R_j=U_j-L_j^*Q_j^\dagger L_j,
 \qquad
 \ell_*:=\max(\norm{L_0},\norm{L_1}).
 \label{eq:Rj-direct}
\end{equation}
Then
\begin{equation}
 \boxed{
 \norm{R_1-R_0}
 \le
 \norm{U_1-U_0}
 +\frac{2\ell_*}{q_*}\norm{L_1-L_0}
 +\frac{\ell_*^2}{q_*^2}\norm{Q_1-Q_0}.}
 \label{eq:direct-stability-bound}
\end{equation}
\end{proposition}

\begin{proof}
By \Cref{lem:positive-block-support}, $L_j=PL_j$.  On $P$ the operators $Q_j$
are invertible with inverse norm at most $q_*^{-1}$, and
\[
 \norm{Q_1^\dagger-Q_0^\dagger}
 =\norm{Q_1^{-1}-Q_0^{-1}}
 \le q_*^{-2}\norm{Q_1-Q_0}.
\]
Write
\begin{align*}
 L_1^*Q_1^\dagger L_1-L_0^*Q_0^\dagger L_0
 ={}&(L_1-L_0)^*Q_1^\dagger L_1\\
 &+L_0^*Q_1^\dagger(L_1-L_0)\\
 &+L_0^*(Q_1^\dagger-Q_0^\dagger)L_0.
\end{align*}
Taking norms and using the displayed bounds gives
\eqref{eq:direct-stability-bound} after adding $\norm{U_1-U_0}$.
\end{proof}

\begin{assessment}[First compression gained]
\label{ass:first-compression}
The direct full-primitive residual has only one support transition and one
pseudoinverse.  This does not replace the five-matrix Duhamel attribution; it
prevents the main continuum transport from inheriting two unnecessary
support-chart conditions when only the final ancestry residual is required.
\end{assessment}

% =============================================================================
\section{The physical source metric and the reserve carried by the same card}
\label{sec:source-metric}
% =============================================================================

Let $Q=B^*B$ be the source-source corner of the static card and let
$\Msrc\succeq0$ be the physical source-action form.

\begin{theorem}[Matrix-native null-cost test and reserve spectrum]
\label{thm:metric-reserve-matrix}
The source/action pair is null-cost consistent exactly when
\begin{equation}
 \boxed{\Ker\Msrc\subseteq\Ker Q.}
 \label{eq:null-cost-Q}
\end{equation}
On this branch define
\begin{equation}
 \boxed{
 \Gres_{B,\Msrc}
 :=\Msrc^{\dagger/2}Q\Msrc^{\dagger/2}\succeq0.}
 \label{eq:reserve-matrix}
\end{equation}
Then:
\begin{enumerate}[label=\textup{(R\arabic*)},leftmargin=3.2em]
\item $\Ran(B\Msrc^{\dagger/2})=\Ran B$;
\item the nonzero spectrum of \eqref{eq:reserve-matrix} equals the nonzero
      spectrum of the ambient unit-action covariance
      \begin{equation}
       \Sigma_{B,\Msrc}=B\Msrc^\dagger B^*;
       \label{eq:ambient-covariance}
      \end{equation}
\item it is the generalized positive spectrum of the pencil
      \begin{equation}
       Qv=\gamma\Msrc v
       \quad\text{on }E/\Ker\Msrc;
       \label{eq:generalized-reserve}
      \end{equation}
\item under any invertible source-coordinate change
      \begin{equation}
       \widetilde B=BS,
       \qquad
       \widetilde Q=S^*QS,
       \qquad
       \widetilde{\Msrc}=S^*\Msrc S,
       \label{eq:source-coordinate-change}
      \end{equation}
      the generalized reserve spectrum and the ambient covariance are
      unchanged.
\end{enumerate}
\end{theorem}

\begin{proof}
Since $Q=B^*B$, one has $\Ker Q=\Ker B$, proving
\eqref{eq:null-cost-Q}.  If $u=u_1+u_0$ with
$u_0\in\Ker\Msrc$ and $u_1\perp\Ker\Msrc$, null-cost consistency gives
$Bu=Bu_1$.  Since $\Msrc^{\dagger/2}$ is onto
$(\Ker\Msrc)^\perp$, the ranges of $B$ and $B\Msrc^{\dagger/2}$ agree.

Put $C=B\Msrc^{\dagger/2}$.  Then
$C^*C=\Gres_{B,\Msrc}$ and $CC^*=\Sigma_{B,\Msrc}$, so the two positive
operators have identical nonzero singular-value squares.  Restricting to
$(\Ker\Msrc)^\perp$ and writing $w=\Msrc^{1/2}v$ converts
$Qv=\gamma\Msrc v$ into
\[
 \Msrc^{-1/2}Q\Msrc^{-1/2}w=\gamma w,
\]
which is \eqref{eq:reserve-matrix} on the supported quotient.

For \eqref{eq:source-coordinate-change}, the generalized eigenvalue equation is
carried bijectively by $v=S\widetilde v$.  The ambient covariance is also
invariant because the physical variational problem
\[
 \sup_u\{2\operatorname{Re}\langle x,Bu\rangle-\langle u,\Msrc u\rangle\}
\]
is unchanged after the substitution $u=S\widetilde u$.
\end{proof}

\begin{definition}[Supported source reserve window]
\label{def:reserve-window}
The lower and upper unit-action reserve are
\begin{equation}
 g_-(B,\Msrc)=\lambda_{\min}^{+}(\Gres_{B,\Msrc}),
 \qquad
 g_+(B,\Msrc)=\norm{\Gres_{B,\Msrc}}.
 \label{eq:reserve-window}
\end{equation}
A cofinal family has a \emph{stable reserve window} when
\begin{equation}
 0<g_*\le g_-(B_n,\Msrc_n)
 \le g_+(B_n,\Msrc_n)\le g^*<\infty.
 \label{eq:uniform-reserve-window}
\end{equation}
\end{definition}

\begin{corollary}[Reserve and Duhamel ancestry are distinct outputs]
\label{cor:reserve-ancestry-separate}
On the null-cost branch, the full-primitive Duhamel residual
\begin{equation}
 R^{\mathrm{full}}=U-L^*Q^\dagger L
 \label{eq:residual-independent-M}
\end{equation}
depends only on the occurring source range and the target incidence.  The
reserve window depends on $(Q,\Msrc)$.  Consequently changing the physical
source cost while retaining the same source range can change the reserve
without changing the dynamic-birth rank, and changing the target incidence can
change the dynamic residual while leaving the reserve fixed.
\end{corollary}

\begin{proof}
The first statement is \Cref{thm:direct-residual}; the second is
\Cref{thm:metric-reserve-matrix}.  The two formulas involve different data.
\end{proof}

\begin{theorem}[Target-metric normalization]
\label{thm:target-normalization}
Let $\Ntar=\Ntar^*\succ0$ be the physically transported metric on the target
coefficient space $D$.  For every residual form $R\succeq0$ define
\begin{equation}
 R^{\sharp}=\Ntar^{-1/2}R\Ntar^{-1/2}.
 \label{eq:target-normalized-R}
\end{equation}
If target coordinates are changed by an invertible map $T$ and
\begin{equation}
 \widetilde R=T^*RT,
 \qquad
 \widetilde{\Ntar}=T^*\Ntar T,
 \label{eq:target-coordinate-change}
\end{equation}
then $R^{\sharp}$ and $\widetilde R^{\sharp}$ are unitarily equivalent on their
supports.  In particular rank, generalized eigenvalues, and every physically
normalized floor are coordinate invariant.
\end{theorem}

\begin{proof}
The pairs $(R,\Ntar)$ and $(\widetilde R,\widetilde{\Ntar})$ are congruent.
Their generalized eigenvalue equations are carried bijectively by
$v=T\widetilde v$, exactly as in the proof of
\Cref{thm:metric-reserve-matrix}.  The normalized representatives are positive
realizations of the same generalized form and are therefore unitarily
equivalent on their supported ranges.
\end{proof}

\begin{firewall}[Do not confuse the two multipliers]
\label{fire:two-Ms}
The physical source-action metric $\Msrc$ prices source coefficients.  The
right-energy Duhamel multiplier $\mathfrak m_\lambda(H^R_\lambda)$ generates the
specific target $Y_\lambda$.  A small generalized eigenvalue of
$(Q,\Msrc)$ is reserve loss; a failure of
$\mathfrak m_\lambda(H^R_\lambda)\Ran A_\lambda$ to lie in the full primitive
range is dynamic dispersion.  Neither is renamed as the other.
\end{firewall}

% =============================================================================
\section{The normalized and physical horizon cards from raw delayed Grams}
\label{sec:horizon-card}
% =============================================================================

The main V37 theorem defines the range-normalized horizon operator through the
ambient projection $P_B$.  The next result shows that its complete Tikhonov
card is executable directly from the raw delayed source/source--target blocks
and the static source Gram.  No ambient source basis or freely chosen whitening
is required.

Let
\begin{equation}
 Q=B^*B,
 \qquad
 J_B:=BQ^{\dagger/2}:\supp(Q)E\to\mathcal K.
 \label{eq:JB}
\end{equation}

\begin{lemma}[Canonical polar source isometry]
\label{lem:polar-source-isometry}
The map $J_B$ satisfies
\begin{equation}
 \boxed{J_B^*J_B=P_Q,
 \qquad J_BJ_B^*=P_B.}
 \label{eq:JB-supports}
\end{equation}
Hence it is the source-minimal isometry from the supported static source
coefficient quotient onto the physical source range.
\end{lemma}

\begin{proof}
The first identity is
$Q^{\dagger/2}QQ^{\dagger/2}=P_Q$.  The second is the Moore--Penrose range
formula $BQ^\dagger B^*=P_{\Ran B}$.
\end{proof}

Define, from the raw delayed blocks \eqref{eq:raw-delayed-blocks},
\begin{equation}
 \widehat K(t)=Q^{\dagger/2}K_{BB}(t)Q^{\dagger/2},
 \qquad
 \widehat C(t)=Q^{\dagger/2}K_{BY}(t).
 \label{eq:normalized-delayed-blocks}
\end{equation}
On $L^2((0,T);P_QE)$ set
\begin{align}
 (\widehat{\mathbf G}_Tf)(s)
 &=\int_0^T\widehat K(s+t)f(t)\dd t,
 \label{eq:Ghat-kernel}\\
 (\widehat{\mathbf C}_Td)(s)
 &=\widehat C(s)d.
 \label{eq:Chat-kernel}
\end{align}

\begin{theorem}[Raw-Gram reconstruction of the normalized horizon compiler]
\label{thm:raw-normalized-horizon}
Let
\begin{equation}
 \widehat{\mathscr S}_Tf
 =\int_0^T e^{-tH}J_Bf(t)\dd t.
 \label{eq:normalized-control-synthesis}
\end{equation}
Then
\begin{equation}
 \boxed{
 \widehat{\mathbf G}_T
 =\widehat{\mathscr S}_T^*\widehat{\mathscr S}_T,
 \qquad
 \widehat{\mathbf C}_T
 =\widehat{\mathscr S}_T^*Y.}
 \label{eq:normalized-Hankel-identities}
\end{equation}
Moreover,
\begin{equation}
 \widehat{\mathscr S}_T\widehat{\mathscr S}_T^*
 =\What_{T,B}
 :=\int_0^T e^{-tH}P_Be^{-tH}\dd t.
 \label{eq:normalized-W-reconstructed}
\end{equation}
For every $\delta>0$ and every bounded control synthesis
\(
 F:D\to L^2((0,T);P_QE)
\), define the positive target form
\begin{equation}
 \mathcal J_{T,\delta}(F)
 :=(Y-\widehat{\mathscr S}_TF)^*(Y-\widehat{\mathscr S}_TF)
   +\delta F^*F.
 \label{eq:normalized-control-functional}
\end{equation}
This family has a least element in the Loewner order, and that minimum equals
either of the exactly equivalent matrices
\begin{equation}
 \boxed{
 \begin{aligned}
 \Rhat_{T,\delta}
 &=Y^*Y-\widehat{\mathbf C}_T^*
   (\widehat{\mathbf G}_T+\delta I)^{-1}
   \widehat{\mathbf C}_T\\
 &=\delta Y^*(\What_{T,B}+\delta I)^{-1}Y.
 \end{aligned}}
 \label{eq:normalized-two-representations}
\end{equation}
As $\delta\downarrow0$,
\begin{equation}
 \boxed{
 \Rhat_{T,\delta}\searrow
 Y^*(I-P_H^B)Y=\Rdyn.}
 \label{eq:normalized-zero-limit}
\end{equation}
Thus $(Q,K_{BB},K_{BY},Y^*Y)$ is a complete coefficient-side acquisition of
the range-normalized horizon card.
\end{theorem}

\begin{proof}
By \Cref{lem:polar-source-isometry},
\begin{align*}
 J_B^*e^{-tH}J_B
 &=Q^{\dagger/2}B^*e^{-tH}BQ^{\dagger/2}
 =\widehat K(t),\\
 J_B^*e^{-tH}Y
 &=Q^{\dagger/2}B^*e^{-tH}Y
 =\widehat C(t).
\end{align*}
Taking the adjoint of \eqref{eq:normalized-control-synthesis} gives
$(\widehat{\mathscr S}_T^*x)(s)=J_B^*e^{-sH}x$.  Composition yields the
integral kernel $\widehat K(s+t)$ and the cross loading $\widehat C(s)$,
proving \eqref{eq:normalized-Hankel-identities}.  The other composition is
\[
 \widehat{\mathscr S}_T\widehat{\mathscr S}_T^*
 =\int_0^T e^{-tH}J_BJ_B^*e^{-tH}\dd t
 =\What_{T,B}.
\]

Put
\[
 F_\delta
 =(\widehat{\mathbf G}_T+\delta I)^{-1}\widehat{\mathbf C}_T.
\]
Using \eqref{eq:normalized-Hankel-identities}, direct expansion gives the exact
operator-valued completion of the square
\[
 \mathcal J_{T,\delta}(F)
 =\Rhat_{T,\delta}
 +(F-F_\delta)^*(\widehat{\mathbf G}_T+\delta I)(F-F_\delta).
\]
Thus the first line of \eqref{eq:normalized-two-representations} is the unique
Loewner minimum.  The push-through identity
\[
 I-S(S^*S+\delta I)^{-1}S^*
 =\delta(SS^*+\delta I)^{-1}
\]
with $S=\widehat{\mathscr S}_T$ gives the second.  Spectral calculus shows that
$\delta(\What+\delta I)^{-1}$ decreases to the kernel projection of $\What$.
The inherited positive-horizon support theorem identifies that kernel with
$I-P_H^B$, proving \eqref{eq:normalized-zero-limit}.
\end{proof}

\begin{theorem}[The physical-metric horizon card is reconstructed from the same raw blocks]
\label{thm:raw-physical-horizon}
Assume $\Ker\Msrc\subseteq\Ker Q$ and define
\begin{equation}
 K^{\Msrc}(t)
 =\Msrc^{\dagger/2}K_{BB}(t)\Msrc^{\dagger/2},
 \qquad
 C^{\Msrc}(t)
 =\Msrc^{\dagger/2}K_{BY}(t).
 \label{eq:physical-delayed-blocks}
\end{equation}
Let $\mathbf G_T^{\Msrc}$ and $\mathbf C_T^{\Msrc}$ be the corresponding
Hankel and cross operators, and define
\[
 \mathscr S_T^{\Msrc}f
 :=\int_0^T e^{-tH}B\Msrc^{\dagger/2}f(t)\dd t
 \quad\text{on }L^2((0,T);P_{\Msrc}E),
 \qquad P_{\Msrc}:=\Msrc\Msrc^\dagger.
\]
For every bounded $F:D\to L^2((0,T);P_{\Msrc}E)$, set
\[
 \mathcal J_{T,\varepsilon}^{\Msrc}(F)
 :=(Y-\mathscr S_T^{\Msrc}F)^*(Y-\mathscr S_T^{\Msrc}F)
   +\varepsilon F^*F.
\]
Then the least element of this family in the Loewner order is
\begin{equation}
 \boxed{
 \begin{aligned}
 \Rphys_{T,\varepsilon}
 &=Y^*Y-(\mathbf C_T^{\Msrc})^*
   (\mathbf G_T^{\Msrc}+\varepsilon I)^{-1}
   \mathbf C_T^{\Msrc}\\
 &=\varepsilon Y^*(\Wphys_{T,B}+\varepsilon I)^{-1}Y,
 \end{aligned}}
 \label{eq:physical-two-representations}
\end{equation}
where
\begin{equation}
 \Wphys_{T,B}
 =\int_0^T e^{-tH}B\Msrc^\dagger B^*e^{-tH}\dd t.
 \label{eq:physical-W}
\end{equation}
Hence the physical reserve and both horizon compilers are deterministic
functions of one same-card packet
\begin{equation}
 \boxed{(\Msrc,Q,K_{BB}(\cdot),K_{BY}(\cdot),Y^*Y).}
 \label{eq:one-horizon-packet}
\end{equation}
\end{theorem}

\begin{proof}
Put $B_{\Msrc}=B\Msrc^{\dagger/2}$.  Then
\[
 (B_{\Msrc})^*e^{-tH}B_{\Msrc}=K^{\Msrc}(t),
 \qquad
 (B_{\Msrc})^*e^{-tH}Y=C^{\Msrc}(t).
\]
Apply the proof of \Cref{thm:raw-normalized-horizon} to $B_{\Msrc}$.  In
unit-action coordinates, the $L^2$ norm is exactly the original integrated
source action, and the ambient Gramian is \eqref{eq:physical-W}.
\end{proof}

\begin{corollary}[Reserve--shape comparison on the paired raw card]
\label{cor:paired-reserve-shape}
If $g_-$ and $g_+$ are the supported reserve bounds of
\Cref{def:reserve-window}, then
\begin{equation}
 g_-\What_{T,B}\preceq\Wphys_{T,B}\preceq g_+\What_{T,B}
 \label{eq:paired-W-window}
\end{equation}
and, for every $\varepsilon>0$,
\begin{equation}
 \boxed{
 \Rhat_{T,\varepsilon/g_+}
 \preceq\Rphys_{T,\varepsilon}
 \preceq\Rhat_{T,\varepsilon/g_-}.}
 \label{eq:paired-R-window}
\end{equation}
All three matrices have the same zero-regularization limit $\Rdyn$.
\end{corollary}

\begin{proof}
The polar reserve factorization gives
$g_-P_B\preceq B\Msrc^\dagger B^*\preceq g_+P_B$.  Conjugate by the semigroup
and integrate to obtain \eqref{eq:paired-W-window}.  The scalar function
$x\mapsto\varepsilon/(x+\varepsilon)$ is operator monotone decreasing; applying
it to the two inequalities and compressing by $Y$ gives
\eqref{eq:paired-R-window}.  The common kernel is the orthogonal complement of
the $H$-cyclic source space.
\end{proof}

\begin{assessment}[Second compression gained]
\label{ass:second-compression}
The actual horizon acquisition need not print an ambient polar source frame.
The static source Gram normalizes the raw delayed blocks canonically, while the
physical source-action matrix generates the unit-action version.  The two
compilers are therefore tied to one raw same-history experiment and cannot be
aligned independently after the residual is known.
\end{assessment}

% =============================================================================
\section{The exact Duhamel--horizon two-tail squeeze}
\label{sec:two-tail}
% =============================================================================

This is the main new bridge.  Fix one compatible paired card and let
$\Lambda_0$ be a finite set of left energies.  Define the finite Duhamel head
\begin{equation}
 L_{\Lambda_0}
 :=\sum_{\lambda\in\Lambda_0}R_\lambda^{\mathrm{Duh}}\succeq0
 \label{eq:Duhamel-head}
\end{equation}
and the omitted-left-energy tail
\begin{equation}
 \ThetaD_{\Lambda_0}
 :=\sum_{\lambda\notin\Lambda_0}R_\lambda^{\mathrm{Duh}}\succeq0.
 \label{eq:Duhamel-tail}
\end{equation}
The inherited norm-convergent fibre theorem gives
\begin{equation}
 \Rdyn=L_{\Lambda_0}+\ThetaD_{\Lambda_0}.
 \label{eq:Rdyn-head-tail}
\end{equation}
For a physical-metric horizon card put
\begin{equation}
 \PhiLev_{T,\varepsilon;\Msrc}
 :=\Rphys_{T,\varepsilon}-\Rdyn\succeq0,
 \label{eq:physical-leverage-tail}
\end{equation}
and for the range-normalized card put
\begin{equation}
 \widehat\PhiLev_{T,\delta}
 :=\Rhat_{T,\delta}-\Rdyn\succeq0.
 \label{eq:normalized-leverage-tail}
\end{equation}

\begin{theorem}[Exact two-tail gap identity]
\label{thm:two-tail-gap}
For every finite left-energy head, every $T>0$, and every positive
regularization parameter,
\begin{equation}
 \boxed{
 \Rphys_{T,\varepsilon}-L_{\Lambda_0}
 =\ThetaD_{\Lambda_0}
  +\PhiLev_{T,\varepsilon;\Msrc}.}
 \label{eq:physical-two-tail-gap}
\end{equation}
Likewise,
\begin{equation}
 \boxed{
 \Rhat_{T,\delta}-L_{\Lambda_0}
 =\ThetaD_{\Lambda_0}
  +\widehat\PhiLev_{T,\delta}.}
 \label{eq:normalized-two-tail-gap}
\end{equation}
Every term on the right-hand sides is positive.  Consequently
\begin{equation}
 \boxed{
 L_{\Lambda_0}
 \preceq\Rdyn
 \preceq\Rphys_{T,\varepsilon},
 \qquad
 L_{\Lambda_0}
 \preceq\Rdyn
 \preceq\Rhat_{T,\delta}.}
 \label{eq:ancestry-sandwich}
\end{equation}
The lower error is exactly the unresolved left-energy tail; the upper error is
exactly the unresolved positive-leverage tail.  No third hidden error is
present.
\end{theorem}

\begin{proof}
Subtract \eqref{eq:Rdyn-head-tail} from
\eqref{eq:physical-leverage-tail}:
\begin{align*}
 \Rphys_{T,\varepsilon}-L_{\Lambda_0}
 &=\bigl(\Rphys_{T,\varepsilon}-\Rdyn\bigr)
   +\bigl(\Rdyn-L_{\Lambda_0}\bigr)\\
 &=\PhiLev_{T,\varepsilon;\Msrc}+\ThetaD_{\Lambda_0}.
\end{align*}
The normalized identity is identical.  Positivity of the Duhamel tail is the
inherited fibre theorem, and positivity of the leverage tail is the spectral
Tikhonov decomposition.  The order sandwich follows.
\end{proof}

\begin{corollary}[One cross-gap controls both unresolved tails]
\label{cor:one-gap-controls-two}
Let
\begin{equation}
 \mathcal G_{\Lambda_0,T,\varepsilon}^{\Msrc}
 :=\Rphys_{T,\varepsilon}-L_{\Lambda_0}\succeq0.
 \label{eq:cross-gap}
\end{equation}
Then
\begin{equation}
 0\preceq\ThetaD_{\Lambda_0}\preceq\mathcal G_{\Lambda_0,T,\varepsilon}^{\Msrc},
 \qquad
 0\preceq\PhiLev_{T,\varepsilon;\Msrc}
 \preceq\mathcal G_{\Lambda_0,T,\varepsilon}^{\Msrc}.
 \label{eq:tail-dominated-by-gap}
\end{equation}
In particular, if a same-card calculation certifies
\begin{equation}
 \norm{\mathcal G_{\Lambda_0,T,\varepsilon}^{\Msrc}}\le\eta,
 \label{eq:gap-eta}
\end{equation}
then
\begin{equation}
 \boxed{
 \norm{\Rdyn-L_{\Lambda_0}}\le\eta,
 \qquad
 \norm{\Rphys_{T,\varepsilon}-\Rdyn}\le\eta.}
 \label{eq:squeeze-eta}
\end{equation}
Thus a direct bound on one observable cross-gap replaces two separately
allocated tail estimates.
\end{corollary}

\begin{proof}
Both tails are positive summands of the gap in
\eqref{eq:physical-two-tail-gap}; hence each is bounded above by their sum in
the Loewner order.  The operator norm is monotone on positive operators.  The
two differences in \eqref{eq:squeeze-eta} are precisely the two tails.
\end{proof}

\begin{corollary}[Robust birth and follower certificates]
\label{cor:birth-follower-certificates}
Let $P$ be a physically transported target projection.
\begin{enumerate}[label=\textup{(C\arabic*)},leftmargin=3.2em]
\item If
      \begin{equation}
       PL_{\Lambda_0}P\succeq cP,
       \qquad c>0,
       \label{eq:positive-head-certificate}
      \end{equation}
      then $P\Rdyn P\succeq cP$.  This is a dynamic-birth certificate which
      requires neither the omitted left-energy tail nor the small-leverage
      tail.
\item If
      \begin{equation}
       \norm{\Rphys_{T,\varepsilon}}\le\eta,
       \label{eq:upper-follower-certificate}
      \end{equation}
      then $\norm{\Rdyn}\le\eta$.  A cofinal family with such upper bounds
      tending to zero is a complete follower certificate.
\item If \eqref{eq:gap-eta} holds, then every normalized target vector $d$
      obeys the sharp interval
      \begin{equation}
       \ip{d}{L_{\Lambda_0}d}
       \le\ip{d}{\Rdyn d}
       \le\ip{d}{L_{\Lambda_0}d}+\eta\norm{d}^2.
       \label{eq:scalar-squeeze}
      \end{equation}
\item The rank inequality
      \begin{equation}
       \rank\Rdyn\ge\rank L_{\Lambda_0}
       \label{eq:rank-lower}
      \end{equation}
      holds for every finite head.
\end{enumerate}
\end{corollary}

\begin{proof}
Items (C1)--(C3) follow directly from
\eqref{eq:ancestry-sandwich} and \Cref{cor:one-gap-controls-two}.  For (C4),
$0\preceq L_{\Lambda_0}\preceq\Rdyn$ implies
$\Ker\Rdyn\subseteq\Ker L_{\Lambda_0}$, hence the rank inequality.
\end{proof}

\begin{assessment}[Why the two compilers are complementary]
\label{ass:lower-upper-complementarity}
The finite Duhamel card is naturally a lower compiler: every computed positive
fibre block is already part of the exact residual.  The physical-horizon
Tikhonov card is naturally an upper compiler: its only excess is positive
small-leverage mass.  Asking either compiler to perform the other's role caused
the earlier programme to circle.  The exact gap identity gives them a single
meeting interface.
\end{assessment}

% =============================================================================
\section{Cofinal squeeze and the exact borderline survivor}
\label{sec:cofinal-squeeze}
% =============================================================================

Residual matrices at different cutoffs may be compared only after the physical
target coefficient spaces and their metrics have been transported.  In this
section all forms have already been pulled to one fixed finite target screen
and normalized as in \Cref{thm:target-normalization}.

For cutoff $n$, choose a finite left-energy head $\Lambda_n$, a physical horizon
$T_*>0$, and a regularization $\varepsilon_n>0$.  Write
\begin{equation}
 L_n=L_{\Lambda_n,n}^{\sharp},
 \qquad
 R_n=R_{\mathrm{dyn},n}^{\sharp},
 \qquad
 U_n=(\Rphys_{n,T_*,\varepsilon_n})^{\sharp}.
 \label{eq:cofinal-LRU}
\end{equation}
Then
\begin{equation}
 0\preceq L_n\preceq R_n\preceq U_n.
 \label{eq:cofinal-order}
\end{equation}

\begin{theorem}[Cofinal paired-card squeeze]
\label{thm:cofinal-squeeze}
Assume:
\begin{enumerate}[label=\textup{(S\arabic*)},leftmargin=3.2em]
\item the source, source-action metric, target metric, completed clock, left
      spectral head, and static/delayed Grams are transported on one physical
      route;
\item $(L_n)$ is Cauchy in operator norm;
\item the paired cross-gap closes:
      \begin{equation}
       \norm{U_n-L_n}\longrightarrow0.
       \label{eq:cross-gap-closes}
      \end{equation}
\end{enumerate}
Then there is a unique positive form $R_\infty$ on the target screen such that
\begin{equation}
 \boxed{
 L_n\longrightarrow R_\infty,
 \qquad
 R_n\longrightarrow R_\infty,
 \qquad
 U_n\longrightarrow R_\infty.}
 \label{eq:triple-convergence}
\end{equation}
Moreover the omitted-left-energy and positive-leverage tails both converge to
zero in norm, without requiring a separate estimate for either one.
\end{theorem}

\begin{proof}
By assumption $L_n\to R_\infty$ for some positive $R_\infty$.  Since
$0\preceq R_n-L_n\preceq U_n-L_n$, positivity gives
\[
 \norm{R_n-L_n}\le\norm{U_n-L_n}\to0.
\]
Similarly $0\preceq U_n-R_n\preceq U_n-L_n$, so
$\norm{U_n-R_n}\to0$.  This proves \eqref{eq:triple-convergence}.  The two
positive tails are exactly $R_n-L_n$ and $U_n-R_n$.
\end{proof}

\begin{corollary}[Cofinal follower or persistent dynamic source]
\label{cor:cofinal-two-branches}
Under the hypotheses of \Cref{thm:cofinal-squeeze}:
\begin{enumerate}[label=\textup{(B\arabic*)},leftmargin=3.2em]
\item if $R_\infty=0$, the target has no cofinal dynamic source on the selected
      screen;
\item if $R_\infty\ne0$, then with $P_\infty=\supp R_\infty$,
      \begin{equation}
       R_\infty\succeq
       \lambda_{\min}^{+}(R_\infty)P_\infty,
       \label{eq:limit-positive-floor}
      \end{equation}
      and the dynamic source has persistent rank $\rank R_\infty$ on that
      screen.
\end{enumerate}
\end{corollary}

\begin{proof}
The first statement is the definition of followerhood on the transported target
screen.  In finite dimension every nonzero positive operator has a positive
least supported eigenvalue, giving the second statement.
\end{proof}

\begin{theorem}[Typed survivor when the paired gap does not close]
\label{thm:gap-survivor}
Suppose the cards are null-cost consistent and physically aligned, and assume
for some $c>0$ that
\begin{equation}
 \limsup_{n\to\infty}\norm{U_n-L_n}\ge c.
 \label{eq:gap-not-close}
\end{equation}
Then there is a subsequence and unit target vectors $d_n$ such that at least one
of the following two alternatives holds along a further subsequence:
\begin{align}
 \ip{d_n}{(\ThetaD_{\Lambda_n,n})^{\sharp}d_n}
 &\ge c/4,
 \label{eq:Duhamel-tail-survivor}\\
 \ip{d_n}{(\PhiLev_{n,T_*,\varepsilon_n;\Msrc_n})^{\sharp}d_n}
 &\ge c/4.
 \label{eq:leverage-tail-survivor}
\end{align}
Thus failure of the squeeze returns either a positive left-energy survivor or a
positive compact-horizon leverage survivor; it does not justify another
untyped abstraction.
\end{theorem}

\begin{proof}
Choose a subsequence for which
$\norm{U_n-L_n}\ge c/2$.  Since $U_n-L_n\succeq0$, choose a unit vector $d_n$
with
\[
 \ip{d_n}{(U_n-L_n)d_n}\ge c/2.
\]
The exact two-tail identity writes this scalar as the sum of two nonnegative
numbers.  At least one is at least $c/4$.  One of the two alternatives occurs
infinitely often; pass to that subsequence.
\end{proof}

\begin{theorem}[Reserve--geometry--ancestry alternative]
\label{thm:cofinal-master-alternative}
For a cofinal family of proposed paired cards on a fixed target screen, execute
the rows in the following order.
\begin{enumerate}[label=\textup{(A\arabic*)},leftmargin=3.2em]
\item If $\Ker\Msrc_n\nsubseteq\Ker Q_n$, return a zero-action nonzero-source
      vector.
\item If null-cost consistency holds but the supported reserve window
      \eqref{eq:uniform-reserve-window} fails, return source-reserve collapse or
      blow-up before interpreting normalized clock geometry.
\item If the reserve window holds and the paired cross-gap closes, return the
      unique $R_\infty$ of \Cref{thm:cofinal-squeeze}, classified by
      \Cref{cor:cofinal-two-branches}.
\item If the reserve window holds and the paired cross-gap does not close,
      return one of the normalized survivors in
      \Cref{thm:gap-survivor}.
\end{enumerate}
These outcomes are exhaustive after physical occurrence, support stability, and
route alignment have been checked.
\end{theorem}

\begin{proof}
The null-cost and reserve branches are disjoint by definition.  On the stable
reserve branch, the physical and range-normalized Tikhonov tails are equivalent
by \Cref{cor:paired-reserve-shape}.  The cross-gap either converges to zero or
has a positive limsup.  The first case is \Cref{thm:cofinal-squeeze}; the second
is \Cref{thm:gap-survivor}.
\end{proof}

\begin{firewall}[A reserve collapse is not a birth]
\label{fire:reserve-not-birth}
The squeeze theorem remains an ancestry statement, while the reserve window is
a physical-cost statement.  A card may have a stable source range and a fixed
zero dynamic residual while its unit-action reserve collapses.  Such a family
fails the continuum source-control handoff even though it does not generate a
new dynamic source.
\end{firewall}

% =============================================================================
\section{Why every row of the paired card is necessary}
\label{sec:irredundancy}
% =============================================================================

The following finite examples show that the source metric, mixed static block,
and delayed clock block answer different questions.

\begin{theorem}[Irredundancy of metric, mixed incidence, and clock]
\label{thm:paired-card-irredundancy}
No one of the following rows is determined by the other two.
\begin{enumerate}[label=\textup{(I\arabic*)},leftmargin=3.2em]
\item \emph{Physical source metric.}  Let $\mathcal K=E=D=\C$, $B=Y=1$, and
      $H=0$.  The same static and delayed Grams are compatible with
      $\Msrc_1=1$ and $\Msrc_2=4$.  The dynamic residual is zero in both cases,
      but the unit-action reserve is respectively $1$ and $1/4$.
\item \emph{Mixed primitive--target incidence.}  Positive blocks
      \begin{equation}
       \Gamma_c=\begin{pmatrix}1&c\\\overline c&1\end{pmatrix},
       \qquad |c|\le1,
       \label{eq:mixed-countercards}
      \end{equation}
      have the same source and target marginals and residual
      $R_c=1-|c|^2$.  Thus separate marginals permit both a complete follower
      $(|c|=1)$ and a full birth $(c=0)$.
\item \emph{Completed clock.}  Let $\mathcal K=\C^2$, $B1=e_1$, $Y1=e_2$,
      and $\Msrc=1$.  For $H_0=0$, the source-cyclic space is
      $\Span\{e_1\}$ and $\Rdyn=1$.  For
      \begin{equation}
       H_1=P_v,
       \qquad v=2^{-1/2}(e_1+e_2),
       \label{eq:clock-countercard}
      \end{equation}
      one has $H_1e_1=\tfrac12(e_1+e_2)$, so the source-cyclic space is all of
      $\C^2$ and $\Rdyn=0$.  The static source/target Gram is identical in the
      two cases.
\end{enumerate}
\end{theorem}

\begin{proof}
For (I1), $Q=1$ and
$\Gres=\Msrc^{-1}Q$, giving the two reserve values.  Since $Y\in\Ran B$, the
residual vanishes independently of the metric.

For (I2), the supported Schur complement of the source corner is
$1-\overline c c$.  The matrices are positive exactly for $|c|\le1$.

For (I3), the $H_0$-cyclic source space is immediate.  For $H_1$, the vectors
$e_1$ and $H_1e_1$ are linearly independent, so polynomial and semigroup cyclic
closures equal $\C^2$.  The zero atom is therefore respectively the projection
onto $\Span\{e_2\}$ and zero.
\end{proof}

\begin{assessment}[The minimum common acquisition]
\label{ass:minimum-common-acquisition}
The source metric is needed for physical reserve; the mixed static block is
needed for the finite Duhamel lower card; and the delayed kernel is needed for
the compact-horizon upper card.  These rows may be acquired in one experiment,
but none may be inferred from the other two.
\end{assessment}

% =============================================================================
\section{Canonical ancestry export after the squeeze lands}
\label{sec:ancestry-export}
% =============================================================================

A positive residual should not be represented by an arbitrary new basis.  Its
polar range is the unique source-minimal export.

\begin{theorem}[Canonical minimum dynamic-source bank]
\label{thm:canonical-dynamic-bank}
Let $P=P_H^B$ and
\begin{equation}
 Z=(I-P)Y,
 \qquad
 \Rdyn=Z^*Z.
 \label{eq:Z-Rdyn}
\end{equation}
Define
\begin{equation}
 J_{\rm dyn}:=Z\Rdyn^{\dagger/2}.
 \label{eq:Jdyn}
\end{equation}
Then
\begin{equation}
 \boxed{
 J_{\rm dyn}^*J_{\rm dyn}=P_{\Rdyn},
 \qquad
 J_{\rm dyn}J_{\rm dyn}^*=P_{\Ran Z},
 \qquad
 Z=J_{\rm dyn}\Rdyn^{1/2}.}
 \label{eq:Jdyn-polar}
\end{equation}
The minimum number of additional Hilbert directions required by the dynamic
source is $\rank\Rdyn$.  If $\Rdyn=0$, no dynamic source is exported.
\end{theorem}

\begin{proof}
This is the polar decomposition of the finite-rank map $Z$.  Since
$Z^*Z=\Rdyn$, the Moore--Penrose polar factor is
$Z\Rdyn^{\dagger/2}$.  Its initial and final projections are the support of
$Z^*Z$ and the projection onto $\Ran Z$.  Every realization of $Z$ has range
dimension at least $\rank Z=\rank\Rdyn$, while the displayed factor uses
exactly that dimension.
\end{proof}

\begin{definition}[Cofinal ancestry packet]
\label{def:ancestry-packet}
Once the paired card and route have landed, the packet exported to later ESM
compilers is
\begin{equation}
 \boxed{
 \mathfrak A_{\rm ESM}
 =\bigl(
 \Sigma_{B,\Msrc},\Gres_{B,\Msrc},P_B,
 \What_{T,B},\Rdyn,J_{\rm dyn}
 \bigr),}
 \label{eq:ancestry-packet}
\end{equation}
together with its physical target metric, support projections, reserve window,
left-energy provenance split, and adjacent/direct transport certificates.
\end{definition}

\begin{firewall}[An ancestry export is not yet current, stress, or Einstein]
\label{fire:ancestry-not-Einstein}
The map $J_{\rm dyn}$ says which response direction is not generated by the
prior source-cyclic bank.  It does not establish that this direction occurs as
a local electromagnetic current, lies in the common current graph domain,
has the required contact subtraction, descends to the determinant/Pfaffian
line, defines a represented stress tensor, or has a cutoff-uniform covariant
Einstein residual.  Those claims require their own same-history incidence,
locality, domain, line, Ward, and transport rows.
\end{firewall}

% =============================================================================
\section{A terminating executable card for the next physical calculation}
\label{sec:compiler}
% =============================================================================

\begin{acquisition}[Paired reserve--Duhamel--horizon card]
\label{card:paired-acquisition}
At cutoff $n$, the minimum payload is:
\begin{enumerate}[label=\textup{(E\arabic*)},leftmargin=3.2em]
\item one named target-independent unquenched witness and its physical target
      metric;
\item the routed full primitive source synthesis, represented by its static
      source Gram $Q_n$, and the physical source-action matrix $\Msrc_n$ on the
      same coefficient list;
\item a finite left-energy head, the role-resolved positive Grams
      $\Gamma_{A_\lambda F_\lambda Y_\lambda}$, and the literal Duhamel
      incidence residual
      \begin{equation}
       \Delta_\lambda
       :=Y_\lambda-\mathfrak m_\lambda(H^R_\lambda)A_\lambda;
       \label{eq:Duhamel-incidence-defect}
      \end{equation}
\item the raw delayed blocks $(K_{BB,n},K_{BY,n})$ on $[0,2T_*]$;
\item supported floors for the full static source block and the physical
      reserve, plus either exact kernels or validated interval enclosures;
\item one adjacent-cutoff route and one independently assembled direct/staged
      route for $B_n,\Msrc_n,H_n,Y_n$ and the target metric.
\end{enumerate}
No source whitening, role average, target-shaped selector, or independently
fitted mixed block is part of the payload.
\end{acquisition}

\begin{theorem}[Sound terminating paired-card compiler]
\label{thm:terminating-compiler}
On the payload of \Cref{card:paired-acquisition}, execute the following order.
\begin{enumerate}[label=\textup{(C\arabic*)},leftmargin=3.2em]
\item Test common-history positivity and the exact Duhamel incidence
      $\Delta_\lambda=0$ on every retained fibre.
\item Test $\Ker\Msrc_n\subseteq\Ker Q_n$.  On failure return a zero-action
      nonzero-source vector.
\item Compute the generalized reserve spectrum and retain source rank,
      $g_{-,n}$, and $g_{+,n}$ separately.
\item Compute each finite-head residual by the direct full-primitive Schur
      complement.  Compute the nested five-matrix split only when variance
      versus prior-primitive absorption is required.
\item Normalize the raw delayed blocks by $Q_n$ and by $\Msrc_n$ to construct
      the range-normalized and physical horizon cards.
\item Form the lower card $L_n$, upper card $U_n$, and positive cross-gap
      $U_n-L_n$.
\item Transport the primitive source, metric, target, and clock first; recompute
      every reserve, projector, horizon Gramian, and residual afterward.
\end{enumerate}
The compiler terminates with exactly one of the following statuses:
\begin{equation}
 \boxed{
 \begin{gathered}
 \mathsf{PASS\!:\ DYNAMIC\ BIRTH},\qquad
 \mathsf{PASS\!:\ COMPLETE\ FOLLOWER},\\
 \mathsf{PASS\!:\ COFINAL\ SQUEEZE},\\
 \mathsf{TYPED\ METRIC\ OBSTRUCTION},\qquad
 \mathsf{TYPED\ RESERVE\ COLLAPSE},\\
 \mathsf{TYPED\ DUHAMEL\ INCIDENCE\ FAILURE},\\
 \mathsf{TYPED\ SUPPORT/SPECTRAL\ TRANSITION},\qquad
 \mathsf{TYPED\ CLOCK/ROUTE\ FAILURE},\\
 \mathsf{STOP\!:\ UNSQUEEZED\ TWO\!\!\!-TAIL\ GAP},\\
 \mathsf{STOP\!:\ WAITING\ FOR\ PHYSICAL\ CARD}.
 \end{gathered}}
 \label{eq:compiler-statuses}
\end{equation}
Every pass and obstruction is accompanied by its positive matrix, support,
vector, or route witness.
\end{theorem}

\begin{proof}
Step (C1) is an occurrence and typing test.  Step (C2) is exact by
\Cref{thm:metric-reserve-matrix}.  Step (C3) uses the generalized reserve
pencil.  Step (C4) is sound by
\Cref{thm:direct-residual,thm:nested-attribution}.  Step (C5) is sound by
\Cref{thm:raw-normalized-horizon,thm:raw-physical-horizon}.  Step (C6) is the
exact identity of \Cref{thm:two-tail-gap}; the birth, follower, and squeeze
passes are certified by
\Cref{cor:birth-follower-certificates,thm:cofinal-squeeze}.  If the gap does not
close, \Cref{thm:gap-survivor} returns one of the two typed tail survivors.
Step (C7) is required because the derived cards are nonlinear functions of the
primitive source, metric, and clock.

Every operation is finite-dimensional linear algebra on a finite head or a
bounded integral-operator calculation on a fixed compact horizon.  A failed
premise is returned before its inverse or transport is used.  Hence the
compiler terminates and the listed statuses are exhaustive.
\end{proof}

\begin{firewall}[Independent route is a held-out test]
\label{fire:held-out-route}
The direct/staged route is not fitted to make the cross-gap small.  Source,
metric, clock, left-energy projections, and mixed Grams are transported from
the primitive card.  Agreement of already shorted residuals without agreement
of those primitive rows is not a cutoff certificate.
\end{firewall}

% =============================================================================
\section{Version V1 theorem, proof-status ledger, and the Version V2 attack}
\label{sec:V1-conclusion}
% =============================================================================

\begin{theorem}[Integrated Version V1 result]
\label{thm:V1-integrated}
Preserving the imported line/clock, source-metric, horizon, and terminal-order
theorems at their stated hypotheses, Version~V1 proves:
\begin{enumerate}[label=\textup{(V1.\arabic*)},leftmargin=4.6em]
\item one full primitive/target Gram computes the exact static ancestry
      residual by one supported Schur complement;
\item the V15 five-matrix formula is the orthogonal nested attribution of that
      same complement;
\item direct residual transport needs only one full primitive support floor;
\item null-cost consistency and the supported reserve spectrum are executable
      from the matrix pair $(\Msrc,Q)$ and are invariant under physical source
      reparametrization;
\item the range-normalized and physical-metric horizon cards are reconstructed
      from the same raw delayed source/source--target blocks;
\item a finite Duhamel head is a positive lower bound and every physical or
      normalized Tikhonov card is a positive upper bound for $\Rdyn$;
\item the upper-minus-lower cross-gap is exactly the sum of the omitted
      left-energy tail and the positive-leverage tail;
\item one small cross-gap controls both tails, while a positive Duhamel head is
      already a rigorous birth witness;
\item cofinal closure of the cross-gap squeezes the exact residual to a unique
      transported limit;
\item failure of the squeeze returns a normalized left-energy or leverage-tail
      survivor;
\item the polar residual bank is the unique minimum new dynamic-source export;
\item the resulting compiler has the terminating outcomes
      \eqref{eq:compiler-statuses}.
\end{enumerate}
No item evaluates the missing physical ESM card or promotes the ancestry packet
to current, stress, determinant-line orientation, interacting quantum
continuum, Lorentzian reconstruction, Einstein backreaction, or pure-colour
closure.
\end{theorem}

\begin{proof}
Items (V1.1)--(V1.3) are
\Cref{thm:direct-residual,thm:nested-attribution,prop:direct-stability}.
Items (V1.4)--(V1.5) are
\Cref{thm:metric-reserve-matrix,thm:raw-normalized-horizon,thm:raw-physical-horizon}.
Items (V1.6)--(V1.8) are
\Cref{thm:two-tail-gap,cor:one-gap-controls-two,cor:birth-follower-certificates}.
Items (V1.9)--(V1.10) are
\Cref{thm:cofinal-squeeze,thm:gap-survivor}.  Item (V1.11) is
\Cref{thm:canonical-dynamic-bank}, and the final item is
\Cref{thm:terminating-compiler}.  The scope statement is
\Cref{fire:ancestry-not-Einstein}.
\end{proof}

\begingroup\small
\begin{longtable}{>{\raggedright\arraybackslash}p{0.28\textwidth}
                  >{\raggedright\arraybackslash}p{0.20\textwidth}
                  >{\raggedright\arraybackslash}p{0.44\textwidth}}
\caption{Version V1 proof-status ledger.}
\label{tab:V1-ledger}\\
\toprule
\textbf{Row}&\textbf{Status after V1}&\textbf{Exact conclusion}\\
\midrule
\endfirsthead
\toprule
\textbf{Row}&\textbf{Status after V1}&\textbf{Exact conclusion}\\
\midrule
\endhead
Line/reflection/domain/\newline $H_Q$/OS--KMS compiler
&\textbf{IMPORTED / PRESERVED}
&One target-independent witness carries the finite line and physical clocks.\\[0.35em]
Fibre Duhamel five-matrix theorem
&\textbf{IMPORTED / PRESERVED}
&Every left fibre has right-reduction, prior-primitive absorption, or positive unabsorbed dispersion.\\[0.35em]
Physical source metric and reserve
&\textbf{IMPORTED, THEN PAIRED}
&The physical metric precedes leverage; V1 ties its reserve to the same static and delayed card used by the Duhamel compiler.\\[0.35em]
Direct full-primitive Schur card
&\textbf{PROVED}
&One positive $(B,Y)$ Gram gives the final residual with one support and one pseudoinverse.\\[0.35em]
Raw delayed-Gram horizon compiler
&\textbf{PROVED}
&Static $Q$ canonically normalizes the raw delayed blocks; $\Msrc$ gives the physical unit-action version.\\[0.35em]
Exact Duhamel--horizon two-tail gap
&\textbf{PROVED}
&Upper minus lower equals omitted-left-energy tail plus positive-leverage tail, with both summands positive.\\[0.35em]
Cofinal paired-card squeeze
&\textbf{PROVED}
&Closing one cross-gap transports the exact residual and kills both unresolved tails.\\[0.35em]
Borderline failure object
&\textbf{PROVED}
&A nonclosing gap yields a normalized left-energy or leverage-tail survivor.\\[0.35em]
Actual routed ESM source-action matrix
&\textbf{OPEN PHYSICAL VALUE}
&The attached corpus does not print $\Msrc$ on one selected unquenched card.\\[0.35em]
Actual static primitive/target Grams
&\textbf{OPEN PHYSICAL VALUE}
&The cardwise $\Gamma_{A_\lambda F_\lambda Y_\lambda}$ fields and support floors are not supplied.\\[0.35em]
Actual row-summed weak delayed kernel
&\textbf{OPEN PHYSICAL VALUE}
&The complete $(K_{BB},K_{BY})$ block on one physical horizon is not supplied.\\[0.35em]
Current, stress, determinant/Pfaffian, Einstein
&\textbf{BLOCKED DOWNSTREAM}
&These gates remain closed until the ancestry packet lands and their independent incidence/domain rows are populated.\\
\bottomrule
\end{longtable}
\endgroup

\begin{assessment}[Version V1 physical verdict]
\label{ass:V1-verdict}
Version~V1 is a mathematical \passstatus{} for the combined paired-card
compiler and a physical \stopstatus{} at the first named acquisition.  The
source files contain the structural formulas and the exact finite reductions,
but not the actual common-history matrices needed to choose between follower,
prior-primitive absorption, robust dynamic birth, reserve collapse, or a
positive two-tail survivor on one selected ESM thread.  No numerical branch is
fabricated.
\end{assessment}

\begin{openproblem}[Version V2: execute one literal paired reserve--Duhamel--horizon card]
\label{op:V2}
Preserve Version~V1 verbatim and do not add another abstract ancestry compiler.
On one named target-independent unquenched ESM thread:
\begin{enumerate}[label=\textup{(V2.\arabic*)},leftmargin=5.2em]
\item Freeze the electromagnetic profile bank, deterministic Higgs/electroweak
      router, Ward quotient, full primitive coefficient list, target metric,
      and physical source-action matrix before inspecting any leverage or
      residual.
\item Acquire the static positive Gram of
      $(A_\lambda,F_\lambda,Y_\lambda)$ on the first predeclared finite
      left-energy head and verify the calibrated Duhamel relation.
\item Compute the direct full-primitive head residual and, on every positive
      fibre, its V15 variance/absorption attribution.
\item Acquire $(K_{BB},K_{BY})$ on one compact physical horizon and construct
      both the physical and polar-normalized upper cards from the raw blocks.
\item Compute the paired cross-gap.  Return immediately on a positive lower
      head, a vanishing transported upper card, or a rigorously small cross-gap.
\item If the gap remains positive, separate its left-energy and leverage
      summands and return the first normalized survivor.
\item Transport the primitive source, metric, target, clock, static Gram, and
      delayed kernel through two adjacent cutoffs and one independent
      direct/staged route.
\item Export the cofinal ancestry packet only after the reserve window and
      route checks pass.  Reopen current/stress/Einstein rows only through their
      own physical incidence and domain cards.
\end{enumerate}
The accepted endpoint is
\begin{equation}
 \boxed{
 \begin{gathered}
 \text{one literal paired source-metric/static-Duhamel/horizon card,}\\
 \text{one reserve window and one lower--exact--upper ancestry sandwich,}\\
 \text{one cofinal follower, robust birth, or typed two-tail survivor,}\\
 \text{and one primitive adjacent/direct route certificate.}
 \end{gathered}}
 \label{eq:V2-target}
\end{equation}
Absent those matrices, the exact terminal status is
\begin{center}
\small\ttfamily
STOP\_WAITING\_FOR\_PHYSICAL\_SOURCE\_METRIC\_STATIC\_DUHAMEL\_AND\_HORIZON\_CARD.
\end{center}
\end{openproblem}

\begin{firewall}[Claims not made in Version V1]
\label{fire:V1-claims}
Version~V1 does not evaluate the actual source reserve, normalized leverage
law, finite fibre matrices, positive tails, or dynamic residual on a selected
ESM regulator.  It does not orient a determinant or Pfaffian line, prove a
fermion measure phase, construct a local current graph domain, establish a
Drude or vacuum Kubo coefficient, close an interacting Higgs/gauge/fermion
continuum, reconstruct Lorentzian fields, prove a Yang--Mills mass gap, or
derive the Einstein equations from the finite ancestry card.  Its new theorem
is the exact meeting interface and squeeze between the already established
Duhamel lower compiler and physical-horizon upper compiler.
\end{firewall}

% =============================================================================
\section*{Version V1 append-only record}
\addcontentsline{toc}{section}{Version V1 append-only record}
% =============================================================================

\begin{description}[leftmargin=3.8cm,style=nextline]
\item[Version V1 -- 2 September 2026.]
Started the consolidated Physical Source Metric, Duhamel Ancestry and Quantum
Continuum Programme.  Imported, without re-proving, the completed
sector-patched line/clock Duhamel theorem, the main source-metric/reserve and
finite-horizon theorem, and the Grand-Tensor terminal promotion order.

Compressed the full primitive ancestry residual to one supported Schur
complement of the common static primitive/target Gram.  Proved that the V15
five-matrix variance/absorption formula is the orthogonal nested factorization
of the same complement, and gave a one-support perturbation bound for direct
transport.

Proved that the physical source-action matrix and static source Gram determine
the coordinate-invariant reserve spectrum, while the static Gram and raw
delayed source/source--target blocks reconstruct the polar-normalized horizon
compiler.  The same raw blocks, combined with the physical source metric,
reconstruct the unit-action physical horizon compiler.

Established the exact Duhamel--horizon two-tail identity: finite Duhamel head
below, exact dynamic residual in the middle, physical or normalized Tikhonov
card above.  The upper-minus-lower gap is exactly the sum of the omitted
left-energy tail and the positive-leverage tail.  Proved robust birth, follower,
and cofinal squeeze certificates, and showed that a nonclosing gap returns a
normalized survivor in one of those two tails.

Constructed the source-minimal polar dynamic bank, stated the exact downstream
firewall, and reduced Version~V2 to literal population and transport of one
paired physical card.  The attached corpus does not contain those actual
matrices, so no numerical follower or birth verdict is asserted.
\end{description}

\begin{thebibliography}{9}

\bibitem{GrandTensorV062}
A.~P\'elissier,
\emph{The Renewal Grand Tensor: Relative Primitive Minimality, Operational
Spectralization, and Domain-Specific SM--Spacetime and Arithmetic Loadings},
Version V062, 2026.

\bibitem{ESMMainV37}
A.~P\'elissier,
\emph{Renewal Dedicated Einstein--Standard--Model Quantum-State,
Dynamic-Source and Continuum Closure Programme}, Version V37, 2026.

\bibitem{ESMLineClockV15}
A.~P\'elissier,
\emph{Renewal Einstein--Standard--Model Sector-Patched Fermion Line and
Physical Clock Programme}, Version V15, 2026.

\end{thebibliography}

\begin{firewall}[Append-only instruction after Version V1]
\label{fire:append-after-V1}
Any Version~V2 must preserve the complete present source verbatim and append
new material immediately before the sole final \verb|\end{document}|.  Its
filename is
\begin{center}
\path{Renewal_Einstein_Standard_Model_Physical_Source_Metric_Duhamel_Ancestry_and_Quantum_Continuum_Programme_V2.tex}.
\end{center}
It may not replace the physical source-action metric by a raw coefficient norm,
fit the primitive/target mixed block, whiten before same-history assembly,
count right-clock variance before prior-primitive absorption, infer a complete
follower from a finite zero Duhamel head, infer a birth from reserve collapse,
transport already shorted residuals instead of primitive rows, or promote the
ancestry packet to current, stress, determinant-line, continuum, Lorentzian,
Einstein, or pure-colour closure without the corresponding physical cards.
\end{firewall}


% =============================================================================
% BEGIN APPENDED VERSION V2 MATERIAL
% =============================================================================
\clearpage
\providecommand{\Dom}{\operatorname{Dom}}
\providecommand{\HS}{\mathrm{HS}}
\providecommand{\hD}{\mathsf h}
\providecommand{\bjet}{\mathcal J_{\partial,\lambda}}

\section*{Version V2: native clock incidence, hard-wall leakage, and an evaluated control card}
\addcontentsline{toc}{section}{Version V2: clock audit, hard-wall obstruction, and evaluated control}

\begin{abstract}
The paired-card identities of V1 are retained.  An exact audit of the literal
finite-rate RP resampling clock first proves an infinite-detail heat-trace
obstruction and a finite-lag contradiction on an already retained thermal
weight source.  A source-specific increasing-rate approximation bound separates
that exact obstruction from controlled approximation.  The main thermal
calculation then concerns a
specific feature of the inherited source formulas: primitive columns have a
fixed right thermal factor, whereas the derivative of the Gibbs square root
contains a Duhamel endpoint.  On a finite left-energy fibre, every primitive
form column belongs to every right-energy power domain.  A nonzero Dirichlet
boundary compatibility defect in the unsmeared action row, however, makes the
Duhamel target fail the third right-energy domain.  It therefore cannot be
absorbed by any finite bank of those primitive columns.  For a smooth fixed-wall
Schr\"odinger realization, the defect is the explicit normal jet
$(\partial_n\psi)^*(\partial_nV)$.  Its rank is a lower bound for the fibre
ancestry rank.  This is a sufficient, source-specific obstruction, not a
converse or a cutoff-uniform floor.

The mechanism is evaluated on a fully specified scalar hard-wall diagnostic:
$\hD(u)=-\partial_x^2+u(x-\pi/2)$ on $(0,\pi)$ with Dirichlet boundary,
$\beta=1$, and entrance delay $1/4$.  The source cost is fixed independently by
the normalized Lebesgue profile metric.  The complete sine-basis insertion
table, five Duhamel matrices on the first left fibre, full source reserve,
raw delayed kernels, and positive residual tails are calculated.  An exact
two-right-mode annihilator certifies a target-normalized residual greater than
$10^{-5}$.  Integer outward-interval arithmetic gives a complete residual
interval and a source-reserve window.  Both persist under the explicitly
specified spectral truncations.

The diagnostic is not identified with the full RPESM witness or with a physical
spacetime cutoff.  On the actual ESM branch the new task is to evaluate the
boundary/domain incidence of its already declared action derivatives, or to
rule out this mechanism and return to the five-matrix calculation.  No
boundary condition, source metric, electromagnetic command, or continuum
transport missing from the attached card is silently supplied by the control.
\end{abstract}

\section{V2 direction, source scope, and one necessary V1 qualification}
\label{sec:V2-direction}

\paragraph{Imported definitions.}
The line/clock programme \cite{ESMLineClockV15}, equations
\textup{(6.4)--(6.13)} and its V15 fibre theorem, fixes the following order:
one self-adjoint quantum Hamiltonian $\hD$, its Gibbs state, primitive
weight/action/line roles, and then the derivative of the Gibbs square root.
The main programme \cite{ESMMainV37} fixes a physical coefficient metric before
reserve or leverage is interpreted.  Neither construction is repeated here.
In particular, the clock used below is left multiplication by the quantum
Hamiltonian, not the auxiliary global-refresh Markov generator.

\paragraph{What is and is not supplied by the sources.}
The attached main programme calls its selected branch hard-wall and
sector-superselected, and imports compact resolvent and trace-class heat
semigroups.  Those statements alone do not specify a smooth Dirichlet
realization, the boundary of an overlap sector, or the conormal derivative of
an electromagnetic action writer.  The boundary theorem below explicitly
states those extra analytic hypotheses.  The scalar interval calculation is a
new diagnostic realization of the inherited \emph{source formulas}; it is not
claimed to be the first-rational full-coupling witness named in line/clock V14.
This distinction prevents a solvable replacement model from being mistaken
for the missing ESM physical card.

\begin{remark}[Cauchy premise in the V1 cofinal summary]
\label{rem:V2-Cauchy-correction}
In the pass branch of \Cref{thm:cofinal-master-alternative}, the Cauchy premise
of \Cref{thm:cofinal-squeeze} must also be retained explicitly.  A vanishing
cross-gap alone does not imply convergence: the positive scalar triples
$L_n=R_n=U_n=1+(-1)^n/4$ have zero gap and no limit.  The valid cofinal statement
is therefore the theorem with both Cauchy lower heads and vanishing gap, not
an implication from the gap alone.  No previous source text is altered.
\end{remark}

\section{Literal resampling: an exact clock audit before thermal substitution}
\label{sec:V2R-native-clock}

This independent calculation uses the actual projective Markov skeleton
\textup{(RP.13)} in \cite{GrandTensorV062}.  The distinction between path,
source-cyclic, modular, and ambient thermal clocks is already part of
line/clock V4 and is not a new theorem here.  What is evaluated here is the
literal detail spectrum and a finite-lag obstruction on an \emph{existing}
primitive source, before this particular path clock is used as a substitute
for the quantum clock.

Let $\nu(dx,dy)=\mu(dx)K_x(dy)$ and let
$J:L^2(\mu)\to L^2(\nu)$ be $Jf(x,y)=f(x)$.  Then
$J^*g(x)=\int g(x,y)K_x(dy)$.  Set
$P=JJ^*$, $Q=I-P$, and $\mathcal D=\Ran Q$.  Let $P_0$ be a reversible
Markov contraction on $L^2(\mu)$ and retain the literal fine skeleton
\[
 P_1((x,y),d(x',y'))=P_0(x,dx')K_{x'}(dy').
\]
For one calibrated Poisson rate $\omega>0$, put
$A_i=\omega(I-P_i)$.  This Poissonization is not the logarithm of the skeleton.

\begin{theorem}[Exact native detail spectrum]
\label{thm:V2R-native-spectrum}
On $L^2(\nu)=JL^2(\mu)\oplus\mathcal D$ one has
\begin{align}
 P_1&=JP_0J^*,&P_1Q&=0,\\
 A_1&=JA_0J^*+\omega Q,& A_1|_{\mathcal D}&=\omega I_{\mathcal D},
 \label{eq:V2R-generator}\\
 e^{-tA_1}&=Je^{-tA_0}J^*+e^{-\omega t}Q,
 \label{eq:V2R-semigroup}\\
 (A_1+zI)^{-1}&=J(A_0+zI)^{-1}J^*+(\omega+z)^{-1}Q
 \quad(z>0).
 \label{eq:V2R-resolvent}
\end{align}
The full energy-$\omega$ projection is
$JE_{A_0}(\{\omega\})J^*+Q$; its two summands must be assembled before any
source short.  If $\mathcal D\ne0$, $P_1$ is not $e^{-hK}$ for any
nonnegative self-adjoint $K$ and finite $h>0$.
\end{theorem}

\begin{proof}
Integrating the displayed kernel against $g(x',y')$ first performs $J^*$,
then $P_0$, and then $J$.  Thus $P_1=JP_0J^*$ and it annihilates $\mathcal D$.
Since $J^*J=I$, $A_1$ is the orthogonal direct sum of $A_0$ and
$\omega I_{\mathcal D}$.  Functional calculus proves every formula and the
resonant projection statement.  Finally $e^{-hK}$ has zero kernel: for any
nonzero $v$, its squared image norm is the integral of the strictly positive
function $e^{-2hE}$ against the nonzero spectral measure of $v$.  The nonzero
kernel of $P_1$ contradicts this property.
\end{proof}

\begin{lemma}[A continuous conditional detail is infinite dimensional]
\label{lem:V2R-infinite-detail}
Suppose the conditional detail contains a real coordinate with a jointly
measurable conditional distribution function $F_x$ which is continuous for
$\mu$-almost every $x$.  Then $\mathcal D$ is infinite dimensional.
\end{lemma}

\begin{proof}
The conditional probability-integral transform
$U=F_x(y)$ is uniform on $[0,1]$ given $x$.  For $j\ge1$, the bounded
vectors $d_j=e^{2\pi i jU}$ have conditional mean zero and satisfy
$\langle d_j,d_k\rangle=\delta_{jk}$ by direct integration on $[0,1]$.
Thus they form an infinite orthonormal family in $\mathcal D$.  In particular,
a conditional Higgs coordinate with a positive Lebesgue density meets this
hypothesis; no inference from the number of lattice sites to Hilbert-space
dimension is needed.
\end{proof}

\begin{theorem}[Fixed-cutoff thermal obstruction for the literal finite-rate clock]
\label{thm:V2R-heat-obstruction}
Let the literal RP generator at one cutoff be
\[
 L=\rho(R-I)+\sum_{j=1}^{N}r_j(E_j-I),
 \qquad \rho,r_j\ge0,\quad N<\infty,
\]
where $R$ and $E_j$ are the actual stationary and conditional-expectation
projections.  Then, with $C=\rho+\sum_jr_j<\infty$,
\[
 0\preceq A:=-L\preceq CI.
\]
If its $L^2$ carrier is infinite dimensional, then for every $\beta>0$,
\begin{equation}
 \boxed{e^{-\beta A}\succeq e^{-\beta C}I,
 \qquad \operatorname{Tr}e^{-\beta A}=\infty.}
 \label{eq:V2R-infinite-heat-trace}
\end{equation}
Its resolvent is not compact.  More precisely, on the infinite-detail branch
of \Cref{thm:V2R-native-spectrum}, every compact operator $K_t$ satisfies
\begin{equation}
 \|e^{-tA_1}-K_t\|\ge e^{-\omega t}
 \quad(t>0).
 \label{eq:V2R-essential-heat}
\end{equation}
\end{theorem}

\begin{proof}
Each $I-E_j$ and $I-R$ is an orthogonal projection, hence lies between zero and
$I$.  Summing gives the finite upper bound on $A$.  Spectral calculus gives
\eqref{eq:V2R-infinite-heat-trace}; its trace against the first $m$ vectors of
any orthonormal family is at least $m e^{-\beta C}$ and diverges.
Similarly $(A+I)^{-1}\succeq(C+1)^{-1}I$, which excludes compactness on an
infinite-dimensional space.  For the sharper statement take the orthonormal
detail family of \Cref{lem:V2R-infinite-detail}.  It converges weakly to zero,
so $K_td_j\to0$ in norm.  Since $e^{-tA_1}d_j=e^{-\omega t}d_j$, evaluating
the operator difference on $d_j$ gives \eqref{eq:V2R-essential-heat}.
\end{proof}

\begin{corollary}[Actual native detail reserve and horizon coefficient]
\label{cor:V2R-detail-values}
For finite path syntheses $B:E\to\mathcal D$, $Y:D\to\mathcal D$, put
$Q_B=B^*B$.  On the \emph{explicitly Dirichlet-priced} native subcard
$M_{\rm Dir}=B^*A_1B=\omega Q_B$, one has
\begin{equation}
 \boxed{BM_{\rm Dir}^{\dagger}B^*=\omega^{-1}P_B,
 \qquad
 \widehat W_{T,B}=c_{\omega,T}P_B,
 \quad c_{\omega,T}=\frac{1-e^{-2\omega T}}{2\omega}.}
 \label{eq:V2R-detail-values}
\end{equation}
For $B\ne0$, the supported reserve is exactly $1/\omega$; the physical horizon is
$(c_{\omega,T}/\omega)P_B$, and
\[
 R_{T,\varepsilon}^{\rm Dir}
 =Y^*(I-P_B)Y+
   \frac{\varepsilon}{\varepsilon+c_{\omega,T}/\omega}Y^*P_BY.
\]
\end{corollary}

\begin{proof}
The source lies entirely at the scalar energy $\omega$, so its cyclic range is
$\Ran B$ and its range projection commutes with the semigroup.  Integrating
$e^{-2\omega t}$ gives $c_{\omega,T}$.  The pseudoinverse identity
$(\omega Q_B)^\dagger=\omega^{-1}Q_B^\dagger$ gives the covariance; inversion
on $\Ran B$ and its orthogonal complement gives the final expression.
\end{proof}

\begin{firewall}[Precisely what this native calculation does not identify]
\label{fire:V2R-native-scope}
A perfectly normalizable stationary probability $\mu$ is compatible with the
infinite \emph{operator} heat trace above.  No failure of the stationary RP
path is inferred.  Nor is $M_{\rm Dir}$ identified with the physical routed
electromagnetic source-action matrix.  The native detail target $Y$ has not
been identified with the Gibbs Duhamel half-source.  A separately constructed
quantum Hamiltonian can remain perfectly valid.  The rejected step is using
this bounded finite-rate generator as that quantum Hamiltonian without an
additional, source-valid clock-incidence theorem.
\end{firewall}

\section{A finite-lag separator in an already retained thermal source}
\label{sec:V2R-source-separator}

The preceding heat-trace obstruction is ambient.  It alone would not exclude
agreement on a small source-cyclic subspace.  The following additional
calculation uses precisely the sector-weight primitive that line/clock V15
already retains, rather than adding a target-shaped probe.

Let $\mathsf h\ge0$ be an infinite-dimensional compact-resolvent thermal
Hamiltonian with $Z(s)=\operatorname{Tr}e^{-s\mathsf h}<\infty$ for all
$s>0$.  For fixed $\beta,\tau>0$, its retained weight column is
\[
 b=Z(\beta)^{-1/2}e^{-(\tau+\beta/2)\mathsf h}\in\HS(\mathcal H),
 \qquad H_LX=\mathsf hX.
\]
No Duhamel target is used in the next theorem.

\begin{theorem}[Exact weight-source energy law]
\label{thm:V2R-weight-law}
The left-energy measure and delayed source Gram of $b$ are
\begin{equation}
 \mu_b(dE)=\frac{e^{-(\beta+2\tau)E}}{Z(\beta)}
                \operatorname{Tr}E_{\mathsf h}(dE),
 \qquad
 k(t)=\langle b,e^{-tH_L}b\rangle
     =\frac{Z(\beta+2\tau+t)}{Z(\beta)}.
 \label{eq:V2R-weight-law}
\end{equation}
Every finite energy cap has strictly positive omitted source mass, with
\[
 0<\mu_b((C,\infty))\le e^{-2\tau C}\quad(C<\infty),
 \qquad
 \int E^m\mu_b(dE)<\infty\quad(m\ge0).
\]
\end{theorem}

\begin{proof}
In the spectral basis the squared weight in an energy-$E$ eigenspace is its
multiplicity times $Z(\beta)^{-1}e^{-(\beta+2\tau)E}$.  This proves
\eqref{eq:V2R-weight-law}.  Compact resolvent on an infinite-dimensional space
implies unbounded eigenvalues of finite multiplicity, and every displayed
weight is positive, proving the strict tail statement.  Bounding
$e^{-2\tau E}$ by $e^{-2\tau C}$ in the tail and summing the remaining Gibbs
weights gives the upper bound.  The function $E^m e^{-2\tau E}$ is bounded,
which gives every moment.
\end{proof}

For a scalar kernel $k$, an integer $m\ge1$, a lag $h>0$, and a proposed
energy cap $C\ge0$, define the finite-lag expression
\begin{equation}
 D_{m,h,C}(k)=
 (1-e^{-hC})^m k(0)-\sum_{j=0}^{m}(-1)^j\binom mj k(jh).
 \label{eq:V2R-lag-separator}
\end{equation}

\begin{theorem}[Finite-horizon witness against bounded-clock incidence]
\label{thm:V2R-finite-lag}
If $k(t)=\langle s,e^{-tA}s\rangle$ with $0\preceq A\preceq CI$, then
$D_{m,h,C}(k)\ge0$ for every $m,h$.  For the retained thermal weight kernel
\eqref{eq:V2R-weight-law}, every finite $C$ and every $T_*>0$ admit a finite
$m$ and $h>0$ with $mh\le T_*$ and
\[
 \boxed{D_{m,h,C}(k)<0.}
\]
For $C>0$ an explicit choice is as follows.  Retain one eigenvalue $E>C$
with weight $w=\mu_b(\{E\})>0$, and choose $m$ such that
\[
 \Delta:=m\log(E/C)-\log(k(0)/w)>0.
\]
Every $0<h\le\min\{T_*/m,\Delta/(2mE)\}$ gives
\begin{equation}
 D_{m,h,C}(k)\le-\gamma,
 \quad
 \gamma=w(1-e^{-hE})^m-k(0)(1-e^{-hC})^m>0.
 \label{eq:V2R-margin}
\end{equation}
\end{theorem}

\begin{proof}
The binomial formula writes the finite sum in
\eqref{eq:V2R-lag-separator} as
$\langle s,(I-e^{-hA})^ms\rangle$.  On $[0,C]$ the spectral multiplier is
at most $(1-e^{-hC})^m$, giving nonnegativity.  For the thermal kernel,
positivity of its spectral measure gives the lower bound
$w(1-e^{-hE})^m$ for that same sum.  The scalar inequalities
$1-e^{-hE}\ge hE e^{-hE}$ and $1-e^{-hC}\le hC$ imply
\[
 \log\frac{w(1-e^{-hE})^m}{k(0)(1-e^{-hC})^m}
 \ge\Delta-mhE\ge\Delta/2>0.
\]
This proves the stated margin.  If $C=0$, take $m=1$ and any $0<h\le T_*$:
$D_{1,h,0}(k)=k(h)-k(0)<0$ because there is occupied positive energy.
\end{proof}

\begin{corollary}[A source-level obstruction, not merely an ambient one]
\label{cor:V2R-source-obstruction}
No bounded finite-rate RP clock can reproduce the full delayed V15 primitive
kernel on a bank retaining the weight column of an infinite-dimensional
thermal sector.  In particular there is no source-fixing semigroup unitary on
that complete primitive bank.  An invertible source-coordinate change cannot
remove this witness: the matrix form of \eqref{eq:V2R-lag-separator}
transforms by congruence, and the retained vector $b$ remains in the bank.
\end{corollary}

\begin{proof}
The two putatively identical kernels would give opposite signs for the same
finite expression of \Cref{thm:V2R-finite-lag}.  Congruence preserves positivity
and preserves the existence of a negative direction, proving the last claim.
\end{proof}

\begin{proposition}[Error budget for the finite-lag witness]
\label{prop:V2R-lag-errors}
Suppose $|\widetilde k(jh)-k(jh)|\le\eta_j$.  With
$a=(1-e^{-hC})^m$, the resulting separator error is at most
\[
 (1-a)\eta_0+\sum_{j=1}^m\binom mj\eta_j.
\]
For a common radius $\eta$, this is $(2^m-a)\eta$.  An outward negative
upper endpoint proves the obstruction; failure to resolve its sign proves
neither incidence nor a positive result.
\end{proposition}

\begin{proof}
The coefficient of $k(0)$ is $a-1$ and every other coefficient has absolute
value $\binom mj$.  Apply the triangle inequality and sum the binomial
coefficients.  This also shows why the existence of a finite witness does not
supply a uniformly cheap numerical acquisition.
\end{proof}

\section{Increasing-rate approximation is not excluded by the exact obstruction}
\label{sec:V2R-approximation}

The exact obstruction must not be turned into an approximation no-go.  Keep
the thermal state and source $b$ fixed and define the bounded positive
Poissonized heat-step clock
\[
 A_h=\frac{I-e^{-hH_L}}h,\qquad 0\preceq A_h\preceq h^{-1}I.
\]
This uses the \emph{actual thermal heat step}, not the conditionally
resampling skeleton; it is a source-kernel comparison and not a new Gibbs
state.

\begin{theorem}[Explicit source error and necessary rate divergence]
\label{thm:V2R-Poisson-approximation}
For the weight source above,
\begin{equation}
 \boxed{
 \sup_{t\ge0}\|(e^{-tA_h}-e^{-tH_L})b\|
       \le\frac{h}{e^2\tau},\qquad
 0\le\langle b,(e^{-tA_h}-e^{-tH_L})b\rangle
       \le\frac{h}{2e^2\tau}.}
 \label{eq:V2R-Poisson-error}
\end{equation}
Conversely, if $k_n(t)=\langle s_n,e^{-tA_n}s_n\rangle$ converges uniformly
to $k(t)$ on some $[0,T_*]$, and $0\preceq A_n\preceq C_nI$, then
$C_n\to\infty$.  Uniformly bounded rates cannot approximate this complete
source kernel arbitrarily accurately even on a fixed nontrivial horizon.
\end{theorem}

\begin{proof}
For $E>0$, let $q=(1-e^{-hE})/h\le E$.  Integrating in the logarithmic
spectral variable yields, uniformly in $t\ge0$,
\[
 0\le e^{-tq}-e^{-tE}
 =\int_{\log q}^{\log E} t e^r e^{-te^r}\,dr
 \le e^{-1}\log(E/q)\le hE/e.
\]
The last inequality uses $1-e^{-hE}\ge hE e^{-hE}$.  The estimate is zero
at $E=0$ by continuity.  Moreover
\[
 \|H_Lb\|^2=\int E^2\mu_b(dE)\le(e\tau)^{-2},\qquad
 \int E\mu_b(dE)\le(2e\tau)^{-1},
\]
by taking the suprema of $E^2e^{-2\tau E}$ and $Ee^{-2\tau E}$ against the
normalized Gibbs law.  Spectral integration proves
\eqref{eq:V2R-Poisson-error}.  If $C_n$ did not tend to infinity, a subsequence
would have one finite upper bound $C$.  On that subsequence the finite
separator of \Cref{thm:V2R-finite-lag} is nonnegative.  Convergence at its
finitely many lags contradicts the strictly negative limit.
\end{proof}

\begin{assessment}[Consequence for the remaining thermal calculation]
\label{ass:V2R-thermal-route}
The native path-clock spectrum, detail reserve, and finite-lag contradiction
are now explicit.  They do not revoke any separately populated quantum
Hamiltonian or a comparison which never asserted this forbidden equality.
The hard-wall calculation below uses that thermal source formula on its own
specified Hamiltonian, not the finite-rate RP generator.  Kernel approximation
with increasing rates remains possible, but an ancestry-limit claim still
needs V1's physical reserve, same-source transport, and positive-tail controls;
small kernel error alone does not make a cyclic projection continuous.
\end{assessment}

\section{The right-energy regularity carried by the literal source formulas}
\label{sec:V2-right-regularity}

Let $\hD\ge0$ be self-adjoint with compact resolvent on $\mathcal H$ and
$\operatorname{Tr}e^{-s\hD}<\infty$ for every $s>0$.  Fix
$b=\beta/2>0$, $\tau>0$, and $Z=\operatorname{Tr}e^{-2b\hD}$.
Write $E_\lambda$ for a finite-rank left spectral projection.  On
$\mathcal K=\HS(\mathcal H)$, let $H_LX=\hD X$ and $H_RX=X\hD$.
The left fibre is $\mathcal K_\lambda=E_\lambda\mathcal K$.

A finite primitive bank has columns of the inherited form
\begin{equation}
 B_j=Z^{-1/2}e^{-\tau\hD}W_je^{-b\hD},
 \qquad
 \widehat W_j=(1+\hD)^{-1/2}W_j(1+\hD)^{-1/2}\in\mathcal B(\mathcal H).
 \label{eq:V2-primitive}
\end{equation}
The equality for a form writer is understood by the bounded-factor
realization, not by multiplying unbounded operators outside their domains.
Scalar sector weights and the identity column are included; a finite
line-current bank is included whenever its columns satisfy the same form
hypothesis.  Additional previously occurring columns are not omitted merely
because they are inconvenient for the short.

\begin{lemma}[All right power domains on a finite left fibre]
\label{lem:V2-right-smoothing}
For every $k\in\mathbb N_0$ and every primitive column,
\begin{equation}
 E_\lambda B_j\in\Dom(H_R^k),
 \qquad
 \|H_R^kE_\lambda B_j\|_{\HS}
 \le \frac{e^{-\tau\lambda}\sqrt{(1+\lambda)\rank E_\lambda}}{\sqrt Z}
 \|\widehat W_j\|
 \sup_{x\ge0}(1+x)^{1/2}x^ke^{-bx}.
 \label{eq:V2-right-smoothing}
\end{equation}
Consequently every vector in the finite primitive fibre range belongs to
$\bigcap_{k\ge0}\Dom(H_R^k)$.
\end{lemma}

\begin{proof}
On spectral cores the right-regularized row equals
\[
 Z^{-1/2}e^{-\tau\lambda}(1+\lambda)^{1/2}
 E_\lambda\widehat W_j(1+\hD)^{1/2}\hD^ke^{-b\hD}.
\]
The last factor is bounded by spectral calculus, and $E_\lambda$ has
Hilbert--Schmidt norm $\sqrt{\rank E_\lambda}$.  The ideal inequality gives
\eqref{eq:V2-right-smoothing}.  Spectral approximation, followed by closedness
of right multiplication by $\hD^k$, identifies this bounded-factor expression
with $H_R^kE_\lambda B_j$.  A finite linear combination remains in each linear
domain.  No assertion is made that an infinite Hilbert closure of such vectors
retains this regularity.
\end{proof}

For a bounded common-action deformation $V(u)$, put
$V^\circ(u)=V(u)-\operatorname{Tr}(\rho V(u))I$, where
$\rho=Z^{-1}e^{-2b\hD}$.  The inherited internal half-source is
\begin{equation}
 Y(u)=Z^{-1/2}\int_0^1
       -b e^{-b(1-s)\hD}V^\circ(u)e^{-bs\hD}\,ds.
 \label{eq:V2-half-source}
\end{equation}
For a general form deformation the same statements below apply whenever this
Duhamel row and the indicated unsmeared finite-left row are defined.
Introduce
\begin{equation}
 d_\lambda(x)=
 \begin{cases}
 \displaystyle\frac{e^{-bx}-e^{-b\lambda}}{x-\lambda},&x\ne\lambda,\\[0.4em]
 -be^{-b\lambda},&x=\lambda.
 \end{cases}
 \label{eq:V2-divided-difference}
\end{equation}
The diagonal value is removable and is never divided by zero.  If
$G_\lambda(u)=E_\lambda V^\circ(u)$ is a Hilbert--Schmidt row, then
\begin{equation}
 Y_\lambda(u)=Z^{-1/2}d_\lambda(H_R)G_\lambda(u),
 \qquad
 (H_R-\lambda)\sqrt ZY_\lambda(u)
   =(e^{-bH_R}-e^{-b\lambda})G_\lambda(u).
 \label{eq:V2-row-resolvent}
\end{equation}
These are the literal Duhamel coefficients; the source metric does not enter
them.

\begin{theorem}[Duhamel endpoint preserves a right-domain defect]
\label{thm:V2-domain-defect}
Let $k\ge0$ and suppose
\[
 G_\lambda(u)\in\Dom(H_R^k)\setminus\Dom(H_R^{k+1}).
\]
Then
\begin{equation}
 Y_\lambda(u)\in\Dom(H_R^{k+1})
                  \setminus\Dom(H_R^{k+2}).
 \label{eq:V2-one-power}
\end{equation}
If the complete prior primitive fibre range is finite and satisfies
\eqref{eq:V2-primitive}, then
\begin{equation}
 \langle u,R_\lambda^{\rm Duh}u\rangle
   =\|(I-P_{B_\lambda})Y_\lambda(u)\|^2>0.
 \label{eq:V2-domain-birth}
\end{equation}
The conclusion includes absorption by all the primitive columns, not just the
action column.
\end{theorem}

\begin{proof}
The function $(1+x)|d_\lambda(x)|$ is bounded on $[0,\infty)$: it is continuous
near $\lambda$, and tends to $e^{-b\lambda}$ at infinity.  Hence the first
membership in \eqref{eq:V2-one-power} follows from spectral calculus.
If $Y_\lambda(u)$ also belonged to $\Dom(H_R^{k+2})$, the left side of the
second identity in \eqref{eq:V2-row-resolvent} would belong to
$\Dom(H_R^{k+1})$.  The vector $e^{-bH_R}G_\lambda(u)$ belongs to every power
domain.  Subtracting it from that identity, and dividing by the nonzero scalar
$e^{-b\lambda}$, would put $G_\lambda(u)$ in $\Dom(H_R^{k+1})$, a
contradiction.  By \Cref{lem:V2-right-smoothing}, $Y_\lambda(u)$ is outside the
finite primitive range.  That range is closed, so its Hilbert distance is
strictly positive.  The imported V15 full-fibre formula gives
\eqref{eq:V2-domain-birth}.
\end{proof}

\begin{remark}[Why the left clock cannot remove this defect]
On the $\lambda$ fibre the left semigroup is the scalar $e^{-t\lambda}$.
It does not change the right-energy regularity of a primitive row.  The
argument is therefore about the actual left-clock ancestry short, not about
failure to include enough positive left times.  Allowing an independently
occurring right-clock module would be a different primitive declaration.
\end{remark}

\section{A concrete boundary-jet test on fixed Dirichlet sector charts}
\label{sec:V2-boundary-jet}

We now impose explicit hypotheses, rather than deducing a boundary condition
from the word ``hard-wall''.  Let $\Omega\subset\mathbb R^d$ be bounded with
smooth boundary, let $\kappa>0$, and let
\begin{equation}
 \hD=-\kappa\Delta+W(x)\ge0
 \quad\hbox{on }L^2(\Omega;\mathbb C^r),\qquad
 \Dom\hD=H^2(\Omega;\mathbb C^r)\cap H_0^1(\Omega;\mathbb C^r),
 \label{eq:V2-Dirichlet-chart}
\end{equation}
where $W$ is smooth and Hermitian up to the boundary.  Assume the selected
eigenvectors are smooth up to the boundary.  These standard smooth-elliptic
domain hypotheses are stated explicitly; only their integer-power recursion
is used below.\footnote{For the broader distinction between ordinary
regularity and domains of powers of elliptic boundary realizations, see
G.~Grubb, \emph{Regularity of spectral fractional Dirichlet and Neumann
problems}, Math. Nachr. \textbf{289} (2016), 831--844,
\href{https://doi.org/10.1002/mana.201500041}{doi:10.1002/mana.201500041}.
The interval calculation below is proved directly and uses no fractional-domain
result.}
The deformation $V(u)$ is smooth matrix multiplication, Hermitian for real
coefficients, and the wall is fixed under this deformation.  Domain motion,
singular coefficients, and non-Dirichlet walls are outside this theorem.

Choose an orthonormal basis $\psi_1,\ldots,\psi_s$ of $E_\lambda\mathcal H$.
The normal derivative is the ordinary outward derivative in
\eqref{eq:V2-Dirichlet-chart}.  Define the linear boundary-row map
\begin{equation}
 \bjet u
   =\bigl((\partial_n\psi_a)^*(\partial_nV(u))\bigr)_{a=1}^s
   \in\bigoplus_{a=1}^s L^2(\partial\Omega;\mathbb C^{1\times r}).
 \label{eq:V2-boundary-map}
\end{equation}
Its zero set and rank are independent of the eigenbasis.  Subtracting a
constant scalar to form $V^\circ$ leaves this boundary map unchanged.

\begin{theorem}[Normal action jet forces a full-primitive fibre residual]
\label{thm:V2-boundary-birth}
On \eqref{eq:V2-Dirichlet-chart}, if $\bjet u\ne0$, then
\begin{equation}
 G_\lambda(u)\in\Dom(H_R)\setminus\Dom(H_R^2),\qquad
 Y_\lambda(u)\in\Dom(H_R^2)\setminus\Dom(H_R^3).
 \label{eq:V2-boundary-domains}
\end{equation}
For every finite complete primitive bank satisfying
\eqref{eq:V2-primitive},
\begin{equation}
 \boxed{
 \Ker R_\lambda^{\rm Duh}\subseteq\Ker\bjet,
 \qquad
 \rank R_\lambda^{\rm Duh}\ge\rank\bjet.}
 \label{eq:V2-boundary-rank}
\end{equation}
In particular a nonzero boundary jet certifies a positive dynamic residual
without resolving the entire right spectrum or the compact-horizon inverse.
\end{theorem}

\begin{proof}
The adjoint-vector representatives of the rows of $G_\lambda(u)$ are
$g_a=V^\circ(u)^*\psi_a$.  They are smooth and vanish on the boundary, so
$g_a\in\Dom\hD$.  Direct product differentiation gives
\[
 (\hD-\lambda)g_a
 =-\kappa(\Delta V(u)^*)\psi_a
   -2\kappa\sum_j(\partial_jV(u)^*)\partial_j\psi_a
   +[W,V(u)^*]\psi_a.
\]
At the boundary, $\psi_a=0$ and its tangential derivatives vanish.  Therefore
\begin{equation}
 \gamma_0((\hD-\lambda)g_a)
 =-2\kappa(\partial_nV(u)^*)\partial_n\psi_a.
 \label{eq:V2-boundary-commutator}
\end{equation}
If this trace is nonzero, $(\hD-\lambda)g_a$ is not in
$H_0^1$, and hence is not in $\Dom\hD$.  By the recursive definition
$\Dom\hD^2=\{g\in\Dom\hD:\hD g\in\Dom\hD\}$, one has
$g_a\notin\Dom\hD^2$.  Finite row assembly gives the first statement in
\eqref{eq:V2-boundary-domains}.  Its second statement and strict residual
follow from \Cref{thm:V2-domain-defect} with $k=1$.

Thus $R_\lambda^{\rm Duh}u=0$ implies $\bjet u=0$.  Both are linear kernel
conditions after the usual conjugate-vector identification of rows, so the
inclusion holds on the complexified coefficient space.  Rank--nullity gives
\eqref{eq:V2-boundary-rank}.  Changing the left eigenbasis is a unitary mixing
of the boundary rows and preserves its kernel and rank.
\end{proof}

\begin{remark}[Sufficiency, not a complete boundary classification]
A zero first boundary jet does not prove followerhood.  A higher compatibility
jet can fail, and membership in all power domains does not alone give the
right heat-analyticity of a primitive column.  Conversely, a nonzero jet proves
strict positivity at one fixed card but gives no universal Hilbert-norm floor:
finite spectral sums are in every power domain and are dense.  A growing
predeclared spectral dictionary can therefore approximate a nonsmooth target
arbitrarily well.  A numerical lower witness or a quantitative comparison is
still required for a cofinal floor.
\end{remark}

\section{A frozen scalar card with a physical profile cost}
\label{sec:V2-scalar-card}

The following control is fixed by its domain, action command, temperature,
delay, and source-cost experiment, before any residual is inspected:
\begin{equation}
 \boxed{
 \begin{gathered}
 \mathcal H=L^2(0,\pi),\qquad
 \hD(u)=-\partial_x^2+u\,v(x),\qquad v(x)=x-\pi/2,\\
 \Dom\hD(u)=H^2(0,\pi)\cap H_0^1(0,\pi),\qquad |u|\le1/4,\\
 u_*=0,\qquad \beta=1,\qquad b=1/2,\qquad\tau=1/4.
 \end{gathered}}
 \label{eq:V2-frozen-control}
\end{equation}
There is one sector, a trivial constant line, no line-current column, and the
complete declared prior bank is $(F,A)$, with
\begin{equation}
 F=Z^{-1/2}e^{-3\hD/4},\qquad
 A=Z^{-1/2}e^{-\hD/4}v e^{-\hD/2},\qquad
 Z=\operatorname{Tr}e^{-\hD},\quad\hD=\hD(0).
 \label{eq:V2-control-sources}
\end{equation}
The target is $Y=D_u(Z(u)^{-1/2}e^{-\hD(u)/2})|_{u=0}$.
The identity column is retained as the inherited weight primitive even though
its target derivative is zero on this one-sector centered command.

\begin{definition}[Independent profile-cost calibration]
\label{def:V2-profile-cost}
The cost of the source command $(c,u)\in\mathbb C^2$ is declared to be the
normalized physical profile action
\begin{equation}
 \frac1\pi\int_0^\pi |c+u(x-\pi/2)|^2\,dx
   =|c|^2+\frac{\pi^2}{12}|u|^2.
 \label{eq:V2-profile-action}
\end{equation}
Thus, in the order $(F,A)$,
\begin{equation}
 \Msrc=\diag(1,\pi^2/12),\qquad \Ntar=\pi^2/12.
 \label{eq:V2-profile-metric}
\end{equation}
This is part of the diagnostic experiment, not a value inferred for the actual
ESM electromagnetic metric.  It is neither the thermal source Gram nor a
normalization fitted to the residual.
\end{definition}

\begin{proposition}[The frozen control carries the inherited thermal source formulas]
\label{prop:V2-control-validity}
The family \eqref{eq:V2-frozen-control} is self-adjoint on one domain,
nonnegative, and has compact resolvent and trace-class heat operators.
Complex conjugation is an antiunitary reflection commuting with the family.
At the base point,
\begin{equation}
 \psi_m(x)=\sqrt{2/\pi}\sin(mx),\qquad\lambda_m=m^2,
 \qquad Z=\sum_{m\ge1}e^{-m^2}.
 \label{eq:V2-sine-data}
\end{equation}
The Gibbs state is faithful, $\operatorname{Tr}(\rho v)=0$, and the target is
exactly the centered Duhamel half-source of the inherited one-sector formula.
The source metric is positive, so null-cost consistency holds.
\end{proposition}

\begin{proof}
The sine functions give the complete positive Dirichlet realization of
$-\partial_x^2$.  Bounded multiplication by $uv$ preserves self-adjointness
and the domain; the resolvent identity preserves compact resolvent.
Since $\|v\|_\infty=\pi/2$,
$\hD(u)\ge1-\pi/8>1/2$.  The min--max principle, or the variational
characterization on the sine subspaces, bounds its eigenvalues below by
$m^2-\pi/8$, proving heat-trace finiteness.  All coefficients are real,
which proves reflection covariance.

At $u=0$, the change $x\mapsto\pi-x$ fixes $|\psi_m|^2$ and reverses $v$.
Hence every diagonal $v_{mm}$ is zero.  It follows that
$\operatorname{Tr}(\rho v)=0$ and $D_u Z(u)|_0=0$.  Duhamel differentiation
therefore gives \eqref{eq:V2-half-source} with no omitted normalization term.
The positive heat eigenvalues give faithfulness.  Finally,
\eqref{eq:V2-profile-action} is a direct integral and its two diagonal entries
are strictly positive.
\end{proof}

\begin{theorem}[Complete insertion table and literal target coefficients]
\label{thm:V2-insertion-table}
In the sine basis,
\begin{equation}
 v_{mn}=\langle\psi_m,v\psi_n\rangle
 =\begin{cases}
 \displaystyle-\frac{8mn}{\pi(m^2-n^2)^2},&m+n\text{ odd},\\[0.4em]
 0,&m+n\text{ even}.
 \end{cases}
 \label{eq:V2-v-table}
\end{equation}
For $m\ne n$ let
\begin{equation}
 d_{mn}=\frac{e^{-n^2/2}-e^{-m^2/2}}{n^2-m^2},
 \qquad d_{mm}=-\tfrac12e^{-m^2/2}.
 \label{eq:V2-d-table}
\end{equation}
Then the exact source and target matrices are
\begin{equation}
 \boxed{
 \begin{aligned}
 F_{mn}&=Z^{-1/2}e^{-3m^2/4}\delta_{mn},\\
 A_{mn}&=Z^{-1/2}e^{-m^2/4-n^2/2}v_{mn},\\
 Y_{mn}&=Z^{-1/2}v_{mn}d_{mn}.
 \end{aligned}}
 \label{eq:V2-all-coefficients}
\end{equation}
This includes all aliases and both right-energy directions; no nearest-neighbor
spectral approximation is used.
\end{theorem}

\begin{proof}
The product-to-sum identity gives
$2\sin(mx)\sin(nx)=\cos((m-n)x)-\cos((m+n)x)$.
For nonzero integer $k$,
\[
 \int_0^\pi x\cos(kx)\,dx=\frac{(-1)^k-1}{k^2}.
\]
The off-diagonal integral is therefore zero for equal parity, and for opposite
parity equals
$-2\pi^{-1}[(m-n)^{-2}-(m+n)^{-2}]$, which is
\eqref{eq:V2-v-table}.  The diagonal vanishes by the same reflection symmetry
used above.  Multiplying by the two heat factors gives $F$ and $A$.
Integration of the Duhamel coefficient gives $d_{mn}$, including its continuous
diagonal value.  This proves \eqref{eq:V2-all-coefficients}.
\end{proof}

\section{An evaluated five-matrix fibre, reserve window, and bounded annihilator}
\label{sec:V2-five-matrix-values}

For the left fibre $m^2$, set
\begin{align}
 s_m&=\sum_{n\ge1}e^{-n^2}|v_{mn}|^2,
 &t_m&=\sum_{n\ge1}e^{-n^2/2}|v_{mn}|^2d_{mn},
 &u_m&=\sum_{n\ge1}|v_{mn}|^2d_{mn}^2.
 \label{eq:V2-scaled-stu}
\end{align}
Each $s_m>0$ because there is an opposite-parity right mode.
The exact V15 matrices, in the role order $(A,F,Y)$, are scalars:
\begin{equation}
 \boxed{
 S_m=Z^{-1}e^{-m^2/2}s_m,\quad
 T_m=Z^{-1}e^{-m^2/4}t_m,\quad
 U_m=Z^{-1}u_m,\quad
 G_m=Z^{-1}e^{-3m^2/2},\quad C_m=0.}
 \label{eq:V2-five-values}
\end{equation}
Here $A_m^*F_m=0$ and $F_m^*Y_m=0$ by diagonal versus off-diagonal support.
Consequently the complete fibre Gram and residual are
\begin{equation}
 \Gamma_m=
 \begin{pmatrix}S_m&0&T_m\\0&G_m&0\\T_m&0&U_m\end{pmatrix},
 \qquad
 R_m=\frac1Z\left(u_m-\frac{t_m^2}{s_m}\right).
 \label{eq:V2-fibre-Gram}
\end{equation}
Primitive absorption is exactly zero on this declared bank; it has not merely
been neglected.  These formulas use the inherited Schur theorem rather than
reproving it.

\begin{theorem}[Two right modes give an exact positive lower certificate]
\label{thm:V2-two-mode-witness}
On the first left fibre, write $a_n=A_{1n}$ and $y_n=Y_{1n}$.
The unit vector
\begin{equation}
 w=\frac{a_4\,|1\rangle\langle2|-a_2\,|1\rangle\langle4|}
          {\sqrt{a_2^2+a_4^2}}
 \label{eq:V2-dual-vector}
\end{equation}
is orthogonal to the \emph{complete} left-cyclic primitive bank of this control.
It gives
\begin{equation}
 \boxed{
 R_{\rm dyn}\ge R_1\ge
 \mathfrak d:=\frac{(a_2y_4-a_4y_2)^2}{a_2^2+a_4^2}>0.}
 \label{eq:V2-dual-lower}
\end{equation}
The certificate has the following expression without square roots:
\begin{equation}
 \mathfrak d=
 \frac{|v_{12}|^2|v_{14}|^2}{Z}
 \frac{(e^{-2}d_{14}-e^{-8}d_{12})^2}
      {e^{-4}|v_{12}|^2+e^{-16}|v_{14}|^2}.
 \label{eq:V2-dual-exact}
\end{equation}
With the fixed target metric \eqref{eq:V2-profile-metric},
\begin{equation}
 \boxed{R_{\rm dyn}^{\sharp}\ge\frac{12\mathfrak d}{\pi^2}>\frac1{100000}.}
 \label{eq:V2-rational-floor}
\end{equation}
\end{theorem}

\begin{proof}
The vector $w$ is off-diagonal, so it is orthogonal to every delayed weight
column.  Against $A$ its two entries give $a_4a_2-a_2a_4=0$.
Every left delay multiplies both entries by $e^{-t}$, so the cancellation
holds at every time.  It also holds against every other left fibre by
orthogonality.  Thus $w$ is an exact invariant-complement witness.
Bessel's inequality gives the first lower bound in
\eqref{eq:V2-dual-lower}, and substitution gives
\eqref{eq:V2-dual-exact}.

For strict positivity, on the ground-left fibre the multiplier relating $y_n$
to $a_n$ is
\[
 -e^{1/4}\frac{e^{(n^2-1)/2}-1}{n^2-1}.
\]
The function $(e^{z/2}-1)/z=\int_0^{1/2}e^{sz}\,ds$ is strictly increasing
for $z>0$.  Its values at $3$ and $15$ differ, and both $a_2,a_4$ are nonzero.
Hence the determinant in \eqref{eq:V2-dual-exact} is nonzero.
The rational numerical inequality in \eqref{eq:V2-rational-floor} is certified
in \Cref{prop:V2-outward-certificate}; its proof and reproducing exact-arithmetic
code are supplied below.
\end{proof}

\begin{theorem}[Physical source reserve of the same control]
\label{thm:V2-reserve}
The global static source Gram is diagonal in the physical order $(F,A)$:
\begin{equation}
 Q=\diag(q_F,q_A),\qquad
 q_F=\frac{\sum_m e^{-3m^2/2}}Z,\qquad
 q_A=\frac{\sum_{m,n}e^{-m^2/2-n^2}|v_{mn}|^2}{Z}.
 \label{eq:V2-global-Q}
\end{equation}
The supported unit-action reserve eigenvalues are exactly
\begin{equation}
 g_F=q_F,\qquad g_A=12q_A/\pi^2,
 \label{eq:V2-reserve-eigenvalues}
\end{equation}
and they satisfy
\begin{equation}
 \boxed{0.0616\,P_B\preceq B\Msrc^{-1}B^*\preceq0.5841\,P_B.}
 \label{eq:V2-reserve-window-numeric}
\end{equation}
Thus the positive residual in \eqref{eq:V2-rational-floor} is not caused by
source-reserve collapse or coefficient rescaling.
\end{theorem}

\begin{proof}
Diagonal support of $F$ and zero diagonal of $A$ give $F^*A=0$.
The squared Hilbert--Schmidt norms give \eqref{eq:V2-global-Q}.
The physical metric is the independent diagonal matrix
\eqref{eq:V2-profile-metric}, so the V1 reserve formula gives
\eqref{eq:V2-reserve-eigenvalues}.  The two nonzero eigenvalues are enclosed in
\Cref{prop:V2-outward-certificate}; both lie strictly between the displayed
rational endpoints.  The ambient covariance has these eigenvalues on its
range, giving \eqref{eq:V2-reserve-window-numeric}.
\end{proof}

\begin{remark}[The infinite right tail is visible, not silently dropped]
For even $n\to\infty$,
\begin{equation}
 A_{1n}\sim-\frac{8e^{-1/4}}{\pi\sqrt Z}\,n^{-3}e^{-n^2/2},
 \qquad
 Y_{1n}\sim\frac{8e^{-1/2}}{\pi\sqrt Z}\,n^{-5}.
 \label{eq:V2-tail-asymptotics}
\end{equation}
These follow by dividing the explicit rational and exponential factors by
their displayed leading terms.  Since the right energies are $n^2$,
$Y_1\in\Dom(H_R^2)$ but $Y_1\notin\Dom(H_R^3)$, whereas $A_1$ belongs to
every right power domain.  This checks the boundary theorem directly:
$v'=1$ and both endpoint derivatives of $\psi_1$ are nonzero.  The bounded
two-mode vector in \eqref{eq:V2-dual-vector}, rather than a singular boundary
functional, is the quantitative lower witness.
\end{remark}

\section{Complete positive tails and reproducible outward intervals}
\label{sec:V2-tails}

All infinite sums in this calculation have explicit error bounds.  For $c>0$
and an integer $N\ge1$ define
\begin{equation}
 \mathcal E_c(N)=
 \frac{e^{-c(N+1)^2}}{1-e^{-c(2N+3)}}.
 \label{eq:V2-Gaussian-tail}
\end{equation}
Then $\sum_{n>N}e^{-cn^2}\le\mathcal E_c(N)$, because the ratio of successive
terms is at most $e^{-c(2N+3)}$.  Parseval and
$\|v\|_\infty^2=\pi^2/4$ give
\begin{equation}
 \sum_n|v_{mn}|^2\le\pi^2/4,
 \qquad
 \sum_m|v_{mn}|^2\le\pi^2/4.
 \label{eq:V2-Parseval-cap}
\end{equation}

\begin{lemma}[Finite right-sum enclosures for every retained left fibre]
\label{lem:V2-right-tails}
Let $N>m$ and put $\alpha_{m,N}=1-m^2/(N+1)^2>0$.
If $s_{m,N},t_{m,N},u_{m,N}$ denote the sums in
\eqref{eq:V2-scaled-stu} restricted to $n\le N$, then
\begin{align}
 0\le s_m-s_{m,N}
 &\le\frac{\pi^2}{4}e^{-(N+1)^2},
 \label{eq:V2-s-tail}\\
 0\le t_{m,N}-t_m
 &\le\frac{\pi^2 e^{-((N+1)^2+m^2)/2}}
               {4((N+1)^2-m^2)},
 \label{eq:V2-t-tail}\\
 0\le u_m-u_{m,N}
 &\le\frac{64m^2e^{-m^2}}
 {9\pi^2\alpha_{m,N}^{6}N^9}.
 \label{eq:V2-u-tail}
\end{align}
The global unnormalized action-source sum restricted to $m,n\le N$ has omitted
mass at most
\begin{equation}
 \frac{\pi^2}{4}\bigl(\mathcal E_{1/2}(N)+\mathcal E_1(N)\bigr).
 \label{eq:V2-Q-tail}
\end{equation}
\end{lemma}

\begin{proof}
The first bound is \eqref{eq:V2-Parseval-cap} after pulling out the largest
omitted exponential.  For $n>N>m$, $d_{mn}<0$ and
$|d_{mn}|\le e^{-m^2/2}/(n^2-m^2)$, which gives
\eqref{eq:V2-t-tail}.  The explicit insertion table yields
\[
 |v_{mn}|^2d_{mn}^2
 \le\frac{64m^2n^2e^{-m^2}}{\pi^2(n^2-m^2)^6}
 \le\frac{64m^2e^{-m^2}}{\pi^2\alpha_{m,N}^6}\,n^{-10}.
\]
Since $\sum_{n>N}n^{-10}\le\int_N^\infty x^{-10}dx=N^{-9}/9$,
\eqref{eq:V2-u-tail} follows.  For \eqref{eq:V2-Q-tail}, cover the omitted
rectangle by $m>N$ or $n>N$, discard the other exponential, and apply
\eqref{eq:V2-Parseval-cap} and \eqref{eq:V2-Gaussian-tail}.
Possible overlap is overcounted only in an upper bound.
\end{proof}

\begin{theorem}[Complete left-energy tail and the total dynamic residual]
\label{thm:V2-total-tail}
Put $K=\sum_{n\ge1}n^2e^{-n^2}$.  The full residue is
\begin{equation}
 R_{\rm dyn}=\sum_{m\ge1}R_m,
 \qquad R_m=\frac1Z\left(u_m-\frac{t_m^2}{s_m}\right)\ge0.
 \label{eq:V2-total-residue}
\end{equation}
For every integer $M\ge1$, its positive omitted-left tail obeys
\begin{equation}
 \boxed{
 0\le\sum_{m>M}R_m\le\sum_{m>M}U_m\le\mathcal T(M):=
 \frac{64(4/3)^6K}{9\pi^2ZM^9}
 +\frac{\pi^2}{16Z}\mathcal E_{1/4}(M).}
 \label{eq:V2-total-tail}
\end{equation}
In particular the series in \eqref{eq:V2-total-residue} is a norm-convergent
positive sum and supplies a complete left-energy tail, not just a zero-head
diagnostic.
\end{theorem}

\begin{proof}
For each omitted $m$, split the right modes into $n\le m/2$ and $n>m/2$.
In the first region $m^2-n^2\ge3m^2/4$, and the same calculation as above gives
\[
 Z^{-1}|v_{mn}|^2d_{mn}^2
 \le\frac{64(4/3)^6}{\pi^2Z}\,\frac{n^2e^{-n^2}}{m^{10}}.
\]
Summing over $n$ and then $m>M$ gives the first term in
\eqref{eq:V2-total-tail}.  In the second region the mean-value formula for the
exponential gives
\[
 |d_{mn}|\le\tfrac12e^{-\min(m^2,n^2)/2},
 \qquad |d_{mn}|^2\le\tfrac14e^{-m^2/4}.
\]
Summing $|v_{mn}|^2$ with \eqref{eq:V2-Parseval-cap} gives the second term.
The same estimates and a finite head prove $Y\in\HS$.
The imported fibre formula now applies: every left fibre contributes its
positive full-primitive residual, and $R_m\le U_m$ controls the positive sum.
\end{proof}

\begin{proposition}[Executed exact outward certificate]
\label{prop:V2-outward-certificate}
The following intervals are outward enclosures.  Every printed endpoint is a
terminating decimal, hence a rational number.
\begin{center}\small
\begin{tabular}{@{}lll@{}}
\toprule
Quantity & Lower endpoint & Upper endpoint\\
\midrule
$Z$ & $0.386318602413326076$ & $0.386318602413326077$\\
$S_1$ & $0.009208405346085056$ & $0.009208405346085057$\\
$T_1$ & $-0.013722351317441194$ & $-0.013722351317441193$\\
$U_1$ & $0.020457652758072111$ & $0.020457652758072150$\\
$G_1$ & $0.577580677592379239$ & $0.577580677592379240$\\
$C_1$ & $0$ & $0$\\
$R_1$ & $0.000008626178313994$ & $0.000008626178314033$\\
$\mathfrak d$ & $0.000008498137176950$ & $0.000008498137176951$\\
$12\mathfrak d/\pi^2$ & $0.000010332495810282$ & $0.000010332495810283$\\
$q_F$ & $0.584000568216363821$ & $0.584000568216363822$\\
$q_A$ & $0.050691181987332573$ & $0.050691181987332574$\\
$12q_A/\pi^2$ & $0.061633086710228258$ & $0.061633086710228259$\\
$R_{\rm dyn}$ & $0.000478229699853324$ & $0.000478229767293226$\\
$R_{\rm dyn}^{\sharp}$ & $0.000581457590904705$ & $0.000581457672901793$\\
\bottomrule
\end{tabular}
\end{center}
The source reserve window and the strict rational lower certificate
\eqref{eq:V2-rational-floor} follow from these intervals.
\end{proposition}

\begin{proof}
Use the grid $10^{-70}\mathbb Z$.  For an interval with integer endpoints
$(a,b)$ in grid units, addition is endpoint addition; multiplication takes
all four endpoint products and rounds the lower down and upper up after
dividing by $10^{70}$; inversion takes reciprocal endpoints and rounds outward.
Thus induction through a finite expression preserves containment.

The constants are enclosed without floating-point arithmetic.  Machin's
identity
$\pi=16\arctan(1/5)-4\arctan(1/239)$ follows from the tangent addition formula
and the angle range.  For each arctangent retain terms $k=0,\ldots,159$ of
its alternating series, with term $160$ bounding the remainder.
For $e^{-1/4}$ retain terms $k=0,\ldots,119$ of
$\sum_k(-1/4)^k/k!$, with term $120$ bounding the remainder.
All exponentials used are nonnegative integer powers of this enclosed
$e^{-1/4}$ and are evaluated by repeated squaring with outward rounding.

For $Z$ and the heat-trace sums retain $16$ terms and add
\eqref{eq:V2-Gaussian-tail}.  For $(S_1,T_1,U_1)$ retain right indices
$n\le64$ and add \eqref{eq:V2-s-tail}--\eqref{eq:V2-u-tail}, with the
negative tail added on the correct side of $T_1$.  For $q_A$ retain the
$10\times10$ rectangle and add \eqref{eq:V2-Q-tail}.  For the total residue,
retain $m\le16$, using $n\le128$ and the complete right tails on each row,
and then add \eqref{eq:V2-total-tail}.  The constant $K$ is enclosed by its
first $16$ terms and $\mathcal E_{1/2}(16)$, since
$n^2\le e^{n^2/2}$.  Independently proved positivity permits replacing a
negative lower interval endpoint of a residual by zero.

The exact arithmetic implementation in \Cref{sec:V2-code} evaluates these
finite rational expressions and produces the displayed intervals.  It also
asserts, by integer comparisons, the strict rational lower bound and the
reserve endpoints.  No floating-point diagonalization, fitted support
threshold, empirical truncation test, or statistical confidence assertion
enters this certificate.
\end{proof}

\section{Raw horizon kernels and what the spectral regulator actually transports}
\label{sec:V2-kernels-transport}

\begin{proposition}[Complete raw delayed packet on the control]
\label{prop:V2-raw-kernels}
In source order $(F,A)$ the raw delayed source Gram is diagonal and, for every
$t\ge0$,
\begin{align}
 K_{FF}(t)&=Z^{-1}\sum_m e^{-(t+3/2)m^2},\\
 K_{AA}(t)&=Z^{-1}\sum_{m,n}e^{-(t+1/2)m^2-n^2}|v_{mn}|^2,\\
 K_{FA}(t)&=K_{FY}(t)=0,\\
 K_{AY}(t)&=Z^{-1}\sum_{m,n}
 e^{-(t+1/4)m^2-n^2/2}|v_{mn}|^2d_{mn},\\
 Y^*Y&=Z^{-1}\sum_{m,n}|v_{mn}|^2d_{mn}^2.
 \label{eq:V2-raw-kernels}
\end{align}
The sums and their common finite-rectangle approximations converge uniformly
on every compact nonnegative time interval.  Together with
\eqref{eq:V2-profile-metric}, they are a literal raw input to both V1 horizon
constructions.
\end{proposition}

\begin{proof}
Insert \eqref{eq:V2-all-coefficients} into the Hilbert--Schmidt inner product;
left delay contributes $e^{-tm^2}$.  The vanishing mixed terms follow from
diagonal/off-diagonal orthogonality.  The source-source tails are bounded by
\eqref{eq:V2-Q-tail} and the heat-trace tail, uniformly because $e^{-tm^2}\le1$.
The mixed source-target tail outside a common rectangle is bounded by the
product of the omitted source and target Hilbert norms, by Cauchy--Schwarz.
Those norms vanish by the source bounds and \Cref{thm:V2-total-tail}, using
$Y_{mn}=Y_{nm}$ for the second target index.  The same bounds prove uniform
absolute convergence.  The final assertion is the direct substitution into
\Cref{thm:raw-normalized-horizon,thm:raw-physical-horizon}.
\end{proof}

\begin{remark}[The positive-head stopping rule is actually used]
The compact-horizon Tikhonov inverse is not numerically evaluated in V2.
It is unnecessary for the positive branch, since
\eqref{eq:V2-dual-vector} is already a bounded annihilator of the complete
cyclic bank.  V2 supplies exact raw horizon kernels and their tails, but does
not claim to have printed a numerical Tikhonov curve or a numerical
Duhamel--horizon cross-gap.  The complete residue interval above comes from
left-fibre sums and positive tails instead.
\end{remark}

Let $P_N$ project onto $\psi_1,\ldots,\psi_N$, and define the diagnostic
\emph{energy} regulator
\begin{equation}
 \hD_N(u)=P_N\hD P_N+uP_NvP_N\quad\hbox{on }P_N\mathcal H,
 \qquad Z_N=\sum_{m=1}^N e^{-m^2}.
 \label{eq:V2-spectral-regulator}
\end{equation}
Construct $(F_N,A_N,Y_N)$ using the same $\beta,\tau,\Msrc,\Ntar$.
At $u=0$ their embedded matrices satisfy exactly
\begin{equation}
 X_N=\sqrt{Z/Z_N}\,P_NXP_N,
 \qquad X\in\{F,A,Y\}.
 \label{eq:V2-spectral-incidence}
\end{equation}
This identity holds for the displayed base point; it is not an assertion that
$e^{-\hD_N(u)}$ is the compressed full heat operator at every nonzero $u$.

\begin{theorem}[Persistent birth and reserve along the declared energy truncations]
\label{thm:V2-spectral-persistence}
For every $N\ge4$, the same two-right-mode witness gives
\begin{equation}
 R_{{\rm dyn},N}^{\sharp}
 \ge\frac{12}{\pi^2}\frac Z{Z_N}\mathfrak d
 >10^{-5}.
 \label{eq:V2-spectral-lower}
\end{equation}
The full source bank has rank two and obeys, uniformly for $N\ge4$,
\begin{equation}
 0.0616\,P_{B_N}\preceq B_N\Msrc^{-1}B_N^*
                   \preceq0.5841\,P_{B_N}.
 \label{eq:V2-spectral-reserve}
\end{equation}
Moreover, in the stated common Hilbert embedding,
\begin{equation}
 B_N\longrightarrow B,\qquad Y_N\longrightarrow Y,
 \qquad R_{{\rm dyn},N}\longrightarrow R_{\rm dyn}
 \label{eq:V2-spectral-limit}
\end{equation}
in Hilbert--Schmidt synthesis norm and, for the last scalar, in ordinary norm.
\end{theorem}

\begin{proof}
The state centering remains zero because every $v_{mm}=0$.
Thus \eqref{eq:V2-spectral-incidence} follows entry by entry, including the
Duhamel derivative.  The $n=2,4$ ground-row witness is unchanged after
normalization of its direction, while its target pairing is multiplied by
$\sqrt{Z/Z_N}$.  This proves \eqref{eq:V2-spectral-lower}.

For the reserve, let $Q_{A,4}^{\rm raw}$ be the unnormalized action-source sum
on the $4\times4$ rectangle and $Q_A^{\rm raw}=Zq_A$ its complete sum.  Then
\[
 \frac{12Q_{A,4}^{\rm raw}}{\pi^2Z}
 \le\frac{12q_{A,N}}{\pi^2}
 \le\frac{12Q_A^{\rm raw}}{\pi^2Z_4},
 \qquad
 \frac{e^{-3/2}}Z\le q_{F,N}
 \le\frac{\sum_m e^{-3m^2/2}}{Z_4}.
\]
The same exact arithmetic used above encloses the first lower endpoint above
$0.061633085990638666$, the action upper below $0.061633086712443972$,
the weight lower above $0.577580677592379239$, and the weight upper below
$0.584000568237358676$.  This proves \eqref{eq:V2-spectral-reserve}.

The first two convergences follow from \eqref{eq:V2-spectral-incidence},
$Z_N\to Z$, and Hilbert--Schmidt rectangle exhaustion.  For the residue,
fix a finite left head $m\le M$.  The finite right-row sums converge to
$s_m,t_m,u_m$, with $s_m>0$, so each residual converges by
\eqref{eq:V2-fibre-Gram}.  The remaining positive tail is uniformly bounded by
\[
 \sum_{m>M}R_{m,N}\le\sum_{m>M}U_{m,N}
 \le (Z/Z_1)\sum_{m>M}U_m\le (Z/Z_1)\mathcal T(M).
\]
Let $N\to\infty$ at fixed $M$, then $M\to\infty$.  This proves the last
convergence without assuming continuity of an infinite cyclic projection.
\end{proof}

\begin{proposition}[Exact direct/staged incidence, with its normalization retained]
\label{prop:V2-spectral-routes}
For $M\le N$ set
\begin{equation}
 \mathcal T_{M\leftarrow N}(X)
   =\sqrt{Z_N/Z_M}\,P_MXP_M.
 \label{eq:V2-spectral-route}
\end{equation}
At the base point this maps every primitive and target row on card $N$ to the
corresponding row on card $M$, and
\[
 \mathcal T_{L\leftarrow M}\mathcal T_{M\leftarrow N}
   =\mathcal T_{L\leftarrow N}\qquad(L\le M\le N).
\]
It intertwines left delays on the retained rows.  It is generally not an
isometry, and no equality of the shorted residues at different $N$ is asserted.
\end{proposition}

\begin{proof}
Nested spectral projections multiply to the smaller projection, and
$\sqrt{Z_M/Z_L}\sqrt{Z_N/Z_M}=\sqrt{Z_N/Z_L}$.
The source and target identities follow from
\eqref{eq:V2-spectral-incidence}.  Spectral projections commute with the base
Hamiltonian, proving the delayed identity.  Compression discards right rows
and the scalar normalization is not one, explaining the stated limitation.
\end{proof}

\begin{firewall}[Energy exhaustion is not an ESM spacetime continuum]
\label{fire:V2-energy-not-continuum}
The sequence \eqref{eq:V2-spectral-regulator} removes an auxiliary spectral
truncation of one fixed interval Hamiltonian.  It is not a spatial lattice
refinement, thermodynamic limit, overlap-sector crossing, Wilsonian
renormalization, or the source transport of the full ESM branch.  Its exact
spectral routes do not populate an independent physical direct/staged cutoff
experiment.  The result is a complete infinite-dimensional diagnostic with a
persistent ancestry residual, not an interacting four-dimensional continuum.
\end{firewall}

\section{What V2 changes on the actual ESM programme}
\label{sec:V2-ESM-export}

\begin{corollary}[Conditional export to a genuine fixed-wall ESM card]
\label{cor:V2-ESM-export}
Suppose an occupied reflection-closed ESM card, with its original source metric
and complete prior primitive bank, has a fixed Dirichlet chart satisfying
\eqref{eq:V2-Dirichlet-chart}.  Suppose its actual action derivatives are the
smooth multipliers used in \eqref{eq:V2-boundary-map}, all prior fibre columns
satisfy \eqref{eq:V2-primitive}, and the inherited Duhamel identification holds
on that same card.  Then evaluation of its native boundary map gives
\[
 \rank R_{{\rm dyn},\lambda}\ge\rank\mathcal J_{\partial,\lambda}.
\]
A nonzero jet is a valid positive-head stop after the physical source-metric
checks; no separate complete horizon evaluation is needed to establish that
finite-card positive branch.
\end{corollary}

\begin{proof}
Apply \Cref{thm:V2-boundary-birth} to the complete prior bank on that card and
then the imported positive left-fibre sum.  This implication changes neither
the physical metric nor the source declaration.  It is conditional precisely
on the displayed native chart and incidence hypotheses.
\end{proof}

\paragraph{The next upstream calculation.}
Before requesting another complete five-matrix field, locate the actual
hard-wall operator domain and differentiate the fixed-chart common action on
one predeclared occupied left eigenspace.  Test its boundary/domain incidence.
A nonzero normal jet supplies a finite-card positive branch by the preceding
corollary.  On a zero-jet branch, examine the next compatibility condition or
compute the literal fibre matrices; do not infer followerhood.  If the wall
moves, the operator is not a Dirichlet realization, or a prior primitive has
unsmoothed boundary support, retain that fact and use the appropriate native
domain calculation instead.  The source documents do not specify these data
well enough to report their ESM value here.

\paragraph{Quantitative cofinal requirement.}
Strict positivity from a domain mismatch is not a regulator-uniform number.
The actual ESM programme still needs an independently normalized bounded
annihilator or a positive finite-head floor, plus the physical source-reserve
and primitive cutoff routes.  The two-mode control demonstrates one way such
a lower witness can be extracted; its sine modes, metric, and numerical floor
are not exported to ESM by analogy.

\begin{assessment}[Version V2 result]
\label{ass:V2-result}
\textbf{Proved:} a literal finite-rate path/thermal clock-incidence obstruction,
its finite-lag source witness and increasing-rate approximation bounds, a
Duhamel right-domain obstruction, an explicit Dirichlet
boundary rank test, and a complete evaluated scalar hard-wall diagnostic with
positive source reserve, a certified positive full-primitive residual,
complete spectral tails, and exact diagnostic energy routes.

\textbf{Physical ESM value not yet evaluated:} the native wall/domain jet,
actual routed metric and reserve, full ESM fibre residual, and physical
cofinal ancestry route.  Current, stress, determinant/Pfaffian orientation,
interacting continuum, Lorentzian reconstruction, Einstein backreaction, and
pure-colour closure remain downstream.

The native clock audit returns an exact obstruction to a specific path/thermal
identification, not to either independently constructed clock.  The thermal
calculation returns a \obsstatus{} to complete followerhood on the frozen
hard-wall control, with a bounded two-mode witness.  Its applicability to the
actual ESM card is the explicit conditional corollary, not a fabricated
physical numerical verdict.  Unlike V1's unevaluated acquisition, this note
contains executed source-specific infinite-dimensional calculations.  The
conditional application to the physical ESM branch remains unevaluated.
\end{assessment}

\section{Reproducing the exact rational certificate}
\label{sec:V2-code}

The following Python 3 program uses only standard-library integer and rational
arithmetic.  Its output is the interval table above, including the uniform
energy-regulator checks.  The companion source file is
\begin{center}\small
\path{Renewal_SM_Einstein_Hard_Wall_Certificate_V2.py}.
\end{center}
The exact script SHA--256 is
\begin{center}\small\ttfamily
2d902811de44485f6d43b79fd4f5af63cb48746c8ddf440d590439143215dde8.
\end{center}
Running the script is a replay of the finite arithmetic proof, not a physical
acquisition or a replacement for the analytic tail bounds.

\begingroup\scriptsize
\begin{verbatim}
#!/usr/bin/env python3
"""Exact outward-rational certification of the V2 hard-wall control.

No floating-point arithmetic is used in any certified result. Intervals are
integer multiples of 10**(-70); operations round outward. pi and exp(-1/4)
are enclosed by alternating series with explicit next-term remainders.
The card is a scalar Dirichlet diagnostic, not the full ESM regulator.
"""
from __future__ import annotations
from dataclasses import dataclass
from fractions import Fraction as F
from functools import lru_cache
import json

DIGITS = 70
D = 10**DIGITS

def floorfrac(x: F) -> int:
    return x.numerator // x.denominator

def ceilfrac(x: F) -> int:
    return -((-x.numerator) // x.denominator)

@dataclass(frozen=True)
class I:
    lo: int
    hi: int

    def __post_init__(self):
        if self.lo > self.hi:
            raise ValueError('Reversed interval')

    @classmethod
    def q(cls, a=0, b=1):
        x=F(a,b)
        return cls(floorfrac(x*D),ceilfrac(x*D))

    @classmethod
    def bounds(cls, lo: F, hi: F):
        return cls(floorfrac(lo*D),ceilfrac(hi*D))

    def __add__(self,o):
        if not isinstance(o,I): o=I.q(o)
        return I(self.lo+o.lo,self.hi+o.hi)
    __radd__=__add__

    def __neg__(self): return I(-self.hi,-self.lo)
    def __sub__(self,o):
        if not isinstance(o,I): o=I.q(o)
        return self+(-o)
    def __rsub__(self,o): return I.q(o)-self

    def __mul__(self,o):
        if not isinstance(o,I): o=I.q(o)
        ps=[a*b for a in (self.lo,self.hi) for b in (o.lo,o.hi)]
        return I(min(ps)//D, -((-max(ps))//D))
    __rmul__=__mul__

    def recip(self):
        if self.lo <= 0 <= self.hi: raise ZeroDivisionError('Interval includes zero')
        return I.bounds(F(D,self.hi),F(D,self.lo))
    def __truediv__(self,o):
        if not isinstance(o,I): o=I.q(o)
        return self*o.recip()
    def __rtruediv__(self,o): return I.q(o)/self

    def __pow__(self,n:int):
        if n<0: return (self.recip())**(-n)
        z=I.q(1); x=self
        while n:
            if n&1: z=z*x
            x=x*x; n//=2
        return z

    def positive(self):
        # Use only when positivity has been proved independently.
        return I(max(0,self.lo), max(0,self.hi))

    def plus_positive_tail(self,t):
        if t.hi<0: raise ValueError('Negative tail bound')
        return I(self.lo,self.hi+max(0,t.hi))

    def plus_negative_tail(self,t):
        if t.hi<0: raise ValueError('Negative tail bound')
        return I(self.lo-max(0,t.hi),self.hi)

    def decimal(self,n=18):
        scale=10**(DIGITS-n)
        lo=self.lo//scale; hi=-((-self.hi)//scale)
        def fmt(k):
            sg='-' if k<0 else ''; k=abs(k)
            return f'{sg}{k//10**n}.{k%10**n:0{n}d}'
        return [fmt(lo),fmt(hi)]

    def assert_between(self,a:F,b:F):
        assert F(self.lo,D)>a and F(self.hi,D)<b, (self.decimal(),a,b)


def alternating(terms):
    vals=list(terms)
    s=sum(vals[:-1],F(0)); nxt=vals[-1]
    return I.bounds(min(s,s+nxt),max(s,s+nxt))

def atan_recip(q:int,N=160):
    return alternating(F((-1)**k,(2*k+1)*q**(2*k+1)) for k in range(N+1))

PI=16*atan_recip(5)-4*atan_recip(239)
terms=[]; term=F(1)
for k in range(121):
    if k: term=term*F(-1,4*k)
    terms.append(term)
BASE=alternating(terms)
PI2=PI*PI
VBOUND2=PI2/4

@lru_cache(None)
def e4(k:int):
    """Enclose exp(-k/4), k nonnegative integer."""
    if k<0: raise ValueError('Negative exponent index')
    return BASE**k


def gauss_tail(c4:int,N:int):
    """Sum_{n>N} exp(-(c4/4)n^2), by a geometric ratio bound."""
    return e4(c4*(N+1)**2)/(1-e4(c4*(2*N+3)))


def gauss_sum(c4:int,N=16):
    return sum((e4(c4*n*n) for n in range(1,N+1)),I.q(0)).plus_positive_tail(gauss_tail(c4,N))

Z=gauss_sum(4)

@lru_cache(None)
def v2(m:int,n:int):
    if (m+n)%2==0: return I.q(0)
    return 64*m*m*n*n/(PI2*(m*m-n*n)**4)

@lru_cache(None)
def dd(m:int,n:int):
    if m==n: return -e4(2*m*m)/2
    return (e4(2*n*n)-e4(2*m*m))/(n*n-m*m)


def fibre(m:int,N=64):
    if N<=m: raise ValueError('Right cutoff must exceed left index')
    sr=I.q(0); tr=I.q(0); ur=I.q(0)
    for n in range(1,N+1):
        if (m+n)%2==0: continue
        vv=v2(m,n); d=dd(m,n)
        sr += e4(4*n*n)*vv
        tr += e4(2*n*n)*vv*d
        ur += vv*d*d
    sr=sr.plus_positive_tail(VBOUND2*e4(4*(N+1)**2))
    tr=tr.plus_negative_tail(VBOUND2*e4(2*((N+1)**2+m*m))/((N+1)**2-m*m))
    alpha=I.q(1)-I.q(m*m,(N+1)**2)
    utail=64*m*m*e4(4*m*m)/(PI2*alpha**6*9*N**9)
    ur=ur.plus_positive_tail(utail)
    r=((ur-tr*tr/sr)/Z).positive()
    S=e4(2*m*m)*sr/Z
    T=e4(m*m)*tr/Z
    U=ur/Z
    G=e4(6*m*m)/Z
    return {'S':S,'T':T,'U':U,'G':G,'C':I.q(0),'R':r}


def compute():
    f1=fibre(1,64)
    d2,d4=dd(1,2),dd(1,4)
    z2,z4=v2(1,2),v2(1,4)
    determinant=e4(8)*d4-e4(32)*d2
    dual=z2*z4*determinant**2/(Z*(e4(16)*z2+e4(64)*z4))
    dualnorm=12*dual/PI2
    Qw=gauss_sum(6)/Z
    N=10
    raw=sum((e4(2*m*m+4*n*n)*v2(m,n)
             for m in range(1,N+1) for n in range(1,N+1)),I.q(0))
    raw=raw.plus_positive_tail(VBOUND2*(gauss_tail(2,N)+gauss_tail(4,N)))
    Qa=raw/Z
    reserve=12*Qa/PI2
    # Full dynamic residual: sixteen complete left rows plus an explicit tail.
    M=16
    head=sum((fibre(m,128)['R'] for m in range(1,M+1)),I.q(0))
    K=sum((n*n*e4(4*n*n) for n in range(1,17)),I.q(0))
    # n^2 <= exp(n^2/2) gives this tail enclosure for K.
    K=K.plus_positive_tail(gauss_tail(2,16))
    lefttail=64*I.q(4,3)**6*K/(PI2*Z*9*M**9)+PI2*gauss_tail(1,M)/(16*Z)
    full=head.plus_positive_tail(lefttail)
    # Spectral regulators N>=4 retain the dual witness exactly up to Z/Z_N.
    Z4=sum((e4(4*n*n) for n in range(1,5)),I.q(0))
    raw4=sum((e4(2*m*m+4*n*n)*v2(m,n)
              for m in range(1,5) for n in range(1,5)),I.q(0))
    reserveNlower=12*raw4/(PI2*Z)
    reserveNupper=12*raw/(PI2*Z4)
    weightNlower=e4(6)/Z
    weightNupper=gauss_sum(6)/Z4
    out={'Z':Z,**{k+'1':v for k,v in f1.items()},
         'dual_lower':dual,'dual_lower_normalized':dualnorm,
         'Q_weight':Qw,'Q_action':Qa,'reserve_action':reserve,
         'R_total':full,'left_tail_bound':lefttail,
         'R_total_normalized':12*full/PI2,
         'regulator_action_lower':reserveNlower,
         'regulator_action_upper':reserveNupper,
         'regulator_weight_lower':weightNlower,
         'regulator_weight_upper':weightNupper}
    dualnorm.assert_between(F(1,100000),F(104,10**7))
    reserve.assert_between(F(61633,10**6),F(61634,10**6))
    Qw.assert_between(F(584000,10**6),F(584001,10**6))
    f1['R'].assert_between(F(862617,10**11),F(862619,10**11))
    assert F(reserveNlower.lo,D)>F(616,10000)
    assert F(weightNlower.lo,D)>F(616,10000)
    assert F(reserveNupper.hi,D)<F(5841,10000)
    assert F(weightNupper.hi,D)<F(5841,10000)
    return {k:v.decimal() for k,v in out.items()}

if __name__=='__main__':
    print(json.dumps({'arithmetic':'exact integer outward intervals','decimal_grid_digits':DIGITS,
                      'card':'scalar Dirichlet hard-wall control; not full ESM',
                      'certified_intervals':compute()},indent=2))
\end{verbatim}
\endgroup

\clearpage
\section*{Version V2 append-only record}
\addcontentsline{toc}{section}{Version V2 append-only record}

\begin{description}[leftmargin=3.8cm,style=nextline]
\item[Version V2 -- 2 September 2026.]
Preserved V1 verbatim before its terminal document-closing command.  Retained
explicitly the Cauchy hypothesis needed by its cofinal summary.  Evaluated the
literal conditional-resampling spectrum and its native detail reserve.  Proved a source-level finite-lag obstruction to fixed-rate thermal
identification and a quantitative increasing-rate approximation theorem.
Went upstream of another unevaluated horizon inverse to the actual right
thermal dressing of the primitive columns and the Duhamel endpoint of the quantum half-source.

Proved that a fixed-wall Dirichlet normal action jet gives a right-domain
mismatch, survives every finite complete thermally smoothed primitive bank,
and bounds the fibre birth rank from below.  Evaluated a predeclared scalar
interval diagnostic: exact insertion table, five matrices, physical profile
metric, source reserve, bounded two-mode annihilator, total residual interval,
raw delayed kernels, and complete positive tails.  Proved persistence and
convergence on its declared energy regulator without calling that regulator a
spacetime continuum.  Supplied the complete exact-arithmetic replay and isolated
the native ESM boundary/domain incidence as the next actual calculation.
\end{description}

The complete V1 file has SHA--256
\begin{center}\small\ttfamily
3c9f497796f5f85e492b743bd418bfef1004a68b3995cfda41f19a4d1a3e98cc.
\end{center}

\begin{firewall}[Append-only instruction after V2]
\label{fire:V2-append}
V3 must preserve this complete V2 source verbatim, append immediately before
the sole final document-closing command, and use
\begin{center}\small
\path{Renewal_Einstein_Standard_Model_Physical_Source_Metric_Duhamel_Ancestry_and_Quantum_Continuum_Programme_V3.tex}.
\end{center}
Do not identify the bounded finite-rate RP generator with the full thermal
clock; use the actual quantum operator and its own calibrated source packet.
Do not repeat the V1 squeeze, the V15 five-matrix factorization, or the control
calculation.  Evaluate the native fixed-wall boundary/domain incidence or an
actual ESM fibre card.  A failure of the Dirichlet hypotheses is not a negative
physical verdict.  A positive diagnostic residue is not evidence of an
interacting continuum, and an energy-truncation route is not a physical
spacetime-cutoff route.
\end{firewall}

% =============================================================================
% END APPENDED VERSION V2 MATERIAL
% =============================================================================


% =============================================================================
% BEGIN APPENDED VERSION V3 MATERIAL
% =============================================================================
\clearpage
\hypersetup{pdftitle={Renewal Einstein--Standard--Model Physical Source Metric, Duhamel Ancestry and Quantum Continuum Programme V3},pdfsubject={Native overlap-wall motion, complete material action, and a covariant boundary-jet obstruction}}
\providecommand{\Rea}{\operatorname{Re}}
\providecommand{\diver}{\operatorname{div}}
\providecommand{\Jmat}{\mathcal J^{\rm mat}_{\partial,\lambda}}
\providecommand{\Rreal}{R^{\mathbb R}_{\lambda}}

\section*{Version V3: native overlap-wall motion, material action, and the first thermal boundary jet}
\addcontentsline{toc}{section}{Version V3: native wall motion and the material-action boundary jet}

\begin{abstract}
V2 found a fixed-wall Dirichlet obstruction and evaluated one diagnostic
Hamiltonian.  The present continuation returns to the actual domain specified
in the ESM corpus: the Friedrichs--Dirichlet realization on a Wilson-gap
sector.  The existing simple-crossing and Pauli theorems are imported rather
than reproved.  Their finite Wilson derivatives determine the normal velocity
of a regular sector wall.  An exact cancellation shows that this velocity can
be calculated from the zero-energy Feshbach matrix without identifying its
nonzero eigenvalues with those of the full Wilson operator.

The first evaluated matrix identity completes the previously calculated Pauli
polarization tensor.  The two signed wall-defining derivatives supply exactly
the missing diagonal part.  On one real Pauli source jet, their positive Gram
plus the real transition Gram is the full Pauli metric, pointwise in the
angular direction.  In particular, a scalar Pauli derivative is invisible to
polarization but not to motion of the sector domain.  The identity is a finite
normal-jet calculation, not a numerical value for the physical source reserve.

For a moving Dirichlet domain the common-action derivative must include the
material kinetic variation.  We derive its volume form, its exact finite
spectral Hadamard matrix, and its Gibbs boundary pressure.  The isolated
boundary pressure form is not bounded on the energy domain; it must not be
appended as a separate primitive after discarding the cancelling transport
term.  The complete material derivative is the admissible form source.

The main new thermal result evaluates the boundary trace of that form source.
For the predeclared normal, half-density collar transport, all mean-curvature
terms cancel.  With fixed principal metric and a varying unitary connection,
the resulting jet is
\[
 \kappa\bigl((\Delta_{\partial}v_u)f_a
      +2\nabla_{\partial}v_u\cdot\nabla^{\mathcal A}_{\partial}f_a
      -2\dot{\mathcal A}_{\nu,u}f_a\bigr),
 \qquad f_a=\nabla_\nu\psi_a.
\]
A nonzero total jet is incompatible with a complete finite thermally dressed
primitive fibre.  This conclusion remains valid when the unsmeared global
source is only an energy-dual form: the Duhamel divided difference improves
its spectral order by exactly one.  For real source coefficients, a weighted
boundary divergence identity prevents a unitary connection variation from
cancelling a nonconstant normal velocity under the stated boundary conditions.

All physical coefficient norms, Hilbert transports, prior primitives, and
cutoff routes remain fixed.  The normal-velocity and material-jet formulas are
now explicit; their numerical ESM values, a uniform residual floor, and the
cofinal ancestry packet are not supplied by these identities.  No second
scalar diagnostic or downstream continuum claim is introduced.
\end{abstract}

\section{The native domain is specified more precisely than V2 used}
\label{sec:V3-native-domain}

The main ESM manuscript \cite{ESMMainV37}, V13, explicitly places the finite
unquenched Hamiltonian on the Friedrichs--Dirichlet realization at the overlap
singular divisor.  Its V18 simple-crossing theorem writes the negative Wilson
projector locally as a smooth background plus one smooth eigenline whose
occupancy changes across a simple zero.  Thus its one-sided projector and
finite line coefficients can be smooth at a regular wall even though the
smallest absolute Wilson eigenvalue vanishes there.  The V19--V20 Feshbach,
Pauli, gauge-quotient, and transition-tensor calculations are retained.

The statement in V2 that the word ``hard-wall'' alone is insufficient remains
correct.  The source audit now supplies more than that word.  It does not,
however, supply a single globally smooth wall, a cutoff-uniform collar, the
particular electromagnetic wall velocity, or the physical Hilbert-space
transport used when a command moves the divisor.  Those are not inferred from
the local spectral splitting.  In particular, the published simple-crossing
excision concerns the inverse-sum-gap response tail; it does not delete the
Dirichlet boundary from the quantum Hamiltonian.

\paragraph{Fixed source and metric conventions.}
The actual source coefficient space is denoted by $E$, with its independently
fixed source-action form $\Msrc$.  The original target metric $\Ntar$ is also
retained.  A parameter $s$ is a command parameter, not physical time.  The
configuration-space Riemannian metric is denoted by $g$ and is not $\Msrc$.
The Hermitian Wilson matrix is $W_s(x)$; the quantum boson--fermion Hamiltonian
is $\hD_s$.  None of the finite eigenvalues of $W_s$ is identified with a
quantum energy of $\hD_s$.  Physical exact-gauge columns are retained in the
Ward identities and quotiented only where the existing programme permits it.

\paragraph{Nonduplication.}
This continuation does not reprove V1's squeeze, V15's five-matrix short,
V18's simple-crossing regularity, or V20's Pauli transition tensor.  It
calculates the missing domain-motion contribution and the boundary trace of
the complete quantum action source.  Smooth chart assumptions below describe
an explicit regular-wall branch.  Multiple-zero tips, changing orbit types,
nonregular walls, and principal-metric variations are returned separately.

\section{Normal velocity from the actual Wilson matrix and its Schur reduction}
\label{sec:V3-wall-velocity}

Let $\mathscr X$ be a smooth configuration chart with its fixed metric $g$.
Suppose $W_s(x)$ is a $C^2$ Hermitian finite matrix family.  Near a regular
wall point, let $\lambda_c(s,x)$ be its simple crossing eigenvalue and
$E_c(s,x)$ the corresponding rank-one orthogonal projection.  Choose
$\epsilon\in\{1,-1\}$ so that
\[
 d_s(x)=\epsilon\lambda_c(s,x)>0
\]
is the chosen side of the sector.  Assume $\nabla_gd_0\ne0$ on the wall patch.
Its outward unit normal is $\nu=-\nabla_gd_0/|\nabla_gd_0|_g$.

\begin{theorem}[Native simple-wall velocity]
\label{thm:V3-Wilson-velocity}
For a physical command $u$, let $K_u(x)=\partial_sW_s(x)|_{s=0}$ on its
specified command path.  Every differentiable transport $\Phi_s$ carrying
the old wall to $d_s=0$ has the same normal velocity
\begin{equation}
 \boxed{
 v_u=X_u\cdot\nu
 =\frac{\epsilon\Tr(E_cK_u)}{|\nabla_g\lambda_c|_g},
 \qquad X_u=\partial_s\Phi_s|_0.}
 \label{eq:V3-Wilson-velocity}
\end{equation}
The numerator and the differential of $\lambda_c$ are computable from the
same finite Wilson matrix and its first derivative bank.  Replacing $d_s$ by
$a_s d_s$ with $a_s>0$ leaves \eqref{eq:V3-Wilson-velocity} unchanged on the
wall.  A pure matrix-conjugation command $K_u=[G_u,W]$, $G_u^*=-G_u$, has
zero normal velocity.
\end{theorem}

\begin{proof}
For a locally normalized eigenvector $\chi_s$, differentiate
$W_s\chi_s=\lambda_c\chi_s$ and pair with $\chi_s$.  The two differentiated
eigenvector terms cancel, giving
$\partial_s\lambda_c=\chi^*K_u\chi=\Tr(E_cK_u)$.  The same calculation in
every configuration direction gives $\dd_x\lambda_c$.
Differentiating $d_s(\Phi_s(x))=0$ on the old wall gives
$\dot d_u+\dd d_0(X_u)=0$.  Since
$\dd d_0(X_u)=-|\nabla_gd_0|_g v_u$, the asserted sign and denominator
follow.  On $d_0=0$, both derivatives of $a_sd_s$ acquire the same positive
factor $a_0$.  Finally, $W\chi=0$ at the wall, so
$\chi^*[G_u,W]\chi=0$.
\end{proof}

\begin{remark}[What zero velocity means]
Zero velocity fixes the wall to first order.  It does not make the entire
quantum action derivative zero, prove a line-current identity, or close the
Duhamel residual.  Conversely, a nonzero velocity is a domain-incidence
value, not already a new source birth.
\end{remark}

For computation, retain the exact block decomposition
\begin{equation}
 W=\begin{pmatrix} A&B\\B^*&C\end{pmatrix},\qquad
 C\text{ invertible},\qquad
 F=A-BC^{-1}B^*,\qquad
 L=\binom{I}{-C^{-1}B^*}.
 \label{eq:V3-Schur-wall}
\end{equation}
These are the zero-energy Feshbach objects of the inherited programme.

\begin{theorem}[Schur normal velocity without spectral substitution]
\label{thm:V3-Schur-velocity}
Suppose $e$ is a unit vector spanning a simple kernel of $F$ at a wall
point.  Then $\chi=Le/\|Le\|$ spans $\ker W$.  For every configuration or
command derivative $\partial$, let $f$ be the simple zero eigenvalue branch
of $F$.  At the wall,
\begin{equation}
 \boxed{
 \partial\lambda_c
 =\frac{e^*(\partial F)e}{\|Le\|^2}
 =\frac{\partial f}{\|Le\|^2}.}
 \label{eq:V3-Schur-derivative}
\end{equation}
Consequently the normal velocity computed with the signed defining function
$\epsilon f$ equals \eqref{eq:V3-Wilson-velocity}.  The factor $\|Le\|^{-2}$
cancels between its numerator and physical normal denominator.

No equality of the nonzero eigenvalues of $F$ and $W$ is used or asserted.
In particular, a Feshbach Pauli matrix may be used to calculate a regular
zero wall without being substituted for the quantum clock, or even for the
full Wilson spectrum away from zero.
\end{theorem}

\begin{proof}
Block multiplication gives $WL=(F,0)^{\mathsf T}$ and $L^*WL=F$.
Since $Fe=0$, differentiation of the second identity yields
\[
 e^*(\partial F)e=(Le)^*(\partial W)(Le);
\]
the terms containing $\partial L$ vanish because $WLe=0$.  Divide by
$\|Le\|^2$ and apply the simple-eigenvalue differentiation already proved.
The spatial differential is multiplied by the same positive factor, so its
$g$-dual norm is multiplied by that factor as well.  Inertia preservation
under the Schur congruence chooses the same sector side locally.  This proves
the velocity statement without an isospectral claim.
\end{proof}

\section{An evaluated Pauli wall--polarization completion}
\label{sec:V3-Pauli-completion}

Use the inherited Pauli coordinates for a two-dimensional exact Schur head:
\[
 F=\tau I+\mathbf z\cdot\boldsymbol\sigma,\qquad
 r=|\mathbf z|>0,\quad n=\mathbf z/r,\qquad
 K(u)=D_uF=\alpha(u)I+\mathbf k(u)\cdot\boldsymbol\sigma.
\]
The indefinite chamber is $-r<\tau<r$.  Its two positive defining functions
are $d_+=r-\tau$ and $d_-=r+\tau$.  Introduce their calibrated differential
numerators
\begin{equation}
 b_{+,n}(u)=\frac{n\cdot\mathbf k(u)-\alpha(u)}{\sqrt2},\qquad
 b_{-,n}(u)=\frac{n\cdot\mathbf k(u)+\alpha(u)}{\sqrt2}.
 \label{eq:V3-Pauli-b}
\end{equation}
At the respective physical wall these give the actual speed through
\begin{equation}
 v_{\pm,u}=\frac{\sqrt2\,b_{\pm,n}(u)}{\eta_\pm},\qquad
 \eta_\pm=|\nabla_g d_\pm|_g.
 \label{eq:V3-Pauli-speed}
\end{equation}
The physical $\eta_\pm$ is retained; setting it to $\sqrt2$ would require an
independently verified Euclidean Pauli normal metric.  The $b_\pm$ are
normal-jet coordinates, not freely whitened physical speeds.

Let $p_\pm(n)=(I\pm n\cdot\boldsymbol\sigma)/2$ and
$\kappa_n(u)=p_-(n)K(u)p_+(n)$.  The real part of its mixed Gram has already
been computed in main V20 and is imported here.

\begin{theorem}[Pointwise Pauli wall--polarization Pythagoras]
\label{thm:V3-Pauli-Pythagoras}
For real source directions $u,w$,
\begin{equation}
 \boxed{
 \begin{aligned}
 &b_{+,n}(u)b_{+,n}(w)+b_{-,n}(u)b_{-,n}(w)
       +\Rea\Tr\bigl(\kappa_n(u)^*\kappa_n(w)\bigr)\\
 &\hspace{2cm}=\alpha(u)\alpha(w)+\mathbf k(u)\cdot\mathbf k(w)
       =\tfrac12\Tr(K(u)K(w)).
 \end{aligned}}
 \label{eq:V3-Pauli-Pythagoras}
\end{equation}
Thus the two wall-defining derivatives together with the polarization
transition have exactly the full real Pauli source metric, with no angular
null direction.  A scalar source $K(u)=\alpha(u)I$ has zero polarization
transition and wall numerator pair
$(-\alpha(u),\alpha(u))/\sqrt2$.

On a frozen affine normal jet, normalized angular integration in the ordered
Pauli basis $(I,\sigma_1,\sigma_2,\sigma_3)$ gives the exact matrices
\begin{equation}
 \boxed{
 \mathsf Q_{\rm wall}=\diag(1,\tfrac13,\tfrac13,\tfrac13),\qquad
 \mathsf Q_{\rm pol}=\diag(0,\tfrac23,\tfrac23,\tfrac23),\qquad
 \mathsf Q_{\rm wall}+\mathsf Q_{\rm pol}=I_4.}
 \label{eq:V3-Pauli-matrices}
\end{equation}
The normalization is $(4\pi)^{-1}\int_{S^2}\dd\Omega$, not an assumption
about a physical Gibbs angular law.
\end{theorem}

\begin{proof}
The sum of the first two products in \eqref{eq:V3-Pauli-Pythagoras} is
$\alpha(u)\alpha(w)+(n\cdot\mathbf k(u))(n\cdot\mathbf k(w))$.
The imported transition Gram is
$\mathbf k(u)\cdot\mathbf k(w)
 -(n\cdot\mathbf k(u))(n\cdot\mathbf k(w))$.
Addition proves the pointwise identity.  The Pauli trace identity gives its
last expression.  Integrating $n_in_j$ gives $\delta_{ij}/3$ and proves
\eqref{eq:V3-Pauli-matrices}.  No derivative of a fitted eigenframe occurs.
\end{proof}

\begin{firewall}[What this finite matrix evaluation does not replace]
\label{fire:V3-Pauli-metric}
Equation~\eqref{eq:V3-Pauli-Pythagoras} concerns one common differential
$K$; it does not identify different physical source derivatives at two
separated sheet points.  At a given regular wall only its active signed
velocity is a wall velocity there.  The two-sheet affine germ uses the same
frozen $K$ by its stated definition.  For a nonlinear physical family the
actual $D_uF$, metric, densities, and mixed rows must be evaluated at their
actual configurations.

These are source-jet metrics, not $B\Msrc^\dagger B^*$, an entrance reserve,
a Gibbs boundary pressure, or a dynamic-birth Gram.  In particular, main
V20's removal of $\alpha I$ from the cross-gap capacity is retained, but it
cannot be extended to removal of $\alpha I$ from the domain-motion audit.
At the multiple-zero tip $r=0$ the individual sheets and $n$ are unresolved;
no simple-wall formula is continued through that point.
\end{firewall}

\section{The complete material action on a moving Dirichlet sector}
\label{sec:V3-material-action}

First work on a bounded smooth domain $\Omega$ with a smooth finite-rank
Hermitian bundle.  The global smooth compact hypothesis in this section
makes the action and spectral formulae classical operator expressions.  A
local version needed for singular global ESM cards is proved below.
Let
\begin{equation}
 \hD_s=-\kappa\Delta_{\mathcal A_s}+\mathcal V_s\ge0
 \quad\hbox{on }\Omega_s=\Phi_s(\Omega),
 \label{eq:V3-Laplace-family}
\end{equation}
with Dirichlet boundary, fixed $\kappa>0$, fixed ambient principal metric,
unitary connection $\mathcal A_s$, and Hermitian potential $\mathcal V_s$.
Assume smooth coefficients, smooth $\Phi_s$ and their parameter derivatives,
and the common $H^2\cap H_0^1$ operator realization after pullback.  Scalar
energy shifts, including normalization conventions, may be included in
$\mathcal V_s$.  The original thermal and Duhamel admissibility hypotheses
are retained.

The Hilbert pullback $U_s$ includes the positive Jacobian square root and the
baseline-connection parallel lift along the specified configuration flow.
Any additional parameter-dependent unitary lift is retained in the physical
Hilbert connection of \Cref{sec:V3-router-covariance}.  Write
\[
 \widetilde h_s=U_s\hD_sU_s^*,\qquad
 a_u=\partial_s\widetilde h_s|_0,
 \qquad
 D_X=\nabla_X^{\mathcal A}+\tfrac12\diver_gX.
\]
Here $a_u$ means the complete derivative as an energy form.  On the smooth
interior core it is
\begin{equation}
 \boxed{a_u=a_u^{\rm int}+[D_{X_u},\hD],}
 \label{eq:V3-material-commutator}
\end{equation}
where $a_u^{\rm int}$ differentiates all fixed-coordinate coefficients.
In particular the connection variation is included in $a_u^{\rm int}$.

\begin{theorem}[Volume form and material kinetic contribution]
\label{thm:V3-material-form}
In a flat chart with trivial bundle connection, put
$\hD_s=-\kappa\Delta+\mathcal V_s$ and
$\theta=\diver X$.  The exact pulled-back form derivative is
\begin{equation}
 \begin{aligned}
 a_u[\phi,\psi]
 ={}&-\kappa\int_\Omega
    (\nabla\phi)^*(DX+DX^{\mathsf T})\nabla\psi\,dx\\
 &-\frac\kappa2\int_\Omega\left[
       \phi^*\nabla\theta\cdot\nabla\psi
       +(\nabla\phi)^*\cdot\nabla\theta\,\psi\right]dx\\
 &+\int_\Omega\phi^*(\dot{\mathcal V}_u+X\cdot\nabla\mathcal V)\psi\,dx.
 \end{aligned}
 \label{eq:V3-material-volume}
 \end{equation}
It is a continuous form on $H_0^1(\Omega)$, and agrees with
\eqref{eq:V3-material-commutator} on compactly supported smooth sections.
Smooth connection and metric coefficients give the analogous bounded
energy-form derivative by differentiating their pulled-back quadratic form.
A derivative of the multiplication potential alone is therefore not the
complete source of a moving wall.
\end{theorem}

\begin{proof}
Write $J_s=\det D\Phi_s$.  In the scalar flat chart the pulled-back kinetic
form is
\[
 \kappa\int_\Omega J_s
 \left(D\Phi_s^{-\mathsf T}\nabla(J_s^{-1/2}\phi)\right)^*
 \left(D\Phi_s^{-\mathsf T}\nabla(J_s^{-1/2}\psi)\right)dx.
\]
At zero, $\dot J=\theta$, $\partial_sJ_s^{-1/2}=-\theta/2$,
and $\partial_sD\Phi_s^{-\mathsf T}=-(DX)^{\mathsf T}$.  The zeroth-order
Jacobian contributions to the product of the two gradients cancel; its
remaining derivative is the first two lines of
\eqref{eq:V3-material-volume}.  The potential becomes
$\mathcal V_s(\Phi_s(x))$, giving the last line.  Cauchy--Schwarz and bounded
smooth coefficients bound each term by
$C\|\phi\|_{H^1}\|\psi\|_{H^1}$.  Direct differentiation on the interior
core gives $\dot{\mathcal V}+[X\cdot\nabla+\theta/2,\hD]$.
For a smooth bundle and principal metric the same change of variables is
performed on the covariant energy form.  Every differentiated coefficient
multiplies at most one derivative of each test section; bounded coefficients
and uniform ellipticity give the same energy-form continuity.
\end{proof}

\begin{remark}[Source cost is not set by this estimate]
The constant in the energy-form bound verifies admissibility.  It does not
redefine $\Msrc$ or establish a lower source reserve.  On an unbounded Higgs
carrier, the corresponding global polynomial and energy-form bounds must be
retained from the original action, not replaced by a bounded-domain estimate.
\end{remark}

Let $(\psi_j,\lambda_j)$ be an orthonormal Dirichlet eigensystem at the
base card, and set $f_j=\nabla_\nu^{\mathcal A}\psi_j$.

\begin{theorem}[Complete finite spectral Hadamard matrix]
\label{thm:V3-Hadamard-matrix}
For the preceding material derivative,
\begin{equation}
 \boxed{
 \langle\psi_m,a_u\psi_n\rangle
 =a_{u,mn}^{\rm int}
  +(\lambda_n-\lambda_m)\langle\psi_m,D_{X_u}\psi_n\rangle
  -\kappa\int_{\partial\Omega}v_u f_m^*f_n\,dS.}
 \label{eq:V3-Hadamard-matrix}
\end{equation}
In a degenerate eigenspace $E_\lambda$, the transport-energy difference
vanishes.  Its first-order perturbation matrix is exactly the intrinsic
coefficient block minus the displayed boundary-pressure block.  All blocks
are calculated on the same eigensystem and material transport.
\end{theorem}

\begin{proof}
The local identity \eqref{eq:V3-material-commutator} gives
$[D_X,\hD]\psi_n=(\lambda_n-\hD)D_X\psi_n$ in the interior.
Green's identity, with $\psi_m=0$ on the boundary, gives
\[
 \langle\psi_m,\hD\zeta\rangle
 =\lambda_m\langle\psi_m,\zeta\rangle
   +\kappa\int_{\partial\Omega}f_m^*\zeta\,dS
\]
for smooth $\zeta$, without imposing a Dirichlet condition on $\zeta$.
Take $\zeta=D_X\psi_n$.  Its boundary trace is $v_u f_n$, because the
half-density term and the tangential derivatives of $\psi_n$ vanish there.
This proves \eqref{eq:V3-Hadamard-matrix}.  The skew-adjoint identity for
$D_X$ between Dirichlet test functions shows that the middle matrix, including
its energy difference, is Hermitian.  Restriction to $\lambda_m=\lambda_n$
proves the degenerate-block assertion.  The usual first-order eigenvalue
interpretation follows by differentiating the finite spectral projection
block, or is unnecessary when only that block is retained as a form.
\end{proof}

\begin{remark}[Classical formula versus the new use]
The diagonal Dirichlet formula is the classical Hadamard variation.  For
background on spectral boundary variations see P.~Grinfeld,
\emph{Hadamard's Formula Inside and Out}, J. Optim. Theory Appl.
\textbf{146} (2010), 654--690,
\href{https://doi.org/10.1007/s10957-010-9681-6}{doi:10.1007/s10957-010-9681-6}.
The complete material-source boundary trace and its thermal ancestry use below
are proved here.  No result about nonsmooth or multiple-zero ESM walls is
imported from the smooth-domain reference.
\end{remark}

\begin{proposition}[Thermal boundary pressure with its signed incidence]
\label{prop:V3-thermal-pressure}
Under the smooth compact hypotheses and differentiable heat trace,
$Z_\beta=\Tr e^{-\beta\hD}$ satisfies
\begin{equation}
 \boxed{
 D_u\log Z_\beta
 =-\beta\Tr(\rho_\beta a_u^{\rm int})
  +\beta\int_{\partial\Omega}v_u\,p_\beta\,dS,\qquad
 p_\beta=\frac\kappa{Z_\beta}
       \sum_j e^{-\beta\lambda_j}|f_j|^2\ge0.}
 \label{eq:V3-Gibbs-pressure}
\end{equation}
A sign-changing velocity is paired with this positive density before any
positive part or source Gram is taken.  Two moving sheets can cancel in this
scalar expectation even when their velocity Gram is positive.
\end{proposition}

\begin{proof}
Differentiation of the trace gives
$D_uZ_\beta=-\beta\sum_j e^{-\beta\lambda_j}a_{u,jj}$; within a repeated
eigenvalue the sum is the trace of its perturbation matrix.  Insert
\eqref{eq:V3-Hadamard-matrix}.  Standard smooth elliptic estimates give a
polynomial bound on $\|f_j\|_{L^2(\partial\Omega)}$ in $1+\lambda_j$;
it also follows by the $H^2$ Dirichlet estimate and the continuous normal
trace from $H^2$.  Trace-class heat at every positive time therefore makes
the boundary series absolutely summable in $L^1$.  Interchange sum and
boundary integral, divide by $Z_\beta$, and obtain
\eqref{eq:V3-Gibbs-pressure}.  Positivity of $p_\beta$ is termwise.  The
last assertion is the signed linearity of the displayed formula.
\end{proof}

\begin{theorem}[The isolated wall-pressure form is not an energy-form primitive]
\label{thm:V3-boundary-not-form}
Suppose $v$ is real and continuous and is nonzero on an open regular boundary
patch.  The form
\begin{equation}
 \mathfrak b_v[\phi,\psi]
 =\kappa\int_{\partial\Omega}v
      (\nabla_\nu\phi)^*\nabla_\nu\psi\,dS
 \label{eq:V3-isolated-boundary}
\end{equation}
on smooth Dirichlet sections has no continuous extension to the $H_0^1$
energy domain.  Thus the boundary term in
\eqref{eq:V3-Hadamard-matrix} cannot be detached from the material action
and declared an additional form-bounded source.  The complete derivative in
\eqref{eq:V3-material-volume} is nevertheless form bounded.
\end{theorem}

\begin{proof}
Shrink the patch until $v$ has one sign and $|v|\ge v_0>0$.  In a smooth
inward normal collar choose a nonzero smooth tangential section $\eta$ with
support in that patch and a smooth cutoff $\chi$ equal to one near zero and
zero for $r\ge2$.  Put
\[
 \psi_\varepsilon(z,r)
   =\varepsilon^{-1/2}r\chi(r/\varepsilon)\eta(z).
\]
Every $\psi_\varepsilon$ is smooth, has zero boundary trace, and lies in the
operator domain.  Its normal energy is $O(1)$, its tangential energy is
$O(\varepsilon^2)$, and its $L^2$ norm squared is $O(\varepsilon^2)$.
Smooth collar densities and bounded connection coefficients preserve these
bounds.  But $\nabla_\nu\psi_\varepsilon|_{\partial}
=-\varepsilon^{-1/2}\eta$, so
\[
 |\mathfrak b_v[\psi_\varepsilon,\psi_\varepsilon]|
 \ge\kappa v_0\varepsilon^{-1}\|\eta\|_{L^2(\partial)}^2.
\]
This contradicts energy-form boundedness.  The complete material derivative
is already bounded by \Cref{thm:V3-material-form}; the other spectral term in
\eqref{eq:V3-Hadamard-matrix} therefore cannot be discarded when forming its
source.  The construction also works locally inside an otherwise singular
global card whenever its form has the stated regular collar.
\end{proof}

\section{Duhamel action rows when only an energy form is available}
\label{sec:V3-form-Duhamel}

This section uses the original Hilbert space, not a new source carrier.  It
extends V2's integer-domain test to the already admitted energy-dual action
rows.  Let $S=1+\hD$, and let a Hermitian source form $a_u$ satisfy
\begin{equation}
 \widehat a_u=S^{-1/2}a_uS^{-1/2}\in\mathcal B(\mathcal H).
 \label{eq:V3-form-source}
\end{equation}
Its centered form is $a_u^\circ=a_u-\Tr(\rho a_u)I$.  The thermal expectation
exists: its absolute diagonal series is bounded by
$\|\widehat a_u\|\Tr(S e^{-2b\hD})/Z$.  For a finite left eigenspace,
$G_\lambda(u)=E_\lambda a_u^\circ$ is generally an energy-dual row.  Explicitly,
\begin{equation}
 G_\lambda(u)S^{-1/2}
 =(1+\lambda)^{1/2}E_\lambda\widehat a_u^\circ
 \quad\hbox{is Hilbert--Schmidt}.
 \label{eq:V3-dual-row}
\end{equation}
The notation $H_R$ and the divided difference $d_\lambda$ are those of V2.
For a distributional row the statement ``spectral order $s$'' means that its
coefficients have finite squared norm with weight $(1+x)^{2s}$.

\begin{lemma}[Exact one-order Duhamel isomorphism]
\label{lem:V3-Duhamel-orders}
For every fixed $b>0$ and finite $\lambda\ge0$ there are constants
$0<c_{b,\lambda}\le C_{b,\lambda}<\infty$ such that
\begin{equation}
 c_{b,\lambda}\le (1+x)|d_\lambda(x)|\le C_{b,\lambda}
 \quad(x\ge0).
 \label{eq:V3-dd-two-sided}
\end{equation}
Consequently multiplication by $d_\lambda(H_R)$ maps rows of spectral order
$s$ bijectively to rows of order $s+1$.  In particular the form row
\eqref{eq:V3-dual-row} gives a well-defined target
\begin{equation}
 Y_\lambda(u)=Z^{-1/2}d_\lambda(H_R)G_\lambda(u)
 \quad\hbox{of order at least }\tfrac12.
 \label{eq:V3-form-target}
\end{equation}
If $Y_\lambda(u)$ belongs to every right power domain, then $G_\lambda(u)$
belongs to every right power domain as well.
\end{lemma}

\begin{proof}
The divided difference of the strictly decreasing function $e^{-bx}$ is
strictly negative, including its removable diagonal value.  Hence
$(1+x)|d_\lambda(x)|$ is positive and continuous on $[0,\infty)$, and its
limit at infinity is $e^{-b\lambda}>0$.  Compactness of a finite interval
and that positive limit prove both bounds.  Multiply the weighted squared
spectral coefficients by these pointwise inequalities.  Their inverse gives
the bijection, also on the energy-dual completion.  Taking $s=-1/2$ gives
\eqref{eq:V3-form-target}; taking successively larger $s$ proves the last
assertion.
\end{proof}

\begin{proposition}[The complete form-Duhamel synthesis is Hilbert--Schmidt]
\label{prop:V3-form-HS}
For one source direction, the matrix defined by
\[
 Y_{mn}=Z^{-1/2}d_b(\lambda_m,\lambda_n)(a_u^\circ)_{mn}
\]
is Hilbert--Schmidt, where $d_b$ is the symmetric divided difference of
$e^{-bx}$.  It satisfies
\begin{equation}
 \boxed{
 \|Y\|_{\HS}^2
 \le\frac{8(1+b^2)}Z\|\widehat a_u^\circ\|^2
       \Tr\bigl(S^2e^{-2b\hD}\bigr)<\infty.}
 \label{eq:V3-form-HS}
\end{equation}
Whenever the inherited physical Duhamel identification holds, this synthesis
is the covariant derivative of the normalized Gibbs square root.  The bound
alone does not identify a proposed parameter transport with the physical one.
\end{proposition}

\begin{proof}
Put $m=\min(x,y)$ and $M=\max(x,y)$.  If $M\le2m+1$, the mean-value
estimate $|d_b(x,y)|\le b e^{-bm}$ gives
\[
 (1+x)(1+y)|d_b(x,y)|^2
 \le4b^2(1+m)^2e^{-2bm}.
\]
If $M>2m+1$, then $M-m>(M+1)/2$ and
$|d_b(x,y)|\le2e^{-bm}/(M+1)$, giving an upper bound
$4e^{-2bm}$.  Thus in both cases the expression is at most
$4(1+b^2)(1+m)^2e^{-2bm}$.
Now $(a_u^\circ)_{mn}=\sqrt{1+\lambda_m}\,
(\widehat a_u^\circ)_{mn}\sqrt{1+\lambda_n}$.
On the half-plane $\lambda_n\ge\lambda_m$, sum first in $n$ and use
$\sum_n|(\widehat a_u^\circ)_{mn}|^2\le\|\widehat a_u^\circ\|^2$.
Use column sums on the other half-plane.  This proves
\eqref{eq:V3-form-HS}, possibly counting the diagonal twice.  The trace is
finite because a polynomial times $e^{-2bx}$ is bounded by a constant times
$e^{-bx}$.  The final identification is exactly the inherited physical
Duhamel premise, now with the complete material form rather than a truncated
potential derivative.
\end{proof}

\begin{theorem}[A local material trace defect survives the complete primitive short]
\label{thm:V3-local-trace-birth}
Let the global Hamiltonian be the inherited nonnegative Dirichlet realization
with compact resolvent and trace-class heat at every positive time.  It may
have other singular walls.  On one regular open boundary patch $\Gamma$,
suppose the selected eigenvectors and the complete material form's local
differential expression are smooth up to the boundary.  Suppose every
$\Dom\hD$ vector has zero trace on $\Gamma$, and the complete prior finite
primitive bank satisfies V2's form-dressed hypothesis.  Define
\begin{equation}
 \Jmat u
 =\bigl(\gamma_\Gamma(a_u\psi_a)\bigr)_{a=1}^{\rank E_\lambda}.
 \label{eq:V3-material-jet-definition}
\end{equation}
Centering $a_u$ does not change this jet.  For every real physical direction,
\begin{equation}
 \boxed{
 \Jmat u\ne0
 \quad\Longrightarrow\quad
 \|(I-P_{B_\lambda})Y_\lambda(u)\|^2>0.}
 \label{eq:V3-local-positive}
\end{equation}
Equivalently, on the real coefficient space,
\begin{equation}
 \boxed{\ker\Rreal\subseteq\ker\Jmat,\qquad
  \rank_{\R}\Rreal\ge\rank_{\R}\Jmat,}
 \label{eq:V3-material-rank}
\end{equation}
where $\Rreal=\Rea\bigl(Y_\lambda^*(I-P_{B_\lambda})Y_\lambda\bigr)$.
If the global unsmeared row is in $L^2$, then a nonzero jet gives the sharper
statement $Y_\lambda(u)\in\Dom H_R\setminus\Dom H_R^2$.
\end{theorem}

\begin{proof}
For real $u$, the adjoint-vector representatives of $G_\lambda(u)$ are the
form distributions $a_u^\circ\psi_a$.  If such a distribution belonged to
$\Dom\hD$, its restriction to $\Gamma$ would have zero trace.  Equality
with its smooth local differential representative forces that same trace to
be the corresponding component of \eqref{eq:V3-material-jet-definition}.
Thus a nonzero jet implies $G_\lambda(u)\notin\Dom H_R$, regardless of
whether it is globally $L^2$.

If $Y_\lambda(u)$ were in the finite primitive range, V2's right-smoothing
lemma would put it in every right power domain.
\Cref{lem:V3-Duhamel-orders} would then put $G_\lambda(u)$ in every right
power domain, a contradiction.  The finite primitive range is closed, so
the distance is strictly positive.  The kernel implication and real
rank--nullity prove \eqref{eq:V3-material-rank}.  A scalar centering term
has zero trace because every $\psi_a$ is Dirichlet.  Finally, a globally
$L^2$ row has spectral order zero, so its target has order one; failure of
order one in $G_\lambda$ is exactly failure of order two in $Y_\lambda$.
\end{proof}

\begin{remark}[Real rank and genuinely complex coefficients]
The physical command directions and unitary-connection variations in the next
theorem are real.  The rank claim is therefore deliberately for the real
symmetric form $\Rreal$.  It is not a claim that this real rank equals the
rank of an arbitrary complex Hermitian Gram.  In a real-preserving
realization the same result complexifies, but such a real structure is not
assumed for the full chiral ESM card.  The complex, source-normalized residual
itself remains the original V15 object.
\end{remark}

\section{The native normal-collar jet: curvature cancels, connection variation cannot hide normal motion}
\label{sec:V3-collar-jet}

Fix the Hilbert transport before inspecting any target.  On a regular patch
choose inward distance $r\ge0$ and boundary coordinates $z$, and write
\[
 g=dr^2+g_r(z),\qquad d\mu=J(z,r)\,dS(z)dr,\qquad
 a(z,r)=\partial_r\log J(z,r).
\]
Use normal parallel transport for the bundle, so $\mathcal A_r=0$.  The
baseline operator has the local expression
\begin{equation}
 \hD=-\kappa\bigl(\partial_r^2+a\partial_r+
                 \Delta_{g_r,\mathcal A}\bigr)+\mathcal V.
 \label{eq:V3-collar-h}
\end{equation}
The declared material velocity is the constant-normal collar extension
$X_u=-v_u(z)\partial_r$ near $r=0$, tapered only away from the boundary.
Its half-density generator there is
\begin{equation}
 D_{X_u}=-v_u\partial_r-\tfrac12v_ua.
 \label{eq:V3-collar-D}
\end{equation}
Allow the fixed-coordinate unitary connection to vary with derivative
$\dot{\mathcal A}_u$, and allow an arbitrary smooth Hermitian multiplication
variation of $\mathcal V$.  The ambient principal metric and $\kappa$ are
fixed in this theorem.  For real $u$, every
$\dot{\mathcal A}_{\nu,u}$ is skew-adjoint.

\begin{theorem}[Exact complete material boundary jet]
\label{thm:V3-collar-jet}
For $\hD\psi_a=\lambda\psi_a$, put
$f_a=\nabla_\nu\psi_a=-\partial_r\psi_a|_0$.  The complete material action
satisfies
\begin{equation}
 \boxed{
 (\Jmat u)_a
 =\kappa\left[
       (\Delta_{\partial}v_u)f_a
       +2\nabla_{\partial}v_u\cdot
                      \nabla^{\mathcal A}_{\partial}f_a
       -2\dot{\mathcal A}_{\nu,u}f_a\right].}
 \label{eq:V3-native-jet}
\end{equation}
All mean-curvature and normal-Jacobian derivatives cancel.  The smooth
multiplication-potential variation has zero first material boundary trace.
A principal-metric or leading kinetic variation, if physically present, has
its own additional trace and is not omitted from $\Jmat$.
\end{theorem}

\begin{proof}
Write $f=\partial_r\psi|_0=-f_a$, use primes for $r$ derivatives, and put
$\Delta_\partial=\Delta_{g_0,\mathcal A}$.  Since $\psi|_0=0$, all its
tangential derivatives vanish at the boundary.  The eigenvalue equation and
its first normal derivative give
\begin{align}
 \psi''|_0&=-af,\label{eq:V3-collar-jet2}\\
 \psi'''|_0&=(a^2-a')f-\Delta_\partial f
                 +\kappa^{-1}(\mathcal V-\lambda)f.
 \label{eq:V3-collar-jet3}
\end{align}
Indeed $\partial_r(\Delta_{g_r,\mathcal A}\psi)|_0
=\Delta_\partial f$: the derivative of its coefficients acts on tangential
jets of the zero boundary section and hence vanishes.

From \eqref{eq:V3-collar-D},
\[
 D_X\psi|_0=-vf,\quad
 (D_X\psi)'|_0=\tfrac12va f,\quad
 (D_X\psi)''|_0=-v\psi'''-va'f+\tfrac12va^2f.
\]
Insert these in \eqref{eq:V3-collar-h} and use
\eqref{eq:V3-collar-jet3}.  The $a^2$, $a'$, and potential terms cancel,
leaving
\[
 \hD D_X\psi|_0
  =\kappa\{\Delta_\partial(vf)-v\Delta_\partial f\}-v\lambda f.
\]
On the other hand $D_X\hD\psi|_0=-v\lambda f$.  Therefore
\[
 [D_X,\hD]\psi|_0
 =-\kappa\{(\Delta_\partial v)f
                +2\nabla_\partial v\cdot\nabla_\partial^{\mathcal A}f\}.
\]
Replace $f$ by $-f_a$.  Differentiating the covariant Laplacian at fixed
coordinates contributes
$-\kappa(2\dot{\mathcal A}^{\,j}_u\nabla_j+
\nabla^j\dot{\mathcal A}_{j,u})\psi$.  Its boundary trace is
$-2\kappa\dot{\mathcal A}_{\nu,u}f_a$.  The differentiated multiplication
potential has zero trace.  Adding the two contributions proves
\eqref{eq:V3-native-jet}.  A varied principal coefficient produces an
additional second-order differentiated expression, so it is outside this
fixed-principal calculation and must be included separately.
\end{proof}

Define the nonnegative boundary weight
\begin{equation}
 w_\lambda(z)=\sum_{a=1}^{\rank E_\lambda}|f_a(z)|^2.
 \label{eq:V3-boundary-weight}
\end{equation}
It is invariant under unitary changes of eigenbasis and bundle frame.

\begin{theorem}[Weighted divergence and real wall-motion rank]
\label{thm:V3-boundary-divergence}
On the preceding normal-collar branch and for real source coefficients,
\begin{equation}
 \boxed{
 \Rea\sum_a f_a^*(\Jmat u)_a
       =\kappa\diver_\partial(w_\lambda\nabla_\partial v_u).}
 \label{eq:V3-boundary-divergence}
\end{equation}
Suppose either:
\begin{enumerate}[label=\textup{(D\arabic*)},leftmargin=3em]
\item $\Gamma$ is connected, every $v_u$ has compact support in a proper
      open subpatch of $\Gamma$, and $w_\lambda>0$ on $\Gamma$; or
\item $\Gamma$ is a closed connected regular boundary component and
      $w_\lambda>0$ on it.
\end{enumerate}
In case \textup{(D1)},
\begin{equation}
 \boxed{
 \ker\Rreal\subseteq\ker(v:E\to C_c^\infty(\Gamma)),\qquad
 \rank_\R\Rreal\ge\rank_\R v.}
 \label{eq:V3-velocity-rank}
\end{equation}
In case \textup{(D2)}, the same statement holds for $v$ modulo the constant
functions.  On a family of closed components one takes constants on each
component separately.  In particular, a unitary connection variation cannot
cancel the entire first material jet of an admitted nonzero compactly
supported normal velocity.
\end{theorem}

\begin{proof}
Metric compatibility of the unitary connection gives
$\nabla w_\lambda=2\Rea\sum_a f_a^*\nabla^{\mathcal A}f_a$.
For real $u$, skew-adjointness gives
$\Rea f_a^*\dot{\mathcal A}_{\nu,u}f_a=0$.  Pair
\eqref{eq:V3-native-jet} with $f_a$, sum, and use these identities to obtain
\eqref{eq:V3-boundary-divergence}.

If $\Jmat u=0$, multiply the resulting divergence equation by the real
function $v_u$ and integrate.  Compact support in \textup{(D1)} or absence
of a boundary in \textup{(D2)} eliminates the integration boundary term:
\[
 0=-\int_\Gamma w_\lambda|\nabla_\partial v_u|^2\,dS.
\]
Thus $v_u$ is constant on the connected patch.  Compact support in a proper
open patch makes that constant zero in \textup{(D1)}; on a closed component
it is retained as the stated exceptional constant mode.  Combine this kernel
inclusion with \eqref{eq:V3-material-rank} and apply real rank--nullity.
If \textup{(D1)} is used on a compact boundary component itself, replace its
compact-support condition by support in a proper open subpatch, or impose
zero boundary trace on a patch with boundary; a nonzero constant may not be
erased by notation.
\end{proof}

\begin{corollary}[One occupied eigenvector can detect all admitted local wall directions]
\label{cor:V3-local-velocity-bank}
Suppose one selected eigenvector has a nonvanishing normal derivative on a
regular connected patch, all velocities in a finite predeclared real command
bank have compact support in a proper subpatch, and the full source and
transport satisfy the hypotheses above.  Then every direction with nonzero
velocity has a strictly positive full-primitive fibre residual, and
\eqref{eq:V3-velocity-rank} holds.  This is a finite-card strict-positivity
result; no eigenvalue or residual number is needed to prove that branch.
\end{corollary}

\begin{proof}
The single nonvanishing normal derivative makes $w_\lambda>0$ on that patch.
Apply \Cref{thm:V3-boundary-divergence,thm:V3-local-trace-birth}.
\end{proof}

\begin{firewall}[No new physical commands from a proof cutoff]
\label{fire:V3-local-support}
A compactly supported configuration-space wall velocity is an assumption on
an already occurring command, not permission to multiply a physical
spacetime source by an arbitrary configuration cutoff.  On an actual ESM
command without that support, test the total jet
\eqref{eq:V3-native-jet} directly or retain the boundary flux in the weighted
divergence identity.  Nonvanishing of $w_\lambda$, the chosen collar
transport, and the fixed-principal hypothesis are physical incidence rows.
They are not deduced from a Pauli coefficient count.  The local theorem can
be used in the presence of other singular global walls because it tests a
trace on one regular patch, but its global source form and physical Duhamel
identification are still required.
\end{firewall}

\section{Hilbert transport covariance and the exceptional motion-without-birth branch}
\label{sec:V3-router-covariance}

Two flattenings of a moving domain need not have identical ordinary
Hamiltonian derivatives.  Their difference is not a free new source.
In a reference Hilbert chart let $\Gamma_u^*=-\Gamma_u$ be the specified
connection for comparing physical Hilbert fibres.  On its compatible common
core use the convention
\begin{equation}
 \nabla_u h=\dot h_u-[\Gamma_u,h],\qquad
 \nabla_u\sigma=\dot\sigma_u-[\Gamma_u,\sigma].
 \label{eq:V3-Hilbert-connection}
\end{equation}
The preceding material formulae use the declared material pullback as
reference connection, so $\Gamma_u=0$ in that reference chart.  Choosing a
different flattening does not reset the physical connection to zero.

\begin{theorem}[Covariant Duhamel source under a chart change]
\label{thm:V3-chart-covariance}
Let $R_s$ be a differentiable unitary change of reference chart, with the
following identities defined on the compatible core and with the inherited
heat-sandwiched products well defined.  Transform
\begin{equation}
 h'_s=R_sh_sR_s^*,\qquad
 \Gamma'_u=\dot R_uR^*+R\Gamma_uR^*.
 \label{eq:V3-Gamma-transform}
\end{equation}
Then
\begin{equation}
 \boxed{\nabla'_uh'=R(\nabla_uh)R^*,\qquad
        \nabla'_u\sigma'=R(\nabla_u\sigma)R^*.}
 \label{eq:V3-covariant-source}
\end{equation}
Thus the correctly transported primitive Grams, physical reserve, Duhamel
residual, and the boundary-jet kernel with its physical boundary trace also transported
are unchanged.  If $\Gamma'$ is incorrectly
reset to zero, the apparent action gains a term $[\dot R R^*,h']$ and the
apparent target gains $[\dot R R^*,\sigma']$.
\end{theorem}

\begin{proof}
Differentiate $RhR^*$ and use $\dot R R^*+R\dot R^*=0$; the two
inhomogeneous commutator terms cancel against
\eqref{eq:V3-Gamma-transform}.  Functional calculus gives
$\sigma'=R\sigma R^*$ and exactly the same cancellation proves its
covariant derivative identity.  The Duhamel divided difference transforms by
unitary covariance, as do the complete prior primitive range and spectral
projections.  Their Gram and residual conclusions follow on the same source
and target metrics.  The last two terms are the uncancelled derivatives when
the transformed connection is discarded.  For unbounded differential chart
generators this argument asserts only the displayed compatible-core and
heat-sandwiched identities, not a new unrestricted commutator domain.
\end{proof}

\begin{proposition}[Moving walls can lie on a complete follower branch]
\label{prop:V3-moving-follower}
Nonzero wall velocity alone does not force $R_{\rm dyn}>0$.  For example,
let $\Omega_s=e^s\Omega$ and let the physical family be the scalar
Dirichlet Laplacian on $\Omega_s$.  The predeclared unitary dilation gives
\begin{equation}
 \widetilde h_s=e^{-2s}h,\qquad a=-2h.
 \label{eq:V3-dilation}
\end{equation}
On the complete declared weight/action bank the Duhamel target is a left-clock
follower: its residual is zero.  The nonzero geometric speed is
$v(x)=x\cdot\nu(x)$.  For a translated domain with translated potential,
the corresponding translation pullback gives $a=0$.
\end{proposition}

\begin{proof}
The unitary dilation is $U_sf(x)=e^{ds/2}f(e^sx)$; differentiating its
conjugated Laplacian gives \eqref{eq:V3-dilation}.  The centered $a$ commutes
with $h$.  The normalized Gibbs square-root derivative is
$-b(a-\Tr(\rho a)I)\sigma$, a spectral function of $h$ times the weight
column.  In every left eigenspace it is therefore proportional to the
weight row already present.  The complete left-cyclic short has zero
residual.  The velocity follows by differentiating $e^sx$, and the
translation identity follows by direct change of variables.  These are
explicit separating mechanisms, not a second numerical diagnostic replacing
the ESM card.
\end{proof}

\begin{remark}[Consistency with the normal-jet theorem]
The theorem about admitted compactly supported velocities uses a specific
constant-normal collar transport with fixed principal ambient coefficients.
A global dilation has its specified dilation transport and does not satisfy
that compact-support hypothesis.  On a sphere its normal velocity is
constant, which is also an explicit exceptional mode of the boundary
divergence test.  A chart change cannot move an example into the passing
hypotheses while silently changing its physical Hilbert connection.
\end{remark}

\section{The concrete ESM handoff and the exact remaining numerical task}
\label{sec:V3-handoff}

The native domain calculation now has an explicit sequence on the original
finite card:
\begin{equation}
 \boxed{
 \begin{gathered}
 (W,D_uW,\dd_xW,g)\ \longrightarrow\ (v_u,\hbox{regular-wall chart})\\
 \longrightarrow\ (\hbox{physical }U_s,\ a_u^{\rm int}+[D_{X_u},h])\\
 \longrightarrow\ \Jmat u
 \ \longrightarrow\ \hbox{positive full-primitive fibre or unresolved zero jet}.
 \end{gathered}}
 \label{eq:V3-handoff-order}
\end{equation}
The first arrow is the finite matrix calculation
\eqref{eq:V3-Wilson-velocity} or its exact Schur version.  The second uses
all kinetic, unitary-connection, line, and potential contributions of the
same common action.  The third is evaluated by
\eqref{eq:V3-native-jet} on its fixed-principal normal-collar branch, with
any actual principal variation added rather than ignored.  Only the fourth
uses the inherited thermal ancestry short.

\begin{theorem}[Version V3 finite-card export]
\label{thm:V3-export}
On one original ESM card, suppose the source metric and target metric, the
physical Hilbert transport, the complete prior source bank, and the
Duhamel identification have been fixed independently of the residual.
Suppose a regular sector-wall patch satisfies the native analytic hypotheses
above.  Then:
\begin{enumerate}[label=\textup{(E\arabic*)},leftmargin=3em]
\item its normal velocity is determined by the actual finite Wilson derivative;
\item its complete material-action first boundary jet is a sufficient positive
      full-primitive Duhamel certificate whenever nonzero;
\item under the weighted-divergence hypotheses, the real residual rank is at
      least the real dimension of admitted wall velocities modulo the stated
      constant modes;
\item a zero jet, a singular wall, a non-native flattening, or a failed
      form/source hypothesis does not prove followerhood or a negative physical
      result; it selects the corresponding unresolved or ill-typed branch.
\end{enumerate}
No part of this theorem supplies a cutoff-uniform number or a cofinal ESM
ancestry limit.
\end{theorem}

\begin{proof}
Items \textup{(E1)--(E3)} are respectively
\Cref{thm:V3-Wilson-velocity,thm:V3-Schur-velocity},
\Cref{thm:V3-local-trace-birth,thm:V3-collar-jet}, and
\Cref{thm:V3-boundary-divergence}.  Every implication is one-sided and
contains the displayed analytic and source hypotheses.  A failed premise is
therefore not its contrapositive conclusion.  V2 already distinguishes strict
finite-card positivity from a numerical uniform floor, and V1 requires a
transported Cauchy or paired-card certificate for a limit; neither premise is
added by the local trace argument.
\end{proof}

\paragraph{What has actually been evaluated.}
The new finite coefficients are the exact normal velocity formula, the
Schur-normalization cancellation, the Pauli wall and polarization matrices
\eqref{eq:V3-Pauli-matrices}, the complete material spectral matrix
\eqref{eq:V3-Hadamard-matrix}, and the curvature-free boundary differential
\eqref{eq:V3-native-jet}.  The positive-rank implication is proved for the
stated native regular-wall branch.  No Wilson zero configuration, original
ESM eigenfunction, numerical source reserve, or numerical ESM residual has
been manufactured to assign those formulae a physical value.

\paragraph{What still has to be read on the original card.}
The source corpus supplies the Dirichlet sector type, the smooth
simple-crossing reduction, the quotient Wilson jets by formula, and the
thermal source construction.  The current calculation does not supply the
selected command's actual $v_u$, the occupied normal eigenfunction rows
$f_a$, their unitary-connection variation, the complete physical material
router, or a bounded annihilator with a physical-metric lower margin.  It
also does not supply the original spacetime-cutoff comparison.  These are
explicit values and incidence checks, not a request to rebuild a general
clock or shorting theorem.

\paragraph{The next iteration.}
Evaluate the first predeclared physical source on one regular occupied wall
patch using \eqref{eq:V3-Wilson-velocity}.  If the wall is fixed, return to
V2's fixed-wall jet with the actual action; if it moves, retain the physical
material transport and evaluate the total jet
\eqref{eq:V3-material-jet-definition}.  A nonzero result warrants extraction
of a bounded, target-normalized finite spectral annihilator and its route
comparison.  A zero result requires the next compatibility jet or the
literal five-matrix evaluation; it is not a follower certificate.  An actual
metric/coframe source must first include its principal-operator variation.
Do not add a detached wall-pressure primitive or substitute the angular
Pauli metric for the frozen source action.

\begin{assessment}[Version V3 verdict]
\label{ass:V3-verdict}
\passstatus{} for the finite Wilson/Schur domain-motion identities, the
Pauli metric completion, the material-action formulae, and the conditional
regular-wall thermal rank theorem.

\obsstatus{} for treating the isolated nonzero Hadamard wall form as a
bounded energy-form primitive.  Its boundary-layer witness is explicit.
A nonzero native total material jet is also a typed obstruction to complete
followerhood, but no such numerical jet is claimed on the original ESM card
in this version.

\stopstatus{} for the actual cofinal ESM reserve and residual until the
original physical values and their routes are evaluated.  Determinant/Pfaffian
orientation, global line and graph domain, interacting continuum, current,
stress, Lorentzian reconstruction, Einstein backreaction, and pure-colour
mass claims are not promoted by a local wall calculation.
\end{assessment}

\paragraph{Exact algebraic replay.}
The accompanying script
\path{Renewal_SM_Einstein_Material_Wall_Identities_V3.py}
checks the mixed Pauli identity (including the Berry sign), the Schur
null-vector normalization identity, the complete normal-Jacobian cancellation,
and the vanishing real pairing of a skew-adjoint connection coefficient.
It uses symbolic and exact rational arithmetic.  Its JSON report is an
algebra check, not a physical ESM acquisition or a substitute for the domain
proofs above.

\section*{Version V3 append-only record}
\addcontentsline{toc}{section}{Version V3 append-only record}

Preserved the complete V2 byte prefix before its sole terminal
\verb|\end{document}|.  The SHA--256 of the complete V2 input is
\begin{center}\small\ttfamily
ed5e994cebd1b8ea9ab2c6cff449e6c34953d00eb62e92dd03d32e2f541ab0f1.
\end{center}

Returned to the actual Friedrichs--Dirichlet Wilson-gap sector and imported
its existing simple-crossing, Feshbach, and Pauli results.  Calculated the
native wall speed and proved exact cancellation of the Schur lift's
normalization.  Completed the Pauli transition metric by the two wall-defining
numerators, including the previously invisible scalar component.  Derived
the complete material energy form, finite spectral Hadamard matrix, and
thermal boundary pressure; proved that the detached wall form fails the
energy-domain bound.  Extended the Duhamel domain test to energy-dual form
rows.  Evaluated the normal-collar material jet with exact curvature
cancellation and proved its weighted-divergence real-rank consequence.
Retained physical Hilbert-connection covariance and an explicit
motion-without-birth branch.  No prior diagnostic was numerically rerun or
identified with an ESM spacetime regulator.

\begin{firewall}[Append-only instruction after V3]
\label{fire:V3-append}
V4 must preserve this complete V3 source verbatim, append immediately before
the sole final document-closing command, and use
\begin{center}\small
\path{Renewal_Einstein_Standard_Model_Physical_Source_Metric_Duhamel_Ancestry_and_Quantum_Continuum_Programme_V4.tex}.
\end{center}
Do not repeat the Pauli polarization calculation, the V1 paired squeeze,
the V2 interval control, or the present material-domain identities.  Evaluate
the native total jet or a literal original ESM fibre.  Do not treat wall
velocity as a birth, scalar polarization blindness as zero domain motion,
a Schur matrix as an isospectral Wilson replacement, a boundary-pressure
piece as a form primitive, a reference flattening as the physical router,
a real-rank statement as an unrestricted complex-rank theorem, or strict
finite-card positivity as a cofinal margin.
\end{firewall}

% =============================================================================
% END APPENDED VERSION V3 MATERIAL
% =============================================================================


% =============================================================================
% BEGIN APPENDED VERSION V4 MATERIAL
% =============================================================================
\clearpage
\providecommand{\Vraw}{\mathcal V}
\providecommand{\HL}{H_{\mathrm L}}
\providecommand{\HR}{H_{\mathrm R}}
\providecommand{\Pcyc}{P_{H}^{B}}
\providecommand{\Kfrac}{\mathsf K}
\providecommand{\Hfrac}{\mathsf K^{\mathrm{head}}}
\providecommand{\qtherm}{q_{\beta}}
\providecommand{\btherm}{b_{\beta,\tau}}
\providecommand{\dtherm}{d_{\beta}}

\section*{Version V4: thermal-source rigidity, reserve normalization, and a corrected cofinal ancestry card}
\addcontentsline{toc}{section}{Version V4: thermal-source rigidity and corrected cofinal ancestry}

\begin{abstract}
Version~V3 reduced a moving overlap-sector wall to an exact material boundary
jet, but the newly supplied population report changes the first acquisition
priority.  That report places the ancestry compiler on a concrete reversible
refresh-plus-heat-bath slab, identifies collapse of a global compact-horizon
floor as cyclic depth grows, and reports stable normalized dynamic fractions.
The present version audits whether its finite matrices are the physical
sector-patched Duhamel card already fixed by the line/clock programme.

The audit finds one decisive mismatch.  The physical primitive action column is
thermally dressed,
\[
 B_{\rm act}=Z^{-1/2}e^{-\tau H_{\rm L}}
               e^{-\beta H_{\rm R}/2}\Vraw,
\]
whereas the population model inserts an undressed multiplication column and
then applies the quotient multiplier $m_{\beta,\tau}$.  The apparently growing
factor $m_{\beta,\tau}$ is not the target itself: it is exactly the quotient of
the fixed, decaying Fr\'echet derivative of $e^{-\beta H/2}$ by the thermally
dressed source.  On every occupied bi-energy block that quotient is unique.
The sign-reversed ``decaying calibration'' proposed in the population report
therefore defines a different target, not a reparametrization of the inherited
quantum half-source.

This correction has three useful consequences.  First, an explicit same-left-
energy two-block card proves that omitting the thermal source dressing can turn
an exponentially decaying physical fibre residual into an exponentially
growing surrogate residual.  Second, the true Duhamel divided difference has
one full left-energy order of decay for every bounded raw insertion.  This
gives a regulator-independent positive omitted-head bound of order
$(1+\Lambda)^{-2}$, with no sign reversal and no condition
$2\tau>\beta$.  The one-order statement is optimal from bounded insertion data
alone.  Third, the compact-horizon obstruction is proved more directly by a
trace budget:
\[
 \lambda_{\min}^{+}(\widehat{\mathscr W}_{T,B})
 \le \frac{T\rank P_B}{\rank\widehat{\mathscr W}_{T,B}}.
\]
Thus a global horizon floor is indeed impossible at unbounded cyclic depth per
source rank, but it was never needed by the target-relative Tikhonov or
Duhamel-head compiler.

The source reserve requires a separate correction.  Ambient trace
normalization is not harmless unless the physical source-action metric is
scaled with it.  The inherited thermal weight column gives the exact Rayleigh
upper bound
\[
 g_-\le \frac{Z(\beta+2\tau)}{Z(\beta)}
      \le e^{-2\tau\rho}
\]
when the centered clock has gap $\rho$ and the weight coefficient has unit
physical cost.  Dividing the Hilbert--Schmidt inner product by carrier
dimension divides this bound by the same dimension if the source metric is
left fixed.  Hence the large reserve values in the raw-source population table
cannot be values of the inherited thermally dressed card.

Finally, a normalization-free cofinal theorem is reformulated without turning
asymptotic smallness into an exact finite follower.  A physical Duhamel head
which exhausts the target-normalized residual yields exactly three outcomes:
a robust target-normalized survivor, asymptotic follower fraction, or a typed
support/metric/tail/route obstruction.  Exact followerhood remains the exact
zero branch at each cutoff.  A new deterministic replay of the Ising slab with
the inherited thermal formulas still displays positive dynamic fractions on
all tested fibres, but very different reserve values and no ultraviolet target
explosion.  It is a floating-point diagnostic of the regulator class, not the
missing Dirac--Yukawa card and not an Einstein continuum result.
\end{abstract}

% =============================================================================
\section{Direction update after the supplied population and terminal-concentration reports}
\label{sec:V4-direction}
% =============================================================================

The population report contains a useful literal model and one correct warning:
a smallest positive eigenvalue of the complete range-normalized horizon
Gramian cannot remain uniform when a fixed-rank local source generates an
increasing number of cyclic directions.  It also proposes two further steps:
replace the inherited multiplier by a sign-reversed decaying multiplier, and
read a raw multiplication-source computation as a population of the physical
ancestry card.  Those steps require an incidence audit before they can enter
this programme.

The inherited line/clock theorem does not begin with an arbitrary vector called
$A$.  It fixes, on one same-history standard-form carrier, the three primitive
roles
\[
 B=(B_{\rm wt},B_{\rm act},B_{\rm line})
\]
and the internal square-root target $Y_{\rm int}$.  In a bi-energy block the
calibrated multiplier is merely the scalar quotient satisfying
$Y_{\rm int}=m_{\beta,\tau}B_{\rm act}$.  The source and target must therefore
be changed together or not at all.

The terminal-concentration report concerns the Navier--Stokes source-paid
caloric scalar.  Its transferable lesson is an audit discipline rather than an
SM coefficient: raw marginal size, a source-relative normalized record, and an
exact positive residual must not be interchanged.  None of its PDE identities
is imported as an Einstein--Standard--Model estimate.

\begin{firewall}[What Version V4 does not repeat]
\label{fire:V4-no-repeat}
The complete mixed sector--quantum Schur short, the five-matrix fibre compiler,
the V1 paired lower/upper squeeze, the V2 hard-wall interval certificate, and
the V3 material boundary jet are retained verbatim.  Version~V4 audits the
new population card at the earlier source/target incidence and metric stages.
The wall jet remains an independent one-sided annihilator when the actual
Dirac--Yukawa card supplies its domain data; it is not used to repair an
incorrectly normalized population model.
\end{firewall}

% =============================================================================
\section{The physical thermal factorization and rigidity of the calibrated multiplier}
\label{sec:V4-thermal-rigidity}
% =============================================================================

Let $H=H^*\succeq0$ have discrete spectrum on a finite cutoff carrier
$\mathcal H$, and let $\mathfrak S_2(\mathcal H)$ carry its declared
Hilbert--Schmidt inner product.  On this carrier put
\[
 \HL X=HX,
 \qquad
 \HR X=XH.
\]
The two operators commute.  Let $D$ be the finite common-action coefficient
space and let
\[
 \Vraw:D\longrightarrow\mathfrak S_2(\mathcal H)
\]
be the partition-normalized raw centered insertion synthesis.  Sector weights
and direct sums may be included in $\Vraw$; the argument is blockwise and is
unchanged by them.  Write
\[
 a:=\frac\beta2>0,
 \qquad \tau>0.
\]
The inherited action source and internal square-root target are
\begin{align}
 B_{\beta,\tau}
 &=e^{-\tau\HL}e^{-a\HR}\Vraw,
 \label{eq:V4-physical-B}\\
 Y_{\beta}
 &=-a\int_0^1e^{-a(1-q)\HL}e^{-aq\HR}\Vraw\,\dd q.
 \label{eq:V4-physical-Y}
\end{align}
The partition scalar has been absorbed into $\Vraw$ so that it cannot be
mistaken for a new source metric.

For left and right energies $x,y\ge0$, define
\begin{align}
 \btherm(x,y)&:=e^{-\tau x-ay},
 \label{eq:V4-b-factor}\\
 \dtherm(x,y)&:=-a\int_0^1e^{-a[(1-q)x+qy]}\,\dd q
 \notag\\
 &=\begin{cases}
 \displaystyle\frac{e^{-ax}-e^{-ay}}{x-y},&x\ne y,\\[0.8em]
 -ae^{-ax},&x=y,
 \end{cases}
 \label{eq:V4-d-factor}\\
 m_{\beta,\tau}(x,y)&:=\frac{\dtherm(x,y)}{\btherm(x,y)}.
 \label{eq:V4-m-ratio}
\end{align}

\begin{theorem}[Thermal quotient identity and blockwise rigidity]
\label{thm:V4-thermal-rigidity}
Let $P_x,P_y$ be spectral projections of $H$.  On every bi-energy block,
\begin{equation}
 \boxed{
 P_xY_{\beta}(u)P_y
 =m_{\beta,\tau}(x,y)
  P_xB_{\beta,\tau}(u)P_y.}
 \label{eq:V4-block-factorization}
\end{equation}
Moreover:
\begin{enumerate}[label=\textup{(T\arabic*)},leftmargin=3em]
\item $\btherm(x,y)>0$ and $\dtherm(x,y)<0$ for all $x,y\ge0$;
\item if a scalar field $\widetilde m(x,y)$ satisfies
      $P_xY_\beta(u)P_y=\widetilde m(x,y)P_xB_{\beta,\tau}(u)P_y$
      on one nonzero physical insertion block, then
      $\widetilde m(x,y)=m_{\beta,\tau}(x,y)$;
\item an invertible coefficient reparametrization $u=S\widetilde u$ changes
      $B$ and $Y$ by the same precomposition and leaves the quotient
      $m_{\beta,\tau}$ unchanged;
\item the sign-reversed function
      \[
       m_{-\beta,-\tau}(x,y)
       =a e^{-\tau x}\int_0^1e^{-aq(y-x)}\,\dd q>0
      \]
      cannot be the quotient of the inherited physical source and target on
      any nonzero block.
\end{enumerate}
Thus the growing factor in $m_{\beta,\tau}$ is exactly compensated by the
thermal factor already present in $B_{\beta,\tau}$.  Replacing
$m_{\beta,\tau}$ by $m_{-\beta,-\tau}$ changes the target rather than merely
calibrating it.
\end{theorem}

\begin{proof}
On the $(x,y)$ block, \eqref{eq:V4-physical-B} multiplies the raw insertion by
$\btherm(x,y)$.  Equation \eqref{eq:V4-physical-Y} multiplies it by
$\dtherm(x,y)$.  Their quotient is \eqref{eq:V4-m-ratio}, proving
\eqref{eq:V4-block-factorization}.  The integral formula for $\dtherm$ is
strictly negative and $\btherm$ is strictly positive.  If the insertion block
is nonzero, equality of two scalar multiples of that block forces equality of
the scalars.  Precomposition by $S$ affects both maps identically, so the
block quotient is unchanged.  Substitution of $(-\beta,-\tau)$ into the
integral definition gives the displayed positive expression, which cannot
equal the strictly negative physical quotient.
\end{proof}

\begin{corollary}[Delay appears in the primitive source, not in the Duhamel target]
\label{cor:V4-delay-cancellation}
The target multiplier $\dtherm(x,y)$ is independent of $\tau$.  Consequently
no condition comparing $2\tau$ with $\beta$ is required to control the
left-energy tail of the physical target.  Such a condition arises only after
the source factor $\btherm$ has been removed while its quotient multiplier is
retained.
\end{corollary}

\begin{proof}
This is immediate from \eqref{eq:V4-d-factor}.  The identity
$\dtherm=m_{\beta,\tau}\btherm$ shows exactly where the cancellation occurs.
\end{proof}

% =============================================================================
\section{A same-left-energy counterexample to the undressed-source surrogate}
\label{sec:V4-raw-surrogate}
% =============================================================================

The preceding theorem is an incidence statement.  The next calculation shows
that violating it can reverse the ultraviolet behavior of a fibre residual.
It uses one left-energy fibre, so the distinction is not erased by subsequent
left-clock cyclic saturation.

\begin{countertheorem}[Omitting source dressing can reverse fibre-tail asymptotics]
\label{cth:V4-raw-surrogate}
Fix $R>0$.  Let one left-energy fibre be $\mathcal K_R=\C^2$ with
$\HL=RI$ and right energies $0,R$.  Let the raw insertion coefficient have
unit amplitudes in both right-energy cells.  Set
\[
 s_R:=\frac{1-e^{-aR}}{R}.
\]
The physical source and target vectors, up to their common partition scalar,
are
\begin{equation}
 B_R=e^{-\tau R}(1,e^{-aR}),
 \qquad
 Y_R=(-s_R,-ae^{-aR}).
 \label{eq:V4-two-block-physical}
\end{equation}
Their complete one-source fibre residual is
\begin{equation}
 \boxed{
 r_{\rm phys}(R)
 =\dist(Y_R,\Span\{B_R\})^2
 =\frac{e^{-2aR}(a-s_R)^2}{1+e^{-2aR}}.}
 \label{eq:V4-physical-residual}
\end{equation}

If instead the undressed raw source
$\widetilde B_R=(1,1)$ is retained while the quotient multiplier
$m_{\beta,\tau}(R,\cdot)$ is applied, then
\begin{equation}
 \widetilde Y_R
 =-e^{\tau R}(s_R,a),
 \qquad
 \boxed{
 r_{\rm raw}(R)
 =\frac{e^{2\tau R}(a-s_R)^2}{2}.}
 \label{eq:V4-raw-residual}
\end{equation}
Consequently
\begin{equation}
 r_{\rm phys}(R)\sim a^2e^{-2aR},
 \qquad
 r_{\rm raw}(R)\sim\frac{a^2}{2}e^{2\tau R}.
 \label{eq:V4-opposite-asymptotics}
\end{equation}
The undressed-source computation is therefore not a coordinate version of the
physical card.
\end{countertheorem}

\begin{proof}
For two vectors $b,y\in\C^2$,
\[
 \dist(y,\Span\{b\})^2
 =\frac{|\det(b,y)|^2}{\|b\|^2}.
\]
Using \eqref{eq:V4-two-block-physical},
\[
 \det(B_R,Y_R)
 =e^{-(\tau+a)R}(s_R-a),
 \qquad
 \|B_R\|^2=e^{-2\tau R}(1+e^{-2aR}),
\]
which gives \eqref{eq:V4-physical-residual}.  The physical quotient values are
$m_{\beta,\tau}(R,0)=-e^{\tau R}s_R$ and
$m_{\beta,\tau}(R,R)=-ae^{\tau R}$, giving the surrogate target in
\eqref{eq:V4-raw-residual}.  Its distance from the diagonal line is one half
the squared coordinate difference.  Since $s_R\to0$, the asymptotics follow.
\end{proof}

\begin{remark}[Scope of the counterexample]
Adding independently occurring primitive columns can only decrease a residual
on either card.  It cannot turn the two source syntheses into the same physical
object or license transfer of a tail verdict between them.  The example is a
minimal exact obstruction to the asserted source/target identification, not a
numerical verdict on the complete ESM primitive bank.
\end{remark}

% =============================================================================
\section{The physical divided difference has an optimal one-order left-energy tail}
\label{sec:V4-physical-tail}
% =============================================================================

Write
\begin{equation}
 \qtherm(x,y):=|\dtherm(x,y)|
 =a\int_0^1e^{-a[(1-q)x+qy]}\,\dd q.
 \label{eq:V4-q-def}
\end{equation}

\begin{lemma}[Uniform divided-difference envelope]
\label{lem:V4-divided-envelope}
For fixed $x\ge0$, $y\mapsto\qtherm(x,y)$ is decreasing, and
\begin{equation}
 \boxed{
 \qtherm(x,y)
 \le\begin{cases}
 (1-e^{-ax})/x,&x>0,\\
 a,&x=0,
 \end{cases}
 \le\min\{a,x^{-1}\}}
 \label{eq:V4-q-envelope}
\end{equation}
with the last expression interpreted by omitting $x^{-1}$ at $x=0$.
For every $0\le s\le1$,
\begin{equation}
 \boxed{
 C_{a,s}:=\sup_{x,y\ge0}(1+x)^s\qtherm(x,y)
 \le2\max\{a,1\}.}
 \label{eq:V4-Cas}
\end{equation}
\end{lemma}

\begin{proof}
Differentiation under the integral gives
\[
 \partial_y\qtherm(x,y)
 =-a^2\int_0^1q e^{-a[(1-q)x+qy]}\,\dd q\le0.
\]
The maximum in $y$ is therefore attained at $y=0$, where direct integration
gives $(1-e^{-ax})/x$.  The elementary inequality
$1-e^{-z}\le\min\{z,1\}$ yields the second bound.  If $0\le x\le1$, then
$(1+x)^s\qtherm\le2a$.  If $x\ge1$, then
$(1+x)^s\qtherm\le2x^{s-1}\le2$.  Taking the larger constant proves
\eqref{eq:V4-Cas}.
\end{proof}

Let $E_{>\Lambda}^{\rm L}$ denote the left-energy spectral projection of
$\HL$ on $(\Lambda,\infty)$.  Let $P_x^{\rm prim}$ be the orthogonal
projection onto the complete prior primitive range in the left-energy-$x$
fibre.  Define the omitted physical Duhamel residual tail by
\begin{equation}
 \ThetaD_{>\Lambda}
 :=\sum_{x>\Lambda}Y_{\beta,x}^*
      (I-P_x^{\rm prim})Y_{\beta,x}.
 \label{eq:V4-tail-definition}
\end{equation}
The sum is finite at a finite cutoff and is understood as a monotone form sum
on a discrete compact-resolvent card.

\begin{theorem}[Physical target regularity and positive residual-tail bound]
\label{thm:V4-physical-tail}
For every $0\le s\le1$,
\begin{equation}
 \boxed{
 Y_\beta^*(I+\HL)^{2s}Y_\beta
 \preceq C_{a,s}^2\Vraw^*\Vraw.}
 \label{eq:V4-target-regularity}
\end{equation}
Consequently, for every $\Lambda\ge0$,
\begin{equation}
 \boxed{
 0\preceq\ThetaD_{>\Lambda}
 \preceq Y_\beta^*E_{>\Lambda}^{\rm L}Y_\beta
 \preceq C_{a,s}^2(1+\Lambda)^{-2s}\Vraw^*\Vraw.}
 \label{eq:V4-tail-bound}
\end{equation}
In particular, at $s=1$,
\begin{equation}
 \boxed{
 \ThetaD_{>\Lambda}
 \preceq4\max\{\beta/2,1\}^2(1+\Lambda)^{-2}\Vraw^*\Vraw.}
 \label{eq:V4-tail-s-one}
\end{equation}
No horizon floor, source-delay inequality, or sign-reversed multiplier occurs.
\end{theorem}

\begin{proof}
Resolve the commuting pair $(\HL,\HR)$ spectrally.  For $u\in D$,
\begin{align*}
 \|(I+\HL)^sY_\beta u\|_{\HS}^2
 &=\sum_{x,y}(1+x)^{2s}\qtherm(x,y)^2
   \|P_x\Vraw(u)P_y\|_{\HS}^2\\
 &\le C_{a,s}^2\sum_{x,y}\|P_x\Vraw(u)P_y\|_{\HS}^2
 =C_{a,s}^2\|\Vraw(u)\|_{\HS}^2.
\end{align*}
Polarization gives the operator inequality
\eqref{eq:V4-target-regularity}.  Orthogonal projection can only decrease a
norm, so the first two inequalities in \eqref{eq:V4-tail-bound} hold.  On the
range of $E_{>\Lambda}^{\rm L}$,
$(I+\HL)^{-2s}\preceq(1+\Lambda)^{-2s}I$.  Applying
\eqref{eq:V4-target-regularity} gives the final bound.  Insert
\eqref{eq:V4-Cas} at $s=1$.
\end{proof}

\begin{corollary}[Target-normalized head exhaustion]
\label{cor:V4-relative-tail}
Let $M_D\succeq0$ be the physical common-action form on $D$, let
$U:=Y_\beta^*Y_\beta$, and assume on their common supported quotient that
\begin{equation}
 \Vraw^*\Vraw\preceq c_V M_D,
 \qquad
 M_D\preceq c_TU
 \label{eq:V4-incidence-window}
\end{equation}
for finite constants $c_V,c_T$.  Then
\begin{equation}
 \boxed{
 \bigl\|U^{\dagger/2}\ThetaD_{>\Lambda}U^{\dagger/2}\bigr\|
 \le4\max\{\beta/2,1\}^2c_Vc_T(1+\Lambda)^{-2}.}
 \label{eq:V4-relative-tail-bound}
\end{equation}
Thus increasing physical left-energy heads exhaust the dynamic fraction under
a source/action/target incidence window, without a compact-horizon minimum
eigenvalue.
\end{corollary}

\begin{proof}
Combine \eqref{eq:V4-tail-s-one} with
\eqref{eq:V4-incidence-window}.  If $0\preceq A\preceq cU$, then supported
congruence by $U^{\dagger/2}$ gives
$0\preceq U^{\dagger/2}AU^{\dagger/2}\preceq cP_U$.
\end{proof}

\begin{countertheorem}[One left-energy order is optimal for bounded insertions]
\label{cth:V4-tail-sharp}
For every $s>1$ there is no constant depending only on $a$ such that
\[
 Y_\beta^*(I+\HL)^{2s}Y_\beta\preceq C\Vraw^*\Vraw
\]
for all finite nonnegative Hamiltonians and bounded raw insertions.
\end{countertheorem}

\begin{proof}
On a two-level carrier take one raw Hilbert--Schmidt insertion supported only
from right energy $0$ to left energy $R$, with unit Hilbert--Schmidt norm.  Its
physical target norm is
\[
 \frac{1-e^{-aR}}{R}.
\]
After applying $(1+\HL)^s$ the norm is
$(1+R)^s(1-e^{-aR})/R$, which diverges as $R\to\infty$ whenever $s>1$.
\end{proof}

\begin{remark}[What sharpness does and does not say]
The counterexample is sharp for a bound derived from bounded raw insertion
alone.  Additional locality, commutator, form-domain, or primitive-absorption
information may improve the residual tail on a particular regulator.  Such an
improvement must be read from that physical card; it is not obtained by
changing the Duhamel target.
\end{remark}

% =============================================================================
\section{A trace-budget obstruction to a global compact-horizon floor}
\label{sec:V4-horizon-trace}
% =============================================================================

Let $P_B$ be the orthogonal projection onto the entrance source range and let
$r:=\rank P_B<\infty$.  For $T>0$ define the range-normalized horizon Gramian
\begin{equation}
 \What_{T,B}=\int_0^Te^{-tH}P_Be^{-tH}\,\dd t.
 \label{eq:V4-W-hat}
\end{equation}
At a finite card let $d:=\rank\What_{T,B}$ and
$\widehat w_T:=\lambda_{\min}^{+}(\What_{T,B})$.

\begin{theorem}[Horizon trace budget]
\label{thm:V4-horizon-trace-budget}
For every nonnegative self-adjoint $H$,
\begin{equation}
 \boxed{
 \Tr\What_{T,B}
 =\int_0^T\Tr(P_Be^{-2tH})\,\dd t
 \le rT.}
 \label{eq:V4-W-trace}
\end{equation}
If $d<\infty$, then
\begin{equation}
 \boxed{
 \widehat w_T\le\frac{rT}{d}.}
 \label{eq:V4-horizon-floor-bound}
\end{equation}
If the $H$-cyclic source space is infinite dimensional, $\What_{T,B}$ is
trace class and has no strictly positive lower bound on its support.  Along a
cofinal family, $d_n/r_n\to\infty$ therefore forces
$\widehat w_{T,n}\to0$ at every fixed $T$.
\end{theorem}

\begin{proof}
Cyclicity of the trace and positivity give
\[
 \Tr(e^{-tH}P_Be^{-tH})=\Tr(P_Be^{-2tH})\le\Tr P_B=r.
\]
Integration proves \eqref{eq:V4-W-trace}.  If the $d$ positive eigenvalues are
all at least $\widehat w_T$, then
$d\widehat w_T\le\Tr\What_{T,B}$, proving
\eqref{eq:V4-horizon-floor-bound}.  In infinite dimension a positive lower
bound $\What_{T,B}\succeq cP_{\supp\What}$ would give infinite trace, contrary
to \eqref{eq:V4-W-trace}.  The cofinal conclusion follows from the finite
bound.
\end{proof}

\begin{assessment}[What the cyclic-depth report gets right]
\label{ass:V4-horizon-assessment}
The supplied population report is correct that a global supported horizon
floor is unattainable when cyclic depth grows faster than source rank.  The
trace budget proves this without a spectral-window or Cauchy-matrix estimate.
The conclusion does not obstruct this programme: V1 and the inherited V37
compiler use a target-visible Tikhonov tail, not a minimum eigenvalue over every
cyclic direction.  Small horizon eigenvalues which carry vanishing target
spectral mass are harmless; a positive global floor was an unnecessarily
strong proposed population premise.
\end{assessment}

% =============================================================================
\section{Physical reserve, trace normalization, and the thermal weight-ray test}
\label{sec:V4-reserve-normalization}
% =============================================================================

Let $Q=B^*B$ be the primitive source Gram and let $M$ be its physical
source-action form.  On the supported quotient the reserve spectrum is the
positive spectrum of
\begin{equation}
 G_{B,M}=M^{\dagger/2}QM^{\dagger/2}.
 \label{eq:V4-reserve-Gram}
\end{equation}

\begin{proposition}[Ambient normalization is physical unless the source metric follows]
\label{prop:V4-normalization-scaling}
Multiply the target Hilbert inner product by $c>0$.  If the coefficient metric
$M$ is held fixed, then
\begin{equation}
 Q\mapsto cQ,
 \qquad
 G_{B,M}\mapsto cG_{B,M}.
 \label{eq:V4-reserve-scale}
\end{equation}
If instead the physical metric is simultaneously changed to $cM$, the reserve
spectrum is unchanged.  For a target Gram $U$ and residual $R$, simultaneous
ambient scaling sends $(U,R)$ to $(cU,cR)$ and leaves the dynamic fraction
\begin{equation}
 U^{\dagger/2}RU^{\dagger/2}
 \label{eq:V4-fraction-invariant}
\end{equation}
unchanged.
\end{proposition}

\begin{proof}
The source and target Grams are inner products and therefore acquire the
factor $c$.  Congruence by $M^{\dagger/2}$ gives
\eqref{eq:V4-reserve-scale}.  Since $(cM)^{\dagger/2}=c^{-1/2}M^{\dagger/2}$,
scaling both forms cancels.  The same pseudoinverse scaling proves
\eqref{eq:V4-fraction-invariant}.
\end{proof}

\begin{theorem}[Exact thermal weight-ray reserve bound]
\label{thm:V4-weight-ray}
Consider one occupied sector with
\[
 Z(s):=\Tr e^{-sH},
 \qquad
 B_{\rm wt}=Z(\beta)^{-1/2}e^{-(\tau+\beta/2)H}.
\]
Under the unnormalized Hilbert--Schmidt trace,
\begin{equation}
 \boxed{
 \|B_{\rm wt}\|_{\HS}^2
 =\frac{Z(\beta+2\tau)}{Z(\beta)}.}
 \label{eq:V4-weight-partition-ratio}
\end{equation}
Suppose the weight coefficient has physical cost $M_{\rm wt,wt}=m_{\rm wt}>0$
and the metric-quotient source map is injective, equivalently
$\supp G_{B,M}=\supp M$.  Then the smallest supported reserve satisfies
\begin{equation}
 \boxed{
 g_-(B,M)
 \le\frac{1}{m_{\rm wt}}
       \frac{Z(\beta+2\tau)}{Z(\beta)}.}
 \label{eq:V4-weight-Rayleigh}
\end{equation}
If $H\succeq\rho I$ on the carrier, then
\begin{equation}
 \boxed{
 g_-(B,M)\le m_{\rm wt}^{-1}e^{-2\tau\rho}.}
 \label{eq:V4-weight-gap-bound}
\end{equation}
If the ambient Hilbert--Schmidt inner product is divided by its dimension $D$
while $M$ is fixed, the right sides of
\eqref{eq:V4-weight-Rayleigh}--\eqref{eq:V4-weight-gap-bound} are divided by
$D$.
\end{theorem}

\begin{proof}
Squaring the standard-form vector and taking the trace gives
\[
 \|B_{\rm wt}\|_{\HS}^2
 =Z(\beta)^{-1}\Tr e^{-2(\tau+\beta/2)H}
 =\frac{Z(\beta+2\tau)}{Z(\beta)}.
\]
Under the injectivity hypothesis the smallest supported reserve is the minimum
of the generalized Rayleigh quotient $u^*Qu/(u^*Mu)$ over the physical
coefficient quotient.  Testing the weight coefficient gives
\eqref{eq:V4-weight-Rayleigh}.  If every energy is at least $\rho$, then
$e^{-(\beta+2\tau)H}\preceq e^{-2\tau\rho}e^{-\beta H}$; taking traces gives
\eqref{eq:V4-weight-gap-bound}.  Normalizing the target trace scales $Q$ by
$D^{-1}$ and the claim follows from
\Cref{prop:V4-normalization-scaling}.
\end{proof}

\begin{corollary}[Exact audit of the supplied raw-source reserve table]
\label{cor:V4-population-reserve-audit}
In the supplied Ising population model the declared four-column source has
rank four, the weight writer is the identity, its coefficient cost is one,
$\tau=1/4$, and the refresh theorem
gives the centered spectral floor $\rho=1/2$.  Any reserve computed from the
inherited thermally dressed primitive bank must therefore satisfy
\begin{equation}
 \boxed{
 g_-^{\rm unnorm}\le e^{-1/4}=0.778800\ldots,
 \qquad
 g_-^{\rm norm}\le D^{-1}e^{-1/4}.}
 \label{eq:V4-population-reserve-bound}
\end{equation}
The reported raw-source values approaching $0.9$ under a claimed normalized
trace cannot be values of that physical bank.  This contradiction is exact
and precedes any diagonalization of the target residual.
\end{corollary}

\begin{proof}
Insert the declared parameters into
\eqref{eq:V4-weight-gap-bound}.  The identity writer has unit
$L^2(\mu)$ cost.  The last assertion follows because a generalized minimum
cannot exceed the weight-coordinate Rayleigh quotient.
\end{proof}

\begin{theorem}[Thermodynamic pressure alternative for the weight reserve]
\label{thm:V4-pressure-reserve}
Let $H_n$ be a volume-growing sequence with volume $V_n\to\infty$, and suppose
for $s\in\{\beta,\beta+2\tau\}$ that
\[
 \frac1{V_n}\log Z_n(s)\longrightarrow p(s).
\]
Then
\begin{equation}
 \boxed{
 \frac1{V_n}\log
 \frac{Z_n(\beta+2\tau)}{Z_n(\beta)}
 \longrightarrow p(\beta+2\tau)-p(\beta).}
 \label{eq:V4-pressure-ratio}
\end{equation}
If the metric-quotient source is injective, the difference is strictly
negative, and the physical weight cost has a positive lower bound, then
$g_{-,n}\to0$ exponentially.  A noncollapsing
weight reserve on such a branch therefore requires a compensating physical
metric scaling, a vanishing delay, or a different declared source quotient;
it cannot be produced by ambient trace normalization alone.
\end{theorem}

\begin{proof}
Subtract the two assumed pressure limits.  The reserve conclusion is
\eqref{eq:V4-weight-Rayleigh}.  The listed alternatives are the only ways to
remove the resulting Rayleigh upper bound while preserving the stated source
role.
\end{proof}

% =============================================================================
\section{A corrected normalization-free cofinal ancestry alternative}
\label{sec:V4-cofinal-alternative}
% =============================================================================

Let the finite target coefficient spaces be identified with one fixed
finite-dimensional space $D$.  At cutoff $n$, write
\begin{equation}
 U_n=Y_n^*Y_n,
 \qquad
 0\preceq L_n\preceq R_n\preceq U_n,
 \label{eq:V4-LRU}
\end{equation}
where $R_n$ is the complete dynamic residual and $L_n$ is a physical finite
left-energy Duhamel head.  Assume the supports of $U_n$ have been transported
to one fixed projection $P$ and are positive on $P$.  Define
\begin{equation}
 \Kfrac_n:=U_n^{\dagger/2}R_nU_n^{\dagger/2},
 \qquad
 \Hfrac_n:=U_n^{\dagger/2}L_nU_n^{\dagger/2},
 \qquad
 \epsilon_n:=\|\Kfrac_n-\Hfrac_n\|.
 \label{eq:V4-normalized-fractions}
\end{equation}
Then $0\preceq\Hfrac_n\preceq\Kfrac_n\preceq P$.

\begin{theorem}[Target-relative survivor/follower alternative under head exhaustion]
\label{thm:V4-cofinal-fraction-alternative}
Assume $\epsilon_n\to0$.  Exactly one of the following two asymptotic branches
occurs.
\begin{enumerate}[label=\textup{(A\arabic*)},leftmargin=3em]
\item \emph{Asymptotic follower fraction:}
      $\|\Kfrac_n\|\to0$.  Equivalently,
      \begin{equation}
       v^*R_nv=o(1)\,v^*U_nv
       \quad\text{uniformly in }v\in D.
       \label{eq:V4-asymptotic-follower}
      \end{equation}
      If $\sup_n\|U_n\|<\infty$, then $\|R_n\|\to0$.
      This is an asymptotic statement.  Exact followerhood at cutoff $n$ is
      still the exact condition $R_n=0$.
\item \emph{Target-normalized survivor:}
      $\limsup_n\|\Kfrac_n\|>0$.  Then there are $\kappa>0$, a subsequence, and
      coefficients $v_n\in PD$ such that
      \begin{equation}
       v_n^*U_nv_n=1,
       \qquad
       v_n^*R_nv_n\ge\kappa,
       \qquad
       v_n^*L_nv_n\ge\frac\kappa2
       \label{eq:V4-survivor}
      \end{equation}
      for all sufficiently large indices in that subsequence.  Hence the
      positive branch is already visible in the finite physical head.
\end{enumerate}
If support identification, target incidence, head exhaustion, or route
comparison fails before this alternative is formed, the output is the
corresponding typed obstruction rather than either branch.
\end{theorem}

\begin{proof}
The two cases $\|\Kfrac_n\|\to0$ and
$\limsup\|\Kfrac_n\|>0$ are exhaustive.  Since
$R_n=U_n^{1/2}\Kfrac_nU_n^{1/2}$ on $P$, the first gives
\eqref{eq:V4-asymptotic-follower} and
$\|R_n\|\le\|U_n\|\|\Kfrac_n\|$.  Exact vanishing of $R_n$ is equivalent to
exact vanishing of $\Kfrac_n$, but convergence to zero does not imply eventual
vanishing.

In the second case choose $\kappa>0$ and a subsequence with
$\|\Kfrac_n\|\ge2\kappa$.  Positivity gives a unit vector $z_n\in PD$ with
$z_n^*\Kfrac_nz_n\ge2\kappa$.  Put
$v_n=U_n^{\dagger/2}z_n$.  Then
$v_n^*U_nv_n=1$ and $v_n^*R_nv_n\ge2\kappa$.  Since
$\|\Kfrac_n-\Hfrac_n\|=\epsilon_n$, the head value is at least
$2\kappa-\epsilon_n\ge\kappa$ eventually.  Renaming $2\kappa$ as $\kappa$
gives the displayed constants.  A failed premise has not entered this
positive-operator dichotomy and is therefore retained with its own type.
\end{proof}

\begin{countertheorem}[Vanishing fraction is not eventual exact followerhood]
\label{cth:V4-asymptotic-not-exact}
Let $D=\C$ and put
\[
 U_n=1,
 \qquad
 R_n=L_n=\frac1n.
\]
Then $\epsilon_n=0$ and $\|\Kfrac_n\|=1/n\to0$, but every finite card has a
strictly positive dynamic residual.  Thus a theorem which turns
$\|\Kfrac_n\|\to0$ into an eventually exact follower verdict is false.
\end{countertheorem}

\begin{theorem}[Physical Duhamel population theorem with no global horizon floor]
\label{thm:V4-corrected-population}
Consider a cofinal family of physical thermal cards
$(H_n,B_n,Y_n,M_{D,n})$ with complete primitive fibre projectors.  Assume:
\begin{enumerate}[label=\textup{(P\arabic*)},leftmargin=3em]
\item the source, target, physical metric, left/right clocks, and primitive
      roles occur on one transported history and satisfy
      \Cref{thm:V4-thermal-rigidity};
\item the target supports are stably identified and the raw insertions obey
      uniformly
      \[
       \Vraw_n^*\Vraw_n\preceq c_VM_{D,n},
       \qquad M_{D,n}\preceq c_TU_n;
      \]
\item $\Lambda_n\to\infty$ and the physical finite head $L_n$ contains every
      residual fibre with left energy at most $\Lambda_n$;
\item the transported row defects, including support and direct/staged route
      defects, tend to zero in the target-normalized metric.
\end{enumerate}
Then
\begin{equation}
 \boxed{
 \epsilon_n
 \le4\max\{\beta/2,1\}^2c_Vc_T(1+\Lambda_n)^{-2}
     +o(1)
 \longrightarrow0.}
 \label{eq:V4-corrected-epsilon}
\end{equation}
The cofinal ancestry verdict is therefore the terminating alternative of
\Cref{thm:V4-cofinal-fraction-alternative}.  No hypothesis
$\inf_n\widehat w_{T,n}>0$ and no replacement of
$m_{\beta,\tau}$ by $m_{-\beta,-\tau}$ is admissible or required.

The physical source reserve remains a separate reported row.  Its collapse is
not a dynamic birth, and a positive dynamic fraction does not repair a failed
null-cost, weight-ray, or cutoff-metric test.
\end{theorem}

\begin{proof}
Before route error, the positive difference $R_n-L_n$ is bounded by the
omitted Duhamel tail.  Apply
\Cref{cor:V4-relative-tail} with $\Lambda=\Lambda_n$.  Transporting the rows to
one target metric adds the stated vanishing defect.  Thus
$\epsilon_n\to0$, and
\Cref{thm:V4-cofinal-fraction-alternative} applies.  The horizon-floor
statement is excluded by \Cref{thm:V4-horizon-trace-budget}; the multiplier
statement is \Cref{thm:V4-thermal-rigidity}.  Reserve is logically separate
because it depends on the physical coefficient metric whereas the cyclic
support and exact dynamic residual do not.
\end{proof}

% =============================================================================
\section{Independent thermally dressed replay of the reversible Ising slab}
\label{sec:V4-Ising-replay}
% =============================================================================

The supplied population report studies an $N$-site Ising ring with
$J=0.4$, $h=0.15$, refresh rate $\rho=0.5$, and generator
\begin{equation}
 L=\rho(R-I)+\sum_i(E_i-I).
 \label{eq:V4-Ising-generator}
\end{equation}
On the centered carrier,
\begin{equation}
 H=-L=\rho I+\sum_i(I-E_i)\succeq\rho I
 \label{eq:V4-Ising-gap}
\end{equation}
exactly, because every $E_i$ is an orthogonal conditional-expectation
projection in $L^2(\mu)$.  The report takes two local action writers, the
identity weight writer, and the global parity line writer.

The accompanying V4 replay retains this model but replaces its raw ancestry
pair by the inherited standard-form pair:
\begin{align}
 B_{\rm wt}&=Z^{-1/2}e^{-\tau H}Ie^{-\beta H/2},\notag\\
 B_{{\rm act},j}&=Z^{-1/2}e^{-\tau H}V_j^\circ e^{-\beta H/2},\notag\\
 B_{\rm line}&=Z^{-1/2}e^{-\tau H}J_{\rm line}e^{-\beta H/2},
 \label{eq:V4-Ising-B}\\
 Y_j&=-\frac\beta2Z^{-1/2}
 \int_0^1e^{-\beta(1-q)H/2}V_j^\circ e^{-\beta qH/2}\,\dd q.
 \label{eq:V4-Ising-Y}
\end{align}
Here $V_j^\circ$ is centered in the Gibbs state of the centered clock, and the
physical coefficient metric is the $L^2(\mu)$ Gram of the four declared scalar
writers, exactly as in the supplied model.  Left-energy degeneracies are
resolved and $Y$ is shorted against the complete four-column primitive bank in
each fibre.  The target parameters are $\beta=1$, $\tau=1/4$, and the displayed
head uses $\Lambda=2$.

\begin{center}
\small
\begin{tabular}{rrrrrr}
\toprule
$N$&$D$&gap&$g_-^{\rm Tr}$&$g_+^{\rm Tr}$&$g_-^{\rm Tr/D}$\\
\midrule
3&7&0.849753&0.142917&0.702925&$2.041678\times10^{-2}$\\
4&15&0.853555&0.178670&0.633380&$1.191134\times10^{-2}$\\
5&31&0.855929&0.165075&0.576976&$5.325015\times10^{-3}$\\
6&63&0.857288&0.162862&0.538167&$2.585117\times10^{-3}$\\
7&127&0.858020&0.135132&0.499109&$1.064034\times10^{-3}$\\
8&255&0.858400&0.110179&0.462438&$4.320727\times10^{-4}$\\
\bottomrule
\end{tabular}
\end{center}

\begin{center}
\small
\begin{tabular}{rrrrrrr}
\toprule
$N$&fibres&head fibres&$\spec\Kfrac_n$&$\spec\Hfrac_n$&tail$/\|U\|$&nonzero fibres\\
\midrule
3&5&3&$[.013486,.015078]$&$[.006692,.011905]$&.005962&5\\
4&11&3&$[.022500,.026928]$&$[.006786,.012127]$&.015206&11\\
5&19&3&$[.025518,.038227]$&$[.004916,.009647]$&.027310&19\\
6&41&6&$[.029350,.044440]$&$[.010819,.015587]$&.026359&41\\
7&73&6&$[.031068,.051760]$&$[.007558,.011432]$&.035739&73\\
8&160&8&$[.032766,.055532]$&$[.006286,.010615]$&.038981&160\\
\bottomrule
\end{tabular}
\end{center}

\begin{assessment}[Exact and numerical content of the replay]
\label{ass:V4-Ising-status}
The generator lower bound \eqref{eq:V4-Ising-gap}, the thermal factorization,
and the weight-ray inequalities are exact.  The table is deterministic
floating-point linear algebra, not outward interval arithmetic.  Its maximum
structural residuals are of order $10^{-13}$ or smaller through $N=8$, but the
strictly positive fibre values are therefore diagnostic rather than certified
inequalities.

Within that scope the replay gives four useful conclusions.
\begin{enumerate}[label=\textup{(R\arabic*)},leftmargin=3em]
\item The thermally dressed unnormalized reserve is roughly $0.11$--$0.18$,
      not the $0.6$--$0.9$ raw-source reserve of the supplied table.
\item Dividing the Hilbert--Schmidt trace by $D$ while retaining the same
      coefficient metric collapses the reserve, exactly as
      \Cref{prop:V4-normalization-scaling,thm:V4-weight-ray} require.
\item The physical target has no exponentially growing ultraviolet factor.
      At the fixed small head $\Lambda=2$, the omitted residual is a few
      percent of $\|U\|$ rather than a calibration-induced exponential
      explosion.
\item Every tested fibre has a numerically nonzero residual and the target-normalized
      dynamic fraction remains at the percent level.  Thus correcting the
      source/target pair does not erase the toy model's evidence for dynamic
      birth; it changes the reserve and tail interpretation.  No convergence
      of these six cards is asserted.
\end{enumerate}
\end{assessment}

The replay source is
\begin{center}\small
\path{Renewal_SM_Einstein_Physical_Thermal_Card_V4.py}
\end{center}
with SHA--256
\begin{center}\small\ttfamily
fa0a0971276d1b49f402be551d3fec4a352dfcaac3cd6776d90642ff534d86df.
\end{center}
Its JSON output has SHA--256
\begin{center}\small\ttfamily
f13f9370a0c143c311cfa06330f878c888ffe61aa5b3e6d0f325dc9ecc2a86d2.
\end{center}

\begin{firewall}[The Ising replay remains a class diagnostic]
\label{fire:V4-Ising-not-ESM}
The model has neither the chiral Dirac--Yukawa coefficient nor the physical
sector-patched determinant/Pfaffian packet, gauge links, Higgs quartic,
BRST/ghost bank, regional current atlas, or Einstein stress insertion.  Its
positive numerical fractions do not populate the original RPESM card.  Its
purpose is to test the ancestry compiler and to expose which conclusions
survive restoration of the inherited thermal source.
\end{firewall}

% =============================================================================
\section{Revised physical acquisition and the Version V5 handoff}
\label{sec:V4-handoff}
% =============================================================================

The supplied population report makes a literal finite evaluation preferable
to another extension of the wall calculus, but only after the physical pair is
restored.  The next original-card acquisition is therefore:
\begin{equation}
 \boxed{
 \begin{gathered}
 (H_n,\text{sector weights},V_{j,n}^\circ,J_{{\rm line},n})\\
 \longrightarrow
 (B_{{\rm wt},n},B_{{\rm act},n},B_{{\rm line},n},Y_{\beta,n})\\
 \longrightarrow
 (M_n,Q_n,U_n,\text{five fibre matrices})\\
 \longrightarrow
 (g_{\pm,n},\Kfrac_n,\Hfrac_{n,\Lambda},
  \ThetaD_{>\Lambda},\text{route defect})\\
 \longrightarrow
 \text{robust survivor, asymptotic follower, or typed obstruction}.
 \end{gathered}}
 \label{eq:V4-next-card}
\end{equation}
The required finite output must retain both unnormalized and physically chosen
Hilbert metrics, state the weight coefficient cost, and pass
\Cref{thm:V4-weight-ray} before reporting a reserve.  It must form $Y$ directly
from the Duhamel derivative, using $m_{\beta,\tau}$ only as an exact quotient
check.  The finite head is then closed by \Cref{thm:V4-physical-tail} or by a
stronger card-native bound.  V3's wall jet may supply an independent positive
annihilator, but it is no longer the only available route to a literal fibre
calculation.

\begin{theorem}[Version V4 export to the main continuum programme]
\label{thm:V4-export}
The main SM--Einstein continuum programme may replace the proposed population
premises
\[
 \inf_n\widehat w_{T,n}>0,
 \qquad
 m_{\beta,\tau}\rightsquigarrow m_{-\beta,-\tau}
\]
by the following exact packet:
\begin{enumerate}[label=\textup{(E\arabic*)},leftmargin=3em]
\item the physical thermally dressed source and the unique Duhamel target;
\item the physical source-action metric and the weight-ray reserve audit;
\item a target-normalized finite Duhamel head;
\item the universal $(1+\Lambda)^{-2}$ omitted-head estimate, or a stronger
      regulator-native tail;
\item target support and direct/staged cutoff transport;
\item the exact finite residual branch, followed by the asymptotic alternative
      of \Cref{thm:V4-cofinal-fraction-alternative}.
\end{enumerate}
This packet is sufficient for the linear ancestry handoff.  It does not open
current, stress, charge, regional locality, Lorentzian, or Einstein rows.
\end{theorem}

\begin{proof}
Items \textup{(E1)--(E2)} are
\Cref{thm:V4-thermal-rigidity,thm:V4-weight-ray}.  Items
\textup{(E3)--(E4)} are
\Cref{thm:V4-physical-tail,cor:V4-relative-tail}.  With item~\textup{(E5)},
\Cref{thm:V4-corrected-population} gives item~\textup{(E6)}.  The downstream
scope is the inherited terminal promotion order and is unchanged by these
linear positive-form identities.
\end{proof}

\begin{assessment}[Version V4 verdict]
\label{ass:V4-verdict}
\passstatus{} for the exact thermal quotient rigidity, the undressed-source
counterexample, the optimal physical divided-difference tail, the horizon
trace-budget theorem, the reserve normalization law, the thermal weight-ray
test, and the corrected target-normalized cofinal alternative.

\obsstatus{} for treating $m_{-\beta,-\tau}$ as a calibration of the inherited
quantum half-source, for reporting a raw identity/multiplication source reserve
as the reserve of the thermally dressed primitive bank, or for promoting
$\|\Kfrac_n\|\to0$ to eventual exact followerhood.  Each obstruction has an
explicit sign, Rayleigh, or scalar counterexample.

\stopstatus{} for the actual Dirac--Yukawa ancestry value until one named
RPESM card exports its finite Hamiltonian, sector weights, centered action and
line writers, physical coefficient metric, and cutoff routes.  The corrected
Ising replay is evidence about the compiler class only.  No current, stress,
Kubo, interacting continuum, Lorentzian reconstruction, Einstein
backreaction, or pure-colour mass conclusion is promoted.
\end{assessment}

\section*{Version V4 append-only record}
\addcontentsline{toc}{section}{Version V4 append-only record}

Preserved the complete V3 source verbatim before its sole terminal
\verb|\end{document}|.  The SHA--256 of the complete V3 input is
\begin{center}\small\ttfamily
9214d8e52b1177ab9788f1899bf8e1bab18de57391c38a154e8fe6bbcf57ca43.
\end{center}

Audited the newly supplied ancestry population report against the inherited
standard-form source formulas.  Proved that its ultraviolet-growing multiplier
is a unique quotient rather than the physical target, and that the proposed
sign-reversed calibration changes the target.  Constructed an exact
same-left-energy counterexample showing opposite physical and raw-surrogate
tail asymptotics.  Proved the optimal one-order physical Duhamel target bound
and its positive omitted-residual corollary.  Replaced the proposed global
horizon floor by an exact trace-budget obstruction and retained the
source-visible Tikhonov tail.  Proved the metric covariance of reserve,
calculated the exact thermal weight-ray bound, and exposed the raw reserve
normalization mismatch.  Corrected the cofinal theorem so that vanishing
fraction means asymptotic, not eventual exact, followerhood.  Replayed the
supplied Ising slab with the physical thermal source and target; retained its
numerical positivity only as a floating-point class diagnostic.

\begin{firewall}[Append-only instruction after V4]
\label{fire:V4-append}
Any Version~V5 must preserve this complete V4 source verbatim, append new
mathematical material immediately before the sole final
\verb|\end{document}|, and use
\begin{center}\small
\path{Renewal_Einstein_Standard_Model_Physical_Source_Metric_Duhamel_Ancestry_and_Quantum_Continuum_Programme_V5.tex}.
\end{center}
It must acquire one original RPESM Hamiltonian/writer card or return the first
missing physical row.  It may not use an undressed source with the inherited
quotient multiplier, reverse the Duhamel sign to obtain a tail, infer reserve
invariance from ambient normalization while holding the coefficient metric
fixed, demand a positive global horizon floor at unbounded cyclic depth,
identify asymptotic smallness with exact followerhood, promote floating-point
Ising values to the Dirac--Yukawa regulator, or reopen current/stress/Einstein
before the linear ancestry packet and its physical routes land.
\end{firewall}

% =============================================================================
% END APPENDED VERSION V4 MATERIAL
% =============================================================================



% =============================================================================
% BEGIN APPENDED VERSION V5 MATERIAL
% =============================================================================
\clearpage
\providecommand{\Trap}{\mathsf T}
\providecommand{\Qsplit}{\mathsf Q}
\providecommand{\Dend}{\mathsf D^{\rm end}}
\providecommand{\Rquad}{\mathscr R^{\rm quad}}
\providecommand{\Rhead}{\mathscr R^{\rm head}}
\providecommand{\LagL}{\mathsf L}
\providecommand{\LagR}{\mathsf R}

\section*{Version V5: mixed-clock Duhamel quadrature, finite card acquisition, and the first missing RPESM row}
\addcontentsline{toc}{section}{Version V5: mixed-clock Duhamel quadrature and finite card acquisition}

\begin{abstract}
Version~V4 restored the physical thermally dressed source and showed that the
calibrated multiplier is only the quotient of the fixed square-root Duhamel
target by that source.  It also removed the unattainable global horizon-floor
premise.  The remaining request was literal acquisition of the selected
RPESM source metric and Duhamel fibre matrices.  The newly supplied population
material proves existence of a large finite same-history packet, but it leaves
the common action, local reversible generator, Dirac--Yukawa coefficient, and
therefore the numerical mixed source--target matrices inside an open class.
This version identifies the smallest finite read which closes that gap to any
prescribed resolution.

Put $a=\beta/2$.  On the commuting left/right standard-form clocks the physical
internal target is
\[
 Y=-\int_0^a e^{-(a-s)H_{\rm L}}e^{-sH_{\rm R}}\Vraw\,\dd s.
\]
On a left-energy fibre $\lambda$, the composite trapezoidal rule with
$h=a/m$ uses only the $m+1$ split-clock vectors
\[
 e^{-(a-jh)\lambda}e^{-jhH^R_\lambda}\Vraw_\lambda,
 \qquad 0\le j\le m.
\]
A scalar monotonicity argument gives the sharp, spectrum-independent estimate
\[
 \left|\int_0^a e^{-(a-s)x-sy}\,\dd s
 -h\!\left[\frac{f(0)+f(a)}2+\sum_{j=1}^{m-1}f(jh)\right]\right|
 \le \frac h2\,|e^{-ay}-e^{-ax}|\le\frac h2.
\]
Thus the quadrature error is paid only by the endpoint thermal mismatch
$e^{-aH_{\rm R}}-e^{-aH_{\rm L}}$ and vanishes exactly on equal-energy blocks.
The constant $h/2$ is sharp without additional bi-energy regularity.  A
uniform fourth bi-energy moment improves the target error to order $h^2$.

For each retained left fibre, one static complete-primitive Gram, $m+1$ mixed
primitive/raw-action lag matrices, and $2m+1$ raw-action lag Grams reconstruct
the quadrature target, its full-primitive Schur residual, and the inherited
five-matrix Duhamel packet.  No right spectral diagonalization, global horizon
inverse, sign-reversed multiplier, or flat reconstruction of the complete
clock is required.  The exact number is $3m+3$ finite matrices per fibre,
plus the already required physical source and target metrics.

Combining the endpoint-sensitive quadrature error with the positive left-energy
tail of Version~V4 gives an explicit target-metric interval around the full
dynamic residual.  A support-stable perturbation theorem adds certified input
radii.  This produces a terminating finite compiler: a strict lower edge gives
a robust dynamic birth; a small upper edge gives quantified followerhood; an
exact follower still requires an exact reducing relation or an exhausting
zero upper sequence.  A vanishing asymptotic fraction is not silently promoted
to exact finite followerhood.

A companion replay on the thermally corrected reversible Ising slab verifies
the finite formulas.  On its largest tested card, the target-normalized residual
error decreases from $5.12\times10^{-3}$ at $m=1$ to
$3.88\times10^{-6}$ at $m=32$.  This is a floating-point test of the acquisition
mechanism, not the missing Dirac--Yukawa value.  The first unpopulated original
RPESM row is now exact: the mixed split-clock matrices between the complete
primitive fibre and the raw centered action insertion, together with their
physical metric and cutoff route.
\end{abstract}

% =============================================================================
\section{Why the next acquisition is a mixed-clock packet rather than another horizon theorem}
\label{sec:V5-direction}
% =============================================================================

The source programme has reached a finite physical-value boundary.  The line
and clock compiler has already reduced the dynamic question to a complete
primitive fibre, a canonical Duhamel target, five finite matrices, and a
positive left-energy tail.  Version~V4 further showed that the target must be
formed before the multiplier is interpreted and that a global minimum
horizon eigenvalue is neither available nor required.  The remaining problem
is therefore not to enlarge the Hilbert-space compiler.  It is to acquire the
mixed matrices of one named physical card.

The RPESM population theorem supplies a literal event law on which action,
line, source, current, stress, and held-out writers occur.  Its cutoff
definition nevertheless chooses a finite-range common action, massive line
sections, local conditional-expectation generators, and nonzero mixed
coefficients from a nonempty open set.  No displayed table fixes the numerical
Dirac--Yukawa coefficient or the resulting standard-form matrices.  Existence
of a writer on a cylinder and evaluation of its mixed thermal Gram are
therefore different statements.

The present version exploits the integral defining the already fixed target.
It asks for finitely many physical split-clock Grams and returns an explicit
outward interval.  This is weaker than reconstruction of the complete
source-cyclic clock and stronger than a diagonal source or target norm.

\begin{firewall}[Rows retained and rows not reopened]
\label{fire:V5-scope}
Version~V5 does not reprove the sector-patched line, the complete mixed
entrance short, the semigroup-cyclic ancestry theorem, the five-matrix
Pythagoras, the material wall jet, or the physical left-energy tail.  It
constructs a finite acquisition of the canonical Duhamel target from the
raw centered action insertion and the already populated left/right clocks.
Current, contact, stress, charge, regional locality, determinant/Pfaffian
orientation, Lorentzian reconstruction, and Einstein backreaction remain
closed.
\end{firewall}

% =============================================================================
\section{The physical target is a split-clock integral}
\label{sec:V5-split-integral}
% =============================================================================

Let $\mathcal K$ be one finite-cutoff standard-form Hilbert space.  Let
\[
 H_{\rm L}=H_{\rm L}^*\succeq0,
 \qquad
 H_{\rm R}=H_{\rm R}^*\succeq0,
 \qquad
 [H_{\rm L},H_{\rm R}]=0,
\]
and let
\[
 \Vraw:E_Q\longrightarrow\mathcal K
\]
be the partition-normalized raw centered action-insertion synthesis.  Write
$a=\beta/2>0$.  The physical square-root target of Version~V4 is
\begin{equation}
 \boxed{
 Y=-\int_0^a
 e^{-(a-s)H_{\rm L}}e^{-sH_{\rm R}}\Vraw\,\dd s.}
 \label{eq:V5-split-integral}
\end{equation}
The Bochner integral exists at finite cutoff.  The same formula defines a
bounded synthesis for arbitrary nonnegative self-adjoint commuting clocks,
because the integrand is a contraction times the finite-rank map $\Vraw$.

Let $E_\lambda^{\rm L}$ be a left-energy projection.  Put
\[
 \Vraw_\lambda=E_\lambda^{\rm L}\Vraw,
 \qquad
 H_\lambda^R=H_{\rm R}|_{E_\lambda^{\rm L}\mathcal K}.
\]
Then
\begin{equation}
 \boxed{
 Y_\lambda=-\int_0^a
 e^{-(a-s)\lambda}e^{-sH_\lambda^R}\Vraw_\lambda\,\dd s.}
 \label{eq:V5-fibre-integral}
\end{equation}
This is exactly the direct divided-difference target; the source delay $\tau$
does not occur.

\begin{lemma}[Sharp scalar split-clock trapezoidal bound]
\label{lem:V5-trapezoid}
Let $x,y\ge0$, $m\ge1$, $h=a/m$, $s_j=jh$, and
\[
 f_{x,y}(s)=e^{-(a-s)x-sy}.
\]
Define the composite trapezoidal value
\begin{equation}
 \Trap_m(x,y)
 =h\left[
 \frac12 f_{x,y}(0)
 +\sum_{j=1}^{m-1}f_{x,y}(s_j)
 +\frac12f_{x,y}(a)
 \right].
 \label{eq:V5-trapezoid-value}
\end{equation}
Then
\begin{equation}
 \boxed{
 0\le
 \Trap_m(x,y)-\int_0^af_{x,y}(s)\,\dd s
 \le\frac h2\,|e^{-ay}-e^{-ax}|
 \le\frac h2.}
 \label{eq:V5-trapezoid-bound}
\end{equation}
The uniform constant $h/2$ is sharp:
\begin{equation}
 \sup_{x,y\ge0}
 \left|
 \Trap_m(x,y)-\int_0^af_{x,y}(s)\,\dd s
 \right|=\frac h2.
 \label{eq:V5-trapezoid-sharp}
\end{equation}
\end{lemma}

\begin{proof}
The function $f_{x,y}$ is positive, monotone, and convex on $[0,a]$:
\[
 f'_{x,y}=(x-y)f_{x,y},
 \qquad
 f''_{x,y}=(x-y)^2f_{x,y}\ge0.
\]
Let $L_m$ and $R_m$ be the left and right rectangular sums.  Monotonicity
places the integral between them, while
\[
 |R_m-L_m|
 =h|f_{x,y}(a)-f_{x,y}(0)|
 =h|e^{-ay}-e^{-ax}|.
\]
The trapezoidal value is $(L_m+R_m)/2$, so its distance from any point between
$L_m$ and $R_m$ is at most half their separation.  Convexity gives the first
sign in \eqref{eq:V5-trapezoid-bound}.  Both endpoint values lie in $[0,1]$.
For sharpness take $x=0$ and let $y\to\infty$.  The integral tends to zero,
all positive-node terms vanish, and the half-weight at $s=0$ tends to $h/2$.
\end{proof}

For $0\le j\le m$ put
\begin{equation}
 \omega_{m,j}=
 \begin{cases}
 1/2,&j=0\text{ or }j=m,\\
 1,&1\le j\le m-1,
 \end{cases}
 \qquad
 Q_{m,j}=e^{-(a-s_j)H_{\rm L}}e^{-s_jH_{\rm R}}.
 \label{eq:V5-Qmj}
\end{equation}
Define the finite physical quadrature target
\begin{equation}
 \boxed{
 Y^{(m)}=-h\sum_{j=0}^{m}\omega_{m,j}Q_{m,j}\Vraw.}
 \label{eq:V5-Ym-global}
\end{equation}
Its left-fibre restriction is denoted $Y_\lambda^{(m)}$.

\begin{theorem}[Endpoint-sensitive finite target approximation]
\label{thm:V5-target-approximation}
Let
\begin{equation}
 \boxed{
 \Dend_a
 :=\Vraw^*
 \bigl(e^{-aH_{\rm R}}-e^{-aH_{\rm L}}\bigr)^2
 \Vraw\succeq0.}
 \label{eq:V5-endpoint-discrepancy}
\end{equation}
Then
\begin{equation}
 \boxed{
 (Y-Y^{(m)})^*(Y-Y^{(m)})
 \preceq\frac{h^2}{4}\Dend_a
 \preceq\frac{h^2}{4}\Vraw^*\Vraw.}
 \label{eq:V5-target-error-form}
\end{equation}
The same inequalities hold after compression to any left-energy head or one
left fibre.  In particular,
\begin{equation}
 \|Y-Y^{(m)}\|
 \le\frac h2\|\Dend_a\|^{1/2}
 \le\frac h2\|\Vraw\|.
 \label{eq:V5-target-error-norm}
\end{equation}
The endpoint discrepancy vanishes exactly on raw insertion components
supported on equal left and right energies.
\end{theorem}

\begin{proof}
Use the joint spectral theorem for the commuting pair
$(H_{\rm L},H_{\rm R})$.  The scalar multiplier of $Y-Y^{(m)}$ is the negative
of the error in \eqref{eq:V5-trapezoid-bound}.  Its square is bounded pointwise
by
\[
 \frac{h^2}{4}(e^{-ay}-e^{-ax})^2.
\]
Functional calculus and compression by $\Vraw$ give the first inequality.
Since $e^{-aH_{\rm L}}$ and $e^{-aH_{\rm R}}$ are commuting positive
contractions, their scalar difference has absolute value at most one, giving
the second inequality.  The zero statement is the zero set of the scalar
endpoint difference on the support of the raw insertion.
\end{proof}

\begin{corollary}[Second-order branch under a bi-energy curvature moment]
\label{cor:V5-second-order}
Put $\Delta_H=H_{\rm L}-H_{\rm R}$.  If
$\Delta_H^2\Vraw$ is bounded, then
\begin{equation}
 \boxed{
 \|Y-Y^{(m)}\|
 \le\frac{ah^2}{12}\|\Delta_H^2\Vraw\|.}
 \label{eq:V5-second-order-bound}
\end{equation}
Equivalently,
\begin{equation}
 (Y-Y^{(m)})^*(Y-Y^{(m)})
 \preceq\frac{a^2h^4}{144}
 \Vraw^*\Delta_H^4\Vraw.
 \label{eq:V5-second-order-form}
\end{equation}
Thus an independently bounded fourth bi-energy moment upgrades the uniform
$O(m^{-1})$ target estimate to $O(m^{-2})$.
\end{corollary}

\begin{proof}
For a twice continuously differentiable scalar function, the composite
trapezoidal error on an interval of length $a$ and mesh $h$ is bounded by
$ah^2\sup|f''|/12$.  Here
$f''_{x,y}(s)=(x-y)^2f_{x,y}(s)$ and $0\le f_{x,y}\le1$.  Hence the joint
spectral multiplier error is bounded by
$(ah^2/12)|x-y|^2$.  Functional calculus gives
\eqref{eq:V5-second-order-bound}; squaring gives the form statement.
\end{proof}

\begin{countertheorem}[No spectrum-free improvement of the endpoint bound]
\label{cth:V5-uniform-rate-sharp}
Without an additional bi-energy moment there is no constant $c<1/2$ such that
\[
 \|Y-Y^{(m)}\|\le ch\|\Vraw\|
\]
for every nonnegative commuting clock pair and every bounded raw insertion.
\end{countertheorem}

\begin{proof}
Take one left energy $x=0$, one right energy $y=R$, and a unit insertion on
that bi-energy block.  By the sharpness part of
\Cref{lem:V5-trapezoid}, the ratio of the target error to $h\|\Vraw\|$ tends
to $1/2$ as $R\to\infty$.
\end{proof}

% =============================================================================
\section{The finite mixed-clock matrix card}
\label{sec:V5-finite-card}
% =============================================================================

Fix a retained left energy $\lambda$.  Let
\[
 B_\lambda:E_B\longrightarrow\mathcal K_\lambda
\]
be the complete prior primitive fibre synthesis, including sector weight,
common action, line current, and every independently occurring primitive.
Put
\[
 S_\lambda=B_\lambda^*B_\lambda,
 \qquad
 P_\lambda^B=B_\lambda S_\lambda^\dagger B_\lambda^*.
\]
The exact fibre residual is
\[
 \mathscr R_\lambda
 =Y_\lambda^*(I-P_\lambda^B)Y_\lambda.
\]

For $0\le j\le m$ and $0\le n\le2m$, define the finite mixed-clock matrices
\begin{align}
 C_{\lambda,j}^{(m)}
 &:=B_\lambda^*
 e^{-(a-jh)\lambda}e^{-jhH_\lambda^R}\Vraw_\lambda,
 \label{eq:V5-C-lag}\\
 K_{\lambda,n}^{(m)}
 &:=\Vraw_\lambda^*
 e^{-(2a-nh)\lambda}e^{-nhH_\lambda^R}\Vraw_\lambda.
 \label{eq:V5-K-lag}
\end{align}
Every exponent in \eqref{eq:V5-K-lag} is nonnegative.  Let
\begin{equation}
 \gamma_{m,n}
 :=\sum_{\substack{0\le i,j\le m\\i+j=n}}
 \omega_{m,i}\omega_{m,j},
 \qquad 0\le n\le2m.
 \label{eq:V5-gamma}
\end{equation}
These are fixed positive rational numbers determined by $m$.

\begin{definition}[Mixed-clock Duhamel acquisition packet]
\label{def:V5-acquisition-packet}
At mesh depth $m$, the fibre packet is
\begin{equation}
 \boxed{
 \mathfrak D_{\lambda,m}
 =\left(
 S_\lambda;
 (C_{\lambda,j}^{(m)})_{j=0}^{m};
 (K_{\lambda,n}^{(m)})_{n=0}^{2m}
 \right).}
 \label{eq:V5-packet}
\end{equation}
It contains exactly $3m+3$ finite matrices.  Source and target coefficient
metrics, support projections, uncertainty radii, and cutoff routes are carried
as typed metadata and are not included in this count.
\end{definition}

\begin{theorem}[Exact reconstruction of the quadrature residual]
\label{thm:V5-card-reconstruction}
From \eqref{eq:V5-packet} define
\begin{align}
 T_{\lambda,m}
 &:=-h\sum_{j=0}^{m}\omega_{m,j}C_{\lambda,j}^{(m)},
 \label{eq:V5-Tm}\\
 U_{\lambda,m}
 &:=h^2\sum_{n=0}^{2m}\gamma_{m,n}K_{\lambda,n}^{(m)},
 \label{eq:V5-Um}\\
 \Rquad_{\lambda,m}
 &:=U_{\lambda,m}
 -T_{\lambda,m}^*S_\lambda^\dagger T_{\lambda,m}.
 \label{eq:V5-Rm}
\end{align}
Then
\begin{equation}
 \boxed{
 T_{\lambda,m}=B_\lambda^*Y_\lambda^{(m)},
 \qquad
 U_{\lambda,m}=(Y_\lambda^{(m)})^*Y_\lambda^{(m)},
 \qquad
 \Rquad_{\lambda,m}
 =(Y_\lambda^{(m)})^*(I-P_\lambda^B)Y_\lambda^{(m)}\succeq0.}
 \label{eq:V5-reconstruction-identities}
\end{equation}
No right-energy spectral decomposition and no full one-lag flatness test is
needed.
\end{theorem}

\begin{proof}
The first identity is the definition of the quadrature target.  For the
second, the commuting semigroups give
\[
 \left(e^{-(a-ih)\lambda}e^{-ihH_\lambda^R}\right)^*
 \left(e^{-(a-jh)\lambda}e^{-jhH_\lambda^R}\right)
 =e^{-(2a-(i+j)h)\lambda}e^{-(i+j)hH_\lambda^R}.
\]
Grouping the double sum by $n=i+j$ yields \eqref{eq:V5-Um}.  The
Moore--Penrose range formula makes $P_\lambda^B$ the orthogonal projection
onto $\Ran B_\lambda$, and expanding the last norm gives
\eqref{eq:V5-Rm}.
\end{proof}

\begin{corollary}[Direct population of the inherited five matrices]
\label{cor:V5-five-matrix-population}
Write the complete primitive fibre as $B_\lambda=(A_\lambda,F_\lambda)$ and
put
\[
 S^A_\lambda=A_\lambda^*A_\lambda,
 \qquad
 D_\lambda=A_\lambda^*F_\lambda,
 \qquad
 E_\lambda=F_\lambda^*F_\lambda.
\]
Split $T_{\lambda,m}$ into its action and auxiliary rows,
\[
 T_{\lambda,m}
 =\binom{T^A_{\lambda,m}}{T^F_{\lambda,m}}.
\]
Then the five-matrix packet of the line/clock compiler for the quadrature
target is
\begin{equation}
 \boxed{
 \begin{aligned}
 S&=S^A_\lambda,
 &T&=T^A_{\lambda,m},
 &U&=U_{\lambda,m},\\
 G&=E_\lambda-D_\lambda^*(S^A_\lambda)^\dagger D_\lambda,
 &C&=T^F_{\lambda,m}
      -D_\lambda^*(S^A_\lambda)^\dagger T^A_{\lambda,m}.
 \end{aligned}}
 \label{eq:V5-five-matrices}
\end{equation}
Consequently
\begin{equation}
 \boxed{
 \Rquad_{\lambda,m}
 =U-T^*S^\dagger T-C^*G^\dagger C.}
 \label{eq:V5-five-matrix-residual}
\end{equation}
Thus the mixed-clock packet populates the existing compiler rather than
replacing it.
\end{corollary}

\begin{proof}
The residualized auxiliary synthesis is
$N=(I-P_A)F$.  Taking Grams gives
$N^*N=E-D^*(S^A)^\dagger D$.  Likewise
\[
 N^*Y^{(m)}
 =F^*Y^{(m)}-F^*A(S^A)^\dagger A^*Y^{(m)}
 =T^F-D^*(S^A)^\dagger T^A.
\]
The orthogonal decomposition
$\Ran(A,F)=\Ran A\oplus\Ran N$ gives
\eqref{eq:V5-five-matrix-residual}.
\end{proof}

\begin{countertheorem}[The auxiliary mixed row is irredundant even on a canonical Duhamel fibre]
\label{cth:V5-C-irredundant}
Let $M=\diag(\mu_1,\mu_2,\mu_3)$ on $\C^3$ with
$\mu_1\ne\mu_2$, and put
\[
 A1=2^{-1/2}(e_1+e_2),
 \qquad
 Y=MA.
\]
Consider the two unit auxiliary primitives
\[
 F_01=2^{-1/2}(e_1-e_2),
 \qquad
 F_11=e_3.
\]
For both cards the first four matrices $(S,T,U,G)$ are identical:
\begin{equation}
 S=1,
 \quad
 T=\frac{\mu_1+\mu_2}{2},
 \quad
 U=\frac{\mu_1^2+\mu_2^2}{2},
 \quad
 G=1.
 \label{eq:V5-same-four}
\end{equation}
Their mixed rows and full residuals are
\begin{equation}
 \boxed{
 C_0=\frac{\mu_1-\mu_2}{2},
 \quad \mathscr R_0=0;
 \qquad
 C_1=0,
 \quad \mathscr R_1=\frac{(\mu_1-\mu_2)^2}{4}>0.}
 \label{eq:V5-opposite-C}
\end{equation}
Taking $M=f_\lambda(H_\lambda^R)$ for three right energies makes this a
canonical Duhamel fibre.  Hence source and target self-Grams, action variance,
and the auxiliary norm do not determine ancestry; the mixed primitive--target
row must be physically acquired.
\end{countertheorem}

\begin{proof}
The vectors $A1,F_01$ are an orthonormal basis of
$\Span\{e_1,e_2\}$, whereas $F_11$ is orthogonal to that plane.  Direct
calculation gives \eqref{eq:V5-same-four}.  In the first card the entire
component of $Y$ orthogonal to $A1$ is carried by $F_01$; in the second it is
orthogonal to both primitive columns.  The displayed values follow from the
projection Pythagoras.  The calibrated multiplier is a real strictly monotone
function of right energy, so distinct energies may be chosen with values
$\mu_1\ne\mu_2$.
\end{proof}

% =============================================================================
\section{A certified target-metric interval for the full dynamic residual}
\label{sec:V5-certified-interval}
% =============================================================================

Let $N=N^*\succeq0$ be the physical target coefficient metric and assume the
raw target incidence is null-cost consistent,
\begin{equation}
 \Ker N\subseteq\Ker\Vraw.
 \label{eq:V5-target-null-cost}
\end{equation}
All following inequalities are read on $P_N E_Q$, where
$P_N=NN^\dagger$.  Fix a finite left-energy head $\lambda\le\Lambda$ and put
\begin{equation}
 \Rquad_{\Lambda,m}
 :=\sum_{\lambda\le\Lambda}\Rquad_{\lambda,m}.
 \label{eq:V5-total-quadrature-head}
\end{equation}
Let
\begin{align}
 \Dend_{a,\Lambda}
 &:=\sum_{\lambda\le\Lambda}
 \Vraw_\lambda^*(e^{-aH_\lambda^R}-e^{-a\lambda})^2\Vraw_\lambda,
 \label{eq:V5-head-endpoint}\\
 e_{a,\Lambda;N}
 &:=\left\|
 N^{\dagger/2}\Dend_{a,\Lambda}N^{\dagger/2}
 \right\|,
 \label{eq:V5-e-endpoint}\\
 u_{\Lambda,m;N}
 &:=\left\|
 N^{\dagger/2}
 \left(\sum_{\lambda\le\Lambda}U_{\lambda,m}\right)
 N^{\dagger/2}
 \right\|^{1/2},
 \label{eq:V5-u-m}
\end{align}
and
\begin{equation}
 \eta_{\Lambda,m;N}
 :=\frac h2\sqrt{e_{a,\Lambda;N}},
 \qquad
 \delta^{\rm quad}_{\Lambda,m;N}
 :=2u_{\Lambda,m;N}\eta_{\Lambda,m;N}
   +\eta_{\Lambda,m;N}^2.
 \label{eq:V5-delta-quad}
\end{equation}
The endpoint matrix is itself part of the acquired $K$-bank:
\begin{equation}
 \boxed{
 \Dend_{a,\lambda}
 =K_{\lambda,2m}^{(m)}-2K_{\lambda,m}^{(m)}
  +K_{\lambda,0}^{(m)}.}
 \label{eq:V5-endpoint-from-K}
\end{equation}

\begin{theorem}[Quadrature interval on a fixed left-energy head]
\label{thm:V5-head-interval}
Let
\[
 \Rhead_\Lambda
 =\sum_{\lambda\le\Lambda}
 Y_\lambda^*(I-P_\lambda^B)Y_\lambda.
\]
Then
\begin{equation}
 \boxed{
 -\delta^{\rm quad}_{\Lambda,m;N}N
 \preceq
 \Rhead_\Lambda-\Rquad_{\Lambda,m}
 \preceq
 \delta^{\rm quad}_{\Lambda,m;N}N.}
 \label{eq:V5-head-interval}
\end{equation}
A universal fallback is obtained from
$\Vraw^*\Vraw\preceq c_VN$:
\begin{equation}
 \boxed{
 \delta^{\rm quad}_{\Lambda,m;N}
 \le \left(ah+\frac{h^2}{4}\right)c_V.}
 \label{eq:V5-universal-delta}
\end{equation}
\end{theorem}

\begin{proof}
On the orthogonal direct sum of retained left fibres let
$Q_\Lambda=\bigoplus_{\lambda\le\Lambda}(I-P_\lambda^B)$,
$Y_\Lambda=\bigoplus Y_\lambda$, and
$Y_{\Lambda,m}=\bigoplus Y_\lambda^{(m)}$.  Put
$E=Y_\Lambda-Y_{\Lambda,m}$.  Then
\[
 \Rhead_\Lambda-\Rquad_{\Lambda,m}
 =Y_{\Lambda,m}^*Q_\Lambda E
  +E^*Q_\Lambda Y_{\Lambda,m}+E^*Q_\Lambda E.
\]
After target whitening, \Cref{thm:V5-target-approximation} gives
$\|EN^{\dagger/2}\|\le\eta_{\Lambda,m;N}$, while
$\|Y_{\Lambda,m}N^{\dagger/2}\|=u_{\Lambda,m;N}$.
The norm of the displayed difference is therefore at most
$2u\eta+\eta^2$.  A self-adjoint form of target-normalized norm at most
$\delta$ lies between $-\delta N$ and $\delta N$ on the supported quotient.

For the fallback, $\Dend_{a,\Lambda}\preceq\Vraw^*\Vraw\preceq c_VN$, so
$\eta\le(h/2)\sqrt{c_V}$.  Also
$\|Y_\Lambda N^{\dagger/2}\|\le a\sqrt{c_V}$ and hence the direct estimate
$\|R(Y_\Lambda)-R(Y_{\Lambda,m})\|
\le\eta(2a\sqrt{c_V}+\eta)$ gives
\eqref{eq:V5-universal-delta}.
\end{proof}

Let $C_a=2\max\{a,1\}$.  Under the physical incidence bound
\begin{equation}
 \Vraw^*\Vraw\preceq c_VN,
 \label{eq:V5-incidence-N}
\end{equation}
Version~V4 gives the positive omitted-left-energy tail bound
\begin{equation}
 0\preceq\ThetaD_{>\Lambda}
 \preceq C_a^2c_V(1+\Lambda)^{-2}N.
 \label{eq:V5-left-tail-N}
\end{equation}
Put
\begin{equation}
 \delta^{\rm tail}_{\Lambda;N}
 :=C_a^2c_V(1+\Lambda)^{-2}.
 \label{eq:V5-delta-tail}
\end{equation}

\begin{theorem}[Finite mixed-clock interval for the full dynamic residual]
\label{thm:V5-full-interval}
The exact dynamic residual satisfies
\begin{equation}
 \boxed{
 \Rquad_{\Lambda,m}
 -\delta^{\rm quad}_{\Lambda,m;N}N
 \preceq R_{\rm dyn}
 \preceq
 \Rquad_{\Lambda,m}
 +\bigl(
   \delta^{\rm quad}_{\Lambda,m;N}
   +\delta^{\rm tail}_{\Lambda;N}
  \bigr)N.}
 \label{eq:V5-full-interval}
\end{equation}
Consequently:
\begin{enumerate}[label=\textup{(I\arabic*)},leftmargin=3em]
\item if
      $\Rquad_{\Lambda,m}\succeq
      (\delta^{\rm quad}_{\Lambda,m;N}+c)N$ for some $c>0$, then
      $R_{\rm dyn}\succeq cN$;
\item if
      $\Rquad_{\Lambda,m}\preceq\varepsilon N$, then
      $R_{\rm dyn}\preceq
      (\varepsilon+\delta^{\rm quad}_{\Lambda,m;N}
      +\delta^{\rm tail}_{\Lambda;N})N$;
\item an exact zero branch follows if there is an exhausting sequence
      $(\Lambda_k,m_k)$ for which the upper coefficient in item~\textup{(I2)}
      tends to zero, or from an independent exact reducing relation.
\end{enumerate}
No compact-horizon minimum eigenvalue occurs.
\end{theorem}

\begin{proof}
Write
$R_{\rm dyn}=\Rhead_\Lambda+\ThetaD_{>\Lambda}$.
Combine \Cref{thm:V5-head-interval} with
$0\preceq\ThetaD_{>\Lambda}\preceq
\delta^{\rm tail}_{\Lambda;N}N$.  The three consequences are immediate from
Loewner order.  In item~\textup{(I3)}, the fixed positive operator
$R_{\rm dyn}$ is squeezed to zero in target-normalized norm.
\end{proof}

% =============================================================================
\section{Input uncertainty and a domain-safe finite verdict}
\label{sec:V5-uncertainty}
% =============================================================================

The preceding interval assumes exact finite matrices.  The next statement
adds outward radii without confusing support loss with a negative verdict.
Let two finite Schur packets be
\[
 (S_i,T_i,U_i),
 \qquad
 R_i=U_i-T_i^*S_i^\dagger T_i,
 \qquad i=1,2.
\]
Assume their source supports have been aligned to one projection $P$ and
\begin{equation}
 S_i\succeq s_*P,
 \qquad s_*>0.
 \label{eq:V5-source-floor}
\end{equation}

\begin{proposition}[Supported Schur perturbation bound]
\label{prop:V5-Schur-perturbation}
Let $N$ be the target metric and put
\begin{align}
 \varepsilon_S&=\|S_2-S_1\|,
 \notag\\
 \varepsilon_T&=\|(T_2-T_1)N^{\dagger/2}\|,
 \notag\\
 \varepsilon_U&=\|N^{\dagger/2}(U_2-U_1)N^{\dagger/2}\|,
 \label{eq:V5-input-radii}\\
 t_i&=\|T_iN^{\dagger/2}\|.
 \notag
\end{align}
Then
\begin{equation}
 \boxed{
 \|N^{\dagger/2}(R_2-R_1)N^{\dagger/2}\|
 \le
 \varepsilon_U
 +s_*^{-1}(t_1+t_2)\varepsilon_T
 +s_*^{-2}t_1t_2\varepsilon_S.}
 \label{eq:V5-Schur-perturbation}
\end{equation}
\end{proposition}

\begin{proof}
On the common support,
\[
 S_2^\dagger-S_1^\dagger
 =-S_2^\dagger(S_2-S_1)S_1^\dagger.
\]
Insert and subtract the mixed terms to obtain
\[
 \begin{aligned}
 R_2-R_1={}&U_2-U_1
 -(T_2-T_1)^*S_2^\dagger T_2
 -T_1^*S_1^\dagger(T_2-T_1)\\
 &+T_1^*S_2^\dagger(S_2-S_1)S_1^\dagger T_2.
 \end{aligned}
\]
Congruence by $N^{\dagger/2}$, the triangle inequality, and
$\|S_i^\dagger\|\le s_*^{-1}$ prove the claim.
\end{proof}

For submitted lag matrices let $\widetilde T_{\Lambda,m}$ and
$\widetilde U_{\Lambda,m}$ be formed by the exact linear formulas
\eqref{eq:V5-Tm}--\eqref{eq:V5-Um}.  Their radii may be assembled directly:
\begin{align}
 \varepsilon_T^{\rm lag}
 &\le h\sum_{\lambda\le\Lambda}
       \sum_{j=0}^{m}\omega_{m,j}
       \|\Delta C_{\lambda,j}N^{\dagger/2}\|,
 \label{eq:V5-lag-T-radius}\\
 \varepsilon_U^{\rm lag}
 &\le h^2\sum_{\lambda\le\Lambda}
       \sum_{n=0}^{2m}\gamma_{m,n}
       \|N^{\dagger/2}\Delta K_{\lambda,n}N^{\dagger/2}\|.
 \label{eq:V5-lag-U-radius}
\end{align}
Applying \Cref{prop:V5-Schur-perturbation} gives a total data radius
$\delta^{\rm data}_{\Lambda,m;N}$.

\begin{theorem}[Domain-safe finite mixed-clock compiler]
\label{thm:V5-compiler}
A submitted packet consists of:
\begin{enumerate}[label=\textup{(C\arabic*)},leftmargin=3em]
\item one named physical card, inverse temperature, left-energy head, complete
      primitive fibre bank, and raw centered action insertion;
\item physical source and target coefficient metrics with their null-cost
      tests;
\item the matrices \eqref{eq:V5-packet}, their canonical supports, and
      outward radii;
\item a common supported source floor $s_*$;
\item the incidence constant $c_V$, left-tail cutoff, and direct/staged route
      labels.
\end{enumerate}
If a metric, support, positivity, same-card incidence, or route equation fails,
the compiler returns the corresponding typed input or physical obstruction.
On a passing packet put
\[
 \delta_-:=\delta^{\rm quad}_{\Lambda,m;N}
             +\delta^{\rm data}_{\Lambda,m;N},
 \qquad
 \delta_+:=\delta_-+\delta^{\rm tail}_{\Lambda;N}.
\]
Let
\[
 \widetilde K_{\Lambda,m}
 =N^{\dagger/2}\widetilde\Rquad_{\Lambda,m}N^{\dagger/2}.
\]
Then:
\begin{enumerate}[label=\textup{(V\arabic*)},leftmargin=3em]
\item if, on a declared protected target screen $P$, the compressed
      matrix satisfies
      \[
       P\widetilde K_{\Lambda,m}P
       \succeq(\delta_-+c)P
      \]
      for some $c>0$, the compiler returns a robust dynamic birth with margin
      $c$ on that screen;
\item if $\|\widetilde K_{\Lambda,m}\|+\delta_+\le\varepsilon$, it returns
      an $\varepsilon$-follower certificate;
\item an exact follower is returned only from an exact reducing relation or an
      exhausting validated sequence with upper edge tending to zero;
\item otherwise the result is \emph{unresolved at the submitted resolution},
      not a negative theorem verdict.
\end{enumerate}
\end{theorem}

\begin{proof}
The support and metric checks are necessary for the Moore--Penrose Schur and
target whitening operations to have their declared meaning.  On the passing
branch, \Cref{prop:V5-Schur-perturbation} replaces
$\Rquad_{\Lambda,m}$ in \eqref{eq:V5-full-interval} by its submitted value and
adds $\delta^{\rm data}$ to both edges.  Compression to a protected target
screen and the min--max principle prove item~\textup{(V1)}.  The upper edge
proves item~\textup{(V2)}.  Item~\textup{(V3)} is
\Cref{thm:V5-full-interval}(I3), and item~\textup{(V4)} records that failure of
a sufficient strict inequality is not a dual obstruction.
\end{proof}

% =============================================================================
\section{Cofinal transport without a global horizon floor}
\label{sec:V5-cofinal}
% =============================================================================

\begin{theorem}[Cofinal landing from finite mixed-clock packets]
\label{thm:V5-cofinal}
Let $n\mapsto\mathcal C_n$ be a cofinal family of physical Duhamel cards, with
all target coefficient spaces transported to one fixed finite screen.  Let
$N_n$ be their physical target metrics and assume:
\begin{enumerate}[label=\textup{(H\arabic*)},leftmargin=3em]
\item $a_n=\beta_n/2$ stays in a compact subset of $(0,\infty)$;
\item $\Vraw_n^*\Vraw_n\preceq c_VN_n$ with one finite $c_V$;
\item $\Lambda_n\to\infty$, $m_n\to\infty$, and the V4 physical left-tail
      bound is uniform;
\item the submitted source supports stabilize, their supported floors are
      uniform, and the accumulated finite-data radii tend to zero;
\item after direct/staged transport, the reconstructed normalized matrices
      \[
       \widetilde K_n
       =N_n^{\dagger/2}
         \widetilde\Rquad_{\Lambda_n,m_n,n}
         N_n^{\dagger/2}
      \]
      converge in norm to a positive matrix $K_\infty$.
\end{enumerate}
Then the exact target-normalized dynamic residuals converge to the same limit:
\begin{equation}
 \boxed{
 N_n^{\dagger/2}R_{{\rm dyn},n}N_n^{\dagger/2}
 \longrightarrow K_\infty.}
 \label{eq:V5-cofinal-limit}
\end{equation}
If $K_\infty\ne0$, every sufficiently fine card has a target-normalized
survivor on a convergent protected direction.  If $K_\infty=0$, the family is
an asymptotic follower in target fraction.  The latter conclusion does not
assert exact followerhood at any finite cutoff.
\end{theorem}

\begin{proof}
The universal quadrature coefficient satisfies
\[
 \delta^{\rm quad}_{n}
 \le\left(a_nh_n+\frac{h_n^2}{4}\right)c_V
 \longrightarrow0,
 \qquad h_n=a_n/m_n.
\]
The left-tail coefficient tends to zero because $\Lambda_n\to\infty$, and the
finite-data radius tends to zero by hypothesis.  The interval of
\Cref{thm:V5-compiler} therefore shows that the exact normalized residual and
$\widetilde K_n$ differ in norm by $o(1)$.  This proves
\eqref{eq:V5-cofinal-limit}.  If the limit is nonzero, choose a unit eigenvector
of a positive eigenvalue and transport it back; norm convergence preserves a
positive Rayleigh quotient.  If the limit is zero, the normalized residual
norm tends to zero.  Exact finite zero is a separate algebraic statement.
\end{proof}

\begin{corollary}[No cyclic-depth hypothesis]
\label{cor:V5-no-depth}
The conclusion of \Cref{thm:V5-cofinal} requires neither a lower bound on the
smallest positive eigenvalue of a compact-horizon Gramian nor bounded cyclic
depth per source rank.  Increasing cyclic depth can make complete clock
reconstruction ill conditioned while the finite target-specific quadrature
remains uniformly accurate.
\end{corollary}

\begin{proof}
The proof uses only contraction of the split semigroups, the physical
source/target incidence window, the left-energy tail, finite support stability,
and transport of the listed mixed matrices.  No inverse of a horizon Gramian
appears.
\end{proof}

% =============================================================================
\section{A finite replay of the acquisition mechanism}
\label{sec:V5-replay}
% =============================================================================

The V4 reversible Ising diagnostic was rerun with the same thermally dressed
source and direct Duhamel target, but the target was reconstructed only from
the finite trapezoidal mixed-clock packet.  The following table reports
\[
 \left\|
 U^{-1/2}(R_{\rm dyn}-\Rquad_m)U^{-1/2}
 \right\|
\]
on the $N=8$ card, where $U=Y^*Y$ is the physical target Gram.

\begin{center}
\small
\begin{tabular}{rcc}
\toprule
$m$&$h=(\beta/2)/m$&target-normalized residual error\\
\midrule
1&$5.000000\times10^{-1}$&$5.1160643\times10^{-3}$\\
2&$2.500000\times10^{-1}$&$1.0677418\times10^{-3}$\\
4&$1.250000\times10^{-1}$&$2.5317395\times10^{-4}$\\
8&$6.250000\times10^{-2}$&$6.2424303\times10^{-5}$\\
16&$3.125000\times10^{-2}$&$1.5551608\times10^{-5}$\\
32&$1.562500\times10^{-2}$&$3.8844956\times10^{-6}$\\
\bottomrule
\end{tabular}
\end{center}

The observed $O(m^{-2})$ behavior is consistent with the finite-card
bi-energy curvature branch of \Cref{cor:V5-second-order}.  The universal
endpoint theorem only claims the spectrum-free $O(m^{-1})$ rate.
All matrix identities, positivity checks, and the endpoint error bound pass to
floating precision.  The replay source is
\begin{center}\small
\path{Renewal_SM_Einstein_Mixed_Clock_Duhamel_V5.py}
\end{center}
and its JSON output is
\begin{center}\small
\path{Renewal_SM_Einstein_Mixed_Clock_Duhamel_V5.json}.
\end{center}

\begin{firewall}[Replay scope]
\label{fire:V5-replay-scope}
The replay uses the Ising member of the reversible refresh-plus-local-generator
class.  It does not contain the selected chiral Dirac--Yukawa coefficient,
sector-patched determinant/Pfaffian line, gauge/Higgs/ghost field shell,
physical current atlas, stress insertion, or Einstein route.  It verifies the
finite acquisition formulas and their numerical conditioning only.
\end{firewall}

% =============================================================================
\section{The first missing original RPESM row and the Version V6 attack}
\label{sec:V5-handoff}
% =============================================================================

The attached RPESM construction establishes that every required finite writer
may occur on one event law.  It does not print one selected numerical tuple
\[
 (H_{\rm L,n},H_{\rm R,n},B_{\lambda,n},\Vraw_{\lambda,n},M_n,N_n)
\]
or the mixed split-clock matrices derived from it.  The separate population
report likewise leaves the Dirac--Yukawa card open after evaluating a reversible
surrogate.  Therefore the first missing physical object is not another
abstract source space.  It is the following finite card.

\begin{definition}[Original-card acquisition manifest]
\label{def:V5-original-manifest}
For one named RPESM cutoff, one finite left-energy head, and one mesh $m$, the
minimum submitted payload is
\begin{equation}
 \boxed{
 \begin{gathered}
 \texttt{RPESM\_THERMAL\_DUHAMEL\_MIXED\_CLOCK\_HEAD\_PACKET}:\\
 \beta,\ \tau,\ M,\ N,\ (E_\lambda^{\rm L})_{\lambda\le\Lambda},
 \ (S_\lambda)_{\lambda\le\Lambda},\\
 (C_{\lambda,j}^{(m)})_{\lambda\le\Lambda,\ 0\le j\le m},
 \quad
 (K_{\lambda,n}^{(m)})_{\lambda\le\Lambda,\ 0\le n\le2m},\\
 \text{canonical supports, supported floors, outward radii,}\\
 \text{one adjacent route and one independent direct/staged route.}
 \end{gathered}}
 \label{eq:V5-original-manifest}
\end{equation}
The target is reconstructed by \Cref{thm:V5-card-reconstruction}; it is not
supplied by changing the sign of a multiplier.  The reserve is evaluated in
$M$, and the residual interval in $N$.
\end{definition}

\begin{theorem}[Sufficiency and exact stopping boundary of the manifest]
\label{thm:V5-manifest-sufficiency}
A passing packet \eqref{eq:V5-original-manifest} determines:
\begin{enumerate}[label=\textup{(M\arabic*)},leftmargin=3em]
\item the quadrature approximation to every retained physical Duhamel target;
\item the complete five-matrix fibre packet and full-primitive residual;
\item the endpoint-sensitive quadrature debit and positive left-energy tail;
\item a robust birth, quantified follower, exact follower, or unresolved
      finite verdict according to \Cref{thm:V5-compiler};
\item the target-normalized cofinal handoff under the hypotheses of
      \Cref{thm:V5-cofinal}.
\end{enumerate}
Conversely, diagonal source metrics, target norms, source reserves, or
horizon eigenvalues without the mixed matrices $C_{\lambda,j}^{(m)}$ cannot
determine the full-primitive ancestry branch, as shown by
\Cref{cth:V5-C-irredundant}.  If the named RPESM card does not expose these
mixed same-history rows, the correct output is the missing-row stop rather
than a numerical ancestry claim.
\end{theorem}

\begin{proof}
Items \textup{(M1)--(M3)} are
\Cref{thm:V5-card-reconstruction,cor:V5-five-matrix-population,thm:V5-full-interval}.
Items \textup{(M4)--(M5)} are
\Cref{thm:V5-compiler,thm:V5-cofinal}.  Irredundancy is the pair of exact cards
in \Cref{cth:V5-C-irredundant}.
\end{proof}

\begin{openproblem}[Version V6: execute the first original mixed-clock packet]
\label{op:V6-original-card}
Preserve Versions~V1--V5 verbatim.  Freeze one named finite RPESM cutoff and
one physical common-action command before inspecting the residual.  Then:
\begin{enumerate}[label=\textup{(V6.\arabic*)},leftmargin=5.1em]
\item export the actual finite $H_{\rm L}$ and $H_{\rm R}$ or their exact
      split-clock semigroups on the selected standard-form carrier;
\item export the complete primitive fibre bank and the raw centered
      Dirac--Yukawa/common-action insertion on one left-energy head;
\item freeze the source-action metric $M$ and target metric $N$ and pass their
      null-cost and support tests;
\item choose $m$ from a predeclared target tolerance, acquire
      \eqref{eq:V5-original-manifest}, and verify its block positivity;
\item compute the five matrices, endpoint debit, data radius, and V4
      left-energy tail;
\item return the first strict robust-birth fibre, an exact reducing follower,
      a quantified unresolved interval, or the first support/metric/route
      obstruction;
\item repeat at one adjacent cutoff and one independent direct route;
\item keep current, stress, Kubo, charge, locality, Lorentzian, Einstein, and
      pure-colour conclusions closed until this packet lands.
\end{enumerate}
No further general horizon or source-cyclic theorem is the accepted substitute
for this acquisition.
\end{openproblem}

\begin{theorem}[Version V5 export to the main continuum programme]
\label{thm:V5-export}
The main SM--Einstein continuum programme may replace a demand for either the
complete one-lag kernel field or a uniform global horizon floor by the finite
mixed-clock packet \eqref{eq:V5-original-manifest}.  At mesh $m$ and left head
$\Lambda$ it carries the explicit target-metric debit
\begin{equation}
 \boxed{
 \delta_-\le
 \delta^{\rm data}_{\Lambda,m;N}
 +\left(ah+\frac{h^2}{4}\right)c_V,
 \qquad
 \delta_+\le\delta_-
 +4\max\{a,1\}^2c_V(1+\Lambda)^{-2}.}
 \label{eq:V5-export-debits}
\end{equation}
A strict lower margin beyond $\delta_-$ is a robust dynamic birth.  An upper
margin tending to zero is an asymptotic or exact follower according to whether
the zero is limiting or algebraic.  Stable route transport and vanishing
debits give the cofinal ancestry packet.  Nothing in this export promotes the
linear target to current, stress, interaction, Lorentzian, or Einstein closure.
\end{theorem}

\begin{proof}
The finite reconstruction and quadrature debit are
\Cref{thm:V5-card-reconstruction,thm:V5-head-interval}.  The left-tail term is
\Cref{thm:V4-physical-tail}.  Data radii are
\Cref{prop:V5-Schur-perturbation}.  The branch and cofinal conclusions are
\Cref{thm:V5-compiler,thm:V5-cofinal}.
\end{proof}

\begin{assessment}[Version V5 verdict]
\label{ass:V5-verdict}
\passstatus{} for the exact split-clock integral, sharp trapezoidal endpoint
bound, finite $3m+3$-matrix reconstruction, direct population of the inherited
five matrices, target-metric residual interval, supported-data perturbation
bound, domain-safe finite compiler, and cofinal landing without a global
horizon floor.

\obsstatus{} for claiming that the first four fibre matrices, diagonal source
and target Grams, raw reserve, or compact-horizon spectrum determine the
full-primitive ancestry branch.  The two canonical Duhamel cards of
\Cref{cth:V5-C-irredundant} have identical $(S,T,U,G)$ and opposite residuals.

\stopstatus{} for the numerical Dirac--Yukawa ancestry value.  The supplied
population material proves finite occurrence but does not expose the named
mixed split-clock matrices, physical coefficient metrics, and route-stable
supports required by \eqref{eq:V5-original-manifest}.  The Ising replay is a
compiler diagnostic only.  Current, stress, Kubo, charge, interacting
continuum, Lorentzian reconstruction, Einstein backreaction, and pure-colour
mass remain downstream.
\end{assessment}

\section*{Version V5 append-only record}
\addcontentsline{toc}{section}{Version V5 append-only record}

Preserved the complete V4 source verbatim before its sole terminal
\verb|\end{document}|.  The SHA--256 of the complete V4 input is
\begin{center}\small\ttfamily
33be0e4430902e3d2ba35beabaf74c48d9f23483bd4363b47d276f43d67ffe01.
\end{center}

Moved from the unattainable complete-horizon floor to a target-specific finite
acquisition.  Rewrote the physical square-root derivative as a split-clock
integral and proved the sharp composite-trapezoid error, its endpoint thermal
mismatch refinement, and a second-order bi-energy-curvature branch.  Constructed
the $3m+3$ mixed-lag matrix packet, proved exact reconstruction of the
quadrature target and full-primitive residual, and derived the inherited five
matrices from the same data.  Proved an exact irredundancy pair with identical
$(S,T,U,G)$ and opposite ancestry.  Combined quadrature, input, and left-tail
debits into a target-metric interval and a domain-safe finite compiler.  Proved
cofinal convergence of the exact residual from transported finite packets,
without a global horizon floor or bounded cyclic depth.  Replayed the finite
acquisition on the thermally corrected reversible slab and isolated the first
missing original RPESM row as the mixed primitive/raw-action split-clock Gram.

\begin{firewall}[Append-only instruction after V5]
\label{fire:V5-append}
Any Version~V6 must preserve this complete V5 source verbatim, append new
mathematical material immediately before the sole final
\verb|\end{document}|, and use
\begin{center}\small
\path{Renewal_Einstein_Standard_Model_Physical_Source_Metric_Duhamel_Ancestry_and_Quantum_Continuum_Programme_V6.tex}.
\end{center}
It must execute one named original RPESM mixed-clock packet or return its first
absent physical matrix.  It may not replace the complete primitive fibre by the
action source alone, omit the mixed primitive--target lag row, infer exact
followerhood from one finite quadrature, use an unverified support in a
Moore--Penrose short, treat finite population existence as a numerical card,
demand a global horizon floor, or reopen current/stress/Einstein before the
linear ancestry packet and its routes land.
\end{firewall}

% =============================================================================
% END APPENDED VERSION V5 MATERIAL
% =============================================================================


% =============================================================================
% BEGIN APPENDED VERSION V6 MATERIAL
% =============================================================================
\clearpage
\providecommand{\Dh}{\mathsf D_h}

\section*{Version V6: finite-lag absorption moments, a physically corrected Dirac--Yukawa replay, and the first cofinal route obstruction}
\addcontentsline{toc}{section}{Version V6: finite-lag absorption moments and corrected Dirac--Yukawa replay}

\begin{abstract}
Version~V5 reduced the physical square-root Duhamel target to a finite
mixed-clock packet and gave a target-metric interval for the complete dynamic
residual.  The newly supplied common cofinal-survivor note has already proved
the generic concentration, proper-moment, and summable-transport theorems and
has correctly stopped and exported its remaining Standard-Model rows.  Those
rows are sector specific: on a collapsing Duhamel residual island one must
locate the noncollapsed prior-primitive absorption carrier and control its
right-energy and primitive-depth tails.  The five unweighted Duhamel matrices
do not determine those weighted tails.

This version supplies the missing finite acquisition bridge.  If $A$ is the
action source, $F$ the complete prior-primitive bank, and $N=(I-P_A)F$, then
for every nonnegative physical label operator $K$ the residualized primitive
lag $G_K(s)=N^*e^{-sK}N$ is reconstructed from one joint action/prior-primitive
lag panel.  The bounded semigroup-difference label
\[
 D_h(K)=\frac{I-e^{-hK}}h
\]
has the exact three-lag moment formula
\[
 N^*(I+D_h(K))^2N
 =\left(1+\frac2h+\frac1{h^2}\right)G_K(0)
  -\left(\frac2h+\frac2{h^2}\right)G_K(h)
  +\frac1{h^2}G_K(2h).
\]
Thus the two weighted primitive Grams requested by the cofinal note can be
approached monotonically using only the static panel and two positive-time lag
panels.  At $h=R^{-1}$ this finite matrix controls the spectral tail above $R$
with order $R^{-2}$.  A low-island generalized edge therefore gives a
summable $O(R^{-1})$ carrier-norm tail without a global unbounded moment on
source directions which do not participate in the failure.  Exact
countercards prove that the mixed primitive lag and the $2h$ panel are both
irredundant.  The route estimate exposes the unavoidable $h^{-2}$ amplification
of shrinking-lag defects.

The supplied consolidated session report also contains a concrete finite
$\mathbb Z_3$-link, three-valued-Higgs Dirac--Yukawa model.  Its printed
ancestry calculation uses the sign-reversed multiplier which Version~V4 proved
is a different target.  We replay the same finite model with the inherited
thermally dressed primitive bank and the unique Gibbs-square-root Duhamel
derivative.  At $N=2$ and $N=3$ sites every retained fibre has a numerically
positive full-primitive residual; the target-normalized dynamic fractions lie
in $[0.0166,0.0336]$ and $[0.0280,0.0688]$, respectively; and the V5 and V6
finite-lag identities pass to floating-point accuracy.  This is evidence that
the finite birth pattern survives the physical source/target correction on two
cards, not an outward interval certificate, a cofinal sequence, or the
four-dimensional RPESM regulator.  The first unresolved row is now exact:
transport the physical coefficient metric and V5 card, and, only on a
collapsing residual island, transport the two-lag action/prior-primitive panels
for the right-energy and primitive-depth labels.  No current, stress, Kubo,
interacting-continuum, Lorentzian, or Einstein claim is opened.
\end{abstract}

% =============================================================================
\section{Direction after the supplied cofinal and session reports}
\label{sec:V6-direction}
% =============================================================================

The common cofinal-survivor programme has already converted collapse of a
positive effective short under a noncollapsed direct window into a normalized
physical carrier, decomposed that carrier into compact, delocalized, and
escaping stocks, and promoted a selected survivor under summable transport.
Its Version~V3 stops and exports because every remaining Standard-Model
quantity is a literal fibre matrix or route defect.  Repeating that abstract
machinery here would move this programme backward.

The exported task has two ordered layers.  First, the V5 mixed-clock packet
decides whether the full-primitive dynamic residual is strictly positive,
exactly zero, or unresolved within its outward interval.  Only on a
low-residual branch with a noncollapsed Duhamel-variance window is there a
second question: whether the prior-primitive absorption carrier stays in a
compact physical right-energy/primitive-depth head and transports cofinally.
The second question is not an additional condition for a strict first-layer
birth certificate.

\begin{firewall}[No target inflation after the V5 verdict]
\label{fire:V6-no-target-inflation}
If the V5 lower interval is strictly positive on a protected target screen,
the dynamic-birth branch has landed and no weighted absorption panel is
required.  The panels introduced below are acquired only when a residual
island is being classified or promoted.  Conversely, a bounded absorption
moment does not decide the full-primitive residual without the V5 mixed
primitive--target rows.  The two finite cards answer different questions.
\end{firewall}

The consolidated session report supplies a useful finite model, but its final
numerical branch uses a sign-reversed ``decaying calibration.''  Version~V4
proved that this changes the physical target rather than merely its
coordinates.  The model definition is reused only after replacing both source
and target by the inherited thermal formulas; see
\Cref{sec:V6-explicit-replay}.

% =============================================================================
\section{Residualized prior-primitive lags are one joint positive panel}
\label{sec:V6-residualized-lags}
% =============================================================================

Fix one left-energy fibre and suppress its label.  Let
\[
 A:E_A\longrightarrow\mathcal K,
 \qquad F:E_F\longrightarrow\mathcal K
\]
be the physical action-source synthesis and the complete bank of prior
primitives which occur before the dynamic short.  Put
\begin{equation}
 \begin{aligned}
 S&=A^*A,
 &D&=A^*F,
 &T&=S^\dagger D,
 &P_A&=AS^\dagger A^*,\\
 N&=(I-P_A)F=F-AT,
 &G&=N^*N,
 &P_N&=NG^\dagger N^*.
 \end{aligned}
 \label{eq:V6-static-residualization}
\end{equation}
Positivity of the static joint Gram implies $\Ran D\subseteq\Ran S$, so
$AT=P_AF$ and the residualization is exact.

Let $K=K^*\succeq0$ be a physical label operator on $\mathcal K$.  In the two
applications below, $K$ is the right Hamiltonian $H^R$ or a separately
declared primitive-depth operator $D_{\rm pr}$.  For $s\ge0$ define
\begin{equation}
 \mathbb L_K(s)
 :=\begin{pmatrix}Q_K(s)&R_K(s)\\R_K(s)^*&P_K(s)\end{pmatrix}
 :=\begin{pmatrix}A^*\\F^*\end{pmatrix}
 e^{-sK}\begin{pmatrix}A&F\end{pmatrix}\succeq0.
 \label{eq:V6-joint-lag-panel}
\end{equation}
The static panel at $s=0$ is already contained in the complete primitive Gram.

\begin{theorem}[Exact residualized primitive-lag formula]
\label{thm:V6-residualized-lag}
For every $s\ge0$,
\begin{equation}
 \boxed{
 G_K(s):=N^*e^{-sK}N
 =P_K(s)-R_K(s)^*T-T^*R_K(s)+T^*Q_K(s)T.}
 \label{eq:V6-residualized-lag}
\end{equation}
Thus the right-energy or primitive-depth law of the absorbing primitive is
obtained by residualizing one same-card joint lag panel.
\end{theorem}

\begin{proof}
Insert $N=F-AT$ on both sides of $e^{-sK}$ and expand.  The resulting four
terms are the four blocks in \eqref{eq:V6-joint-lag-panel}.  Positivity follows
independently because $G_K(s)$ is the Gram of $e^{-sK/2}N$.
\end{proof}

\begin{countertheorem}[Lag marginals do not determine the residualized lag]
\label{cth:V6-lag-marginals}
There are two physical cards with the same static data and identical action and
prior-primitive lag marginals, but different residualized lags.
\end{countertheorem}

\begin{proof}
Fix $s>0$ and take
\[
 \mathcal K=\C^3,
 \quad K=\diag\!\left(0,\frac{\log2}{s},\frac{\log4}{s}\right),
 \quad A1=\frac{e_1+e_2+e_3}{\sqrt3}.
\]
Use
\[
 F_01=\frac{e_1+e_2-e_3}{\sqrt3},
 \qquad F_11=\frac{e_1-e_2+e_3}{\sqrt3}.
\]
Both static cards have $S=E=1$, $D=1/3$, and $T=1/3$.  Both have
$Q_K(s)=P_K(s)=7/12$, whereas
\[
 R_{K,0}(s)=\frac5{12},
 \qquad R_{K,1}(s)=\frac14.
\]
Formula \eqref{eq:V6-residualized-lag} gives
$G_{K,0}(s)=10/27$ and $G_{K,1}(s)=13/27$.  Every quantity comes from the same
nonnegative semigroup, so the mixed lag is a genuine physical row.
\end{proof}

% =============================================================================
\section{Three-lag recovery of the weighted absorption moment}
\label{sec:V6-three-lag}
% =============================================================================

For $h>0$ define
\begin{equation}
 \boxed{D_h(K):=\frac{I-e^{-hK}}h.}
 \label{eq:V6-Dh}
\end{equation}
This is a bounded positive Borel function of $K$, even when $K$ is unbounded.

\begin{theorem}[Exact three-lag weighted primitive Gram]
\label{thm:V6-three-lag}
Define
\begin{equation}
 \Gamma_{K,h}:=N^*(I+D_h(K))^2N\succeq0.
 \label{eq:V6-Gamma-definition}
\end{equation}
Then
\begin{equation}
 \boxed{
 \begin{aligned}
 \Gamma_{K,h}
 ={}&\left(1+\frac2h+\frac1{h^2}\right)G_K(0)\\
 &-\left(\frac2h+\frac2{h^2}\right)G_K(h)
 +\frac1{h^2}G_K(2h).
 \end{aligned}}
 \label{eq:V6-three-lag-formula}
\end{equation}
Beyond the static Gram, two positive-time joint lag panels at $h$ and $2h$
therefore reconstruct the bounded weighted moment exactly.  For every
coefficient vector $c$,
\begin{equation}
 c^*\Gamma_{K,h}c\uparrow\|(I+K)Nc\|^2
 \quad(h\downarrow0),
 \label{eq:V6-monotone-form-limit}
\end{equation}
with value $+\infty$ when $Nc\notin\operatorname{Dom}K$.
\end{theorem}

\begin{proof}
Writing $E_h=e^{-hK}$,
\[
 (I+D_h(K))^2
 =I+\frac2h(I-E_h)+\frac1{h^2}(I-2E_h+E_h^2).
\]
Since $E_h^2=e^{-2hK}$, compression by $N$ proves
\eqref{eq:V6-three-lag-formula}.  For scalar $x\ge0$,
$d_h(x)=(1-e^{-hx})/h$ increases to $x$ as $h\downarrow0$, because
$t\mapsto(1-e^{-t})/t$ is decreasing.  Spectral monotone convergence proves
\eqref{eq:V6-monotone-form-limit}.
\end{proof}

\begin{theorem}[Matched finite-lag tail certificate]
\label{thm:V6-three-lag-tail}
Let $R>0$, put $h=R^{-1}$, and set $c_0:=1-e^{-1}$.  Then
\begin{equation}
 \boxed{
 N^*\one_{[R,\infty)}(K)N
 \preceq\frac{1}{(1+c_0R)^2}\Gamma_{K,1/R}.}
 \label{eq:V6-three-lag-tail}
\end{equation}
\end{theorem}

\begin{proof}
For $x\ge R$,
$D_{1/R}(x)=R(1-e^{-x/R})\ge c_0R$.  Spectral calculus gives
$(I+D_{1/R}(K))^2\succeq(1+c_0R)^2\one_{[R,\infty)}(K)$; compress by $N$.
\end{proof}

\begin{countertheorem}[The $2h$ lag is irredundant]
\label{cth:V6-two-h-irredundant}
For every $h>0$ there are two one-column primitive cards with identical
$G_K(0)$ and $G_K(h)$ but different $G_K(2h)$ and hence different
$\Gamma_{K,h}$.
\end{countertheorem}

\begin{proof}
Take first $\mathcal K_0=\C$, $N_01=1$, and $K_0=(\log2)/h$.  Then the three
lags are $1,1/2,1/4$.  Take second
\[
 \mathcal K_1=\C^2,
 \quad K_1=\diag\!\left(0,\frac{\log4}{h}\right),
 \quad N_11=\frac1{\sqrt3}e_1+\sqrt{\frac23}e_2.
\]
Its three lags are $1,1/2,3/8$.  Formula
\eqref{eq:V6-three-lag-formula} gives
$\Gamma_{K_1,h}-\Gamma_{K_0,h}=1/(8h^2)>0$.
\end{proof}

% =============================================================================
\section{Low-island finite-lag properness for the absorption carrier}
\label{sec:V6-low-island}
% =============================================================================

Retain the five-matrix notation on one fibre:
\[
 \mathbb V=Y^*(I-P_A)Y,
 \qquad \mathbb A=C^*G^\dagger C,
 \qquad \mathbb R=\mathbb V-\mathbb A,
 \qquad C=N^*Y.
\]
Whiten by the physical action-source Gram $S=A^*A$ and define
\begin{equation}
 X=(I-P_A)YS^{\dagger/2},
 \quad Z=NG^\dagger CS^{\dagger/2},
 \quad R=(I-P_A-P_N)YS^{\dagger/2}.
 \label{eq:V6-whitened-syntheses}
\end{equation}
Then
\[
 X^*X=\mathbb V^{\rm wh},
 \quad Z^*Z=\mathbb A^{\rm wh},
 \quad R^*R=\mathbb R^{\rm wh},
 \quad X=Z+R,
 \quad Z\perp R.
\]

Let $P=\one_{[0,\rho]}(\mathbb R^{\rm wh})$ be a nonzero residual island and
assume
\begin{equation}
 \boxed{
 0<\ell_-P\preceq P\mathbb V^{\rm wh}P\preceq\ell_+P,
 \qquad 0\le\rho<\ell_-.}
 \label{eq:V6-island-window}
\end{equation}
Since $\mathbb A^{\rm wh}=\mathbb V^{\rm wh}-\mathbb R^{\rm wh}$,
\begin{equation}
 (\ell_- -\rho)P
 \preceq P\mathbb A^{\rm wh}P
 \preceq\ell_+P.
 \label{eq:V6-absorption-island-window}
\end{equation}

For a physical label $K$, define
\begin{equation}
 \begin{aligned}
 W_{K,h}
 &:=Z^*(I+D_h(K))^2Z\\
 &=S^{\dagger/2}C^*G^\dagger
   \Gamma_{K,h}G^\dagger CS^{\dagger/2}.
 \end{aligned}
 \label{eq:V6-low-island-W}
\end{equation}

\begin{definition}[Finite-lag low-island edge]
\label{def:V6-low-island-edge}
On $\Ran P$ put
\begin{equation}
 \boxed{
 \kappa_{K,h;P}
 :=\lambda_{\max}\!\left[
 (P\mathbb A^{\rm wh}P)^{\dagger/2}
 PW_{K,h}P
 (P\mathbb A^{\rm wh}P)^{\dagger/2}
 \right].}
 \label{eq:V6-low-island-edge}
\end{equation}
The inverse is ordinary on $\Ran P$ by
\eqref{eq:V6-absorption-island-window}.
\end{definition}

\begin{theorem}[Finite-lag low-island tail]
\label{thm:V6-low-island-tail}
Let $R_0>0$ and $h=R_0^{-1}$.  For every unit $v\in\Ran P$, put
$w_v=Zv/\norm{Zv}$.  Then $Zv\ne0$ and
\begin{equation}
 \boxed{
 \norm{\one_{[R_0,\infty)}(K)w_v}^2
 \le\frac{\kappa_{K,1/R_0;P}}{(1+c_0R_0)^2}.}
 \label{eq:V6-low-island-tail}
\end{equation}
Moreover,
\begin{equation}
 \boxed{
 \norm{\one_{[R_0,\infty)}(K)ZP}
 \le\frac{\sqrt{\kappa_{K,1/R_0;P}\ell_+}}{1+c_0R_0}.}
 \label{eq:V6-low-island-operator-tail}
\end{equation}
A uniformly bounded finite-lag edge therefore gives a summable carrier-norm
tail along any $R_n\uparrow\infty$ with $\sum_nR_n^{-1}<\infty$.
\end{theorem}

\begin{proof}
Equation \eqref{eq:V6-absorption-island-window} gives
$\norm{Zv}^2\ge\ell_- -\rho>0$.  By the generalized-edge definition,
\[
 v^*W_{K,1/R_0}v
 \le\kappa_{K,1/R_0;P}\,v^*\mathbb A^{\rm wh}v
 =\kappa_{K,1/R_0;P}\norm{Zv}^2.
\]
Apply \Cref{thm:V6-three-lag-tail} to $Zv$ and divide by $\norm{Zv}^2$.
Before dividing, use $\norm{Zv}^2\le\ell_+\norm v^2$ to obtain the operator
bound.
\end{proof}

\begin{corollary}[Two-label compact absorption head]
\label{cor:V6-two-label-head}
Suppose $H^R$ and $D_{\rm pr}$ are strongly commuting nonnegative physical
label operators and their joint bounded sublevels have finite rank.  If their
finite-lag island edges are bounded uniformly at schedules
$R_n,Q_n\uparrow\infty$ satisfying
\[
 \sum_n(R_n^{-1}+Q_n^{-1})<\infty,
\]
then the part of $Z_nP_n$ outside
$\one_{[0,R_n]}(H_n^R)\one_{[0,Q_n]}(D_{{\rm pr},n})$ has summable norm.
Together with the V5 left-energy tail and the imported summable survivor
transport theorem, this supplies a nonzero cofinal prior-absorption survivor
whenever the residual-island and same-head route defects are summable.
\end{corollary}

\begin{proof}
Strong commutation makes the product of the sublevel projections the joint
sublevel projection.  Its complement is bounded by the sum of the two
individual complements.  Apply \Cref{thm:V6-low-island-tail} to each label and
sum.  The V5 positive left-energy tail controls the remaining label.  These
are precisely the compact-head and summable-tail hypotheses of the imported
cofinal-survivor theorem.
\end{proof}

\begin{firewall}[Finite-lag positivity is not a uniform edge]
\label{fire:V6-edge-not-automatic}
Positivity of the lag panels proves $\Gamma_{K,h}\succeq0$ but does not bound
$\kappa_{K,h;P}$.  A diverging edge is a typed right-energy or primitive-depth
polarization inside the failing island.  The theorem removes the need to
measure an unbounded moment directly; it does not manufacture the physical lag
panels or their generalized-eigenvalue bound.
\end{firewall}

% =============================================================================
\section{Transport cost of the three-lag moment}
\label{sec:V6-three-lag-transport}
% =============================================================================

The exact formula contains cancellations with coefficients of order $h^{-2}$.
Its route transport must therefore be controlled on the raw positive lag
panels before the final combination is interpreted.

\begin{theorem}[Explicit three-lag transport amplification]
\label{thm:V6-three-lag-transport}
After pulling two adjacent cards to one coefficient chart, suppose
\[
 \norm{\Delta G_K(0)}\le\varepsilon_0,
 \quad\norm{\Delta G_K(h)}\le\varepsilon_h,
 \quad\norm{\Delta G_K(2h)}\le\varepsilon_{2h}.
\]
Then
\begin{equation}
 \boxed{
 \begin{aligned}
 \norm{\Delta\Gamma_{K,h}}
 \le{}&\left(1+\frac2h+\frac1{h^2}\right)\varepsilon_0\\
 &+\left(\frac2h+\frac2{h^2}\right)\varepsilon_h
 +\frac1{h^2}\varepsilon_{2h}.
 \end{aligned}}
 \label{eq:V6-three-lag-transport}
\end{equation}
At $h=R^{-1}$,
\begin{equation}
 \boxed{
 \norm{\Delta\Gamma_{K,1/R}}
 \le(1+2R)^2
 \max\{\varepsilon_0,\varepsilon_{1/R},\varepsilon_{2/R}\}.}
 \label{eq:V6-R2-amplification}
\end{equation}
Thus fixed-lag convergence does not suffice on a shrinking-lag diagonal; the
raw route defects must beat the $R^2$ amplification.
\end{theorem}

\begin{proof}
Subtract the two copies of \eqref{eq:V6-three-lag-formula} and apply the
triangle inequality.  For $h=R^{-1}$, the sum of the absolute coefficients is
\[
 (1+2R+R^2)+(2R+2R^2)+R^2=(1+2R)^2.
\]
\end{proof}

\begin{corollary}[Raw joint-panel route budget]
\label{cor:V6-raw-panel-budget}
If the static residualization maps and the joint blocks in
\eqref{eq:V6-joint-lag-panel} have support-stable pulled-back defects, then
\eqref{eq:V6-residualized-lag} first gives defects of the three residualized
lags.  Summability of the right side of
\eqref{eq:V6-three-lag-transport} implies summable transport of the bounded
weighted primitive Gram.  The variance, absorption, final residual, and
low-island projection remain separate positive-term route rows.
\end{corollary}

\begin{proof}
The residualized-lag formula is a finite polynomial in the joint blocks and
$T$.  Uniform norm ceilings and repeated product estimates transport its three
values.  Apply \Cref{thm:V6-three-lag-transport}.  The final sentence is the
positive-term transport discipline already used by V5.
\end{proof}

% =============================================================================
\section{The conditional target-minimal extension of the V5 packet}
\label{sec:V6-manifest}
% =============================================================================

\begin{definition}[Conditional absorption-moment extension]
\label{def:V6-extended-manifest}
For $K\in\{H^R,D_{\rm pr}\}$ and a predeclared matched lag $h_K=R_K^{-1}$,
acquire
\begin{equation}
 \boxed{
 \begin{gathered}
 \mathfrak A_{K,h_K}
 =\bigl(\mathbb L_K(h_K),\mathbb L_K(2h_K)\bigr),\\
 \text{with canonical supports and supported floors,}\\
 \text{outward radii, and physical route tags.}
 \end{gathered}}
 \label{eq:V6-absorption-extension}
\end{equation}
Together with the static complete-primitive Gram, these two panels reconstruct
\[
 G_K(0),\qquad G_K(h_K),\qquad G_K(2h_K),\qquad\Gamma_{K,h_K},
\]
as well as the low-island edge and the spectral-tail certificate.  The
primitive-depth extension is omitted when no physical depth operator has been
declared; its absence is a typed loading obstruction, not an invitation to
invent an index on the auxiliary columns.
\end{definition}

\begin{theorem}[Target-minimality of the V6 extension]
\label{thm:V6-extension-minimality}
The two positive-time panels in \eqref{eq:V6-absorption-extension} are
sufficient by \Cref{thm:V6-three-lag}.  The $2h$ row is necessary by
\Cref{cth:V6-two-h-irredundant}, and the mixed action/prior-primitive block in
each panel is necessary by \Cref{cth:V6-lag-marginals}.  Hence the extension is
minimal among universal linear reconstructions from the semigroup lags
$0,h,2h$.
\end{theorem}

\begin{proof}
Sufficiency is residualization followed by the three-lag identity.  Each
countertheorem holds all proposed reduced data fixed while changing the desired
output, proving the two independent necessity statements.
\end{proof}

\begin{assessment}[Revised packet order]
\label{ass:V6-packet-order}
The executable order is:
\begin{enumerate}[label=\textup{(P\arabic*)},leftmargin=3.2em]
\item run the V5 mixed-clock packet and target-metric interval;
\item stop on a strict dynamic birth or exact reducing follower;
\item on an unresolved or collapsing residual island, verify a noncollapsed
      Duhamel-variance window in the physical action metric;
\item only then acquire the right-energy and, when physically declared,
      primitive-depth extensions \eqref{eq:V6-absorption-extension};
\item compute the finite-lag generalized edges and matched tail bounds;
\item transport the raw lag panels and all three positive Duhamel terms through
      one adjacent and one independent direct route.
\end{enumerate}
This is the sector-specific implementation of the exported cofinal theorem,
not a new generic survivor programme.
\end{assessment}

% =============================================================================
\section{Physically corrected replay of the explicit Dirac--Yukawa representative}
\label{sec:V6-explicit-replay}
% =============================================================================

The consolidated session report specifies a finite one-dimensional member of
the broad RPESM class: $N$ sites, a $\mathbb Z_3$ link and a Higgs value in
$\{-1,0,1\}$ at each site, a finite-range bosonic action, a two-component
chiral Dirac--Yukawa coefficient, a complex determinant line, a real Pfaffian
sign line, and a refresh-plus-local-heat-bath generator.  We retain its frozen
couplings and writer list, but replace the report's sign-reversed target by the
inherited physical card.

At $\beta_Q=1$, $a=1/2$, and $\tau_0=1/4$, every primitive writer $W$ is
represented by
\begin{equation}
 B(W)=Z^{-1/2}e^{-\tau_0H_{\rm L}}W e^{-aH_{\rm R}},
 \label{eq:V6-replay-source}
\end{equation}
and every centered action insertion $V$ has target
\begin{equation}
 Y(V)=-Z^{-1/2}\int_0^a
 e^{-(a-s)H_{\rm L}}V e^{-sH_{\rm R}}\,\dd s.
 \label{eq:V6-replay-target}
\end{equation}
The complete primitive bank contains the sector weight, five action writers,
the determinant phase, and the Pfaffian sign.  The V5 quadrature uses $m=8$;
the retained left head is $\lambda\le\lambda_{\min}+1.5$, and the V6 matched
right-energy tail level is $R=\lambda_{\min}+3$.

\begin{table}[ht]
\centering
\small
\begin{tabular}{@{}lrr@{}}
\toprule
Quantity & $N=2$ & $N=3$\\
\midrule
Configuration count & $81$ & $729$\\
Centered carrier dimension & $80$ & $728$\\
Minimum $|\det D_\chi|$ & $0.575754$ & $0.493066$\\
Centered spectral gap & $1.110655$ & $1.294584$\\
All fibres / retained fibres & $65/26$ & $445/46$\\
Numerically positive retained residuals & $26/26$ & $46/46$\\
Dynamic-fraction eigenvalue interval & $[0.016565,0.033616]$ & $[0.028038,0.068804]$\\
Head-fraction eigenvalue interval & $[0.008560,0.020368]$ & $[0.010555,0.019766]$\\
Omitted left residual divided by $\|U\|$ & $0.020092$ & $0.042830$\\
Exact right-energy relative-edge range & $[9.533,17.350]$ & $[9.459,18.779]$\\
Finite-lag relative-edge range & $[6.623,10.220]$ & $[6.790,10.787]$\\
\bottomrule
\end{tabular}
\caption{Floating-point replay with the physical thermal source and unique
Duhamel target.  ``Positive'' means operator norm above $10^{-12}$ in the
dense numerical calculation; it is not an outward interval statement.}
\label{tab:V6-corrected-replay}
\end{table}

The coefficient metric is the stationary-law Gram declared by the explicit
model.  Because ambient Hilbert--Schmidt normalization is a separate physical
choice, both reserve scalings are recorded:
\begin{align}
 [g_-^{\rm HS},g_+^{\rm HS}]_{N=2}
 &= [0.11237,1.21634],
 & [g_-^{\rm HS},g_+^{\rm HS}]_{N=3}
 &= [0.06069,0.59253],
 \label{eq:V6-replay-reserves-a}\\
 [g_-^{\rm nHS},g_+^{\rm nHS}]_{N=2}
 &= [1.4046\!\times\!10^{-3},1.5204\!\times\!10^{-2}],
 & [g_-^{\rm nHS},g_+^{\rm nHS}]_{N=3}
 &= [8.3360\!\times\!10^{-5},8.1391\!\times\!10^{-4}].
 \label{eq:V6-replay-reserves-b}
\end{align}
Here ${\rm HS}$ denotes the unnormalized Hilbert--Schmidt product and
${\rm nHS}$ normalized trace with the same fixed coefficient metric.  The
second row warns against calling a raw source Gram a cutoff-uniform physical
reserve.  The source-action metric must be transported with its declared
ambient normalization.

\begin{proposition}[Finite replay checks]
\label{prop:V6-replay-checks}
For the two cards in \Cref{tab:V6-corrected-replay}, the maximum numerical
residuals are
\begin{equation}
 \begin{array}{c|cc}
  &N=2&N=3\\ \hline
 \text{five-matrix reconstruction}
  &1.06\times10^{-13}&3.76\times10^{-14}\\
 \text{V5 mixed-clock lag reconstruction}
  &1.10\times10^{-15}&1.88\times10^{-15}\\
 \text{V6 three-lag weighted-Gram reconstruction}
  &4.68\times10^{-16}&9.23\times10^{-16}
 \end{array}
 \label{eq:V6-replay-errors}
\end{equation}
The least eigenvalue of the V6 tail slack over all tested fibres is
$3.96\times10^{-7}$ for $N=2$ and $5.67\times10^{-8}$ for $N=3$.  The complete
target-minus-residual Grams have least eigenvalues $2.42\times10^{-2}$ and
$3.16\times10^{-2}$.
\end{proposition}

\begin{proof}
The accompanying script enumerates all $9^N$ configurations, constructs the
reversible heat-bath conditional expectations and centered Hamiltonian,
diagonalizes its left/right standard-form clocks, builds
\eqref{eq:V6-replay-source}--\eqref{eq:V6-replay-target}, and evaluates every
fibre using thin singular-vector bases for the action, residualized prior, and
complete primitive ranges.  It compares the direct projection residual with
the five-matrix and V5 lag formulas, compares the direct finite-difference
weighted Gram with \eqref{eq:V6-three-lag-formula}, and diagonalizes the
Loewner slack in \eqref{eq:V6-three-lag-tail}.  The displayed values are the
resulting floating-point extrema.
\end{proof}

\begin{assessment}[What the corrected replay proves and does not prove]
\label{ass:V6-replay-status}
The replay supplies evidence that the explicit finite model's dynamic-birth
pattern survives replacement of the sign-reversed target by the inherited
physical source and Duhamel derivative.  It also populates the right-energy
three-lag panel on two finite cards.  It does not provide certified outward
radii, a primitive-depth operator, a physical old/fine cutoff isometry, a
summable adjacent route, or the four-dimensional chiral gauge--Higgs--ghost
RPESM regulator.  The robustness ratios printed in the supplied report are not
inherited because they were computed for a different target.  No cofinal
Standard-Model branch is asserted here.
\end{assessment}

% =============================================================================
\section{Version V6 endpoint and next physical acquisition}
\label{sec:V6-endpoint}
% =============================================================================

\begin{theorem}[Version V6 finite-lag export]
\label{thm:V6-export}
The main SM--Einstein continuum programme may use the following terminating
extension of V5.
\begin{enumerate}[label=\textup{(E\arabic*)},leftmargin=3.2em]
\item The first-stage branch is decided by the physical V5 $3m+3$-matrix
      mixed-clock packet and its target-metric interval.
\item On a collapsing residual island with a noncollapsed variance window, the
      right-energy or primitive-depth weighted absorption moment is replaced by
      two joint primitive lag panels at $h$ and $2h$.
\item The three-lag identity reconstructs a positive bounded moment, and
      \eqref{eq:V6-low-island-tail} gives an explicit matched spectral tail.
\item Uniform low-island edges and summable inverse head levels reduce the
      absorption carrier to summably transported finite physical heads.
\item The raw lag route must beat the $h^{-2}$ amplification in
      \eqref{eq:V6-three-lag-transport}; otherwise the returned object is a
      clock-depth transport obstruction, not a cofinal survivor.
\end{enumerate}
No global horizon floor, global unbounded source moment, or source direction
outside the attained residual island is required.
\end{theorem}

\begin{proof}
Items \textup{(E2)--(E3)} are
\Cref{thm:V6-residualized-lag,thm:V6-three-lag,thm:V6-three-lag-tail}.
Item~\textup{(E4)} is
\Cref{thm:V6-low-island-tail,cor:V6-two-label-head}, and
item~\textup{(E5)} is \Cref{thm:V6-three-lag-transport}.  Item~\textup{(E1)}
is the preserved V5 compiler order.
\end{proof}

\begin{openproblem}[Version V7: transport one physical card before enlarging the target]
\label{op:V7-physical-route}
Preserve Versions~V1--V6 verbatim.  On one named original RPESM cutoff and one
adjacent cutoff:
\begin{enumerate}[label=\textup{(V7.\arabic*)},leftmargin=5.5em]
\item acquire the physical source-action metric, target metric, complete
      thermally dressed primitive bank, and V5 mixed-clock packet;
\item provide outward matrix radii and canonical supported singular-value
      floors, then decide the strict V5 branch;
\item if and only if a noncollapsed low-residual island remains, declare the
      physical primitive-depth operator and acquire the right-energy and
      primitive-depth panels at $h$ and $2h$;
\item evaluate the two low-island generalized edges and matched tail
      certificates;
\item transport every raw lag panel, source support, low-island projection,
      variance, absorption, and residual through one adjacent and one
      independent direct/staged route, with the $h^{-2}$ defect budget;
\item return a robust dynamic birth, an exact or asymptotic follower, a
      cofinal prior-absorption survivor, or the first metric, support, lag,
      depth, tail, or route obstruction;
\item keep determinant/Pfaffian orientation, represented current/stress,
      locality, charge, Kubo, interacting continuum, Lorentzian reconstruction,
      and Einstein backreaction closed.
\end{enumerate}
The finite model replay is not a substitute for these original-card and route
rows.
\end{openproblem}

\begin{assessment}[Version V6 verdict]
\label{ass:V6-verdict}
\passstatus{} for the exact residualized-lag formula, positive three-lag
reconstruction, monotone approach to the unbounded weighted moment, matched
$R^{-2}$ spectral-tail certificate, low-island generalized-edge reduction,
and explicit $h^{-2}$ transport budget.  The supplied common cofinal theorem
is thereby implemented by finite Standard-Model lag panels rather than
repeated abstract machinery.

\obsstatus{} for treating the five unweighted matrices, separate primitive lag
marginals, or the static plus one-lag panel as sufficient to locate the
absorbing primitive.  Exact countercards show that the mixed lag and the $2h$
row are irredundant.  A lag defect which vanishes slower than $h^2$ also cannot
be promoted automatically along a shrinking-lag cofinal schedule.

\stopstatus{} for a cofinal physical Dirac--Yukawa verdict.  The corrected
$N=2,3$ replay has positive finite residuals and passes the new lag identities,
but lacks outward intervals, a physical primitive-depth panel, and an actual
adjacent/direct RPESM route.  These are the first missing physical rows.
\end{assessment}

\section*{Version V6 append-only record}
\addcontentsline{toc}{section}{Version V6 append-only record}

Preserved the complete V5 source verbatim before its sole terminal
\verb|\end{document}|.  The SHA--256 of the complete V5 input is
\begin{center}\small\ttfamily
9affb46f528e41e1242891bfffa5e3f4c1eee01c59a528239a2d2b8bb14ed79d.
\end{center}

Reviewed the terminal shared cofinal-survivor programme and did not duplicate
its failure-carrier, concentration-stock, or direct-limit theorem.  Imported
its sector-specific export: the Standard-Model low residual island requires
right-energy and primitive-depth weighted prior-absorption panels.  Replaced
the unbounded query by an exact finite three-lag reconstruction from joint
action/prior-primitive semigroup panels, proved the matched spectral-tail
bound, low-island generalized-edge theorem, panel minimality, and shrinking-lag
transport amplification.  Replayed the explicit finite Dirac--Yukawa
representative with the inherited thermal source and unique square-root
Duhamel target, separating dynamic-fraction evidence from the still-untransported
physical reserve.  The next step is one original RPESM card and its
adjacent/direct route, not another compiler theorem.

\begin{firewall}[Append-only instruction after V6]
\label{fire:V6-append}
Any Version~V7 must preserve this complete V6 source verbatim, append new
mathematical material immediately before the sole final
\verb|\end{document}|, and use
\begin{center}\small
\path{Renewal_Einstein_Standard_Model_Physical_Source_Metric_Duhamel_Ancestry_and_Quantum_Continuum_Programme_V7.tex}.
\end{center}
It may not require the weighted absorption panels before the V5 residual branch
is known, substitute a sign-reversed target, call an ambient normalization a
physical reserve without transporting the coefficient metric, infer cofinality
from two finite cards, omit the mixed action/prior lag, reconstruct a quadratic
weight from one lag, ignore the $h^{-2}$ route amplification, invent a primitive
depth label, or reopen current/stress/Einstein before the linear ancestry packet
and its routes land.
\end{firewall}

% =============================================================================
% END APPENDED VERSION V6 MATERIAL
% =============================================================================


% =============================================================================
% BEGIN APPENDED VERSION V7 MATERIAL
% =============================================================================
\clearpage
\section*{Version V7: the actual Hamiltonian tangent, an exact incidence obstruction, and the sharp Dirac mass threshold}
\addcontentsline{toc}{section}{Version V7: actual Hamiltonian tangent and sharp Dirac threshold}

\begin{abstract}
This continuation goes upstream of the V5--V6 lag and tail estimates.  The
inherited line/clock definition is $V(u)=D_u h$, not multiplication by the
classical action derivative.  The V6 finite Dirac--Yukawa replay thermally
regularized the latter but used a heat-bath generator as its Hamiltonian.
These choices specify a legitimate inserted-potential experiment; they do not
establish that its Duhamel target differentiates the stated heat-bath family.
We compute that missing derivative, including the moving stationary measure
and centered carrier, rather than assuming their identification.

For an exponential deformation of a finite positive law, the derivative of
each square-root-conjugated conditional expectation is an explicit rank-two
conditional-score operator.  Its compression gives the actual centered
Hamiltonian tangent.  The vacuum row of the full derivative reconstructs the
stationary Fisher metric exactly through the squared Green operator; the
native stationary half-density source is an exact follower of the original
score bank.  The finite-temperature Gibbs state of the centered kinetic
Hamiltonian remains a different, explicitly declared target.

On the original 81-record diagnostic cylinder, an outward integer-interval
calculation proves that the old inserted gauge writer cannot equal the actual
kinetic tangent, even modulo any infinitesimal unitary connection and scalar
energy shift.  Its normalized squared incidence error exceeds $0.01627$.
Independent finite differences validate the repaired Gibbs derivative, and
new ancestry calculations retain all old primitive columns.  Separately, the
homogeneous Fourier sector of the same frozen Dirac coefficient gives the
exact all-size singular-value threshold
\[
 \inf_{N\ge2,U,\phi}\sigma_{\min}(D_{m_0}(U,\phi))
 =\left(m_0+\frac15-\frac{3\sqrt2}{5}\right)_+,
 \qquad m_0\ge0.
\]
Thus its frozen value $m_0=1/2$ has no uniform Dirac singular-value floor,
although every homogeneous finite-size determinant at that value is nonzero.
This is a theorem about the declared finite-state representative, not a
Standard-Model mass-gap or continuum conclusion.  All V1--V6 source text is
preserved; the interpretation of its numerical inserted-source replay is
qualified here explicitly.
\end{abstract}

\section{Upstream question and inherited source roles}
\label{sec:V7-upstream}

The architecture in the supplied programme description requests actual source
coefficients rather than another generic compiler.  The line/clock definition
\cite{ESMLineClockV15}, equations (6.3)--(6.14), fixes
\[
 V(u)=D_u h,\qquad
 B_{\rm act}(u)=Z_\beta^{-1/2}e^{-\tau h}V^\circ(u)e^{-\beta h/2},
 \qquad Y(u)=D_u\bigl(Z_\beta^{-1/2}e^{-\beta h/2}\bigr).
\]
The derivative is taken in the physical Hilbert-bundle connection.  Thermal
multipliers alone cannot supply its source incidence.

\begin{assessment}[Explicit qualification of the V6 replay]
\label{ass:V7-V6-qualification}
The V6 replay forms $h$ from the reversible heat-bath generator, but its raw
action columns are $Q M_{w_a}Q$, with $w_a$ classical action-density writers.
Consequently its algebraic Duhamel formulas differentiate the auxiliary family
$h+\sum_a t_aQM_{w_a}Q$.  They do not, without a further identity, differentiate
the heat-bath generator obtained by varying the underlying law.  The former
is a valid finite inserted-potential card, and its calculated Grams and the
V5--V6 matrix theorems remain valid on that card.  The description of that
replay as the physical coupling derivative must be restricted accordingly.
This is a source-incidence issue before any weighted absorption or cofinal
transport test.
\end{assessment}

The calculation below keeps the finite cylinder of the supplied
\emph{SM--Einstein session consolidated report}, Part IV.  It does not replace
its gauge group, field set, action, or fixed couplings.  For the derivative
experiment we declare an actual path through that cylinder before evaluating
any residual.  The usual square-root conjugation of reversible stochastic
operators is established background \cite{V7SMF}; all derivative identities
used here are proved below.

\section{Differentiating the conditional heat bath on its native carrier}
\label{sec:V7-heatbath}

Let $\Omega=\prod_{j=1}^J\Omega_j$ be finite and let $\mu(x)>0$ be a
probability law.  Fix real writers $v_1,\ldots,v_p$ and the exponential family
\begin{equation}
 \mu_t(x)=\frac{\mu(x)\exp[-\sum_a t_av_a(x)]}
 {\sum_y\mu(y)\exp[-\sum_a t_av_a(y)]}.
 \label{eq:V7-law-path}
\end{equation}
At $t=0$, put $w_a=v_a-\mu(v_a)$, $s=\sqrt\mu$, and
$R=|s\rangle\langle s|$.  The scalar metric
\begin{equation}
 M_{ab}=\sum_x\mu(x)w_a(x)w_b(x)
 \label{eq:V7-Fisher}
\end{equation}
is the Fisher metric of this declared law path, not an arbitrarily chosen
Hilbert--Schmidt normalization.  Let $E_{j,t}$ be conditional expectation
under $\mu_t$ given all variables outside coordinate $j$.  For fixed
$c_j\ge0$ and $\rho>0$, set
\begin{equation}
 -L_t=\rho(I-R_t^{\rm prob})+\sum_jc_j(I-E_{j,t}),
 \qquad R_t^{\rm prob}f=\mu_t(f).
 \label{eq:V7-kinetic-law}
\end{equation}
The unitary $U_tf=\sqrt{\mu_t}\,f$ from $L^2(\mu_t)$ to counting-measure
$\ell^2(\Omega)$ gives
\begin{equation}
 \mathsf H_t=U_t(-L_t)U_t^*
 =\rho(I-R_t)+\sum_jc_j(I-P_{j,t}),\qquad
 R_t=|s_t\rangle\langle s_t|.
 \label{eq:V7-H}
\end{equation}
For an outside configuration $z$, the block of $P_{j,t}$ is the rank-one
projection onto $r_{j,z,t}(x_j)=\sqrt{\mu_t(x_j\mid z)}$.  In particular
$\mathsf H_ts_t=0$, and $\mathsf H_t\succeq\rho(I-R_t)$.
All dots below denote a fixed real parameter direction; write $v,w$ for its
writer and centered writer.

\begin{theorem}[Conditional-score derivative of the actual kinetic Hamiltonian]
\label{thm:V7-conditional-derivative}
Define the conditional innovation
$\delta_jw=w-E_jw=v-E_jv$, viewed as a diagonal multiplication operator.
Then
\begin{align}
 \dot s&=-\tfrac12M_ws,
 &\dot R&=-\tfrac12(M_wR+RM_w),
 \label{eq:V7-root-derivative}\\
 \dot P_j&=-\tfrac12(M_{\delta_jw}P_j+P_jM_{\delta_jw}),
 \label{eq:V7-projection-derivative}\\
 \dot{\mathsf H}
 &=\tfrac\rho2(M_wR+RM_w)
   +\tfrac12\sum_jc_j(M_{\delta_jw}P_j+P_jM_{\delta_jw}).
 \label{eq:V7-actual-full-tangent}
\end{align}
The operator derivative is therefore a signed conditional-score insertion,
not in general $M_w$.  On each conditional block its projection derivative
has rank at most two and norm
\begin{equation}
 \|\dot P_j\|
 =\tfrac12\max_z
 \sqrt{\operatorname{Var}_{\mu(\cdot\mid z)}(v)}.
 \label{eq:V7-projection-norm}
\end{equation}
\end{theorem}

\begin{proof}
Differentiating \eqref{eq:V7-law-path} gives $\dot\mu=-w\mu$, hence the first
line.  Conditional normalization gives
$\partial_t\log\mu_t(x_j\mid z)|_0=-\delta_jw(x_j,z)$.  Therefore
$\dot r=-\frac12\delta_jw\,r$ and $\langle r,\dot r\rangle=0$ on each block.
Differentiating $rr^*$ proves \eqref{eq:V7-projection-derivative}, and then
\eqref{eq:V7-H} gives \eqref{eq:V7-actual-full-tangent}.
On the span of $r,\dot r$, when $\dot r\ne0$, the derivative $\dot r r^*+
r\dot r^*$ has eigenvalues $\pm\|\dot r\|$; it is zero on the orthogonal
complement.  Its norm is $\|\dot r\|=\frac12\sqrt{\operatorname{Var}(v\mid z)}$.
Taking the largest block norm proves \eqref{eq:V7-projection-norm}.
\end{proof}

\begin{remark}[Bounded-score continuous configuration version]
\label{rem:V7-continuous}
The same result holds on standard Borel product probability spaces for
bounded writers $v$ and the exponential deformation
\eqref{eq:V7-law-path}.  Use the fixed reference $L^2(\mu)$ and the unitary
multiplier $(d\mu_t/d\mu)^{1/2}$.  Disintegration over $z$ represents the
conditional projections as direct integrals of the same rank-one operators.
The bounded exponential has a uniform operator-norm Taylor remainder;
conditional normalizers stay bounded above and away from zero locally in
$t$.  Thus the differentiated rank-one formula passes in operator norm, and
the maximum in \eqref{eq:V7-projection-norm} becomes an essential supremum.
This applies to bounded native conditional writers without a thermal trace.
Unbounded Higgs or determinant-score directions still require their own
form-domain or integrability argument.
\end{remark}

\section{The moving centered subspace and the exact Fisher--Green bridge}
\label{sec:V7-centered}

The constant function in $L^2(\mu_t)$ becomes the moving vector $s_t$ after
square-root conjugation.  It must not be deleted using a fixed arbitrary
basis.  Put $Q_t=I-R_t$ and choose the minimal ground-line transport
\begin{equation}
 \dot W_t=[\dot R_t,R_t]W_t,\quad W_0=I,
 \qquad h_t=Q_0W_t^*\mathsf H_tW_tQ_0\big|_{\Ran Q_0}.
 \label{eq:V7-ground-transport}
\end{equation}
This is a declared connection for the diagnostic.  It does not purport to
identify every physical ESM connection with this one.

\begin{theorem}[Centered material derivative and conditional locality]
\label{thm:V7-centered-tangent}
The transport is unitary, $W_ts_0=s_t$, and $W_tQ_0W_t^*=Q_t$.  In its fixed
centered carrier the actual tangent is
\begin{equation}
 \boxed{V:=\dot h_0=Q\dot{\mathsf H}Q
 =\tfrac12\sum_jc_jQ(M_{\delta_jw}P_j+P_jM_{\delta_jw})Q.}
 \label{eq:V7-centered-actual}
\end{equation}
It satisfies
\begin{equation}
 \operatorname{Tr}V=0,
 \qquad
 \|V\|\le\tfrac12\sum_jc_j\max_z\sqrt{\operatorname{Var}(v\mid z)}
 \le\tfrac14\sum_jc_j\operatorname{osc}_j(v),
 \label{eq:V7-local-norm}
\end{equation}
where $\operatorname{osc}_j(v)$ is the largest range of $v$ when only $x_j$
varies.  In particular, only coordinates on which the writer depends
contribute, regardless of the strength of correlations in $\mu$.
\end{theorem}

\begin{proof}
The commutator $K=[\dot R,R]$ is skew-adjoint, so its evolution is unitary.
Differentiating $R^2=R$ gives $R\dot RR=0$ and
$[K,R]=\dot R$.  Hence $W_tR_0W_t^*=R_t$.  More directly,
$Ks=\dot s$, since $\langle s,\dot s\rangle=0$, so the specified positive
root is transported without an extra phase.
Differentiating $W_t^*\mathsf H_tW_t$ gives
$\dot{\mathsf H}+[\mathsf H,K]$.  The operator $K$ is off-diagonal between
$\Ran R$ and $\Ran Q$, while $\mathsf H=Q\mathsf HQ$.  Thus
$Q[\mathsf H,K]Q=0$.  The refresh derivative also vanishes after this
compression, giving \eqref{eq:V7-centered-actual}.

Every projection $P_{j,t}$ has constant rank, and $R_t$ has rank one, so
$\operatorname{Tr}\mathsf H_t$ is constant.  Also
$\langle s,\dot{\mathsf H}s\rangle=0$, by differentiating the zero ground
energy.  It follows that $\operatorname{Tr}V=0$.
Compression and the triangle inequality with
\eqref{eq:V7-projection-norm} give the first norm bound.  A real random
variable supported in an interval of length $b$ has variance at most
$b^2/4$: subtract the midpoint before minimizing over constants.  This gives
the second bound.  If $v$ does not depend on coordinate $j$, then
$E_jv=v$ and the corresponding term is exactly zero.
\end{proof}

\begin{corollary}[Literal volume-independent local coefficients in the finite model]
\label{cor:V7-three-local-bounds}
For the model's three local writers
\[
 v_g=-\operatorname{Re}u_0,\qquad
 v_4=\phi_0^4,\qquad
 v_{\rm hop}=-\operatorname{Re}u_0\,\phi_0\phi_1,
\]
where $u_j\in\mathbb Z_3$ and $\phi_j\in\{-1,0,1\}$, all local rates equal
one, and $N\ge2$, the actual centered tangents obey
\begin{equation}
 \boxed{\|V_g\|\le\frac38,\qquad
 \|V_4\|\le\frac14,\qquad
 \|V_{\rm hop}\|\le\frac{11}{8}.}
 \label{eq:V7-literal-three-bounds}
\end{equation}
These constants hold for every strictly positive law on that record set,
not merely for the two numerically tested sizes.
\end{corollary}

\begin{proof}
The oscillations of $v_g$ and $v_4$ are respectively $3/2$ and $1$ in one
coordinate.  The hopping writer has oscillations at most $3/2,2,2$ in its
link and two endpoint Higgs coordinates, and zero in the others.  Insert
these values in \eqref{eq:V7-local-norm}.
\end{proof}

\begin{theorem}[Exact native source metric from the Hamiltonian vacuum row]
\label{thm:V7-Fisher-Green}
For the full, uncompressed derivative define the source map
$\mathcal J e_a=\dot{\mathsf H}_as$.  Then
\begin{equation}
 \boxed{
 \mathcal J e_a=\tfrac12\mathsf H M_{w_a}s,\qquad
 \mathsf H^\dagger\mathcal J e_a=\tfrac12M_{w_a}s,\qquad
 4\mathcal J^*(\mathsf H^\dagger)^2\mathcal J=M.}
 \label{eq:V7-Fisher-Green}
\end{equation}
Thus $\Ker\mathcal J=\Ker M$.  For the native stationary half-density
experiment, let $B^{\rm stat}e_a=M_{w_a}s$ and
$Y^{\rm stat}e_a=\dot s_a$.  Then
\begin{equation}
 \boxed{Y^{\rm stat}=-\tfrac12B^{\rm stat},\qquad
 (Y^{\rm stat})^*(I-P_{\mathsf H}^{B^{\rm stat}})Y^{\rm stat}=0.}
 \label{eq:V7-stationary-follower}
\end{equation}
The same identities hold in the bounded-score continuous version, where the
refresh provides a bounded inverse on $\Ran Q$; no Gibbs trace is required.
\end{theorem}

\begin{proof}
Differentiate $\mathsf H_ts_t=0$ and use
\eqref{eq:V7-root-derivative}.  Since $\mu(w_a)=0$, the vector $M_{w_a}s$
lies in $\Ran Q$.  Applying $\mathsf H^\dagger$ and taking inner products
proves \eqref{eq:V7-Fisher-Green}; invertibility on $\Ran Q$ gives the
kernel statement.  The first identity in
\eqref{eq:V7-stationary-follower} is the differentiated square root of the
law.  Its range is already in the entrance source bank, hence also in the
semigroup-cyclic bank.  The continuous case uses exactly these bounded
operator and Hilbert-space identities.
\end{proof}

\begin{firewall}[The stationary root is not a Gibbs state of the kinetic operator]
\label{fire:V7-two-states}
The native Markov law is represented by the stationary vector $s=\sqrt\mu$.
The state $e^{-\beta h}/\operatorname{Tr}(e^{-\beta h})$ on the centered
kinetic carrier is an additional finite-temperature experiment.  Its
square-root derivative may have a positive dynamic residual even though
\eqref{eq:V7-stationary-follower} holds.  These are different states and
source syntheses, not contradictory ancestry verdicts.  The latter Gibbs
experiment is evaluated below as declared; it is not promoted to the
physical chiral ESM state simply because the kinetic path is reversible.
\end{firewall}

\section{An incidence error cannot be repaired by thermal dressing or a connection}
\label{sec:V7-incidence}

On a fixed $d$-dimensional centered carrier let $h=h^*\succeq0$, $a=\beta/2$,
$Z_\beta=\operatorname{Tr}e^{-\beta h}$, and
$\varrho_\beta=e^{-\beta h}/Z_\beta$.  Define the real-linear tangent map
\begin{equation}
 \mathcal T_{\beta,h}(X)
 =-Z_\beta^{-1/2}\int_0^a e^{-(a-t)h}
 (X-\operatorname{Tr}(\varrho_\beta X)I)e^{-th}\,dt.
 \label{eq:V7-thermal-tangent}
\end{equation}
This recalls, rather than re-proves, the physical Duhamel derivative already
used in V1--V6.

\begin{proposition}[Exact source-incidence test for the Gibbs root]
\label{prop:V7-tangent-injective}
For Hermitian $V,W$,
\begin{equation}
 \mathcal T_{\beta,h}(W)=\mathcal T_{\beta,h}(V)
 \quad\Longleftrightarrow\quad W-V\in\R I.
 \label{eq:V7-injective-mod-scalar}
\end{equation}
Furthermore
\begin{equation}
 \|\mathcal T_{\beta,h}(V)\|_{\rm HS}^2
 \le a^2\operatorname{Tr}(\varrho_\beta(V-\operatorname{Tr}(\varrho_\beta V)I)^2)
 \le a^2\|V\|^2.
 \label{eq:V7-thermal-ceiling}
\end{equation}
For $\tau\ge0$, the corresponding thermal action column satisfies
$\|Z_\beta^{-1/2}e^{-\tau h}V^\circ e^{-ah}\|_{\rm HS}\le\|V\|$.
\end{proposition}

\begin{proof}
In an eigenbasis of $h$, every off-diagonal entry of
\eqref{eq:V7-thermal-tangent} is a strictly nonzero real divided-difference
coefficient times $X_{ij}$.  This remains true inside a repeated eigenspace.
The diagonal coefficient is $-a\sqrt{\varrho_{\beta,ii}}$ times
$X_{ii}-\operatorname{Tr}(\varrho_\beta X)$.  The kernel therefore consists
exactly of the scalar matrices, proving \eqref{eq:V7-injective-mod-scalar}.

Let $q(x,y)=\int_0^a e^{-(a-t)x-ty}\,dt$.  Convexity of the exponential gives
$q(x,y)\le a(e^{-ax}+e^{-ay})/2$, and hence
$q(x,y)^2\le a^2(e^{-2ax}+e^{-2ay})/2$.  Sum this inequality against
$|V^\circ_{ij}|^2/Z_\beta$.  Hermiticity identifies the two sums and yields
the first inequality in \eqref{eq:V7-thermal-ceiling}.  The variance is at
most $\operatorname{Tr}(\varrho_\beta V^2)\le\|V\|^2$.  Finally,
$e^{-2\tau h}\preceq I$ gives the action-column bound by the same thermal
variance.  No raw Hilbert--Schmidt dimension bound is used.
\end{proof}

\begin{corollary}[Populated local source ceilings without raw trace growth]
\label{cor:V7-local-source-ceilings}
For the three local directions in \eqref{eq:V7-literal-three-bounds}, the
Gibbs targets at $\beta=1$ satisfy
\[
 \|Y_g\|_{\rm HS}\le\frac3{16},\qquad
 \|Y_4\|_{\rm HS}\le\frac18,\qquad
 \|Y_{\rm hop}\|_{\rm HS}\le\frac{11}{16}
\]
at every size on the fixed positive-support branch.  The thermal action-column bounds are respectively
$3/8,1/4,11/8$.  These are upper source bounds, not reserve floors, and do not
alone control a target-normalized omitted-energy tail.
\end{corollary}

\begin{proof}
Apply \eqref{eq:V7-thermal-ceiling} with $a=1/2$ and
\eqref{eq:V7-literal-three-bounds}.
\end{proof}

\begin{theorem}[Connection-independent spectral-diagonal incidence]
\label{thm:V7-connection-obstruction}
Let $D=W-V$, $h=\sum_\lambda\lambda E_\lambda$, and
$\mathcal D_h(D)=\sum_\lambda E_\lambda DE_\lambda$.  Then
\begin{equation}
 \boxed{
 \inf_{K^*=-K,\,c\in\R}\|D-[h,K]-cI\|_{\rm HS}^2
 =\left\|\mathcal D_h(D)-\frac{\operatorname{Tr}D}{d}I\right\|_{\rm HS}^2.}
 \label{eq:V7-connection-infimum}
\end{equation}
In particular, with $h_c=h-\operatorname{Tr}(h)I/d$ and $h_c\ne0$,
\begin{equation}
 \inf_{K,c}\|D-[h,K]-cI\|_{\rm HS}^2
 \ge\frac{|\operatorname{Tr}(h_cD)|^2}{\|h_c\|_{\rm HS}^2}.
 \label{eq:V7-trace-witness}
\end{equation}
Thus a nonzero trace witness cannot be erased by an infinitesimal unitary
connection or scalar energy recalibration.
\end{theorem}

\begin{proof}
All commutators $[h,K]$ have zero diagonal eigenspace blocks.  Conversely,
for any Hermitian off-diagonal $D$, setting
$E_\lambda K E_\mu=(\lambda-\mu)^{-1}E_\lambda DE_\mu$ for
$\lambda\ne\mu$ gives a skew-adjoint $K$ with $[h,K]=D$.  Orthogonality of
the diagonal and off-diagonal blocks in Hilbert--Schmidt norm reduces the
minimization to the diagonal part, where the best scalar is
$\operatorname{Tr}D/d$.  This proves \eqref{eq:V7-connection-infimum}.
Also $\operatorname{Tr}(h_c[h,K])=0$ by cyclicity and
$\operatorname{Tr}h_c=0$.  Cauchy--Schwarz then gives
\eqref{eq:V7-trace-witness}.
\end{proof}

\begin{countertheorem}[Exact three-record source mismatch]
\label{cth:V7-three-record}
Take $\Omega=\{-1,0,1\}$, $\mu_t(x)\propto e^{-tx}$, and a single full
heat bath.  Then $\mathsf H_t=(\rho+c)(I-R_t)$ and the transported centered
Hamiltonian is $(\rho+c)I_2$ for every $t$.  Its actual tangent and Gibbs
root derivative vanish.  At $t=0$, however, $W=QM_xQ|_{\Ran Q}$ is traceless
and $\|W\|_{\rm HS}^2=2/3$, so the inserted-potential target has
\begin{equation}
 \boxed{\|\mathcal T_{\beta,h}(W)\|_{\rm HS}^2=\frac{a^2}{3}>0.}
 \label{eq:V7-three-record-mismatch}
\end{equation}
At $\beta=1$ the squared mismatch is exactly $1/12$.
\end{countertheorem}

\begin{proof}
The unitary ground-line transport sends $Q_0$ to $Q_t$, leaving the scalar
centered Hamiltonian constant.  At the uniform law,
$Q=I-\mathbf1\mathbf1^*/3$.  Direct multiplication gives
$\operatorname{Tr}(QM_xQ)=0$ and
$\operatorname{Tr}((QM_xQ)^2)=2/3$.  The centered Gibbs density is $I_2/2$,
so its artificial tangent is $-aW/\sqrt2$.  This proves the formula.
\end{proof}

\section{An outward-certified obstruction on the original 81-record cylinder}
\label{sec:V7-exact-cylinder}

We now test the unchanged $N=2$ representative of the supplied report, with
\[
 (\beta_g,\kappa_H,\lambda_4,m^2,m_0,r,\kappa,y,m_M,y',\rho)
 =(3/5,2/5,7/20,1/2,1/2,1,3/5,4/5,3/5,1,1/2).
\]
There are $n=81$ records and $d=80$ centered states.  Use the local gauge
potential $v=-\operatorname{Re}u_0$ in \eqref{eq:V7-law-path}.  This is a
local coupling deformation, not a simultaneous change of every link
coefficient.  All four single-coordinate heat-bath rates are one.
The base measure is exactly
\begin{equation}
 \mu(U,\phi)=Z^{-1}e^{-S_B(U,\phi)}
 |\det D_\chi(U,\phi)|^2
 \prod_{j=0}^1(3/5+\phi_j)^2,
 \label{eq:V7-exact-law}
\end{equation}
with the printed action and coefficient, retaining both oriented
contributions when the two sites share their neighbours.

Write $p_j(x)=\mu(x_j\mid x_{j^c})$,
$\delta_jv(x)=v(x)-E_jv(x)$, $C_0=\rho+J$, and $J=4$.  The following
formulas evaluate the trace witness with no eigenvectors or square roots of
probabilities:
\begin{align}
 T_1&:=\operatorname{Tr}\mathsf H
       =\rho(n-1)+J(n-n/3)=256,\quad \alpha=T_1/d=16/5,
 \label{eq:V7-trace-constant}\\
 \mathsf H_{xx}&=C_0-\rho\mu(x)-\sum_jp_j(x),
 \label{eq:V7-diagonal}\\
 T_{\dot H}&:=\operatorname{Tr}(\mathsf H\dot{\mathsf H})
 =-\sum_{j<k}\sum_xp_j(x)p_k(x)(\delta_jv(x)+\delta_kv(x)),
 \label{eq:V7-trace-derivative}\\
 I_g&:=\operatorname{Tr}(h_c(W-V))
 =\sum_x(\mathsf H_{xx}-\alpha)(v(x)-\mu(v))-T_{\dot H},
 \label{eq:V7-Ig}\\
 T_2&:=\operatorname{Tr}\mathsf H^2
 =C_0^2n-2C_0(\rho+Jn/3)+\rho^2+Jn/3+2\rho J
    +2\sum_{j<k,x}p_j(x)p_k(x),
 \label{eq:V7-trace-square}\\
 \|h_c\|_{\rm HS}^2&=T_2-T_1^2/d,
 \qquad M_g=\operatorname{Var}_\mu(v).
 \label{eq:V7-trace-norm}
\end{align}

\begin{lemma}[Literal finite-sum trace formulas]
\label{lem:V7-trace-formulas}
Equations \eqref{eq:V7-trace-constant}--\eqref{eq:V7-trace-norm} hold exactly.
\end{lemma}

\begin{proof}
Every $P_j$ has rank $n/3$, and $RP_j=P_jR=R$.  For $j\ne k$, an entry
$(P_j)_{xy}(P_k)_{yx}$ can be nonzero only if $x=y$, because the two updates
change distinct single coordinates.  Hence
$\operatorname{Tr}(P_jP_k)=\sum_xp_j(x)p_k(x)$.  Expand
$\mathsf H=C_0I-\rho R-\sum_jP_j$ and its square to get
\eqref{eq:V7-trace-constant}, \eqref{eq:V7-diagonal}, and
\eqref{eq:V7-trace-square}.  Since $\dot p_j=-\delta_jv\,p_j$ and the other
terms of $T_2$ are constant, differentiating $T_2/2$ proves
\eqref{eq:V7-trace-derivative}.

The identities $\mathsf HR=0$ and $\operatorname{Tr}V=0$ give
$\operatorname{Tr}(hW)=\operatorname{Tr}(\mathsf H M_w)$,
$\operatorname{Tr}(hV)=\operatorname{Tr}(\mathsf H\dot{\mathsf H})$, and
$\operatorname{Tr}W=\sum_xw(x)$, because $\mu(w)=0$.
This proves \eqref{eq:V7-Ig}.  Removing the zero eigenvalue of $\mathsf H$
and subtracting its centered-carrier mean gives
\eqref{eq:V7-trace-norm}.
\end{proof}

\begin{theorem}[Certified failure of source--Hamiltonian incidence at $N=2$]
\label{thm:V7-certified-incidence}
For the cylinder \eqref{eq:V7-exact-law} and the gauge command above, one has
\begin{align}
 7.17401696931958014789070610155
 &<I_g<7.17401696931958014789070610157,
 \label{eq:V7-Ig-interval}\\
 70.2833518624719820252535532341
 &<\|h_c\|_{\rm HS}^2<70.2833518624719820252535532343,
 \label{eq:V7-Hc-interval}\\
 0.56249416528692981933478406615
 &<M_g<0.56249416528692981933478406617.
 \label{eq:V7-Mg-interval}
\end{align}
Consequently
\begin{equation}
 \boxed{
 \inf_{K^*=-K,\,c\in\R}
 \frac{\|W-V-[h,K]-cI\|_{\rm HS}^2}{80M_g}>0.01627.}
 \label{eq:V7-certified-mismatch}
\end{equation}
In particular the inserted V6 gauge column cannot represent the actual
heat-bath coupling tangent by a change of connection or scalar origin.
\end{theorem}

\begin{proof}[Computer-assisted finite proof with outward rational bounds]
Enumerate the 81 tuples in $\{0,1,2\}^4$ and substitute
$u\in\{1,(-1+i\sqrt3)/2,(-1-i\sqrt3)/2\}$ and
$\phi\in\{-1,0,1\}$.  Every bosonic action is rational with absolute value
at most $23/10$.  Exact four-by-four determinant arithmetic gives
$|\det D_\chi|^2\in\mathbb Q(\sqrt3)$; the replay verifies a strictly positive
lower bound on every record.  All Pfaffian factors are nonzero rationals.
Thus \eqref{eq:V7-exact-law} is a strictly positive finite law.

The accompanying exact certificate uses a common grid $10^{-50}$, rounding
both endpoints outwards after every addition, multiplication, and division.
For rational $x\ge0$, it encloses $e^x$ between the 121-term sum
$\sum_{k=0}^{120}x^k/k!$ and that sum plus
\[
 \frac{x^{121}}{121!}\frac{1}{1-x/122}.
\]
The bound follows because every subsequent term ratio is at most $x/122<1$.
Negative exponents are enclosed by reciprocal intervals.  The square root
of three is enclosed by the integer square root of $3\cdot10^{100}$ divided
by $10^{50}$ and its next grid point.  Exact determinant coefficients,
normalized weights, conditional probabilities, and all the finite sums in
\eqref{eq:V7-trace-constant}--\eqref{eq:V7-trace-norm} are then evaluated by
these interval operations, never by floating-point spectral decomposition.
They give the stated outward enclosures.  In particular their interval
quotient gives
\[
 \frac{I_g^2}{80\|h_c\|_{\rm HS}^2M_g}
 \in(0.01627287637853917887614960810,
       0.01627287637853917887614960811).
\]
Equation \eqref{eq:V7-trace-witness} proves
\eqref{eq:V7-certified-mismatch}.  The executable proof object consists of
\path{Renewal_SM_Einstein_Exact_Incidence_Certificate_V7.py} and its JSON
output, which includes the complete 81-record exact determinant/action table.
The claim is relative to exact integer/rational arithmetic and the displayed
finite algebra; it is not an interval claim about the separate eigenvalue
replay below.
\end{proof}

\section{Repaired finite-temperature experiment with all old primitives retained}
\label{sec:V7-replay}

For all five base-point writers of V6, take
$\mu_t\propto\mu\exp(-\sum_at_aw_a)$.  This path has the original local
bosonic first variations.  For the two determinant writers its first-order
orientation is the negative of increasing the printed global $\kappa$ or
$y$, since those V6 writers were $+\partial\log|\det D_\chi|^2$.
The exponential path is declared explicitly; no nonlinear coupling family is
inferred from one base-point writer.  Use the actual kinetic tangents
\eqref{eq:V7-centered-actual}, the ground-line connection
\eqref{eq:V7-ground-transport}, and $\beta=1,\tau=1/4$.

Importantly, the corrected ancestry bank does not delete the old columns.
If $W_a=QM_{w_a}Q$, and $F_\theta,F_{\rm Pf}$ are the two old scalar prior
operators, let $\Xi(X)=Z_\beta^{-1/2}e^{-\tau h}Xe^{-\beta h/2}$ and use
\begin{equation}
 \boxed{
 B^{\rm keep}=(\Xi(I),\ (\Xi(V_a^\circ))_{a=1}^5,
       (\Xi(W_a^\circ))_{a=1}^5,\Xi(F_\theta),\Xi(F_{\rm Pf})),
 \quad Y_a=\mathcal T_{\beta,h}(V_a).}
 \label{eq:V7-kept-bank}
\end{equation}
These thirteen columns form a comparison dictionary for the same five
commands; no new independent action debit is assigned to their deterministic
responses.  The scalar phase and Pfaffian sign are retained exactly as
previously declared.  They are not thereby proved to be covariant
fermion-line current operators.  No physical reserve of the enlarged
thirteen-column bank is reported without a corresponding coefficient metric.

\begin{proposition}[Finite numerical replay, not interval-certified spectral margins]
\label{prop:V7-numerical}
The deterministic binary64 replay of \eqref{eq:V7-kept-bank} gives the
following rounded values.  Dynamic fractions are the generalized
$R_{\rm dyn}/(Y^*Y)$ eigenvalues, and the left-energy head is
$\lambda\le\operatorname{gap}(h)+3/2$.
\begin{center}
\small
\begin{tabular}{@{}lcc@{}}
\toprule
Quantity & $N=2$ & $N=3$\\
\midrule
Configurations / centered dimension & $81/80$ & $729/728$\\
Centered gap & $1.110655$ & $1.294584$\\
Original inserted-target fraction range & $[0.016565,0.033616]$ & $[0.028038,0.068804]$\\
Actual tangent, old columns retained & $[0.017273,0.029116]$ & $[0.039740,0.066784]$\\
Actual target head-fraction range & $[0.009626,0.020232]$ & $[0.014784,0.029439]$\\
Omitted residual / $\|Y^*Y\|$ & $0.008713$ & $0.037608$\\
Retained fibres with positive computed trace & $26/26$ & $46/46$\\
Relative Fisher--Green matrix error & $1.73\cdot10^{-15}$ & $9.53\cdot10^{-16}$\\
\bottomrule
\end{tabular}
\end{center}
No finite positive margin in this table is asserted to be an outward
spectral certificate.  The exact incidence obstruction is instead
\Cref{thm:V7-certified-incidence}.
\end{proposition}

\begin{proof}[Reproducible numerical evaluation]
The standalone V7 script constructs the original finite cylinder and each
conditional square-root projection, differentiates it by
\eqref{eq:V7-actual-full-tangent}, and compresses to the centered spectral
basis.  It forms \eqref{eq:V7-kept-bank}, evaluates the divided difference
without subtractive exponential cancellation, and orthogonally shorts the
complete retained bank on every left eigenspace.  The script also separately
reproduces the old eight-column inserted-potential card; it does not compare
residual fractions belonging to an accidentally reduced prior bank.
The output contains the complete target and residual matrices, all fibre
ranks, and the numerical tolerances.
\end{proof}

An independent derivative check uses the command
$(1,-2,1/2,1/4,-1/8)$, normalized to unit $M$-length.
For a perturbed stationary root $s_t$, a centered orthonormal basis $B$
at zero transports to
\[
 B_t=B-\frac{(s_0+s_t)(s_t^*B)}{1+s_0^*s_t}.
\]
This is the minimal orthogonal rotation, and its derivative agrees with
\eqref{eq:V7-ground-transport} at zero.  Differentiating the actual matrices
$B_t^*\mathsf H_tB_t$ and their Gibbs square roots by symmetric secants gives:
\begin{center}
\small
\begin{tabular}{@{}cccc@{}}
\toprule
Size & Secant step & Relative true-root derivative error & Relative old-target mismatch\\
\midrule
$N=2$ & $0.004$ & $6.38\cdot10^{-6}$ & $1.215922$\\
$N=2$ & $0.0005$ & $9.98\cdot10^{-8}$ & $1.215923$\\
$N=3$ & $0.004$ & $9.10\cdot10^{-6}$ & $1.232238$\\
$N=3$ & $0.0005$ & $1.42\cdot10^{-7}$ & $1.232238$\\
\bottomrule
\end{tabular}
\end{center}
The reference norm in both last columns is the norm of the true Gibbs
root derivative.  The finite differences decrease quadratically for the
correct tangent but approach a nonzero discrepancy for the old insertion.
This is a numerical cross-check of source incidence, not the proof of
\eqref{eq:V7-certified-mismatch}.

\section{A sharp all-size obstruction in the frozen Dirac coefficient}
\label{sec:V7-Dirac}

A second upstream question is whether the frozen Dirac coefficient supplies
the uniform inverse bound used to control its global determinant writers.
This can be answered exactly on its own configuration set, without running
larger heat-bath matrices.

Keep $r=1$, $\kappa=3/5$, and $y=4/5$, but expose the bare scalar parameter
$m_0\ge0$:
\begin{equation}
 D_{m_0}(U,\phi)=(m_0+1)I+\tfrac35T(U)+\tfrac45Y(\phi).
 \label{eq:V7-D-family}
\end{equation}
These are the very operators of the consolidated report: $T$ has forward
block $P_L(u_j)(-I+i\sigma_1)/2$, backward block
$P_L(\bar u_j)(-I-i\sigma_1)/2$, and
$Y(\phi)=\operatorname{diag}_j(\phi_j\sigma_1)$.

\begin{theorem}[Exact uniform singular-value threshold on the declared record set]
\label{thm:V7-Dirac-threshold}
Let $u_j\in\mathbb Z_3$, $\phi_j\in\{-1,0,1\}$, and $N\ge2$ with periodic
boundary conditions.  For every $m_0\ge0$,
\begin{equation}
 \boxed{
 \inf_{N\ge2}\ \min_{U,\phi}\sigma_{\min}(D_{m_0}(U,\phi))
 =\max\left\{0,m_0+\frac15-\frac{3\sqrt2}{5}\right\}.}
 \label{eq:V7-sharp-threshold}
\end{equation}
In particular the frozen point $m_0=1/2$ has zero all-size infimum.  At that
point the homogeneous configurations $u_j=\phi_j=1$ satisfy
\begin{equation}
 \sigma_{\min}(D_{1/2}(\mathbf1,\mathbf1))
 \le\frac{3\sqrt2\pi}{5N}\quad\text{for every }N\ge2.
 \label{eq:V7-Dirac-rate}
\end{equation}
\end{theorem}

\begin{proof}
Each oriented hopping operator is a cyclic shift times unitary link factors
and a fixed matrix $(-I\pm i\sigma_1)/2$ of norm $1/\sqrt2$.  Thus
$\|T(U)\|\le\sqrt2$ for all links and sizes, and $\|Y(\phi)\|\le1$.
The reverse triangle inequality gives
\[
 \sigma_{\min}(D_{m_0})\ge
 \left(m_0+1-\frac{3\sqrt2}{5}-\frac45\right)_+.
\]

For homogeneous links and Higgs, the Fourier vector $e^{ikj}$ reduces the
coefficient to
\[
 D_{m_0}(k)=\left(m_0+1-\frac35\cos k\right)I
           +\left(\frac45+\frac35\sin k\right)\sigma_1,
 \qquad k=\frac{2\pi l}{N}.
\]
It is Hermitian, with eigenvalues
\begin{equation}
 d_+(k)=m_0+\frac95-\frac35\cos k+\frac35\sin k,
 \qquad
 d_-(k)=m_0+\frac15-\frac35(\cos k+\sin k).
 \label{eq:V7-Fourier-bands}
\end{equation}
When $m_0+1/5\ge3\sqrt2/5$, $k=\pi/4$ attains the lower bound, and is
available whenever $8$ divides $N$.  When
$0\le m_0+1/5<3\sqrt2/5$, the continuous minus band has a zero and the
periodic grids approximate it, so the infimum is zero.  This proves
\eqref{eq:V7-sharp-threshold}, including equality at the threshold.

At $m_0=1/2$, one zero $k_*$ is specified by
\begin{equation}
 \cos k_*=\frac{7+\sqrt{23}}{12},\qquad
 \sin k_*=\frac{7-\sqrt{23}}{12}.
 \label{eq:V7-exact-zero-angle}
\end{equation}
The nearest grid point is within $\pi/N$, and
$|d_-'(k)|\le3\sqrt2/5$.  The mean-value theorem proves
\eqref{eq:V7-Dirac-rate}.
\end{proof}

\begin{corollary}[Every homogeneous finite determinant is nonzero at the frozen point]
\label{cor:V7-nonzero-no-uniform}
At $m_0=1/2$ the homogeneous determinant is nonzero for every finite $N$, but
its inverse norm is at least $5N/(3\sqrt2\pi)$.  These configurations have
strictly positive weight in the finite cylinder.
\end{corollary}

\begin{proof}
The plus band is at least $23/10-3\sqrt2/5>0$.  A zero of the minus band at
$k=2\pi l/N$ would imply
$2(\cos k+\sin k)=7/3$.  With $\zeta=e^{ik}$ this reads
\[
 (1-i)\zeta+(1+i)\zeta^{-1}=7/3.
\]
The left side is an algebraic integer, because $\zeta$ and $i$ are roots of
unity.  A rational algebraic integer is an integer: substituting a reduced
fraction into a monic integer polynomial proves its denominator must be one.
Thus $7/3$ is impossible.  All band eigenvalues are nonzero, while
\eqref{eq:V7-Dirac-rate} gives the inverse bound.  The bosonic exponential is
positive and the homogeneous Pfaffian factor is $(8/5)^N\ne0$, so these
records are not removed by a zero weight.
\end{proof}

\begin{corollary}[An explicit repair threshold, not an adopted new action]
\label{cor:V7-mass-repair}
The strict sufficient and sharp uniform mass condition for this record set is
$m_0>(3\sqrt2-1)/5$.  For example, $m_0=1$ gives the exact all-size floor
$(6-3\sqrt2)/5>0.3514$.  At the threshold, the homogeneous $N=8$ coefficient
already has a zero mode.  Any such change of $m_0$ changes the source law and
requires a fresh ancestry evaluation; it is not a normalization of the frozen
$m_0=1/2$ card.
\end{corollary}

\begin{proof}
This is \eqref{eq:V7-sharp-threshold} and the attaining mode $k=\pi/4$.
\end{proof}

\begin{firewall}[What the inverse obstruction does and does not establish]
\label{fire:V7-Dirac-scope}
The theorem rules out a uniform all-record singular-value floor for the
specified finite-state representative.  It does not rule out source-weighted
integrability: a rare nearly singular record can have a small unquenched
weight, and the determinant/Pfaffian zero-safe calculus must retain that
weight before any conclusion about the source metric.  Nor does the theorem
assert a physical particle mass or a four-dimensional Standard-Model
obstruction.  Its role is to replace an extrapolated inverse bound by a
literal sharp spectral calculation at the unchanged coupling point.
\end{firewall}

\section{V7 conclusion: the first unresolved arrow has moved upstream}
\label{sec:V7-conclusion}

The established output is not another conditionally sufficient lag packet.
It is the actual derivative of the declared kinetic law, its exact
stationary Fisher--Green incidence, three volume-independent local source
ceilings, a certified nonzero mismatch in the prior finite replay, and the
sharp all-size inverse threshold of its Dirac coefficient.

The original numerical birth evidence must now be named by its experiment:
the inserted-potential card, the repaired kinetic-Gibbs card, and the native
stationary-root card have different sources.  Their respective outcomes are
positive finite residuals in the first two tested examples and the exact
stationary follower identity \eqref{eq:V7-stationary-follower}.  None is
silently substituted for the original chiral ESM source.

For the original continuum route, the remaining upstream calculation is to
identify and differentiate its actual common-action Hamiltonian in its
physical connection, and to compare that derivative with the action and
line-source bank before applying V5 or V6.  On the Markov-conditional branch,
\eqref{eq:V7-actual-full-tangent} already supplies the derivative from the
same conditional law.  On a separately reconstructed chiral transfer or
Friedrichs Hamiltonian branch, its own derivative must be used instead.
The common source-metric, line/clock, Duhamel, and cofinal theorems remain
available without being rebuilt.  The frozen toy mass is not changed in this
version, and a uniform inverse is not claimed where
\eqref{eq:V7-sharp-threshold} excludes it.

\begin{assessment}[Status of V7]
\label{ass:V7-status}
\passstatus{} for the proved conditional kinetic tangent, moving-centered
compression, vacuum Fisher--Green identity, stationary follower, local source
ceilings, exact source-incidence and connection tests, and sharp Dirac
threshold.  The $N=2$ trace witness has an outward exact-arithmetic
certificate.  The corrected $N=2,3$ spectral and finite-difference tables are
explicitly numerical diagnostics only.

\obsstatus{} for identifying the old multiplication-source replay with the
actual heat-bath deformation, even after rechoosing a connection or scalar
energy origin; also for a uniform all-record Dirac inverse at the frozen
couplings.  The native stationary response is fully source-accounted and
does not create a new source merely by being represented again.

The full ESM state/clock incidence, source-weighted near-divisor estimates,
physical regulator transport, represented current/stress, and interacting
Einstein--Standard--Model continuum are not established by this version.
No later weighted-Gram certificate is used to repair the earlier incidence
failure.
\end{assessment}

\section*{V7 provenance and append-only record}
\addcontentsline{toc}{section}{V7 provenance and append-only record}
The complete input V6 source has SHA--256
\begin{center}\small\ttfamily
0d7920a915a8cfa83ba2e2be47a9bb3bc235a3415d517a8a365543c9b2fc3c29.
\end{center}
Every byte before its sole terminal document-closing command is preserved.
The accompanying two standalone Python programs distinguish exact algebra
and outward integer intervals from binary64 diagnostics.  Their JSON files
retain the complete exact 81-record table, interval endpoints, source and
residual matrices, and the independent finite-difference comparison.  The
new qualification of V6 is append-only and does not alter its original text.

\begin{thebibliography}{9}
\bibitem{V7SMF}
C.~Castelnovo, C.~Chamon, C.~Mudry, and P.~Pujol,
\emph{From quantum mechanics to classical statistical physics: generalized
Rokhsar--Kivelson Hamiltonians and the stochastic matrix form decomposition},
Annals of Physics \textbf{318} (2005), 316--344,
arXiv:cond-mat/0502068.  Cited only for the established stochastic-Hamiltonian
background; the conditional derivative and the certificate used here are
proved in this version.
\end{thebibliography}

% =============================================================================
% END APPENDED VERSION V7 MATERIAL
% =============================================================================


% =============================================================================
% BEGIN APPENDED VERSION V8 MATERIAL
% =============================================================================
\clearpage
\section*{Version V8: centered Hamiltonian-tangent rigidity, physical connection, and a gain-free transfer-log acquisition}
\addcontentsline{toc}{section}{Version V8: centered Hamiltonian tangent and transfer-log acquisition}

\begin{abstract}
The supplied Navier--Stokes and Yang--Mills continuations identify the same
upstream failure of order in two different forms.  In the Navier--Stokes
card, an Eulerian cross-time pairing can have the wrong sign until the actual
material transport is inserted.  In the reduced-fibre Yang--Mills card,
normalized laws, centered scores, Witten currents, and all of their mixed
Grams can remain fixed while a fibre-independent unnormalized action germ
changes the physical Hessian arbitrarily.  The boundary calculation further
shows that a scalar singlet source cannot be substituted for an adjoint
half-score merely because both admit positive covariance cards.

Version~V7 found the corresponding Standard-Model issue: the derivative of
the stationary law fixes only the vacuum row of the Hamiltonian derivative.
This continuation proves the exact missing statement.  For a self-adjoint
Hamiltonian family with a simple transported ground state, every tangent has
the orthogonal block decomposition
\[
 \dot H=|r\rangle\langle s|+|s\rangle\langle r|+A,
 \qquad r=-H\dot s,
 \qquad A=Q\dot H Q.
\]
The stationary Fisher--Green card determines $r$ and is completely blind to
the centered block $A$.  This is not only abstract nonidentifiability: every
self-adjoint $A$ annihilating the constant vector occurs as the tangent of a
strictly positive reversible finite-state Markov generator whose stationary
law is unchanged.  An explicit four-state pair has the same stationary path,
the same base Hamiltonian, and equal raw tangent Hilbert--Schmidt norm, but
one physical thermally dressed Duhamel card is a fibrewise follower and the
other has a strictly positive final residual.

The remaining ambiguity is split correctly.  Changing the physical Hilbert-
bundle transport changes $A$ by $[h,K]$.  The block-diagonal resonant part
$\sum_\lambda E_\lambda A E_\lambda$ is the exact quotient invariant, while
every off-resonant self-adjoint block is a connection commutator.  Thus a
physical connection is required before the full Duhamel source is meaningful;
without it only the resonant class is intrinsic.

Finally, the centered tangent is acquired directly from positive finite-time
transfer matrices.  On a physically transported finite carrier,
$-\ell^{-1}\log P_{\theta,\ell}$ is the pulled Hamiltonian for every positive
lag.  A symmetric parameter secant therefore reconstructs the tangent with a
sharp $O(\varepsilon^2)$ error.  Unknown positive scalar gains cancel after
passing to the scalar quotient, which is sufficient because the Gibbs source
uses the centered tangent.  A second lag is a held-out clock-incidence test.
A counterexample proves that one may not take a matrix logarithm after an
unproved source compression.  The resulting tangent radius propagates
explicitly into the already established V5 supported-Schur compiler.

The original Einstein--Standard--Model continuum still lacks one same-history
common-action transfer family and its physical connection on the active
chiral reflected carrier.  Version~V8 identifies that finite acquisition and
stops there; it does not add another lag, concentration, current, stress, or
Einstein compiler.
\end{abstract}

% =============================================================================
\section{Direction after the supplied NS and Yang--Mills continuations}
\label{sec:V8-direction}
% =============================================================================

The new sector notes do not provide numerical Standard-Model matrices, and no
Navier--Stokes or Yang--Mills coefficient is imported into the present card.
They do, however, sharpen the order of attack.

\begin{enumerate}[label=\textup{(U\arabic*)},leftmargin=3.4em]
\item The fixed-datum Navier--Stokes continuation moves upstream from a
      reflected cross Gram to the actual material--caloric localization.  A
      later two-sided Gram cannot repair a missing same-history spatial route.
\item The terminal Navier--Stokes continuation constructs smooth finite-record
      examples with negative untransported pairings and shows that the actual
      material pullback restores the physical comparison.  Metric near-isometry
      alone does not align the two states.
\item The reduced-fibre Yang--Mills continuation proves that an arbitrary
      fibre-independent quadratic action changes the physical partition
      Hessian while leaving every normalized score, Witten-current, return,
      and boundary-incidence packet unchanged.
\item The Yang--Mills boundary continuation eliminates a scalar proxy for an
      adjoint current source by an exact colour-selection rule.  Source type
      and same-history incidence precede a covariance short.
\end{enumerate}

The Standard-Model analogue is therefore not another refinement of the V5
mixed-clock quadrature or V6 weighted absorption moment.  It is the centered
operator tangent
\begin{equation}
 \boxed{A=Q\dot H Q}
 \label{eq:V8-centered-block}
\end{equation}
on the physically transported Hilbert bundle.  Version~V7 computes this block
for one conditional heat-bath law, but its stationary Fisher identity uses
only $\dot Hs$.  The next theorem separates those two pieces exactly.

\begin{firewall}[No cross-sector coefficient transfer]
\label{fire:V8-no-cross-sector-transfer}
The supplied NS and YM notes are used only to identify the logical order:
actual transport and unnormalized or operator-valued incidence precede a
normalized residual.  Their material router, pressure current, Wilson
partition ratio, colour current, and radial channel are not Standard-Model
sources.  Every formula below is proved independently on the SM Hamiltonian
card.
\end{firewall}

% =============================================================================
\section{The ground-state row does not determine the centered Hamiltonian tangent}
\label{sec:V8-ground-block}
% =============================================================================

Let $E$ be a finite-dimensional real command space.  For $u\in E$, let
$\dot H(u)$ denote the derivative at zero of a $C^1$ family of self-adjoint
Hamiltonians after transport to one fixed finite Hilbert space $\mathcal K$.
Assume that $H=H_0\succeq0$ has a simple normalized ground vector $s$,
\begin{equation}
 Hs=0,
 \qquad \|s\|=1,
 \qquad P=|s\rangle\langle s|,
 \qquad Q=I-P,
 \qquad h=QHQ.
 \label{eq:V8-ground-data}
\end{equation}
Choose the phase of the transported ground section so that
$\langle s,\dot s(u)\rangle=0$.

\begin{theorem}[Exact vacuum-row and centered-block decomposition]
\label{thm:V8-ground-decomposition}
For every command $u$, put
\begin{equation}
 r(u):=-H\dot s(u)=\dot H(u)s,
 \qquad
 A(u):=Q\dot H(u)Q.
 \label{eq:V8-rA}
\end{equation}
Then $r(u)\in\Ran Q$, $A(u)=A(u)^*$ on $\Ran Q$, and
\begin{equation}
 \boxed{
 \dot H(u)=|r(u)\rangle\langle s|
          +|s\rangle\langle r(u)|+A(u).}
 \label{eq:V8-block-decomposition}
\end{equation}
Conversely, after $H$, $s$, and $\dot s(u)$ are fixed, every self-adjoint
$A(u)=QA(u)Q$ gives a self-adjoint first-order Hamiltonian tangent satisfying
the differentiated ground-state equation.

The Hilbert--Schmidt source Gram has the orthogonal split
\begin{equation}
 \boxed{
 \Tr\bigl(\dot H(u)\dot H(v)\bigr)
 =2\operatorname{Re}\langle r(u),r(v)\rangle
  +\Tr\bigl(A(u)A(v)\bigr).}
 \label{eq:V8-HS-source-split}
\end{equation}
In particular,
\begin{equation}
 \boxed{\|\dot H(u)\|_{\rm HS}^2
 =2\|r(u)\|^2+\|A(u)\|_{\rm HS}^2.}
 \label{eq:V8-HS-Pythagoras}
\end{equation}
\end{theorem}

\begin{proof}
Differentiate $H_ts_t=0$ at zero:
\[
 \dot H(u)s+H\dot s(u)=0.
\]
Thus $\dot H(u)s=-H\dot s(u)=r(u)$.  Since $H\dot s(u)$ is orthogonal to
$s$, one has $r(u)\in\Ran Q$.  Moreover,
\[
 P\dot H(u)P
 =\langle s,\dot H(u)s\rangle P
 =-\langle Hs,\dot s(u)\rangle P=0.
\]
The $QP$ block is $|r(u)\rangle\langle s|$ and self-adjointness fixes the
$PQ$ block.  The remaining $QQ$ block is $A(u)$, proving
\eqref{eq:V8-block-decomposition}.

Conversely, the right side of \eqref{eq:V8-block-decomposition} is
self-adjoint and sends $s$ to $r(u)=-H\dot s(u)$; therefore it satisfies the
linearized ground-state equation.  The three block subspaces $QP$, $PQ$, and
$QQ$ are mutually orthogonal in the Hilbert--Schmidt product.  The two rank-one
blocks each have inner product $\langle r(u),r(v)\rangle$ with their
corresponding block, while the mixed block products have zero trace.  This
proves \eqref{eq:V8-HS-source-split} and
\eqref{eq:V8-HS-Pythagoras}.
\end{proof}

\begin{corollary}[Exact limit of the V7 stationary Fisher card]
\label{cor:V8-Fisher-only-root-row}
The V7 source map $\mathcal J u=\dot H(u)s$ is exactly $r(u)$.  Hence its
identity
\[
 4\mathcal J^*(H^\dagger)^2\mathcal J=\mathsf M
\]
reconstructs the stationary score metric from the vacuum row, but contains no
information about the positive Gram
\begin{equation}
 \boxed{\mathsf G_A(u,v)=\Tr(A(u)A(v)).}
 \label{eq:V8-centered-tangent-Gram}
\end{equation}
If the stationary law is constant, then $\dot s=0$ and $\mathcal J=0$, while
$A$ may still be arbitrary.
\end{corollary}

\begin{proof}
The first statement is \eqref{eq:V8-rA}.  The V7 Fisher--Green identity is a
quadratic identity in $\mathcal J$ and therefore in $r$ alone.  The final
statement follows from the converse part of
\Cref{thm:V8-ground-decomposition}.
\end{proof}

\begin{assessment}[The first missing source row]
\label{ass:V8-first-missing-row}
A stationary-law derivative, even together with the complete base Hamiltonian,
does not determine the finite-temperature Gibbs-root Duhamel source.  That
source depends on $A=Q\dot H Q$.  The centered block must therefore be acquired
from the actual common-action dynamics or an equivalent same-history transfer
experiment before any fibre residual is interpreted.
\end{assessment}

% =============================================================================
\section{Every centered tangent occurs in a reversible Markov family}
\label{sec:V8-Markov-realization}
% =============================================================================

The preceding freedom is realized inside the same finite reversible class used
by the conditional-law construction; it is not an arbitrary matrix extension.
Let $\mathbf 1=(1,\ldots,1)^\mathsf T\in\R^d$ and, for $i<j$, put
\begin{equation}
 L_{ij}=(e_i-e_j)(e_i-e_j)^\mathsf T.
 \label{eq:V8-edge-Laplacian}
\end{equation}

\begin{theorem}[Complete conductance realization of the centered tangent]
\label{thm:V8-conductance-realization}
Let $A=A^\mathsf T\in M_d(\R)$ satisfy $A\mathbf1=0$.  Then there is a unique
edge-coefficient family
\begin{equation}
 a_{ij}=-A_{ij},\qquad i<j,
 \label{eq:V8-conductance-coefficients}
\end{equation}
for which
\begin{equation}
 \boxed{A=\sum_{i<j}a_{ij}L_{ij}.}
 \label{eq:V8-Laplacian-span}
\end{equation}
Let
\begin{equation}
 H_0=\sum_{i<j}c_{ij}L_{ij},
 \qquad c_{ij}>0.
 \label{eq:V8-positive-complete-Laplacian}
\end{equation}
Then, for all sufficiently small $|t|$,
\begin{equation}
 \boxed{H_t=H_0+tA
 =\sum_{i<j}(c_{ij}+ta_{ij})L_{ij}}
 \label{eq:V8-Markov-family}
\end{equation}
is the square-root representation of an irreducible symmetric continuous-time
Markov generator with the same uniform stationary law for every such $t$.
Its Hamiltonian tangent is $A$ and its stationary root is constant.
\end{theorem}

\begin{proof}
The $(i,j)$ off-diagonal entry of $L_{ij}$ is $-1$, and no other edge
Laplacian contributes to that entry.  Thus any representation must have
$a_{ij}=-A_{ij}$.  With this choice, the $i$th diagonal entry of the sum is
\[
 \sum_{j\ne i}a_{ij}=-\sum_{j\ne i}A_{ij}=A_{ii},
\]
because $A\mathbf1=0$.  This proves existence and uniqueness in
\eqref{eq:V8-Laplacian-span}.

There are finitely many edges, so $c_{ij}+ta_{ij}>0$ for all $i<j$ when
$|t|$ is small enough.  The matrix $H_t$ is then a connected weighted graph
Laplacian: it is positive semidefinite, has one-dimensional kernel
$\Span\{\mathbf1\}$, and $-H_t$ has nonnegative off-diagonal entries and zero
row sums.  Symmetry gives reversibility for the uniform law, which is
independent of $t$.
\end{proof}

\begin{countertheorem}[Stationary probability data do not determine the Gibbs Hamiltonian source]
\label{cth:V8-stationary-no-centered-tangent}
Fix any strictly positive complete-graph Hamiltonian $H_0$.  The complete
stationary law, every stationary-law derivative, the V7 stationary
Fisher--Green source metric, and the base Hamiltonian can all be identical,
while the centered tangent $A=Q\dot H Q$ is any prescribed real self-adjoint
operator on the real centered carrier $\Ran Q$.  Consequently those stationary data do not determine
any Duhamel card whose target is the Gibbs-root derivative of $H_t$.
\end{countertheorem}

\begin{proof}
Use \Cref{thm:V8-conductance-realization} for two distinct choices of $A$.
Every family has the same constant stationary root and the same value $H_0$ at
$t=0$, so all stationary derivatives and their Fisher--Green Grams agree and
vanish.  The centered Hamiltonian tangents are the independently prescribed
matrices.
\end{proof}

% =============================================================================
\section{Physical transport and the exact connection quotient}
\label{sec:V8-connection}
% =============================================================================

The centered block is physical only after the Hilbert-bundle transport has
been fixed.  Let $U_t:\mathcal K_0\to\mathcal K_t$ be the declared unitary
transport and let $\widetilde H_t=U_t^*H_tU_t$.  A different transport with
the same value at zero has the form $U_t'=U_tV_t$, where $V_0=I$ and
$K=\dot V_0$ is anti-self-adjoint.  We restrict to transports preserving the
ground line, so $K=QKQ$.

\begin{theorem}[Resonant tangent is the exact connection-independent quotient]
\label{thm:V8-connection-quotient}
Let
\begin{equation}
 h=\sum_{\lambda\in\spec h}\lambda E_\lambda
 \label{eq:V8-h-spectral}
\end{equation}
be the centered base Hamiltonian.  If $A=Q\dot{\widetilde H}_0Q$, then changing
the transport as above sends
\begin{equation}
 \boxed{A\longmapsto A+[h,K].}
 \label{eq:V8-connection-action}
\end{equation}
Define the resonant projection
\begin{equation}
 \boxed{\mathcal R_h(A)=\sum_\lambda E_\lambda A E_\lambda.}
 \label{eq:V8-resonant-projection}
\end{equation}
Then $\mathcal R_h(A)$ is invariant under
\eqref{eq:V8-connection-action}.  Conversely, with
\begin{equation}
 K_A=\sum_{\lambda\ne\mu}
       \frac{E_\lambda A E_\mu}{\lambda-\mu},
 \label{eq:V8-KA}
\end{equation}
one has $K_A^*=-K_A$ and
\begin{equation}
 \boxed{A-\mathcal R_h(A)=[h,K_A].}
 \label{eq:V8-offresonant-commutator}
\end{equation}
Hence
\begin{equation}
 \boxed{
 \inf_{K^*=-K}\|A-[h,K]\|_{\rm HS}
 =\|\mathcal R_h(A)\|_{\rm HS}.}
 \label{eq:V8-connection-distance}
\end{equation}
The infimum is attained, and the quotient of centered self-adjoint tangents by
connection commutators is canonically the self-adjoint commutant of $h$.
\end{theorem}

\begin{proof}
Differentiating
\[
 (U_t')^*H_tU_t'=V_t^*\widetilde H_tV_t
\]
at zero gives $\dot{\widetilde H}_0+[H_0,K]$; compression by $Q$ proves
\eqref{eq:V8-connection-action}.  Every commutator $[h,K]$ has zero
$(\lambda,\lambda)$ spectral block, so
$\mathcal R_h([h,K])=0$.

For \eqref{eq:V8-KA}, self-adjointness of $A$ gives
\[
 (E_\mu K_AE_\lambda)^*
 =\frac{E_\lambda A E_\mu}{\mu-\lambda}
 =-E_\lambda K_AE_\mu,
\]
so $K_A$ is anti-self-adjoint.  Its off-diagonal spectral block satisfies
\[
 E_\lambda[h,K_A]E_\mu
 =(\lambda-\mu)E_\lambda K_AE_\mu
 =E_\lambda A E_\mu,
 \qquad \lambda\ne\mu,
\]
and its diagonal blocks vanish, proving
\eqref{eq:V8-offresonant-commutator}.  Resonant and off-resonant matrix blocks
are Hilbert--Schmidt orthogonal.  Every $[h,K]$ is off-resonant and $K_A$
removes the entire off-resonant part, which proves
\eqref{eq:V8-connection-distance}.
\end{proof}

\begin{firewall}[Changing a connection is not a coordinate-free proof of followerhood]
\label{fire:V8-connection-not-follower}
A mere change of coordinates transports both the Hamiltonian family and its
already declared connection and leaves the covariant tangent unchanged.  The
freedom in \Cref{thm:V8-connection-quotient} compares different physical
transport rules when no connection has yet been supplied.  Once the active
field/OS construction fixes its common-action transport, its full $A$ is the
physical source and may not be replaced by the resonant representative after
reading the residual.  Without that transport, only
$\mathcal R_h(A)$ is licensed.
\end{firewall}

% =============================================================================
\section{A four-state exact branch separator}
\label{sec:V8-four-state}
% =============================================================================

We now show that the missing centered block changes the actual fibrewise
Duhamel verdict, even when the stationary path and a natural raw tangent size
are fixed.  In $\R^4$ put
\begin{equation}
 \begin{aligned}
 s&=\tfrac12(1,1,1,1)^\mathsf T,\\
 e_1&=\tfrac12(1,1,-1,-1)^\mathsf T,\\
 e_2&=\tfrac12(1,-1,1,-1)^\mathsf T,\\
 e_3&=\tfrac12(1,-1,-1,1)^\mathsf T.
 \end{aligned}
 \label{eq:V8-Hadamard-basis}
\end{equation}
These vectors form an orthonormal basis.  Let
\begin{equation}
 \begin{aligned}
 H_0&=4|e_1\rangle\langle e_1|
     +5|e_2\rangle\langle e_2|
     +6|e_3\rangle\langle e_3|,\\
 A^{(0)}&=|e_1\rangle\langle e_1|-|e_2\rangle\langle e_2|,\\
 A^{(1)}&=\frac1{\sqrt3}
  \sum_{\substack{1\le i,j\le3\\i\ne j}}
   |e_i\rangle\langle e_j|.
 \end{aligned}
 \label{eq:V8-two-tangents}
\end{equation}
In the record basis,
\begin{equation}
 H_0=\frac14
 \begin{pmatrix}
 15&-7&-5&-3\\
 -7&15&-3&-5\\
 -5&-3&15&-7\\
 -3&-5&-7&15
 \end{pmatrix}.
 \label{eq:V8-H0-matrix}
\end{equation}
Thus $H_0$ is a complete-graph Markov Hamiltonian with edge conductances
$7/4,5/4,3/4,3/4,5/4,7/4$.

Fix $\beta>0$, $a=\beta/2$, and $\tau>0$, and put
\begin{equation}
 Z_\beta=1+e^{-4\beta}+e^{-5\beta}+e^{-6\beta}.
 \label{eq:V8-Z4}
\end{equation}
For a centered tangent $A$, define its physical action primitive and Gibbs-root
target in the Hilbert--Schmidt carrier by
\begin{equation}
 \begin{aligned}
 B_A&=Z_\beta^{-1/2}e^{-\tau H_0}Ae^{-aH_0},\\
 Y_A&=-Z_\beta^{-1/2}\int_0^a
       e^{-(a-q)H_0}Ae^{-qH_0}\,\dd q,\\
 B_{\rm wt}&=Z_\beta^{-1/2}e^{-(\tau+a)H_0}.
 \end{aligned}
 \label{eq:V8-four-state-BY}
\end{equation}
Resolve these vectors by left energy.  On the left-energy-$x$ row let
$P_x^A$ be the projection onto the span of the two retained primitive rows
$(B_{{\rm wt},x},B_{A,x})$, and define the final fibre residual
\begin{equation}
 \mathfrak R_\beta(A)
 =\sum_{x\in\{4,5,6\}}
   \|\,(I-P_x^A)Y_{A,x}\,\|_{\rm HS}^2.
 \label{eq:V8-four-state-residual}
\end{equation}
There is no line column in this one-sector control.

\begin{countertheorem}[Same stationary path and tangent norm, opposite Duhamel branches]
\label{cth:V8-four-state-branch}
For all sufficiently small $|t|$, both
\[
 H_t^{(j)}=H_0+tA^{(j)},\qquad j=0,1,
\]
are irreducible symmetric Markov Hamiltonians with the same constant uniform
stationary law.  Moreover,
\begin{equation}
 \boxed{\|A^{(0)}\|_{\rm HS}^2
       =\|A^{(1)}\|_{\rm HS}^2=2.}
 \label{eq:V8-equal-tangent-size}
\end{equation}
Nevertheless,
\begin{equation}
 \boxed{\mathfrak R_\beta(A^{(0)})=0,}
 \label{eq:V8-follower-control}
\end{equation}
whereas $\mathfrak R_\beta(A^{(1)})>0$ for every $\beta,\tau>0$.
More explicitly, set $\lambda_1=4$, $\lambda_2=5$, $\lambda_3=6$ and
\begin{equation}
 m_x(y)=-e^{\tau x}\int_0^a e^{q(y-x)}\,\dd q,
 \qquad
 w_{ij}=\frac{e^{-2\tau\lambda_i-2a\lambda_j}}{3Z_\beta}.
 \label{eq:V8-mw}
\end{equation}
If $\{i,j,k\}=\{1,2,3\}$, then
\begin{equation}
 \boxed{
 \mathfrak R_\beta(A^{(1)})
 =\sum_{i=1}^3
   \frac{w_{ij}w_{ik}}{w_{ij}+w_{ik}}
   \abs{m_{\lambda_i}(\lambda_j)
        -m_{\lambda_i}(\lambda_k)}^2>0.}
 \label{eq:V8-positive-four-state-residual}
\end{equation}
At $\beta=1$ and $\tau=1/4$ the three fibre contributions are, respectively,
\begin{equation}
 \begin{aligned}
 2.60280177327448694\times10^{-5},\qquad
 4.61455993261561298\times10^{-5},\qquad
 9.57517261832257374\times10^{-6},
 \end{aligned}
 \label{eq:V8-four-state-values}
\end{equation}
with total
\begin{equation}
 \boxed{\mathfrak R_1(A^{(1)})
 =8.17487896772235729\times10^{-5}.}
 \label{eq:V8-four-state-total}
\end{equation}
The displayed decimal is a high-precision diagnostic; positivity is proved by
\eqref{eq:V8-positive-four-state-residual}.
\end{countertheorem}

\begin{proof}
Both tangents are symmetric and annihilate $s$.  The Markov realization follows
from \Cref{thm:V8-conductance-realization}; the positive base conductances
remain positive for small $|t|$.  Orthogonality of the matrix units in the
$e_i$ basis gives \eqref{eq:V8-equal-tangent-size}.

The tangent $A^{(0)}$ is diagonal in the right-energy basis.  On every nonzero
left row, $B_{A^{(0)},x}$ and $Y_{A^{(0)},x}$ have the same single right-energy
support, so the target belongs to the action-primitive line.  This proves
\eqref{eq:V8-follower-control}.

For $A^{(1)}$, the left row $i$ has exactly the two right components $j,k$,
each with matrix coefficient $1/\sqrt3$.  The weight primitive is diagonal
and hence orthogonal to those components.  The squared component magnitudes of
$B_{A^{(1)},\lambda_i}$ are $w_{ij},w_{ik}$, while the corresponding target
components are obtained by multiplication by
$m_{\lambda_i}(\lambda_j)$ and $m_{\lambda_i}(\lambda_k)$.  Orthogonal
projection of the latter two-vector onto the primitive line gives the weighted
variance identity
\[
 \|Y_i\|^2-\frac{|\langle B_i,Y_i\rangle|^2}{\|B_i\|^2}
 =\frac{w_{ij}w_{ik}}{w_{ij}+w_{ik}}|m_i(\lambda_j)-m_i(\lambda_k)|^2.
\]
The weights are positive.  Moreover,
\[
 \partial_y m_x(y)
 =-e^{\tau x}\int_0^a q e^{q(y-x)}\,\dd q<0,
\]
so the two multiplier values are unequal.  Summation proves strict positivity.
The numerical values are a direct high-precision evaluation of the exact
formula.
\end{proof}

\begin{assessment}[What the finite separator proves]
\label{ass:V8-four-state-meaning}
The stationary path and its Fisher source do not select the Gibbs Duhamel
branch.  Equal raw tangent size does not do so either.  The branch is decided
by the physical centered operator tangent and its right-energy incidence.
For $A^{(1)}$ the identity record connection makes the off-resonant tangent a
physical perturbation of edge conductances.  If that connection were not part
of the experiment, \Cref{thm:V8-connection-quotient} would allow its
off-resonant part to be moved by a different transport; that is precisely why
the connection must be frozen before the residual is read.
\end{assessment}

% =============================================================================
\section{Exact acquisition from a positive transfer logarithm}
\label{sec:V8-transfer-log}
% =============================================================================

Let $H_\theta$ be a $C^3$ family on finite Hilbert fibres and let
$U_\theta:\mathcal K_0\to\mathcal K_\theta$ be the declared physical
transport, with $U_0=I$.  Put
\begin{equation}
 \widetilde H_\theta=U_\theta^*H_\theta U_\theta,
 \qquad
 P_{\theta,\ell}=U_\theta^*e^{-\ell H_\theta}U_\theta
                =e^{-\ell\widetilde H_\theta},
 \qquad \ell>0.
 \label{eq:V8-transfer-family}
\end{equation}
Every $P_{\theta,\ell}$ is strictly positive and therefore has a unique
self-adjoint principal logarithm.

\begin{theorem}[Lag-independent transfer-log tangent]
\label{thm:V8-transfer-log}
For every $\varepsilon>0$ in the parameter interval and every $\ell>0$,
define
\begin{equation}
 V_{\varepsilon,\ell}
 :=-\frac{1}{2\varepsilon\ell}
   \left(\log P_{\varepsilon,\ell}
        -\log P_{-\varepsilon,\ell}\right).
 \label{eq:V8-transfer-log-secant}
\end{equation}
Then
\begin{equation}
 \boxed{
 V_{\varepsilon,\ell}
 =\frac{\widetilde H_\varepsilon-
        \widetilde H_{-\varepsilon}}{2\varepsilon},}
 \label{eq:V8-log-exact}
\end{equation}
so the result is independent of the positive lag.  If
\begin{equation}
 M_3(\varepsilon)
 :=\sup_{|\theta|\le\varepsilon}
   \|\widetilde H_\theta^{(3)}\|<\infty,
 \label{eq:V8-M3}
\end{equation}
then
\begin{equation}
 \boxed{
 \|V_{\varepsilon,\ell}-\dot{\widetilde H}_0\|
 \le\frac{\varepsilon^2}{6}M_3(\varepsilon).}
 \label{eq:V8-central-tangent-error}
\end{equation}
The constant $1/6$ is sharp for a scalar cubic family.
\end{theorem}

\begin{proof}
Functional calculus gives
$\log P_{\theta,\ell}=-\ell\widetilde H_\theta$, proving
\eqref{eq:V8-log-exact}.  Taylor's formula with integral remainder gives
\[
 \widetilde H_\varepsilon-
 \widetilde H_{-\varepsilon}
 =2\varepsilon\dot{\widetilde H}_0+R,
 \qquad
 \|R\|\le\frac{\varepsilon^3}{3}M_3(\varepsilon).
\]
Division by $2\varepsilon$ proves
\eqref{eq:V8-central-tangent-error}.  Equality of the scalar remainder
constant is attained by a pure cubic.
\end{proof}

The Duhamel source is invariant under adding a scalar to the Hamiltonian
tangent after the prescribed Gibbs centering.  Therefore the transfer may be
known only up to a positive scalar gain.  For $d=\dim\mathcal K_0$ define
\begin{equation}
 \operatorname{ctr}_d(X)=X-\frac{\Tr X}{d}I,
 \qquad
 \mathcal C_{\rho_0}(X)=X-\Tr(\rho_0X)I,
 \label{eq:V8-two-centerings}
\end{equation}
where $\rho_0=e^{-\beta H_0}/\Tr(e^{-\beta H_0})$ is the base Gibbs state.

\begin{corollary}[Gain-free centered tangent]
\label{cor:V8-gain-free-log}
Let $g_{\theta,\ell}>0$ be arbitrary scalar gains and suppose the acquired
positive matrices are
\begin{equation}
 \overline P_{\theta,\ell}=g_{\theta,\ell}P_{\theta,\ell}.
 \label{eq:V8-gained-transfer}
\end{equation}
Then
\begin{equation}
 \boxed{
 -\frac{1}{2\varepsilon\ell}
 \operatorname{ctr}_d\!
 \left(\log\overline P_{\varepsilon,\ell}
      -\log\overline P_{-\varepsilon,\ell}\right)
 =\operatorname{ctr}_d(V_{\varepsilon,\ell}).}
 \label{eq:V8-gain-free-secant}
\end{equation}
Moreover,
\begin{equation}
 \boxed{
 \mathcal C_{\rho_0}
 \bigl(\operatorname{ctr}_d(V_{\varepsilon,\ell})\bigr)
 =\mathcal C_{\rho_0}(V_{\varepsilon,\ell}).}
 \label{eq:V8-centering-compatible}
\end{equation}
Thus an unknown, even parameter-dependent, positive common gain does not
obstruct acquisition of the centered Gibbs source.  This differs from the
Yang--Mills partition-curvature target, where parameter-dependent absolute
weight is itself part of the desired Hessian.
\end{corollary}

\begin{proof}
Since a scalar commutes with every matrix,
\[
 \log(gP)=(\log g)I+\log P.
\]
The trace-centering map annihilates the scalar terms, and
\Cref{thm:V8-transfer-log} gives
\eqref{eq:V8-gain-free-secant}.  The two arguments of
$\mathcal C_{\rho_0}$ in \eqref{eq:V8-centering-compatible} differ by a
scalar multiple of the identity, which Gibbs centering removes.
\end{proof}

\begin{corollary}[Held-out lag-incidence test]
\label{cor:V8-two-lag-test}
For two positive lags $\ell_1,\ell_2$, a well-typed exact transfer card obeys
\begin{equation}
 \boxed{
 \operatorname{ctr}_d\!
 \left(-\ell_1^{-1}\log\overline P_{\theta,\ell_1}\right)
 =
 \operatorname{ctr}_d\!
 \left(-\ell_2^{-1}\log\overline P_{\theta,\ell_2}\right).}
 \label{eq:V8-two-lag-identity}
\end{equation}
Failure of this identity is a clock, transport, gain-type, or common-carrier
incidence defect; it is not a negative Duhamel residual.
\end{corollary}

\begin{proof}
Both sides equal $\operatorname{ctr}_d(\widetilde H_\theta)$ by functional
calculus and \Cref{cor:V8-gain-free-log}.
\end{proof}

% =============================================================================
\section{Outward matrix radii and the no-compressed-log obstruction}
\label{sec:V8-log-radii}
% =============================================================================

\begin{lemma}[Operator-Lipschitz logarithm on a positive floor]
\label{lem:V8-log-Lipschitz}
If $X,Y\succeq p_-I$ with $p_->0$, then
\begin{equation}
 \boxed{\|\log X-\log Y\|\le p_-^{-1}\|X-Y\|.}
 \label{eq:V8-log-Lipschitz}
\end{equation}
\end{lemma}

\begin{proof}
The positive-operator logarithm has the integral identity
\[
 \log X-\log Y
 =\int_0^\infty
   \bigl[(Y+tI)^{-1}-(X+tI)^{-1}\bigr]\,\dd t.
\]
The resolvent identity turns the integrand into
$(X+tI)^{-1}(X-Y)(Y+tI)^{-1}$.  Its norm is at most
$\|X-Y\|(p_-+t)^{-2}$.  Integration gives $p_-^{-1}\|X-Y\|$.
\end{proof}

It is useful to state the gain-free error in the quotient of self-adjoint
matrices by scalars.  Put
\begin{equation}
 \|[X]\|_{\rm sc}:=\inf_{c\in\R}\|X-cI\|.
 \label{eq:V8-scalar-quotient-norm}
\end{equation}

\begin{theorem}[Validated transfer-log tangent interval]
\label{thm:V8-log-error}
Let the exact gain-scaled transfers be
$\overline P_{\pm}=g_\pm P_{\pm\varepsilon,\ell}$, and let submitted positive
matrices $\widehat P_\pm$ satisfy
\begin{equation}
 \overline P_\pm\succeq p_-I,
 \qquad
 \widehat P_\pm\succeq p_-I,
 \qquad
 \|\widehat P_\pm-\overline P_\pm\|\le\delta_\pm.
 \label{eq:V8-transfer-input-radii}
\end{equation}
Define the submitted scalar-quotient secant from
\eqref{eq:V8-gain-free-secant}.  Then
\begin{equation}
 \boxed{
 \left\|
 \left[\widehat V_{\varepsilon,\ell}
       -\dot{\widetilde H}_0\right]
 \right\|_{\rm sc}
 \le
 \frac{\varepsilon^2}{6}M_3(\varepsilon)
 +\frac{\delta_++\delta_-}{2\varepsilon\ell p_-}.}
 \label{eq:V8-quotient-tangent-radius}
\end{equation}
For the Gibbs-centered representatives,
\begin{equation}
 \boxed{
 \|\mathcal C_{\rho_0}(\widehat V_{\varepsilon,\ell})
   -\mathcal C_{\rho_0}(\dot{\widetilde H}_0)\|
 \le
 \frac{\varepsilon^2}{3}M_3(\varepsilon)
 +\frac{\delta_++\delta_-}{\varepsilon\ell p_-}.}
 \label{eq:V8-centered-tangent-radius}
\end{equation}
\end{theorem}

\begin{proof}
By \Cref{lem:V8-log-Lipschitz}, the difference between the exact and submitted
logarithmic secants is at most
$(\delta_++\delta_-)/(2\varepsilon\ell p_-)$ before taking the scalar quotient.
The quotient norm cannot increase this error.  Add the central-difference
remainder from \Cref{thm:V8-transfer-log} to obtain
\eqref{eq:V8-quotient-tangent-radius}.  If $c$ is any scalar, then
\[
 \|\mathcal C_{\rho_0}(X)\|
 =\|X-cI-\Tr(\rho_0(X-cI))I\|
 \le2\|X-cI\|.
\]
Take the infimum over $c$ and apply this inequality to the quotient error,
which proves \eqref{eq:V8-centered-tangent-radius}.
\end{proof}

A smaller source panel is attractive computationally, but the logarithm does
not commute with compression unless the panel is already reducing.

\begin{countertheorem}[A compressed transfer does not determine the compressed generator]
\label{cth:V8-compressed-log}
Let
\begin{equation}
 P_0=\begin{pmatrix}1/2&0\\0&1/2\end{pmatrix},
 \qquad
 P_1=\begin{pmatrix}1/4&0\\0&3/4\end{pmatrix},
 \qquad
 J=\frac1{\sqrt2}\binom11.
 \label{eq:V8-compressed-log-example}
\end{equation}
Then
\begin{equation}
 \boxed{J^*P_0J=J^*P_1J=\frac12,}
 \label{eq:V8-same-compressed-transfer}
\end{equation}
but
\begin{equation}
 \boxed{
 J^*(-\log P_0)J=\log2,
 \qquad
 J^*(-\log P_1)J=\frac12\log\frac{16}{3}.}
 \label{eq:V8-different-compressed-generators}
\end{equation}
Thus $-\log(J^*PJ)$ cannot be substituted for $J^*(-\log P)J$ on an
unproved source head.  If $\Ran J$ reduces $P$, then functional calculus does
give
\begin{equation}
 J^*f(P)J=f(J^*PJ)
 \label{eq:V8-reducing-functional-calculus}
\end{equation}
for every continuous $f$ on the spectrum, and the compressed logarithm is
licensed.
\end{countertheorem}

\begin{proof}
The two scalar compressions are immediate.  Since both matrices are diagonal,
compression of their logarithms gives the arithmetic means of the two diagonal
logarithms, yielding \eqref{eq:V8-different-compressed-generators}; the two
numbers are unequal because $16/3\ne4$.  On a reducing subspace, $P$ is block
diagonal relative to $\Ran J\oplus(\Ran J)^\perp$, so the same is true for
$f(P)$ by functional calculus, proving
\eqref{eq:V8-reducing-functional-calculus}.
\end{proof}

\begin{firewall}[No logarithm after an ancestry fit]
\label{fire:V8-no-posthoc-log}
A source-cyclic or finite-word head may be used for the transfer-log card only
after its reducing property has been proved independently of the desired
tangent.  Choosing a head because its compressed logarithm produces the
preferred residual is a fit.  When reduction is unavailable, the full finite
transfer or a larger already certified reducing carrier is the required
input.
\end{firewall}

% =============================================================================
\section{Propagation of the tangent radius into the existing Duhamel compiler}
\label{sec:V8-Duhamel-debit}
% =============================================================================

Return to the base Hamiltonian $H=H_0$, put $a=\beta/2$, and let
$V:E_Q\to\mathsf{HS}(\mathcal K)$ be the physically centered tangent synthesis.
For any approximation $\widehat V$ on the same command space define
\begin{equation}
 \begin{aligned}
 B_V&=Z_\beta^{-1/2}e^{-\tau H_{\rm L}}V e^{-aH_{\rm R}},\\
 Y_V&=-Z_\beta^{-1/2}\int_0^a
       e^{-(a-s)H_{\rm L}}V e^{-sH_{\rm R}}\,\dd s,
 \end{aligned}
 \label{eq:V8-linear-BY}
\end{equation}
and similarly $B_{\widehat V},Y_{\widehat V}$.  Here the exponentials act on
the Hilbert--Schmidt carrier and are contractions.

\begin{theorem}[Tangent-acquisition debit for the static ancestry packet]
\label{thm:V8-tangent-to-Schur}
Assume
\begin{equation}
 \|(V-\widehat V)\mathsf N^{\dagger/2}\|
 \le\eta_V
 \label{eq:V8-tangent-input-error}
\end{equation}
in the target-normalized operator norm.  Then
\begin{equation}
 \boxed{
 \|(Y_V-Y_{\widehat V})\mathsf N^{\dagger/2}\|
 \le aZ_\beta^{-1/2}\eta_V.}
 \label{eq:V8-Y-error}
\end{equation}
If the same coefficient norm is used for the action-source column, then also
\begin{equation}
 \boxed{\|B_V-B_{\widehat V}\|
 \le Z_\beta^{-1/2}\|V-\widehat V\|.}
 \label{eq:V8-B-error}
\end{equation}

Let $B=(B_{\rm fix},B_V)$ and
$\widehat B=(B_{\rm fix},B_{\widehat V})$ be the complete primitive banks,
with every weight, line, and independently occurring prior column held fixed.
Put
\[
 S=B^*B,\quad T=B^*Y_V,\quad U=Y_V^*Y_V,
\]
and use hats for the submitted packet.  With
\begin{equation}
 \delta_B=\|\widehat B-B\|,
 \qquad
 \delta_Y=\|(Y_{\widehat V}-Y_V)\mathsf N^{\dagger/2}\|,
 \label{eq:V8-BY-radii}
\end{equation}
one has
\begin{align}
 \varepsilon_S&:=\|\widehat S-S\|
 \le(\|B\|+\|\widehat B\|)\delta_B,
 \label{eq:V8-S-radius}\\
 \varepsilon_T&:=\|(\widehat T-T)\mathsf N^{\dagger/2}\|
 \le\delta_B\|Y_V\mathsf N^{\dagger/2}\|
     +\|\widehat B\|\delta_Y,
 \label{eq:V8-T-radius}\\
 \varepsilon_U&:=\|\mathsf N^{\dagger/2}
            (\widehat U-U)\mathsf N^{\dagger/2}\|
 \le
 (\|Y_V\mathsf N^{\dagger/2}\|
  +\|Y_{\widehat V}\mathsf N^{\dagger/2}\|)\delta_Y.
 \label{eq:V8-U-radius}
\end{align}
If $S$ and $\widehat S$ have one verified common support $P$ and
$S,\widehat S\succeq s_*P$, then the exact and submitted full-primitive
residuals
\[
 R=U-T^*S^\dagger T,
 \qquad
 \widehat R=\widehat U-\widehat T^*\widehat S^\dagger\widehat T
\]
satisfy
\begin{equation}
 \boxed{
 \|\mathsf N^{\dagger/2}(\widehat R-R)
       \mathsf N^{\dagger/2}\|
 \le
 \varepsilon_U+s_*^{-1}(t+\widehat t)\varepsilon_T
 +s_*^{-2}t\widehat t\,\varepsilon_S,}
 \label{eq:V8-tangent-Schur-radius}
\end{equation}
where
$t=\|T\mathsf N^{\dagger/2}\|$ and
$\widehat t=\|\widehat T\mathsf N^{\dagger/2}\|$.
\end{theorem}

\begin{proof}
The contraction bounds in \eqref{eq:V8-linear-BY} give
\eqref{eq:V8-Y-error}--\eqref{eq:V8-B-error}; the integral has length $a$.
The three Gram estimates follow by the identities
\[
 \widehat B^*\widehat B-B^*B
 =(\widehat B-B)^*\widehat B+B^*(\widehat B-B),
\]
\[
 \widehat B^*Y_{\widehat V}-B^*Y_V
 =(\widehat B-B)^*Y_V
  +\widehat B^*(Y_{\widehat V}-Y_V),
\]
and the analogous two-term expansion of
$Y_{\widehat V}^*Y_{\widehat V}-Y_V^*Y_V$.
Equation~\eqref{eq:V8-tangent-Schur-radius} is then exactly the supported
Schur perturbation estimate of \Cref{prop:V5-Schur-perturbation}.
\end{proof}

\begin{corollary}[Transfer-log debit in the finite V5--V6 certificate]
\label{cor:V8-transfer-to-compiler}
Insert the centered tangent radius from
\eqref{eq:V8-centered-tangent-radius} into
\eqref{eq:V8-Y-error}--\eqref{eq:V8-tangent-Schur-radius}.  Add the result to
the already separate V5 quadrature, left-energy tail, measured-matrix, and
route radii, and to the V6 weighted-absorption radii when that target is used.
A protected-screen lower edge which remains positive after this sum is a
robust finite Duhamel birth.  A small upper edge is an approximate follower.
An exact follower still requires the exact reducing/range identity or an
exhausting zero upper sequence.

If the common source support or its floor is not certified, the packet stops
at a support transition; the perturbation bound is not evaluated with a
counterfeit pseudoinverse.
\end{corollary}

\begin{proof}
All listed radii bound distinct arrows and are added by the triangle
inequality.  The branch semantics and the support qualification are those of
\Cref{thm:V5-compiler}; no new residual calculus is required.
\end{proof}

% =============================================================================
\section{The target-minimal physical acquisition and V8 endpoint}
\label{sec:V8-endpoint}
% =============================================================================

\begin{acquisition}[Common-action centered-tangent transfer card]
\label{acq:V8-transfer-card}
On one predeclared finite active gravity--gauge--Higgs--fermion card, freeze
before reading any residual:
\begin{enumerate}[label=\textup{(A\arabic*)},leftmargin=3.4em]
\item the physical common-action parameter $\theta$, the three values
      $0,\pm\varepsilon$, and the same field-shell, line, divisor, and support
      convention at all three values;
\item the physical Hilbert-bundle transport $U_\theta$ and the base Gibbs
      centering state $\rho_0$;
\item the positive transfer matrices
      $\overline P_{\pm\varepsilon,\ell}$ on the full finite carrier or on an
      independently proved reducing carrier, with a positive spectral floor
      and outward matrix radii;
\item a second predeclared lag for the held-out identity
      \eqref{eq:V8-two-lag-identity};
\item a bound for $M_3(\varepsilon)$, or any sharper valid central-secant
      remainder;
\item the complete fixed primitive bank, including weight and line columns,
      and the source and target metrics used by the V5--V6 compiler.
\end{enumerate}
The output is the Gibbs-centered tangent and the certified radius in
\eqref{eq:V8-centered-tangent-radius}, followed by the already established
Duhamel packet.  No current, stress, charge, or Einstein row is included in
this acquisition.
\end{acquisition}

\begin{theorem}[Version V8 upstream tangent and acquisition theorem]
\label{thm:V8-integrated}
Preserving every theorem and qualification of Versions~V1--V7, Version~V8
proves the following.
\begin{enumerate}[label=\textup{(V8.\arabic*)},leftmargin=5.0em]
\item A transported ground-state derivative fixes only the vacuum row
      $r=\dot Hs$; the centered Hamiltonian tangent $A=Q\dot H Q$ is an
      orthogonal, independent source block.
\item Every centered self-adjoint block occurs as the tangent of a strictly
      positive reversible finite-state Markov family with unchanged stationary
      law.  The V7 stationary Fisher--Green card therefore cannot determine
      the Gibbs Duhamel source.
\item The physical connection acts by $A\mapsto A+[h,K]$.  The resonant block
      $\sum_\lambda E_\lambda AE_\lambda$ is the exact quotient invariant,
      and every off-resonant block is a connection commutator.
\item The four-state card
      \eqref{eq:V8-Hadamard-basis}--\eqref{eq:V8-two-tangents} has identical
      stationary paths and equal tangent Hilbert--Schmidt norms, but realizes
      an exact follower and a strict positive Duhamel residual, respectively.
\item One positive transfer matrix at each of $\pm\varepsilon$ reconstructs
      the pulled Hamiltonian secant by a principal logarithm.  The result is
      independent of lag and has the sharp tangent error
      $\varepsilon^2M_3/6$.
\item Unknown positive scalar gains cancel in the scalar quotient and therefore
      do not obstruct the Gibbs-centered source.  A second lag is an exact
      held-out clock/transport incidence test.
\item A positive transfer floor gives the explicit logarithmic data radius
      \eqref{eq:V8-quotient-tangent-radius}.  The centered radius propagates
      through the thermal source maps and the V5 supported Schur complement by
      \eqref{eq:V8-tangent-Schur-radius}.
\item A compressed transfer cannot be logarithmically promoted unless its
      source carrier is already reducing.  Equal compressed lag data can have
      different compressed generators.
\end{enumerate}
No item identifies the active chiral common-action transfer of the original
ESM regulator, or supplies its physical connection, line incidence, cofinal
transport, represented current/stress, or quantum Einstein residual.
\end{theorem}

\begin{proof}
Items~\textup{(V8.1)--(V8.2)} are
\Cref{thm:V8-ground-decomposition,thm:V8-conductance-realization,cth:V8-stationary-no-centered-tangent}.
Item~\textup{(V8.3)} is \Cref{thm:V8-connection-quotient}.  Item
\textup{(V8.4)} is \Cref{cth:V8-four-state-branch}.  Items
\textup{(V8.5)--(V8.7)} are
\Cref{thm:V8-transfer-log,cor:V8-gain-free-log,cor:V8-two-lag-test,thm:V8-log-error,thm:V8-tangent-to-Schur}.
The last item is \Cref{cth:V8-compressed-log}.  The limitations are the input
rows of \Cref{acq:V8-transfer-card} which have not been populated on the
original active ESM carrier.
\end{proof}

\begin{assessment}[Status after V8]
\label{ass:V8-status}
\passstatus{} for the vacuum-row/centered-block Pythagoras, reversible
conductance realization, connection quotient, exact four-state branch
separator, gain-free transfer-log reconstruction, positive-floor error bound,
held-out two-lag identity, no-compressed-log theorem, and tangent-to-Schur
debit.

\obsstatus{} for inferring the centered Hamiltonian tangent from stationary
probability/Fisher data, for interpreting an off-resonant derivative without a
physical connection, and for taking a matrix logarithm after an unproved
source compression.

\stopstatus{} on the original ESM continuum at the first unpopulated
same-history common-action transfer and connection card.  V5 and V6 remain the
correct downstream compilers once that card lands.  Further generic lag,
low-island, or concentration machinery would be circular at this stage.
\end{assessment}

\begin{openproblem}[Version V9: populate one physical transfer-log tangent card]
\label{op:V8-next}
Preserve Versions~V1--V8 verbatim and proceed only in the following order.
\begin{enumerate}[label=\textup{(V9.\arabic*)},leftmargin=5.0em]
\item Freeze one actual common-action deformation of the active reflected
      Dirac--Yukawa field shell and its physical Hilbert-bundle connection.
\item Acquire the transported positive transfer at $\theta=\pm\varepsilon$
      and one physical lag on the full finite carrier or an independently
      reducing source-cyclic carrier.  Scalar gain may remain unknown.
\item Acquire a second lag as a held-out incidence check and a valid central
      secant remainder.  Reject, rather than reinterpret, a failed common-
      carrier or logarithmic positivity test.
\item Form the Gibbs-centered tangent, then the complete thermally dressed
      weight/action/line bank and Duhamel target.  Only then evaluate the V5
      mixed-clock matrices and the V6 weighted absorption panels.
\item On a separated finite verdict, acquire adjacent/direct transfer cards
      before any current/stress promotion.  If the first missing row is the
      reflected hopping/line packet rather than the tangent, stop and export
      that earlier gate.
\end{enumerate}
The admissible outputs are
\[
 \boxed{
 \begin{gathered}
 \text{robust finite ancestry branch}
 \quad\big|\quad
 \text{connection, transfer, support, or line obstruction}\\[-0.15em]
 \big|\\[-0.45em]
 \text{STOP at the first earlier terminal-order gate}
 \end{gathered}}
\]
No replacement toy mass, fitted connection, post-hoc reducing head, or new
abstract source space is admissible.
\end{openproblem}

\section*{V8 verification and append-only record}
\addcontentsline{toc}{section}{V8 verification and append-only record}
The exact finite replay accompanying this continuation checks the four-state
Markov Hamiltonian, both conductance tangents, equality of their raw
Hilbert--Schmidt norms, the positive Duhamel formula, the connection-commutator
decomposition, the compressed-log counterexample, and exact scalar-gain
cancellation in a centered logarithmic secant.  Transcendental residual values
are evaluated at one hundred decimal digits and are diagnostic; every sign
claim used by the theorem has an analytic proof above.

The complete Version~V7 source preceding its sole final document-closing
command is preserved byte for byte.  Its SHA--256 digest is
\begin{center}\small\ttfamily
03162cbd39b9f01f2c9d02ec0f708815d46366ebbe178eb0878b731e5cb51af2.
\end{center}
A later continuation must remove only the sole final
\verb|\end{document}|, append new material immediately before it, restore one
terminal command, and use the filename ending in \texttt{\_V9.tex}.

% =============================================================================
% END APPENDED VERSION V8 MATERIAL
% =============================================================================


% =============================================================================
% BEGIN APPENDED VERSION V9 MATERIAL
% =============================================================================
\clearpage
\hypersetup{
 pdftitle={Renewal Einstein--Standard--Model Physical Source Metric, Duhamel Ancestry and Quantum Continuum Programme V9},
 pdfsubject={The physical time-split occurrence gate and a target-minimal crossing-action germ}
}
\section*{Version V9: the physical time-split occurrence gate and a target-minimal crossing-action germ}
\addcontentsline{toc}{section}{Version V9: physical time split and crossing-action germ}

\begin{abstract}
Version~V8 reconstructed a centered Hamiltonian tangent from a physically
transported positive transfer and proved that stationary Fisher data cannot
supply the missing centered block.  The terminal quantum order in the attached
Grand-Tensor and sector-patched line programmes places one stricter gate before
that tangent is used: on one fixed bosonic history the complete chiral
Dirac--Yukawa precision must first carry an occurring positive/negative-time
split, a physical reflection, the Berezin-order sign, and the completely
assembled crossing-hopping/Yukawa block.  The already established
relation-descended hopping theorem then tests whether that block factors through
the one-letter word relations and lies in the Hermitian positive cone.
Version~V9 does not reprove that theorem.  It constructs the minimum direct
acquisition which supplies its missing input.

The first result is an exact nonidentifiability theorem.  One fixed strictly
positive precision, its complete positive transfer at every lag, its
determinant modulus, and even its two diagonal time-sector compressions admit
two valid orthogonal time splits whose reflection-adapted crossing coefficients
are respectively $12/25$ and $-12/25$.  Thus a perfect transfer-log card cannot
select the hopping branch before the physical time split and reflection have
occurred.  Reversing the Berezin-order sign changes the branch again while
leaving every positive modulus unchanged.

The missing crossing row can nevertheless be acquired without reconstructing a
larger state, covariance, or path law.  If $q(z)=\langle z,\mathcal Kz\rangle$
is the deterministic complex quadratic Read of the completely assembled finite
precision, the arbitrary sesquilinear polarization identity
\[
 \langle x,\mathcal Ky\rangle
 =\frac14\sum_{r=0}^{3}i^{-r}q(x+i^ry)
\]
recovers every oriented positive/negative-time matrix element, without assuming
that $\mathcal K$ is Hermitian.  Applied only to the occurring one-letter bank,
this gives the raw reflected crossing matrix required by the inherited
relation and hopping compiler.  Scalar energy shifts and arbitrary block-
diagonal same-time terms cancel exactly; they are not charged merely to read
the crossing target.  One held-out mixed superposition tests same-action
assembly.

Outward query radii propagate first to the raw crossing matrix, then through the
supported Moore--Penrose word quotient, and finally to a one-Lipschitz interval
for the distance from the reflected-hopping cone.  A nonzero lower distance is
an explicit obstruction; an exact reflection-Hermiticity identity together
with a positive eigenvalue floor is a pass; every remaining interval is a
stop.  An adjacent-cutoff rectangle separates crossing-block transport from
reflection transport and treats a Berezin-sign change as a $\mathbb Z_2$ line
obstruction rather than a small metric defect.

An exact rational replay verifies the split counterexample, all four-phase
polarization identities for a non-Hermitian precision, a nontrivial three-word
relation quotient, the source-minimal represented crossing, its positive
factor, and a relation-breaking mutation.  The active Einstein--Standard--Model
regulator still does not export the physical time split, reflection/Berezin
sign, or same-history crossing-action values on one named bosonic card.
Version~V9 therefore moves the stopping boundary upstream of V8 rather than
adding another Duhamel, lag, or concentration compiler.
\end{abstract}

% =============================================================================
\section{Why the crossing gate precedes the V8 transfer tangent}
\label{sec:V9-direction}
% =============================================================================

The cumulative programme now contains two mathematically valid but differently
ordered finite compilers.
\begin{enumerate}[label=\textup{(O\arabic*)},leftmargin=3.4em]
\item Versions~V5--V8 reconstruct, approximate, and transport a thermally
      dressed Duhamel target after a physical Hamiltonian and its differentiated
      common-action source have been identified.
\item The attached terminal fermion theorem first asks whether the completely
      assembled time-split quadratic precision has a relation-descended,
      reflection-positive one-sided hopping factorization.  This requires the
      finite tuple
      \begin{equation}
       \boxed{
       (\mathcal H_+,\mathcal H_-,R_\Theta,
        \varepsilon_\Theta,X_\times,W_1)}
       \label{eq:V9-crossing-tuple}
      \end{equation}
      on one fixed bosonic history.
\end{enumerate}
The second row is logically earlier.  A later transfer generator, source metric,
or Duhamel residual cannot create a time split, a reflection convention, or an
absent crossing coefficient of the fermion precision.

\begin{proposition}[Imported crossing compiler and the V9 nonduplication rule]
\label{prop:V9-imported-crossing}
For the tuple in \eqref{eq:V9-crossing-tuple}, put
\begin{equation}
 K_\times=-\varepsilon_\Theta X_\times R_\Theta
 \label{eq:V9-Kcross}
\end{equation}
on the positive-time carrier, let $W_1:\mathcal F_1\to\mathcal H_+$ be the
occurring one-letter synthesis, and put $P_W=W_1^\dagger W_1$.  The attached
terminal theorem already proves the following.
\begin{enumerate}[label=\textup{(C\arabic*)},leftmargin=3.4em]
\item The raw coefficient form $K_F=W_1^*K_\times W_1$ descends through the
      one-letter relations exactly when
      \begin{equation}
       K_F=P_WK_FP_W.
       \label{eq:V9-imported-relation}
      \end{equation}
\item On that branch the minimum-norm represented operator is
      \begin{equation}
       K_{\rm rep}=(W_1^\dagger)^*K_FW_1^\dagger.
       \label{eq:V9-imported-represented}
      \end{equation}
\item Writing $H=(K_{\rm rep}+K_{\rm rep}^*)/2$ and
      $A=(K_{\rm rep}-K_{\rm rep}^*)/2$, the exact squared distance to the
      Hermitian positive cone is
      \begin{equation}
       \dist_{\rm HS}(K_{\rm rep},\operatorname{Herm}_+)^2
       =\norm{A}_{\rm HS}^2+\norm{H_-}_{\rm HS}^2.
       \label{eq:V9-imported-cone}
      \end{equation}
      Its zero branch is the source-minimal one-sided factorization
      $K_{\rm rep}=R_\times^*R_\times$, of minimum source rank
      $\rank K_{\rm rep}$.
\end{enumerate}
Version~V9 treats these statements as an imported theorem layer.  Its new task
is only to construct and certify $K_F$ from the first physically occurring
quadratic action row.
\end{proposition}

\begin{proof}
This proposition records the already established relation-descended reflected-
hopping theorem in the notation needed below.  No part of its proof is repeated
in this continuation.
\end{proof}

\begin{assessment}[Correction of the V8 open-problem order]
\label{ass:V9-V8-order}
The transfer-log reconstruction of Version~V8 remains valid for a physically
identified positive Hamiltonian transfer.  It is now classified as a downstream
or independently held-out common-action row.  It cannot be used to infer the
fermionic tuple \eqref{eq:V9-crossing-tuple}.  The first active quantum datum is
the time-split crossing action itself.
\end{assessment}

% =============================================================================
\section{The physical split packet and its exact domain audit}
\label{sec:V9-split}
% =============================================================================

Let $\mathcal H$ be one finite one-particle coefficient carrier.  A candidate
linear time reflection is represented, after the anti-linear conjugation and
Berezin transpose convention have been fixed, by a partial isometry on
$\mathcal H$.

\begin{definition}[Finite physical time-split packet]
\label{def:V9-split-packet}
A candidate split packet is
\begin{equation}
 \mathfrak T_\Theta=(P_+,P_-,R_\Theta,\varepsilon_\Theta),
 \qquad \varepsilon_\Theta\in\{+1,-1\},
 \label{eq:V9-split-packet}
\end{equation}
where $P_\pm\in\operatorname{End}(\mathcal H)$ and
$R_\Theta\in\operatorname{End}(\mathcal H)$.  Define the nonnegative domain
residual
\begin{equation}
\boxed{
\begin{aligned}
 \Delta_{\rm split}:={}&
 \sum_{\sigma\in\{+,-\}}
 \left(\norm{P_\sigma-P_\sigma^*}_{\rm HS}^2
      +\norm{P_\sigma^2-P_\sigma}_{\rm HS}^2\right)\\
 &+\norm{P_+P_-}_{\rm HS}^2
  +\norm{P_++P_--I}_{\rm HS}^2\\
 &+\norm{R_\Theta-P_-R_\Theta P_+}_{\rm HS}^2
  +\norm{R_\Theta^*R_\Theta-P_+}_{\rm HS}^2
  +\norm{R_\Theta R_\Theta^*-P_-}_{\rm HS}^2.
\end{aligned}}
\label{eq:V9-split-residual}
\end{equation}
This is an input-domain residual.  Its failure rejects the proposed split; it
is not a negative hopping mode.
\end{definition}

\begin{theorem}[Exact split and reflection criterion]
\label{thm:V9-split-criterion}
One has $\Delta_{\rm split}=0$ if and only if
\begin{equation}
 \mathcal H=\mathcal H_+\oplus\mathcal H_-,
 \qquad \mathcal H_\pm=\Ran P_\pm,
 \label{eq:V9-orthogonal-split}
\end{equation}
are orthogonal complementary subspaces and
\begin{equation}
 R_\Theta:\mathcal H_+\longrightarrow\mathcal H_-
 \label{eq:V9-unitary-reflection}
\end{equation}
is unitary, extended by zero on $\mathcal H_-$.  In particular the two sector
dimensions agree.
\end{theorem}

\begin{proof}
Every summand in \eqref{eq:V9-split-residual} is nonnegative.  If the sum
vanishes, $P_\pm$ are self-adjoint idempotents, hence orthogonal projections;
$P_+P_-=0$ and $P_++P_-=I$ give the orthogonal direct sum.  The support identity
puts $R_\Theta$ from $\mathcal H_+$ into $\mathcal H_-$.  The remaining two
identities say that it is isometric on $\mathcal H_+$ and onto
$\mathcal H_-$.  Conversely, an orthogonal complementary split and a unitary
map between the two summands satisfy every displayed identity, so the residual
vanishes.
\end{proof}

\begin{firewall}[The split is an occurrence row]
\label{fire:V9-split-occurrence}
A spectral decomposition, a convenient checkerboard, or a basis chosen after
viewing the crossing sign is not a physical time split.  The projections,
reflection, boundary convention, and Berezin ordering must come from the
predeclared reflected slab.  The criterion above verifies a submitted packet;
it does not manufacture one from the unsplit precision.
\end{firewall}

% =============================================================================
\section{A fixed positive precision admits opposite physical crossing branches}
\label{sec:V9-split-counterexample}
% =============================================================================

The next finite example proves that the V8 transfer-log result cannot be moved
in front of the split packet.

\begin{countertheorem}[Unsplit precision and diagonal time marginals do not determine the crossing sign]
\label{cth:V9-unsplit-no-crossing}
On $\R^2$ let
\begin{equation}
 \mathcal K_0=\begin{pmatrix}1&0\\0&2\end{pmatrix}.
 \label{eq:V9-K0}
\end{equation}
For $\sigma\in\{+1,-1\}$ put
\begin{equation}
 u_\sigma=\frac15\binom{4}{3\sigma},
 \qquad
 v_\sigma=\frac15\binom{-3\sigma}{4},
 \label{eq:V9-uv}
\end{equation}
let $P_{+,\sigma}=u_\sigma u_\sigma^{\mathsf T}$,
$P_{-,\sigma}=v_\sigma v_\sigma^{\mathsf T}$, and let
$R_{\Theta,\sigma}=v_\sigma u_\sigma^{\mathsf T}$.  Both candidate packets
pass \Cref{thm:V9-split-criterion}.  They have identical same-time precision
compressions,
\begin{equation}
 \boxed{
 u_\sigma^{\mathsf T}\mathcal K_0u_\sigma=\frac{34}{25},
 \qquad
 v_\sigma^{\mathsf T}\mathcal K_0v_\sigma=\frac{41}{25},}
 \label{eq:V9-same-diagonal-marginals}
\end{equation}
but opposite mixed coefficients,
\begin{equation}
 \boxed{
 u_\sigma^{\mathsf T}\mathcal K_0v_\sigma
 =\sigma\frac{12}{25}.}
 \label{eq:V9-opposite-crossing}
\end{equation}
With the same Berezin convention $\varepsilon_\Theta=-1$, the represented
one-dimensional matrices are therefore
\begin{equation}
 K_{\times,\sigma}=\sigma\frac{12}{25}.
 \label{eq:V9-two-crossing-branches}
\end{equation}
The $\sigma=+1$ packet has a strict one-sided hopping factorization, while the
$\sigma=-1$ packet has an explicit negative vector.

Nevertheless both packets have the same complete unsplit precision, the same
positive transfer $e^{-\ell\mathcal K_0}$ for every $\ell>0$, the same
precision determinant $2$, and the same two diagonal time-sector Reads in
\eqref{eq:V9-same-diagonal-marginals}.  Hence none of those data determines the
crossing branch without the physical split and mixed action row.
\end{countertheorem}

\begin{proof}
The two vectors in \eqref{eq:V9-uv} form an orthonormal basis.  The projector
and partial-isometry identities are immediate.  Direct multiplication gives
\[
 \frac{16+18}{25}=\frac{34}{25},
 \qquad
 \frac{9+32}{25}=\frac{41}{25},
\]
and
\[
 \frac{-12\sigma+24\sigma}{25}
 =\sigma\frac{12}{25}.
\]
For $\varepsilon_\Theta=-1$, equation~\eqref{eq:V9-Kcross} leaves this scalar
unchanged.  A positive scalar is its own positive square; a negative scalar is
outside the Hermitian positive cone.  All remaining listed objects depend only
on the fixed matrix $\mathcal K_0$ or on the squared split coordinates, so they
are identical for the two signs.
\end{proof}

\begin{corollary}[A positive determinant modulus does not fix the Berezin-order sign]
\label{cor:V9-line-sign-independent}
Fix the $\sigma=+1$ split.  Replacing $\varepsilon_\Theta=-1$ by
$\varepsilon_\Theta=+1$ changes $K_\times=12/25$ to $K_\times=-12/25$ while
leaving $\mathcal K_0$, its determinant modulus, every singular value, and both
sector compressions unchanged.  Thus the Berezin-order/reflection sign is an
independent orientation row.  The later real Pfaffian sign and determinant-line
holonomy remain separate still.
\end{corollary}

\begin{proof}
Equation~\eqref{eq:V9-Kcross} is linear in
$-\varepsilon_\Theta$.  Positive moduli are insensitive to this sign.
\end{proof}

% =============================================================================
\section{The arbitrary sesquilinear polarization identity}
\label{sec:V9-polarization}
% =============================================================================

The complete precision matrix is a deterministic coefficient of the assembled
quadratic action.  It need not be Hermitian.  The following polarization is
therefore stated for an arbitrary sesquilinear form rather than for a quadratic
energy.

\begin{lemma}[Four-phase polarization without Hermiticity]
\label{lem:V9-four-phase}
Let $b:\mathcal H\times\mathcal H\to\C$ be conjugate linear in the first
variable and linear in the second, with no symmetry assumption, and put
$q(z)=b(z,z)$.  Then for all $x,y\in\mathcal H$,
\begin{equation}
 \boxed{
 b(x,y)=\frac14\sum_{r=0}^{3}i^{-r}q(x+i^ry).}
 \label{eq:V9-four-phase}
\end{equation}
Equivalently, if the diagonal values $q(x)$ and $q(y)$ have already been read,
then
\begin{equation}
 \boxed{
 \begin{aligned}
 b(x,y)=\frac12\bigl[&q(x+y)-q(x)-q(y)\\
 &-i\{q(x+iy)-q(x)-q(y)\}\bigr].
 \end{aligned}}
 \label{eq:V9-two-mixed-probes}
\end{equation}
\end{lemma}

\begin{proof}
For $\alpha=i^r$,
\[
 q(x+\alpha y)
 =q(x)+\alpha b(x,y)+\overline\alpha b(y,x)+q(y).
\]
Multiply by $\alpha^{-1}=\overline\alpha$ and sum over the four fourth roots
of unity.  The sums of $\overline\alpha$ and $\overline\alpha^2$ vanish,
whereas $\sum\overline\alpha\alpha=4$, proving
\eqref{eq:V9-four-phase}.  Expanding the two differences in
\eqref{eq:V9-two-mixed-probes} gives respectively
$b(x,y)+b(y,x)$ and $i b(x,y)-i b(y,x)$; the displayed linear combination
isolates $b(x,y)$.
\end{proof}

\begin{remark}[This is a coefficient Read, not a fermion configuration]
\label{rem:V9-commuting-probe}
The vectors in \Cref{lem:V9-four-phase} are commuting test spinors used to read
the finite coefficient matrix of an already occurring quadratic Grassmann
action.  No probability law for commuting fermions is introduced.  If the
complete precision matrix is submitted directly, the polarization is merely an
independent reconstruction check.
\end{remark}

% =============================================================================
\section{Target-compressed acquisition of the raw reflected crossing matrix}
\label{sec:V9-crossing-acquisition}
% =============================================================================

Assume the split packet passes.  Let $\mathcal K$ be the completely assembled
finite precision on $\mathcal H_-\oplus\mathcal H_+$, with every gauge,
Higgs/Yukawa, boundary, regulator, ghost, and admitted counterterm contribution
included before any short or positive part is taken.  Put
\begin{equation}
 X_\times=P_+\mathcal KP_-:\mathcal H_-\longrightarrow\mathcal H_+,
 \qquad
 q_{\mathcal K}(z)=\ip{z}{\mathcal Kz}.
 \label{eq:V9-crossing-from-precision}
\end{equation}
Let $W_1:\mathcal F_1\to\mathcal H_+$ have coefficient basis
$e_1,\ldots,e_d$ and columns $\xi_j=W_1e_j$.

\begin{theorem}[Four-phase crossing-germ card]
\label{thm:V9-crossing-germ}
For $1\le i,j\le d$ define the four same-action scalar Reads
\begin{equation}
 Q_{ji}^{(r)}
 :=q_{\mathcal K}(\xi_j+i^rR_\Theta\xi_i),
 \qquad r=0,1,2,3.
 \label{eq:V9-crossing-queries}
\end{equation}
Then the complete raw reflected crossing matrix is
\begin{equation}
 \boxed{
 (K_F)_{ji}
 =-\frac{\varepsilon_\Theta}{4}
   \sum_{r=0}^{3}i^{-r}Q_{ji}^{(r)},
 \qquad
 K_F=W_1^*K_\times W_1.}
 \label{eq:V9-KF-from-queries}
\end{equation}
This acquisition uses only the one-letter target bank.  It does not require the
same-time diagonal blocks of $\mathcal K$, its inverse, a fermion covariance,
a transfer Hamiltonian, a determinant value, or a complete reflected word
table.
\end{theorem}

\begin{proof}
The positive- and negative-time subspaces are orthogonal.  With
$x=\xi_j\in\mathcal H_+$ and $y=R_\Theta\xi_i\in\mathcal H_-$,
\[
 \ip{x}{\mathcal Ky}
 =\ip{\xi_j}{X_\times R_\Theta\xi_i}
 =-\varepsilon_\Theta\ip{\xi_j}{K_\times\xi_i}.
\]
Apply \Cref{lem:V9-four-phase} to the sesquilinear form
$b(x,y)=\ip{x}{\mathcal Ky}$ and rearrange the sign.  The last matrix identity
is the definition of coefficient compression by $W_1$.
\end{proof}

\begin{corollary}[Exact insensitivity to scalar and same-time action gauges]
\label{cor:V9-block-diagonal-invariance}
Let $D=P_-DP_-+P_+DP_+$ be any block-diagonal operator and $c\in\C$.  Replacing
\begin{equation}
 \mathcal K\longmapsto\mathcal K+D+cI
 \label{eq:V9-block-gauge}
\end{equation}
leaves $X_\times$, $K_\times$, every entry in
\eqref{eq:V9-KF-from-queries}, and the inherited hopping verdict unchanged.
Thus an unknown scalar energy origin and unrequested same-time precision blocks
are not charged to the crossing target.
\end{corollary}

\begin{proof}
Orthogonality gives $P_+(D+cI)P_-=0$.  Equivalently,
$\ip{x}{(D+cI)y}=0$ for the mixed vectors used in
\Cref{thm:V9-crossing-germ}, so the four-phase sum cancels their diagonal
contributions exactly.
\end{proof}

\begin{corollary}[One held-out mixed action checks assembly]
\label{cor:V9-heldout-crossing}
Let $a,b\in\mathcal F_1$ be fixed before the entries of $K_F$ are populated,
and put $x=W_1a$, $y=R_\Theta W_1b$.  Define
\begin{equation}
 \mathcal H_{a,b}^{\times}
 :=-\frac{\varepsilon_\Theta}{4}
   \sum_{r=0}^{3}i^{-r}q_{\mathcal K}(x+i^ry)
   -a^*K_Fb.
 \label{eq:V9-heldout-residual}
\end{equation}
On one common assembled action card,
\begin{equation}
 \boxed{\mathcal H_{a,b}^{\times}=0.}
 \label{eq:V9-heldout-zero}
\end{equation}
A nonzero held-out value rejects the proposed common-action assembly or its
routing to the one-letter bank.  It is not repaired by projecting $K_F$ onto
the positive cone.
\end{corollary}

\begin{proof}
The four-phase term is
$\ip{W_1a}{K_\times W_1b}=a^*K_Fb$ by
\Cref{thm:V9-crossing-germ} and linearity.  The qualification follows because
the positive-cone projection would alter the action rather than explain a
failure of the defining equality.
\end{proof}

\begin{countertheorem}[The mixed crossing row is independent of both time marginals]
\label{cth:V9-marginals-no-mixed}
Even complete knowledge of the two restricted quadratic forms
$q_{\mathcal K}|_{\mathcal H_+}$ and
$q_{\mathcal K}|_{\mathcal H_-}$ does not determine the Reads in
\eqref{eq:V9-crossing-queries}.  The two valid splits in
\Cref{cth:V9-unsplit-no-crossing} have identical restricted scalar forms and
opposite mixed crossing coefficients.  Hence the mixed action row is an
independent same-history acquisition.
\end{countertheorem}

\begin{proof}
Equation~\eqref{eq:V9-same-diagonal-marginals} gives equality of the two
one-dimensional restricted forms, while
\eqref{eq:V9-opposite-crossing} gives the distinct mixed value.
\end{proof}

% =============================================================================
\section{Exact input packet and row irredundancy}
\label{sec:V9-packet}
% =============================================================================

\begin{acquisition}[Time-split crossing-germ packet]
\label{acq:V9-crossing-card}
The finite packet required before the imported hopping theorem is
\begin{equation}
\boxed{
\begin{aligned}
 \mathfrak C_{\times}=\bigl(&\omega_{\rm bos},
 P_+,P_-,R_\Theta,\varepsilon_\Theta,
 W_1,G_1,\\
 &\{Q_{ji}^{(r)}:1\le i,j\le d,\ 0\le r\le3\},
 (a,b),\mathcal H_{a,b}^{\times},\\
 &q_{\rm line},s_{\rm Pf},
 \text{supports and outward radii}\bigr).
\end{aligned}}
\label{eq:V9-acquisition-card}
\end{equation}
Here $\omega_{\rm bos}$ names the fixed bosonic history,
$G_1=W_1^*W_1$, $q_{\rm line}$ is the absolute determinant-line weight on the
nonzero branch, and $s_{\rm Pf}$ is the separately occurring real Pfaffian sign
when present.  The last two rows are retained for the subsequent exterior
kernel; they are not used to change the crossing matrix.
\end{acquisition}

\begin{theorem}[Finite termination of the crossing-germ card]
\label{thm:V9-card-termination}
Assume the following input-domain checks pass:
\begin{enumerate}[label=\textup{(D\arabic*)},leftmargin=3.4em]
\item all scalar action Reads in \eqref{eq:V9-crossing-queries} occur on the
      same bosonic history and use one fully assembled precision;
\item $\Delta_{\rm split}=0$, $\varepsilon_\Theta\in\{\pm1\}$, and the
      one-letter Gram has its canonical Moore--Penrose support;
\item the held-out equality \eqref{eq:V9-heldout-zero} passes.
\end{enumerate}
Then \eqref{eq:V9-KF-from-queries} uniquely determines the raw crossing matrix
and the imported compiler in \Cref{prop:V9-imported-crossing} returns exactly
one of:
\begin{equation}
 \boxed{
 \begin{gathered}
 \text{relation-descended one-sided hopping factor and its minimum rank},\\
 \text{an explicit one-letter relation, skew, or negative-part obstruction},\\
 \text{an unresolved outward interval at the same finite gate}.
 \end{gathered}}
 \label{eq:V9-terminating-outcomes}
\end{equation}
No covariance, determinant phase, Duhamel residual, current, stress, or
continuum argument is needed to obtain these outcomes.
\end{theorem}

\begin{proof}
The split check makes every vector and reflection in
\Cref{thm:V9-crossing-germ} well typed.  That theorem uniquely reconstructs
$K_F$ entry by entry.  The held-out row is an independent linear combination
of the same matrix and rejects an inconsistent assembly.  The three outcomes
are exactly the relation and positive-cone alternatives of the imported
compiler, with outward uncertainty handled below.  All listed downstream
objects are absent from the construction.
\end{proof}

\begin{theorem}[Row irredundancy of the first crossing gate]
\label{thm:V9-row-irredundancy}
The following rows cannot be deleted from
\eqref{eq:V9-acquisition-card} by any of the remaining rows.
\begin{enumerate}[label=\textup{(I\arabic*)},leftmargin=3.4em]
\item The physical split is not determined by the unsplit precision, its
      positive transfer, determinant modulus, or diagonal time marginals.
\item The Berezin-order sign is not determined by positive line or singular-
      value data.
\item The mixed crossing action is not determined by the two same-time action
      marginals.
\item The one-letter relation support is not determined by the represented
      positive matrix: a raw coefficient form may couple to a zero word.
\item Same-action incidence is not established by fitting all matrix entries;
      at least one independent mixed superposition or an equivalent direct
      common-cylinder check remains held out.
\end{enumerate}
Thus the card is row irredundant relative to the declared crossing target.
\end{theorem}

\begin{proof}
Items~\textup{(I1)--(I3)} are
\Cref{cth:V9-unsplit-no-crossing,cor:V9-line-sign-independent,cth:V9-marginals-no-mixed}.
For item~\textup{(I4)}, take
\begin{equation}
 W_1=\begin{pmatrix}1&0&1\\0&1&1\end{pmatrix},
 \qquad
 v=\begin{pmatrix}-1\\-1\\1\end{pmatrix},
 \qquad W_1v=0,
 \label{eq:V9-W-control}
\end{equation}
and let $K=\diag(2,1)$.  The physical raw form
$K_F=W_1^*KW_1$ annihilates $v$, whereas
$K_F+vv^*$ has the same represented compression on $\Ran W_1^*$ only after
an illicit deletion of its nonzero relation coupling.  Its relation residual
is positive.  Thus the canonical support is independent input.  Item~\textup{(I5)}
is the ordinary distinction between an interpolating table and a held-out
identity; two data sources may populate the fitted entries and disagree on a
new linear combination.  Equation~\eqref{eq:V9-heldout-residual} is the finite
same-action test which excludes that possibility.
\end{proof}

\begin{firewall}[The held-out row is not selected after the fit]
\label{fire:V9-heldout-selection}
The coefficient pair $(a,b)$, the physical split, and every action parameter
are fixed before the entrywise matrix is read.  Choosing $(a,b)$ from a
numerical nullspace of the observed discrepancy would test only the fit and
would not establish common-action occurrence.
\end{firewall}

% =============================================================================
\section{Outward radii from scalar action Reads to the hopping cone}
\label{sec:V9-radii}
% =============================================================================

The target-compressed acquisition is useful only if scalar query errors can be
propagated without losing the exact support and cone semantics.

\begin{theorem}[Four-phase interval propagation]
\label{thm:V9-query-radius}
Suppose the query values have complex outward disks
\begin{equation}
 \abs{\widehat Q_{ji}^{(r)}-Q_{ji}^{(r)}}
 \le\eta_{ji}^{(r)}.
 \label{eq:V9-query-disks}
\end{equation}
Define
\begin{equation}
 \widehat K_{F,ji}
 =-\frac{\varepsilon_\Theta}{4}
   \sum_{r=0}^{3}i^{-r}\widehat Q_{ji}^{(r)},
 \qquad
 \eta_{ji}=\frac14\sum_{r=0}^{3}\eta_{ji}^{(r)},
 \label{eq:V9-entry-radii}
\end{equation}
and
\begin{equation}
 \eta_F=\left(\sum_{i,j=1}^{d}\eta_{ji}^2\right)^{1/2}.
 \label{eq:V9-F-radius}
\end{equation}
Then
\begin{equation}
 \boxed{
 \abs{\widehat K_{F,ji}-K_{F,ji}}\le\eta_{ji},
 \qquad
 \norm{\widehat K_F-K_F}_{\rm HS}\le\eta_F.}
 \label{eq:V9-KF-radius}
\end{equation}
The constants are sharp when the four scalar errors have aligned phases.
\end{theorem}

\begin{proof}
Apply the triangle inequality to
\eqref{eq:V9-KF-from-queries}; every fourth-root coefficient has modulus one.
Squaring and summing the entrywise bounds gives the Hilbert--Schmidt estimate.
Aligned errors attain the entrywise triangle bound.
\end{proof}

\begin{theorem}[Supported relation and represented-operator radii]
\label{thm:V9-supported-radius}
Let $G_1=W_1^*W_1$, $P_W=G_1G_1^\dagger$, and assume the exact supported floor
\begin{equation}
 G_1\succeq g_*P_W,
 \qquad g_*>0.
 \label{eq:V9-word-floor}
\end{equation}
Put
\begin{equation}
 \widehat\rho_{\rm rel}
 =\left(
  \norm{(I-P_W)\widehat K_F}_{\rm HS}^2
  +\norm{\widehat K_F(I-P_W)}_{\rm HS}^2
 \right)^{1/2}.
 \label{eq:V9-estimated-relation}
\end{equation}
Then the exact relation residual $\rho_{\rm rel}$ satisfies
\begin{equation}
 \boxed{
 \abs{\widehat\rho_{\rm rel}-\rho_{\rm rel}}
 \le\sqrt2\,\eta_F.}
 \label{eq:V9-relation-radius}
\end{equation}
If exact relation descent is independently certified, define
\begin{equation}
 \widehat K_{\rm rep}
 =(W_1^\dagger)^*\widehat K_FW_1^\dagger.
 \label{eq:V9-estimated-represented}
\end{equation}
Then
\begin{equation}
 \boxed{
 \norm{\widehat K_{\rm rep}-K_{\rm rep}}_{\rm HS}
 \le \eta_K:=g_*^{-1}\eta_F.}
 \label{eq:V9-represented-radius}
\end{equation}
\end{theorem}

\begin{proof}
The map
$E\mapsto((I-P_W)E,E(I-P_W))$ has norm at most $\sqrt2$ from the
Hilbert--Schmidt space to its orthogonal product, and the norm on a Hilbert
space is one-Lipschitz.  This proves \eqref{eq:V9-relation-radius}.  On the
supported word range,
$\norm{W_1^\dagger}=g_*^{-1/2}$.  Two-sided multiplication of the raw error by
$W_1^\dagger$ therefore gives
\eqref{eq:V9-represented-radius}.
\end{proof}

\begin{corollary}[Validated distance to the hopping cone]
\label{cor:V9-cone-radius}
Let
\begin{equation}
 \widehat d_+
 :=\dist_{\rm HS}(\widehat K_{\rm rep},\operatorname{Herm}_+).
 \label{eq:V9-estimated-cone-distance}
\end{equation}
Then
\begin{equation}
 \boxed{
 \max\{0,\widehat d_+-\eta_K\}
 \le
 \dist_{\rm HS}(K_{\rm rep},\operatorname{Herm}_+)
 \le
 \widehat d_++\eta_K.}
 \label{eq:V9-cone-distance-interval}
\end{equation}
Consequently:
\begin{enumerate}[label=\textup{(V\arabic*)},leftmargin=3.4em]
\item if $\widehat\rho_{\rm rel}>\sqrt2\eta_F$, the exact card has a
      relation obstruction;
\item if $\widehat d_+>\eta_K$, it has a hopping-cone obstruction with a
      strictly positive certified distance;
\item if exact relation descent and exact reflection Hermiticity are supplied,
      put
      \begin{equation}
       \widehat H_{\rm rep}
       =\frac12(\widehat K_{\rm rep}+\widehat K_{\rm rep}^*).
       \label{eq:V9-estimated-Hermitian-part}
      \end{equation}
      If
      \begin{equation}
       \lambda_{\min}(\widehat H_{\rm rep})>\eta_K,
       \label{eq:V9-positive-floor-test}
      \end{equation}
      then $K_{\rm rep}\succ0$ and the hopping factorization passes with a
      certified floor;
\item every valid remaining interval is unresolved at this same finite gate.
\end{enumerate}
\end{corollary}

\begin{proof}
Distance to a nonempty closed subset of a normed space is one-Lipschitz, which
proves \eqref{eq:V9-cone-distance-interval}.  The first two claims use the
strict lower endpoints of \eqref{eq:V9-relation-radius} and
\eqref{eq:V9-cone-distance-interval}.  Under exact Hermiticity, taking Hermitian parts does not increase the
operator error.  Weyl's inequality and
\eqref{eq:V9-represented-radius} therefore give
$\lambda_{\min}(K_{\rm rep})>0$ from
\eqref{eq:V9-positive-floor-test}.  If no strict pass or obstruction margin is
proved, an interval touching the boundary does not decide the branch.
\end{proof}

\begin{remark}[Exact semidefinite passes]
\label{rem:V9-exact-semidefinite}
A singular positive matrix may still pass when exact arithmetic or an exact
factor $R_\times$ proves $K_{\rm rep}=R_\times^*R_\times$.  A floating interval
which merely touches zero does not prove this identity.
\end{remark}

% =============================================================================
\section{Adjacent-cutoff transport of the first crossing gate}
\label{sec:V9-transport}
% =============================================================================

Let two adjacent cards have positive/negative-time isometries
\begin{equation}
 J_\pm:\mathcal H_{\pm,n}\longrightarrow\mathcal H_{\pm,n+1},
 \qquad J_\pm^*J_\pm=I.
 \label{eq:V9-split-isometries}
\end{equation}
Write
$K_n=-\varepsilon_nX_nR_n$ and
$K_{n+1}=-\varepsilon_{n+1}X_{n+1}R_{n+1}$.

\begin{theorem}[Crossing/reflection transport rectangle]
\label{thm:V9-crossing-transport}
Assume the Berezin-order sign is stable,
$\varepsilon_{n+1}=\varepsilon_n=\varepsilon$, and define
\begin{align}
 \delta_{X,n}
 &=\norm{J_+^*X_{n+1}J_- -X_n}_{\rm op},
 \label{eq:V9-X-defect}\\
 \delta_{R,n}
 &=\norm{R_{n+1}J_+-J_-R_n}_{\rm op}.
 \label{eq:V9-R-defect}
\end{align}
Then
\begin{equation}
 \boxed{
 \norm{J_+^*K_{n+1}J_+-K_n}_{\rm op}
 \le
 \delta_{X,n}+\norm{X_{n+1}}_{\rm op}\delta_{R,n}.}
 \label{eq:V9-K-transport}
\end{equation}
The same bound holds in Hilbert--Schmidt norm with the corresponding norms.
In particular, the distance to the Hermitian positive cone changes by at most
the right side.

If $\varepsilon_{n+1}\ne\varepsilon_n$, the discrepancy is a discrete line-
orientation transition.  It is not absorbed into $\delta_{X,n}$ or
$\delta_{R,n}$.
\end{theorem}

\begin{proof}
Using $R_n$ unitary on $\mathcal H_{+,n}$, add and subtract
$J_+^*X_{n+1}J_-R_n$:
\[
\begin{aligned}
 J_+^*K_{n+1}J_+-K_n
 =-\varepsilon\bigl[&
 (J_+^*X_{n+1}J_--X_n)R_n\\
 &+J_+^*X_{n+1}(R_{n+1}J_+-J_-R_n)
 \bigr].
\end{aligned}
\]
Submultiplicativity gives \eqref{eq:V9-K-transport}.  Cone distance is one-
Lipschitz.  A sign change lies in the disconnected set $\{\pm1\}$ and is a
separate orientation/holonomy event.
\end{proof}

\begin{corollary}[Cofinal strict-margin stability at the first gate]
\label{cor:V9-cofinal-crossing}
Suppose the canonical one-letter supports also transport exactly, the relation
rows pass, the signs are eventually constant, and
\begin{equation}
 \sum_n\left(
 \delta_{X,n}+\norm{X_{n+1}}_{\rm op}\delta_{R,n}
 \right)<\infty.
 \label{eq:V9-crossing-summable}
\end{equation}
Then the transported represented crossing matrices are Cauchy on every fixed
finite screen.  A uniform strict positive floor persists to the limit.  A
uniform positive lower distance from the hopping cone persists as an explicit
crossing obstruction.  A collapsing distance with no positive floor remains a
soft first-gate sequence and is not promoted to a Duhamel follower.
\end{corollary}

\begin{proof}
The bound in \Cref{thm:V9-crossing-transport} and summability give the Cauchy
property.  Eigenvalue and distance-to-cone functions are continuous in
operator norm on a fixed finite screen.  The last sentence records that the
hopping gate precedes the dynamic-source short; zero distance only places the
limit in the passing cone and says nothing yet about its subsequent Duhamel
ancestry.
\end{proof}

% =============================================================================
\section{Exact rational control and mutation replay}
\label{sec:V9-control}
% =============================================================================

The accompanying verifier contains two independent exact controls.  The first
is \Cref{cth:V9-unsplit-no-crossing}.  The second uses block order
$\mathcal H_-\oplus\mathcal H_+$ and the non-Hermitian precision
\begin{equation}
 \mathcal K_{\rm ctl}
 =\begin{pmatrix}
  3&0&0&1\\
  0&4&-1&0\\
  2&0&5&0\\
  0&1&0&6
 \end{pmatrix}.
 \label{eq:V9-control-precision}
\end{equation}
Its oriented crossing block is
\begin{equation}
 X=\begin{pmatrix}2&0\\0&1\end{pmatrix},
 \label{eq:V9-control-X}
\end{equation}
while the reverse crossing block is deliberately different.  Take
$R_\Theta=I_2$, $\varepsilon_\Theta=-1$, and the synthesis in
\eqref{eq:V9-W-control}.  Then
\begin{equation}
 G_1=W_1^*W_1
 =\begin{pmatrix}
 1&0&1\\0&1&1\\1&1&2
 \end{pmatrix},
 \label{eq:V9-control-G}
\end{equation}
with relation vector $v=(-1,-1,1)^{\mathsf T}$, and
\begin{equation}
 \boxed{
 K_F=W_1^*XW_1
 =\begin{pmatrix}
 2&0&2\\0&1&1\\2&1&3
 \end{pmatrix}.}
 \label{eq:V9-control-KF}
\end{equation}

\begin{proposition}[Exact control verdict]
\label{prop:V9-control-verdict}
For the control above:
\begin{enumerate}[label=\textup{(R\arabic*)},leftmargin=3.4em]
\item the four-phase formula recovers every entry of the full non-Hermitian
      matrix $\mathcal K_{\rm ctl}$ and every entry of
      \eqref{eq:V9-control-KF} exactly;
\item $K_F=P_WK_FP_W$, and the Moore--Penrose reconstruction in
      \eqref{eq:V9-imported-represented} returns exactly
      $K_{\rm rep}=X$;
\item $X=\diag(\sqrt2,1)^*\diag(\sqrt2,1)$, so the passing hopping rank is two;
\item the mutation
      \begin{equation}
       K_F^{\rm bad}=K_F+vv^*
       \label{eq:V9-bad-KF}
      \end{equation}
      has squared two-sided relation residual exactly $18$ and is rejected
      before any cone projection.
\end{enumerate}
\end{proposition}

\begin{proof}
Item~\textup{(R1)} is \Cref{lem:V9-four-phase} applied to the standard basis
and to the columns of $W_1$.  Since $W_1v=0$, the support of $G_1$ is
$v^\perp$ and the physical form $W_1^*XW_1$ annihilates $v$ on both sides,
proving relation descent.  Direct computation gives
\begin{equation}
 W_1^\dagger
 =\frac13\begin{pmatrix}2&-1\\-1&2\\1&1\end{pmatrix}
 \label{eq:V9-control-pinv}
\end{equation}
and substitution returns $X$.  Its displayed square factor proves item
\textup{(R3)}.  For the mutation, $I-P_W=vv^*/3$ and therefore
both $(I-P_W)K_F^{\rm bad}$ and
$K_F^{\rm bad}(I-P_W)$ equal $vv^*$.  Since
$\norm{vv^*}_{\rm HS}^2=\norm v^4=9$, the sum of the two squared residuals is
$18$.
\end{proof}

\begin{firewall}[Control scope]
\label{fire:V9-control-scope}
The rational matrices above replay the finite algebra and its mutation
semantics.  They are not the active chiral Dirac--Yukawa precision, do not
select a bosonic RPESM history, and do not supply a determinant/Pfaffian line,
current, stress tensor, or continuum estimate.
\end{firewall}

% =============================================================================
\section{Integrated V9 theorem and the revised stopping boundary}
\label{sec:V9-integrated}
% =============================================================================

\begin{theorem}[Version V9 physical time-split and crossing-germ theorem]
\label{thm:V9-integrated}
Preserving every theorem and limitation of Versions~V1--V8, Version~V9 proves:
\begin{enumerate}[label=\textup{(V9.\arabic*)},leftmargin=5.4em]
\item the reflected hopping gate is earlier than the V8 common-Hamiltonian
      transfer tangent on the active terminal quantum branch;
\item the finite split residual \eqref{eq:V9-split-residual} is an exact domain
      test for complementary time sectors and their unitary linear reflection;
\item one fixed positive precision, positive transfer, determinant modulus,
      and pair of diagonal time marginals admit valid splits with opposite
      crossing signs, so none determines the physical hopping branch;
\item the Berezin-order sign is an independent orientation row;
\item arbitrary non-Hermitian quadratic coefficients obey the exact four-phase
      identity \eqref{eq:V9-four-phase};
\item one four-phase mixed action orbit per ordered one-letter pair reconstructs
      the complete raw reflected crossing matrix by
      \eqref{eq:V9-KF-from-queries};
\item scalar energy origins and all same-time block terms cancel from this
      target-native acquisition;
\item one predeclared held-out mixed pair tests same-action assembly before the
      relation or cone short is used;
\item outward scalar radii propagate through the exact word support to a
      validated interval for the distance from the hopping cone;
\item adjacent-cutoff crossing and reflection motion obey the separate rectangle
      \eqref{eq:V9-K-transport}, while a sign change is a discrete line event;
\item the exact rational control recovers a rank-two hopping factor and rejects
      a relation-breaking mutation with squared residual $18$.
\end{enumerate}
No item populates the active ESM split/crossing tuple, proves its determinant or
Pfaffian line sign, constructs the exterior covariance, identifies the
common-action transfer Hamiltonian, or opens the current/stress/Einstein rows.
\end{theorem}

\begin{proof}
The order statement is \Cref{prop:V9-imported-crossing,ass:V9-V8-order}.
The split criterion is \Cref{thm:V9-split-criterion}.  The nonidentifiability and
sign claims are
\Cref{cth:V9-unsplit-no-crossing,cor:V9-line-sign-independent}.  The
polarization and acquisition statements are
\Cref{lem:V9-four-phase,thm:V9-crossing-germ,cor:V9-block-diagonal-invariance,cor:V9-heldout-crossing}.
The interval and transport claims are
\Cref{thm:V9-query-radius,thm:V9-supported-radius,cor:V9-cone-radius,thm:V9-crossing-transport}.
The replay is \Cref{prop:V9-control-verdict}.  Every excluded conclusion belongs
to a later gate than the unpopulated packet
\eqref{eq:V9-acquisition-card}.
\end{proof}

\begin{assessment}[Status after V9]
\label{ass:V9-status}
\passstatus{} for the exact split-domain criterion, the split and sign
nonidentifiability theorems, arbitrary sesquilinear polarization, the
one-letter crossing-germ acquisition, block-diagonal invariance, the held-out
same-action row, outward support/cone radii, adjacent crossing/reflection
transport, and the exact rational replay.

\obsstatus{} for inferring a crossing branch from an unsplit transfer or from
separate time-sector marginals, for inferring the Berezin sign from a positive
line modulus, for deleting a one-letter relation, and for repairing a failed
physical crossing by nearest-cone projection.

\stopstatus{} on the original Einstein--Standard--Model regulator at the first
unpopulated active tuple
\begin{equation}
 (\omega_{\rm bos},P_+,P_-,R_\Theta,
   \varepsilon_\Theta,\mathcal K,W_1)
 \label{eq:V9-active-stop-tuple}
\end{equation}
and its same-history mixed action values.  V8 is retained as an optional
subsequent common-Hamiltonian acquisition, while V5--V6 remain the subsequent
Duhamel/absorption compilers.
\end{assessment}

\begin{openproblem}[Version V10 is licensed only by one active time-split card]
\label{op:V9-next}
Preserve Versions~V1--V9 verbatim.  No new generic source, lag, localizer,
concentration, or continuum theorem is admissible.  A Version~V10 continuation
may proceed only after one named active bosonic history supplies:
\begin{enumerate}[label=\textup{(V10.\arabic*)},leftmargin=5.5em]
\item the complete uneliminated chiral precision and its physically declared
      positive/negative-time projections, reflection, boundary convention,
      and Berezin-order sign;
\item the occurring one-letter synthesis and its canonical support;
\item the four-phase mixed action panel or the direct crossing block, with one
      independent held-out mixed pair;
\item the absolute determinant-line weight, divisor data, and real Pfaffian
      sign on the same sector cover;
\item one adjacent time-split/reflection route if a cofinal statement is sought.
\end{enumerate}
The first action is to run the inherited relation and hopping residuals.  A
passing factor then opens the already established hard-tail covariance and
exterior-word tests.  A relation, skew, negative-part, line, divisor, or route
witness is exported immediately.  If the tuple is not supplied, the programme
stops at \eqref{eq:V9-active-stop-tuple}; it does not reopen V8 or invent a
replacement toy slab.
\end{openproblem}

\section*{V9 verification and append-only record}
\addcontentsline{toc}{section}{V9 verification and append-only record}
The accompanying exact script uses symbolic rational arithmetic.  It verifies
both orthogonal split packets in \Cref{cth:V9-unsplit-no-crossing}, their equal
diagonal marginals and opposite $\pm12/25$ crossings, every four-phase matrix
entry of the non-Hermitian control precision, the raw three-word matrix
\eqref{eq:V9-control-KF}, its Moore--Penrose represented recovery, the rank-two
positive factor, and the relation mutation with squared residual $18$.

The complete Version~V8 source preceding its sole final document-closing
command is preserved byte for byte.  Its preterminal SHA--256 digest is
\begin{center}\small\ttfamily
321375e96b219f6462c541d982c6e14d3eaf4fbe23aeb334dd005de6215a4c3b
\end{center}
A later licensed continuation must remove only the sole final
\verb|\end{document}|, append new material immediately before it, restore one
terminal command, and increment the filename to end in \texttt{\_V10.tex}.

% =============================================================================
% END APPENDED VERSION V9 MATERIAL
% =============================================================================


% =============================================================================
% START APPENDED VERSION V10 MATERIAL -- complete V1--V9 preterminal source preserved
% =============================================================================

\clearpage
\hypersetup{
 pdftitle={Renewal Einstein--Standard--Model Physical Source Metric, Duhamel Ancestry and Quantum Continuum Programme V10},
 pdfsubject={The common-history Grassmann quadratic germ and the reflection-fixed boundary Schur gate}
}
\section*{Version V10: the common-history Grassmann germ and the reflection-fixed boundary gate}
\addcontentsline{toc}{section}{Version V10: Grassmann germ and reflection-fixed boundary gate}

\begin{abstract}
Version~V9 correctly moved the active quantum stopping boundary ahead of the
Hamiltonian/Duhamel compilers: a reflected hopping verdict requires an
occurring time split, reflection and Berezin sign, complete crossing block, and
one-letter synthesis on one bosonic history.  Its stopping tuple still contains
two logically prior constructions which were left implicit.  First, the
``complete precision'' must be the bilinear connected coefficient of one
unnormalized same-history fermion cylinder, rather than a matrix inferred from
a determinant modulus, a covariance, or separately sampled action pieces.
Second, an actual reflection plane may carry fixed boundary modes.  Until those
modes are assigned or integrated, the positive- and negative-time spaces are
not complementary and the crossing queried in Version~V9 is not yet the
physical crossing.

Version~V10 gives finite exact compilers for both rows.  The nilpotent logarithm
of a finite Grassmann cylinder has a unique degree-$(1,1)$ coefficient.  It
recovers the complete chiral precision directly, is insensitive to a common
bosonic scalar weight, remains defined on determinant divisors, and separates
the quadratic crossing germ from all later connected four- and higher-fermion
terms.  A positive same-history assembly residual then tests whether all
reported gauge, Higgs, Yukawa, boundary, and counterterm contributions really
sum to that coefficient on the same bosonic support.

For a chronological decomposition
$\mathcal H=\mathcal H_-\oplus\mathcal H_0\oplus\mathcal H_+$, an exact
three-sector residual decides whether the linear part of reflection exchanges
the two open half-spaces and preserves the fixed boundary sector.  If the
boundary block $D$ of the complete precision is invertible, finite Berezin
Gaussian elimination gives
\[
 \det D\,\exp[-\bar\psi(A-BD^{-1}C)\psi].
\]
Consequently the crossing used by the reflected-hopping theorem is the direct
crossing minus one signed boundary-return term.  Only the action of $D^{-1}$ on
the occurring negative-time one-letter columns is required.  A three-mode
positive-definite example has identical diagonal time and boundary blocks,
zero direct crossing, equal spectrum and determinant, but opposite effective
crossing signs because the mixed boundary incidence has opposite orientation.

On a divisor, Moore--Penrose replacement is forbidden.  The boundary marginal
has zero scalar body and can remain a nonzero nilpotent polynomial through
zero-mode saturation; it is not a nonzero Gaussian line times an effective
precision.  The determinant/adjugate hierarchy must then remain explicit.
Outward block errors give a closed singular-value-controlled radius for the
boundary-return crossing and feed directly into the Version~V9 relation and
cone intervals.  The active regulator still does not populate the common
Grassmann cylinder, chronological fixed-plane record, boundary block, or its
same-history solves, so the physical verdict remains stopped one arrow earlier
than Version~V9's four-phase crossing Read.
\end{abstract}

% =============================================================================
\section{The first missing object is an unnormalized connected fermion action}
\label{sec:V10-Grassmann-germ}
% =============================================================================

Let $E\cong\C^d$ be a finite one-particle coefficient space with ordered basis.
Write
\[
 \mathfrak G_E
 =\Lambda_{\C}(\bar e_1,\ldots,\bar e_d,e_1,\ldots,e_d)
\]
for the associated Grassmann algebra.  The coefficient of the unit is called
the \emph{body}.  Denote by $[\bar e_i e_j]X$ the coefficient of the ordered
monomial $\bar e_i e_j$ in $X$.

\begin{definition}[Connected action and quadratic Grassmann germ]
\label{def:V10-connected-action}
Let $\mathscr W\in\mathfrak G_E^{\rm even}$ have positive body $w_0>0$.
Writing $w_0^{-1}\mathscr W=1+N$, define
\begin{equation}
 \operatorname{Log}_0\mathscr W
 :=\sum_{m=1}^{2d}\frac{(-1)^{m+1}}mN^m,
 \qquad
 \mathscr A_{\mathscr W}:=-\operatorname{Log}_0\mathscr W.
 \label{eq:V10-nilpotent-log}
\end{equation}
The \emph{quadratic Grassmann germ} is the matrix
\begin{equation}
 \boxed{
 \mathcal K_{\mathscr W,ij}
 :=[\bar e_i e_j]\mathscr A_{\mathscr W}.}
 \label{eq:V10-quadratic-germ}
\end{equation}
The ordering of the generators is part of the Berezin-orientation packet and
is fixed before any coefficient is read.
\end{definition}

\begin{theorem}[Exact zero-safe recovery of the complete precision]
\label{thm:V10-Grassmann-log}
The sum in \eqref{eq:V10-nilpotent-log} is finite and is the unique logarithm
of $w_0^{-1}\mathscr W$ in the nilpotent ideal.  If
\begin{equation}
 \mathscr W
 =w_0\exp\!\left[-\sum_{i,j}\bar e_iK_{ij}e_j
                   -\mathscr S_{\ge4}\right],
 \label{eq:V10-general-fermion-weight}
\end{equation}
where every monomial of $\mathscr S_{\ge4}$ has Grassmann degree at least four,
then
\begin{equation}
 \boxed{
 \mathscr A_{\mathscr W}
 =\sum_{i,j}\bar e_iK_{ij}e_j+\mathscr S_{\ge4},
 \qquad
 \mathcal K_{\mathscr W}=K.}
 \label{eq:V10-log-recovers-K}
\end{equation}
In particular:
\begin{enumerate}[label=\textup{(G\arabic*)},leftmargin=3.5em]
\item multiplication of $\mathscr W$ by any positive bosonic scalar leaves
      $\mathcal K_{\mathscr W}$ unchanged;
\item all connected terms of degree at least four leave the quadratic germ
      unchanged;
\item multiplication by
      $\exp[-\sum\bar e_iB_{ij}e_j]$ changes the germ by the physical amount
      $B$;
\item the construction remains defined when $\det K=0$ and requires neither
      $K^{-1}$ nor a normalized fermionic Gaussian law.
\end{enumerate}
\end{theorem}

\begin{proof}
The augmentation ideal of $\mathfrak G_E$ is nilpotent of index at most
$2d+1$, so $N^{2d+1}=0$ and the displayed logarithm is finite.  The ordinary
formal identities $\log(\exp X)=X$ and $\exp(\log(1+N))=1+N$ hold exactly in a
nilpotent algebra because all power series truncate.  Applying them to
\eqref{eq:V10-general-fermion-weight} gives
\eqref{eq:V10-log-recovers-K}.  A scalar body is removed before the logarithm,
whereas an additional bilinear exponential adds its bilinear coefficient to
the connected action.  No determinant or inverse enters any step.
\end{proof}

The preceding theorem extracts a matrix from one occurring unnormalized
fermion cylinder.  It does not by itself prove that separately submitted action
terms occurred on that same cylinder.

\begin{definition}[Same-history quadratic assembly residual]
\label{def:V10-assembly-residual}
Let $\Omega_B$ be a finite bosonic history set with strictly positive
unnormalized weight $\Lambda_B(\omega)$ and
$Z_B=\sum_\omega\Lambda_B(\omega)$.  Suppose the full cylinder gives the germ
$\mathcal K(\omega)$ and the declared component actions give germs
$\mathcal K_\alpha(\omega)$ on the same history.  Put
\begin{equation}
 \boxed{
 \Delta_{\rm quad}^{\rm occ}
 :=\frac1{Z_B}\sum_{\omega\in\Omega_B}
 \Lambda_B(\omega)
 \left\|\mathcal K(\omega)-\sum_\alpha
                 \mathcal K_\alpha(\omega)\right\|_{\rm HS}^2.}
 \label{eq:V10-assembly-residual}
\end{equation}
\end{definition}

\begin{theorem}[Exact common-history assembly criterion]
\label{thm:V10-common-history-assembly}
One has $\Delta_{\rm quad}^{\rm occ}=0$ if and only if
\begin{equation}
 \boxed{
 \mathcal K(\omega)=\sum_\alpha\mathcal K_\alpha(\omega)
 \quad\text{for every }\omega\in\supp\Lambda_B.}
 \label{eq:V10-pointwise-assembly}
\end{equation}
Thus separately correct gauge, Higgs, Yukawa, boundary, or counterterm
marginals do not populate the full precision unless a common cylinder or a
direct same-event Read supplies \eqref{eq:V10-pointwise-assembly}.
\end{theorem}

\begin{proof}
Every summand of \eqref{eq:V10-assembly-residual} is nonnegative and every
weight on the support is strictly positive.  Their weighted sum vanishes
exactly when every squared norm vanishes.  The final sentence is the domain
meaning of the pointwise identity: terms indexed by unrelated histories do not
even define its left-hand side.
\end{proof}

\begin{countertheorem}[The crossing germ does not decide conditional Gaussianity]
\label{cth:V10-germ-no-Gaussianity}
For $d\ge2$, fix any matrix $K$ and any nonzero $g\in\C$.  The two cylinders
\begin{align}
 \mathscr W_0
 &=\exp[-\bar eKe],
 \label{eq:V10-W0}\\
 \mathscr W_g
 &=\exp[-\bar eKe
       -g\,\bar e_1e_1\bar e_2e_2]
 \label{eq:V10-Wg}
\end{align}
have the same body and the same quadratic germ $K$, but their connected
four-fermion coefficients differ by $g$.  Hence a passing Version~V9 hopping
factorization does not establish the later conditional exterior-kernel or
quasi-free branch.
\end{countertheorem}

\begin{proof}
Apply \Cref{thm:V10-Grassmann-log}.  The bilinear coefficient is $K$ in both
cases, while the coefficient of the displayed degree-four monomial in the
connected action is respectively zero and $g$.
\end{proof}

\begin{assessment}[What is acquired before determinants and covariances]
\label{ass:V10-germ-order}
The quadratic germ is the target-native occurrence row for the complete
precision.  On an invertible card, a determinant together with the complete
complex covariance could reconstruct the same matrix, but that is a stronger
downstream packet and fails on the divisor.  A determinant modulus, a
phase-quenched law, or a selected covariance compression is not substituted
for the same-history coefficient in \eqref{eq:V10-quadratic-germ}.
\end{assessment}

% =============================================================================
\section{Chronological three-sector occurrence and the fixed reflection plane}
\label{sec:V10-three-sector}
% =============================================================================

Version~V9 used complementary positive- and negative-time projections.  An
actual lattice reflection can also have a fixed one-particle boundary sector.
That sector must be retained until a boundary convention has occurred.

\begin{definition}[Chronological fixed-plane packet]
\label{def:V10-chronological-packet}
A candidate packet is
\begin{equation}
 \mathfrak T_{0,\Theta}
 =(P_-,P_0,P_+,\Theta_0),
 \label{eq:V10-chronological-packet}
\end{equation}
where the $P_\sigma$ and $\Theta_0$ act on one finite carrier.  Define
\begin{equation}
\boxed{
\begin{aligned}
 \Delta_{0,\Theta}:={}&
 \sum_{\sigma\in\{-,0,+\}}
 \left(\|P_\sigma-P_\sigma^*\|_{\rm HS}^2
      +\|P_\sigma^2-P_\sigma\|_{\rm HS}^2\right)\\
 &+\sum_{\sigma<\tau}\|P_\sigma P_\tau\|_{\rm HS}^2
  +\|P_-+P_0+P_+-I\|_{\rm HS}^2\\
 &+\|\Theta_0^*\Theta_0-I\|_{\rm HS}^2
  +\|\Theta_0\Theta_0^*-I\|_{\rm HS}^2
  +\|\Theta_0^2-I\|_{\rm HS}^2\\
 &+\|\Theta_0P_+-P_-\Theta_0\|_{\rm HS}^2
  +\|\Theta_0P_--P_+\Theta_0\|_{\rm HS}^2
  +\|\Theta_0P_0-P_0\Theta_0\|_{\rm HS}^2.
\end{aligned}}
\label{eq:V10-chronological-residual}
\end{equation}
The anti-linear conjugation has already been separated, exactly as in the
Version~V9 convention; $\Theta_0$ is its linear part.
\end{definition}

\begin{theorem}[Exact fixed-plane chronology criterion]
\label{thm:V10-chronology-criterion}
One has $\Delta_{0,\Theta}=0$ if and only if
\begin{equation}
 \boxed{
 \mathcal H=\mathcal H_-\oplus\mathcal H_0\oplus\mathcal H_+}
 \label{eq:V10-threeway-split}
\end{equation}
is an orthogonal decomposition, $\Theta_0$ is a unitary involution, it exchanges
$\mathcal H_+$ and $\mathcal H_-$, and it preserves $\mathcal H_0$.
On this branch
\begin{equation}
 R_\Theta:=P_-\Theta_0P_+:
 \mathcal H_+\longrightarrow\mathcal H_-
 \label{eq:V10-restricted-reflection}
\end{equation}
is unitary.  After a physically declared removal or assignment of
$\mathcal H_0$, the resulting two-sector packet satisfies the Version~V9 split
criterion exactly.
\end{theorem}

\begin{proof}
All summands in \eqref{eq:V10-chronological-residual} are nonnegative.  Their
vanishing makes the three $P_\sigma$ mutually orthogonal complementary
projections.  The next three terms make $\Theta_0$ a unitary involution, and
the final intertwining identities give the stated action on the three ranges.
Conversely those geometric conditions annihilate every residual term.  On the
zero branch,
\[
 R_\Theta^*R_\Theta
 =P_+\Theta_0^*P_-\Theta_0P_+=P_+,
 \qquad
 R_\Theta R_\Theta^*=P_-,
\]
so the restriction is unitary.  Removing the orthogonal fixed sector leaves
complementary exchanged half-spaces and therefore the two-sector identities
of Version~V9.
\end{proof}

\begin{firewall}[A reflection-fixed mode is not assigned by a sign convention]
\label{fire:V10-fixed-mode}
The discrete Berezin-order sign or an orientation of the determinant line does
not decide whether a fixed-plane one-particle mode is retained, constrained,
or integrated.  That is a boundary occurrence/domain row.  Deleting the mode
because its later Schur return is inconvenient is an ill-typed input, not a
positive hopping certificate.
\end{firewall}

% =============================================================================
\section{Exact boundary Gaussian elimination and the physical crossing}
\label{sec:V10-boundary-Schur}
% =============================================================================

Use the chronological head
$\mathcal H_{\rm h}=\mathcal H_-\oplus\mathcal H_+$ and boundary space
$\mathcal H_0$.  Relative to this decomposition write the complete quadratic
germ as
\begin{equation}
 \mathcal K
 =\begin{pmatrix}A&B\\ C&D\end{pmatrix},
 \qquad
 A:\mathcal H_{\rm h}\to\mathcal H_{\rm h},
 \quad D:\mathcal H_0\to\mathcal H_0.
 \label{eq:V10-K-block}
\end{equation}
No Hermiticity is assumed.  Fix the Berezin measure convention by
\begin{equation}
 \int \mathcal D(\bar\chi,\chi)
       \exp(-\bar\chi D\chi)=\det D.
 \label{eq:V10-Berezin-convention}
\end{equation}

\begin{theorem}[Boundary-complete fermion Schur identity]
\label{thm:V10-Grassmann-Schur}
If $D$ is invertible, then
\begin{equation}
\boxed{
\begin{aligned}
 &\int \mathcal D(\bar\chi,\chi)
 \exp\!\left[-\bar\psi A\psi-\bar\psi B\chi
              -\bar\chi C\psi-\bar\chi D\chi\right]\\
 &\qquad=\det D\,
 \exp\!\left[-\bar\psi
       \left(A-BD^{-1}C\right)\psi\right].
\end{aligned}}
\label{eq:V10-Grassmann-Schur}
\end{equation}
Thus the boundary-complete head precision is
\begin{equation}
 \boxed{\mathcal K_{\rm eff}=A-BD^{-1}C.}
 \label{eq:V10-Keff}
\end{equation}
Writing $P_\pm$ now for the two head projections, its oriented crossing is
\begin{equation}
 \boxed{
 X_\times^{\rm eff}
 =P_+AP_- - P_+BD^{-1}CP_-.}
 \label{eq:V10-effective-crossing}
\end{equation}
The determinant factor $\det D$ belongs to the line/divisor packet and is not
absorbed into the crossing matrix.
\end{theorem}

\begin{proof}
Make the translation
\[
 \chi'=\chi+D^{-1}C\psi,
 \qquad
 \bar\chi'=\bar\chi+\bar\psi BD^{-1}.
\]
Berezin translation invariance preserves the measure.  Direct expansion gives
\[
 \bar\psi A\psi+\bar\psi B\chi+\bar\chi C\psi+\bar\chi D\chi
 =\bar\psi(A-BD^{-1}C)\psi+\bar\chi'D\chi'.
\]
The integral factors and \eqref{eq:V10-Berezin-convention} proves
\eqref{eq:V10-Grassmann-Schur}.  Compressing the effective matrix from
negative to positive time gives \eqref{eq:V10-effective-crossing}.
\end{proof}

\begin{corollary}[Target-minimal boundary solve on the one-letter bank]
\label{cor:V10-minimal-boundary-solve}
Let $W_1:E_1\to\mathcal H_+$ be the occurring one-letter synthesis and let
$R_\Theta:\mathcal H_+\to\mathcal H_-$ be the restricted reflection.  It is
sufficient to solve
\begin{equation}
 \boxed{
 DZ_-=CP_-R_\Theta W_1}
 \label{eq:V10-boundary-solve}
\end{equation}
for the minimum collection of boundary columns $Z_-$.  The complete raw
reflected one-letter crossing is then
\begin{equation}
\boxed{
\begin{aligned}
 K_F^{\rm eff}
 =-\varepsilon_\Theta\bigl(&
 W_1^*P_+AP_-R_\Theta W_1\\
 &-W_1^*P_+BZ_-\bigr).
\end{aligned}}
\label{eq:V10-KF-effective}
\end{equation}
No full inverse of $D$, boundary covariance outside the occurring column
range, or exterior-word state is required.  One held-out coefficient vector
$u$ checks the solve through
\begin{equation}
 \mathcal H_{0,u}
 :=\|DZ_-u-CP_-R_\Theta W_1u\|^2.
 \label{eq:V10-boundary-heldout}
\end{equation}
\end{corollary}

\begin{proof}
Invertibility makes \eqref{eq:V10-boundary-solve} equivalent to
$Z_-=D^{-1}CP_-R_\Theta W_1$.  Substitute this identity into
\eqref{eq:V10-effective-crossing}, multiply by the fixed reflection/Berezin
sign, and compress with $W_1$.  The held-out quantity vanishes for the same
linear solve and is nonnegative by definition.
\end{proof}

\begin{assessment}[Assembly before the Version V9 cone short]
\label{ass:V10-assembly-before-cone}
The two terms of \eqref{eq:V10-KF-effective} are signed and may cancel or
reverse sign.  They are assembled before relation descent, Hermitian/skew
splitting, negative-part extraction, or positive factorization.  The direct
head crossing alone is the correct Version~V9 input only when the fixed sector
is absent, is physically constrained so that the return is zero, or has
already been integrated by \eqref{eq:V10-Grassmann-Schur}.
\end{assessment}

% =============================================================================
\section{A positive-definite separator for the boundary-return sign}
\label{sec:V10-boundary-counterexample}
% =============================================================================

\begin{countertheorem}[Identical direct cards, opposite effective crossings]
\label{cth:V10-boundary-return-sign}
Let $\sigma\in\{+1,-1\}$ and, in block order
$\mathcal H_-\oplus\mathcal H_0\oplus\mathcal H_+$, put
\begin{equation}
 \mathcal K_\sigma
 =\begin{pmatrix}
 3&\sigma&0\\
 \sigma&2&1\\
 0&1&3
 \end{pmatrix}.
 \label{eq:V10-Ksigma}
\end{equation}
Then:
\begin{enumerate}[label=\textup{(S\arabic*)},leftmargin=3.4em]
\item both matrices are positive definite, have determinant $12$, and are
      orthogonally similar, hence have the same spectrum and every
      unitary-invariant positive modulus;
\item their negative-time, fixed-plane, and positive-time diagonal blocks are
      identical, their direct negative-to-positive crossing is zero, and the
      two boundary-coupling vectors have identical separate Gram data;
\item after integrating the fixed mode,
      \begin{equation}
       \boxed{
       \mathcal K_{\rm eff,\sigma}
       =\begin{pmatrix}
         5/2&-\sigma/2\\
         -\sigma/2&5/2
        \end{pmatrix},}
       \label{eq:V10-Keff-sigma}
      \end{equation}
      so the effective negative-to-positive crossing is $-\sigma/2$.
\end{enumerate}
Thus the direct crossing, all three diagonal blocks, the determinant, the
spectrum, and the separate coupling norms do not determine the reflected
hopping branch.  The signed mixed boundary-return incidence is an independent
same-history row.
\end{countertheorem}

\begin{proof}
The leading principal minors are $3$, $5$, and
\[
 \det\mathcal K_\sigma
 =3(6-1)-\sigma(3\sigma)=12,
\]
so Sylvester's criterion gives positivity.  The two matrices are conjugate by
$\diag(-1,1,1)$, proving spectral equality.  Their direct corner and diagonal
claims are immediate, and the coupling magnitudes are both one.

For the head order $(-,+)$ one has
\[
 A=3I_2,
 \qquad
 B=\binom{\sigma}{1},
 \qquad
 D=(2),
 \qquad C=B^*.
\]
Substitution in $A-BD^{-1}C$ gives
\eqref{eq:V10-Keff-sigma}.  Its off-diagonal coefficient is the asserted
crossing.
\end{proof}

\begin{corollary}[A fitted boundary orientation can manufacture a hopping sign]
\label{cor:V10-no-fitted-boundary}
Choosing the sign of a boundary coupling, a fixed-plane basis, or a boundary
identification after reading the desired effective crossing changes the
physical same-history return in \eqref{eq:V10-effective-crossing}.  Such a
choice is not a harmless coordinate fit unless the complete precision,
reflection, and Berezin orientation are transported by the same predeclared
unitary.
\end{corollary}

\begin{proof}
The pair \eqref{eq:V10-Ksigma} has all listed downstream positive invariants in
common and differs only in the signed mixed boundary incidence which controls
the effective crossing.  A common predeclared unitary transports the whole
packet and preserves the represented cone verdict; changing only the mixed
incidence selects a different physical cylinder.
\end{proof}

% =============================================================================
\section{The divisor branch: why Moore--Penrose elimination is not Berezin marginalization}
\label{sec:V10-divisor}
% =============================================================================

\begin{theorem}[Zero body on a singular boundary block]
\label{thm:V10-singular-boundary}
Let $\mathscr M(\bar\psi,\psi)$ be the boundary marginal on the left side of
\eqref{eq:V10-Grassmann-Schur}, without assuming that $D$ is invertible.  Its
scalar body is
\begin{equation}
 \boxed{[1]\mathscr M=\det D.}
 \label{eq:V10-marginal-body}
\end{equation}
Consequently, if $\det D=0$, the marginal cannot have the form
\begin{equation}
 c\exp(-\bar\psi K_{\rm eff}\psi)
 \label{eq:V10-nonzero-Gaussian-form}
\end{equation}
with $c\ne0$.  If the marginal is not identically zero, it is a nilpotent
zero-mode-saturation polynomial and the pure Gaussian head compiler is not
well typed.  Replacing $D^{-1}$ by $D^\dagger$ does not in general compute the
Berezin marginal.
\end{theorem}

\begin{proof}
Set all head generators to zero.  The boundary integral reduces to
\eqref{eq:V10-Berezin-convention}, so the body is $\det D$.  A nonzero scalar
multiple of an exponential has body $c$, proving the exclusion.  A nonzero
Grassmann polynomial with zero body has strictly positive degree and therefore
cannot be represented by \eqref{eq:V10-nonzero-Gaussian-form}.  The following
explicit example disproves Moore--Penrose substitution.
\end{proof}

\begin{countertheorem}[One zero mode defeats the pseudoinverse shortcut]
\label{cth:V10-pseudoinverse-not-marginal}
Take one head pair $(\bar h,h)$ and one boundary pair
$(\bar\chi,\chi)$, with $A=D=0$ and $B=C=1$.  Under the ordered coefficient
convention
\begin{equation}
 \int\mathcal D(\bar\chi,\chi)\,\bar\chi\chi=-1,
 \label{eq:V10-one-mode-convention}
\end{equation}
one has
\begin{equation}
 \boxed{
 \int\mathcal D(\bar\chi,\chi)
 \exp[-\bar h\chi-\bar\chi h]
 =\bar h h\ne0.}
 \label{eq:V10-zero-mode-marginal}
\end{equation}
But $D^\dagger=0$ and the formal Moore--Penrose Schur matrix is
$A-BD^\dagger C=0$.  Multiplication by $\det D=0$ gives the zero polynomial,
whereas assigning a pseudo-determinant one gives a nonzero scalar body.  Both
outputs are wrong.
\end{countertheorem}

\begin{proof}
Put $x=\bar h\chi$ and $y=\bar\chi h$.  They are even nilpotents and commute,
so
\[
 \exp[-(x+y)]=1-x-y+xy.
\]
In the canonical order $\bar h,h,\bar\chi,\chi$ one has
$xy=-\bar h h\bar\chi\chi$.  The orientation in \eqref{eq:V10-one-mode-convention} changes the coefficient sign and gives
\eqref{eq:V10-zero-mode-marginal}.  The pseudoinverse statements are direct.
\end{proof}

\begin{firewall}[The divisor is a branch, not an inconvenience]
\label{fire:V10-divisor}
On $\det D=0$, the correct outputs are the boundary zero-mode exterior factor,
its divisor order and orientation, and the induced higher head polynomial.
They may subsequently be tested by the already retained exterior-adjugate
hierarchy.  A Moore--Penrose effective precision is not inserted merely to
keep the Version~V9 one-particle cone test available.
\end{firewall}

% =============================================================================
\section{Outward boundary-return radii and cofinal transport}
\label{sec:V10-boundary-radii}
% =============================================================================

The inverse branch has a closed finite perturbation calculus.  Write
\begin{equation}
 X=A_\times-B_+D^{-1}C_-,
 \label{eq:V10-X-error-notation}
\end{equation}
where $A_\times=P_+AP_-$, $B_+=P_+B$, and $C_-=CP_-$.  Use hats for a
submitted packet and suppose
\begin{equation}
 \|\widehat A_\times-A_\times\|\le\delta_A,
 \quad
 \|\widehat B_+-B_+\|\le\delta_B,
 \quad
 \|\widehat C_--C_-\|\le\delta_C,
 \quad
 \|\widehat D-D\|\le\delta_D.
 \label{eq:V10-block-radii}
\end{equation}

\begin{theorem}[Validated boundary-Schur crossing radius]
\label{thm:V10-boundary-radius}
Assume
\begin{equation}
 s_{\min}(D)\ge d_*>0,
 \qquad
 0\le\delta_D<d_*.
 \label{eq:V10-D-floor}
\end{equation}
Then $\widehat D$ is invertible,
\begin{equation}
 \|\widehat D^{-1}\|
 \le(d_*-\delta_D)^{-1},
 \qquad
 \|\widehat D^{-1}-D^{-1}\|
 \le\frac{\delta_D}{d_*(d_*-\delta_D)},
 \label{eq:V10-inverse-radius}
\end{equation}
and
\begin{equation}
\boxed{
\begin{aligned}
 \|\widehat X-X\|
 \le\delta_X:={}&\delta_A
 +\frac{\delta_B\|\widehat C_-\|}{d_*-\delta_D}\\
 &+\frac{\|B_+\|\,\delta_D\|\widehat C_-\|}
        {d_*(d_*-\delta_D)}
 +\frac{\|B_+\|\,\delta_C}{d_*}.
\end{aligned}}
\label{eq:V10-X-radius}
\end{equation}
If $W_1$ and $R_\Theta$ are exact, the raw one-letter matrix satisfies
\begin{equation}
 \boxed{
 \|\widehat K_F^{\rm eff}-K_F^{\rm eff}\|
 \le\|W_1\|^2\delta_X.}
 \label{eq:V10-KF-radius}
\end{equation}
This radius is inserted into the Version~V9 support, relation, Hermiticity and
positive-cone intervals; it is not combined with the determinant-line radius.
\end{theorem}

\begin{proof}
Weyl's singular-value inequality gives
$s_{\min}(\widehat D)\ge d_*-\delta_D$.  The resolvent identity
\[
 \widehat D^{-1}-D^{-1}
 =D^{-1}(D-\widehat D)\widehat D^{-1}
\]
proves \eqref{eq:V10-inverse-radius}.  Expand
\[
 \widehat B_+\widehat D^{-1}\widehat C_-
 -B_+D^{-1}C_-
 =(
 \widehat B_+-B_+)\widehat D^{-1}\widehat C_-
 +B_+(\widehat D^{-1}-D^{-1})\widehat C_-
 +B_+D^{-1}(\widehat C_--C_-)
\]
and apply the norm bounds, then add the direct-block radius.  Compression by
$W_1^*$ and $R_\Theta W_1$ has norm at most $\|W_1\|^2$, since
$R_\Theta$ is unitary.
\end{proof}

\begin{corollary}[Boundary-return cutoff rectangle]
\label{cor:V10-boundary-transport}
Let adjacent cards be transported to common head and boundary screens.  If the
transported block defects obey \eqref{eq:V10-block-radii}, the boundary blocks
have a common singular-value floor $d_*>0$, and
\begin{equation}
 \sum_n\bigl(
 \delta_{A,n}+\delta_{B,n}+\delta_{C,n}+\delta_{D,n}
 \bigr)<\infty
 \label{eq:V10-summable-boundary-defects}
\end{equation}
with uniformly bounded transported $B_+,C_-$, then the effective crossings are
Cauchy on every fixed one-letter screen.  Adding the already separate
Version~V9 reflection and one-letter transport defects gives a unique cofinal
crossing limit.
\end{corollary}

\begin{proof}
For all late $n$, the denominators in \eqref{eq:V10-X-radius} are bounded below
by a positive constant.  Uniform bounds on $B_+,C_-$ turn the right side into
a fixed multiple of the summand in
\eqref{eq:V10-summable-boundary-defects}.  The telescoping series therefore
converges in operator norm.  Version~V9 then transports the reflection and
one-letter compressions separately.
\end{proof}

\begin{firewall}[No fixed-plane inverse without a physical floor]
\label{fire:V10-no-boundary-inverse-fit}
A small singular value of $D$ is not hidden inside a large crossing-error bar.
If no positive floor is certified, the finite packet is stopped at the
boundary divisor/conditioning row.  Selecting a boundary subspace after
removing its soft vector is the same forbidden post hoc support fit already
identified for the Duhamel and hopping shorts.
\end{firewall}

% =============================================================================
\section{Exact finite replay and the revised active acquisition}
\label{sec:V10-replay}
% =============================================================================

\begin{proposition}[Exact replay identities]
\label{prop:V10-replay}
The accompanying exact-rational verifier performs the following independent
checks.
\begin{enumerate}[label=\textup{(R\arabic*)},leftmargin=3.4em]
\item In a two-mode Grassmann algebra it forms an arbitrary rational bilinear
      precision, adds a nonzero quartic connected term and a positive scalar
      body, exponentiates, takes the finite nilpotent logarithm, and recovers
      every bilinear and quartic coefficient exactly.
\item For both matrices \eqref{eq:V10-Ksigma}, it verifies all leading
      principal minors, determinant $12$, equality of characteristic
      polynomials, orthogonal similarity, zero direct crossing, and the
      effective crossings $\mp1/2$.
\item For the singular one-boundary-mode card, exact Grassmann coefficient
      extraction gives $-\bar h h$, while the Moore--Penrose shortcut gives
      neither its body nor its nilpotent coefficient.
\item Mutating one same-history component matrix makes
      \eqref{eq:V10-assembly-residual} strictly positive and is rejected.
\end{enumerate}
All output is exact rational or integer data.  No entry is a physical RPESM
measurement.
\end{proposition}

\begin{proof}
Each item is a direct evaluation of the finite identities proved above.  The
verifier implements Grassmann monomials as ordered bit masks and computes all
signs from inversion parity.  Matrix determinants, characteristic
polynomials, Schur complements, and residuals use exact rational arithmetic.
\end{proof}

\begin{definition}[First active Version V10 physical card]
\label{def:V10-active-card}
The target-minimal upstream packet is
\begin{equation}
\boxed{
\begin{gathered}
 \texttt{COMMON\_HISTORY\_GRASSMANN\_QUADRATIC\_GERM}\\[-0.15em]
 \texttt{\_CHRONOLOGICAL\_FIXED\_PLANE\_AND}\\[-0.15em]
 \texttt{\_BOUNDARY\_COMPLETE\_CROSSING\_CARD}.
\end{gathered}}
\label{eq:V10-active-card-name}
\end{equation}
Its minimum payload is:
\begin{enumerate}[label=\textup{(A\arabic*)},leftmargin=3.5em]
\item one named positive unnormalized bosonic field-shell history and the
      complete ordered conditional Grassmann weight;
\item the quadratic germ extracted by \eqref{eq:V10-quadratic-germ}, together
      with the pointwise component assembly residual
      \eqref{eq:V10-assembly-residual};
\item the physical three-sector chronological record
      $(P_-,P_0,P_+,\Theta_0)$, its boundary convention, and the independently
      oriented Berezin line;
\item if $P_0=0$, the direct Version~V9 time-split tuple;
\item if $P_0\ne0$ and $D$ is invertible, a singular-value floor, the direct
      crossing, the boundary solve \eqref{eq:V10-boundary-solve}, the signed
      return, and the determinant-line factor $\det D$;
\item if $D$ is singular, the divisor order, boundary zero-mode exterior
      factor, and its induced head polynomial, with no pseudoinverse Gaussian
      replacement;
\item the occurring one-letter synthesis, one held-out boundary solve, and the
      four-phase or direct effective crossing values required by Version~V9;
\item one adjacent/direct chronological and boundary-block route if a cofinal
      verdict is sought.
\end{enumerate}
\end{definition}

\begin{theorem}[Version V10 upstream quadratic-germ and boundary theorem]
\label{thm:V10-integrated}
Preserving every theorem and limitation of Versions~V1--V9, Version~V10 proves:
\begin{enumerate}[label=\textup{(V10.\arabic*)},leftmargin=5.8em]
\item the complete finite fermion precision is the degree-$(1,1)$ coefficient
      of the connected logarithm of one unnormalized same-history Grassmann
      cylinder;
\item this coefficient is zero-safe, independent of a common bosonic scalar,
      and independent of later connected terms of degree at least four;
\item one positive assembly residual decides whether all declared physical
      action terms sum to the same precision on every occupied history;
\item a three-sector residual reconstructs the open half-spaces, fixed
      reflection plane, and restricted unitary reflection before the
      complementary Version~V9 split is used;
\item an invertible fixed-plane block gives the exact Berezin Schur marginal,
      its determinant-line factor, and the direct-minus-return effective
      crossing;
\item one boundary solve on the occurring one-letter columns is sufficient to
      populate that return;
\item positive-definite full precisions with identical direct cards and all
      separate positive invariants can have opposite effective crossing signs;
\item a singular boundary block has zero marginal body and may yield a nonzero
      zero-mode polynomial, so Moore--Penrose elimination is not a zero-safe
      fermion marginal;
\item a certified boundary singular-value floor and outward block radii give
      the explicit effective-crossing error
      \eqref{eq:V10-X-radius} and a summable cutoff rectangle;
\item the exact replay verifies the connected-action, boundary-return, divisor,
      and mutation identities.
\end{enumerate}
No item populates the active RPESM Grassmann cylinder, fixes its chronological
boundary convention, or supplies the physical determinant/Pfaffian line,
exterior kernel, Duhamel packet, current, stress, charge, continuum, or
Einstein residual.
\end{theorem}

\begin{proof}
Items~\textup{(V10.1)--(V10.3)} are
\Cref{thm:V10-Grassmann-log,thm:V10-common-history-assembly,cth:V10-germ-no-Gaussianity}.
Item~\textup{(V10.4)} is \Cref{thm:V10-chronology-criterion}.  Items
\textup{(V10.5)--(V10.6)} are
\Cref{thm:V10-Grassmann-Schur,cor:V10-minimal-boundary-solve}.  Item
\textup{(V10.7)} is \Cref{cth:V10-boundary-return-sign}.  Item
\textup{(V10.8)} is
\Cref{thm:V10-singular-boundary,cth:V10-pseudoinverse-not-marginal}.  Item
\textup{(V10.9)} is
\Cref{thm:V10-boundary-radius,cor:V10-boundary-transport}, and the replay is
\Cref{prop:V10-replay}.  Every excluded output lies downstream of the still
unpopulated card in \eqref{eq:V10-active-card-name}.
\end{proof}

\begin{assessment}[Status after Version V10]
\label{ass:V10-status}
\passstatus{} for the finite connected Grassmann logarithm, zero-safe quadratic
germ, same-history assembly residual, fixed-plane chronology criterion,
invertible boundary Berezin--Schur identity, target-minimal boundary solve,
exact separator, divisor obstruction, outward crossing radius, cutoff
rectangle, and exact replay.

\obsstatus{} for inferring the complete precision from determinant modulus or a
normalized proxy; for applying the Version~V9 direct crossing before a
reflection-fixed sector is assigned; for shorting direct and boundary-return
terms separately; and for replacing a singular fermion marginal by a
Moore--Penrose Gaussian.

\stopstatus{} on the original Einstein--Standard--Model regulator at the first
unpopulated same-history packet
\begin{equation}
 (\Omega_B,\Lambda_B,\mathscr W_\omega,
   P_-,P_0,P_+,\Theta_0,
   \varepsilon_\Theta,W_1)
 \label{eq:V10-physical-stop}
\end{equation}
and its extracted quadratic/boundary blocks.  Version~V9 is run immediately
once this packet lands on its boundary-free or invertible-boundary branch.
\end{assessment}

\begin{openproblem}[Version V11: evaluate the one occurring quadratic boundary card]
\label{op:V10-next}
Preserve Versions~V1--V10 verbatim.  Do not add another source, lag, localizer,
word, concentration, or continuum compiler.  Proceed only in this order.
\begin{enumerate}[label=\textup{(V11.\arabic*)},leftmargin=5.8em]
\item Freeze one named active RPESM bosonic history, ordered Grassmann cylinder,
      chronological reflection record, and boundary convention before reading
      any crossing sign.
\item Extract the connected bilinear germ and run the same-history component
      assembly residual.  A positive residual is the first occurrence stop.
\item Evaluate the fixed-plane chronology residual.  Reject a split selected
      from the desired hopping verdict.
\item On $P_0=0$, run Version~V9 directly.  On $P_0\ne0$, first decide the
      boundary divisor.  If $D$ is invertible, acquire its floor and the one-
      letter boundary solve; if it is singular, export the zero-mode exterior
      polynomial and line/divisor packet without pseudoinverse replacement.
\item Assemble the direct and boundary-return crossing once, then apply the
      Version~V9 relation and positive-cone compiler with its held-out mixed
      action Read.
\item Only a passing hopping and line packet opens the conditional exterior
      kernel and later V8/V5/V6 common-Hamiltonian and Duhamel rows.
\end{enumerate}
If the physical cylinder is still absent, the exact endpoint is
\eqref{eq:V10-physical-stop}.  Further generic mathematics would now repeat a
compiler whose first required field-shell coefficient has not occurred.
\end{openproblem}

\section*{V10 verification and append-only record}
\addcontentsline{toc}{section}{V10 verification and append-only record}
The accompanying verifier uses only exact integer and rational arithmetic.  It
implements the finite Grassmann product, exponential, logarithm, coefficient
extraction and boundary integration; checks the two positive-definite boundary
cards and their effective crossings; verifies the singular zero-mode marginal;
and rejects a same-history action mutation.

The complete Version~V9 source preceding its sole final document-closing
command is preserved byte for byte.
The complete Version~V9 preterminal byte string has SHA--256 digest
\begin{center}\small\ttfamily
671e4760b4123df32770c41327b11e23e3b453767c34d18c413f15a05222c5e2
\end{center}
The exact verifier and its JSON output have SHA--256 digests
\begin{center}\small\ttfamily
9d0a1c0e7f85e11975a19349924a3b5963f749e8079d16b7e544b900cd2b0652\\
488b193e748de1bd041d1d0da7fa427d6123b13be5d6f809aa1ced882470a8d6
\end{center}
  A later continuation must remove only the
sole final \verb|\end{document}|, append new material immediately before it,
restore one terminal command, and increment the filename to end in
\texttt{\_V11.tex}.

% =============================================================================
% END APPENDED VERSION V10 MATERIAL
% =============================================================================


% =============================================================================
% START APPENDED VERSION V11 MATERIAL -- complete V1--V10 preterminal source preserved
% =============================================================================

\clearpage
\hypersetup{
 pdftitle={Renewal Einstein--Standard--Model Physical Source Metric, Duhamel Ancestry and Quantum Continuum Programme V11},
 pdfsubject={Positive-time bosonic readability, record-relative reflected positivity, and the hidden-history innovation}
}
\section*{Version V11: positive-time bosonic readability and the record-relative reflected kernel}
\addcontentsline{toc}{section}{Version V11: positive-time bosonic readability and the record-relative reflected kernel}

\begin{abstract}
Version~V10 correctly identified the unnormalized Grassmann quadratic germ,
the chronological fixed-plane decomposition, and the boundary-complete
crossing as inputs which precede the reflected-hopping and Duhamel compilers.
Its active card nevertheless froze one named full bosonic history before
asking which part of that history is actually selectable by the positive-time
observable algebra.  That distinction is essential.  A positive mixture of
conditional fermion kernels can be reflection positive even when some
conditional histories are not; conversely, a conditional failure cannot be
hidden once its bosonic selector is itself an occurring positive-time Read.
The conditional Gaussian statement, the integrated one-letter statement, and
the statement on an intermediate readable record are therefore three
different targets.

This continuation moves to that earlier record-incidence row.  For a finite
bosonic cylinder and a predeclared readable record, the reflected one-letter
form is the direct sum of one \emph{unnormalized conditional matrix} per record
cell.  Positivity on the represented positive-time algebra is equivalent to
positivity of exactly those matrices and requires no division by a rare-cell
probability.  Refining the record can expose a negative direction or a
zero-word coupling which cancels after coarse mixing.  An exact two-history
example has positive coarse kernel $I_2$ although both history-resolved kernels
have a negative eigenvalue.  A second example has zero coarse relation
residual although both history-resolved matrices couple to a represented zero
word.

The conditional matrices form an operator-valued martingale under record
refinement.  Their Hilbert--Schmidt innovations obey exact Pythagoras, while
the squared distance to the relation-descended positive cone is monotone under
refinement by conditional Jensen.  A strict coarse floor survives to the full
history whenever the operator-norm record innovations have total size below
that floor.  Without such a bound, integrated positivity is only an
integrated verdict.

One selector-weighted quadratic Read per basis direction and polarization
pair reconstructs every unnormalized cell matrix.  A held-out refinement
identity tests that coarse and fine cells came from one cylinder, and outward
scalar radii give a direct matrix radius and a terminating pass,
obstruction, or unresolved interval.  The Version~V10 historywise
Grassmann/boundary card remains mandatory when conditional Gaussianity, a
historywise hopping factor, or exterior second quantization is claimed.  It is
not charged merely to prove a one-letter Osterwalder--Schrader statement on a
coarser physically readable record.
\end{abstract}

% =============================================================================
\section{The bosonic Read precedes conditioning on a named history}
\label{sec:V11-readable-record}
% =============================================================================

Let $(\Omega,\rho)$ be a finite positive bosonic cylinder with
$\rho(\omega)>0$ on its declared support.  Fix a finite complex coefficient
space $F$ which has already been transported to one common positive-time
one-letter screen.  For every $\omega\in\Omega$, let
\begin{equation}
 H(\omega)\in\operatorname{End}(F)
 \label{eq:V11-history-kernel}
\end{equation}
be the unnormalized conditional reflected one-letter form on that screen.
The determinant/Pfaffian line weight, the physical reflection convention, the
one-letter synthesis, and any Version~V10 boundary elimination which has
actually been performed are included in $H(\omega)$.  No Hermiticity or
positivity is assumed at this point.  Alternatively, $H(\omega)$ may denote a
direct same-event one-letter Read without a conditional Gaussian
factorization.

A finite bosonic record is a map
\begin{equation}
 r:\Omega\longrightarrow\mathcal R
 \label{eq:V11-record-map}
\end{equation}
fixed before the matrices in \eqref{eq:V11-history-kernel} are inspected.  It
is called \emph{positive-time readable} when its cell indicators are actual
admitted positive-time bosonic Reads.  Mere presence of the full history in a
global event ledger does not by itself make every singleton indicator an
admitted positive-time observable.

For $a\in\mathcal R$ put
\begin{equation}
 p_a=\sum_{r(\omega)=a}\rho(\omega),
 \qquad
 M_a=\sum_{r(\omega)=a}\rho(\omega)H(\omega).
 \label{eq:V11-cell-matrices}
\end{equation}
Zero-mass record labels are removed.  The normalized conditional mean is
$\overline H_a=p_a^{-1}M_a$, but the unnormalized matrix $M_a$ is the native
entry of the reflected form.

\begin{definition}[Record-relative reflected one-letter form]
\label{def:V11-record-form}
On the direct-sum coefficient space
\begin{equation}
 F_r=\bigoplus_{a\in\mathcal R}F,
 \label{eq:V11-record-carrier}
\end{equation}
define
\begin{equation}
 \boxed{
 \mathfrak q_r(v,w)
 =\sum_{a\in\mathcal R}\langle v_a,M_aw_a\rangle,
 \qquad v=(v_a),\quad w=(w_a).}
 \label{eq:V11-record-form}
\end{equation}
Equivalently, if a positive-time word uses coefficient $v_a$ whenever the
record says $a$, then \eqref{eq:V11-record-form} is its mixed reflected
expectation on the original cylinder.
\end{definition}

\begin{theorem}[Exact record-relative Osterwalder--Schrader criterion]
\label{thm:V11-record-OS}
The form $\mathfrak q_r$ is Hermitian and nonnegative if and only if
\begin{equation}
 \boxed{M_a=M_a^*\succeq0\quad\text{for every }a\in\mathcal R.}
 \label{eq:V11-cell-positive}
\end{equation}
In particular, no lower bound on $p_a$ and no division by $p_a$ are needed for
this positivity statement.
\end{theorem}

\begin{proof}
Insert \eqref{eq:V11-cell-matrices} into the mixed expectation and group the
finite sum by the value of $r$.  This gives \eqref{eq:V11-record-form}.
If every $M_a$ is Hermitian positive semidefinite, each summand is a positive
Hermitian form and so is their direct sum.  Conversely, fix $a$ and choose a
vector in $F_r$ supported only in its $a$-component.  Hermiticity and
nonnegativity of $\mathfrak q_r$ then say
$\langle v,M_av\rangle\ge0$ for every $v\in F$.  Complex polarization first
forces $M_a=M_a^*$ and the same inequality gives $M_a\succeq0$.
\end{proof}

\begin{corollary}[The exact hierarchy of mixed, readable, and conditional targets]
\label{cor:V11-record-hierarchy}
Three cases are distinct.
\begin{enumerate}[label=\textup{(\roman*)},leftmargin=3em]
\item If $r$ is the constant record, the sole matrix is
      \begin{equation}
       M_{\rm mix}=\sum_{\omega}\rho(\omega)H(\omega).
       \label{eq:V11-mixed-matrix}
      \end{equation}
      Its positivity is exactly the integrated one-letter verdict.
\item If $r$ is the full history record, then
      \begin{equation}
       M_\omega=\rho(\omega)H(\omega),
       \label{eq:V11-full-history-matrix}
      \end{equation}
      and positivity is equivalent to the historywise conditional verdict.
\item If $s:\Omega\to\mathcal S$ refines $r$ through
      $r=\pi\circ s$, then
      \begin{equation}
       \boxed{M_a=\sum_{\pi(b)=a}M_b^{s}.}
       \label{eq:V11-record-additivity}
      \end{equation}
      Hence positivity on the finer readable record implies positivity on the
      coarser record.  The converse is false.
\end{enumerate}
\end{corollary}

\begin{proof}
The first two claims are the corresponding choices of the fibres in
\eqref{eq:V11-cell-matrices}.  The refinement identity is finite additivity of
the same positive cylinder over the disjoint subcells.  A sum of positive
matrices is positive.  The counterexample below proves failure of the reverse
implication.
\end{proof}

\begin{countertheorem}[Integrated positivity can hide conditional negativity]
\label{cth:V11-hidden-negativity}
Let $\Omega=\{0,1\}$ with weights $1/2$, let $F=\C^2$, and set
\begin{equation}
 H(0)=\begin{pmatrix}-1&0\\0&3\end{pmatrix},
 \qquad
 H(1)=\begin{pmatrix}3&0\\0&-1\end{pmatrix}.
 \label{eq:V11-hidden-negative-matrices}
\end{equation}
For the constant record,
\begin{equation}
 \boxed{M_{\rm mix}=\frac12H(0)+\frac12H(1)=I_2\succ0.}
 \label{eq:V11-hidden-negative-average}
\end{equation}
For the full history record, each cell has a negative direction.  Thus the
integrated one-letter kernel passes with floor one while the conditional
history statement fails on both histories.
\end{countertheorem}

\begin{proof}
The displayed average is immediate.  The first coordinate is negative for
$H(0)$ and the second coordinate is negative for $H(1)$.  Multiplication by
the positive history weights does not change those signs.
\end{proof}

\begin{firewall}[A global event record is not automatically a positive-time selector]
\label{fire:V11-event-versus-OS-read}
The population theorem for a common field-shell cylinder may retain the full
bosonic history as provenance.  The historywise implication in
\eqref{eq:V11-full-history-matrix} is licensed only when the claimed
conditional law is itself the target, or when the corresponding selectors are
admitted in the positive-time test algebra.  If the physical target contains
only a coarser record, replacing that record by singleton histories is target
inflation.  Conversely, when singleton selectors really occur, the coarse
matrix \eqref{eq:V11-mixed-matrix} cannot hide a failing conditional cell.
\end{firewall}

% =============================================================================
\section{Record resolution and the relation-descended positive cone}
\label{sec:V11-cone-resolution}
% =============================================================================

Fix an exact one-letter relation support $P=P^*=P^2$ on $F$, transported before
any history value is read.  Define the closed convex cone
\begin{equation}
 \mathfrak C_P
 =\{A\in\operatorname{End}(F):A=PAP,\ A=A^*,\ A\succeq0\}.
 \label{eq:V11-supported-cone}
\end{equation}
It combines the Version~V9 zero-word relation and reflected-hopping cone on one
fixed represented screen.  The explicit Version~V9 relation, skew, and
negative-part residuals remain the finite evaluator of membership; they are
not reproved here.

For a readable record $r$, regard
\begin{equation}
 \overline H_r(\omega)=\overline H_{r(\omega)}
 =\mathbb E_\rho[H\mid r](\omega)
 \label{eq:V11-conditional-matrix-field}
\end{equation}
as an $L^2(\rho;\operatorname{End}(F))$ matrix field and define
\begin{equation}
 \boxed{
 \mathfrak D_r
 =\mathbb E_\rho\,
   \dist_{\rm HS}\bigl(\overline H_r,\mathfrak C_P\bigr)^2.}
 \label{eq:V11-record-cone-defect}
\end{equation}

\begin{theorem}[Cone obstruction is monotone under record refinement]
\label{thm:V11-cone-monotone}
If $s$ refines $r$, then
\begin{equation}
 \boxed{0\le\mathfrak D_r\le\mathfrak D_s.}
 \label{eq:V11-cone-monotone}
\end{equation}
Consequently a fine relation/hopping pass implies the coarse pass, whereas a
coarse pass proves only the verdict on the coarse readable carrier.
\end{theorem}

\begin{proof}
Conditional expectation gives
\begin{equation}
 \overline H_r=\mathbb E_\rho[\overline H_s\mid r].
 \label{eq:V11-tower-matrix}
\end{equation}
Distance to a nonempty closed convex subset of a Hilbert space is convex.  On
each $r$-cell, therefore,
\[
 \dist_{\rm HS}(\overline H_r,\mathfrak C_P)
 \le
 \mathbb E_\rho[
   \dist_{\rm HS}(\overline H_s,\mathfrak C_P)\mid r].
\]
Squaring and applying conditional Jensen gives the same inequality with the
squared distance on the right.  Integration proves
\eqref{eq:V11-cone-monotone}.
\end{proof}

\begin{countertheorem}[Coarse mixing can erase a zero-word incidence defect]
\label{cth:V11-hidden-relation}
Let
\begin{equation}
 P=\begin{pmatrix}1&0\\0&0\end{pmatrix},
 \qquad
 L_0=\begin{pmatrix}1&1\\1&0\end{pmatrix},
 \qquad
 L_1=\begin{pmatrix}1&-1\\-1&0\end{pmatrix}
 \label{eq:V11-hidden-relation-matrices}
\end{equation}
with equal history weights.  Then
\begin{equation}
 \boxed{\frac12L_0+\frac12L_1=P\in\mathfrak C_P.}
 \label{eq:V11-hidden-relation-average}
\end{equation}
However the two-sided Version~V9 relation residual is
\begin{equation}
 \boxed{
 \|(I-P)L_j\|_{\rm HS}^2+\|L_j(I-P)\|_{\rm HS}^2=2,
 \qquad j=0,1.}
 \label{eq:V11-hidden-relation-residual}
\end{equation}
Thus an averaged zero-word relation need not be a historywise relation.
\end{countertheorem}

\begin{proof}
The off-support entries have opposite signs and cancel in the average.  For
each $j$, $(I-P)L_j$ consists of the second row and $L_j(I-P)$ of the second
column.  Each has squared Hilbert--Schmidt norm one, proving
\eqref{eq:V11-hidden-relation-residual}.
\end{proof}

\begin{assessment}[Exact meaning of a coarse hopping pass]
\label{ass:V11-coarse-pass}
A vanishing $\mathfrak D_r$ says that the \emph{record-relative} one-letter
functional descends and is positive on every physically selectable $r$-cell.
It does not produce a conditional precision, a historywise one-sided hopping
factor, or a Gaussian exterior hierarchy.  Those stronger claims require the
full-history Version~V10 card or an independently complete higher-word
occurrence packet.  The difference is a missing physical selector/conditioning
row, not a weakness of the positive-cone calculation.
\end{assessment}

% =============================================================================
\section{The hidden-history matrix martingale}
\label{sec:V11-martingale}
% =============================================================================

Let
\begin{equation}
 \sigma(r_0)\subseteq\sigma(r_1)\subseteq\cdots
 \label{eq:V11-record-filtration}
\end{equation}
be an increasing sequence of physically declared readable records on one
positive cylinder, and put
\begin{equation}
 H_n=\mathbb E_\rho[H\mid r_n].
 \label{eq:V11-matrix-martingale}
\end{equation}
The matrix space carries the Hilbert--Schmidt inner product and is regarded as
a real Hilbert space when Hermitian and skew parts are discussed.

\begin{theorem}[Exact matrix-valued record Pythagoras]
\label{thm:V11-martingale-Pythagoras}
For every $n$,
\begin{equation}
 \boxed{
 \mathbb E_\rho\|H_{n+1}\|_{\rm HS}^2
 =\mathbb E_\rho\|H_n\|_{\rm HS}^2
  +\mathbb E_\rho\|H_{n+1}-H_n\|_{\rm HS}^2.}
 \label{eq:V11-forward-Pythagoras}
\end{equation}
Moreover,
\begin{equation}
 \boxed{
 \mathbb E_\rho\|H-H_n\|_{\rm HS}^2
 =\mathbb E_\rho\|H-H_{n+1}\|_{\rm HS}^2
  +\mathbb E_\rho\|H_{n+1}-H_n\|_{\rm HS}^2.}
 \label{eq:V11-tail-Pythagoras}
\end{equation}
Thus
\begin{equation}
 \Xi_n^{\rm hist}
 :=\mathbb E_\rho\|H-H_n\|_{\rm HS}^2
 \label{eq:V11-hidden-history-stock}
\end{equation}
is the exact positive one-letter information hidden inside the current
bosonic record.
\end{theorem}

\begin{proof}
Conditional expectation is the orthogonal projection from
$L^2(\rho;\operatorname{End}(F))$ onto the subspace of
$\sigma(r_n)$-measurable matrix fields.  The increment $H_{n+1}-H_n$ is
orthogonal to every $r_n$-measurable field, in particular to $H_n$.  This gives
\eqref{eq:V11-forward-Pythagoras}.  The terminal residual $H-H_{n+1}$ is
orthogonal to every $r_{n+1}$-measurable field, in particular to the same
increment, which proves \eqref{eq:V11-tail-Pythagoras}.
\end{proof}

\begin{corollary}[Finite exhaustion and exact conditional landing]
\label{cor:V11-finite-exhaustion}
If $\Omega$ is finite and the records eventually separate all histories, then
$H_n=H$ for all sufficiently large $n$ and $\Xi_n^{\rm hist}=0$.  More
generally, if the generated record sigma-algebra is the full one and
$H\in L^2$, the matrix martingale converges to $H$ in $L^2$.
\end{corollary}

\begin{proof}
On a finite set, an increasing sequence of partitions stabilizes after finitely
many strict refinements.  Full separation makes conditional expectation the
identity.  The general statement is the ordinary $L^2$ martingale convergence
theorem applied to the finite-dimensional Hilbert-valued random variable
$H$.
\end{proof}

\begin{theorem}[A strict coarse floor survives a summable record refinement]
\label{thm:V11-floor-survival}
Assume every $H_n$ is Hermitian on one fixed supported carrier $P F$ and
\begin{equation}
 H_0\succeq\alpha P,
 \qquad \alpha>0.
 \label{eq:V11-initial-floor}
\end{equation}
Put
\begin{equation}
 \epsilon_n
 =\operatorname*{ess\,sup}_{\omega}
   \|H_{n+1}(\omega)-H_n(\omega)\|_{\rm op}.
 \label{eq:V11-record-increment}
\end{equation}
If
\begin{equation}
 \sum_{n=0}^{\infty}\epsilon_n<\alpha,
 \label{eq:V11-summable-record-increments}
\end{equation}
then, for every $m$,
\begin{equation}
 \boxed{
 H_m\succeq
 \left(\alpha-\sum_{n=0}^{m-1}\epsilon_n\right)P,}
 \label{eq:V11-finite-floor}
\end{equation}
and the limiting matrix field satisfies
\begin{equation}
 \boxed{
 H_\infty\succeq
 \left(\alpha-\sum_{n=0}^{\infty}\epsilon_n\right)P.}
 \label{eq:V11-limit-floor}
\end{equation}
If the records exhaust the full history, this is a historywise conditional
floor.
\end{theorem}

\begin{proof}
For every vector $v\in PF$ and almost every history,
\[
 |\langle v,(H_{n+1}-H_n)v\rangle|
 \le\epsilon_n\|v\|^2.
\]
Telescope from $0$ to $m$ and use \eqref{eq:V11-initial-floor}; this proves
\eqref{eq:V11-finite-floor}.  The summability assumption makes $(H_m)$ Cauchy
in $L^\infty$ operator norm, so it has a uniform limit and the Loewner
inequality passes to that limit.  Full-history exhaustion identifies the limit
with $H$.
\end{proof}

\begin{remark}[The sufficient bridge is genuinely additional]
\label{rem:V11-floor-bridge-sharpness}
In \Cref{cth:V11-hidden-negativity}, the constant-record floor is one, while
the one-step operator-norm innovation from $I_2$ to either history matrix is
two.  Thus the hypothesis of \Cref{thm:V11-floor-survival} fails exactly where
conditional negativity appears.  A small $L^2$ innovation without an
operator-norm or target-visible control is not enough for a uniform
historywise floor.
\end{remark}

% =============================================================================
\section{Target-minimal selector-weighted acquisition}
\label{sec:V11-acquisition}
% =============================================================================

The record matrices can be acquired without reconstructing the conditional
precision or dividing by the cell probability.  For each occurring record
label $a$ and vector $z\in F$, define the unnormalized quadratic Read
\begin{equation}
 Q_a(z)
 =\sum_{r(\omega)=a}\rho(\omega)
    \langle z,H(\omega)z\rangle
 =\langle z,M_az\rangle.
 \label{eq:V11-cell-quadratic-read}
\end{equation}
The selector and the word probe are placed on the same event.  Separate
marginals of the selector and the fermionic word do not define this quantity.

\begin{theorem}[Exact Hermitian polarization of one record cell]
\label{thm:V11-cell-polarization}
Assume $M_a=M_a^*$ and let $e_1,\ldots,e_d$ be the frozen coefficient basis.
Then
\begin{equation}
 \boxed{(M_a)_{ii}=Q_a(e_i),}
 \label{eq:V11-polar-diagonal}
\end{equation}
and, for $i<j$,
\begin{equation}
 \boxed{
 \begin{aligned}
 (M_a)_{ij}=\frac14\bigl(&Q_a(e_i+e_j)-Q_a(e_i-e_j)\\
 &+\mathrm iQ_a(e_i-\mathrm i e_j)
  -\mathrm iQ_a(e_i+\mathrm i e_j)\bigr).
 \end{aligned}}
 \label{eq:V11-polar-offdiagonal}
\end{equation}
Thus $d^2$ real scalar coordinates, obtainable from the displayed quadratic
probes, determine the complete unnormalized cell matrix.
\end{theorem}

\begin{proof}
Expand $Q_a(x+y)$ and $Q_a(x-y)$ to isolate twice the real part of
$\langle x,M_ay\rangle$.  Expanding $Q_a(x\pm\mathrm i y)$ isolates its
imaginary part with the sign displayed in
\eqref{eq:V11-polar-offdiagonal}.  Taking $x=e_i$, $y=e_j$ gives the formula.
A Hermitian $d\times d$ matrix has $d^2$ real coordinates.
\end{proof}

\begin{theorem}[Outward matrix radius and finite cone verdict]
\label{thm:V11-outward-cell}
Suppose each scalar probe in \eqref{eq:V11-polar-diagonal}--
\eqref{eq:V11-polar-offdiagonal} is enclosed by an estimate with nonnegative
radius.  Let $\eta_{a,ii}$ be the diagonal radius and let
\begin{equation}
 \eta_{a,ij}
 =\frac14(\eta_{a,+}+\eta_{a,-}
               +\eta_{a,-\mathrm i}+\eta_{a,+\mathrm i}),
 \qquad i<j,
 \label{eq:V11-entry-radius}
\end{equation}
be the sum of the four corresponding radii.  The Hermitian matrix estimate
$\widehat M_a$ reconstructed from their midpoints satisfies
\begin{equation}
 \boxed{
 \|\widehat M_a-M_a\|_{\rm HS}
 \le
 \eta_a^{\rm mat}
 :=\left(
   \sum_i\eta_{a,ii}^2
   +2\sum_{i<j}\eta_{a,ij}^2
 \right)^{1/2}.}
 \label{eq:V11-matrix-radius}
\end{equation}
Assume the relation support $P$ has been established independently.  Then:
\begin{enumerate}[label=\textup{(\roman*)},leftmargin=3em]
\item if $M_a=P M_aP$ is exact and
      \begin{equation}
       \lambda_{\min}
       (P\widehat M_aP|_{PF})>\eta_a^{\rm mat},
       \label{eq:V11-cell-pass}
      \end{equation}
      then $M_a$ has a strict positive floor on $PF$;
\item if
      \begin{equation}
       \dist_{\rm HS}(\widehat M_a,\mathfrak C_P)
       >\eta_a^{\rm mat},
       \label{eq:V11-cell-obstruction}
      \end{equation}
      then the exact cell has a relation, skew, or negative-part obstruction;
\item every remaining valid interval is unresolved at this finite precision.
\end{enumerate}
\end{theorem}

\begin{proof}
Each polarization entry is a linear combination of the acquired scalars, so
the triangle inequality gives \eqref{eq:V11-entry-radius}; summing the squared
entry bounds gives \eqref{eq:V11-matrix-radius}.  Weyl's inequality proves
item~\textup{(i)}.  Distance to a closed set is one-Lipschitz, hence
\[
 \dist_{\rm HS}(M_a,\mathfrak C_P)
 \ge
 \dist_{\rm HS}(\widehat M_a,\mathfrak C_P)
 -\|\widehat M_a-M_a\|_{\rm HS};
\]
this proves item~\textup{(ii)}.  If neither strict inequality holds, the
uncertainty set meets both the passing cone and its complement or fails to
separate them, so no stronger theorem verdict is licensed.
\end{proof}

\begin{definition}[Held-out record-refinement residual]
\label{def:V11-record-additivity-residual}
For an independently frozen refinement $s$ of $r$, define
\begin{equation}
 \boxed{
 \Delta_{s/r}^{\rm rec}
 =\sum_{a\in\mathcal R}
  \left\|M_a-
   \sum_{\pi(b)=a}M_b^s\right\|_{\rm HS}^2.}
 \label{eq:V11-record-additivity-residual}
\end{equation}
The coarse and fine matrices used here are acquired on disjoint calibration
queries or with one direct held-out cell sum; the refinement is not chosen to
make the residual small.
\end{definition}

\begin{proposition}[Same-cylinder record incidence]
\label{prop:V11-record-incidence}
On one common positive cylinder and one deterministic refinement,
\begin{equation}
 \boxed{\Delta_{s/r}^{\rm rec}=0.}
 \label{eq:V11-record-incidence-zero}
\end{equation}
A strictly positive outward lower interval for this residual is a
same-history selector/word incidence obstruction.  In the exact replay below,
changing one fine-cell entry by $1/7$ while retaining the coarse matrix gives
\begin{equation}
 \boxed{\Delta_{s/r}^{\rm rec}=\frac1{49}.}
 \label{eq:V11-record-mutation}
\end{equation}
\end{proposition}

\begin{proof}
Equation \eqref{eq:V11-record-additivity} proves the zero statement term by
term.  For the mutation, the complete discrepancy is the rank-one diagonal
matrix $\operatorname{diag}(-1/7,0)$, whose squared Hilbert--Schmidt norm is
$1/49$.
\end{proof}

\begin{firewall}[Conditional normalization is a later source-stock question]
\label{fire:V11-no-rare-cell-division}
The one-letter reflected form uses $M_a$, not $p_a^{-1}M_a$.  Dividing by a
small $p_a$ is required only when a normalized conditional law, conditional
source metric, or historywise reserve is itself claimed.  A large conditional
floor on a vanishing-probability history is not a noncollapsing unnormalized
physical source stock; conversely, a positive $M_a$ is a valid OS contribution
without a separate probability floor.
\end{firewall}

% =============================================================================
\section{Record transport across cutoffs}
\label{sec:V11-record-transport}
% =============================================================================

Let cutoff $n$ have a readable finite record $\mathcal R_n$, unnormalized cell
matrices $M_{n,a}$, and common one-letter space $F_n$.  Let
\begin{equation}
 \pi_n:\mathcal R_{n+1}\longrightarrow\mathcal R_n,
 \qquad
 J_n:F_n\longrightarrow F_{n+1}
 \label{eq:V11-record-cutoff-map}
\end{equation}
be the physically frozen record and coefficient transports.  Define the
adjacent record rectangle defect
\begin{equation}
 \boxed{
 \varepsilon_n^{\rm rec}
 =\sum_{a\in\mathcal R_n}
 \left\|M_{n,a}
 -\sum_{\pi_n(b)=a}J_n^*M_{n+1,b}J_n
 \right\|_{\rm op}.}
 \label{eq:V11-record-rectangle}
\end{equation}
The sum norm is chosen so that refinement multiplicity is paid once in the
complete cell partition.

\begin{theorem}[Target-native cofinal transport of a readable record]
\label{thm:V11-record-cofinal}
Assume the $J_n$ compose isometrically and
\begin{equation}
 \sum_{n=0}^{\infty}\varepsilon_n^{\rm rec}<\infty.
 \label{eq:V11-summable-record-rectangle}
\end{equation}
Fix a cutoff $m$ and one cell $a\in\mathcal R_m$.  For $n>m$, aggregate every
$n$-cell descending to $a$ and transport it back to $F_m$:
\begin{equation}
 \widetilde M_{n\downarrow m,a}
 =\sum_{\pi_{n,m}(b)=a}
   J_{n,m}^*M_{n,b}J_{n,m}.
 \label{eq:V11-descendant-aggregate}
\end{equation}
Then $(\widetilde M_{n\downarrow m,a})_{n\ge m}$ is Cauchy in operator norm and
has a unique limit.  More precisely,
\begin{equation}
 \boxed{
 \|\widetilde M_{N\downarrow m,a}-M_{m,a}\|_{\rm op}
 \le\sum_{n=m}^{N-1}\varepsilon_n^{\rm rec}.}
 \label{eq:V11-record-telescope}
\end{equation}
If every fine cell matrix is positive, every exact descendant aggregate is
positive; if $M_{m,a}$ has floor $\alpha$ and the tail on the right of
\eqref{eq:V11-record-telescope} is smaller than $\alpha$, the limiting
aggregate retains the corresponding positive floor.
\end{theorem}

\begin{proof}
Expand the difference between two successive descendant aggregates.  Grouping
the $(n+1)$-cells first by their $n$-parent produces exactly the terms in
\eqref{eq:V11-record-rectangle}, compressed by the composed isometry.  The
operator norm of the contribution descending to one fixed $a$ is bounded by
the complete sum of parent defects.  Telescoping gives
\eqref{eq:V11-record-telescope}; summability gives the Cauchy property.  Sums
and isometric compressions preserve positivity.  The floor statement follows
from the operator-norm perturbation bound.
\end{proof}

\begin{assessment}[What is and is not transported]
\label{ass:V11-transport-scope}
The theorem transports one declared readable record and its one-letter
reflected matrix.  It does not transport a full bosonic path, a conditional
precision, a line connection, or the complete exterior hierarchy unless those
objects are themselves included in the target record.  This is the
record-level analogue of the target-native scalar and response hierarchy: the
stronger history is not charged merely to preserve a weaker reflected matrix.
\end{assessment}

% =============================================================================
\section{Revised order of the fermionic entrance}
\label{sec:V11-revised-order}
% =============================================================================

\begin{theorem}[Record-relative entrance before the historywise Gaussian gate]
\label{thm:V11-revised-order}
The first fermionic reflection gate is ordered by the target being claimed.
\begin{enumerate}[label=\textup{(T\arabic*)},leftmargin=3.5em]
\item For a one-letter reflected statement on an occurring readable bosonic
      record $r$, the target-minimal input is the family of unnormalized cell
      matrices $(M_a)_{a\in\mathcal R}$ and their exact word-relation support.
      \Cref{thm:V11-record-OS,thm:V11-outward-cell} gives the complete finite
      verdict.
\item For a historywise reflected hopping factor, conditional Gaussian law,
      or exterior second quantization, the full-history selector/conditioning
      row must occur.  Each history then enters the Version~V10 Grassmann,
      chronology, boundary, and Version~V9 hopping gates.
\item A positive mixed one-letter kernel cannot be promoted to conditional
      Gaussianity or to a mixed higher-word hierarchy.  One direct held-out
      higher word, the quadratic/Berezin occurrence theorem, or the complete
      higher-word bank remains necessary at that stronger target.
\item The Hamiltonian tangent, Duhamel ancestry, current/stress, charge, and
      Einstein rows remain downstream of whichever reflected packet has
      actually passed.  No coarse one-letter pass retroactively populates
      their common-history or line data.
\end{enumerate}
Equivalently, the corrected order is
\begin{equation}
 \boxed{
 \begin{gathered}
 \text{actual positive-time bosonic Read}\
 \Downarrow\\[-0.25em]
 \text{record-relative unnormalized reflected one-letter matrix}\
 \Downarrow\\[-0.25em]
 \begin{matrix}
 \text{record-level OS verdict},&
 \text{or full-history refinement if conditional Gaussianity is claimed}
 \end{matrix}\\
 \Downarrow\\[-0.25em]
 \text{V10 boundary-complete germ and V9 hopping only on that stronger branch}\\
 \Downarrow\\[-0.25em]
 \text{exterior, Hamiltonian, Duhamel, locality, current/stress, and Einstein gates}.
 \end{gathered}}
 \label{eq:V11-revised-order}
\end{equation}
\end{theorem}

\begin{proof}
Item~\textup{(T1)} is the direct-sum characterization
\Cref{thm:V11-record-OS}; finite acquisition and certification are
\Cref{thm:V11-cell-polarization,thm:V11-outward-cell}.  The hidden-negativity
and hidden-relation examples show that neither a coarse pass nor the mixed
matrix determines the full-history branch.  Conditional Gaussianity and
exterior second quantization use the pointwise quadratic/Berezin structure of
Version~V10 and hence require item~\textup{(T2)}.  The existing positive
bosonic-mixture theorem supplies positivity after conditional passes but also
shows that mixing generally loses quasi-free closure, proving
item~\textup{(T3)}.  The final item is the inherited terminal quantum audit
order.
\end{proof}

\begin{firewall}[V11 qualifies the V10 active card; it does not delete it]
\label{fire:V11-V10-scope}
The Version~V10 open problem remains the correct next calculation whenever one
named conditional history, its hopping factor, or its exterior Gaussian
kernel is the intended output.  Version~V11 adds an earlier question: is that
history actually the positive-time record named by the target?  On a coarser
record the direct cell matrix is the correct minimum.  On the full record the
cell theorem reduces exactly to the V10 historywise card.  No previous proof is
rewritten.
\end{firewall}

% =============================================================================
\section{Exact replay, integrated theorem, and the new physical stopping line}
\label{sec:V11-replay}
% =============================================================================

\begin{proposition}[Exact record-relative replay]
\label{prop:V11-replay}
The accompanying standard-library verifier checks, in exact rational or
Gaussian-rational arithmetic:
\begin{enumerate}[label=\textup{(R\arabic*)},leftmargin=3.4em]
\item the two history matrices in
      \eqref{eq:V11-hidden-negative-matrices}, their positive coarse mean, and
      the weighted full-history negative-part stock one;
\item the matrix-martingale Pythagoras identity
      \begin{equation}
       \boxed{10=2+8;}
       \label{eq:V11-replay-Pythagoras}
      \end{equation}
\item the two relation failures in
      \eqref{eq:V11-hidden-relation-residual} and their exact cancellation on
      the constant record;
\item polarization of
      \begin{equation}
       \begin{pmatrix}2&1+2\mathrm i\\1-2\mathrm i&5\end{pmatrix}
       \label{eq:V11-replay-Hermitian}
      \end{equation}
      from the six scalar values
      $2,5,9,5,3,11$;
\item exact coarse/fine record additivity and rejection of the $1/7$ mutation
      with residual $1/49$.
\end{enumerate}
Normal and optimized Python executions produce byte-identical JSON output.
No displayed matrix is an RPESM measurement.
\end{proposition}

\begin{proof}
Every statement is a finite substitution into the formulas above.  The
verifier represents complex rationals as ordered pairs of
\texttt{Fraction} objects, implements matrix products and Hilbert--Schmidt
norms directly, and asserts every equality before writing the certificate.
\end{proof}

\begin{definition}[First active Version V11 physical card]
\label{def:V11-active-card}
The next target-minimal packet is
\begin{equation}
 \boxed{
 \begin{gathered}
 \texttt{POSITIVE\_TIME\_BOSONIC\_READABILITY}\\[-0.15em]
 \texttt{\_AND\_RECORD\_RELATIVE\_REFLECTED}\\[-0.15em]
 \texttt{\_ONE\_LETTER\_KERNEL\_CARD}.
 \end{gathered}}
 \label{eq:V11-active-card-name}
\end{equation}
Its minimum payload is:
\begin{enumerate}[label=\textup{(A\arabic*)},leftmargin=3.5em]
\item one common positive unnormalized RPESM field-shell cylinder;
\item the actual bosonic record whose cell selectors occur in the
      positive-time observable algebra, fixed before any reflected sign is
      read;
\item one common transported one-letter coefficient screen and its exact word
      relation support;
\item the selector-weighted quadratic Reads in
      \eqref{eq:V11-cell-quadratic-read}, with outward radii, for every cell
      used by the target;
\item one independently retained refinement/additivity check
      \eqref{eq:V11-record-additivity-residual};
\item a declaration whether the target is the mixed one-letter kernel, a
      record-conditioned one-letter kernel, or the full historywise Gaussian
      mechanism;
\item on the full-history target, the complete Version~V10 chronology,
      Grassmann germ, boundary/divisor, and Version~V9 hopping packet;
\item one record rectangle \eqref{eq:V11-record-rectangle} when a cofinal
      statement is sought.
\end{enumerate}
\end{definition}

\begin{theorem}[Version V11 readable-record and hidden-history theorem]
\label{thm:V11-integrated}
Preserving every theorem and limitation of Versions~V1--V10, Version~V11
proves:
\begin{enumerate}[label=\textup{(V11.\arabic*)},leftmargin=5.8em]
\item the physically relevant one-letter reflected matrix is conditioned only
      on the bosonic record actually readable in the positive-time algebra;
\item its complete form is the direct sum of the unnormalized cell matrices
      $M_a$, and positivity is equivalent to their individual positivity;
\item the constant, intermediate, and full-history records give genuinely
      different targets, related by exact matrix additivity under refinement;
\item integrated positivity may hide a negative conditional direction on
      every history;
\item a coarse relation-descended hopping pass may hide nonzero couplings to a
      represented zero word on every history;
\item the squared distance to the supported positive cone is monotone under
      record refinement;
\item the conditional matrices form an exact Hilbert--Schmidt martingale whose
      record innovations obey two Pythagoras identities;
\item a strict coarse floor survives to the full history under a summable
      operator-norm innovation budget smaller than that floor;
\item selector-weighted polarization reconstructs every cell matrix directly,
      outward scalar radii give an exact finite matrix radius, and one held-out
      refinement tests same-cylinder incidence;
\item a summable adjacent record rectangle transports each fixed readable cell
      without transporting an unnecessarily stronger history object.
\end{enumerate}
No item proves that the active RPESM full history is selectable in the
positive-time algebra, populates its record-cell matrices, or supplies its
historywise Grassmann, hopping, line, exterior, Duhamel, current/stress, charge,
continuum, or Einstein values.
\end{theorem}

\begin{proof}
Items~\textup{(V11.1)--(V11.3)} are
\Cref{thm:V11-record-OS,cor:V11-record-hierarchy}.  Items
\textup{(V11.4)--(V11.5)} are
\Cref{cth:V11-hidden-negativity,cth:V11-hidden-relation}.  Item
\textup{(V11.6)} is \Cref{thm:V11-cone-monotone}.  Items
\textup{(V11.7)--(V11.8)} are
\Cref{thm:V11-martingale-Pythagoras,thm:V11-floor-survival}.  Item
\textup{(V11.9)} is
\Cref{thm:V11-cell-polarization,thm:V11-outward-cell,prop:V11-record-incidence},
and the final item is \Cref{thm:V11-record-cofinal}.  The excluded outputs are
exactly the unpopulated rows in \Cref{def:V11-active-card} and the downstream
terminal quantum order.
\end{proof}

\begin{assessment}[Status after Version V11]
\label{ass:V11-status}
\passstatus{} for the record-relative direct-sum theorem, exact positivity
criterion, refinement identity, cone-obstruction monotonicity, matrix
martingale Pythagoras, summable-floor bridge, selector-weighted polarization,
outward finite verdict, record-incidence held-out, cofinal rectangle, and exact
replay.

\obsstatus{} for treating a global provenance history as an automatically
readable positive-time selector; for inferring a conditional hopping or
Gaussian branch from a positive integrated kernel; for allowing zero-word,
skew, or negative conditional defects to cancel before a stronger historywise
claim; and for normalizing a rare record cell before its unnormalized source
stock and probability floor have been separated.

\stopstatus{} on the active Einstein--Standard--Model regulator at the first
unpopulated tuple
\begin{equation}
 \boxed{
 (\Omega_B,\Lambda_B,r_+,F,
   \{Q_a(z)\}_{a,z},P_{\rm word},
   \Delta_{s/r}^{\rm rec})}
 \label{eq:V11-physical-stop}
\end{equation}
where $r_+$ is the actual positive-time readable bosonic record.  If $r_+$ is
the full history, the programme immediately resumes Version~V10.  If it is
coarser, the first physical verdict is the record-relative matrix of this
version and no conditional Gaussian statement is inferred.
\end{assessment}

\begin{openproblem}[Version V12: populate the readable record before the conditional slab]
\label{op:V11-next}
Preserve Versions~V1--V11 verbatim.  Do not add another source, localizer,
exterior, lag, concentration, or continuum compiler.  Proceed only in this
order.
\begin{enumerate}[label=\textup{(V12.\arabic*)},leftmargin=5.8em]
\item On one active RPESM cutoff, list the bosonic marks which are actually
      admitted as positive-time Reads and form their coarsest declared record
      $r_+$.  Do not replace it by the full provenance ledger unless singleton
      selectors occur.
\item On the frozen one-letter screen, acquire the unnormalized cell quadratic
      matrices and the exact word support.  Run the finite cone verdict before
      any conditional history is selected.
\item Retain one independent refinement.  A positive
      $\Delta_{s/r}^{\rm rec}$ is the first same-cylinder incidence stop; a
      vanishing residual licenses the finer target only if those fine selectors
      are themselves physical Reads.
\item If the desired target is only the mixed or $r_+$-conditioned one-letter
      OS form, stop after the separated cell verdict and its cutoff rectangle.
\item If a historywise hopping, Gaussian, line, or exterior statement is
      required and the full history is readable, apply Version~V10 to every
      selected cell, including the fixed-plane divisor branch, and then
      Version~V9.
\item Only after that stronger packet passes may V8/V5/V6 be used for the
      common Hamiltonian and Duhamel ancestry.
\end{enumerate}
If the positive-time selector record or its mixed quadratic Reads are not
physically supplied, the exact endpoint is
\eqref{eq:V11-physical-stop}.  Further historywise algebra without that record
would be a refinement of an unobserved variable rather than an advance on the
active quantum target.
\end{openproblem}

\section*{V11 verification and append-only record}
\addcontentsline{toc}{section}{V11 verification and append-only record}
The accompanying verifier uses only Python standard-library exact rational
arithmetic.  It verifies the hidden conditional negative matrices, hidden
zero-word relation, matrix-martingale Pythagoras, complex Hermitian
polarization, record additivity, and the held-out mutation.

The complete Version~V10 source preceding its sole final document-closing
command is preserved byte for byte.  Its preterminal SHA--256 digest is
\begin{center}\small\ttfamily
ced1fcffa98ed527b31af7804a9a7f5a938a0544613c931e88de8dd2df728207
\end{center}
The exact verifier and its JSON output have SHA--256 digests
\begin{center}\small\ttfamily
0a32ed75da93e28b58ad0dc0fa03f712c08baa7472610df24b29f3075113c5ca\\
67b165613131159507dd358aa7422dba12162c01d3a9422ecdfc70183c5d4204
\end{center}
A later continuation must remove only the sole final
\verb|\end{document}|, append new material immediately before it, restore one
terminal command, and increment the filename to end in \texttt{\_V12.tex}.

% =============================================================================
% END APPENDED VERSION V11 MATERIAL
% =============================================================================


% =============================================================================
% START APPENDED VERSION V12 MATERIAL -- complete V1--V11 preterminal source preserved
% =============================================================================

\clearpage
\hypersetup{
 pdftitle={Renewal Einstein--Standard--Model Physical Source Metric, Duhamel Ancestry and Quantum Continuum Programme V12},
 pdfsubject={The reflection-fixed boundary record, Markov entrance Gram, and the correction of future-cell sectorization}
}
\makeatletter
\immediate\write\@auxout{\detokenize{\gdef\VtwelveTOCwide{\makeatletter\renewcommand*\l@section[2]{\ifnum\c@tocdepth>\z@\addpenalty\@secpenalty\addvspace{1.0em \@plus\p@}\setlength\@tempdima{2.1em}\begingroup\parindent\z@\rightskip\@pnumwidth\parfillskip-\@pnumwidth\leavevmode\bfseries\advance\leftskip\@tempdima\hskip-\leftskip#1\nobreak\hfil\nobreak\hb@xt@\@pnumwidth{\hss#2\kern-\p@\kern\p@}\par\endgroup\fi}\makeatother}}}
\makeatother
\addtocontents{toc}{\protect\VtwelveTOCwide}
\section*{Version V12: the reflection-fixed boundary record and the Markov entrance Gram}
\addcontentsline{toc}{section}{V12: boundary record and Markov entrance Gram}

\begin{abstract}
Version~V11 moved correctly to the physical Read which precedes a
record-conditioned Osterwalder--Schrader statement, but it used one phrase,
``positive-time readable record'', for two different roles.  A record written
at strictly positive time is an admissible positive-time \emph{writer}.  It is
not thereby a reflection-fixed same-cut selector whose cells form orthogonal
Osterwalder--Schrader sectors.  The latter role belongs to the boundary record
at the reflection cut.  The distinction is already present in the inherited
RPESM slab: the full event tuple is point readable, while the reflected Markov
square is obtained by conditioning past and future on the time-zero state.

This continuation goes to that earlier incidence.  For a stationary
reflection-symmetric Markov slab, every finite positive-time writer bank has a
boundary entrance synthesis.  A physically readable boundary record produces
one positive operator-valued measure and an exact direct-integral
Osterwalder--Schrader factorization.  This formulation works for the continuous
compact-gauge/unbounded-Higgs record of the RPESM cylinder and requires no
singleton atoms or rare-cell normalization.

A finite future-segment mark has a separate, target-minimal eligibility test.
The Osterwalder--Schrader Gram of its vacuum-grade cell indicators is the Gram
of their boundary entrance probabilities.  Its deficit
\[
 1-\operatorname{Tr}C
\]
is exactly the conditional one-hot variance left after the boundary state is
known.  It vanishes if and only if the mark is already a deterministic
function of the reflection-fixed boundary record.  Only on that zero branch
does the complete selected-writer Gram become block diagonal and reduce to the
cell direct sum used in Version~V11.  On the positive branch the cell writers
have nonzero cross blocks, and their complete coherent Gram must be assembled
before any relation short or positivity test.

An exact two-arm Markov cylinder gives the separator.  Its boundary record has
zero eligibility defect, while a future bit has defect $3/8$.  The two future
cell indicators have off-diagonal Osterwalder--Schrader overlap $3/16$; deleting
the cross blocks loses exactly $3/8$ of the vacuum norm.  For a correlated
writer, future-cell selection also creates an explicit conditional-covariance
entrance which cannot be recovered by multiplying the unselected boundary
amplitude by the cell probability.

The source audit consequently changes the active status.  The existential
RPESM regulator already supplies the common event tuple, the time-zero Markov
state, the physical boundary record, and all finite same-history rows.  The
missing work is no longer existence of a readable record.  It is the literal
extraction and evaluation of the first boundary-complete time-split crossing
card, together with the cofinal source-energy and line-transport margins.  No
new record, exterior, lag, or concentration compiler is opened.
\end{abstract}

% =============================================================================
\section{The RPESM record occurs, but its three roles are different}
\label{sec:V12-three-record-roles}
% =============================================================================

The inherited RPESM construction has two pieces which must be used together.
First, one physical event is a point-readable tuple containing the relational
state, fast scheduler path, retained route, initial and terminal fields,
update and reverse labels, duration, boundary route, source and Green marks,
current and stress responses, determinant transport, Pfaffian sign, and
calibration data.  Second, reflection positivity of the stationary slab is
proved by conditioning the past and future on the time-zero Markov state.
Thus occurrence of the global event tuple and centrality at the reflection cut
are not the same statement.

Let $(\Omega,\mathscr H,\mathbb P)$ be one finite-cutoff RPESM path cylinder,
let $X_0$ be its state at the reflection cut, and let $\Theta$ be the declared
reflection.  We distinguish:
\begin{equation}
 \boxed{
 \begin{aligned}
 \mathscr H_{\rm prov}
   &:=\sigma(\text{the complete point-readable event tuple}),\\
 \mathscr B_0
   &:=\sigma(X_0,\text{the declared reflection-fixed boundary marks}),\\
 \mathscr F_+
   &:=\sigma(\text{the strictly positive-time path and its future marks}).
 \end{aligned}}
 \label{eq:V12-three-record-algebras}
\end{equation}
The first sigma-algebra is the same-history provenance used to assemble one
common action and to disintegrate its conditional Grassmann germ.  The second
is the same-cut boundary algebra which may supply central
Osterwalder--Schrader selectors.  The third supplies positive-time writers.
A random variable may be point readable in $\mathscr H_{\rm prov}$ and belong
to $\mathscr F_+$ without belonging to $\mathscr B_0$.

\begin{theorem}[Three record roles in a reflected Markov cylinder]
\label{thm:V12-three-record-roles}
On a standard Borel reflected Markov cylinder the following statements are
logically distinct.
\begin{enumerate}[label=\textup{(R\arabic*)},leftmargin=3.6em]
\item A common-history conditional coefficient is a regular conditional
      disintegration with respect to $\mathscr H_{\rm prov}$.  It does not
      require positive-probability singleton histories.
\item A direct-sum Osterwalder--Schrader sector label is a finite or Borel
      record measurable with respect to $\mathscr B_0$.  Its selectors are
      fixed by reflection and multiply both reflected arms at the same cut.
\item A future mark in $\mathscr F_+$ is an admissible writer.  Unless it is
      also $\mathscr B_0$-measurable, selecting its values changes the Markov
      entrance amplitude and does not produce orthogonal central sectors.
\end{enumerate}
Consequently the full provenance record may be used to define a measurable
family of Version~V10 Grassmann germs, while only its reflection-fixed boundary
quotient may be inserted into the Version~V11 direct-sum form.
\end{theorem}

\begin{proof}
The finite-cutoff RPESM field space is a finite product of compact Lie-group
coordinates, Euclidean Higgs and auxiliary coordinates, finite marks, and a
cadlag path cylinder.  It is standard Borel, so every measurable subrecord has
a regular conditional probability.  This proves the first item and shows why
an atomic selector is unnecessary for conditional disintegration.

A same-cut selector $h\in L^\infty(\mathscr B_0)$ satisfies
$h\circ\Theta=h$.  Multiplication by $h$ therefore acts on the two reflected
arms through the same boundary value.  Disjoint boundary cells have disjoint
multipliers, giving central orthogonal sectors.  This is the second item.

For a future mark $A$, multiplication by a cell indicator occurs before the
conditional expectation onto $X_0$.  In general
$\mathbb E[\one_{\{A=a\}}F\mid X_0]$ is not
$\mathbb E[\one_{\{A=a\}}\mid X_0]\mathbb E[F\mid X_0]$.  The exact difference
is computed in \Cref{thm:V12-cell-word-Gram}.  Hence future cells are writers,
not central sectors, unless the indicators are already boundary measurable.
\end{proof}

\begin{firewall}[Global conditioning is not a positive-time singleton Read]
\label{fire:V12-disintegration-not-selector}
Conditioning the common cylinder on its global provenance is a mathematical
same-history disintegration.  Inserting a history selector into the
Osterwalder--Schrader algebra is a stronger operational statement.  On the
continuous RPESM field space most singleton histories have probability zero,
so Version~V11's singleton language is only a finite-atomic shorthand.  The
correct objects are the boundary sigma-algebra and the operator-valued measure
constructed below.
\end{firewall}

% =============================================================================
\section{The boundary entrance synthesis and its positive operator-valued measure}
\label{sec:V12-boundary-entrance}
% =============================================================================

Let $(X_t)_{t\in\mathbb R}$ be a stationary reversible Markov process with
invariant probability $\mu$ on a standard Borel state space $S$.  Let
$F_1,\ldots,F_d\in L^2(\mathscr F_+)$ be a finite positive-time writer bank and
put
\begin{equation}
 F_c=\sum_{j=1}^dc_jF_j,
 \qquad c\in\mathbb C^d.
 \label{eq:V12-writer-bank}
\end{equation}
Assume the declared reflection-pair convention, so that conditional
independence of past and future at $X_0$ gives
\begin{equation}
 \mathbb E\!\left[\overline{F_c(\Theta\omega)}F_e(\omega)\right]
 =\int_S\overline{b_c(x)}b_e(x)\,\dd\mu(x),
 \qquad
 b_c(x):=\mathbb E[F_c\mid X_0=x].
 \label{eq:V12-Markov-polarized}
\end{equation}
Equation~\eqref{eq:V12-Markov-polarized} is the polarized form of the inherited
Markov square.  Write
\begin{equation}
 B(x):\mathbb C^d\longrightarrow\mathbb C,
 \qquad B(x)c=b_c(x),
 \qquad K(x)=B(x)^*B(x)\succeq0.
 \label{eq:V12-entrance-K}
\end{equation}
The same statements hold for a fixed finite-dimensional boundary Hilbert
space after replacing the scalar $B(x)$ by a measurable operator-valued
entrance synthesis.

Let
\begin{equation}
 r_0:S\longrightarrow\mathsf R_0
 \label{eq:V12-boundary-record-map}
\end{equation}
be a physically declared reflection-fixed boundary record, and put
$\nu_0=(r_0)_*\mu$.

\begin{theorem}[Boundary operator-valued measure and direct-integral OS factorization]
\label{thm:V12-boundary-POVM}
The formula
\begin{equation}
 \boxed{
 \mathbf M_{r_0}(E)
 =\int_{\{x:r_0(x)\in E\}}K(x)\,\dd\mu(x),
 \qquad E\subseteq\mathsf R_0\text{ Borel},}
 \label{eq:V12-boundary-POVM}
\end{equation}
defines a finite positive operator-valued measure on $\mathsf R_0$.  It has a
$\nu_0$-almost-everywhere defined Radon--Nikodym density
\begin{equation}
 \boxed{
 K_{r_0}(a)
 =\mathbb E[K(X_0)\mid r_0=a]\succeq0.}
 \label{eq:V12-boundary-density}
\end{equation}
For every bounded simple coefficient field
$c:\mathsf R_0\to\mathbb C^d$, the corresponding boundary-selected writer
satisfies
\begin{equation}
 \boxed{
 \begin{aligned}
 \mathfrak q_{r_0}(c)
 &:=\mathbb E\!\left[
  \overline{F_{c(r_0(X_0))}(\Theta\omega)}
  F_{c(r_0(X_0))}(\omega)\right]\\
 &=\int_S\|B(x)c(r_0(x))\|^2\,\dd\mu(x)\\
 &=\int_{\mathsf R_0}
   \langle c(a),K_{r_0}(a)c(a)\rangle\,\dd\nu_0(a)
 \ge0.
 \end{aligned}}
 \label{eq:V12-direct-integral-form}
\end{equation}
Its null space is exactly
\begin{equation}
 \boxed{
 \mathfrak q_{r_0}(c)=0
 \iff B(x)c(r_0(x))=0\quad\mu\text{-almost everywhere}.}
 \label{eq:V12-direct-integral-null}
\end{equation}
Thus the source-minimal record-relative one-letter OS carrier is the closure of
the range of
\begin{equation}
 \mathcal J_{r_0}c(x)=B(x)c(r_0(x)),
 \label{eq:V12-boundary-factor}
\end{equation}
not an arbitrary direct sum over provenance histories.
\end{theorem}

\begin{proof}
Every matrix entry of $K$ belongs to $L^1(\mu)$ by Cauchy--Schwarz, so
\eqref{eq:V12-boundary-POVM} is a countably additive finite matrix-valued
measure.  For every $v\in\mathbb C^d$,
\[
 v^*\mathbf M_{r_0}(E)v
 =\int_{r_0^{-1}(E)}\|B(x)v\|^2\,\dd\mu(x)\ge0,
\]
which proves positivity.  Componentwise Radon--Nikodym differentiation with
respect to $\nu_0$ gives \eqref{eq:V12-boundary-density}.  Positivity of that
density follows by testing a countable dense set of rational complex vectors
and intersecting the corresponding full-measure sets.

Because $c(r_0(X_0))$ is $X_0$-measurable, it may be pulled through the
conditional expectation defining $B$.  Apply
\eqref{eq:V12-Markov-polarized} and then disintegrate over $r_0$ to obtain
\eqref{eq:V12-direct-integral-form}.  The integrand is nonnegative, so its
integral vanishes exactly when it vanishes almost everywhere, proving
\eqref{eq:V12-direct-integral-null}.  The range statement is the ordinary Gram
factorization by $\mathcal J_{r_0}$.
\end{proof}

\begin{corollary}[Finite boundary cells recover the valid part of Version V11]
\label{cor:V12-finite-boundary-cells}
If $r_0$ has finitely many values, put
\begin{equation}
 M_a=\mathbf M_{r_0}(\{a\}).
 \label{eq:V12-finite-boundary-Ma}
\end{equation}
Then
\begin{equation}
 \boxed{
 \mathfrak q_{r_0}((c_a)_a)
 =\sum_a\langle c_a,M_ac_a\rangle,
 \qquad M_a\succeq0.}
 \label{eq:V12-valid-V11-direct-sum}
\end{equation}
Thus the Version~V11 cell direct sum is exact when its record is strengthened
from an arbitrary positive-time mark to a reflection-fixed boundary record.
For a continuous record, \eqref{eq:V12-direct-integral-form} is the correct
replacement; no singleton cell probability is required.
\end{corollary}

\begin{proof}
Apply \Cref{thm:V12-boundary-POVM} to the finite measurable partition into
record cells.  Distinct cell indicators have disjoint support as functions of
$X_0$, so the off-diagonal sector terms vanish.
\end{proof}

% =============================================================================
\section{A vacuum-grade test for whether a future mark is a boundary selector}
\label{sec:V12-boundary-eligibility}
% =============================================================================

Let $A_+$ be a finite-valued positive-time mark with values
$1,\ldots,m$, let $A_-=A_+\circ\Theta$ be its reflected past copy, and put
\begin{equation}
 \chi_a^+=\one_{\{A_+=a\}},
 \qquad
 \chi_a^-=\one_{\{A_-=a\}},
 \qquad
 p_a(x)=\mathbb P(A_+=a\mid X_0=x).
 \label{eq:V12-future-cells}
\end{equation}
Reversibility gives the same conditional probabilities for $A_-$, and the
Markov property makes the two arms conditionally independent given $X_0$.
Define the vacuum-grade selector entrance Gram
\begin{equation}
 \boxed{
 C^A_{ab}
 :=\mathbb E[\chi_a^-\chi_b^+]
 =\int_Sp_a(x)p_b(x)\,\dd\mu(x).}
 \label{eq:V12-selector-Gram}
\end{equation}
Let
\begin{equation}
 \pi_a=\mathbb P(A_+=a),
 \qquad D_\pi=\operatorname{diag}(\pi_1,\ldots,\pi_m).
 \label{eq:V12-cell-probabilities}
\end{equation}

\begin{theorem}[Boundary-determinism residual]
\label{thm:V12-boundary-determinism}
The matrix
\begin{equation}
 \boxed{
 V^A:=D_\pi-C^A
 =\mathbb E\bigl[(e_{A_+}-p(X_0))(e_{A_+}-p(X_0))^{\mathsf T}\bigr]
 \succeq0}
 \label{eq:V12-record-innovation-Gram}
\end{equation}
is the exact one-hot record innovation left after the boundary state is known.
Its total stock is
\begin{equation}
 \boxed{
 \begin{aligned}
 \Delta_{\partial}(A)
 &:=\operatorname{Tr}V^A\\
 &=1-\operatorname{Tr}C^A\\
 &=\sum_{a=1}^m\int_Sp_a(1-p_a)\,\dd\mu\\
 &=\sum_{a\ne b}C^A_{ab}.
 \end{aligned}}
 \label{eq:V12-boundary-defect}
\end{equation}
The following are equivalent:
\begin{enumerate}[label=\textup{(B\arabic*)},leftmargin=3.6em]
\item $\Delta_{\partial}(A)=0$;
\item $V^A=0$;
\item every $\chi_a^+$ is $\sigma(X_0)$-measurable;
\item $A_+$ is a deterministic function of $X_0$ almost surely;
\item $C^A=D_\pi$, equivalently the vacuum-grade cell writers are mutually
      orthogonal in the Osterwalder--Schrader form.
\end{enumerate}
When the equivalent conditions fail,
$\operatorname{rank}V^A$ is the dimension of the source-minimal linear
future-record innovation beyond the boundary state.  This rank is a record
cost, not a new physical field source.
\end{theorem}

\begin{proof}
Since $\mathbb E[e_{A_+}\mid X_0]=p(X_0)$, conditional-expectation
Pythagoras gives \eqref{eq:V12-record-innovation-Gram}.  Direct calculation
also gives
\[
 \mathbb E[(\chi_a^+-p_a)(\chi_b^+-p_b)]
 =\delta_{ab}\pi_a-C^A_{ab}.
\]
Taking traces yields the first three expressions in
\eqref{eq:V12-boundary-defect}.  Because $\sum_ap_a=1$,
\[
 \sum_{a,b}C^A_{ab}
 =\int\left(\sum_ap_a\right)^2\dd\mu=1,
\]
which gives the last expression.

If $\Delta_{\partial}(A)=0$, every nonnegative integrand
$p_a(1-p_a)$ vanishes almost everywhere, so each $p_a$ takes only the values
zero and one.  Moreover
\[
 \mathbb E|\chi_a^+-p_a(X_0)|^2
 =\mathbb E[p_a(1-p_a)]=0,
\]
so $\chi_a^+=p_a(X_0)$ almost surely.  Thus $A_+$ is determined by $X_0$.
The converse is immediate.  The Gram is diagonal exactly when its sum of
off-diagonal entries, namely $\Delta_{\partial}(A)$, vanishes.  Finally,
$V^A$ is the Gram of the centered one-hot residual bank, so its rank is the
dimension of its linear span and hence the minimum Hilbert dimension of that
innovation.
\end{proof}

\begin{corollary}[Target-minimal eligibility acquisition]
\label{cor:V12-eligibility-acquisition}
To decide whether a finite future mark may be promoted to a boundary-sector
label, it is enough to acquire the single $m\times m$ positive matrix $C^A$ in
\eqref{eq:V12-selector-Gram}.  No fermion precision, determinant, boundary
Schur solve, or Duhamel matrix is needed for this decision.

Suppose outward intervals satisfy
\begin{equation}
 |\widehat C^A_{aa}-C^A_{aa}|\le\eta_{aa}.
 \label{eq:V12-C-diagonal-errors}
\end{equation}
Then
\begin{equation}
 \boxed{
 \begin{aligned}
 \underline\Delta_A
 &=\max\left\{0,
  1-\sum_a(\widehat C^A_{aa}+\eta_{aa})\right\},\\
 \overline\Delta_A
 &=\min\left\{1,
  1-\sum_a(\widehat C^A_{aa}-\eta_{aa})\right\}
 \end{aligned}}
 \label{eq:V12-Delta-interval}
\end{equation}
contains $\Delta_{\partial}(A)$.  Hence
$\underline\Delta_A>0$ proves that $A$ is a genuine future mark and not a
boundary selector.  Exact promotion to a boundary record requires either an
exact decoder $A=d(X_0)$ or an exact zero certificate; an interval merely
meeting zero is unresolved.
\end{corollary}

\begin{proof}
The identity $\Delta_{\partial}(A)=1-\sum_aC^A_{aa}$ and interval arithmetic
give \eqref{eq:V12-Delta-interval}.  The logical conclusions are
\Cref{thm:V12-boundary-determinism}.
\end{proof}

% =============================================================================
\section{Future-cell writers form a coherent Gram, not a direct sum}
\label{sec:V12-future-cell-Gram}
% =============================================================================

Retain the writer bank $F_1,\ldots,F_d$.  For every future cell and writer put
\begin{equation}
 b_{a,j}(x)
 :=\mathbb E[\chi_a^+F_j\mid X_0=x],
 \qquad
 b_j(x):=\mathbb E[F_j\mid X_0=x].
 \label{eq:V12-selected-entrances}
\end{equation}
Define the complete cell--writer entrance Gram by
\begin{equation}
 \boxed{
 \mathbb G^A_{(a,i),(b,j)}
 =\int_S\overline{b_{a,i}(x)}b_{b,j}(x)\,\dd\mu(x).}
 \label{eq:V12-complete-cell-word-Gram}
\end{equation}

\begin{theorem}[Exact coherent cell--writer assembly]
\label{thm:V12-cell-word-Gram}
For every coefficient array $v=(v_{a,j})$,
\begin{equation}
 \boxed{
 \mathbb E\!\left[
  \overline{\Theta\!\left(\sum_{a,j}v_{a,j}\chi_a^+F_j\right)}
  \left(\sum_{b,k}v_{b,k}\chi_b^+F_k\right)
 \right]
 =v^*\mathbb G^Av\ge0.}
 \label{eq:V12-cell-word-OS}
\end{equation}
Moreover:
\begin{enumerate}[label=\textup{(C\arabic*)},leftmargin=3.6em]
\item The unselected entrance and Gram are recovered by complete assembly:
      \begin{equation}
       b_j=\sum_ab_{a,j},
       \qquad
       G_{ij}=\sum_{a,b}\mathbb G^A_{(a,i),(b,j)}.
       \label{eq:V12-cell-assembly}
      \end{equation}
\item Put $\widetilde F_j=F_j-b_j(X_0)$ and
      $\eta_a=\chi_a^+-p_a(X_0)$.  Then
      \begin{equation}
       \boxed{
       b_{a,j}=p_ab_j+c_{a,j},
       \qquad
       c_{a,j}
       =\mathbb E[\eta_a\widetilde F_j\mid X_0],
       \qquad
       \sum_ac_{a,j}=0.}
       \label{eq:V12-cell-incidence-correction}
      \end{equation}
      The second term is the exact future-record/writer incidence missing from
      probability-weighted boundary reweighting.
\item If $\Delta_{\partial}(A)=0$, then
      $b_{a,j}=\chi_a^+b_j$ and
      \begin{equation}
       \boxed{
       \mathbb G^A_{(a,i),(b,j)}=0
       \quad(a\ne b).}
       \label{eq:V12-cell-block-diagonal}
      \end{equation}
      The complete Gram reduces to the Version~V11 direct sum.
\item Conversely, if the constant vacuum writer $F_0=1$ is included and its
      cell sub-Gram is block diagonal, then
      $\Delta_{\partial}(A)=0$.  Thus the vacuum grade alone decides whether
      direct-sum sectorization is legitimate.
\end{enumerate}
On the positive-defect branch, separately positive diagonal cell matrices do
not determine \eqref{eq:V12-cell-word-OS}; every cross-cell block must be
assembled before relation descent, shorting, or a cone verdict.
\end{theorem}

\begin{proof}
Condition the reflected and future arms on $X_0$.  The Markov property makes
them conditionally independent, while reversibility identifies the reflected
entrance convention.  The resulting product of the selected entrance
amplitudes is exactly \eqref{eq:V12-complete-cell-word-Gram}, proving
\eqref{eq:V12-cell-word-OS}.

The first identity in \eqref{eq:V12-cell-assembly} follows from
$\sum_a\chi_a^+=1$; taking Grams gives the second.  To prove
\eqref{eq:V12-cell-incidence-correction}, write
\[
 \chi_a^+F_j=(p_a+\eta_a)(b_j+\widetilde F_j)
\]
and condition on $X_0$.  The terms $p_a\mathbb E[\widetilde F_j\mid X_0]$ and
$b_j\mathbb E[\eta_a\mid X_0]$ vanish.  Summing over $a$ uses
$\sum_a\eta_a=0$.

On the zero-defect branch, \Cref{thm:V12-boundary-determinism} gives
$\chi_a^+=p_a(X_0)$.  Pulling this indicator through conditional expectation
gives $b_{a,j}=\chi_a^+b_j$, and disjoint cells prove
\eqref{eq:V12-cell-block-diagonal}.  Conversely, the vacuum-grade selected
entrances are exactly $p_a$.  Their off-diagonal Gram entries are
$C^A_{ab}$.  Block diagonality therefore gives
$C^A_{ab}=0$ for $a\ne b$, which is equivalent to zero defect by
\eqref{eq:V12-boundary-defect}.
\end{proof}

\begin{firewall}[Do not delete cross-cell coherence]
\label{fire:V12-no-cell-dephasing}
Replacing $\mathbb G^A$ by its block diagonal is a physical dephasing of the
future record.  It is exact only after the boundary-determinism residual has
vanished.  Positivity of the dephased matrix cannot justify that replacement:
both the complete and dephased matrices may be positive while representing
different Osterwalder--Schrader forms and different source ancestry.
\end{firewall}

% =============================================================================
\section{An exact two-arm separator}
\label{sec:V12-exact-separator}
% =============================================================================

Let the boundary state $X$ be uniform on $\{0,1\}$.  Conditional on $X$, let
the past and future bits $A_-,A_+$ be independent, with
\begin{equation}
 \mathbb P(A_+=1\mid X=0)=\frac14,
 \qquad
 \mathbb P(A_+=1\mid X=1)=\frac34,
 \label{eq:V12-example-transition}
\end{equation}
and the same law for $A_-$.  Reflection exchanges the two arms.  This is the
one-cut Markov structure used in the reflected-square argument.

\begin{proposition}[Boundary selectors versus future selectors]
\label{prop:V12-exact-separator}
For the boundary record $X$, the vacuum selector Gram and innovation are
\begin{equation}
 C^X=
 \begin{pmatrix}1/2&0\\0&1/2\end{pmatrix},
 \qquad
 V^X=0,
 \qquad
 \Delta_{\partial}(X)=0.
 \label{eq:V12-example-boundary}
\end{equation}
For the future record $A_+$,
\begin{equation}
 \boxed{
 C^A=\frac1{16}
 \begin{pmatrix}5&3\\3&5\end{pmatrix},
 \qquad
 V^A=\frac1{16}
 \begin{pmatrix}3&-3\\-3&3\end{pmatrix},
 \qquad
 \Delta_{\partial}(A)=\frac38.}
 \label{eq:V12-example-future}
\end{equation}
Thus the two future cell indicators have Osterwalder--Schrader overlap
$3/16$.  Complete assembly gives
\begin{equation}
 \sum_{a,b}C^A_{ab}=1,
 \qquad
 \sum_aC^A_{aa}=\frac58,
 \qquad
 \sum_{a\ne b}C^A_{ab}=\frac38.
 \label{eq:V12-example-coherence-loss}
\end{equation}
A false direct-sum replacement loses exactly the complete boundary-determinism
defect.

For the positive-time writer $F=2A_+-1$, its boundary entrance is
\begin{equation}
 b=(-1/2,1/2).
 \label{eq:V12-example-b}
\end{equation}
Selecting the future cell $A_+=1$ gives
\begin{equation}
 \boxed{
 b_1=(1/4,3/4)
 =p_1b+c_1,
 \qquad
 p_1b=(-1/8,3/8),
 \qquad
 c_1=(3/8,3/8).}
 \label{eq:V12-example-incidence}
\end{equation}
With the uniform boundary measure,
\begin{equation}
 \boxed{
 \|b_1\|_2^2=\frac5{16},
 \qquad
 \|p_1b\|_2^2=\frac5{64},
 \qquad
 \|b_1\|_2^2-\|p_1b\|_2^2=\frac{15}{64}.}
 \label{eq:V12-example-selected-norm}
\end{equation}
The missing $15/64$ is the sum of the signed cross term $6/64$ and the
positive incidence norm $9/64$.
\end{proposition}

\begin{proof}
The boundary indicators are deterministic functions of $X$, giving
\eqref{eq:V12-example-boundary}.  For the future cells the two boundary
entrance functions are
\[
 p_0=(3/4,1/4),
 \qquad p_1=(1/4,3/4).
\]
Taking their Gram in the uniform boundary measure gives $C^A$ in
\eqref{eq:V12-example-future}.  Subtract it from
$D_\pi=\operatorname{diag}(1/2,1/2)$ to obtain $V^A$; its trace is $3/8$.
The sums in \eqref{eq:V12-example-coherence-loss} are immediate.

For $F=2A_+-1$,
$\mathbb E[F\mid X]=2p_1-1=(-1/2,1/2)$.  Since
$\one_{\{A_+=1\}}(2A_+-1)=\one_{\{A_+=1\}}$, its selected entrance is $p_1$.
Subtracting $p_1b$ gives $c_1$.  The three norm and cross-term identities are
exact rational arithmetic.
\end{proof}

\begin{countertheorem}[Positive-time readability does not imply OS sector orthogonality]
\label{cth:V12-readable-not-sector}
There is no theorem which turns every finite positive-time Read into a direct
sum of Osterwalder--Schrader sectors using only the Markov property,
reflection symmetry, and positivity of the global form.  The cylinder in
\Cref{prop:V12-exact-separator} satisfies all three properties, yet its future
cell Gram has nonzero off-diagonal entry $3/16$ and positive boundary defect
$3/8$.
\end{countertheorem}

\begin{proof}
The claimed universal theorem would force the off-diagonal cell overlap to
vanish.  Equation~\eqref{eq:V12-example-future} gives a strictly positive exact
value.
\end{proof}

% =============================================================================
\section{The finite selector-first acquisition and its stopping semantics}
\label{sec:V12-acquisition}
% =============================================================================

The preceding theorems give a smaller and correctly ordered physical
acquisition than Version~V11.

\begin{definition}[Boundary/segment entrance card]
\label{def:V12-active-card}
On one predeclared RPESM cutoff, the packet
\begin{equation}
 \boxed{
 \begin{gathered}
 \texttt{REFLECTION\_FIXED\_BOUNDARY}\\[-0.15em]
 \texttt{\_AND\_SEGMENT\_ENTRANCE\_CARD}
 \end{gathered}}
 \label{eq:V12-active-card-name}
\end{equation}
contains:
\begin{enumerate}[label=\textup{(A\arabic*)},leftmargin=3.5em]
\item the common point-readable RPESM event tuple and its global provenance
      sigma-algebra $\mathscr H_{\rm prov}$;
\item the time-zero state and all declared reflection-fixed boundary marks,
      forming $\mathscr B_0$;
\item every finite future mark proposed as a sector label, together with its
      vacuum selector entrance Gram $C^A$ and
      $\Delta_{\partial}(A)$;
\item the finite positive-time one-letter writer bank and, for every
      non-boundary future mark retained by the target, the complete coherent
      cell--writer Gram $\mathbb G^A$;
\item on the boundary branch, the operator-valued measure
      $\mathbf M_{r_0}$ or its predeclared finite screen, together with the
      exact one-letter relation support;
\item when conditional hopping or Gaussianity is intended, the Version~V10
      measurable Grassmann/boundary germ on the global provenance and the
      Version~V9 time split, reflection, Berezin sign, and crossing Read;
\item one held-out positive-time writer or higher fermion word and, for a
      cofinal claim, one adjacent boundary-entrance rectangle.
\end{enumerate}
\end{definition}

\begin{theorem}[Terminating selector-first decision]
\label{thm:V12-selector-decision}
For each proposed finite record component $A$, the first applicable branch is
an exact terminating output.
\begin{enumerate}[label=\textup{(D\arabic*)},leftmargin=3.6em]
\item \emph{Boundary selector.}  An exact decoder $A=d(X_0)$ is supplied, or
      $\Delta_{\partial}(A)=0$ is proved.  The mark joins the boundary record,
      and the Version~V11 direct-sum cell theorem is licensed.
\item \emph{Future-record innovation.}  A strict lower bound
      $\Delta_{\partial}(A)>0$ is proved.  The polar bank of $V^A$ is retained
      as the source-minimal record innovation, and $A$ remains a writer.  If
      the target uses its cells, the complete Gram $\mathbb G^A$, including
      every cross block, is mandatory.
\item \emph{Unresolved eligibility.}  The outward interval for
      $\Delta_{\partial}(A)$ meets zero and no exact decoder is present.  The
      mark may be used as an ordinary future writer but not as a central
      record sector.
\item \emph{Input rejection.}  The submitted selector Gram is not positive,
      has incompatible row sums, or fails the reflected same-label and common-
      boundary incidence checks.  This is an ill-typed card, not a negative
      Standard-Model mode.
\end{enumerate}
After every component is classified, the one-letter OS form is evaluated on
the boundary POVM plus the coherent future-writer banks.  Only then are the
Version~V10/V9 conditional crossing gates applied.  A nearest block-diagonal
or nearest-positive replacement is not an admissible branch.
\end{theorem}

\begin{proof}
The first three branches are the mutually exclusive interval and exact-zero
alternatives of \Cref{thm:V12-boundary-determinism,cor:V12-eligibility-acquisition}.
The polar range statement follows because $V^A$ is the exact innovation Gram.
The required Gram on the positive branch is
\Cref{thm:V12-cell-word-Gram}.  Positivity, row sums, and common-boundary
incidence are defining identities of \eqref{eq:V12-selector-Gram}; their
failure places the input outside the theorem domain.
\end{proof}

\begin{proposition}[No duplicated source price across future cells]
\label{prop:V12-no-duplicate-cell-source}
The selected entrance columns obey the exact relation
\begin{equation}
 \boxed{b_j=\sum_ab_{a,j}.}
 \label{eq:V12-cell-source-relation}
\end{equation}
Therefore a source-minimal factorization of the enlarged cell--writer Gram
must descend through this relation before any cell is priced as a new source.
On the boundary-selector branch the cells are orthogonal pieces of the same
writer.  On the future-innovation branch only the component orthogonal to the
unselected and already represented entrance bank can enlarge the source atlas.
\end{proposition}

\begin{proof}
Equation~\eqref{eq:V12-cell-source-relation} is the first identity in
\eqref{eq:V12-cell-assembly}.  The remaining statements are the
Moore--Penrose source-minimal interpretation of an exact linear relation and
its orthogonal Schur residual.
\end{proof}

% =============================================================================
\section[RPESM application and corrected stopping line]{Application to the populated RPESM cylinder and the corrected stopping line}
\label{sec:V12-RPESM-landing}
% =============================================================================

The source audit changes two earlier status statements.

\begin{theorem}[Finite RPESM boundary occurrence is already populated]
\label{thm:V12-RPESM-occurrence}
For the existential RPESM regulator constructed in the inherited manuscript:
\begin{enumerate}[label=\textup{(P\arabic*)},leftmargin=3.6em]
\item the common unnormalized field-shell cylinder and the complete
      point-readable event tuple occur at every finite cutoff;
\item the time-zero Markov state and the physical boundary record occur on the
      same tuple;
\item every bounded positive-time writer has the Markov entrance square
      conditioned on that state;
\item every finite word, source, locality, line, operator, stress, and held-out
      grammar can be placed on the same event law, and a retained-terminal
      presentation supplies a cofinal bounded path grammar with summable debit.
\end{enumerate}
Hence the Version~V11 stop at the \emph{existence} of a readable RPESM record
is discharged on this named existential regulator.  The correct remaining
finite task is to separate the boundary subrecord from future marks by
\Cref{thm:V12-boundary-determinism}, and then to evaluate the already occurring
boundary-complete one-letter and crossing matrices.
\end{theorem}

\begin{proof}
The inherited regulator defines the bosonic record to contain the gauge and
Higgs fields, scheduler-derived geometry, auxiliaries, route data, and the
physical boundary record.  Its physical event tuple additionally retains the
initial and terminal fields and all source, response, line, and calibration
marks.  The reflected-slab theorem gives the conditional Markov square at
$X_0$.  The population theorem places every frozen finite grammar on this one
event law, and the retained-terminal theorem supplies the bounded cofinal path
selection.  These are occurrence statements.  The distinction between their
boundary and future components is exactly the decomposition
\eqref{eq:V12-three-record-algebras}.
\end{proof}

\begin{assessment}[What is and is not closed by population]
\label{ass:V12-population-scope}
The inherited population theorem is existential and structural.  It shows
that the finite matrices and mixed rows are functions or Reads on one named
regulator; it does not display the first selected numerical time-split
crossing matrix, its relation residual, its cone margin, the determinant or
Pfaffian transport margin, or the cofinal source-energy window.  Thus the
record-occurrence row is \passstatus{}, while the literal selected coefficient
card remains \stopstatus{}.  This is different from calling the record itself
unpopulated.
\end{assessment}

\begin{theorem}[Version V12 boundary-record and entrance theorem]
\label{thm:V12-integrated}
Preserving every theorem and limitation of Versions~V1--V11, and explicitly
superseding only the unrestricted record interpretation of Version~V11,
Version~V12 proves:
\begin{enumerate}[label=\textup{(V12.\arabic*)},leftmargin=6.2em]
\item global same-history provenance, reflection-fixed boundary records, and
      strictly positive-time writer records have different mathematical roles;
\item standard Borel disintegration defines historywise Grassmann coefficients
      without singleton positive-time selectors;
\item a boundary record produces the positive operator-valued measure
      \eqref{eq:V12-boundary-POVM} and the exact direct-integral OS
      factorization \eqref{eq:V12-direct-integral-form};
\item the finite Version~V11 direct sum is valid exactly on a
      reflection-fixed boundary record;
\item one vacuum selector Gram gives the positive boundary-determinism residual
      \eqref{eq:V12-boundary-defect};
\item its zero branch is exactly boundary readability, while its rank is the
      minimum linear future-record innovation on the positive branch;
\item a non-boundary future record produces the coherent cell--writer Gram
      \eqref{eq:V12-complete-cell-word-Gram}, whose cross blocks cannot be
      dropped before assembly;
\item the exact two-arm cylinder has future-record defect $3/8$, cell overlap
      $3/16$, and a selected-writer incidence correction $15/64$ in squared
      OS norm;
\item outward diagonal intervals decide strict non-boundary occurrence, while
      an exact decoder or exact zero is required for sector promotion;
\item the existential RPESM construction already populates the common event,
      time-zero state, boundary record, and finite grammar, so the next
      physical row is evaluation of its first boundary-complete crossing card,
      not another record or source compiler.
\end{enumerate}
No item evaluates that physical crossing matrix, proves a cofinal hopping or
line floor, establishes conditional Gaussianity on the active coefficient,
or reaches the Duhamel, current/stress, charge, continuum, or Einstein gates.
\end{theorem}

\begin{proof}
Items~\textup{(V12.1)--(V12.2)} are
\Cref{thm:V12-three-record-roles,fire:V12-disintegration-not-selector}.  Items
\textup{(V12.3)--(V12.4)} are
\Cref{thm:V12-boundary-POVM,cor:V12-finite-boundary-cells}.  Items
\textup{(V12.5)--(V12.6)} are
\Cref{thm:V12-boundary-determinism}.  Item~\textup{(V12.7)} is
\Cref{thm:V12-cell-word-Gram}.  Item~\textup{(V12.8)} is
\Cref{prop:V12-exact-separator}.  Item~\textup{(V12.9)} is
\Cref{cor:V12-eligibility-acquisition,thm:V12-selector-decision}.  The final
item is \Cref{thm:V12-RPESM-occurrence,ass:V12-population-scope}.
\end{proof}

\begingroup\small
\begin{longtable}{>{\raggedright\arraybackslash}p{0.27\textwidth}
                  >{\raggedright\arraybackslash}p{0.19\textwidth}
                  >{\raggedright\arraybackslash}p{0.46\textwidth}}
\caption{Version V12 proof-status ledger.}\label{tab:V12-ledger}\\
\toprule
\textbf{Row}&\textbf{Status}&\textbf{Exact conclusion}\\
\midrule
\endfirsthead
\toprule
\textbf{Row}&\textbf{Status}&\textbf{Exact conclusion}\\
\midrule
\endhead
V1--V10 Grassmann, hopping, tangent, Duhamel, and ancestry compilers
&Inherited
&Preserved verbatim; no prior matrix or semigroup theorem is recomputed.\\[0.3em]
V11 finite cell direct sum
&Qualified and corrected
&Exact for reflection-fixed boundary records; false for an arbitrary future
segment record because cross-cell entrance blocks can survive.\\[0.3em]
RPESM common event and provenance
&\passstatus{} finite occurrence
&The inherited point-readable tuple and population theorem place the full
finite grammar on one event law.\\[0.3em]
Historywise conditioning
&\passstatus{} standard Borel
&Regular conditional disintegration needs no positive-mass singleton
selector and is distinct from an OS sector projection.\\[0.3em]
Boundary record
&\passstatus{} structural occurrence
&The time-zero Markov state and physical boundary record are present on the
same RPESM event tuple.\\[0.3em]
Boundary OS kernel
&\passstatus{} exact factorization
&One positive operator-valued measure and the entrance synthesis give the
complete direct-integral reflected form.\\[0.3em]
Future-mark sector eligibility
&\passstatus{} finite compiler
&The vacuum selector Gram gives
$\Delta_{\partial}=1-\operatorname{Tr}C$; zero is exactly boundary
determinism.\\[0.3em]
Future record innovation
&\passstatus{} positive residual
&$D_\pi-C$ is the exact conditional one-hot innovation; its rank is the
minimum additional linear record dimension.\\[0.3em]
Future cell--writer matrix
&\passstatus{} exact coherent Gram
&Non-boundary cells require all cross blocks
$\mathbb G^A_{(a,i),(b,j)}$; diagonal dephasing changes the physical form.\\[0.3em]
Exact separator
&\passstatus{} rational witness
&The future bit has defect $3/8$, overlap $3/16$, and selected-writer norm
correction $15/64$.\\[0.3em]
Selected RPESM crossing value
&\stopstatus{}
&The constructed regulator supplies the writer, but the first explicit
boundary-complete time-split crossing interval and cone verdict are not
evaluated in the supplied sources.\\[0.3em]
Cofinal line, energy, and Duhamel margins
&\stopstatus{}
&Occurrence is finite; uniform source-energy, line orientation, exterior,
weighted-tail, and adjacent transport margins remain physical estimates.\\[0.3em]
Quantum continuum and Einstein equation
&Open
&No downstream conclusion is inferred from the record correction.\\
\bottomrule
\end{longtable}
\endgroup

\begin{assessment}[Status after Version V12]
\label{ass:V12-status}
\passstatus{} for the three-role record audit, standard-Borel disintegration,
boundary positive operator-valued measure, direct-integral entrance
factorization, boundary-determinism residual, coherent future-cell Gram,
finite outward decision, exact separator, and finite RPESM record occurrence.

\obsstatus{} for treating a strictly positive-time mark as a central
reflection-fixed sector merely because it is point readable; for requiring
singleton histories on a continuous field cylinder; for deleting cross-cell
entrance coherence; or for pricing each cell as a separate source before the
relation \eqref{eq:V12-cell-source-relation} is shorted.

\stopstatus{} at the first unevaluated physical tuple
\begin{equation}
 \boxed{
 \bigl(
 r_{0,n},\,\mathcal K_n^{\rm full},\,P_{-,n},P_{0,n},P_{+,n},\,
 R_{\Theta,n},\varepsilon_{\Theta,n},\,
 A_n,B_n,C_n,D_n,\,W_{1,n}
 \bigr),}
 \label{eq:V12-physical-stop}
\end{equation}
with all entries extracted from one already occurring RPESM field-shell
history and with the direct-minus-boundary-return crossing evaluated before
relation descent.  The stop concerns evaluation and uniform margins, not
existence of the common record.
\end{assessment}

\begin{openproblem}[Version V13: evaluate the first populated boundary crossing]
\label{op:V12-next}
Preserve Versions~V1--V12 verbatim.  No new record, source, exterior, lag,
localizer, concentration, or continuum theorem is admissible.  Proceed only
as follows.
\begin{enumerate}[label=\textup{(V13.\arabic*)},leftmargin=6.2em]
\item On the first predeclared RPESM control cutoff, extract $X_0$ and the
      reflection-fixed boundary record $r_{0,n}$ from the already occurring
      point-readable tuple.
\item For every future mark used as a proposed sector label, evaluate its
      selector entrance Gram.  Promote it only on the exact zero-defect branch;
      otherwise retain it inside the coherent writer bank and do not revisit
      the selector theory.
\item Extract the complete same-history quadratic Grassmann germ and its
      chronological blocks $A_n,B_n,C_n,D_n$.  If $D_n$ is invertible, perform
      the Version~V10 boundary solve and form the direct-minus-return crossing.
      If it is singular, export the zero-mode exterior/divisor packet and stop.
\item On the same card, evaluate the Version~V9 four-phase crossing germ,
      relation residual, Hermiticity/skew part, negative spectral part, and
      source-minimal positive factor rank.
\item Retain the determinant/Pfaffian line weight and sign separately.  A
      passing positive modulus does not orient either line.
\item Only after this finite crossing packet passes may the Version~V8
      Hamiltonian tangent and Versions~V5--V6 Duhamel/weighted-absorption cards
      be populated.  Stop at the first unavailable numerical coefficient or
      outward margin rather than opening another compiler.
\end{enumerate}
The acceptable outputs are a positive hopping factor, an explicit relation,
skew, negative-mode, divisor, or line obstruction, or a named unresolved
finite interval on this same physical card.
\end{openproblem}

\section*{V12 exact replay and append-only record}
\addcontentsline{toc}{section}{V12 exact replay and append-only record}
The accompanying standard-library verifier constructs the two-arm Markov
cylinder of \Cref{prop:V12-exact-separator} using exact rational arithmetic.
It verifies the boundary and future selector Grams, the one-hot innovation,
the ranks and determinants, complete-versus-diagonal vacuum assembly, the
selected-writer incidence formula, and every rational norm identity.
Ordinary and optimized Python executions produce byte-identical output.

The complete Version~V11 source preceding its sole final active
document-closing command is preserved byte for byte.  Its preterminal
SHA--256 digest is
\begin{center}\small\ttfamily
007688b735a465fe413bc2a3da4c37c86d9db9b1e8a9bbd0a79bcef648cf262f
\end{center}
The exact verifier and its JSON output have SHA--256 digests
\begin{center}\small\ttfamily
f0eeedbb93a9ad29007d6ed9ce15033185d4b2073aca00267ab49893ea8362e2\\
8a57a673fb72e6e04c4d1945bf4ceb258420f2afe4d14147e84ca1e479180de2
\end{center}
A later continuation must remove only the sole final active
\verb|\end{document}|, append new material immediately before it, restore one
terminal command, and increment the filename to end in \texttt{\_V13.tex}.

% =============================================================================
% END APPENDED VERSION V12 MATERIAL
% =============================================================================

% =============================================================================
% START APPENDED VERSION V13 MATERIAL -- complete V1--V12 preterminal source preserved
% =============================================================================

\clearpage
\hypersetup{
 pdftitle={Renewal Einstein--Standard--Model Physical Source Metric, Duhamel Ancestry and Quantum Continuum Programme V13},
 pdfsubject={The overlap-sign cross-boundary germ, free link-reflection factorization, and the first gauge-active crossing stop}
}
\section*{Version V13: the overlap-sign cross-boundary germ and the first gauge-active crossing stop}
\addcontentsline{toc}{section}{V13: overlap-sign crossing germ and gauge-active stop}

\begin{abstract}
Version~V12 correctly reduced the next quantum step to one populated
boundary-complete crossing card.  The literal finite field-shell source,
however, specifies that card through an overlap chiral coefficient.  Its moving
source space is
\(\Ran\widehat P_-\), with
\(\widehat P_-=(I+\operatorname{sgn}H_W)/2\), rather than an independently
split positive-time coefficient space.  The first coefficient to be read is
therefore the cross-boundary block of the Wilson sign, together with the
physical reflection frame and the one-letter source Gram.

This continuation makes that reduction exact.  The overlap kinetic term
collapses algebraically to \(2a^{-1}P_+\widehat P_-\).  On a link-reflection
split, every vertex-local Yukawa term is same-side, so the complete fermionic
crossing is
\[
 a^{-1}T_+\star_0P_+\operatorname{sgn}(H_W)T_-.
\]
The chronological compression of the moving chiral source is the positive
effect \(A_+=\widehat P_-T_+\widehat P_-\).  Its failure to be a projection is
measured by one exact positive scalar,
\[
 \Delta_{\chi/t}=\Tr(A_+-A_+^2)
 =\frac18\|[T_+,\operatorname{sgn}H_W]\|_{\rm HS}^2.
\]
This is not by itself a negative hopping mode: it says that the chiral source
must enter the inherited relation compiler through its actual nonisometric
one-letter synthesis rather than through a fictitious sharp half-space basis.

The Wilson gap localizes this incidence to the temporal boundary but does not
decide its sign.  An exact two-by-two overlap-sign family has the same Wilson
gap, positive full spectrum, determinant, and chronological leakage for both
members, while its fixed-reflection crossing has opposite signs.  Thus the
actual signed cross block remains indispensable.

On the gauge-trivial link-reflection branch, the free overlap spectral
representation gives the missing positive factor explicitly: the cross action
is an integral of reflected one-sided squares.  The conclusion persists under
reflection-paired vertex-local chiral Yukawa interactions.  A finite
Hilbert--Schmidt perturbation theorem then gives an open gauge-near-free branch
at every fixed cutoff.  It does not prove the arbitrary gauge-active chart.
The active RPESM stop is consequently sharpened from the entire time-split
precision to one finite, directly evaluable Wilson-sign boundary matrix and
its relation, Hermiticity, cone, line, and transport margins.
\end{abstract}

\begin{assessment}[Why Version V13 moves upstream without opening another programme]
\label{ass:V13-direction}
The inherited finite field shell defines its chiral coefficient by a gapped
Hermitian Wilson kernel, its sign, the overlap operator, and a local Yukawa
map.  The later reflected-hopping compiler assumes that the complete crossing
and its one-letter synthesis have already been populated.  These statements
leave one intermediate incidence uncomputed: how the moving overlap source
projector meets the physical positive/negative-time cut.  The present version
computes that incidence and recasts the known free overlap
link-reflection factorization in the exact Version~V9 matrix language.  It
does not introduce a new record, source algebra, exterior hierarchy, clock,
localizer, or continuum compactness theorem.
\end{assessment}

% =============================================================================
\section{The literal overlap coefficient and its cross-boundary reduction}
\label{sec:V13-overlap-reduction}
% =============================================================================

Let \(\mathscr H\) be the finite ambient lattice spin--internal Hilbert space.
Let \(\gamma_5=\gamma_5^*=\gamma_5^{-1}\), put
\begin{equation}
 P_+=\frac12(I+\gamma_5),
 \qquad
 \varepsilon_W=\operatorname{sgn}(H_W),
 \qquad
 \widehat P_-=\frac12(I+\varepsilon_W),
 \label{eq:V13-projectors}
\end{equation}
and assume the protected Wilson window
\begin{equation}
 H_W=H_W^*,
 \qquad
 \spec(H_W)\subset[-M_*,-g_*]\cup[g_*,M_*],
 \qquad g_*>0.
 \label{eq:V13-Wilson-window}
\end{equation}
The overlap operator in the inherited normalization is
\begin{equation}
 D_{\rm ov}=a^{-1}(I+\gamma_5\varepsilon_W),
 \qquad a>0.
 \label{eq:V13-Dov}
\end{equation}

\begin{lemma}[Exact collapse of the chiral overlap kinetic constituent]
\label{lem:V13-chiral-collapse}
One has the operator identity
\begin{equation}
 \boxed{
 P_+D_{\rm ov}\widehat P_-
 =\frac{2}{a}P_+\widehat P_-.}
 \label{eq:V13-chiral-collapse}
\end{equation}
No commutation between \(\gamma_5\) and \(\varepsilon_W\) is assumed.
\end{lemma}

\begin{proof}
Since \(P_+\gamma_5=P_+\),
\[
 P_+(I+\gamma_5\varepsilon_W)=P_+(I+\varepsilon_W).
\]
Moreover \(\varepsilon_W^2=I\), and hence
\[
 (I+\varepsilon_W)\widehat P_-
 =\frac12(I+\varepsilon_W)^2
 =I+\varepsilon_W
 =2\widehat P_-.
\]
Substitution in \eqref{eq:V13-Dov} proves
\eqref{eq:V13-chiral-collapse}.
\end{proof}

Freeze a link-reflection plane.  Let \(T_+\) and \(T_-\) be the ambient
positive- and negative-time projections, so
\begin{equation}
 T_+=T_+^*,\qquad T_-=T_-^*,\qquad
 T_+T_-=0,\qquad T_++T_-=I.
 \label{eq:V13-link-time-split}
\end{equation}
There is no reflection-fixed fermion site on this branch.  Assume that
\(T_\pm\) commute with \(\gamma_5\) and with the positive vertex Hodge factor
\(\star_0\).  Let \(\mathscr Y(H)\) be vertex local, so
\begin{equation}
 T_+\mathscr Y(H)T_-=0=T_-\mathscr Y(H)T_+.
 \label{eq:V13-Yukawa-local}
\end{equation}

\begin{theorem}[The active overlap crossing is one Wilson-sign boundary block]
\label{thm:V13-crossing-reduction}
For the completely assembled chiral coefficient
\begin{equation}
 \mathcal K_\chi
 =\star_0\left(P_+D_{\rm ov}\widehat P_-+\mathscr Y(H)\right),
 \label{eq:V13-Kchi}
\end{equation}
the negative-to-positive link-reflection crossing is
\begin{equation}
 \boxed{
 X_\times^{\rm ov}
 :=T_+\mathcal K_\chi T_-
 =\frac1a\,T_+\star_0P_+\varepsilon_WT_-.}
 \label{eq:V13-sign-crossing}
\end{equation}
Thus the local Yukawa constituent and the identity part of the overlap
operator do not enter the cross-plane matrix.  They remain in the same-side
blocks and in the determinant line.
\end{theorem}

\begin{proof}
By \Cref{lem:V13-chiral-collapse},
\[
 T_+\star_0P_+D_{\rm ov}\widehat P_-T_-
 =\frac2aT_+\star_0P_+\widehat P_-T_-.
\]
Insert \(\widehat P_-=(I+\varepsilon_W)/2\).  The identity contribution
vanishes because \(T_+T_-=0\), leaving the first term on the right of
\eqref{eq:V13-sign-crossing}.  Equation~\eqref{eq:V13-Yukawa-local} removes
the Yukawa cross block.  No inverse, determinant, covariance, or hard-tail
short is used.
\end{proof}

\begin{firewall}[Link reflection is a physical branch choice]
\label{fire:V13-link-reflection}
Equation~\eqref{eq:V13-sign-crossing} is the target-minimal formula on a
predeclared link-reflection branch.  If the active regulator uses a
site-reflection plane with a nonzero fixed fermion sector, the Version~V10
boundary block must first be assembled and eliminated or exported on its
divisor branch.  Choosing a link plane after inspecting the crossing sign is
not permitted.  The bosonic links cut by a link reflection remain part of the
reflection-fixed boundary record even though the fermion site projection
\(T_0\) is zero.
\end{firewall}

% =============================================================================
\section{Chronological incidence of the moving chiral source}
\label{sec:V13-chiral-time-incidence}
% =============================================================================

Let
\begin{equation}
 E_\chi=\Ran\widehat P_-,
 \qquad
 J_\chi:E_\chi\hookrightarrow\mathscr H
 \label{eq:V13-chiral-space}
\end{equation}
be the moving source space and its isometric inclusion.  Its two chronological
one-letter syntheses are
\begin{equation}
 W_{\chi,+}=T_+J_\chi,
 \qquad
 W_{\chi,-}=T_-J_\chi.
 \label{eq:V13-time-syntheses}
\end{equation}
Their source Grams are the complementary effects
\begin{equation}
 A_+=W_{\chi,+}^*W_{\chi,+}
     =J_\chi^*T_+J_\chi,
 \qquad
 A_-=I_{E_\chi}-A_+.
 \label{eq:V13-time-effects}
\end{equation}

\begin{theorem}[Exact overlap-chirality/time-incidence identity]
\label{thm:V13-time-incidence}
The following number is positive and has four equivalent formulas:
\begin{equation}
 \boxed{
 \begin{aligned}
 \Delta_{\chi/t}
 &:=\Tr_{E_\chi}(A_+-A_+^2)\\
 &=\|T_+\widehat P_-T_-\|_{\rm HS}^2\\
 &=\frac12\|[T_+,\widehat P_-]\|_{\rm HS}^2\\
 &=\frac18\|[T_+,\varepsilon_W]\|_{\rm HS}^2.
 \end{aligned}}
 \label{eq:V13-time-incidence}
\end{equation}
Moreover, the following are equivalent:
\begin{enumerate}[label=\textup{(\roman*)},leftmargin=2.8em]
\item \(\Delta_{\chi/t}=0\);
\item \(A_+\) is an orthogonal projection on \(E_\chi\);
\item \([T_+,\widehat P_-]=0\);
\item the moving chiral source splits sharply as
      \[
       E_\chi=(E_\chi\cap\Ran T_+)
       \oplus(E_\chi\cap\Ran T_-).
      \]
\end{enumerate}
On the positive branch the exact number of mixed chronological source
directions is
\begin{equation}
 \boxed{
 r_{\chi/t}
 :=\rank(A_+(I-A_+))
 =\dim E_\chi
  -\dim(E_\chi\cap\Ran T_+)
  -\dim(E_\chi\cap\Ran T_-).}
 \label{eq:V13-mixed-rank}
\end{equation}
\end{theorem}

\begin{proof}
On \(E_\chi\), the compression of \(T_+\) is
\(A_+=\widehat P_-T_+\widehat P_-\).  Hence
\[
 \Tr(A_+-A_+^2)
 =\Tr(\widehat P_-T_+\widehat P_-T_-)
 =\|T_+\widehat P_-T_-\|_{\rm HS}^2.
\]
Relative to \(\mathscr H=\Ran T_+\oplus\Ran T_-\), the commutator
\([T_+,\widehat P_-]\) consists of the two mutually adjoint off-diagonal
blocks \(T_+\widehat P_-T_-\) and
\(-T_-\widehat P_-T_+\).  This proves the second and third formulas.
The last follows from
\(\widehat P_-=(I+\varepsilon_W)/2\).

If the trace is zero, positivity of \(A_+-A_+^2\) gives
\(A_+^2=A_+\).  Equivalently,
\[
 0=\widehat P_-T_+(I-\widehat P_-)T_+\widehat P_-
   =\bigl((I-\widehat P_-)T_+\widehat P_-\bigr)^*
     \bigl((I-\widehat P_-)T_+\widehat P_-\bigr).
\]
Thus \(T_+\widehat P_-=\widehat P_-T_+\widehat P_-\); taking adjoints gives
\(\widehat P_-T_+=\widehat P_-T_+\widehat P_-\), so the projections commute.
The converse is immediate.  Commuting projections give the displayed
orthogonal intersection decomposition, and such a decomposition makes the
compression a projection.

Finally, the eigenvalue-one space of \(A_+\) is exactly
\(E_\chi\cap\Ran T_+\), the eigenvalue-zero space is exactly
\(E_\chi\cap\Ran T_-\), and every eigenvalue in \((0,1)\) contributes one
dimension to \(\rank(A_+(I-A_+))\).  Rank--nullity proves
\eqref{eq:V13-mixed-rank}.
\end{proof}

\begin{corollary}[The actual one-letter source metric]
\label{cor:V13-source-metric}
The positive-time overlap one-letter Gram is \(A_+\), not the identity.  Its
sharp finite visibility floor is
\begin{equation}
 \boxed{
 g_{\chi,+}=\lambda_{\min}^{+}(A_+),}
 \label{eq:V13-source-floor}
\end{equation}
and its canonical support is \(A_+A_+^\dagger\).  The Version~V9 relation
short must use the synthesis \(W_{\chi,+}\) and this support.  Replacing
\(A_+\) by the identity is licensed only on an independently proved isometric
chronological source chart.
\end{corollary}

\begin{proof}
The Gram statement is \eqref{eq:V13-time-effects}.  Finite spectral calculus
gives its supported floor and Moore--Penrose support.  The relation compiler
is invariant under the actual synthesis but not under deletion of its kernel
or rescaling of its positive eigenvalues without transporting the physical
source metric.
\end{proof}

\begin{proposition}[Canonical sharp dilation and the mixed chronological cost]
\label{prop:V13-Naimark}
Put \(S_\chi=(A_+(I-A_+))^{1/2}\).  On
\(E_\chi\oplus E_\chi\),
\begin{equation}
 \boxed{
 \Pi_{\chi,+}
 =\begin{pmatrix}
 A_+&S_\chi\\
 S_\chi&I-A_+
 \end{pmatrix}}
 \label{eq:V13-Naimark}
\end{equation}
is an orthogonal projection whose compression to the first copy is \(A_+\).
After deleting the zero ancillary directions, the minimum additional dilation
dimension is \(r_{\chi/t}\).
\end{proposition}

\begin{proof}
The two entries \(A_+\) and \(S_\chi\) commute.  Direct multiplication gives
\[
 A_+^2+S_\chi^2=A_+,
 \qquad
 A_+S_\chi+S_\chi(I-A_+)=S_\chi,
\]
and the analogous lower-right identity, proving \(\Pi_{\chi,+}^2=\Pi_{\chi,+}\).
On an eigenvector with eigenvalue zero or one, no ancillary direction is
needed.  For each eigenvalue \(\alpha\in(0,1)\), a projection having scalar
compression \(\alpha\) cannot act on a one-dimensional space and needs one
additional dimension; the displayed two-by-two block supplies it.  Summing
over the interior eigenvalues proves minimality and
\eqref{eq:V13-mixed-rank}.
\end{proof}

\begin{firewall}[A positive chronological incidence is not a hopping obstruction]
\label{fire:V13-incidence-not-hop}
A nonzero \(\Delta_{\chi/t}\) says that a global overlap-chiral coefficient
cannot be replaced by two independent sharp chiral half-space coefficient
lists.  It does not imply that the reflected crossing is negative.  The
correct response is to retain \(W_{\chi,+}\), its Gram \(A_+\), and the full
crossing \eqref{eq:V13-sign-crossing} in the inherited relation and cone test.
Promoting \(\Delta_{\chi/t}\) itself to a new fermion species or source birth
would duplicate the original overlap source.
\end{firewall}

% =============================================================================
\section{The Wilson gap localizes the incidence but does not decide the crossing}
\label{sec:V13-gap-bound}
% =============================================================================

\begin{theorem}[Sharp Hilbert--Schmidt commutator bound for the sign]
\label{thm:V13-sign-commutator}
Under \eqref{eq:V13-Wilson-window}, every operator \(T\) satisfies
\begin{equation}
 \boxed{
 \|[T,\varepsilon_W]\|_{\rm HS}
 \le g_*^{-1}\|[T,H_W]\|_{\rm HS}.}
 \label{eq:V13-sign-commutator}
\end{equation}
Consequently
\begin{equation}
 \boxed{
 0\le\Delta_{\chi/t}
 \le\frac{1}{8g_*^2}\|[T_+,H_W]\|_{\rm HS}^2
 =\frac{1}{4g_*^2}\|T_+H_WT_-\|_{\rm HS}^2.}
 \label{eq:V13-boundary-incidence-bound}
\end{equation}
For a finite-range Wilson kernel, the right side is supported on the temporal
boundary collar.  The estimate supplies an upper stock, not an exact split or
a sign.
\end{theorem}

\begin{proof}
Choose an orthonormal eigenbasis \((e_j)\) of \(H_W\), with eigenvalues
\(\lambda_j\).  Then
\[
 [T,\varepsilon_W]_{ij}
 =\bigl(\operatorname{sgn}\lambda_j-
        \operatorname{sgn}\lambda_i\bigr)T_{ij},
\]
while
\[
 [T,H_W]_{ij}=(\lambda_j-\lambda_i)T_{ij}.
\]
If \(\lambda_i,\lambda_j\) have the same sign, the first coefficient is zero.
If they have opposite signs, then
\[
 \frac{|\operatorname{sgn}\lambda_j-
        \operatorname{sgn}\lambda_i|}
      {|\lambda_j-\lambda_i|}
 \le\frac{2}{2g_*}=g_*^{-1}.
\]
Squaring and summing over \(i,j\) proves
\eqref{eq:V13-sign-commutator}.  Insert it into
\eqref{eq:V13-time-incidence}.  The final equality uses the two off-diagonal
blocks of the commutator of a projection with a self-adjoint operator.
\end{proof}

\begin{countertheorem}[Gap, determinant modulus, and leakage do not determine the crossing sign]
\label{cth:V13-two-by-two}
For \(\sigma\in\{+1,-1\}\), let
\begin{equation}
 H_\sigma
 =\begin{pmatrix}3/4&\sigma\\ \sigma&-3/4\end{pmatrix},
 \qquad
 T_+=\begin{pmatrix}1&0\\0&0\end{pmatrix}.
 \label{eq:V13-Hsigma}
\end{equation}
Then
\begin{equation}
 H_\sigma^2=\frac{25}{16}I,
 \qquad
 |H_\sigma|=\frac54I,
 \label{eq:V13-Hsigma-gap}
\end{equation}
and
\begin{equation}
 \varepsilon_\sigma
 =\begin{pmatrix}3/5&4\sigma/5\\4\sigma/5&-3/5\end{pmatrix},
 \qquad
 \widehat P_\sigma
 =\begin{pmatrix}4/5&2\sigma/5\\2\sigma/5&1/5\end{pmatrix}.
 \label{eq:V13-Psigma}
\end{equation}
Both cards have
\begin{equation}
 \boxed{
 \Delta_{\chi/t}=\frac4{25},
 \qquad
 g_{\chi,+}=\frac45,}
 \label{eq:V13-same-incidence}
\end{equation}
and the bound \eqref{eq:V13-boundary-incidence-bound} is an equality.
The massive overlap-sign coefficients
\begin{equation}
 \mathcal K_\sigma=I+2\widehat P_\sigma
 =\begin{pmatrix}13/5&4\sigma/5\\4\sigma/5&7/5\end{pmatrix}
 \label{eq:V13-Ksigma}
\end{equation}
have the same strictly positive spectrum \(\{1,3\}\) and determinant three.
Nevertheless, with one fixed reflection frame and
\(\varepsilon_\Theta=-1\), their crossing matrices are the opposite scalars
\begin{equation}
 \boxed{K_{\times,\sigma}=\frac{4\sigma}{5}.}
 \label{eq:V13-opposite-crossing}
\end{equation}
Thus the \(\sigma=+1\) card has a rank-one positive factor and the
\(\sigma=-1\) card has hopping-cone distance squared \(16/25\).
\end{countertheorem}

\begin{proof}
Direct multiplication gives \eqref{eq:V13-Hsigma-gap}; division by \(5/4\)
and \((I+\varepsilon_\sigma)/2\) give
\eqref{eq:V13-Psigma}.  The positive eigenvector of
\(\widehat P_\sigma\) is proportional to \((2,\sigma)\), so the
positive-time compression has the single eigenvalue \(4/5\) and
\(\Delta_{\chi/t}=(4/5)(1/5)=4/25\).  Moreover
\(\|[T_+,H_\sigma]\|_{\rm HS}^2=2\), and the right side of
\eqref{eq:V13-boundary-incidence-bound} is
\[
 \frac{2}{8(5/4)^2}=\frac4{25}.
\]
The trace and determinant of \(\mathcal K_\sigma\) are four and three,
respectively, proving the common spectrum.  Its upper-right time block is
\(4\sigma/5\); with the stated common reflection and Berezin convention this
is the reflection-adapted scalar.  The Version~V9 cone residual is zero for
the positive scalar and the square of the negative scalar for the other.
\end{proof}

\begin{remark}[Gauge/frame covariance of the separator]
The two cards in \Cref{cth:V13-two-by-two} are not claimed to be distinct when
both the negative-time basis and the physical reflection map are transformed
simultaneously.  Under a legitimate simultaneous frame change, the
reflection-adapted crossing changes by unitary congruence and its inertia is
unchanged.  The countertheorem fixes the reflection record and varies the
relative cross-block incidence.  It proves that the positive Wilson window,
full coefficient spectrum, determinant modulus, and chronological leakage do
not reconstruct that missing relative orientation.
\end{remark}

% =============================================================================
\section{Free overlap link-reflection factorization}
\label{sec:V13-free-factorization}
% =============================================================================

The following finite control is the precise positive branch relevant to the
field-shell overlap constituent.  It is a recasting of the free overlap
link-reflection spectral factorization of Y.~Kikukawa and K.~Usui,
\emph{Phys. Rev. D} \textbf{82} (2010), 114503,
\href{https://arxiv.org/abs/1005.3751}{arXiv:1005.3751}.  The positivity of the
spectral matrices needed by the argument is proved below because it is the
coefficient which interfaces with the present crossing card.

Let \(\Lambda=[-L+1,L]^4\subset\mathbb Z^4\), with anti-periodic temporal and
periodic spatial boundary conditions, and link reflection
\(\theta(t,\mathbf x)=(1-t,\mathbf x)\).  Let the Euclidean gamma matrices be
Hermitian and satisfy the Clifford relations.  For \(0<m\le1\), put
\begin{equation}
 X(p)=i\sum_{\mu=0}^{3}\gamma_\mu\sin p_\mu
      +\sum_{\mu=0}^{3}(1-\cos p_\mu)-m,
 \qquad
 D_{\rm free}=\frac12\left(I+X(X^*X)^{-1/2}\right).
 \label{eq:V13-free-overlap}
\end{equation}

For spatial momentum \(\mathbf p\) and \(E>0\), define
\begin{equation}
 \begin{aligned}
 W(E,\mathbf p)
 &=\sum_{k=1}^{3}(1-\cos p_k)+1-\cosh E-m,\\
 r(E,\mathbf p)^2
 &=\sum_{k=1}^{3}\sin^2p_k+W(E,\mathbf p)^2.
 \end{aligned}
 \label{eq:V13-Wr}
\end{equation}
Let \(E_1(\mathbf p)\) be the first positive solution of
\begin{equation}
 \sinh^2E=r(E,\mathbf p)^2.
 \label{eq:V13-cut-edge}
\end{equation}

\begin{lemma}[Positive spectral matrices on the overlap cut]
\label{lem:V13-spectral-positive}
For every \(E\ge E_1(\mathbf p)\),
\begin{equation}
 \boxed{
 B_+(E,\mathbf p):=-\gamma_0X(iE,\mathbf p)\succeq0,
 \qquad
 B_-(E,\mathbf p):=\gamma_0X(-iE,\mathbf p)\succeq0.}
 \label{eq:V13-Bpm-positive}
\end{equation}
Consequently there are measurable positive square roots
\(Y_\pm=B_\pm^{1/2}\).
\end{lemma}

\begin{proof}
Since \(\sin(iE)=i\sinh E\) and \(\cos(iE)=\cosh E\),
\[
 B_+=\sinh E\,I-i\sum_{k=1}^{3}\gamma_0\gamma_k\sin p_k
      -W(E,\mathbf p)\gamma_0.
\]
The second and third terms form a Hermitian matrix \(C_+\).  Clifford
anticommutation gives
\[
 C_+^2=r(E,\mathbf p)^2I.
\]
Thus the eigenvalues of \(B_+\) are
\(\sinh E\pm r(E,\mathbf p)\), which are nonnegative by
\eqref{eq:V13-cut-edge} and the definition of the cut.  The same calculation
for \(B_-\) changes the sign of \(C_+\) and yields the same two eigenvalues.
Finite-dimensional spectral calculus gives the measurable positive square
roots.
\end{proof}

\begin{theorem}[Free overlap crossing is an integral of one-sided reflected squares]
\label{thm:V13-free-overlap-factor}
For the finite free overlap action with the preceding anti-periodic and
link-reflection conventions, its complete cross-plane action has a convergent
representation
\begin{equation}
 \boxed{
 \mathscr C_{\times,0}
 =\int_{\Sigma_L}\Theta_F(F_s)F_s\,\dd\nu_L(s),
 \qquad \nu_L\ge0,}
 \label{eq:V13-free-square-integral}
\end{equation}
where \(s\) consists of a spatial momentum, an energy
\(E\ge E_1(\mathbf p)\), a spin index, and the two positive anti-periodic image
branches.  The one-sided rows \(F_s\) are obtained by multiplying the
positive-time exponential wave packets by the square roots \(Y_+\) and
\(Y_-\) of \Cref{lem:V13-spectral-positive}.

Equivalently, on every finite physical one-letter screen there is a synthesis
\(R_{\times,0}\) such that
\begin{equation}
 \boxed{
 K_{\times,0}=R_{\times,0}^*R_{\times,0}\succeq0.}
 \label{eq:V13-free-K-positive}
\end{equation}
After the exact one-letter relation short, the Version~V9 residuals satisfy
\begin{equation}
 \boxed{
 \Delta_{\rm rel}^{\times}=0,
 \qquad
 \Delta_{\rm hop}^{\times}=0,}
 \label{eq:V13-free-residuals}
\end{equation}
and the minimum finite one-sided source number is
\(\rank K_{\times,0}\).
\end{theorem}

\begin{proof}
For unequal times, Fourier transformation in the spatial variables writes the
nonlocal part of \(D_{\rm free}\) as the temporal Fourier integral of
\(X(p_0,\mathbf p)/(2\sqrt{X^*X(p_0,\mathbf p)})\).  For
\(x_0>0\ge y_0\), deform the \(p_0\) contour to the positive imaginary cut
\(p_0=iE\), \(E\ge E_1(\mathbf p)\); for the reflected orientation use the
negative cut.  The jump of the square root supplies a positive scalar density,
and anti-periodic temporal images supply the two positive weights
\[
 \frac{1}{1+e^{-2EL}},
 \qquad
 \frac{e^{-2EL}}{1+e^{-2EL}}.
\]
The remaining spin matrices are precisely \(B_+\) and \(B_-\) in
\eqref{eq:V13-Bpm-positive}.  Factor them as
\(Y_\pm^*Y_\pm\), split the exponential
\(e^{-E(x_0-y_0)}\) at the reflection plane, and move the negative-time field
through the frozen order-reversing fermionic reflection.  The Berezin-order
minus sign cancels the contour sign.  Each integrand is then
\(\Theta_F(F_s)F_s\), proving
\eqref{eq:V13-free-square-integral}.  Riemann sums converge coefficientwise in
the finite Grassmann algebra, so the integral belongs to the closure of the
one-sided positive cone.

Compress the rows \(F_s\) to any finite one-letter coefficient screen.  Their
Gram is \(K_{\times,0}\), proving
\eqref{eq:V13-free-K-positive}.  Since this Gram is formed from the actual
one-letter synthesis, it annihilates every zero represented word on both
legs, which is the relation identity.  The positive Gram has zero skew and
negative parts.  The rank and uniqueness statement is the inherited polar
factor theorem.
\end{proof}

\begin{corollary}[Weyl compression and vertex-local chiral Yukawa interactions]
\label{cor:V13-Weyl-Yukawa}
On the gauge-trivial field-shell branch, compressing the free Dirac one-letter
factor through the physical overlap-Weyl source synthesis preserves
positivity.  Adding a reflection-paired vertex-local chiral Yukawa action does
not change the cross matrix \eqref{eq:V13-sign-crossing}; it modifies only the
two same-side actions and the determinant line.  Therefore the free overlap
Weyl/Yukawa control has a passing reflected-hopping gate.
\end{corollary}

\begin{proof}
A compression of \(R^*R\) by any synthesis \(W\) is
\((RW)^*(RW)\succeq0\), and its relation support is automatic.  A vertex-local
Yukawa action on a link-reflected lattice has the form
\(A_Y=A_{Y,+}+\Theta_F(A_{Y,+})\), with no mixed-time bilinear.  Multiplication
of the bare Grassmann functional by
\(e^{A_{Y,+}}\Theta_F(e^{A_{Y,+}})\) preserves the reflected positive cone,
while the crossing factor remains the free one.  This is exactly the nongauge
chiral Yukawa scope of the free overlap factorization.
\end{proof}

\begin{firewall}[The free theorem is not the active gauge theorem]
\label{fire:V13-free-not-gauge}
The finite RPESM coefficient contains a nontrivial \(G_{\rm SM}\) link field.
The preceding contour factorization uses spatial translation invariance and a
momentum-wise positive cut matrix.  It proves the gauge-trivial and nongauge
Yukawa control branch, not an arbitrary gauge-history branch.  The cited free
overlap theorem explicitly leaves the gauge-model reflection-positivity problem
outside its proof.  Version~V13 therefore does not copy the free conclusion
onto the active field law.
\end{firewall}

% =============================================================================
\section{A fixed-cutoff gauge-near-free stability window}
\label{sec:V13-gauge-stability}
% =============================================================================

The free factorization gives a concrete base card.  A finite perturbation
estimate identifies an open branch on which the same represented crossing
remains positive.  It is not a cofinal estimate.

\begin{theorem}[Hilbert--Schmidt Lipschitz control of a gapped sign path]
\label{thm:V13-sign-path}
Let
\begin{equation}
 H_s=H_0+sE,
 \qquad 0\le s\le1,
 \label{eq:V13-Hpath}
\end{equation}
be a self-adjoint finite-dimensional path satisfying
\begin{equation}
 \dist(0,\spec H_s)\ge g_{\rm path}>0
 \quad\text{for every }s.
 \label{eq:V13-path-gap}
\end{equation}
Then
\begin{equation}
 \boxed{
 \|\operatorname{sgn}H_1-\operatorname{sgn}H_0\|_{\rm HS}
 \le g_{\rm path}^{-1}\|E\|_{\rm HS}.}
 \label{eq:V13-sign-Lipschitz}
\end{equation}
In particular, if \(H_0\) has gap \(g_0\) and
\(\|E\|_{\rm op}<g_0\), one may take
\(g_{\rm path}=g_0-\|E\|_{\rm op}\).
\end{theorem}

\begin{proof}
At every parameter where a differentiable spectral frame is chosen, the
Fr\'echet derivative of the sign has divided-difference entries
\[
 (\dot\varepsilon_s)_{ij}
 =\frac{\operatorname{sgn}\lambda_i(s)-
        \operatorname{sgn}\lambda_j(s)}
       {\lambda_i(s)-\lambda_j(s)}E_{ij},
\]
with the same-sign and diagonal coefficients equal to zero.  Opposite-sign
eigenvalues differ by at least \(2g_{\rm path}\), so the absolute divided
difference is at most \(g_{\rm path}^{-1}\).  Summing the squared entries gives
\(\|\dot\varepsilon_s\|_{\rm HS}\le
 g_{\rm path}^{-1}\|E\|_{\rm HS}\).  The formula is basis independent and
extends through eigenvalue multiplicities by spectral projections.  Integrate
from zero to one.  The last statement is Weyl's eigenvalue perturbation
bound along \eqref{eq:V13-Hpath}.
\end{proof}

\begin{corollary}[Finite gauge-near-free crossing window]
\label{cor:V13-gauge-window}
Transport a gauge-history card and the free card to one reflection-adapted
boundary frame, and suppose \(\star_0\), \(P_+\), \(T_\pm\), and the Berezin
sign agree in that frame.  Put
\begin{equation}
 E_U=H_W(U)-H_W(I),
 \qquad
 g_U=\inf_{0\le s\le1}\dist(0,\spec(H_W(I)+sE_U)).
 \label{eq:V13-EU}
\end{equation}
Then the overlap crossing matrices satisfy
\begin{equation}
 \boxed{
 \|K_\times(U)-K_\times(I)\|_{\rm HS}
 \le
 \frac{\|\star_0\|_{\rm op}}{a\,g_U}
 \|E_U\|_{\rm HS}.}
 \label{eq:V13-crossing-perturbation}
\end{equation}
Let \(P\) be a predeclared finite represented one-letter screen on which the
free crossing has floor
\begin{equation}
 PK_\times(I)P\succeq\kappa_{0,P}P,
 \qquad \kappa_{0,P}>0.
 \label{eq:V13-free-screen-floor}
\end{equation}
If reflection covariance makes the transported gauge crossing Hermitian and
\begin{equation}
 \frac{\|\star_0\|_{\rm op}}{a\,g_U}
 \|E_U\|_{\rm HS}<\kappa_{0,P},
 \label{eq:V13-open-gauge-condition}
\end{equation}
then
\begin{equation}
 \boxed{
 PK_\times(U)P\succeq
 \left(
 \kappa_{0,P}-
 \frac{\|\star_0\|_{\rm op}}{a\,g_U}\|E_U\|_{\rm HS}
 \right)P>0.}
 \label{eq:V13-gauge-floor}
\end{equation}
Thus every fixed free positive screen has a nonempty open gauge-near-free
passing neighbourhood.
\end{corollary}

\begin{proof}
Use \eqref{eq:V13-sign-crossing}, contractivity of the chronological and
chiral projections, unitarity of the reflection frame, and
\Cref{thm:V13-sign-path}.  The Hilbert--Schmidt norm dominates the operator
norm.  Weyl's lower-eigenvalue inequality on \(\Ran P\) gives
\eqref{eq:V13-gauge-floor}.
\end{proof}

\begin{remark}[What must be added when the physical frame moves]
If the Hodge factor, chronological identification, or reflection map varies
with the gauge/history card, their independently transported differences add
to the right side of \eqref{eq:V13-crossing-perturbation} by the product rule.
A change of the discrete Berezin sign is not a small norm error; it is a line
orientation event.  Version~V13 keeps these rows separate rather than hiding
them in \(E_U\).
\end{remark}

\begin{firewall}[No uniform gauge theorem from a fixed-cutoff neighbourhood]
\label{fire:V13-no-uniform-gauge}
The free screen floor \(\kappa_{0,P}\), the Wilson path gap \(g_U\), and the
gauge perturbation stock may all deteriorate with the cutoff, volume, or
selected history.  The open condition
\eqref{eq:V13-open-gauge-condition} proves nonemptiness and a terminating
finite branch.  It does not establish a regulator-uniform reflected-hopping
margin on the full gauge-active RPESM law.
\end{firewall}

% =============================================================================
\section{Validated finite reconstruction of the Wilson-sign crossing}
\label{sec:V13-validated-sign}
% =============================================================================

The remaining active card is finite.  It may be evaluated by exact spectral
calculus or by a certified sign approximation.

\begin{theorem}[Polynomial or rational sign enclosure]
\label{thm:V13-sign-enclosure}
Let \(p\) be a real odd polynomial or a pole-free real rational function on
\([-M_*,-g_*]\cup[g_*,M_*]\) satisfying
\begin{equation}
 \sup_{g_*\le|x|\le M_*}
 |p(x)-\operatorname{sgn}x|
 \le\eta_p.
 \label{eq:V13-uniform-sign-error}
\end{equation}
Put \(\widetilde\varepsilon=p(H_W)\), let
\(d=\dim\mathscr H\), and define
\begin{equation}
 \widetilde K_\times
 =-\varepsilon_\Theta a^{-1}
  T_+\star_0P_+\widetilde\varepsilon T_-R_\Theta.
 \label{eq:V13-Ktilde}
\end{equation}
Then
\begin{equation}
 \boxed{
 \|\widetilde\varepsilon-\varepsilon_W\|_{\rm HS}
 \le\sqrt d\,\eta_p,
 \qquad
 \|\widetilde K_\times-K_\times\|_{\rm HS}
 \le\frac{\|\star_0\|_{\rm op}}a\sqrt d\,\eta_p.}
 \label{eq:V13-sign-cross-error}
\end{equation}
If
\begin{equation}
 c_p=\|[T_+,\widetilde\varepsilon]\|_{\rm HS},
 \qquad
 e_p=2\sqrt d\,\eta_p,
 \label{eq:V13-cp-ep}
\end{equation}
then the exact chronological incidence is enclosed by
\begin{equation}
 \boxed{
 \frac18\max\{0,c_p-e_p\}^2
 \le\Delta_{\chi/t}
 \le\frac18(c_p+e_p)^2.}
 \label{eq:V13-incidence-interval}
\end{equation}
\end{theorem}

\begin{proof}
Functional calculus and \eqref{eq:V13-uniform-sign-error} give the operator
error \(\eta_p\), and the Hilbert--Schmidt error is at most
\(\sqrt d\eta_p\).  Formula~\eqref{eq:V13-sign-crossing}, contractivity of the
projections, and unitarity of \(R_\Theta\) prove the second estimate.
Furthermore,
\[
 \bigl|
 \|[T_+,\varepsilon_W]\|_{\rm HS}
 -\|[T_+,\widetilde\varepsilon]\|_{\rm HS}
 \bigr|
 \le2\|\varepsilon_W-\widetilde\varepsilon\|_{\rm HS}
 \le e_p.
\]
Square the resulting interval and use
\eqref{eq:V13-time-incidence}.
\end{proof}

\begin{acquisition}[Target-minimal overlap crossing card]
\label{acq:V13-overlap-card}
At one selected RPESM cutoff and one predeclared bosonic boundary record,
acquire only
\begin{equation}
 \boxed{
 \begin{gathered}
 H_W,\ [g_*,M_*],\ T_\pm,\ P_+,\ \star_0,\ R_\Theta,
 \ \varepsilon_\Theta,\\
 W_{\chi,+}=T_+\widehat P_-|_{E_\chi},\\
 K_\times=-\varepsilon_\Theta a^{-1}
 T_+\star_0P_+\varepsilon_WT_-R_\Theta,\\
 \text{canonical supports and outward spectral radii},\\
 \text{one held-out crossing probe and one adjacent card.}
 \end{gathered}}
 \label{eq:V13-overlap-card}
\end{equation}
The already proved Version~V9 relation and cone residuals are then run without
modification.  The determinant/Pfaffian line weight, divisor, and orientation
remain separate entries.
\end{acquisition}

\begin{theorem}[Finite decision semantics for the overlap card]
\label{thm:V13-decision}
After the common-history, reflection-frame, gap, support, and outward-radius
checks pass, the selected card has exactly one of the following outcomes.
\begin{enumerate}[label=\textup{(O\arabic*)},leftmargin=3.5em]
\item \emph{Positive overlap hopping.}  The relation residual and skew interval
      vanish and the lower eigenvalue interval of the represented crossing is
      positive.  The positive square root gives the source-minimal factor and
      its rank.
\item \emph{Explicit crossing obstruction.}  A strict relation, skew, or
      negative-part lower bound returns its finite witness.
\item \emph{Divisor or line obstruction.}  The crossing may pass while the
      complete coefficient lies on a divisor or its required real/complex line
      orientation fails.
\item \emph{Unresolved finite interval.}  Every domain check passes but at
      least one outward interval meets its threshold.
\end{enumerate}
Failure of the common history, the link/site-reflection declaration, the
Wilson gap, or the canonical source support rejects the input before this
list.  It is not a fifth mathematical outcome.
\end{theorem}

\begin{proof}
The sign enclosure supplies the raw crossing interval and the actual
one-letter Gram.  Apply the Version~V9 supported relation test, Hermitian/skew
split, and exact distance to the positive cone.  The determinant and Pfaffian
rows are independent by the inherited line theorem and by
\Cref{cth:V13-two-by-two}, which has equal positive determinant data and
opposite fixed-reflection crossings.  The four alternatives exhaust strict
positive, strict negative, line/divisor, and threshold-touching valid cards.
\end{proof}

% =============================================================================
\section{Integrated V13 theorem and the revised active stop}
\label{sec:V13-integrated}
% =============================================================================

\begin{theorem}[Version V13 overlap-crossing theorem]
\label{thm:V13-integrated}
Preserving every theorem and limitation of Versions~V1--V12, Version~V13
proves:
\begin{enumerate}[label=\textup{(V13.\arabic*)},leftmargin=6.2em]
\item the literal overlap chiral kinetic constituent is
      \(2a^{-1}P_+\widehat P_-\);
\item on a predeclared link-reflection branch its complete local-Yukawa
      crossing is the single Wilson-sign boundary block
      \eqref{eq:V13-sign-crossing};
\item the positive-time one-letter source metric is the compression
      \(A_+=\widehat P_-T_+\widehat P_-\), not an identity metric;
\item the overlap chirality/time incidence has the four exact positive forms
      \eqref{eq:V13-time-incidence}, and its rank is the exact number of mixed
      chronological source directions;
\item the Wilson gap bounds that incidence by the physical temporal-boundary
      hopping stock through \eqref{eq:V13-boundary-incidence-bound};
\item the Wilson gap, determinant modulus, full positive spectrum, and
      chronological leakage do not determine the reflected crossing sign;
\item the free overlap link-reflection crossing is an integral of reflected
      one-sided squares and passes the Version~V9 relation and hopping gate;
\item reflection-paired vertex-local nongauge chiral Yukawa interactions retain
      that crossing factor;
\item every fixed free positive screen has the explicit open gauge-near-free
      stability window \eqref{eq:V13-open-gauge-condition};
\item exact spectral calculus or one certified sign approximant produces
      outward intervals for both the chronological incidence and the physical
      crossing.
\end{enumerate}
No item proves reflection positivity on an arbitrary gauge history, a
cutoff-uniform overlap hopping floor, determinant/Pfaffian orientation,
conditional Gaussianity, the Duhamel residual, represented current/stress,
charge, the interacting quantum continuum, or the Einstein equation.
\end{theorem}

\begin{proof}
Items~\textup{(V13.1)--(V13.2)} are
\Cref{lem:V13-chiral-collapse,thm:V13-crossing-reduction}.  Items
\textup{(V13.3)--(V13.4)} are
\Cref{thm:V13-time-incidence,cor:V13-source-metric,prop:V13-Naimark}.  Item
\textup{(V13.5)} is \Cref{thm:V13-sign-commutator}.  Item
\textup{(V13.6)} is \Cref{cth:V13-two-by-two}.  Items
\textup{(V13.7)--(V13.8)} are
\Cref{thm:V13-free-overlap-factor,cor:V13-Weyl-Yukawa}.  Item
\textup{(V13.9)} is \Cref{thm:V13-sign-path,cor:V13-gauge-window}.  The final
item is \Cref{thm:V13-sign-enclosure,acq:V13-overlap-card}.
\end{proof}

\begingroup\small
\begin{longtable}{>{\raggedright\arraybackslash}p{0.27\textwidth}
                  >{\raggedright\arraybackslash}p{0.19\textwidth}
                  >{\raggedright\arraybackslash}p{0.46\textwidth}}
\caption{Version V13 proof-status ledger.}\label{tab:V13-ledger}\\
\toprule
\textbf{Row}&\textbf{Status}&\textbf{Exact conclusion}\\
\midrule
\endfirsthead
\toprule
\textbf{Row}&\textbf{Status}&\textbf{Exact conclusion}\\
\midrule
\endhead
V1--V12 record, Grassmann, crossing, tangent, Duhamel, and ancestry results
&Inherited
&Preserved verbatim; no earlier compiler is repeated.\\[0.3em]
Literal overlap kinetic constituent
&\passstatus{} exact
&\(P_+D_{\rm ov}\widehat P_-=2a^{-1}P_+\widehat P_-\), without assuming
commutation of \(\gamma_5\) and the Wilson sign.\\[0.3em]
Link-reflection crossing reduction
&\passstatus{} exact branch
&For a vertex-local Yukawa map, the entire cross-plane coefficient is
\(a^{-1}T_+\star_0P_+\varepsilon_WT_-\).\\[0.3em]
Chiral chronological source metric
&\passstatus{} exact
&The actual Gram is \(A_+=\widehat P_-T_+\widehat P_-\); its Moore--Penrose
support and positive floor replace identity normalization.\\[0.3em]
Chirality/time incidence
&\passstatus{} positive identity
&\(\Delta_{\chi/t}=\Tr(A_+-A_+^2)=\|T_+\widehat P_-T_-\|_{\rm HS}^2\), with
mixed rank \eqref{eq:V13-mixed-rank}.\\[0.3em]
Wilson boundary upper stock
&\passstatus{} quantitative
&The sign commutator is bounded by the Wilson-kernel boundary commutator divided
by the protected gap.\\[0.3em]
Gap/modulus inference
&\obsstatus{} exact separator
&Equal Wilson gaps, positive full spectra, determinants, and leakage can carry
opposite fixed-reflection crossing signs.\\[0.3em]
Free overlap crossing
&\passstatus{} positive control
&Link-reflection spectral factorization writes the cross action as an integral
of one-sided reflected squares.\\[0.3em]
Nongauge local Yukawa extension
&\passstatus{} same crossing
&Reflection-paired vertex-local Yukawa terms modify same-side actions and the
line but not the cross matrix.\\[0.3em]
Gauge-near-free fixed card
&\passstatus{} conditional open set
&The sign-path estimate gives an explicit finite perturbation radius preserving
every predeclared positive free screen.\\[0.3em]
Arbitrary gauge-active RPESM history
&\stopstatus{}
&The source files do not supply a uniform or eventwise evaluated
\(T_+\operatorname{sgn}H_W(U)T_-\) cone interval on the selected gauge card.\\[0.3em]
Determinant/Pfaffian line
&\stopstatus{} separate
&The chiral floor gives a nonzero modulus on its protected chart, but the active
phase transport and real Pfaffian orientation remain independent.\\[0.3em]
Duhamel, current/stress, charge, continuum, Einstein
&Closed downstream
&These rows are not entered before the active crossing and line card passes.\\
\bottomrule
\end{longtable}
\endgroup

\begin{assessment}[Status after Version V13]
\label{ass:V13-status}
\passstatus{} for the exact overlap kinetic reduction, chronological source
effect, time-incidence identity, boundary commutator bound, free
link-reflection one-sided factorization, nongauge local-Yukawa extension,
fixed-cutoff gauge stability theorem, and validated sign/crossing enclosure.

\obsstatus{} for replacing the moving overlap source metric by the identity;
for treating a nonzero chronological incidence as a new species; for inferring
the crossing sign from a Wilson gap, determinant modulus, or positive full
precision; or for applying the free Fourier factorization to arbitrary gauge
histories.

\stopstatus{} on the active RPESM gauge branch at the finite matrix
\begin{equation}
 \boxed{
 K_{\times,n}(U,H)
 =-\varepsilon_{\Theta,n}a_n^{-1}
 T_{+,n}\star_{0,n}P_{+,n}
 \operatorname{sgn}(H_{W,n}(U))
 T_{-,n}R_{\Theta,n},}
 \label{eq:V13-active-stop}
\end{equation}
together with its actual one-letter Gram
\begin{equation}
 \boxed{
 A_{+,n}
 =\widehat P_{-,n}T_{+,n}\widehat P_{-,n}.}
 \label{eq:V13-active-Gram}
\end{equation}
The stop is now one explicit Wilson-sign boundary Read, not an unspecified
fermion cylinder or another abstract OS compiler.
\end{assessment}

\begin{openproblem}[Version V14: execute one gauge-active overlap crossing card]
\label{op:V13-next}
Preserve Versions~V1--V13 verbatim.  Do not add another record, Grassmann,
exterior, Duhamel, localizer, or continuum theorem.  At the first predeclared
RPESM gauge-history card:
\begin{enumerate}[label=\textup{(V14.\arabic*)},leftmargin=6.2em]
\item freeze the link- or site-reflection geometry, the boundary gauge frame,
      \(H_{W,n}(U)\), its protected gap interval, and the Berezin sign before
      reading the crossing;
\item evaluate \(\operatorname{sgn}(H_{W,n}(U))\) by exact finite spectral
      calculus or a certified polynomial/rational approximation and populate
      \eqref{eq:V13-active-stop}--\eqref{eq:V13-active-Gram};
\item run the already proved Version~V9 relation, skew, negative-part, and
      positive-factor-rank tests with outward errors;
\item if the card lies in the open neighbourhood
      \eqref{eq:V13-open-gauge-condition}, certify the positive floor directly;
      otherwise evaluate the cone interval without replacing the action by its
      nearest positive matrix;
\item retain the determinant and Pfaffian line, divisor, orientation, and one
      held-out four-phase crossing Read on the same history;
\item repeat only this finite packet at one adjacent cutoff.  If its line or
      cone interval is the first unresolved row, stop there.  Only a passing
      card returns to the Version~V8 Hamiltonian tangent and Versions~V5--V6
      Duhamel/weighted-absorption programme.
\end{enumerate}
A free or gauge-near-free control may be recorded as a finite positive branch,
but it is not substituted for the selected active gauge history.
\end{openproblem}

\section*{V13 exact replay and append-only record}
\addcontentsline{toc}{section}{V13 exact replay and append-only record}
The accompanying standard-library verifier uses exact rational arithmetic for
\Cref{cth:V13-two-by-two}.  It checks the two Wilson kernels, their common gap,
sign matrices, overlap projectors, chronological effects and leakage, equality
in the sign-commutator bound, common massive coefficient spectrum and
determinant, opposite crossing signs, the positive factor rank, and the
negative cone distance.  It also verifies the canonical binary Naimark
projection.  Ordinary and optimized Python executions produce byte-identical
JSON output.

The complete Version~V12 source preceding its sole final active
document-closing command is preserved byte for byte.  Its preterminal
SHA--256 digest is
\begin{center}\scriptsize\ttfamily
7c3866577fcb4d41dca7f87eca1b417ff8a01779329f2be76c452990bba61a2c.
\end{center}
The exact verifier and its JSON output have SHA--256 digests
\begin{center}\scriptsize\ttfamily
571ecae6f3c39c966edd421693a8f04f54efffac4443ab2c65a2c256243b70b7,\\[-0.1em]
97a7ac29b420c1356ac58ee3355c91c69fa43d60442cbafab6fa9013235a80fd.
\end{center}
A later continuation must remove only the sole final active
\verb|\end{document}|, append new material immediately before it, restore
one terminal command, and increment the filename to end in
\texttt{\_V14.tex}.

% =============================================================================
% END APPENDED VERSION V13 MATERIAL
% =============================================================================


% =============================================================================
% START APPENDED VERSION V14 MATERIAL -- complete V1--V13 preterminal source preserved
% =============================================================================

\clearpage
\hypersetup{
 pdftitle={Renewal Einstein--Standard--Model Physical Source Metric, Duhamel Ancestry and Quantum Continuum Programme V14},
 pdfsubject={Balanced gauge-active Wilson collars, polar reflection incidence, and an exact SU(3) overlap crossing card}
}
\section*{Version V14: balanced gauge-active Wilson collars and the polar reflection criterion}
\addcontentsline{toc}{section}{V14: balanced Wilson collars and polar reflection incidence}

\begin{abstract}
Version~V13 reduced the first active fermionic gate to the cross-boundary block
of the gapped Wilson sign, its moving overlap-source Gram, the physical
reflection, and the already established one-letter relation and positive-cone
tests.  The free overlap theorem populated one control card and a small
neighbourhood of it, but did not identify a nontrivial gauge-active family on
which the crossing could be evaluated in closed form.

This continuation supplies such a family.  On a balanced boundary collar the
Hermitian Wilson kernel has diagonal blocks \(-mI\) and \(mI\), while its
cross-plane hopping block is an arbitrary finite matrix \(B\).  Its square is
block diagonal, so its sign is exact.  If \(B=VS\) is the polar decomposition,
the cross-boundary sign is \(-F_m(S)V^*\), where
\(F_m(S)=S(m^2I+S^2)^{-1/2}\).  After the fixed Hodge--chiral left factor is
included, the reflected crossing has the form
\[
 a^{-1}A_m C_\Theta,
 \qquad A_m=L F_m(S),
 \qquad C_\Theta=\varepsilon_\Theta V^*R_\Theta.
\]
The polar unitary of \(A_mC_\Theta\) is the product of the polar unitary of
\(A_m\) with \(C_\Theta\).  Consequently the crossing is positive if and only
if the physical reflection is the unique signed Hodge--polar transport of the
actual gauge hopping block.  On that branch the crossing floor is the least
singular value of \(L F_m(S)\), divided by the lattice scale.  No determinant,
covariance, hard-tail inverse, or fitted positive projection enters the proof.

The same calculation diagonalizes the chronological effect of the moving
overlap source.  Its nonzero mixing spectrum is
\(s^2/[4(m^2+s^2)]\), so the exact number of genuinely mixed chronological
source directions is \(\rank B\).  Directions in \(\ker B\) are purely
one-sided and carry no crossing source.  A supported version of the polar
criterion treats this singular branch without replacing the zero block by a
pseudoinverse Gaussian.

An exact rational replay uses a nonidentity real rotation in
\(SO(3)\subset SU(3)\), a redundant four-column one-letter synthesis, and one
held-out four-phase probe.  With the polar reflection the represented crossing
is \((4/5)I_3\), has factor rank three, and satisfies every word relation.
Keeping the same gauge link but using the identity reflection leaves all
Wilson singular values unchanged and produces a strict cone distance
\(512/625\).  A relation mutation is independently rejected with residual
\(32\).  A second exact collar gives a nonzero adjacent-card crossing defect,
so finite positivity is not confused with zero cutoff transport.

The result is a genuine gauge-active finite control and an executable
sufficient collar around it.  It is not a proof that the selected RPESM Wilson
history has balanced diagonal blocks or the required polar reflection.  The
first remaining physical Read is now the balance residual, polar-reflection
residual, and direct Wilson-sign crossing on one named gauge history.
\end{abstract}

\begin{assessment}[Why the balanced collar is the next upstream calculation]
\label{ass:V14-direction}
The Version~V13 stop asks for one gauge-history Wilson-sign boundary block.
Repeating the free momentum-space factorization, enlarging the exterior word
bank, or returning to the Duhamel compiler would leave that coefficient
unread.  The balanced collar instead isolates the earliest finite data which
can decide the sign: the actual cross-plane Wilson hopping, its same-side mass
balance, and the physical reflection transport.  These are components of the
already required time-split precision and introduce no new programme or
source space.  The calculation below imports the Version~V9 relation/cone
compiler and the Version~V13 sign perturbation theorem without reproving them.
\end{assessment}

% =============================================================================
\section{Exact sign of a balanced Wilson boundary collar}
\label{sec:V14-balanced-sign}
% =============================================================================

Let \(\mathcal E_-\) and \(\mathcal E_+\) be finite-dimensional complex
Hilbert spaces of the same dimension \(d\).  They are a predeclared
negative/positive-time boundary carrier, not the full lattice Hilbert space.
Let
\[
 T_- = \begin{pmatrix}I&0\\0&0\end{pmatrix},
 \qquad
 T_+ = \begin{pmatrix}0&0\\0&I\end{pmatrix}
\]
relative to \(\mathcal E_-\oplus\mathcal E_+\).  Fix \(m>0\) and a linear map
\(B:\mathcal E_+\to\mathcal E_-\).  Define the balanced Wilson collar
\begin{equation}
 \boxed{
 H_{m,B}
 =\begin{pmatrix}
   -mI_{\mathcal E_-}&-B\\
   -B^*&mI_{\mathcal E_+}
  \end{pmatrix}.}
 \label{eq:V14-balanced-H}
\end{equation}
Write the polar decomposition as
\begin{equation}
 B=VS,
 \qquad S=(B^*B)^{1/2},
 \qquad P_S=V^*V,
 \qquad P_R=VV^*.
 \label{eq:V14-polar-B}
\end{equation}
Here \(V\) is the canonical partial isometry, from \(\Ran P_S\) onto
\(\Ran P_R\).  Put
\begin{equation}
 R_-=(m^2I+BB^*)^{1/2},
 \qquad
 R_+=(m^2I+B^*B)^{1/2},
 \qquad
 F_m(S)=S(m^2I+S^2)^{-1/2}.
 \label{eq:V14-RF}
\end{equation}

\begin{theorem}[Exact balanced-collar sign]
\label{thm:V14-balanced-sign}
The collar is invertible and
\begin{equation}
 \boxed{
 H_{m,B}^2
 =\begin{pmatrix}R_-^2&0\\0&R_+^2\end{pmatrix}.}
 \label{eq:V14-H-square}
\end{equation}
Consequently
\begin{equation}
 \boxed{
 \operatorname{sgn}H_{m,B}
 =\begin{pmatrix}
   -mR_-^{-1}&-BR_+^{-1}\\
   -B^*R_-^{-1}&mR_+^{-1}
  \end{pmatrix},}
 \label{eq:V14-sign-blocks}
\end{equation}
and its negative-to-positive chronological block is
\begin{equation}
 \boxed{
 T_+\operatorname{sgn}(H_{m,B})T_-
 =-F_m(S)V^*.}
 \label{eq:V14-sign-cross-block}
\end{equation}
The nonzero singular values of this block are exactly
\begin{equation}
 \boxed{
 f_m(s_j)=\frac{s_j}{\sqrt{m^2+s_j^2}},}
 \label{eq:V14-cross-singular-values}
\end{equation}
where \(s_j>0\) run through the nonzero singular values of \(B\).
\end{theorem}

\begin{proof}
Block multiplication in \eqref{eq:V14-balanced-H} gives
\[
 H_{m,B}^2
 =\begin{pmatrix}
   m^2I+BB^*&mB-Bm\\
   mB^*-B^*m&m^2I+B^*B
  \end{pmatrix},
\]
and the off-diagonal blocks vanish because \(m\) is scalar.  Since
\(m>0\), both diagonal blocks are strictly positive, proving invertibility and
\eqref{eq:V14-H-square}.  Functional calculus gives
\[
 |H_{m,B}|^{-1}=\diag(R_-^{-1},R_+^{-1}),
\]
so multiplication by \(H_{m,B}\) proves \eqref{eq:V14-sign-blocks}.

The polar identities
\[
 R_-=V R_+V^*+m(I-P_R),
 \qquad
 B^*R_-^{-1}=S(m^2I+S^2)^{-1/2}V^*
\]
hold on \(\Ran P_R\), while both sides vanish on its orthogonal complement.
This proves \eqref{eq:V14-sign-cross-block}.  Finally, \(V^*\) is an isometry
from \(\Ran P_R\) onto \(\Ran P_S\), and \(F_m(S)\) has eigenvalues
\(f_m(s_j)\) on the positive singular subspaces and zero on \(\ker B\).
\end{proof}

\begin{remark}[What is balanced and what is gauge active]
The adjective \emph{balanced} concerns only the two same-side scalar blocks
\(-mI\) and \(mI\).  The cross-plane matrix \(B\) is arbitrary and may carry
a nontrivial gauge link, spin transport, flavour mixing, or boundary route.
Thus \Cref{thm:V14-balanced-sign} is not a gauge-trivial Fourier calculation.
The active RPESM card must still show that its actual same-side blocks are in
this collar, or are close enough to it for the perturbation theorem below.
\end{remark}

% =============================================================================
\section{The physical reflection is the polar completion of the crossing}
\label{sec:V14-polar-reflection}
% =============================================================================

First assume that \(B\) is invertible, so \(V:\mathcal E_+\to\mathcal E_-\)
is unitary.  Let
\begin{equation}
 L=L^*>0
 \label{eq:V14-L}
\end{equation}
be the fixed Hodge--chiral left factor on the selected positive-time boundary
carrier.  In the literal overlap card this is the restriction of
\(T_+\star_0P_+T_+\) after its support has been verified.  Let
\(R_\Theta:\mathcal E_+\to\mathcal E_-\) be the unitary linear part of the
physical reflection and let \(\varepsilon_\Theta\in\{+1,-1\}\) be the frozen
Berezin-order sign.  By \eqref{eq:V13-sign-crossing} and
\eqref{eq:V14-sign-cross-block}, the represented boundary crossing before the
one-letter synthesis is
\begin{equation}
 \boxed{
 K_\times(m,B;L,R_\Theta)
 =\frac{\varepsilon_\Theta}{a}
   L F_m(S)V^*R_\Theta.}
 \label{eq:V14-K-general}
\end{equation}
Put
\begin{equation}
 A_m=L F_m(S)=U_m|A_m|,
 \qquad
 C_\Theta=\varepsilon_\Theta V^*R_\Theta,
 \label{eq:V14-A-C}
\end{equation}
where \(U_m\) is the polar unitary of the invertible matrix \(A_m\).

\begin{theorem}[Hodge--polar reflection criterion]
\label{thm:V14-polar-criterion}
The crossing \eqref{eq:V14-K-general} is Hermitian positive definite if and
only if
\begin{equation}
 \boxed{
 C_\Theta=U_m^*
 \quad\Longleftrightarrow\quad
 R_\Theta=\varepsilon_\Theta VU_m^*.}
 \label{eq:V14-polar-reflection}
\end{equation}
On this branch,
\begin{equation}
 \boxed{
 K_\times
 =\frac1a U_m|A_m|U_m^*\succeq
 \kappa_{m,B,L}I,}
 \label{eq:V14-positive-crossing}
\end{equation}
where
\begin{equation}
 \boxed{
 \kappa_{m,B,L}
 =\frac1a\,s_{\min}(L F_m(S))
 \ge
 \frac{\lambda_{\min}(L)}{a}
 \frac{s_{\min}(B)}{\sqrt{m^2+s_{\min}(B)^2}}.}
 \label{eq:V14-crossing-floor}
\end{equation}

Equivalently, the positive branch is the zero set of the polar-incidence
residual
\begin{equation}
 \boxed{
 \Delta_{\rm pol}^{\times}
 =\|\varepsilon_\Theta V^*R_\Theta-U_m^*\|_{\rm HS}^2.}
 \label{eq:V14-polar-residual}
\end{equation}
No positive determinant modulus or singular-value list determines this
residual.
\end{theorem}

\begin{proof}
The matrix \(C_\Theta\) is unitary and
\[
 K_\times=a^{-1}A_mC_\Theta.
\]
Since \(A_m=U_m|A_m|\), one has the polar factorization
\[
 A_mC_\Theta
 =(U_mC_\Theta)(C_\Theta^*|A_m|C_\Theta).
\]
The first factor is unitary and the second is strictly positive.  Hence the
polar unitary of \(A_mC_\Theta\) is \(U_mC_\Theta\).  An invertible matrix is
Hermitian positive definite exactly when its polar unitary is the identity.
This proves \eqref{eq:V14-polar-reflection}.  Substitution gives
\[
 A_mU_m^*=U_m|A_m|U_m^*,
\]
which is positive and has least eigenvalue \(s_{\min}(A_m)\).  Finally,
\[
 \|(LF_m(S))^{-1}\|_{\rm op}
 \le\|F_m(S)^{-1}\|_{\rm op}\|L^{-1}\|_{\rm op},
\]
while \(f_m(s)\) is increasing for \(s>0\).  This proves the lower estimate
in \eqref{eq:V14-crossing-floor}.  The residual statement is merely the
norm-square form of the same equality.
\end{proof}

\begin{corollary}[Commuting Hodge factor and the ordinary polar reflection]
\label{cor:V14-commuting-polar}
If \([L,S]=0\), then \(A_m=LF_m(S)>0\), so \(U_m=I\), and the criterion
reduces to
\begin{equation}
 \boxed{R_\Theta=\varepsilon_\Theta V.}
 \label{eq:V14-simple-polar-reflection}
\end{equation}
The positive crossing is
\begin{equation}
 \boxed{K_\times=a^{-1}L F_m(S).}
 \label{eq:V14-simple-positive-crossing}
\end{equation}
In particular, a scalar Hodge factor does not erase the gauge-link polar
transport.
\end{corollary}

\begin{proof}
Commuting positive matrices have a positive product.  Its polar unitary is the
identity, so \Cref{thm:V14-polar-criterion} applies.
\end{proof}

\begin{firewall}[The reflection is not fitted to repair the action]
\label{fire:V14-no-fit-reflection}
Equation~\eqref{eq:V14-polar-reflection} is a test of an independently declared
physical reflection against the actual cross-plane Wilson hopping.  It is not
permission to replace the submitted reflection after the crossing has been
read.  If the physical \(R_\Theta\) fails the residual, the mismatch is a
finite crossing obstruction unless an independently occurring boundary frame
or action modification supplies a new common-history card.
\end{firewall}

% =============================================================================
\section{Singular hopping support and the zero-safe polar branch}
\label{sec:V14-singular-polar}
% =============================================================================

The full-rank theorem is the relevant branch for a strict one-particle
crossing floor.  A boundary hopping block may nevertheless have a kernel.
The following supported theorem retains it rather than replacing it by a
pseudoinverse Gaussian.  For clarity, take \(L=\ell I\), \(\ell>0\).

\begin{theorem}[Supported polar criterion]
\label{thm:V14-supported-polar}
Let \(B\) be arbitrary, retain the partial polar decomposition
\eqref{eq:V14-polar-B}, and let \(R_\Theta:\mathcal E_+\to\mathcal E_-\) be
unitary.  Put
\begin{equation}
 C_\Theta=\varepsilon_\Theta V^*R_\Theta.
 \label{eq:V14-supported-C}
\end{equation}
Then
\begin{equation}
 C_\Theta C_\Theta^*=P_S,
 \qquad P_SC_\Theta=C_\Theta,
 \label{eq:V14-C-coisometry}
\end{equation}
and the reflected crossing is
\begin{equation}
 K_\times=\frac\ell aF_m(S)C_\Theta.
 \label{eq:V14-supported-K}
\end{equation}
It is Hermitian positive semidefinite if and only if
\begin{equation}
 \boxed{
 C_\Theta=P_S.}
 \label{eq:V14-supported-polar-condition}
\end{equation}
Equivalently,
\begin{equation}
 \boxed{
 R_\Theta x=\varepsilon_\Theta Vx
 \quad(x\in\Ran P_S),
 \qquad
 R_\Theta(\ker B)\subseteq\ker B^*.}
 \label{eq:V14-supported-geometric-condition}
\end{equation}
On this branch,
\begin{equation}
 \boxed{
 K_\times=\frac\ell aF_m(S),
 \qquad
 \rank K_\times=\rank B.}
 \label{eq:V14-supported-positive}
\end{equation}
If \(B\ne0\), its supported floor is
\begin{equation}
 \boxed{
 \lambda_{\min}^{+}(K_\times)
 =\frac\ell a
   \frac{s_{\min}^{+}(B)}
        {\sqrt{m^2+(s_{\min}^{+}(B))^2}}.}
 \label{eq:V14-supported-floor}
\end{equation}
\end{theorem}

\begin{proof}
The polar identities give
\[
 C_\Theta C_\Theta^*
 =V^*R_\Theta R_\Theta^*V=V^*V=P_S,
 \qquad P_SC_\Theta=C_\Theta.
\]
Suppose first that \(K_\times\succeq0\).  Since it is then self-adjoint,
\[
 K_\times^2=K_\times K_\times^*
 =\frac{\ell^2}{a^2}
  F_m(S)C_\Theta C_\Theta^*F_m(S)
 =\frac{\ell^2}{a^2}F_m(S)^2.
\]
Uniqueness of the positive square root gives
\(K_\times=(\ell/a)F_m(S)\), and hence
\[
 F_m(S)(C_\Theta-P_S)=0.
\]
The range of \(C_\Theta-P_S\) lies in \(\Ran P_S\), where \(F_m(S)\) is
injective.  Therefore \(C_\Theta=P_S\).  The converse follows by direct
substitution.

The operator equality \(\varepsilon_\Theta V^*R_\Theta=P_S\) says that
\(R_\Theta\) agrees with \(\varepsilon_\Theta V\) on the initial polar
support.  On its orthogonal complement it maps into \(\ker V^*=\ker B^*\),
which proves \eqref{eq:V14-supported-geometric-condition}; the converse is
immediate.  Rank and the supported floor follow from the spectrum of
\(F_m(S)\).
\end{proof}

\begin{assessment}[Meaning of the singular branch]
\label{ass:V14-singular-meaning}
A vector in \(\ker B\) carries no cross-plane hopping and no positive crossing
floor.  The theorem neither deletes that vector nor declares it a negative
mode.  It must remain in the same-side or divisor/zero-mode packet of the
complete coefficient.  This is consistent with the existing zero-safe
Grassmann and determinant-line hierarchy: a pseudoinverse of the missing
crossing is not a substitute for an oriented zero-mode insertion.
\end{assessment}

% =============================================================================
\section{Exact spectrum of the moving chronological source}
\label{sec:V14-chronological-spectrum}
% =============================================================================

Let
\begin{equation}
 P_{m,B}=\frac12(I+\operatorname{sgn}H_{m,B}),
 \qquad
 E_\chi=\Ran P_{m,B},
 \label{eq:V14-positive-projector}
\end{equation}
and let \(J_\chi:E_\chi\hookrightarrow\mathcal E_-\oplus\mathcal E_+\) be
inclusion.  The positive-time one-letter synthesis and its effect are
\begin{equation}
 W_{\chi,+}=T_+J_\chi,
 \qquad
 A_+=J_\chi^*T_+J_\chi.
 \label{eq:V14-Aplus}
\end{equation}

\begin{theorem}[Chronological effect spectrum and exact mixed rank]
\label{thm:V14-A-spectrum}
The effect \(A_+\) is unitarily equivalent to
\begin{equation}
 \boxed{
 E_m(S)
 =\frac12\left(I+m(m^2I+S^2)^{-1/2}\right)
 \quad\text{on }\mathcal E_+.}
 \label{eq:V14-E-spectrum}
\end{equation}
Consequently, if \(s_j\) are all singular values of \(B\), including zeros,
then
\begin{equation}
 \boxed{
 \spec A_+
 =\left\{
   \frac12\left(1+\frac{m}{\sqrt{m^2+s_j^2}}\right)
  :1\le j\le d
  \right\}.}
 \label{eq:V14-A-eigenvalues}
\end{equation}
Its exact chronological mixing stock is
\begin{equation}
 \boxed{
 \begin{aligned}
 \Delta_{\chi/t}
 &=\Tr(A_+-A_+^2)\\
 &=\frac14\Tr\left[S^2(m^2I+S^2)^{-1}\right]\\
 &=\frac14\sum_{j=1}^{d}
   \frac{s_j^2}{m^2+s_j^2}.
 \end{aligned}}
 \label{eq:V14-exact-leakage}
\end{equation}
Moreover,
\begin{equation}
 \boxed{
 \rank(A_+(I-A_+))=\rank B.}
 \label{eq:V14-mixed-rank}
\end{equation}
Every vector in \(\ker B\) corresponds to an eigenvalue one of \(A_+\): it is
purely positive-time and has zero crossing stock.
\end{theorem}

\begin{proof}
Let
\[
 X=T_+J_\chi:E_\chi\longrightarrow\mathcal E_+.
\]
Then \(X^*X=A_+\), while
\[
 XX^*=T_+P_{m,B}T_+\big|_{\mathcal E_+}.
\]
The lower-right block of \eqref{eq:V14-sign-blocks} gives
\[
 XX^*=\frac12(I+mR_+^{-1})=E_m(S).
\]
The latter is strictly positive because \(m>0\).  A singular-value basis for
\(B\) decomposes \(H_{m,B}\) into \(d\) two-dimensional blocks
\(\begin{psmallmatrix}-m&-s_j\\-s_j&m\end{psmallmatrix}\), each with one
positive and one negative eigenvalue.  Hence \(\dim E_\chi=d\).  Source and
target of \(X\) have the same dimension, and \(XX^*>0\), so \(X\) is
invertible.  The polar decomposition of \(X\) therefore makes \(X^*X\) and
\(XX^*\) unitarily equivalent, proving
\eqref{eq:V14-E-spectrum}--\eqref{eq:V14-A-eigenvalues}.

For a singular value \(s\), write
\[
 e(s)=\frac12\left(1+\frac{m}{\sqrt{m^2+s^2}}\right).
\]
Then
\[
 e(s)(1-e(s))
 =\frac{s^2}{4(m^2+s^2)}.
\]
Summation proves \eqref{eq:V14-exact-leakage}.  This number is positive
exactly for \(s>0\), which gives \eqref{eq:V14-mixed-rank}.  At \(s=0\),
\(e(s)=1\), proving the final assertion.
\end{proof}

\begin{corollary}[Sharp source visibility window]
\label{cor:V14-A-window}
Let \(s_{\max}=\|B\|_{\rm op}\) and
\(s_{\min}=\min\spec S\).  Then
\begin{equation}
 \boxed{
 \frac12\left(1+\frac{m}{\sqrt{m^2+s_{\max}^2}}\right)I
 \preceq A_+
 \preceq
 \frac12\left(1+\frac{m}{\sqrt{m^2+s_{\min}^2}}\right)I.}
 \label{eq:V14-A-window}
\end{equation}
The upper endpoint equals one precisely when \(B\) has a kernel.  Thus the
positive-time source synthesis is injective on every balanced massive collar,
while strict two-sided chronological mixing occurs exactly on the hopping
support.
\end{corollary}

\begin{proof}
The scalar function in \eqref{eq:V14-A-eigenvalues} decreases with \(s\ge0\).
Apply finite spectral calculus.
\end{proof}

% =============================================================================
\section{Gauge covariance and automatic one-letter relation descent}
\label{sec:V14-gauge-relation}
% =============================================================================

\begin{theorem}[Gauge covariance of the polar card]
\label{thm:V14-gauge-covariance}
Let \(g_\pm\) be unitary boundary gauge-frame changes and transform
\begin{equation}
 \begin{aligned}
 B'&=g_-Bg_+^*,
 &L'&=g_+Lg_+^*,\\
 R_\Theta'&=g_-R_\Theta g_+^*,
 &H_{m,B}'&=(g_-\oplus g_+)H_{m,B}(g_-\oplus g_+)^*.
 \end{aligned}
 \label{eq:V14-gauge-transform}
\end{equation}
Then
\begin{equation}
 \boxed{
 S'=g_+Sg_+^*,
 \quad V'=g_-Vg_+^*,
 \quad U_m'=g_+U_mg_+^*,
 \quad K_\times'=g_+K_\times g_+^*.}
 \label{eq:V14-gauge-covariance}
\end{equation}
In particular, the polar-incidence residual, inertia, rank, cone distance,
source-effect spectrum, and positive floor are gauge-frame invariant.
\end{theorem}

\begin{proof}
The first two identities are uniqueness of the positive square root and polar
partial isometry.  They imply
\(F_m(S')=g_+F_m(S)g_+^*\) and
\(A_m'=g_+A_mg_+^*\).  Uniqueness of polar decomposition gives the formula for
\(U_m'\).  Substitution in \eqref{eq:V14-K-general} proves the crossing
congruence.  Every listed quantity is invariant under unitary conjugation.
\end{proof}

Let \(W_1:\mathcal F_1\to\mathcal E_+\) be the already occurring one-letter
synthesis and put
\begin{equation}
 G_1=W_1^*W_1,
 \qquad
 P_W=W_1^\dagger W_1,
 \qquad
 K_F=W_1^*K_\times W_1.
 \label{eq:V14-word-data}
\end{equation}

\begin{theorem}[The balanced crossing descends through every actual one-letter relation]
\label{thm:V14-relation-descent}
For every \(K_\times\), not only on the positive branch,
\begin{equation}
 \boxed{K_F=P_WK_FP_W.}
 \label{eq:V14-automatic-relation}
\end{equation}
Hence the Version~V9 two-sided relation residual vanishes exactly.  On the
supported polar branch of \Cref{thm:V14-supported-polar},
\begin{equation}
 \boxed{
 K_F=\frac\ell aW_1^*F_m(S)W_1
 \succeq
 \kappa_{m,B}^{+}
 W_1^*P_SW_1,}
 \label{eq:V14-word-floor}
\end{equation}
where \(\kappa_{m,B}^{+}\) is the scalar in
\eqref{eq:V14-supported-floor}.  Moreover,
\begin{equation}
 \boxed{\rank K_F=\rank(P_SW_1).}
 \label{eq:V14-word-rank}
\end{equation}
If \(B\) is invertible, then
\begin{equation}
 \boxed{
 K_F\succeq\kappa_{m,B,L}G_1,
 \qquad
 \rank K_F=\rank W_1,}
 \label{eq:V14-full-word-floor}
\end{equation}
and the represented minimum-norm extension in Version~V9 is a positive
compression of \(K_\times\).
\end{theorem}

\begin{proof}
Since \(W_1=W_1P_W\) and \(W_1^*=P_WW_1^*\),
\eqref{eq:V14-automatic-relation} is immediate.  On the supported positive
branch, functional calculus gives
\[
 F_m(S)\succeq
 \frac{s_{\min}^{+}(B)}
      {\sqrt{m^2+(s_{\min}^{+}(B))^2}}P_S.
\]
Congruence by \(W_1\) proves \eqref{eq:V14-word-floor}.  Since
\(F_m(S)^{1/2}\) is invertible on \(\Ran P_S\),
\[
 K_F=\frac\ell a
 (F_m(S)^{1/2}W_1)^*(F_m(S)^{1/2}W_1)
\]
has rank \(\rank(P_SW_1)\).  The full-rank formulas are the special case
\(P_S=I\), with \Cref{thm:V14-polar-criterion}.  Finally,
\[
 (W_1^\dagger)^*K_FW_1^\dagger
 =P_{\Ran W_1}K_\times P_{\Ran W_1}\succeq0,
\]
which is exactly the Version~V9 represented extension.
\end{proof}

\begin{assessment}[No duplicate source price]
\label{ass:V14-no-duplicate-source}
The matrix \(K_F\) is the pullback of the same crossing through the already
occurring one-letter synthesis.  Its factor rank is therefore the number of
crossing-active represented one-letter directions, not an independent fermion
species count.  A direction in \(\ker W_1\), or in \(\Ran W_1\cap\ker B\),
is not repaired by adding another source; it is respectively a word relation
or a zero-crossing direction retained by the complete same-side/divisor
packet.
\end{assessment}

% =============================================================================
\section{Exact cone residual for an isotropic gauge hopping}
\label{sec:V14-isotropic-cone}
% =============================================================================

The polar criterion is qualitative for arbitrary \(L\) and \(S\).  When the
hopping singular values and the Hodge factor are scalar, the full Version~V9
cone distance is explicit.

\begin{theorem}[Unitary phase formula for the crossing-cone residual]
\label{thm:V14-unitary-phase}
Suppose
\begin{equation}
 B=\kappa V,
 \qquad L=\ell I,
 \qquad \kappa,\ell>0,
 \label{eq:V14-isotropic-data}
\end{equation}
and let the eigenvalues of
\begin{equation}
 Q=\varepsilon_\Theta V^*R_\Theta
 \label{eq:V14-Q}
\end{equation}
be \(e^{i\theta_1},\ldots,e^{i\theta_d}\).  Put
\begin{equation}
 f=\frac{\ell\kappa}{a\sqrt{m^2+\kappa^2}}.
 \label{eq:V14-f}
\end{equation}
Then \(K_\times=fQ\), and its Version~V9 cone residual is
\begin{equation}
 \boxed{
 \Delta_{\rm hop}^{\times}(K_\times)
 =f^2\sum_{j=1}^{d}
 \left[
  \sin^2\theta_j+\bigl(\max\{-\cos\theta_j,0\}\bigr)^2
 \right].}
 \label{eq:V14-phase-cone}
\end{equation}
Furthermore,
\begin{equation}
 \boxed{
 \Delta_{\rm pol}^{\times}
 =\|Q-I\|_{\rm HS}^2
 =2\sum_{j=1}^{d}(1-\cos\theta_j).}
 \label{eq:V14-phase-polar}
\end{equation}
Both residuals vanish exactly when \(Q=I\).
\end{theorem}

\begin{proof}
Here \(F_m(S)=\kappa(m^2+\kappa^2)^{-1/2}I\), so
\eqref{eq:V14-K-general} gives \(K_\times=fQ\).  Diagonalize the unitary
\(Q\).  On the eigenline with angle \(\theta\), the anti-Hermitian part has
squared modulus \(f^2\sin^2\theta\), while the Hermitian part has eigenvalue
\(f\cos\theta\) and contributes \(f^2(\max\{-\cos\theta,0\})^2\) to its
negative spectral part.  Summing proves \eqref{eq:V14-phase-cone}.  The second
formula is the ordinary Hilbert--Schmidt norm of \(Q-I\).  Each summand in
\eqref{eq:V14-phase-cone} vanishes only at \(\theta=0\) modulo \(2\pi\).
\end{proof}

\begin{corollary}[Positive singular values do not fix the reflection branch]
\label{cor:V14-singular-values-no-branch}
For fixed \(m,\kappa,\ell,a\), every physical reflection has the same crossing
singular values, all equal to \(f\), while the cone residual ranges from zero
to a strictly positive value.  Thus a Wilson gap, a crossing singular-value
list, and a determinant modulus do not determine the reflected-hopping branch.
\end{corollary}

\begin{proof}
Right multiplication by the unitary \(Q\) leaves the singular values of
\(fI\) unchanged.  The residual variation follows from
\eqref{eq:V14-phase-cone}.
\end{proof}

% =============================================================================
\section{A gauge-active collar around the exact polar card}
\label{sec:V14-polar-collar}
% =============================================================================

The exact family can serve as a reference for a general Wilson history whose
same-side blocks are not exactly scalar.  Freeze the physical mass parameter,
reflection, Hodge factor, and cross-plane block before evaluating the
residual.

Let \(\widetilde H_W=\widetilde H_W^*\) be a submitted Wilson kernel on the
same boundary carrier.  Define its actual cross block and the balanced
reference by
\begin{equation}
 B=-T_-\widetilde H_WT_+,
 \qquad
 H_{m,B}\text{ as in \eqref{eq:V14-balanced-H}},
 \qquad
 E_{\rm bal}=\widetilde H_W-H_{m,B}.
 \label{eq:V14-balance-decomposition}
\end{equation}
Then \(E_{\rm bal}\) is block diagonal and
\begin{equation}
 \boxed{
 \|E_{\rm bal}\|_{\rm HS}^2
 =\|T_-\widetilde H_WT_-+mT_-\|_{\rm HS}^2
  +\|T_+\widetilde H_WT_+-mT_+\|_{\rm HS}^2.}
 \label{eq:V14-balance-residual}
\end{equation}

\begin{theorem}[Finite gauge-active polar-collar certificate]
\label{thm:V14-polar-collar}
Assume that \(B\) is invertible, that the physical reflection satisfies the
polar criterion \eqref{eq:V14-polar-reflection}, and that
\begin{equation}
 \dist\bigl(0,\spec(H_{m,B}+sE_{\rm bal})\bigr)
 \ge g_{\rm col}>0
 \quad(0\le s\le1).
 \label{eq:V14-collar-gap}
\end{equation}
Let \(\widetilde K_\times\) be the exact overlap crossing formed from
\(\operatorname{sgn}\widetilde H_W\) with the same
\(L,R_\Theta,\varepsilon_\Theta\).  Then
\begin{equation}
 \boxed{
 \|\widetilde K_\times-K_\times^{\rm bal}\|_{\rm HS}
 \le
 \eta_{\rm col}
 :=\frac{\|L\|_{\rm op}}{a g_{\rm col}}
   \|E_{\rm bal}\|_{\rm HS}.}
 \label{eq:V14-collar-radius}
\end{equation}
If reflection covariance independently proves
\(\widetilde K_\times=\widetilde K_\times^*\) and
\begin{equation}
 \boxed{\eta_{\rm col}<\kappa_{m,B,L},}
 \label{eq:V14-collar-pass}
\end{equation}
then
\begin{equation}
 \boxed{
 \widetilde K_\times
 \succeq(\kappa_{m,B,L}-\eta_{\rm col})I>0.}
 \label{eq:V14-collar-floor}
\end{equation}
For every one-letter synthesis \(W_1\),
\begin{equation}
 \boxed{
 W_1^*\widetilde K_\times W_1
 \succeq
 (\kappa_{m,B,L}-\eta_{\rm col})G_1.}
 \label{eq:V14-collar-word-floor}
\end{equation}
If either the polar residual or the strict collar margin fails, the exact
Version~V13 sign card remains executable; failure of this sufficient collar is
not itself a negative crossing theorem.
\end{theorem}

\begin{proof}
Apply \Cref{thm:V13-sign-path} to the path
\(H_{m,B}+sE_{\rm bal}\).  Formula~\eqref{eq:V13-sign-crossing}, unitarity of
the reflection, and the operator norm of the fixed left factor give
\eqref{eq:V14-collar-radius}.  The balanced crossing has floor
\(\kappa_{m,B,L}\) by \Cref{thm:V14-polar-criterion}.  On the independently
Hermitian branch, Weyl's inequality and
\(\|X\|_{\rm op}\le\|X\|_{\rm HS}\) give
\eqref{eq:V14-collar-floor}.  Congruence by \(W_1\) proves the final formula.
\end{proof}

\begin{assessment}[Advance over the free-near-free window]
\label{ass:V14-advance}
The Version~V13 perturbation centre was the free translation-invariant overlap
operator.  The centre here may contain an arbitrary invertible gauge hopping
block \(B\).  The price is explicit: its physical reflection must match the
signed Hodge--polar transport, and the actual same-side Wilson blocks must lie
inside the finite balance radius.  This is a larger gauge-active family, but
it remains a sufficient finite collar rather than a theorem for every gauge
history.
\end{assessment}

% =============================================================================
\section{Exact nontrivial \(SU(3)\) replay with relations and a held-out Read}
\label{sec:V14-SU3-replay}
% =============================================================================

Take \(\mathcal E_-=\mathcal E_+=\C^3\), \(a=\ell=\kappa=1\),
\(m=3/4\), and
\begin{equation}
 \boxed{
 V=
 \begin{pmatrix}
  3/5&-4/5&0\\
  4/5& 3/5&0\\
  0&0&1
 \end{pmatrix}\in SO(3)\subset SU(3),
 \qquad B=V.}
 \label{eq:V14-SU3-V}
\end{equation}
This is a nonidentity gauge-link card.  The balanced Wilson kernel is
\begin{equation}
 H=\begin{pmatrix}-(3/4)I&-V\\-V^*&(3/4)I\end{pmatrix}.
 \label{eq:V14-control-H}
\end{equation}

\begin{proposition}[Exact positive gauge-active control]
\label{prop:V14-SU3-positive}
For \eqref{eq:V14-control-H},
\begin{equation}
 \boxed{
 H^2=\frac{25}{16}I_6,
 \qquad
 \operatorname{sgn}H=\frac45H.}
 \label{eq:V14-control-sign}
\end{equation}
With \(\varepsilon_\Theta=+1\) and the polar reflection
\(R_\Theta=V\),
\begin{equation}
 \boxed{
 K_\times=\frac45I_3,
 \qquad
 A_+\simeq\frac45I_3,
 \qquad
 \Delta_{\chi/t}=\frac{12}{25}.}
 \label{eq:V14-control-values}
\end{equation}
The source-minimal one-sided factor is
\begin{equation}
 \boxed{R_\times=\frac{2}{\sqrt5}I_3,}
 \label{eq:V14-control-factor}
\end{equation}
and its rank is three.
\end{proposition}

\begin{proof}
Orthogonality of \(V\) and \Cref{thm:V14-balanced-sign} give the square and
sign identities.  Since \(L=I\), \Cref{cor:V14-commuting-polar} gives
\(K_\times=F_m(I)=4I/5\).  Formula~\eqref{eq:V14-A-eigenvalues} gives the same
value for \(A_+\), and \eqref{eq:V14-exact-leakage} gives
\(3(4/5)(1/5)=12/25\).  The positive square root of \(4I/5\) is
\(2I/\sqrt5\), and strict positivity gives rank three.
\end{proof}

Use the redundant one-letter synthesis
\begin{equation}
 \boxed{
 W_1=
 \begin{pmatrix}
 1&0&0&1\\
 0&1&0&1\\
 0&0&1&1
 \end{pmatrix},
 \qquad
 v=(-1,-1,-1,1)^{\mathsf T}.}
 \label{eq:V14-control-W}
\end{equation}
Then \(W_1v=0\) and
\begin{equation}
 G_1=W_1^*W_1
 =\begin{pmatrix}
 1&0&0&1\\
 0&1&0&1\\
 0&0&1&1\\
 1&1&1&3
 \end{pmatrix}.
 \label{eq:V14-control-G}
\end{equation}

\begin{proposition}[Exact relation descent and mutation rejection]
\label{prop:V14-control-relation}
The raw one-letter crossing is
\begin{equation}
 \boxed{K_F=\frac45G_1.}
 \label{eq:V14-control-KF}
\end{equation}
Its two-sided Version~V9 relation residual is zero and its represented
extension is \((4/5)I_3\).  For the mutation
\begin{equation}
 K_F^{\rm bad}=K_F+vv^*,
 \label{eq:V14-bad-KF}
\end{equation}
the exact relation residual is
\begin{equation}
 \boxed{
 \|(I-P_W)K_F^{\rm bad}\|_{\rm HS}^2
 +\|K_F^{\rm bad}(I-P_W)\|_{\rm HS}^2
 =32.}
 \label{eq:V14-bad-relation}
\end{equation}
\end{proposition}

\begin{proof}
Equation~\eqref{eq:V14-control-KF} is direct congruence of
\((4/5)I_3\).  Since \(\ker W_1=\Span\{v\}\),
\[
 P_W=I_4-\frac14vv^*.
\]
The good matrix is supported on \(P_W\).  For the mutation,
\((I-P_W)vv^*=vv^*=vv^*(I-P_W)\).  Since \(v^*v=4\),
\(\|vv^*\|_{\rm HS}^2=16\).  The two sides therefore contribute
\(16+16=32\).  The represented-extension statement is
\Cref{thm:V14-relation-descent} and surjectivity of \(W_1\).
\end{proof}

\begin{proposition}[Held-out four-phase crossing Read]
\label{prop:V14-heldout}
Let \(e_1=(1,0,0)^{\mathsf T}\), put
\begin{equation}
 x=(0,e_1)\in\mathcal E_-\oplus\mathcal E_+,
 \qquad
 y=(Ve_1,0)\in\mathcal E_-\oplus\mathcal E_+,
 \label{eq:V14-heldout-vectors}
\end{equation}
and let \(q(z)=\langle z,(\operatorname{sgn}H)z\rangle\).  The four exact
phase probes are
\begin{equation}
 \boxed{
 \bigl(q(x+i^ry)\bigr)_{r=0}^{3}
 =\left(-\frac85,0,\frac85,0\right).}
 \label{eq:V14-four-q}
\end{equation}
Consequently the Version~V9 polarization gives
\begin{equation}
 \boxed{
 \frac14\sum_{r=0}^{3}i^{-r}q(x+i^ry)
 =-\frac45,}
 \label{eq:V14-four-polarized}
\end{equation}
and multiplication by the frozen crossing sign returns \(4/5\), exactly the
\((1,1)\) entry of \eqref{eq:V14-control-KF}.  The held-out value is therefore
assembled from the same Wilson sign and is not fitted by the Gram.
\end{proposition}

\begin{proof}
Equation~\eqref{eq:V14-sign-cross-block} gives
\(
 \langle e_1,T_+\operatorname{sgn}(H)T_-Ve_1\rangle=-4/5
\).
The complex polarization identity of Version~V9 then proves
\eqref{eq:V14-four-polarized}.  Direct substitution, retained in the exact
verifier, yields the four rational values in \eqref{eq:V14-four-q}.
\end{proof}

\begin{countertheorem}[The same gauge link with the wrong reflection]
\label{cth:V14-wrong-reflection}
Keep the complete control above but replace the physical reflection by
\(R_\Theta=I_3\).  Then
\begin{equation}
 \boxed{
 K_\times^{\rm mis}
 =\frac45V^*
 =\begin{pmatrix}
  12/25&16/25&0\\
  -16/25&12/25&0\\
  0&0&4/5
 \end{pmatrix}.}
 \label{eq:V14-mismatch-K}
\end{equation}
Its Hermitian part is strictly positive, but its anti-Hermitian part has
squared Hilbert--Schmidt norm
\begin{equation}
 \boxed{
 \Delta_{\rm hop}^{\times}(K_\times^{\rm mis})
 =\frac{512}{625}.}
 \label{eq:V14-mismatch-distance}
\end{equation}
Thus the card fails before any determinant or covariance calculation, despite
having exactly the same Wilson gap and crossing singular values as the passing
card.
\end{countertheorem}

\begin{proof}
Equation~\eqref{eq:V14-mismatch-K} is
\Cref{thm:V14-unitary-phase} with \(Q=V^*\).  The Hermitian part is
\(\diag(12/25,12/25,4/5)\), so it has no negative part.  The two nonzero skew
entries are \(16/25\) and \(-16/25\), yielding
\(2(16/25)^2=512/625\).
\end{proof}

\begin{proposition}[A nonzero two-card transport replay]
\label{prop:V14-two-card}
Keep \(B=R_\Theta=V\) and \(a=L=I\).  Compare the masses
\begin{equation}
 m_0=\frac34,
 \qquad
 m_1=\frac5{12}.
 \label{eq:V14-two-masses}
\end{equation}
The exact crossing floors and chronological effects are
\begin{equation}
 \boxed{
 \kappa_0=\frac45,
 \quad \kappa_1=\frac{12}{13},
 \qquad
 A_{+,0}=\frac45I_3,
 \quad A_{+,1}=\frac9{13}I_3.}
 \label{eq:V14-two-card-values}
\end{equation}
Hence the direct adjacent differences are
\begin{equation}
 \boxed{
 \|K_{\times,1}-K_{\times,0}\|_{\rm op}=\frac8{65},
 \qquad
 \|A_{+,1}-A_{+,0}\|_{\rm op}=\frac7{65}.}
 \label{eq:V14-two-card-defects}
\end{equation}
Finite positivity therefore does not imply exact cutoff transport.
\end{proposition}

\begin{proof}
For \(m_1=5/12\),
\(\sqrt{m_1^2+1}=13/12\).  The crossing value is its reciprocal
\(12/13\), while \eqref{eq:V14-A-eigenvalues} gives
\((1+5/13)/2=9/13\).  Subtraction yields \(8/65\) and \(7/65\).
\end{proof}

% =============================================================================
\section{Target-minimal physical card and revised stopping line}
\label{sec:V14-card}
% =============================================================================

\begin{acquisition}[Balanced/polar overlap crossing card]
\label{acq:V14-polar-card}
At one predeclared gauge history and chronological boundary screen, acquire
\begin{equation}
 \boxed{
 \begin{gathered}
 T_\pm,\quad m,\quad
 B=-T_-H_WT_+,\quad
 E_{\rm bal}=H_W-H_{m,B},\quad
 [g_{\rm col},M_*],\\
 B=VS,\quad
 L=T_+\star_0P_+T_+\\
 \text{on its verified support},\\
 R_\Theta,\quad\varepsilon_\Theta,\quad
 \Delta_{\rm pol}^{\times},\quad
 \eta_{\rm col},\\
 W_1,\quad G_1,\quad K_F,\quad
 \Delta_{\rm rel}^{\times},\quad
 \Delta_{\rm hop}^{\times},\\
 \text{one held-out four-phase Read and line/divisor/orientation data},\\
 \text{one adjacent copy of the same packet.}
 \end{gathered}}
 \label{eq:V14-acquisition}
\end{equation}
The polar and collar rows are sufficient shortcuts.  If either remains
inconclusive, the direct validated Wilson-sign matrix of Version~V13 is the
controlling value.  No nearest-positive projection, fitted reflection, or
pseudoinverse zero-mode Gaussian is inserted.
\end{acquisition}

\begin{theorem}[Finite decision order after the V14 reduction]
\label{thm:V14-decision}
A well-typed finite card is evaluated in the following order.
\begin{enumerate}[label=\textup{(D14.\arabic*)},leftmargin=5.5em]
\item Verify the common history, chronological split, Wilson gap, Hodge--chiral
      support, physical reflection, and Berezin sign.
\item Form the actual cross-plane block \(B\), its polar data, and the
      same-side balance residual before inspecting the hopping verdict.
\item If the exact polar residual vanishes and the strict collar inequality
      \eqref{eq:V14-collar-pass} holds, certify the positive crossing floor
      \eqref{eq:V14-collar-floor}.
\item Otherwise evaluate the direct Version~V13 sign crossing with outward
      error and run the Version~V9 relation, skew, and negative-part tests.
\item On a singular hopping support, retain the zero crossing in the
      same-side/divisor hierarchy and factor only the supported positive block.
\item Test the held-out four-phase Read and the determinant/Pfaffian line rows
      independently.  A passing crossing does not orient either line.
\item Repeat this finite packet at one adjacent cutoff and keep the crossing,
      source-effect, polar-frame, and line defects separate.
\end{enumerate}
The mathematical outcomes are \(\passstatus\), \(\obsstatus\), or
\(\stopstatus\).  A failed domain or same-history check rejects the input
before that trichotomy.
\end{theorem}

\begin{proof}
Items~\textup{(D14.1)--(D14.2)} are the hypotheses of
\Cref{thm:V14-balanced-sign,thm:V14-polar-criterion}.  Item~\textup{(D14.3)}
is \Cref{thm:V14-polar-collar}.  Item~\textup{(D14.4)} is the inherited
Version~V13 validated sign card followed by the Version~V9 exact cone
compiler.  Item~\textup{(D14.5)} is
\Cref{thm:V14-supported-polar,ass:V14-singular-meaning}.  The line and held-out
independence are already part of the inherited card and are explicitly
replayed in \Cref{prop:V14-heldout}.  Item~\textup{(D14.7)} is the existing
crossing transport rectangle, with \Cref{prop:V14-two-card} showing why the
row cannot be omitted.
\end{proof}

\begin{theorem}[Version V14 balanced gauge-active crossing theorem]
\label{thm:V14-integrated}
Versions~V1--V13 are retained.  Version~V14 additionally proves:
\begin{enumerate}[label=\textup{(E14.\arabic*)},leftmargin=5.5em]
\item the sign of every balanced finite Wilson collar is the explicit block
      matrix \eqref{eq:V14-sign-blocks};
\item its crossing singular values are the bounded transform
      \(s/\sqrt{m^2+s^2}\) of the actual hopping singular values;
\item after the positive Hodge--chiral factor is included, crossing positivity
      is equivalent to one exact signed polar-reflection incidence;
\item the positive floor is the least singular value of
      \(L F_m(S)\), divided by \(a\), and not an identity-normalized link
      eigenvalue;
\item singular hopping directions admit the zero-safe supported criterion
      \eqref{eq:V14-supported-polar-condition} and remain in the divisor packet;
\item the moving overlap-source chronological effect has spectrum
      \eqref{eq:V14-A-eigenvalues}, leakage \eqref{eq:V14-exact-leakage}, and
      mixed rank exactly \(\rank B\);
\item the complete polar card is gauge covariant and descends automatically
      through every actual one-letter relation;
\item on isotropic hopping the complete skew/negative cone distance is the
      phase formula \eqref{eq:V14-phase-cone};
\item the balanced card has an explicit gauge-active perturbation collar whose
      debit is the same-side balance residual divided by the Wilson path gap;
\item the nontrivial \(SO(3)\subset SU(3)\) replay has positive crossing
      \((4/5)I_3\), exact factor rank three, a passing held-out four-phase Read,
      a relation mutation residual \(32\), and a wrong-reflection cone distance
      \(512/625\);
\item a second exact collar has nonzero crossing and source-effect transport
      defects, so cofinal passage remains an independent row.
\end{enumerate}
No item evaluates the selected active RPESM gauge history, its determinant or
Pfaffian line, or a cutoff-uniform crossing floor.
\end{theorem}

\begin{proof}
Items~\textup{(E14.1)--(E14.2)} are
\Cref{thm:V14-balanced-sign}.  Items~\textup{(E14.3)--(E14.4)} are
\Cref{thm:V14-polar-criterion}.  Item~\textup{(E14.5)} is
\Cref{thm:V14-supported-polar}.  Item~\textup{(E14.6)} is
\Cref{thm:V14-A-spectrum}.  Item~\textup{(E14.7)} is
\Cref{thm:V14-gauge-covariance,thm:V14-relation-descent}.  Item~\textup{(E14.8)}
is \Cref{thm:V14-unitary-phase}.  Item~\textup{(E14.9)} is
\Cref{thm:V14-polar-collar}.  Item~\textup{(E14.10)} is
\Cref{prop:V14-SU3-positive,prop:V14-control-relation,prop:V14-heldout,cth:V14-wrong-reflection}.
The final transport statement is \Cref{prop:V14-two-card}.
\end{proof}

\begingroup\small
\begin{longtable}{>{\raggedright\arraybackslash}p{0.27\textwidth}
                  >{\raggedright\arraybackslash}p{0.20\textwidth}
                  >{\raggedright\arraybackslash}p{0.45\textwidth}}
\caption{Version V14 proof-status ledger.}\label{tab:V14-ledger}\\
\toprule
\textbf{Row}&\textbf{Status}&\textbf{Exact scope}\\
\midrule
\endfirsthead
\toprule
\textbf{Row}&\textbf{Status}&\textbf{Exact scope}\\
\midrule
\endhead
V1--V13 source metric, Duhamel, clock, record, Grassmann, and overlap compilers
&Imported exactly
&Preserved verbatim; no prior Schur, record, free-overlap, or sign enclosure is
reproved.\\[0.3em]
Balanced Wilson-collar sign
&\passstatus
&Exact for arbitrary finite cross-plane hopping \(B\) and scalar massive
same-side blocks.\\[0.3em]
Hodge--polar reflection
&\passstatus
&For full-rank \(B\), positivity is equivalent to the unique signed polar
transport of \(L F_m(S)\).\\[0.3em]
Singular hopping support
&\passstatus
&The supported positive block has rank \(\rank B\); kernel directions remain
zero-safe line/divisor data.\\[0.3em]
Chronological source effect
&\passstatus
&Its complete spectrum and leakage are explicit; mixed rank is exactly
\(\rank B\).\\[0.3em]
Gauge covariance and word relations
&\passstatus
&The polar card transforms by unitary conjugation and every raw crossing pulled
through \(W_1\) annihilates \(\ker W_1\) on both legs.\\[0.3em]
Arbitrary gauge-active collar
&Conditional sufficient theorem
&Passes when the physical polar incidence is exact, the Wilson path stays
gapped, the same-side balance debit is below the positive floor, and
Hermiticity is independently certified.\\[0.3em]
Nontrivial \(SU(3)\) control
&\passstatus
&Exact rational card with crossing \((4/5)I_3\), source effect \((4/5)I_3\),
leakage \(12/25\), and factor rank three.\\[0.3em]
Wrong physical reflection
&\obsstatus
&Same gauge link and Wilson singular values, but exact skew/cone residual
\(512/625\).\\[0.3em]
One-letter relation mutation
&\obsstatus
&Adding \(vv^*\) in the zero-word direction gives exact two-sided residual
\(32\).\\[0.3em]
Held-out four-phase Read
&\passstatus on the control
&Four direct quadratic probes reconstruct \(-4/5\) before the frozen reflection
sign and exactly predict the retained crossing entry.\\[0.3em]
Adjacent control transport
&Populated, nonzero
&Crossing and source-effect defects are \(8/65\) and \(7/65\), respectively;
finite positivity is not exact projectivity.\\[0.3em]
Selected RPESM gauge history
&\stopstatus
&Its Wilson blocks, polar-reflection residual, balance radius, direct cone
interval, line orientation, and adjacent copy remain unevaluated.\\[0.3em]
Duhamel, current/stress, charge, continuum, Einstein
&Not reopened
&All remain downstream of a passing crossing and line packet.\\
\bottomrule
\end{longtable}
\endgroup

\begin{assessment}[Physical status after Version V14]
\label{ass:V14-status}
The exact gauge-active control shows that nontrivial parallel transport is not
itself an obstruction: when the reflection carries the signed polar transport
of the cross-plane hopping, the full crossing is positive with an explicit
floor.  It also identifies the first way the control can fail.  A reflection
which has the correct source and target spaces but the wrong polar phase
creates a skew or negative crossing even though the Wilson gap, determinant
modulus, and hopping singular values are unchanged.

The supplied RPESM construction guarantees a finite common field shell and
populates the required finite grammar, but it does not print one selected
history's numerical Wilson boundary blocks or polar incidence.  Accordingly,
the present physical verdict is
\begin{equation}
 \boxed{
 \begin{array}{cl}
 \passstatus:&\text{balanced theorem and exact gauge-active control;}\\[0.25em]
 \obsstatus:&\text{wrong reflection and zero-word mutation;}\\[0.25em]
 \stopstatus:&\text{selected RPESM polar/balance/crossing card.}
 \end{array}}
 \label{eq:V14-status}
\end{equation}
No line, exterior, Duhamel, represented current/stress, charge, interacting
continuum, or Einstein conclusion is inferred.
\end{assessment}

\begin{openproblem}[Version V15: populate one actual Wilson boundary collar]
\label{op:V14-next}
Preserve Versions~V1--V14.  Do not add another general crossing, record,
Grassmann, exterior, Duhamel, concentration, or continuum theorem.  At the
first predeclared active RPESM gauge history:
\begin{enumerate}[label=\textup{(V15.\arabic*)},leftmargin=5.8em]
\item extract the actual chronological blocks of \(H_W(U)\), the physical
      mass parameter, cross-plane hopping \(B\), Hodge--chiral factor \(L\),
      reflection \(R_\Theta\), and Berezin sign on one verified common support;
\item compute \(B=VS\), the polar unitary of \(LF_m(S)\), the residual
      \eqref{eq:V14-polar-residual}, the balance debit
      \eqref{eq:V14-balance-residual}, and the protected path gap;
\item if \eqref{eq:V14-collar-pass} clears, export the explicit positive floor;
      otherwise evaluate the direct Wilson sign with the Version~V13 outward
      enclosure and run the exact Version~V9 cone card;
\item use the actual \(W_1\) and source Gram, retain every zero hopping support,
      and acquire one held-out four-phase value on the same history;
\item read the determinant/Pfaffian line, divisor, and orientation independently;
\item repeat this exact finite tuple at one adjacent cutoff and stop at the first
      nontransported crossing, source-effect, polar-frame, or line row.
\end{enumerate}
Only a passing populated tuple returns to the Version~V8 Hamiltonian tangent
and Versions~V5--V6 Duhamel ancestry.  If the active Wilson kernel is far from
the balanced collar, that is not circularity: the direct sign matrix is then
the required finite physical value.
\end{openproblem}

\section*{V14 exact replay and append-only record}
\addcontentsline{toc}{section}{V14 exact replay and append-only record}
The accompanying verifier uses only Python standard-library exact rational
arithmetic.  It reconstructs the nonidentity
\(SO(3)\subset SU(3)\) link, the six-dimensional balanced Wilson kernel, its
sign and overlap projector, the chronological effect, the polar and mismatched
crossings, the redundant one-letter Gram, the canonical relation support, the
held-out four-phase polarization, and the two-card defects.  It checks all
identities by exact equality.  Ordinary and optimized Python executions
produce byte-identical JSON output.

The complete Version~V13 source preceding its sole final active
\verb|\end{document}| is preserved byte for byte.  Its preterminal SHA--256
digest is
\begin{center}\scriptsize\ttfamily
e22ed4c7cf8b1e3e0416542209eeb4a617c57d6644f70f74223b0a78b5df0ebb.
\end{center}
The exact verifier and its JSON output have SHA--256 digests
\begin{center}\scriptsize\ttfamily
c1476650c9042841c005b394c9037a3c167c21606083329041d118920213cb79,\\[-0.1em]
9fcbc942c38ca8bd442f26fc1e4fcf1edb3529085d804b1f0f53dbb457254320.
\end{center}
A later continuation must remove only the sole final active
\verb|\end{document}|, append new material immediately before it, restore one
terminal command, and increment the filename to end in \texttt{\_V15.tex}.

% =============================================================================
% END APPENDED VERSION V14 MATERIAL
% =============================================================================


% =============================================================================
% START APPENDED VERSION V15 MATERIAL -- complete V1--V14 preterminal source preserved
% =============================================================================

\clearpage
\hypersetup{
 pdftitle={Renewal Einstein--Standard--Model Physical Source Metric, Duhamel Ancestry and Quantum Continuum Programme V15},
 pdfsubject={Channelwise affine Wilson collars, a canonical boundary-channel anchor, and a noncommuting gauge-active certificate}
}
\section*{Version V15: channelwise affine Wilson collars and the canonical boundary-channel anchor}
\addcontentsline{toc}{section}{V15: channelwise affine Wilson collars and the canonical boundary-channel anchor}

\begin{abstract}
Version~V14 gave an exact balanced Wilson collar and a polar-reflection
criterion.  Its perturbative debit compared the same-side blocks with the
special scalar pair $-mI,+mI$.  That comparison is sufficient but not
intrinsic: even in one channel, the sign of
$\begin{psmallmatrix}-\alpha&-s\\-s&\beta\end{psmallmatrix}$ depends on
$(\alpha+\beta)/2$ and $s$, but not on $(\beta-\alpha)/2$, while the two
eigenvalues retain opposite signs.

This continuation identifies the exact larger class.  Freeze, before reading
the Wilson blocks, an abelian boundary-channel algebra
$\mathscr A_\partial$ containing the hopping singular operator.  If $B=VS$
and the transported same-side blocks are $C-Z$ and $C+Z$, with
$C,Z,S\in\mathscr A_\partial$, then the complete collar has sign
\[
 \begin{pmatrix}
 -VC(C^2+S^2)^{-1/2}V^*&-VS(C^2+S^2)^{-1/2}\\
 -S(C^2+S^2)^{-1/2}V^*&C(C^2+S^2)^{-1/2}
 \end{pmatrix}.
\]
It is independent of the channelwise time-asymmetry $Z$.  The reflected
crossing and the moving overlap-source effect therefore depend only on the
channel mean $C$, hopping magnitude $S$, Hodge factor, and physical
reflection.

For arbitrary positive same-side blocks, the trace-preserving conditional
expectation onto $\mathscr A_\partial$ produces the unique
Hilbert--Schmidt-nearest exact collar.  Its debit is exactly the sum of the two
off-channel block energies.  This removes the arbitrary scalar reference in
V14.  A two-algebra counterexample shows that the channel algebra must itself
be a physical boundary record: choosing it after seeing the matrix can change
the debit from zero to four.

The canonical anchor yields a direct sign perturbation certificate.  On an
exact noncommuting two-channel control, its debit is $1/16$ and the true
crossing has certified floor $1/3$, without diagonalizing the perturbed sign.
A second exact three-colour control has nonidentity
$SO(3)\subset SU(3)$ transport, channelwise asymmetry, crossing
$\operatorname{diag}(12/13,3/5,12/37)$, and chronological mixed rank three.
The remaining RPESM input is now the physical channel algebra and its literal
same-history Wilson blocks, not another fitted balanced mass.
\end{abstract}

\begin{assessment}[Why V15 goes upstream of another fitted balance test]
\label{ass:V15-direction}
The supplied field-shell manuscript does not print one selected RPESM Wilson
boundary matrix.  Repeating the V14 collar with a newly chosen scalar $m$
would therefore be circular.  The earlier datum that makes the comparison
canonical is the physical boundary-channel decomposition.  Once that algebra
is frozen, the nearest exactly solvable collar is unique.  No source space,
record hierarchy, exterior algebra, clock, or continuum theorem is enlarged.
\end{assessment}

% =============================================================================
\section{The physical boundary-channel algebra}
\label{sec:V15-channel-algebra}
% =============================================================================

Let $\mathcal E_-$ and $\mathcal E_+$ have the same finite dimension $d$.
Let
\begin{equation}
 B=VS:\mathcal E_+\longrightarrow\mathcal E_-
 \label{eq:V15-polar-B}
\end{equation}
be invertible, with $V$ unitary and $S=(B^*B)^{1/2}>0$.  Singular hopping
supports remain governed by the supported V14 theorem and the divisor packet.

\begin{definition}[Predeclared boundary-channel algebra]
\label{def:V15-channel-algebra}
A boundary-channel algebra for $B$ is an abelian unital $*$-subalgebra
\begin{equation}
 \boxed{\mathscr A_\partial\subseteq\operatorname{End}(\mathcal E_+),
 \qquad S\in\mathscr A_\partial,}
 \label{eq:V15-channel-algebra}
\end{equation}
fixed by the physical site, spin, flavour, or boundary-orbit record before any
crossing value is read.  If $P_1,\ldots,P_q$ are its minimal projections, then
$\mathscr A_\partial=\bigoplus_r\mathbb CP_r$ and its trace-preserving
conditional expectation is
\begin{equation}
 \boxed{
 \mathbb E_\partial(X)
 =\sum_{r=1}^q
 \frac{\operatorname{Tr}(P_rXP_r)}{\operatorname{rank}P_r}P_r.}
 \label{eq:V15-channel-expectation}
\end{equation}
\end{definition}

\begin{lemma}[Positive orthogonal channel projection]
\label{lem:V15-channel-expectation}
The map $\mathbb E_\partial$ is unital, trace preserving, completely positive,
$*$-preserving, and is the Hilbert--Schmidt orthogonal projection onto
$\mathscr A_\partial$.  In particular,
\begin{equation}
 X\succeq x_*I\quad\Longrightarrow\quad
 \mathbb E_\partial(X)\succeq x_*I.
 \label{eq:V15-channel-floor}
\end{equation}
\end{lemma}

\begin{proof}
Each corner map in \eqref{eq:V15-channel-expectation} is a normalized trace
followed by multiplication by $P_r$, hence completely positive.  The corner
projections sum to $I$, giving unitality and trace preservation.  For
$A=\sum_ra_rP_r$,
\[
 \operatorname{Tr}\bigl(A^*(X-\mathbb E_\partial X)\bigr)=0,
\]
so the map is the orthogonal projection.  Applying it to $X-x_*I\succeq0$
proves \eqref{eq:V15-channel-floor}.
\end{proof}

\begin{firewall}[The channel algebra is occurrence data]
\label{fire:V15-channel-not-fit}
When $S$ has degenerate singular values, its spectral projections do not
select the physical channels inside the degenerate carrier.  An abelian
algebra fitted after seeing the same-side blocks can make a residual vanish by
definition.  Such a fit is not a Wilson-sign certificate.  The algebra in
\eqref{eq:V15-channel-algebra} must be exported by the boundary record or an
independently frozen representation convention.
\end{firewall}

% =============================================================================
\section{Exact sign of a channelwise affine Wilson collar}
\label{sec:V15-affine-collar}
% =============================================================================

Take self-adjoint $C,Z\in\mathscr A_\partial$ with
\begin{equation}
 C-Z\succ0,
 \qquad C+Z\succ0,
 \label{eq:V15-CZ-positive}
\end{equation}
and put
\begin{equation}
 R=(C^2+S^2)^{1/2},
 \qquad F_{C,S}=SR^{-1}.
 \label{eq:V15-RF}
\end{equation}
Define
\begin{equation}
 \boxed{
 H_{C,Z,B}
 =\begin{pmatrix}
 -V(C-Z)V^*&-VS\\
 -SV^*&C+Z
 \end{pmatrix}.}
 \label{eq:V15-affine-H}
\end{equation}

\begin{theorem}[Exact channelwise affine sign]
\label{thm:V15-affine-sign}
The collar \eqref{eq:V15-affine-H} is invertible and
\begin{equation}
 \boxed{
 \operatorname{dist}(0,\operatorname{spec}H_{C,Z,B})
 =\lambda_{\min}(R-|Z|)>0.}
 \label{eq:V15-affine-gap}
\end{equation}
Its sign is independent of $Z$:
\begin{equation}
 \boxed{
 \operatorname{sgn}H_{C,Z,B}
 =\begin{pmatrix}
 -VCR^{-1}V^*&-VSR^{-1}\\
 -SR^{-1}V^*&CR^{-1}
 \end{pmatrix}.}
 \label{eq:V15-affine-sign}
\end{equation}
Hence
\begin{equation}
 \boxed{
 T_+\operatorname{sgn}H_{C,Z,B}T_-
 =-F_{C,S}V^*.}
 \label{eq:V15-affine-cross}
\end{equation}
On a joint channel with scalar values $(c,z,s)$, the two collar eigenvalues
are $z\pm\sqrt{c^2+s^2}$ and the crossing singular value is
\begin{equation}
 \boxed{\frac{s}{\sqrt{c^2+s^2}}.}
 \label{eq:V15-channel-cross-value}
\end{equation}
\end{theorem}

\begin{proof}
Conjugation by $\operatorname{diag}(V^*,I)$ gives
\[
 H'=\begin{pmatrix}Z-C&-S\\-S&Z+C\end{pmatrix}
 =\begin{pmatrix}Z&0\\0&Z\end{pmatrix}
  +K,
 \qquad
 K=\begin{pmatrix}-C&-S\\-S&C\end{pmatrix}.
\]
Because $C,Z,S$ lie in one abelian algebra,
$K^2=\operatorname{diag}(R^2,R^2)$.  On each character,
$c\pm z>0$ implies $c>|z|$, while $s>0$ gives
$r=\sqrt{c^2+s^2}>c$.  Thus $z+r>0$ and $z-r<0$, so the sign on that channel
is $K/r$ and the gap is $r-|z|$.  Reassembling the characters and conjugating
back proves all displayed formulas.
\end{proof}

\begin{corollary}[Scalar imbalance is sign invisible]
\label{cor:V15-scalar-imbalance}
For $\alpha,\beta>0$ and arbitrary invertible $B$, the cross block of
\begin{equation}
 \begin{pmatrix}-\alpha I&-B\\-B^*&\beta I\end{pmatrix}
 \label{eq:V15-scalar-unbalanced-H}
\end{equation}
is the V14 balanced formula with mass $(\alpha+\beta)/2$.  It is independent
of $\beta-\alpha$ while the gap remains open.
\end{corollary}

\begin{proof}
Apply \Cref{thm:V15-affine-sign} with
$C=(\alpha+\beta)I/2$ and $Z=(\beta-\alpha)I/2$.
\end{proof}

\begin{assessment}[Qualification of the V14 balance debit]
\label{ass:V15-V14-qualification}
The V14 balanced theorem remains correct.  Its scalar balance debit is a
sufficient upper bound, not the intrinsic minimum debit.  Version~V15 removes
every same-side component already represented by the physical channel
algebra and charges only the remaining off-channel mixing.
\end{assessment}

% =============================================================================
\section{Reflected hopping and the chronological source effect}
\label{sec:V15-cross-source}
% =============================================================================

Let $L=L^*>0$ be the verified Hodge--chiral factor on the represented
positive-time boundary carrier.  Let
$R_\Theta:\mathcal E_+\to\mathcal E_-$ be the unitary linear part of the
physical reflection, and let $\varepsilon_\Theta\in\{+1,-1\}$ be the frozen
Berezin sign.  Write
\begin{equation}
 A_{C,S}=LF_{C,S}=U_{C,S}|A_{C,S}|.
 \label{eq:V15-A-polar}
\end{equation}

\begin{corollary}[Channelwise Hodge--polar reflection criterion]
\label{cor:V15-polar-criterion}
The reflected overlap crossing is
\begin{equation}
 K_\times
 =\frac{\varepsilon_\Theta}{a}LF_{C,S}V^*R_\Theta.
 \label{eq:V15-K-crossing}
\end{equation}
It is Hermitian positive definite exactly when
\begin{equation}
 \boxed{R_\Theta=\varepsilon_\Theta VU_{C,S}^*.}
 \label{eq:V15-reflection-match}
\end{equation}
On this branch,
\begin{equation}
 \boxed{
 K_\times
 =\frac1aU_{C,S}|A_{C,S}|U_{C,S}^*
 \succeq\kappa_{C,S,L}I,}
 \label{eq:V15-positive-K}
\end{equation}
where
\begin{equation}
 \boxed{
 \kappa_{C,S,L}
 =\frac1a s_{\min}\!\left(LS(C^2+S^2)^{-1/2}\right).}
 \label{eq:V15-K-floor}
\end{equation}
Neither the criterion nor the floor depends on $Z$.
\end{corollary}

\begin{proof}
Equation~\eqref{eq:V15-affine-cross} gives
\eqref{eq:V15-K-crossing}.  The right unitary factor of $A_{C,S}$ is
$\varepsilon_\Theta V^*R_\Theta$.  Uniqueness of the polar decomposition of
an invertible matrix makes the product positive exactly when this factor is
$U_{C,S}^*$.  Its least eigenvalue is then the least singular value of
$A_{C,S}/a$.
\end{proof}

Let
\begin{equation}
 P_\chi=\frac12(I+\operatorname{sgn}H_{C,Z,B}),
 \qquad \mathcal E_\chi=\operatorname{Ran}P_\chi,
 \label{eq:V15-Pchi}
\end{equation}
and let $J_\chi$ be the inclusion of $\mathcal E_\chi$ into
$\mathcal E_-\oplus\mathcal E_+$.  The actual positive-time source effect is
\begin{equation}
 A_+^\chi=J_\chi^*T_+J_\chi.
 \label{eq:V15-Aplus}
\end{equation}

\begin{theorem}[Exact channelwise chronological effect]
\label{thm:V15-source-effect}
The source effect is unitarily equivalent to
\begin{equation}
 \boxed{Q_+^\chi=\frac12(I+CR^{-1}).}
 \label{eq:V15-Qplus}
\end{equation}
Consequently
\begin{equation}
 \boxed{
 \operatorname{spec}A_+^\chi
 =\left\{\frac12\left(1+\frac{c}{\sqrt{c^2+s^2}}\right)
 : (c,s)\text{ a channel}\right\},}
 \label{eq:V15-Aplus-spectrum}
\end{equation}
\begin{equation}
 \boxed{
 \operatorname{Tr}\bigl(A_+^\chi-(A_+^\chi)^2\bigr)
 =\frac14\operatorname{Tr}\!\left[S^2(C^2+S^2)^{-1}\right],}
 \label{eq:V15-source-leakage}
\end{equation}
and
\begin{equation}
 \boxed{\operatorname{rank}\bigl(A_+^\chi(I-A_+^\chi)\bigr)=d.}
 \label{eq:V15-source-rank}
\end{equation}
These quantities are independent of $Z$.
\end{theorem}

\begin{proof}
The lower-right block of $P_\chi$ is $Q_+^\chi$.  For
$A=T_+J_\chi:\mathcal E_\chi\to\mathcal E_+$ one has
$A^*A=A_+^\chi$ and $AA^*=Q_+^\chi$.  The collar has $d$ positive
eigenvalues, while $Q_+^\chi\succ I/2$, so $A$ is invertible and the two
positive matrices are unitarily equivalent.  Channelwise,
$q(1-q)=s^2/[4(c^2+s^2)]$.  Since $S>0$, every leakage eigenvalue is positive,
which proves the trace and rank formulas.
\end{proof}

% =============================================================================
\section{The canonical nearest channelwise collar}
\label{sec:V15-canonical-anchor}
% =============================================================================

Let the actual finite Wilson collar be
\begin{equation}
 H=\begin{pmatrix}-M_-&-B\\-B^*&M_+\end{pmatrix},
 \qquad M_-=M_-^*>0,
 \quad M_+=M_+^*>0,
 \label{eq:V15-actual-H}
\end{equation}
with the same $B=VS$.  Transport the negative-side block into the positive
boundary frame and form its mean and asymmetry:
\begin{equation}
 X=V^*M_-V,
 \qquad
 \overline C=\frac12(X+M_+),
 \qquad
 \overline Z=\frac12(M_+-X).
 \label{eq:V15-XCZbar}
\end{equation}
Define
\begin{equation}
 \boxed{
 C_\partial=\mathbb E_\partial(\overline C),
 \qquad
 Z_\partial=\mathbb E_\partial(\overline Z).}
 \label{eq:V15-CZ-anchor}
\end{equation}
Then
\begin{equation}
 C_\partial-Z_\partial=\mathbb E_\partial(X)>0,
 \qquad
 C_\partial+Z_\partial=\mathbb E_\partial(M_+)>0.
 \label{eq:V15-anchor-positive}
\end{equation}
Let $H_\partial$ be the affine collar
\eqref{eq:V15-affine-H} formed from
$(C_\partial,Z_\partial,B)$.

\begin{theorem}[Canonical boundary-channel anchor]
\label{thm:V15-canonical-anchor}
Among all exact affine collars with the same hopping block $B$ and
$C,Z\in\mathscr A_\partial$, the pair
$(C_\partial,Z_\partial)$ uniquely minimizes
\begin{equation}
 \mathcal J(C,Z)
 =\|M_--V(C-Z)V^*\|_{\rm HS}^2
  +\|M_+-(C+Z)\|_{\rm HS}^2.
 \label{eq:V15-JCZ}
\end{equation}
The attained intrinsic channel debit is
\begin{equation}
 \boxed{
 \begin{aligned}
 \Delta_\partial^{\rm ch}
 &:=\|H-H_\partial\|_{\rm HS}^2\\
 &=\|(I-\mathbb E_\partial)X\|_{\rm HS}^2
   +\|(I-\mathbb E_\partial)M_+\|_{\rm HS}^2\\
 &=2\|(I-\mathbb E_\partial)\overline C\|_{\rm HS}^2
   +2\|(I-\mathbb E_\partial)\overline Z\|_{\rm HS}^2.
 \end{aligned}}
 \label{eq:V15-channel-debit}
\end{equation}
Thus no component already measurable in the physical boundary-channel algebra
is charged.
\end{theorem}

\begin{proof}
Unitary invariance changes the first term in \eqref{eq:V15-JCZ} into
$\|X-(C-Z)\|_{\rm HS}^2$.  The change of variables
$Y_-=C-Z$, $Y_+=C+Z$ is a bijection on the self-adjoint pairs in
$\mathscr A_\partial^2$.  The minimization therefore separates into the two
orthogonal projection problems for $X$ and $M_+$.  Their unique minimizers are
$\mathbb E_\partial X$ and $\mathbb E_\partial M_+$.  Recombination gives
\eqref{eq:V15-CZ-anchor}.  The last equality in
\eqref{eq:V15-channel-debit} is the parallelogram identity applied to
$\overline C\mp\overline Z$.
\end{proof}

\begin{proposition}[Gauge covariance of the anchor]
\label{prop:V15-anchor-gauge}
Under boundary frame changes
\begin{equation}
 \begin{aligned}
 B&\mapsto U_-BU_+^*,&M_-&\mapsto U_-M_-U_-^*,\\
 M_+&\mapsto U_+M_+U_+^*,&
 \mathscr A_\partial&\mapsto U_+\mathscr A_\partial U_+^*,
 \end{aligned}
 \label{eq:V15-gauge-transform}
\end{equation}
the anchor transforms by the same unitary congruence and
$\Delta_\partial^{\rm ch}$ is invariant.
\end{proposition}

\begin{proof}
The polar factors transform as
$V\mapsto U_-VU_+^*$ and $S\mapsto U_+SU_+^*$.  Hence
$X,\overline C,\overline Z$ transform by $U_+$-conjugacy.  The conditional
expectation onto the transported algebra is
$Y\mapsto U_+\mathbb E_\partial(U_+^*YU_+)U_+^*$, proving covariance.
Hilbert--Schmidt unitary invariance proves the last assertion.
\end{proof}

\begin{countertheorem}[Without an occurring channel algebra the collar can be fitted]
\label{cth:V15-no-channel-fit}
Let $B=I_2$ and
\begin{equation}
 M_-=M_+=\begin{pmatrix}1&0\\0&3\end{pmatrix}.
 \label{eq:V15-fit-data}
\end{equation}
Both
\begin{equation}
 \mathscr A_z=\operatorname{span}\{I,\sigma_z\},
 \qquad
 \mathscr A_x=\operatorname{span}\{I,\sigma_x\}
 \label{eq:V15-two-algebras}
\end{equation}
contain $S=I_2$.  For $\mathscr A_z$ the canonical debit is zero and the
crossing values are $1/\sqrt2$ and $1/\sqrt{10}$.  For $\mathscr A_x$ the
anchor is $2I_2$ on both sides, the debit is
\begin{equation}
 \boxed{\Delta_x^{\rm ch}=4,}
 \label{eq:V15-fit-debit}
\end{equation}
and the anchor crossing is $I_2/\sqrt5$.  The direct Wilson sign is unchanged;
only the post hoc certificate changes.
\end{countertheorem}

\begin{proof}
The first expectation fixes the diagonal matrix.  The second is
$\mathbb E_x(Y)=(Y+\sigma_xY\sigma_x)/2$, so it sends the matrix to $2I$.
Each of the two residual blocks is $\operatorname{diag}(-1,1)$ and has squared
Hilbert--Schmidt norm two.  Formula
\eqref{eq:V15-channel-cross-value} gives the crossing values.
\end{proof}

% =============================================================================
\section{A canonical collar certificate for the actual Wilson sign}
\label{sec:V15-collar-certificate}
% =============================================================================

Put
\begin{equation}
 R_\partial=(C_\partial^2+S^2)^{1/2},
 \qquad
 g_\partial=\lambda_{\min}(R_\partial-|Z_\partial|)>0,
 \label{eq:V15-anchor-gap}
\end{equation}
and
\begin{equation}
 e_{\rm op}=\|H-H_\partial\|_{\rm op},
 \qquad
 e_{\rm HS}=\|H-H_\partial\|_{\rm HS}
 =\sqrt{\Delta_\partial^{\rm ch}}.
 \label{eq:V15-anchor-errors}
\end{equation}

\begin{theorem}[Canonical channel-collar perturbation theorem]
\label{thm:V15-collar-certificate}
If $e_{\rm op}<g_\partial$, then the straight path from $H_\partial$ to $H$
has gap at least $g_\partial-e_{\rm op}$ and
\begin{equation}
 \boxed{
 \|\operatorname{sgn}H-\operatorname{sgn}H_\partial\|_{\rm HS}
 \le\frac{e_{\rm HS}}{g_\partial-e_{\rm op}}.}
 \label{eq:V15-sign-anchor-error}
\end{equation}
Assume that the physical reflection obeys the anchor polar incidence
\eqref{eq:V15-reflection-match} and define
\begin{equation}
 \eta_\partial
 =\frac{\|L\|_{\rm op}}a
   \frac{e_{\rm HS}}{g_\partial-e_{\rm op}}.
 \label{eq:V15-crossing-radius}
\end{equation}
Then
\begin{equation}
 \boxed{
 \|K_\times(H)-K_\times(H_\partial)\|_{\rm HS}
 \le\eta_\partial.}
 \label{eq:V15-K-anchor-error}
\end{equation}
If the true crossing is independently certified Hermitian and
$\eta_\partial<\kappa_{C_\partial,S,L}$, then
\begin{equation}
 \boxed{
 K_\times(H)\succeq
 (\kappa_{C_\partial,S,L}-\eta_\partial)I>0.}
 \label{eq:V15-collar-positive-floor}
\end{equation}
Without the Hermiticity row, the same radius proves only
\begin{equation}
 \operatorname{dist}_{\rm HS}
 (K_\times(H),\operatorname{Herm}_+)\le\eta_\partial.
 \label{eq:V15-cone-upper}
\end{equation}
\end{theorem}

\begin{proof}
Weyl's inequality gives the stated path gap.  The inherited V13
sign-Lipschitz theorem gives \eqref{eq:V15-sign-anchor-error}.  Chronological
compression, reflection, and the chiral projections are contractions, so
multiplication by $L/a$ gives \eqref{eq:V15-K-anchor-error}.  The anchor
crossing is positive with floor \eqref{eq:V15-K-floor}; Weyl's lower-eigenvalue
bound proves \eqref{eq:V15-collar-positive-floor}.  The last inequality is
distance to the particular positive matrix $K_\times(H_\partial)$.
\end{proof}

\begin{lemma}[Centred same-side blocks make the crossing Hermitian]
\label{lem:V15-centred-Hermitian}
In the polar frame suppose
\begin{equation}
 H'=\begin{pmatrix}-X&-S\\-S&X\end{pmatrix},
 \qquad X=X^*,\quad S=S^*.
 \label{eq:V15-centred-H}
\end{equation}
For
\begin{equation}
 \mathcal J=\begin{pmatrix}0&I\\-I&0\end{pmatrix},
 \label{eq:V15-J}
\end{equation}
one has $\mathcal JH'+H'\mathcal J=0$.  The lower-left block of
$\operatorname{sgn}H'$ is therefore self-adjoint.  Consequently, for scalar
positive $L$ and $R_\Theta=\varepsilon_\Theta V$, the reflected crossing is
Hermitian even when $X$ and $S$ do not commute.
\end{lemma}

\begin{proof}
The anticommutator vanishes by block multiplication and hence also vanishes
with $\operatorname{sgn}H'$ by functional calculus.  If
$\operatorname{sgn}H'=\begin{psmallmatrix}P&Q\\R&T\end{psmallmatrix}$,
anticommutation gives $R=Q$ and $T=-P$.  Self-adjointness gives $Q=Q^*$.
Conjugating back inserts $V^*$ on the negative leg, and the stated reflection
removes it.
\end{proof}

% =============================================================================
\section{Exact gauge-active controls}
\label{sec:V15-controls}
% =============================================================================

\subsection{A channelwise asymmetric three-colour seam}

Let
\begin{equation}
 V=\begin{pmatrix}
 3/5&-4/5&0\\4/5&3/5&0\\0&0&1
 \end{pmatrix}\in SO(3)\subset SU(3),
 \qquad B=12V,
 \label{eq:V15-SU3-V}
\end{equation}
and use the diagonal channel algebra.  Put
\begin{equation}
 C=\operatorname{diag}(5,16,35),
 \qquad
 Z=\operatorname{diag}(1,-2,3).
 \label{eq:V15-SU3-CZ}
\end{equation}
Then $C-Z=\operatorname{diag}(4,18,32)$ and
$C+Z=\operatorname{diag}(6,14,38)$, so this collar is neither scalar balanced
nor channelwise symmetric.

\begin{proposition}[Exact nonidentity $SU(3)$ affine-collar card]
\label{prop:V15-SU3-card}
For $a=L=1$, $\varepsilon_\Theta=1$, and $R_\Theta=V$,
\begin{equation}
 R=\operatorname{diag}(13,20,37),
 \qquad g_\partial=12,
 \label{eq:V15-SU3-gap}
\end{equation}
\begin{equation}
 \boxed{
 K_\times=\operatorname{diag}\left(\frac{12}{13},\frac35,
 \frac{12}{37}\right)\succeq\frac{12}{37}I_3,}
 \label{eq:V15-SU3-K}
\end{equation}
\begin{equation}
 \boxed{
 A_+^\chi\simeq
 \operatorname{diag}\left(\frac9{13},\frac9{10},\frac{36}{37}\right),}
 \label{eq:V15-SU3-Aplus}
\end{equation}
and
\begin{equation}
 \boxed{
 \operatorname{Tr}(A_+^\chi-(A_+^\chi)^2)
 =\frac{7619049}{23136100},
 \qquad
 \operatorname{rank}(A_+^\chi(I-A_+^\chi))=3.}
 \label{eq:V15-SU3-leakage}
\end{equation}
The asymmetry $Z$ changes the full collar eigenvalues but none of these
crossing or source-effect values.
\end{proposition}

\begin{proof}
Use $5^2+12^2=13^2$, $16^2+12^2=20^2$, and
$35^2+12^2=37^2$ in
\Cref{thm:V15-affine-sign,cor:V15-polar-criterion,thm:V15-source-effect}.
The gap is the minimum of $13-1$, $20-2$, and $37-3$.  The leakage is
$36/169+9/100+36/1369=7619049/23136100$.
\end{proof}

\subsection{A noncommuting side block with a certified floor}

Let
\begin{equation}
 V_2=\begin{pmatrix}3/5&-4/5\\4/5&3/5\end{pmatrix},
 \quad S_2=\operatorname{diag}(3,5),
 \quad B_2=V_2S_2,
 \label{eq:V15-V2S2}
\end{equation}
use the diagonal channel algebra, and put
\begin{equation}
 C_2=\operatorname{diag}(4,12),
 \quad J_2=\begin{pmatrix}0&1\\1&0\end{pmatrix},
 \quad X_2=C_2+\frac18J_2.
 \label{eq:V15-X2}
\end{equation}
The polar-frame actual collar is
\begin{equation}
 H_2'=\begin{pmatrix}-X_2&-S_2\\-S_2&X_2\end{pmatrix}.
 \label{eq:V15-H2prime}
\end{equation}

\begin{proposition}[Exact off-channel debit and strict true crossing floor]
\label{prop:V15-noncommuting-control}
For $a=L=1$, $\varepsilon_\Theta=1$, and $R_\Theta=V_2$, the canonical
channel anchor has
\begin{equation}
 \boxed{
 \Delta_\partial^{\rm ch}=\frac1{16},
 \qquad e_{\rm HS}=\frac14,
 \qquad e_{\rm op}=\frac18.}
 \label{eq:V15-control-debit}
\end{equation}
Its gap is five and its crossing is
\begin{equation}
 K_{\times,\partial}
 =\operatorname{diag}\left(\frac35,\frac5{13}\right)
 \succeq\frac5{13}I_2.
 \label{eq:V15-control-anchor-K}
\end{equation}
The path gap is at least $39/8$, the crossing radius is at most $2/39$, and
the true crossing satisfies
\begin{equation}
 \boxed{K_\times(H_2)\succeq\frac13I_2.}
 \label{eq:V15-control-floor}
\end{equation}
This is a gauge-active noncommuting card outside the exact channel algebra;
its positivity is certified without diagonalizing the perturbed sign.
\end{proposition}

\begin{proof}
Both same-side remainders are $(1/8)J_2$.  Since
$\|J_2\|_{\rm HS}^2=2$ and $\|J_2\|_{\rm op}=1$, their total squared debit is
$1/16$.  The anchor has
$R=\operatorname{diag}(5,13)$, so its gap is five and its crossing floor is
$5/13$.  Weyl gives path gap $5-1/8=39/8$, and hence
\[
 \eta_\partial\le\frac{1/4}{39/8}=\frac2{39}.
\]
The centred anticommutation lemma gives Hermiticity.  Finally,
$5/13-2/39=1/3$, so
\Cref{thm:V15-collar-certificate} proves the claim.
\end{proof}

% =============================================================================
\section{The revised RPESM acquisition and termination rule}
\label{sec:V15-acquisition}
% =============================================================================

\begin{acquisition}[Boundary-channel affine-collar card]
\label{acq:V15-card}
At the first active RPESM link-reflection collar, freeze before reading any
crossing eigenvalue
\begin{equation}
 \boxed{
 \begin{gathered}
 (\mathcal E_-,\mathcal E_+,T_-,T_+),\quad B=VS,\quad
 \mathscr A_\partial,\quad\mathbb E_\partial,\\
 M_-,M_+,\quad L,\quad R_\Theta,\quad\varepsilon_\Theta,\quad W_1,\\
 C_\partial,Z_\partial,\quad\Delta_\partial^{\rm ch},\quad
 g_\partial,\quad\Delta_{\rm pol}^{\times},\\
 \text{one held-out four-phase crossing Read, the line/divisor packet, and}\\
 \text{one adjacent copy with the transported channel algebra.}
 \end{gathered}}
 \label{eq:V15-acquisition-card}
\end{equation}
The channel algebra must be exported by physical boundary labels or an
independently fixed representation.  It is not inferred from a diagonalizer of
$M_-$, $M_+$, or the observed crossing.
\end{acquisition}

\begin{theorem}[Terminating affine-collar/direct-sign order]
\label{thm:V15-decision-order}
For a well-typed card \eqref{eq:V15-acquisition-card}, the finite decision order
is the following.
\begin{enumerate}[label=\textup{(V15.\arabic*)},leftmargin=5.6em]
\item Verify the common boundary history, channel algebra, hopping support,
      Hodge factor, reflection, and Berezin sign.
\item Form the unique anchor of \Cref{thm:V15-canonical-anchor}; reject a
      channel algebra fitted from the crossing or same-side diagonalizers.
\item Evaluate its gap, polar incidence, and crossing floor.  If
      \eqref{eq:V15-collar-positive-floor} clears, export that strict floor.
\item If the sufficient collar does not clear, return no negative verdict and
      fit no second anchor.  Evaluate the direct Wilson sign by the V13
      enclosure and run the inherited V9 relation, skew, and negative-part
      residuals.
\item On a positive crossing branch, retain every zero hopping support and run
      the determinant/Pfaffian line, divisor, and held-out higher-word rows
      independently.
\item Repeat the complete packet at one adjacent cutoff with the transported
      channel algebra.  Stop at the first channel, crossing, support, polar, or
      line defect.
\end{enumerate}
Thus the valid outcomes are a strict collar pass, a direct crossing
obstruction, an unresolved direct interval, or rejection of an ill-typed
input before those mathematical verdicts.  No further abstract crossing or
Duhamel theorem is opened.
\end{theorem}

\begin{proof}
The anchor is unique and gauge covariant by
\Cref{thm:V15-canonical-anchor,prop:V15-anchor-gauge}.  Its sufficient branch
is \Cref{thm:V15-collar-certificate}.  Failure of a sufficient inequality is
not a negative-mode witness, so the direct V13 sign and V9 cone card are then
the exact tests.  The line and divisor rows remain independent by the inherited
fermionic line theorem.  The adjacent transport must include the channel
algebra because \Cref{cth:V15-no-channel-fit} shows that numerical blocks alone
do not determine the collar certificate.
\end{proof}

\begin{firewall}[The anchor is not a replacement action]
\label{fire:V15-anchor-not-repair}
The matrix $H_\partial$ is used only to certify the sign of the actual Wilson
kernel.  It is not inserted into the fermion action, and its positive crossing
is not substituted for a failed direct crossing.  Off-channel terms remain in
$H$ and in every determinant, Pfaffian, exterior, current, stress, and Duhamel
row.
\end{firewall}

% =============================================================================
\section{Version V15 theorem, ledger, and next physical coefficient}
\label{sec:V15-terminal}
% =============================================================================

\begin{theorem}[Version V15 channelwise affine-collar theorem]
\label{thm:V15-terminal}
Versions~V1--V14 are retained.  Version~V15 additionally proves:
\begin{enumerate}[label=\textup{(E15.\arabic*)},leftmargin=5.5em]
\item the physical boundary-channel algebra has a canonical positive
      Hilbert--Schmidt expectation;
\item the exact affine-collar sign \eqref{eq:V15-affine-sign} is independent of
      every channelwise time-asymmetry represented by that algebra;
\item scalar imbalance is sign invisible, so the V14 scalar balance debit is
      sufficient but not intrinsic;
\item the reflected crossing and chronological source effect have the exact
      channel formulas \eqref{eq:V15-K-floor} and
      \eqref{eq:V15-Aplus-spectrum}--\eqref{eq:V15-source-rank};
\item arbitrary positive same-side blocks have a unique nearest exact collar,
      with debit \eqref{eq:V15-channel-debit};
\item the anchor and debit are gauge covariant;
\item absent an occurring channel algebra, post hoc choices produce the exact
      incompatible debits zero and four;
\item the canonical gap and debit give the strict perturbative floor
      \eqref{eq:V15-collar-positive-floor};
\item centred same-side blocks provide an exact Hermiticity certificate;
\item the asymmetric three-colour seam has crossing
      $\operatorname{diag}(12/13,3/5,12/37)$ and mixed rank three;
\item the noncommuting two-channel card has debit $1/16$ and true crossing
      floor $1/3$.
\end{enumerate}
No item evaluates the active RPESM boundary-channel algebra, its numerical
same-history Wilson blocks, the direct Wilson-sign matrix, line orientation,
or a cofinal crossing floor.
\end{theorem}

\begin{proof}
The channel expectation and anchor are
\Cref{lem:V15-channel-expectation,thm:V15-canonical-anchor}.  The exact sign,
reflection, and source formulas are
\Cref{thm:V15-affine-sign,cor:V15-polar-criterion,thm:V15-source-effect}.
Gauge covariance is \Cref{prop:V15-anchor-gauge}; the post hoc obstruction is
\Cref{cth:V15-no-channel-fit}; the collar theorem and Hermiticity row are
\Cref{thm:V15-collar-certificate,lem:V15-centred-Hermitian}; and the two
controls are \Cref{prop:V15-SU3-card,prop:V15-noncommuting-control}.
\end{proof}

\begingroup\small
\begin{longtable}{>{\raggedright\arraybackslash}p{0.27\textwidth}
                  >{\raggedright\arraybackslash}p{0.20\textwidth}
                  >{\raggedright\arraybackslash}p{0.45\textwidth}}
\caption{Version V15 proof-status ledger.}\label{tab:V15-ledger}\\
\toprule
\textbf{Row}&\textbf{Status}&\textbf{Exact scope}\\
\midrule
\endfirsthead
\toprule
\textbf{Row}&\textbf{Status}&\textbf{Exact scope}\\
\midrule
\endhead
V1--V14 source, Duhamel, record, Grassmann, overlap, and balanced-collar results
&Imported exactly
&Preserved verbatim; no prior compiler or exact control is recomputed.\\[0.3em]
Boundary-channel algebra
&\passstatus
&Exact finite occurrence interface; it must be physically frozen.\\[0.3em]
Channelwise affine sign
&\passstatus
&Exact for arbitrary channelwise positive side blocks and invertible gauge
hopping.\\[0.3em]
Channelwise asymmetry
&No sign debit
&It changes the full eigenvalues and gap but not the sign, crossing, or source
effect.\\[0.3em]
Canonical nearest collar
&\passstatus
&Unique Hilbert--Schmidt projection; debit equals the two off-channel energies.\\[0.3em]
Post hoc channel choice
&\obsstatus
&The same tuple has fitted debits zero and four; direct sign remains controlling.\\[0.3em]
Canonical perturbation collar
&Conditional sufficient theorem
&Passes when the anchor gap, polar incidence, Hermiticity, and explicit radius
all clear.\\[0.3em]
Three-colour asymmetric seam
&\passstatus
&Exact nonidentity gauge transport; crossing floor $12/37$ and mixed rank three.\\[0.3em]
Noncommuting two-channel control
&\passstatus
&Off-channel debit $1/16$ and certified true crossing floor $1/3$.\\[0.3em]
Selected RPESM channel/block tuple
&\stopstatus
&Its channel algebra, side blocks, direct sign, line orientation, and adjacent
copy remain unpopulated.\\[0.3em]
Exterior, Duhamel, current/stress, charge, continuum, Einstein
&Not reopened
&All remain downstream of a passing crossing and line packet.\\
\bottomrule
\end{longtable}
\endgroup

\begin{assessment}[Physical status after Version V15]
\label{ass:V15-status}
The first RPESM crossing card is now canonical once its physical boundary
channels are named.  A channelwise forward/backward asymmetry is not charged
merely because it is nonzero.  The first debit is the part of the two
same-side Wilson blocks which mixes the declared channels.  If that debit is
small relative to the exact anchor gap and crossing floor, the true sign is
certified; otherwise the direct V13 sign is evaluated.

The supplied regulator establishes finite population of the field shell and
its grammar but does not print the boundary-channel algebra or numerical
same-history blocks required by \eqref{eq:V15-acquisition-card}.  The active
status is
\begin{equation}
 \boxed{
 \begin{array}{cl}
 \passstatus:&\text{canonical affine collar and two exact controls;}\\[0.25em]
 \obsstatus:&\text{post hoc channel fitting;}\\[0.25em]
 \stopstatus:&\text{first physical RPESM channel/block/direct-sign tuple.}
 \end{array}}
 \label{eq:V15-status}
\end{equation}
No exterior, Duhamel, represented current/stress, charge, interacting
continuum, or Einstein conclusion is inferred.
\end{assessment}

\begin{openproblem}[Version V16: read the first physical boundary-channel tuple]
\label{op:V15-next}
Preserve Versions~V1--V15.  Do not vary the channel algebra, reflection,
scalar origin, or one-letter support after reading the Wilson blocks.  On one
predeclared active RPESM boundary record:
\begin{enumerate}[label=\textup{(V16.\arabic*)},leftmargin=5.8em]
\item export the actual site/spin/flavour boundary-channel algebra and prove
      that $S$ belongs to it;
\item populate $M_-$, $M_+$, $B=VS$, $L$, $R_\Theta$,
      $\varepsilon_\Theta$, and $W_1$ on one common support;
\item compute $C_\partial,Z_\partial$, the channel debit, anchor gap, polar
      residual, and crossing floor;
\item if the canonical collar clears, export its strict interval; otherwise
      evaluate the direct Wilson sign and the V9 cone card without changing
      the channel algebra;
\item acquire one held-out four-phase crossing Read and the independent
      determinant/Pfaffian divisor and orientation rows;
\item repeat the tuple at one adjacent cutoff with the transported channel
      algebra.
\end{enumerate}
If no physical channel record or numerical Wilson block is supplied, the exact
verdict remains \stopstatus{} at that row.  Another fitted collar is not an
advance.
\end{openproblem}

\section*{V15 exact replay and append-only record}
\addcontentsline{toc}{section}{V15 exact replay and append-only record}
The accompanying verifier uses Python standard-library exact rational
arithmetic.  It checks the three-colour affine sign, gap, crossing, source
effect, and leakage; the two incompatible channel-algebra certificates; the
canonical off-channel debit; the centred anticommutation identity; and every
rational inequality yielding the $1/3$ lower crossing floor.  No numerical
diagonalization is used in the proof certificate.  Ordinary and optimized
executions produce byte-identical JSON output.

The complete Version~V14 source preceding its sole final active
\verb|\end{document}| is preserved byte for byte.  Its preterminal SHA--256
digest is recorded in the validation manifest.  A later continuation must
remove only the final active document-closing command, append new material
immediately before it, restore one terminal command, and increment the
filename to end in \texttt{\_V16.tex}.

% =============================================================================
% END APPENDED VERSION V15 MATERIAL
% =============================================================================


% =============================================================================
% START APPENDED VERSION V16 MATERIAL
% =============================================================================

\clearpage
\hypersetup{
 pdftitle={Renewal Einstein--Standard--Model Physical Source Metric, Duhamel Ancestry and Quantum Continuum Programme V16},
 pdfsubject={Direct boundary-cyclic Wilson-sign reduction, exact Feshbach crossing, and finite local polynomial certification}
}
\pdfinfo{
 /Title (Renewal Einstein--Standard--Model Physical Source Metric, Duhamel Ancestry and Quantum Continuum Programme V16)
 /Subject (Direct boundary-cyclic Wilson-sign reduction, exact Feshbach crossing, and finite local polynomial certification)
}

\section{Version V16: the direct Wilson-sign boundary germ precedes every fitted collar}
\label{sec:V16-direction}

Versions~V1--V15 are retained above.  Version~V15 identified a useful
channelwise affine collar and proved that its channel algebra must be fixed by
physical labels rather than chosen after the Wilson blocks are observed.  No
such labelled RPESM channel tuple is present in the supplied finite regulator
specification.  Continuing to enlarge the collar family would therefore be
circular.

The object required by the inherited reflected-hopping gate is nevertheless
smaller than a full diagonalization of the Wilson kernel.  It is the single
cross-boundary block
\begin{equation}
 X_{\operatorname{sgn}}
 :=T_+\operatorname{sgn}(H_W)T_-.
 \label{eq:V16-sign-crossing-target}
\end{equation}
The present continuation constructs this block directly from the actual
same-history Wilson kernel.  It uses two equivalent finite entrances:
\begin{equation}
 \boxed{
 \begin{gathered}
 \text{imaginary-axis boundary resolvents and their Feshbach shorts},\\
 \text{or a boundary-cyclic Wilson-word panel with a certified sign tail}.
 \end{gathered}}
 \label{eq:V16-two-direct-entrances}
\end{equation}
Neither entrance requires a boundary-channel algebra.  The V15 collar remains
an optional sufficient certificate when a physical channel record is actually
available.

\begin{firewall}[No third collar and no inferred Wilson matrix]
\label{fire:V16-no-third-collar}
The present section does not select another algebra in which the same-side
blocks are approximately diagonal.  It also does not invent the numerical
entries of the active RPESM Wilson kernel.  It proves that once the local
finite matrix $H_W(U)$, the chronological split, and the physical reflection
are exported on one history, the exact crossing is determined by a finite
boundary-cyclic compression or by a convergent local polynomial card with an
explicit outward radius.
\end{firewall}

% =============================================================================
\section{Exact imaginary-axis Feshbach formula for the crossing sign}
\label{sec:V16-Feshbach-sign}
% =============================================================================

Let $\mathcal K=\mathcal K_+\oplus\mathcal K_-$ be a finite Hilbert space with
orthogonal chronological projections $T_+,T_-$.  Let
\begin{equation}
 H=H^*=\begin{pmatrix}A&B\\B^*&D\end{pmatrix},
 \qquad
 gI\preceq |H|\preceq MI,
 \qquad 0<g\le M<\infty.
 \label{eq:V16-H-blocks}
\end{equation}
The block order is $\mathcal K_+\oplus\mathcal K_-$.  For
$z\in\mathbb C\setminus\mathbb R$, put
\begin{equation}
 \mathscr S_+(z)
 :=A-z-B(D-z)^{-1}B^*.
 \label{eq:V16-Schur-z}
\end{equation}

\begin{theorem}[Exact boundary-resolvent representation of the Wilson sign]
\label{thm:V16-Feshbach-sign}
For every $t>0$, $D\mp it$ and $\mathscr S_+(\pm it)$ are invertible and
\begin{equation}
 \boxed{
 T_+(H-z)^{-1}T_-
 =-\mathscr S_+(z)^{-1}B(D-z)^{-1},
 \qquad z\in\{it,-it\}.}
 \label{eq:V16-resolvent-cross-block}
\end{equation}
Moreover the desired crossing has the norm-convergent integral
\begin{equation}
 \boxed{
 \begin{aligned}
 X_{\operatorname{sgn}}
 =-\frac1\pi\int_0^\infty
 \Big[&\mathscr S_+(it)^{-1}B(D-it)^{-1}\\
      &+\mathscr S_+(-it)^{-1}B(D+it)^{-1}\Big]dt .
 \end{aligned}}
 \label{eq:V16-Feshbach-sign-integral}
\end{equation}
Thus the direct sign crossing is determined by the same-history blocks
$(A,B,D)$ and their supported imaginary-axis solves.  No fitted channel
algebra, determinant phase, or fermion covariance is required.
\end{theorem}

\begin{proof}
For nonreal $z$, both $H-z$ and $D-z$ are invertible.  Standard block Gaussian
elimination gives
\[
 \begin{pmatrix}I&-B(D-z)^{-1}\\0&I\end{pmatrix}
 (H-z)
 \begin{pmatrix}I&0\\-(D-z)^{-1}B^*&I\end{pmatrix}
 =\begin{pmatrix}\mathscr S_+(z)&0\\0&D-z\end{pmatrix}.
\]
The left side is invertible, so $\mathscr S_+(z)$ is invertible.  Reading the
upper-right block of the inverse proves
\eqref{eq:V16-resolvent-cross-block}.

For a nonzero real scalar $\lambda$,
\[
 \frac1\pi\int_0^\infty
 \left((\lambda-it)^{-1}+(\lambda+it)^{-1}\right)dt
 =\frac1\pi\int_0^\infty\frac{2\lambda}{\lambda^2+t^2}dt
 =\operatorname{sgn}\lambda.
\]
The finite-dimensional spectral theorem therefore gives
\begin{equation}
 \operatorname{sgn}H
 =\frac1\pi\int_0^\infty
 \left((H-it)^{-1}+(H+it)^{-1}\right)dt.
 \label{eq:V16-sign-resolvent-full}
\end{equation}
Near $t=0$ the integrand is bounded by the gap, and its off-diagonal block is
$O(t^{-2})$ as $t\to\infty$, because the $t^{-1}$ term of the resolvent is a
scalar multiple of the identity.  Hence the crossing block is norm
integrable.  Insert \eqref{eq:V16-resolvent-cross-block} into
\eqref{eq:V16-sign-resolvent-full}.
\end{proof}

\begin{assessment}[What is now the literal numerical row]
\label{ass:V16-resolvent-row}
At a selected finite RPESM history, a validated family of the two block solves
in \eqref{eq:V16-Feshbach-sign-integral}, together with a quadrature tail, is
already a complete direct crossing card.  The same solves can be restricted to
the boundary-cyclic carrier constructed next.  Failure of a proposed Schur
solve to reproduce one held-out full resolvent entry is an input-incidence
error, not a negative hopping mode.
\end{assessment}

% =============================================================================
\section{The exact boundary-cyclic carrier}
\label{sec:V16-boundary-cyclic}
% =============================================================================

Define the two same-side Krylov carriers
\begin{equation}
 \mathcal C_+
 :=\operatorname{span}\{A^k\operatorname{Ran}B:k\ge0\},
 \qquad
 \mathcal C_-
 :=\operatorname{span}\{D^k\operatorname{Ran}B^*:k\ge0\},
 \label{eq:V16-Cpm}
\end{equation}
and let $Q_+^\partial,Q_-^\partial$ be their orthogonal projections.

\begin{theorem}[Boundary-cyclic reduction of every cross-boundary function]
\label{thm:V16-boundary-cyclic}
The spaces $\mathcal C_+$ and $\mathcal C_-$ reduce $A$ and $D$, respectively,
and
\begin{equation}
 \mathcal C_\partial:=\mathcal C_+\oplus\mathcal C_-
 \label{eq:V16-Cpartial}
\end{equation}
reduces $H$.  Moreover
\begin{equation}
 \boxed{B=Q_+^\partial BQ_-^\partial,}
 \label{eq:V16-B-supported}
\end{equation}
and on $\mathcal C_\partial^\perp$ the operator $H$ is block diagonal.  Hence,
for every bounded Borel function $f$ on $\operatorname{spec}H$,
\begin{equation}
 \boxed{
 T_+f(H)T_-
 =Q_+^\partial T_+f(H)T_-Q_-^\partial.}
 \label{eq:V16-function-boundary-support}
\end{equation}
In particular the complete Wilson-sign crossing is obtained from the finite
compression $H|_{\mathcal C_\partial}$, and
\begin{equation}
 \boxed{
 \operatorname{rank}X_{\operatorname{sgn}}
 \le\min\{\dim\mathcal C_+,\dim\mathcal C_-\}.}
 \label{eq:V16-cross-rank-bound}
\end{equation}
No same-side spectator outside this carrier can affect the reflected hopping
verdict.
\end{theorem}

\begin{proof}
The space $\mathcal C_+$ is invariant under $A$ by construction.  Since $A$ is
self-adjoint, its orthogonal complement is invariant as well, so $\mathcal C_+$
reduces $A$; the same argument applies to $\mathcal C_-$ and $D$.  For every
$u\in\mathcal K_-$, $Bu\in\operatorname{Ran}B\subseteq\mathcal C_+$, and for
every $v\in\mathcal K_+$, $B^*v\in\operatorname{Ran}B^*\subseteq\mathcal C_-$.  Thus
$\mathcal C_\partial$ is invariant under $H$ and therefore reduces it.

Because $\mathcal C_-\supseteq\operatorname{Ran}B^*$, one has
$\mathcal C_-^\perp\subseteq\ker B$.  Similarly
$\mathcal C_+^\perp\subseteq\ker B^*$.  This proves
\eqref{eq:V16-B-supported} and shows that the restriction to
$\mathcal C_\partial^\perp$ has no chronological cross block.  Functional
calculus preserves the reducing decomposition, proving
\eqref{eq:V16-function-boundary-support}.  The rank bound is immediate.
\end{proof}

\begin{corollary}[The carrier is generated by the interface seed]
\label{cor:V16-interface-seed}
Let $J_\partial:E_\partial\to\mathcal K$ be any isometry with
\begin{equation}
 \operatorname{Ran}J_\partial
 =\operatorname{Ran}B\oplus\operatorname{Ran}B^*.
 \label{eq:V16-Jboundary-range}
\end{equation}
Then
\begin{equation}
 \boxed{
 \mathcal C_\partial
 =\operatorname{span}\{H^k\operatorname{Ran}J_\partial:k\ge0\}.}
 \label{eq:V16-interface-cyclic-equality}
\end{equation}
\end{corollary}

\begin{proof}
The right side is contained in $\mathcal C_\partial$ because that space is
$H$-invariant and contains the seed.  Conversely, suppose
$A^k\operatorname{Ran}B$ and $D^k\operatorname{Ran}B^*$ lie in the right side.
Applying $H$ to $(A^kv,0)$ gives
$(A^{k+1}v,B^*A^kv)$.  The second component belongs to the seed
$\operatorname{Ran}B^*$ and may be subtracted, leaving
$(A^{k+1}v,0)$.  The negative-side induction is identical.  Hence all the
generators in \eqref{eq:V16-Cpm} lie in the interface-cyclic span.
\end{proof}

% =============================================================================
\section{A finite exact Krylov-moment acquisition}
\label{sec:V16-Krylov}
% =============================================================================

For $r\ge0$ define
\begin{equation}
 W_r(c_0,\ldots,c_r)
 :=\sum_{j=0}^rH^jJ_\partial c_j,
 \qquad
 G_r=W_r^*W_r,
 \label{eq:V16-Wr-Gr}
\end{equation}
\begin{equation}
 K_r=W_r^*HW_r,
 \qquad
 E_{\pm,r}=W_r^*T_\pm W_r.
 \label{eq:V16-Kr-Er}
\end{equation}
Every block of these matrices is a finite interface moment:
\begin{equation}
 \begin{aligned}
 (G_r)_{ij}&=J_\partial^*H^{i+j}J_\partial,\\
 (K_r)_{ij}&=J_\partial^*H^{i+j+1}J_\partial,\\
 (E_{\pm,r})_{ij}&=J_\partial^*H^iT_\pm H^jJ_\partial.
 \end{aligned}
 \label{eq:V16-moment-blocks}
\end{equation}

\begin{theorem}[First-flat interface moment theorem]
\label{thm:V16-first-flat-Krylov}
Suppose
\begin{equation}
 \boxed{\operatorname{rank}G_{r+1}=\operatorname{rank}G_r.}
 \label{eq:V16-first-flat-rank}
\end{equation}
Then $\operatorname{Ran}W_r=\mathcal C_\partial$.  Put
\begin{equation}
 P_r=G_rG_r^\dagger,
 \qquad
 U_r=W_rG_r^{\dagger/2},
 \label{eq:V16-Ur}
\end{equation}
viewed as a unitary from $\operatorname{Ran}P_r$ onto
$\mathcal C_\partial$, and define
\begin{equation}
 \widehat H_r=G_r^{\dagger/2}K_rG_r^{\dagger/2},
 \qquad
 \widehat T_{\pm,r}=G_r^{\dagger/2}E_{\pm,r}G_r^{\dagger/2}.
 \label{eq:V16-reconstructed-H-T}
\end{equation}
On $\operatorname{Ran}P_r$, the operators $\widehat T_{+,r}$ and
$\widehat T_{-,r}$ are complementary orthogonal projections,
$\widehat H_r$ is self-adjoint and gapped by the same $g$, and
\begin{equation}
 \boxed{
 X_{\operatorname{sgn}}
 =U_r\widehat T_{+,r}
   \operatorname{sgn}(\widehat H_r)
   \widehat T_{-,r}U_r^*.}
 \label{eq:V16-exact-Krylov-sign}
\end{equation}
Thus one first-flat finite moment panel reconstructs the exact direct crossing,
including every same-side return to the interface.  No full-volume spectator
basis is needed.
\end{theorem}

\begin{proof}
Equality of the two ranks is equality of the nested ranges
$\operatorname{Ran}W_r\subseteq\operatorname{Ran}W_{r+1}$.  Hence
$H^{r+1}\operatorname{Ran}J_\partial\subseteq\operatorname{Ran}W_r$, and
therefore $H\operatorname{Ran}W_r\subseteq\operatorname{Ran}W_r$.  The range
contains the interface seed, so
\Cref{cor:V16-interface-seed} gives
$\operatorname{Ran}W_r=\mathcal C_\partial$.

The Moore--Penrose identities give
$U_r^*U_r=P_r$ and $U_rU_r^*=Q_\partial$, the projection onto
$\mathcal C_\partial$.  Consequently
\[
 \widehat H_r=U_r^*HU_r,
 \qquad
 \widehat T_{\pm,r}=U_r^*T_\pm U_r.
\]
The boundary-cyclic carrier reduces $H$ and $T_\pm$, so these are the unitary
compressions of the corresponding operators.  Functional calculus commutes
with unitary equivalence.  Apply
\Cref{thm:V16-boundary-cyclic} to the sign function.
\end{proof}

\begin{remark}[Exact stopping and numerical rank]
\label{rem:V16-rank-stop}
At exact finite precision, \eqref{eq:V16-first-flat-rank} is a terminating
criterion.  With interval data, a submitted support is accepted only after its
positive singular-value floor and the held-out relation
$H\operatorname{Ran}W_r\subseteq\operatorname{Ran}W_r$ are certified.  An
ambiguous numerical rank is \stopstatus{}, not a license to delete the weak
Krylov direction.
\end{remark}

% =============================================================================
\section{A certified local polynomial head for the sign}
\label{sec:V16-polynomial}
% =============================================================================

The exact Krylov rank may be large.  A second finite entrance uses only the
Wilson gap and norm window.  Put
\begin{equation}
 \alpha=\frac{M^2+g^2}{2},
 \qquad
 \rho=\frac{M^2-g^2}{M^2+g^2}\in[0,1),
 \label{eq:V16-alpha-rho}
\end{equation}
\begin{equation}
 b_k=(-1)^k4^{-k}\binom{2k}{k},
 \qquad
 p_N(x)=\frac{x}{\sqrt\alpha}
 \sum_{k=0}^{N}b_k
 \left(\frac{x^2-\alpha}{\alpha}\right)^k.
 \label{eq:V16-pN}
\end{equation}
This is an odd polynomial of degree $2N+1$.

\begin{theorem}[Explicit gap-certified Wilson-sign polynomial]
\label{thm:V16-sign-polynomial}
For every self-adjoint $H$ satisfying \eqref{eq:V16-H-blocks},
\begin{equation}
 \boxed{
 \|\operatorname{sgn}H-p_N(H)\|_{\rm op}
 \le\epsilon_N(g,M)
 :=\frac{M}{\sqrt\alpha}
   \frac{\rho^{N+1}}{1-\rho}.}
 \label{eq:V16-sign-poly-error}
\end{equation}
The bound is uniform over all matrix dimensions and all eigenvector
orientations.  When $g=M$, one has $\rho=0$ and $p_0(H)=\operatorname{sgn}H$
exactly.
\end{theorem}

\begin{proof}
For $x\in[-M,-g]\cup[g,M]$, write
$x^2=\alpha(1+u)$, so $|u|\le\rho$.  The convergent binomial expansion is
\[
 (1+u)^{-1/2}=\sum_{k=0}^\infty b_ku^k,
 \qquad |b_k|=4^{-k}\binom{2k}{k}\le1.
\]
Therefore
\[
 |\operatorname{sgn}x-p_N(x)|
 \le\frac{|x|}{\sqrt\alpha}
      \sum_{k=N+1}^\infty|u|^k
 \le\frac{M}{\sqrt\alpha}
      \frac{\rho^{N+1}}{1-\rho}.
\]
Apply the spectral theorem.
\end{proof}

\begin{corollary}[Finite-range Wilson kernels give finite boundary collars]
\label{cor:V16-finite-range-collar}
Suppose the one-particle carrier is supported on a finite graph and $H$ has
propagation range $R$: $H_{xy}=0$ whenever the graph distance exceeds $R$.
Then
\begin{equation}
 T_+p_N(H)T_-
 \label{eq:V16-local-poly-cross}
\end{equation}
is supported inside the $(2N+1)R$ neighbourhood of the chronological
interface on both sides.  Consequently a degree-$(2N+1)$ direct crossing card
requires only that finite Wilson collar, while
\eqref{eq:V16-sign-poly-error} controls the omitted nonlocal sign tail.
\end{corollary}

\begin{proof}
Every monomial in $p_N$ has degree at most $2N+1$.  A matrix element of $H^k$
can be nonzero only when the two graph sites are joined by a path of at most
$k$ steps of length $R$.  A path crossing from one chronological half to the
other must meet the interface, so no matrix element from a site deeper than
$kR$ on either side can occur.  Sum the finitely many monomials.
\end{proof}

\begin{assessment}[The physical locality gain]
\label{ass:V16-locality-gain}
The exact overlap sign is nonlocal, but the protected Wilson gap turns its
crossing into a controlled sequence of finite boundary collars.  This is the
target-native locality needed by the first hopping gate.  It is not yet the
weighted pair-collar or quasilocal exterior-kernel theorem required later in
the terminal quantum order.
\end{assessment}

% =============================================================================
\section{Direct finite cone certificate and relation descent}
\label{sec:V16-cone-card}
% =============================================================================

Let $R_\Theta:\mathcal K_+\to\mathcal K_-$ be the physical unitary reflection
map, let $L\succeq0$ be the fixed Hodge--chiral factor on $\mathcal K_+$, let
$a>0$, and let $W_1:\mathcal F_1\to\mathcal K_+$ be the occurring one-letter
synthesis.  Define
\begin{equation}
 K_F
 =-\frac{\varepsilon_\Theta}{a}
 W_1^*L T_+\operatorname{sgn}(H)T_-R_\Theta W_1,
 \label{eq:V16-KF-exact}
\end{equation}
\begin{equation}
 K_{F,N}
 =-\frac{\varepsilon_\Theta}{a}
 W_1^*L T_+p_N(H)T_-R_\Theta W_1.
 \label{eq:V16-KF-N}
\end{equation}
Let
\begin{equation}
 P_W=W_1^\dagger W_1,
 \qquad r_W=\operatorname{rank}P_W,
 \label{eq:V16-PW-rW}
\end{equation}
and put
\begin{equation}
 \eta_N^{\rm op}
 =\frac{\|W_1\|_{\rm op}^2\|L\|_{\rm op}}{a}
  \epsilon_N(g,M),
 \qquad
 \eta_N^{\rm HS}=\sqrt{r_W}\,\eta_N^{\rm op}.
 \label{eq:V16-etaN}
\end{equation}

\begin{theorem}[Finite direct Wilson-sign hopping certificate]
\label{thm:V16-direct-hopping-certificate}
The exact and approximate raw crossings automatically descend through the
one-letter relation space:
\begin{equation}
 \boxed{K_F=P_WK_FP_W,
 \qquad K_{F,N}=P_WK_{F,N}P_W.}
 \label{eq:V16-automatic-relation}
\end{equation}
Their errors satisfy
\begin{equation}
 \boxed{
 \|K_F-K_{F,N}\|_{\rm op}\le\eta_N^{\rm op},
 \qquad
 \|K_F-K_{F,N}\|_{\rm HS}\le\eta_N^{\rm HS}.}
 \label{eq:V16-KF-error}
\end{equation}
Let $\operatorname{Herm}_+(P_W)$ be the cone of positive Hermitian matrices
supported on $P_W$.
\begin{enumerate}[label=\textup{(C\arabic*)},leftmargin=4.2em]
\item If the independently verified reflection covariance makes $K_F$
      Hermitian and
      \begin{equation}
       \lambda_{\min}\!\left(
       P_W\frac{K_{F,N}+K_{F,N}^*}{2}P_W
       \big|_{\operatorname{Ran}P_W}\right)
       >\eta_N^{\rm op},
       \label{eq:V16-pass-criterion}
      \end{equation}
      then $K_F\succ0$ on the represented one-letter carrier, with the
      displayed eigenvalue minus $\eta_N^{\rm op}$ as a certified floor.
\item If
      \begin{equation}
       \operatorname{dist}_{\rm HS}
       (K_{F,N},\operatorname{Herm}_+(P_W))
       >\eta_N^{\rm HS},
       \label{eq:V16-obstruction-criterion}
      \end{equation}
      then the physical $K_F$ is outside the passing cone and the inherited V9
      skew/negative-part witness is a strict obstruction.
\item If neither strict inequality holds, the finite direct card is
      \stopstatus{} at its current degree.  Increasing $N$ tightens the same
      physical card; it does not alter the Wilson action or fit a new collar.
\end{enumerate}
\end{theorem}

\begin{proof}
Since $W_1=W_1P_W$ and $W_1^*=P_WW_1^*$, both relation identities are exact.
The operator error is the ideal property of the operator norm together with
\Cref{thm:V16-sign-polynomial}; the Hilbert--Schmidt bound follows because the
difference is supported on a space of rank $r_W$.  The first branch is Weyl's
inequality on $\operatorname{Ran}P_W$.  Distance from a closed convex set in a
Hilbert space is one-Lipschitz, so
\eqref{eq:V16-obstruction-criterion} excludes membership of $K_F$ in the
positive cone.  The remaining interval intersects the cone boundary and is
therefore unresolved.
\end{proof}

\begin{firewall}[Polynomial positivity is a certificate, not a replacement]
\label{fire:V16-polynomial-not-repair}
The polynomial $p_N(H)$ is not substituted for the physical overlap operator.
It supplies an interval containing the exact Wilson-sign crossing.  A negative
or skew head may be hidden by the tail unless
\eqref{eq:V16-obstruction-criterion} is strict; a positive head is promoted
only when \eqref{eq:V16-pass-criterion} clears the full tail radius.
\end{firewall}

% =============================================================================
\section{Adjacent-cutoff transport of the direct sign card}
\label{sec:V16-transport}
% =============================================================================

Expand the fixed polynomial in ordinary monomials,
\begin{equation}
 p_N(x)=\sum_{k=0}^{2N+1}a_{k,N}x^k,
 \qquad
 \mathfrak L_N(M)
 :=\sum_{k=1}^{2N+1}k|a_{k,N}|M^{k-1}.
 \label{eq:V16-poly-Lipschitz}
\end{equation}

\begin{theorem}[Direct-sign cutoff rectangle]
\label{thm:V16-sign-transport}
Let $H_0,H_1$ be self-adjoint Wilson kernels, possibly on different finite
carriers, both satisfying the same window
$gI\preceq|H_j|\preceq MI$.  Let $J$ be an isometry respecting the
chronological split and put
\begin{equation}
 \delta_H=\|H_1J-JH_0\|_{\rm op}.
 \label{eq:V16-deltaH}
\end{equation}
Then
\begin{equation}
 \boxed{
 \|\operatorname{sgn}(H_1)J-J\operatorname{sgn}(H_0)\|_{\rm op}
 \le2\epsilon_N(g,M)+\mathfrak L_N(M)\delta_H.}
 \label{eq:V16-sign-transport-bound}
\end{equation}
The same bound applies to the transported cross block.  If the Hodge factor,
reflection, Berezin sign, and one-letter synthesis are transported exactly,
the represented crossing defect is at most
\begin{equation}
 \boxed{
 \frac{\|W_1\|_{\rm op}^2\|L\|_{\rm op}}a
 \left(2\epsilon_N(g,M)+\mathfrak L_N(M)\delta_H\right).}
 \label{eq:V16-K-transport-bound}
\end{equation}
Their actual transport defects are added separately when they are nonzero.
\end{theorem}

\begin{proof}
For every $k\ge1$,
\[
 H_1^kJ-JH_0^k
 =\sum_{j=0}^{k-1}
 H_1^{k-1-j}(H_1J-JH_0)H_0^j.
\]
Hence
$\|H_1^kJ-JH_0^k\|\le kM^{k-1}\delta_H$, and summing the monomial coefficients
gives
\[
 \|p_N(H_1)J-Jp_N(H_0)\|
 \le\mathfrak L_N(M)\delta_H.
\]
Insert and subtract the two polynomial approximants and use
\Cref{thm:V16-sign-polynomial}.  Split compatibility gives the cross-block
bound, and multiplication by the exactly transported physical factors gives
\eqref{eq:V16-K-transport-bound}.
\end{proof}

\begin{corollary}[Cofinal direct crossing without channel transport]
\label{cor:V16-cofinal-direct-sign}
Along a cofinal sequence, choose degrees $N_n$ using one transported gap/norm
window.  If
\begin{equation}
 \sum_n\left[
 \epsilon_{N_n}(g,M)
 +\mathfrak L_{N_n}(M)\delta_{H,n}
 +\delta_{\rm phys,n}
 \right]<\infty,
 \label{eq:V16-cofinal-summability}
\end{equation}
where $\delta_{\rm phys,n}$ contains the separately measured Hodge,
reflection, sign-orientation, and one-letter defects, then the represented
direct crossings are Cauchy on every fixed physical one-letter screen.  This
conclusion does not require a transported boundary-channel algebra.
\end{corollary}

\begin{proof}
Apply \Cref{thm:V16-sign-transport} at every adjacent pair and telescope.  The
additional physical factors obey the ordinary product-difference inequality.
Summability gives the Cauchy property.
\end{proof}

\begin{assessment}[A conditioning tradeoff]
\label{ass:V16-degree-route-tradeoff}
Increasing $N$ decreases the sign tail but may increase
$\mathfrak L_N(M)$.  A cofinal proof therefore requires route accuracy at the
same degree used for the local sign head.  Convergence of each fixed Wilson
word separately is not enough to justify a degree which grows with the
cutoff.
\end{assessment}

% =============================================================================
\section{Why a finite head needs either a tail or exact cyclic flatness}
\label{sec:V16-no-finite-universal}
% =============================================================================

\begin{countertheorem}[No universal exact finite polynomial sign card]
\label{cth:V16-no-finite-polynomial}
Fix $0<g<M$.  There is no finite-degree polynomial $p$ such that
\begin{equation}
 p(x)=\operatorname{sgn}x
 \quad\text{for every }x\in[-M,-g]\cup[g,M].
 \label{eq:V16-impossible-exact-poly}
\end{equation}
Consequently a fixed finite Wilson-word head cannot determine the sign exactly
for every gapped physical kernel.  Exact Krylov stabilization or a certified
analytic tail is indispensable.
\end{countertheorem}

\begin{proof}
If such a polynomial existed, $p-1$ would vanish on the interval $[g,M]$ and
hence would be the zero polynomial.  Thus $p\equiv1$, contradicting
$p=-1$ on $[-M,-g]$.
\end{proof}

\begin{firewall}[No extrapolation from a few Wilson moments]
\label{fire:V16-no-moment-extrapolation}
A sequence of low-order boundary moments may suggest a stable crossing, but it
is not an exact finite card until the Krylov rank has passed its held-out
flatness test.  On the nonflat branch the polynomial tail must be charged using
the same protected gap and norm window.  Neither visual convergence nor a
fitted spectral model replaces these alternatives.
\end{firewall}

% =============================================================================
\section{Exact finite controls}
\label{sec:V16-controls}
% =============================================================================

\subsection{A six-dimensional carrier with a four-dimensional exact boundary core}

Let
\begin{equation}
 A_\circ=
 \begin{pmatrix}2&1&0\\1&3&0\\0&0&5\end{pmatrix},
 \quad
 D_\circ=
 \begin{pmatrix}-3&1&0\\1&-2&0\\0&0&-6\end{pmatrix},
 \quad
 B_\circ=
 \begin{pmatrix}1&0&0\\0&0&0\\0&0&0\end{pmatrix},
 \label{eq:V16-six-blocks}
\end{equation}
\begin{equation}
 H_\circ=
 \begin{pmatrix}A_\circ&B_\circ\\B_\circ^*&D_\circ\end{pmatrix}.
 \label{eq:V16-Hcirc}
\end{equation}

\begin{proposition}[Exact spectator removal by the first-flat card]
\label{prop:V16-six-control}
The matrix $H_\circ$ is invertible with
\begin{equation}
 \det H_\circ=-930.
 \label{eq:V16-Hcirc-det}
\end{equation}
Its interface carriers are
\begin{equation}
 \boxed{
 \mathcal C_+=\operatorname{span}\{e_1,e_2\},
 \qquad
 \mathcal C_-=\operatorname{span}\{f_1,f_2\}.}
 \label{eq:V16-Hcirc-carriers}
\end{equation}
For the two-column interface seed $(e_1,f_1)$,
\begin{equation}
 \boxed{
 \operatorname{rank}G_0=2,
 \qquad
 \operatorname{rank}G_1=4,
 \qquad
 \operatorname{rank}G_2=4.}
 \label{eq:V16-Hcirc-ranks}
\end{equation}
Thus the $r=1$ moment card reconstructs the complete crossing exactly on a
four-dimensional carrier.  The two spectator vectors $e_3,f_3$ have zero
cross block for every bounded function of $H_\circ$, including its sign.
\end{proposition}

\begin{proof}
Direct elimination gives the determinant.  The seed and its first images are
\[
 e_1,\quad f_1,\quad
 H_\circ e_1=2e_1+e_2+f_1,\quad
 H_\circ f_1=e_1-3f_1+f_2,
\]
which span $e_1,e_2,f_1,f_2$.  That space is invariant, while $e_3$ and $f_3$
are uncoupled eigenvectors with eigenvalues $5$ and $-6$.  The rank statements
and the functional-calculus conclusion follow.
\end{proof}

\subsection{A rational direct sign certificate with a nonzero analytic tail}

Put
\begin{equation}
 H_*=\frac15
 \begin{pmatrix}-11&-48\\-48&61\end{pmatrix}.
 \label{eq:V16-Hstar}
\end{equation}

\begin{proposition}[Exact degree-twenty-one crossing certificate]
\label{prop:V16-rational-poly-control}
The matrix $H_*$ has spectrum
\begin{equation}
 \operatorname{spec}H_*=\{-7,17\},
 \qquad g=7,
 \qquad M=17,
 \label{eq:V16-Hstar-spectrum}
\end{equation}
and
\begin{equation}
 \boxed{
 \operatorname{sgn}H_*
 =\frac1{12}H_* -\frac5{12}I
 =\begin{pmatrix}-3/5&-4/5\\-4/5&3/5\end{pmatrix}.}
 \label{eq:V16-Hstar-sign}
\end{equation}
With the upper coordinate positive-time, $a=L=R_\Theta=1$, and
$\varepsilon_\Theta=1$, the exact reflected crossing is $K_*=4/5$.

For $N=10$, \eqref{eq:V16-alpha-rho} gives
\begin{equation}
 \alpha=169,
 \qquad \rho=\frac{120}{169},
 \label{eq:V16-Hstar-alpha}
\end{equation}
and the degree-twenty-one polynomial card gives
\begin{equation}
 \boxed{
 K_{*,10}
 =\frac{75882723393852778807296}
        {95024818874403997194005}
 =0.7985568853\ldots .}
 \label{eq:V16-Hstar-K10}
\end{equation}
Its actual error is
\begin{equation}
 \left|K_{*,10}-\frac45\right|
 =\frac{137131705670418947908}
        {95024818874403997194005}
 =0.0014431146\ldots,
 \label{eq:V16-Hstar-actual-error}
\end{equation}
while the universal theorem radius is the larger exact number
\begin{equation}
 \epsilon_{10}(7,17)
 =\frac{1263114230169600000000000}
        {12106161924599069242516237}
 =0.1043364724\ldots .
 \label{eq:V16-Hstar-radius}
\end{equation}
Nevertheless the fully outward-certified lower crossing satisfies
\begin{equation}
 \boxed{
 K_{*,10}-\epsilon_{10}(7,17)
 =\frac{42021723651036220100247552}
        {60530809622995346212581185}
 >\frac{69}{100}.}
 \label{eq:V16-Hstar-certified-floor}
\end{equation}
Thus one finite local polynomial head, together with the protected gap and norm
window, certifies a strict physical crossing without any channel fit.
\end{proposition}

\begin{proof}
The characteristic polynomial is
$(\lambda+7)(\lambda-17)$.  The affine polynomial
$\lambda/12-5/12$ takes values $-1$ and $1$ on the two eigenvalues, proving
\eqref{eq:V16-Hstar-sign}.  The remaining identities are exact evaluation of
\eqref{eq:V16-pN} in rational arithmetic.  The final strict inequality is the
positive rational difference
\[
 \frac{5109300223388624271330687}
      {1210616192459906924251623700}>0.
\]
\end{proof}

% =============================================================================
\section{Version V16 theorem, ledger, and next literal acquisition}
\label{sec:V16-terminal}
% =============================================================================

\begin{theorem}[Version V16 direct boundary-cyclic Wilson-sign theorem]
\label{thm:V16-terminal}
Versions~V1--V15 are retained.  Version~V16 additionally proves:
\begin{enumerate}[label=\textup{(E16.\arabic*)},leftmargin=5.7em]
\item the exact cross-boundary Wilson sign is the imaginary-axis Feshbach
      integral \eqref{eq:V16-Feshbach-sign-integral};
\item the smallest carrier capable of affecting that crossing is the pair of
      same-side Krylov spaces generated by $\operatorname{Ran}B$ and
      $\operatorname{Ran}B^*$;
\item every same-side spectator outside this carrier is exactly irrelevant to
      every cross-boundary Borel function;
\item the first flat interface-moment Gram reconstructs the exact compressed
      Wilson kernel, chronological projections, and sign crossing;
\item the explicit odd polynomial \eqref{eq:V16-pN} has the uniform error
      \eqref{eq:V16-sign-poly-error};
\item for a finite-range Wilson kernel its degree-$2N+1$ head is supported on a
      finite $(2N+1)R$ chronological collar;
\item the represented direct crossing automatically descends through the
      occurring one-letter relations and has the finite pass/obstruction/stop
      semantics of \Cref{thm:V16-direct-hopping-certificate};
\item adjacent direct signs obey the weighted polynomial rectangle
      \eqref{eq:V16-sign-transport-bound};
\item no fixed finite polynomial can determine the sign exactly on a nontrivial
      gap window;
\item the six-dimensional control reduces exactly to a four-dimensional
      boundary carrier, and the rational two-dimensional card certifies a
      crossing floor greater than $69/100$.
\end{enumerate}
No item supplies the numerical active RPESM Wilson matrix, its selected gauge
history, physical reflection frame, determinant/Pfaffian orientation, or a
cofinal crossing floor.
\end{theorem}

\begin{proof}
Items~\textup{(E16.1)--(E16.4)} are
\Cref{thm:V16-Feshbach-sign,thm:V16-boundary-cyclic,thm:V16-first-flat-Krylov}.
Items~\textup{(E16.5)--(E16.6)} are
\Cref{thm:V16-sign-polynomial,cor:V16-finite-range-collar}.  Item
\textup{(E16.7)} is \Cref{thm:V16-direct-hopping-certificate}; item
\textup{(E16.8)} is \Cref{thm:V16-sign-transport}; item
\textup{(E16.9)} is \Cref{cth:V16-no-finite-polynomial}; and the controls are
\Cref{prop:V16-six-control,prop:V16-rational-poly-control}.
\end{proof}

\begingroup\small
\begin{longtable}{>{\raggedright\arraybackslash}p{0.27\textwidth}
                  >{\raggedright\arraybackslash}p{0.20\textwidth}
                  >{\raggedright\arraybackslash}p{0.45\textwidth}}
\caption{Version V16 proof-status ledger.}\label{tab:V16-ledger}\\
\toprule
\textbf{Row}&\textbf{Status}&\textbf{Exact scope}\\
\midrule
\endfirsthead
\toprule
\textbf{Row}&\textbf{Status}&\textbf{Exact scope}\\
\midrule
\endhead
V1--V15 source, ancestry, record, Grassmann, overlap, and collar results
&Imported exactly
&Preserved verbatim; no prior theorem or control is recomputed.\\[0.3em]
Imaginary-axis crossing formula
&\passstatus
&Exact same-history Feshbach representation of the direct Wilson sign.\\[0.3em]
Boundary-cyclic carrier
&\passstatus
&Unique smallest same-side Krylov pair containing every cross-boundary
functional-calculus contribution.\\[0.3em]
First-flat moment card
&\passstatus
&Exact finite reconstruction after a certified rank stabilization and held-out
invariance check.\\[0.3em]
Gap-certified polynomial head
&\passstatus
&Uniform operator radius; finite-range heads occupy finite chronological
collars.\\[0.3em]
Direct hopping verdict
&Terminating finite compiler
&Strict positive floor, strict cone obstruction, or unresolved tail on the same
physical card.\\[0.3em]
Universal fixed-degree exactness
&\obsstatus
&Impossible on a nontrivial gap interval; a tail or exact cyclic flatness is
mandatory.\\[0.3em]
Six-dimensional boundary core
&\passstatus
&First-flat rank four; two same-side spectators are exactly removed.\\[0.3em]
Rational degree-twenty-one control
&\passstatus
&Exact physical crossing $4/5$ and outward certified floor above $69/100$.\\[0.3em]
Selected RPESM Wilson history
&\stopstatus
&Its explicit local Wilson matrix, split, reflection, one-letter bank, and line
orientation remain unpopulated.\\[0.3em]
Exterior, Duhamel, current/stress, charge, continuum, Einstein
&Not reopened
&All remain downstream of a passing crossing and line packet.\\
\bottomrule
\end{longtable}
\endgroup

\begin{assessment}[Physical status after Version V16]
\label{ass:V16-status}
The channelwise V15 collar is no longer the only executable route.  Once one
actual finite RPESM Wilson matrix is exported, its direct crossing can be
computed without a channel choice by either the exact boundary-resolvent card
or the boundary-cyclic polynomial card.  The latter is local at every finite
degree and has a fully explicit tail.  The first unresolved datum is therefore
not another abstract positivity theorem but the matrix entries of one selected
same-history Wilson boundary collar and its physical reflection.

The current status is
\begin{equation}
 \boxed{
 \begin{array}{cl}
 \passstatus:&\text{direct Feshbach, cyclic, polynomial, and transport compilers;}\\[0.25em]
 \obsstatus:&\text{a fixed finite head without a certified tail or flat rank;}\\[0.25em]
 \stopstatus:&\text{first populated active RPESM local Wilson-sign card.}
 \end{array}}
 \label{eq:V16-status}
\end{equation}
No conditional exterior kernel, Duhamel birth, represented current/stress,
charge, interacting continuum, or Einstein conclusion is inferred.
\end{assessment}

\begin{openproblem}[Version V17: populate one direct local Wilson-sign card]
\label{op:V16-next}
Preserve Versions~V1--V16.  Do not introduce another fitted collar.  On one
predeclared active RPESM boundary history:
\begin{enumerate}[label=\textup{(V17.\arabic*)},leftmargin=5.8em]
\item export the actual finite local matrix $H_W(U)$, chronological projections,
      physical spin/gauge reflection, Berezin sign, Hodge factor, and one-letter
      synthesis on one common support;
\item certify the Wilson gap and norm window and extract the direct interface
      block $B=T_+H_WT_-$;
\item either populate the imaginary-axis Feshbach solves with a validated
      quadrature tail, or acquire boundary moments until exact Krylov flatness,
      or choose one polynomial degree before reading the crossing and charge
      \eqref{eq:V16-sign-poly-error};
\item run the direct V9 relation, skew, negative-part, and positive-factor-rank
      tests with the complete outward radius;
\item retain one held-out full-resolvent or four-phase crossing Read, and the
      determinant/Pfaffian divisor and orientation packet independently;
\item repeat the same direct card at one adjacent cutoff and evaluate
      \eqref{eq:V16-sign-transport-bound} with the separately transported
      physical factors.
\end{enumerate}
If the explicit local Wilson matrix is not supplied, the exact verdict remains
\stopstatus{} at that physical loading row.  A further abstract collar,
channel fit, exterior kernel, or Duhamel theorem is not an advance.
\end{openproblem}

\section*{V16 exact replay and append-only record}
\addcontentsline{toc}{section}{V16 exact replay and append-only record}
The accompanying verifier uses Python standard-library exact rational
arithmetic.  It checks the six-dimensional boundary-cyclic ranks and spectator
split; the exact quadratic identity and sign of $H_*$; every coefficient of
the degree-twenty-one polynomial; the exact physical crossing, actual error,
universal tail radius, and rational lower floor; and one exact Schur-resolvent
cross-block identity.  Ordinary and optimized executions produce byte-identical
JSON output.  The verifier does not populate an RPESM gauge history.

The complete Version~V15 source preceding its sole final active
\verb|\end{document}| is preserved byte for byte.  Its preterminal SHA--256
digest is recorded in the validation manifest.  A later continuation must
remove only the final active document-closing command, append new material
immediately before it, restore one terminal command, and increment the filename
to end in \texttt{\_V17.tex}.

% =============================================================================
% END APPENDED VERSION V16 MATERIAL
% =============================================================================



% =============================================================================
% START APPENDED VERSION V17 MATERIAL
% =============================================================================

\clearpage
\hypersetup{
 pdftitle={Renewal Einstein--Standard--Model Physical Source Metric, Duhamel Ancestry and Quantum Continuum Programme V17},
 pdfsubject={Wilson admissibility occurrence, divisor tubes, zero-safe sign regularization, and spectral-flow transport}
}
\pdfinfo{
 /Title (Renewal Einstein--Standard--Model Physical Source Metric, Duhamel Ancestry and Quantum Continuum Programme V17)
 /Subject (Wilson admissibility occurrence, divisor tubes, zero-safe sign regularization, and spectral-flow transport)
}

\section{Version V17: admissibility occurrence precedes the Wilson-sign card}
\label{sec:V17-direction}

Versions~V1--V16 are retained above.  Versions~V13--V16 reduced the first
fermionic Osterwalder--Schrader gate to the direct boundary block of the
Wilson sign and supplied exact Feshbach, boundary-cyclic, channelwise, and
polynomial compilers for that block.  All of those compilers begin on a
branch on which the selected same-history Wilson kernel is invertible with a
certified gap.

The literal finite field-shell construction appearing in the supplied
manuscript places the overlap coefficient on a protected compact chart and
assumes a Wilson window there.  It does not, by that assumption alone, prove
that the active full-coupling bosonic cylinder is supported in the protected
chart, that its Wilson divisor has zero mass, or that the divisor tubes have a
cutoff-uniform decay law.  Repeating the sign approximation while leaving this
incidence unresolved would be circular.

The present continuation therefore moves one arrow upstream and separates
four physically different branches:
\begin{equation}
 \boxed{
 \begin{gathered}
 \text{hard admissibility on the actual measure support},\\
 \text{almost-sure invertibility with controlled divisor tubes},\\
 \text{positive divisor mass and zero-mode saturation},\\
 \text{spectral-flow change of the overlap source sector}.
 \end{gathered}}
 \label{eq:V17-four-admissibility-branches}
\end{equation}
Only the first branch licenses the pointwise Version~V16 card unchanged.  The
second may support a record-integrated or mixed reflected form, but does not
by itself prove historywise conditional Gaussianity.  The final two branches
belong to the zero-safe line, exterior, and orientation packet rather than to
a pseudoinverted sign projector.

\begin{firewall}[No soft constraint is silently promoted to a hard overlap chart]
\label{fire:V17-no-soft-to-hard}
A positive Boltzmann weight, a compact gauge space, almost-sure invertibility,
a finite inverse moment, or a small probability of low Wilson gap does not
supply a uniform gap on the physical support.  Conditioning away a divisor
tube changes the unnormalized cylinder unless the corresponding selector is
already part of the physical action or record.  The change may be useful as a
separate regulated experiment, but it is not identified with the original
RPESM history by normalization alone.
\end{firewall}

% =============================================================================
\section{The hard-admissibility occurrence theorem}
\label{sec:V17-hard-occurrence}
% =============================================================================

Let $X$ be a compact metric bosonic history space, let $\lambda$ be a finite
Borel measure with full support, and let $w:X\to(0,\infty)$ be continuous.
Put
\begin{equation}
 Z=\int_Xw\,d\lambda,
 \qquad
 d\mu=Z^{-1}w\,d\lambda.
 \label{eq:V17-positive-cylinder}
\end{equation}
Let $\mathcal E$ be a finite-dimensional Hilbert space and let
\begin{equation}
 H_W:X\longrightarrow\operatorname{Herm}(\mathcal E)
 \label{eq:V17-Wilson-family}
\end{equation}
be continuous.  Define
\begin{equation}
 \delta_W(x):=s_{\min}(H_W(x))=\min\operatorname{spec}|H_W(x)|,
 \label{eq:V17-gap-function}
\end{equation}
\begin{equation}
 \mathscr D_W:=\{x:\delta_W(x)=0\},
 \qquad
 X_g:=\{x:\delta_W(x)\ge g\}.
 \label{eq:V17-divisor-hard-set}
\end{equation}

\begin{theorem}[Hard admissibility is a support statement]
\label{thm:V17-hard-occurrence}
The physical law in \eqref{eq:V17-positive-cylinder} has support $X$, and
\begin{equation}
 \boxed{
 \operatorname*{ess\,inf}_{\mu}\delta_W
 =\min_{x\in X}\delta_W(x).}
 \label{eq:V17-essinf-min}
\end{equation}
Consequently, the following are equivalent:
\begin{enumerate}[label=\textup{(\roman*)},leftmargin=3em]
\item there is $g>0$ such that $|H_W(x)|\succeq gI$ for $\mu$-almost every
      $x$;
\item there is $g>0$ such that $|H_W(x)|\succeq gI$ for every $x\in X$;
\item $\mathscr D_W=\varnothing$.
\end{enumerate}
If $\mathscr D_W\ne\varnothing$, or more generally if
$\min_X\delta_W=0$, then
\begin{equation}
 \boxed{
 \mu\{x:\delta_W(x)<g\}>0
 \quad\text{for every }g>0.}
 \label{eq:V17-every-tube-positive}
\end{equation}
Thus a Wilson window imposed only on a proper chart is not yet an occurrence
statement for the original full-support cylinder.
\end{theorem}

\begin{proof}
Every nonempty open set $O\subseteq X$ has $\lambda(O)>0$.  Continuity and
strict positivity of $w$ on the compact closure of a sufficiently small open
subset of $O$ give $\int_Ow\,d\lambda>0$.  Hence $\operatorname{supp}\mu=X$.
The function $\delta_W$ is continuous because singular values are
one-Lipschitz in the operator norm.  Let $m=\min_X\delta_W$.  Trivially
$\delta_W\ge m$ almost everywhere.  Conversely, for every $a>m$ there is
$x_a$ with $\delta_W(x_a)<a$; continuity produces a nonempty open
neighbourhood on which $\delta_W<a$, and this neighbourhood has positive
$\mu$-measure.  Therefore the essential infimum is at most every $a>m$,
proving \eqref{eq:V17-essinf-min}.

The equivalence of the first two statements follows from
\eqref{eq:V17-essinf-min}; compactness and continuity identify the second with
absence of a zero.  If the minimum is zero, the same neighbourhood argument
with $a=g$ proves \eqref{eq:V17-every-tube-positive}.
\end{proof}

\begin{corollary}[Hard conditioning changes the physical cylinder]
\label{cor:V17-hard-conditioning}
Let $g>0$ and suppose
\begin{equation}
 p_g:=\mu(X_g)>0.
 \label{eq:V17-hard-mass}
\end{equation}
The hard-conditioned law
\begin{equation}
 d\mu_g=p_g^{-1}\one_{X_g}\,d\mu
 \label{eq:V17-conditioned-law}
\end{equation}
is equal to $\mu$ if and only if $p_g=1$.  For every bounded operator-valued
writer $F$,
\begin{equation}
 \boxed{
 \left\|\int F\,d\mu_g-\int F\,d\mu\right\|_{\mathrm{op}}
 \le2\|F\|_{L^\infty(\mu;\mathrm{op})}(1-p_g).}
 \label{eq:V17-conditioning-bound}
\end{equation}
Therefore the selector $\one_{X_g}$ and its normalization are an additional
physical action/record row whenever $p_g<1$.
\end{corollary}

\begin{proof}
Write $A=X_g$ and $\epsilon=1-p_g$.  Then
\[
 \int F\,d\mu
 =p_g\int F\,d\mu_g
  +\epsilon\int F\,d\mu(\cdot\mid A^c).
\]
Subtract the first conditional mean and use the triangle inequality.  Equality
of the two measures forces $\mu(A^c)=0$, and the converse is immediate.
\end{proof}

\begin{remark}[Weights with zeros]
\label{rem:V17-weight-support}
If the line modulus or another physical factor vanishes, the same theorem
applies after replacing $X$ by the actual closed support of the weighted
cylinder.  What matters is the minimum Wilson gap on that support, not on an
ambient configuration space and not on a chart chosen after the measure is
formed.
\end{remark}

% =============================================================================
\section{Analytic divisors, tube probabilities, and inverse-gap moments}
\label{sec:V17-divisor-tubes}
% =============================================================================

Almost-sure invertibility and hard admissibility are separated even on a
finite-dimensional smooth gauge cylinder.

\begin{lemma}[A nontrivial analytic divisor has zero smooth volume]
\label{lem:V17-analytic-zero-set}
Let $X$ be a connected real-analytic manifold with a measure whose local
coordinate densities are positive and locally bounded.  If
$f:X\to\mathbb R$ is real analytic and not identically zero, then
\begin{equation}
 \mu\{f=0\}=0.
 \label{eq:V17-analytic-zero-measure}
\end{equation}
In particular, if the entries of $H_W$ are real analytic and
$\det H_W\not\equiv0$, then $\mu(\mathscr D_W)=0$.
\end{lemma}

\begin{proof}
It is enough to work in one relatively compact analytic coordinate box and
use induction on its dimension.  In one dimension, a nonzero analytic
function has isolated zeros.  In dimension $n$, write the convergent power
series in the last coordinate,
\[
 f(x',t)=\sum_{k\ge0}a_k(x')t^k.
\]
At least one analytic coefficient $a_k$ is not identically zero; otherwise
$f$ would vanish on the box and hence, by connected analytic continuation, on
all of $X$.  By the induction hypothesis, the set of $x'$ on which this
chosen coefficient vanishes has $(n-1)$-dimensional measure zero.  For every
remaining $x'$, the one-variable analytic function $t\mapsto f(x',t)$ is not
identically zero and hence has a zero set of one-dimensional measure zero.
Fubini's theorem proves the claim on the box.  A countable analytic atlas
finishes the proof.  Apply this to $f=\det H_W$.
\end{proof}

\begin{assessment}[Measure-zero divisor is not a Wilson window]
\label{ass:V17-zero-not-gap}
When the density is positive near a nonempty analytic divisor,
\Cref{lem:V17-analytic-zero-set} gives almost-sure invertibility while
\Cref{thm:V17-hard-occurrence} gives essential gap zero.  The overlap sign is
then defined almost everywhere but has no uniform polynomial tail, inverse
bound, or historywise crossing floor without an additional divisor-tube
estimate.
\end{assessment}

Let $0<\delta\le M$ be any positive random variable.  In the Wilson
application, $\delta=\delta_W$ on the almost-surely invertible branch and
$M$ may be any uniform upper bound for $\|H_W\|$.

\begin{theorem}[Exact tube--inverse-moment duality]
\label{thm:V17-tube-moment}
For every $q>0$,
\begin{equation}
 \boxed{
 \mathbb E\delta^{-q}
 =M^{-q}
  +q\int_0^M r^{-q-1}\,\mathbb P(\delta<r)\,dr,}
 \label{eq:V17-layer-cake}
\end{equation}
with the value $+\infty$ allowed.  Consequently:
\begin{enumerate}[label=\textup{(T\arabic*)},leftmargin=4em]
\item if $\mathbb E\delta^{-q}\le C_q$, then
      \begin{equation}
       \boxed{\mathbb P(\delta<r)\le C_qr^q}
       \label{eq:V17-Markov-tube}
      \end{equation}
      for every $r>0$;
\item if, for some $\alpha>0$,
      \begin{equation}
       \mathbb P(\delta<r)\le C_\alpha r^\alpha,
       \qquad 0<r\le M,
       \label{eq:V17-tube-power}
      \end{equation}
      then for every $0<q<\alpha$,
      \begin{equation}
       \boxed{
       \mathbb E\delta^{-q}
       \le M^{-q}
       +\frac{qC_\alpha}{\alpha-q}M^{\alpha-q}.}
       \label{eq:V17-moment-from-tube}
      \end{equation}
\end{enumerate}
Thus a uniform gap is the endpoint $q=\infty$ of a hierarchy whose finite
levels are measured by divisor-tube probabilities or inverse-gap moments.
\end{theorem}

\begin{proof}
For each $s\in(0,M]$,
\[
 s^{-q}=M^{-q}+q\int_s^Mr^{-q-1}\,dr.
\]
Integrate this identity with $s=\delta$ and use Tonelli's theorem; the
indicator $\one_{\{\delta<r\}}$ gives \eqref{eq:V17-layer-cake}.  The first
consequence is Markov's inequality applied to $\delta^{-q}$.  Insert
\eqref{eq:V17-tube-power} into the integral and evaluate the convergent power
integral when $q<\alpha$.
\end{proof}

\begin{firewall}[A finite inverse moment is not a pointwise sign certificate]
\label{fire:V17-moment-not-historywise}
A bound in \eqref{eq:V17-moment-from-tube} controls a measure-weighted low-gap
tail.  It does not imply that every bosonic history has an invertible Wilson
kernel, and it does not turn a positive mixed crossing into a positive
historywise crossing.  Conditional exterior Gaussianity continues to require
the appropriate recordwise or historywise gate.
\end{firewall}

% =============================================================================
\section{Zero-safe regularized signs and a quantitative soft-gap card}
\label{sec:V17-regularized-sign}
% =============================================================================

For $\eta>0$ and an arbitrary finite Hermitian matrix $H$, define
\begin{equation}
 \varepsilon_\eta(H)
 :=H(H^2+\eta^2I)^{-1/2},
 \qquad
 P_\eta(H):=\frac12(I+\varepsilon_\eta(H)).
 \label{eq:V17-regularized-sign}
\end{equation}
This is a bounded diagnostic on the divisor, but it is not an exact overlap
projector.

\begin{theorem}[Exact regularized-sign defect]
\label{thm:V17-regularized-sign}
For every Hermitian $H$,
\begin{equation}
 \boxed{
 I-\varepsilon_\eta(H)^2
 =\eta^2(H^2+\eta^2I)^{-1}\succeq0,}
 \label{eq:V17-involution-defect}
\end{equation}
\begin{equation}
 \boxed{
 P_\eta(H)-P_\eta(H)^2
 =\frac{\eta^2}{4}(H^2+\eta^2I)^{-1}\succeq0.}
 \label{eq:V17-projector-defect}
\end{equation}
If $H$ is invertible and $\delta=s_{\min}(H)$, then
\begin{equation}
 \boxed{
 \|\operatorname{sgn}H-\varepsilon_\eta(H)\|_{\mathrm{op}}
 =1-\frac{\delta}{\sqrt{\delta^2+\eta^2}}.}
 \label{eq:V17-exact-sign-error}
\end{equation}
In particular, on the event $\{\delta\ge r\}$,
\begin{equation}
 \boxed{
 \|\operatorname{sgn}H-\varepsilon_\eta(H)\|_{\mathrm{op}}
 \le1-\frac r{\sqrt{r^2+\eta^2}}
 \le\frac{\eta^2}{2r^2}.}
 \label{eq:V17-sign-error-r}
\end{equation}
\end{theorem}

\begin{proof}
All statements follow from scalar functional calculus.  For a real eigenvalue
$\lambda$,
\[
 \varepsilon_\eta(\lambda)^2
 =\frac{\lambda^2}{\lambda^2+\eta^2},
\]
which gives \eqref{eq:V17-involution-defect}; the identity for $P_\eta$
follows from
$P_\eta-P_\eta^2=(I-\varepsilon_\eta^2)/4$.  On a nonzero eigenvalue, the
regularized sign has the same sign and magnitude
$|\lambda|/\sqrt{\lambda^2+\eta^2}$.  The largest error occurs at the smallest
$|\lambda|$, proving \eqref{eq:V17-exact-sign-error}.  Finally,
\[
 1-\frac r{\sqrt{r^2+\eta^2}}
 =\frac{\eta^2}
 {\sqrt{r^2+\eta^2}(\sqrt{r^2+\eta^2}+r)}
 \le\frac{\eta^2}{2r^2}.
\]
\end{proof}

\begin{corollary}[Divisor-tube control gives an integrated sign rate]
\label{cor:V17-integrated-sign-rate}
Assume $\mu(\mathscr D_W)=0$.  For every $r>0$,
\begin{equation}
 \boxed{
 \int_X
 \|\operatorname{sgn}H_W-\varepsilon_\eta(H_W)\|_{\mathrm{op}}\,d\mu
 \le
 \mu\{\delta_W<r\}+\frac{\eta^2}{2r^2}.}
 \label{eq:V17-integrated-sign-split}
\end{equation}
If, for $0<r\le r_0$,
\begin{equation}
 \mu\{\delta_W<r\}
 \le C\left(\frac r{r_0}\right)^\alpha,
 \label{eq:V17-normalized-tube}
\end{equation}
then for $0<\eta\le r_0$, choosing
\begin{equation}
 r=r_0\left(\frac\eta{r_0}\right)^{2/(\alpha+2)}
 \label{eq:V17-optimal-simple-r}
\end{equation}
gives
\begin{equation}
 \boxed{
 \int_X
 \|\operatorname{sgn}H_W-\varepsilon_\eta(H_W)\|_{\mathrm{op}}\,d\mu
 \le
 \left(C+\frac12\right)
 \left(\frac\eta{r_0}\right)^{2\alpha/(\alpha+2)}.}
 \label{eq:V17-integrated-sign-rate}
\end{equation}
\end{corollary}

\begin{proof}
Split the integral into $\{\delta_W<r\}$ and its complement.  On the first
set the pointwise error is at most one; on the second use
\eqref{eq:V17-sign-error-r}.  This proves
\eqref{eq:V17-integrated-sign-split}.  Insert
\eqref{eq:V17-normalized-tube} and the displayed choice of $r$; both terms
then have power $2\alpha/(\alpha+2)$ in $\eta/r_0$.
\end{proof}

Fix chronological projections, a unitary reflection $R_\Theta$, a positive
Hodge--chiral factor $L$, a one-letter synthesis $W_1$, and $a>0$.  On the
invertible histories define
\begin{align}
 K_F(x)
 &=-\frac{\varepsilon_\Theta}{a}
 W_1^*LT_+\operatorname{sgn}(H_W(x))T_-R_\Theta W_1,
 \label{eq:V17-exact-crossing-writer}\\
 K_{F,\eta}(x)
 &=-\frac{\varepsilon_\Theta}{a}
 W_1^*LT_+\varepsilon_\eta(H_W(x))T_-R_\Theta W_1.
 \label{eq:V17-regularized-crossing-writer}
\end{align}
Put
\begin{equation}
 C_\times:=a^{-1}\|W_1\|_{\mathrm{op}}^2\|L\|_{\mathrm{op}}.
 \label{eq:V17-crossing-prefactor}
\end{equation}

\begin{theorem}[Soft-admissible mixed crossing certificate]
\label{thm:V17-soft-mixed-crossing}
Assume $\mu(\mathscr D_W)=0$ and that the physical reflection makes the
integrated matrices
\begin{equation}
 \overline K_F=\int K_F\,d\mu,
 \qquad
 \overline K_{F,\eta}=\int K_{F,\eta}\,d\mu
 \label{eq:V17-integrated-crossings}
\end{equation}
Hermitian.  Then
\begin{equation}
 \boxed{
 \|\overline K_F-\overline K_{F,\eta}\|_{\mathrm{op}}
 \le C_\times
 \int
 \|\operatorname{sgn}H_W-\varepsilon_\eta(H_W)\|_{\mathrm{op}}\,d\mu.}
 \label{eq:V17-mixed-crossing-radius}
\end{equation}
Consequently, if the supported minimum eigenvalue of
$\overline K_{F,\eta}$ exceeds the right side of
\eqref{eq:V17-mixed-crossing-radius}, then the mixed one-letter reflected form
is strictly positive on that support.

This conclusion is only an integrated or record-relative verdict.  It does
not imply $K_F(x)\succeq0$ for almost every history and therefore does not
license a historywise conditional exterior Gaussian.
\end{theorem}

\begin{proof}
The norm of the difference between the two crossing writers is bounded by
$C_\times$ times the sign-function error.  Integrate and use the triangle
inequality to obtain \eqref{eq:V17-mixed-crossing-radius}.  The eigenvalue
claim is Weyl's perturbation inequality.  The final limitation follows because
positivity of an average does not imply pointwise positivity; the record
separators already proved in Version~V11 give explicit finite examples.
\end{proof}

\begin{firewall}[Regularization is not an overlap replacement]
\label{fire:V17-regularization-not-overlap}
Equations \eqref{eq:V17-involution-defect}--\eqref{eq:V17-projector-defect}
show that $\varepsilon_\eta(H_W)$ is not a sign involution and $P_\eta(H_W)$
is not a chiral projection.  The regularized card may certify an integrated
limit after its complete error is charged.  It is not inserted into the
physical Dirac--Yukawa action, its word relations, determinant line, or
zero-mode count.
\end{firewall}

% =============================================================================
\section{Simple divisors force a rank-one spectral-flow jump}
\label{sec:V17-spectral-flow}
% =============================================================================

The divisor is not merely a region where the inverse estimate becomes large.
At a transverse zero it changes the overlap source sector discontinuously.

\begin{theorem}[Exact one-sided sign limits at a simple Wilson zero]
\label{thm:V17-simple-spectral-flow}
Let $H(t)=H(t)^*$ be $C^1$ for $|t|<t_0$.  Assume
\begin{equation}
 \Ker H(0)=\operatorname{span}\{u\},
 \qquad \|u\|=1,
 \label{eq:V17-simple-kernel}
\end{equation}
and that $H(0)$ restricted to $u^\perp$ is invertible.  Put
\begin{equation}
 \dot\lambda_0:=\langle u,\dot H(0)u\rangle
 \ne0,
 \qquad P_0=|u\rangle\langle u|,
 \label{eq:V17-transverse-speed}
\end{equation}
and let $S_{\rm red}$ be the sign of $H(0)|_{u^\perp}$, extended by zero on
$u$.  Then
\begin{equation}
 \boxed{
 \lim_{t\to0^\pm}\operatorname{sgn}H(t)
 =S_{\rm red}\pm\operatorname{sgn}(\dot\lambda_0)P_0.}
 \label{eq:V17-one-sided-sign-limits}
\end{equation}
Hence
\begin{equation}
 \boxed{
 \operatorname{sgn}H(0^+)-\operatorname{sgn}H(0^-)
 =2\operatorname{sgn}(\dot\lambda_0)P_0,}
 \label{eq:V17-sign-jump}
\end{equation}
\begin{equation}
 \boxed{
 \widehat P_-(0^+)-\widehat P_-(0^-)
 =\operatorname{sgn}(\dot\lambda_0)P_0.}
 \label{eq:V17-projector-jump}
\end{equation}
The positive spectral rank changes by
$\operatorname{sgn}(\dot\lambda_0)$.  The direct chronological crossing has
jump
\begin{equation}
 \boxed{
 2\operatorname{sgn}(\dot\lambda_0)T_+P_0T_-.}
 \label{eq:V17-crossing-jump}
\end{equation}
If the last matrix is nonzero, no norm-continuous crossing transport exists
through the divisor.  If it vanishes, the determinant/zero-mode line still
changes rank and remains a separate event.
\end{theorem}

\begin{proof}
Because zero is a simple isolated eigenvalue, the Riesz projection onto the
small spectral circle is $C^1$ for small $t$.  Finite-dimensional
self-adjoint perturbation theory therefore gives a $C^1$ eigenvalue
$\lambda(t)$ and normalized eigenvector $u(t)$ with
\[
 \lambda(0)=0,
 \qquad
 \lambda'(0)=\langle u,\dot H(0)u\rangle=\dot\lambda_0,
 \qquad
 |u(t)\rangle\langle u(t)|\to P_0.
\]
The remaining spectral cluster stays uniformly separated from zero and its
sign converges to $S_{\rm red}$.  For sufficiently small $t\ne0$,
$\operatorname{sgn}\lambda(t)=\operatorname{sgn}(t\dot\lambda_0)$.
Adding the simple eigenprojection to the stable reduced sign gives
\eqref{eq:V17-one-sided-sign-limits}; subtraction proves the remaining
formulas.  The rank statement follows from the single eigenvalue crossing.
\end{proof}

\begin{theorem}[Inertia change is an order-one cutoff defect]
\label{thm:V17-inertia-transport}
Let $S_0,S_1$ be self-adjoint involutions on the same $d$-dimensional Hilbert
space and put
\begin{equation}
 P_j=\frac12(I+S_j),
 \qquad r_j=\rank P_j.
 \label{eq:V17-inertia-projections}
\end{equation}
If $r_0\ne r_1$, then
\begin{equation}
 \boxed{\|S_1-S_0\|_{\mathrm{op}}=2,}
 \label{eq:V17-inertia-op-jump}
\end{equation}
\begin{equation}
 \boxed{
 \|S_1-S_0\|_{\mathrm{HS}}^2\ge4|r_1-r_0|,}
 \label{eq:V17-inertia-HS-jump}
\end{equation}
\begin{equation}
 \boxed{
 \|P_1-P_0\|_{\mathrm{HS}}^2\ge|r_1-r_0|.}
 \label{eq:V17-projection-HS-jump}
\end{equation}
Thus an adjacent-cutoff overlap-index change is a discrete spectral-flow or
line event.  It cannot be hidden inside a summable small sign-transport
remainder.
\end{theorem}

\begin{proof}
Assume $r_1>r_0$.  Dimension counting gives
\[
 \dim(\Ran P_1\cap\Ker P_0)
 \ge r_1+(d-r_0)-d=r_1-r_0.
\]
On that intersection, $S_1=I$ and $S_0=-I$, so $S_1-S_0=2I$.  This proves the
operator norm lower bound and contributes at least $4(r_1-r_0)$ to the
Hilbert--Schmidt norm squared.  The opposite rank order is symmetric, and the
operator-norm upper bound two follows from the triangle inequality for two
involutions.  Since $P_1-P_0=(S_1-S_0)/2$, the projection estimate follows.
\end{proof}

\begin{corollary}[Topological sector patching must precede cofinal sign transport]
\label{cor:V17-sector-patching}
On a cutoff family, exact or summably small transport of Wilson signs on one
fixed common carrier forces eventual constancy of the positive spectral rank.
Whenever the rank changes, the correct output is the corresponding kernel,
cokernel, divisor order, and orientation transport.  The card may resume the
Version~V16 sign crossing only after passing to a constant-rank sector patch.
\end{corollary}

\begin{proof}
A summable defect tends to zero.  By
\eqref{eq:V17-inertia-op-jump}, every rank-changing adjacent pair has defect
two, so only finitely many such changes can occur.  The zero-mode and line
packet is the zero-safe output already required by the inherited fermion
compiler.
\end{proof}

% =============================================================================
\section{Hard/soft crossing assembly and target-relative verdicts}
\label{sec:V17-hard-soft-crossing}
% =============================================================================

A low-gap tube can be small without being empty.  Its contribution must be
assembled before a mixed crossing is declared positive.

\begin{theorem}[Exact hard/soft lower bound for a mixed crossing]
\label{thm:V17-hard-soft-lower}
Let $P$ be a finite coefficient projection and let $K(x)=K(x)^*$ be a
measurable writer supported on $P$.  Let $A\subseteq X$ have mass $p$ and
suppose
\begin{equation}
 K(x)\succeq\kappa P
 \quad(x\in A),
 \qquad
 K(x)\succeq-LP
 \quad(x\notin A),
 \label{eq:V17-hard-soft-assumptions}
\end{equation}
for $\kappa,L\ge0$.  Then
\begin{equation}
 \boxed{
 \int_XK\,d\mu
 \succeq
 \bigl(p\kappa-(1-p)L\bigr)P.}
 \label{eq:V17-hard-soft-floor}
\end{equation}
In particular, the integrated form is strictly positive when
\begin{equation}
 p\kappa>(1-p)L.
 \label{eq:V17-hard-soft-strict}
\end{equation}
The statement applies inside every occurring boundary-record cell after
replacing $\mu$ by its conditional law.
\end{theorem}

\begin{proof}
Integrate the first lower bound on $A$ and the second on its complement.  The
conditional version is the same calculation with the record-cell law.
\end{proof}

\begin{assessment}[Two different promotion targets]
\label{ass:V17-two-targets}
For an integrated one-letter Osterwalder--Schrader target,
\Cref{thm:V17-hard-soft-lower,thm:V17-soft-mixed-crossing} provide legitimate
soft-gap compilers.  For a historywise conditional Gaussian or exterior
kernel, every occupied conditional card must itself pass, except on a declared
zero-mode saturation branch.  A positive mixed matrix is not promoted to the
stronger historywise semantics.
\end{assessment}

\begin{definition}[Wilson admissibility, divisor, and flow packet]
\label{def:V17-admissibility-packet}
At cutoff $n$, the target-minimal upstream packet is
\begin{equation}
 \boxed{
 \begin{aligned}
 \mathfrak A_{W,n}=\bigl(
 &\Lambda_n^+,\ \operatorname{supp}\Lambda_n^+,\ H_{W,n},\ \delta_{W,n},\\
 &\Lambda_n^+(\mathscr D_{W,n}),\
   \{\Lambda_n^+(0<\delta_{W,n}<r_j)\}_{j=1}^J,\\
 &\rank\one_{(0,\infty)}(H_{W,n})
   \text{ on every occupied invertible component},\\
 &\text{one divisor kernel/orientation card when occupied},\\
 &\text{the hard/soft crossing bounds and one adjacent packet}
 \bigr).
 \end{aligned}}
 \label{eq:V17-admissibility-packet}
\end{equation}
The radii $r_j$, record cells, and adjacent identifications are frozen before
the low-gap masses or crossing values are read.
\end{definition}

\begin{theorem}[Terminating admissibility alternative]
\label{thm:V17-admissibility-alternative}
After the common-cylinder and support checks pass, apply the following
\emph{ordered} alternative to one card and its declared adjacent copy.
\begin{enumerate}[label=\textup{(A17.\arabic*)},leftmargin=5.5em]
\item \textbf{Divisor branch.}  The divisor has positive physical mass at
      either card, or a selected history lies on it.  The sign projector and
      normalized covariance are forbidden there; export the oriented
      zero-mode exterior and reduced determinant/Pfaffian packet.
\item \textbf{Spectral-flow branch.}  Both endpoint divisors have zero mass,
      but the transported positive spectral ranks differ.  Return the
      order-one witness of \Cref{thm:V17-inertia-transport} and patch the line
      by constant-rank sectors.
\item \textbf{Hard-admissible stable branch.}  The transported inertia is
      constant and both actual supports have certified positive Wilson gaps.
      Run the direct Version~V16 Feshbach, first-flat, or polynomial crossing
      card.
\item \textbf{Soft-admissible stable branch.}  The transported inertia is
      constant, both divisors have zero mass, and at least one essential gap is
      zero.  A certified tube law or inverse moment may close an
      integrated/record-relative crossing through
      \Cref{thm:V17-soft-mixed-crossing}; historywise conditional Gaussianity
      remains unproved.
\item \textbf{Unresolved branch.}  The submitted interval does not decide
      divisor mass, inertia, a tube exponent, a hard gap, or the relevant
      crossing margin.  Tighten the same upstream card.
\end{enumerate}
A failed common history, support identity, nonnegative cylinder, or record
incidence is an input error rather than a sixth physical branch.
\end{theorem}

\begin{proof}
First test whether either physical card charges its divisor.  On the remaining
branch both signs exist almost everywhere.  Their transported positive ranks
are either equal or unequal; inequality gives the spectral-flow branch.  On
the equal-rank branch, \Cref{thm:V17-hard-occurrence} classifies each support
as hard exactly when its essential gap is positive.  Both hard cards give the
third branch; every other decided zero-divisor-mass case gives the fourth.
Every valid finite interval not separating these ordered conditions is
unresolved.  The domain failures do not define a physical Hermitian crossing
and therefore precede the mathematical alternative.
\end{proof}

% =============================================================================
\section{Exact controls: a measure-zero divisor, a rank jump, and a mixed pass}
\label{sec:V17-controls}
% =============================================================================

\subsection{An analytic divisor with zero mass but no Wilson window}

Let $X=[-1,1]$ with normalized Lebesgue measure and define
\begin{equation}
 H(t)=\begin{pmatrix}t&0\\0&2\end{pmatrix}.
 \label{eq:V17-interval-H}
\end{equation}

\begin{proposition}[Exact interval divisor control]
\label{prop:V17-interval-control}
For the family \eqref{eq:V17-interval-H},
\begin{equation}
 \mathscr D_W=\{0\},
 \qquad \mu(\mathscr D_W)=0,
 \qquad
 \operatorname*{ess\,inf}\delta_W=0.
 \label{eq:V17-interval-basic}
\end{equation}
For $0\le r\le1$,
\begin{equation}
 \boxed{\mu\{\delta_W<r\}=r.}
 \label{eq:V17-interval-tube}
\end{equation}
For $0<q<1$,
\begin{equation}
 \boxed{\mathbb E\delta_W^{-q}=\frac1{1-q},}
 \label{eq:V17-interval-moment}
\end{equation}
whereas the moment diverges for $q\ge1$.  The sign and overlap projector have
the exact jumps
\begin{equation}
 \operatorname{sgn}H(0^+)-\operatorname{sgn}H(0^-)
 =\begin{pmatrix}2&0\\0&0\end{pmatrix},
 \label{eq:V17-interval-sign-jump}
\end{equation}
\begin{equation}
 \widehat P_-(0^+)-\widehat P_-(0^-)
 =\begin{pmatrix}1&0\\0&0\end{pmatrix}.
 \label{eq:V17-interval-projector-jump}
\end{equation}
Thus almost-sure invertibility and every finite inverse moment below order one
coexist with failure of every hard Wilson window.
\end{proposition}

\begin{proof}
Since $|t|\le1<2$, one has $\delta_W(t)=|t|$.  The tube probability is the
normalized length of $(-r,r)$, namely $r$.  Symmetry gives
\[
 \mathbb E|t|^{-q}=\int_0^1t^{-q}\,dt,
\]
which is $1/(1-q)$ exactly for $q<1$ and diverges otherwise.  The two sign
matrices are $\operatorname{diag}(-1,1)$ and
$\operatorname{diag}(1,1)$.
\end{proof}

\begin{proposition}[The hard chart changes an ordinary observable]
\label{prop:V17-conditioning-control}
For $0<g<1$, the hard set $X_g=\{|t|\ge g\}$ has mass
\begin{equation}
 p_g=1-g.
 \label{eq:V17-interval-hard-mass}
\end{equation}
For the bounded writer $F(t)=t^2$,
\begin{equation}
 \boxed{
 \mathbb E_\mu F=\frac13,
 \qquad
 \mathbb E_{\mu_g}F=\frac{1+g+g^2}{3}.}
 \label{eq:V17-conditioned-second-moment}
\end{equation}
At $g=1/4$, the conditioned value is $7/16$ and differs from the original by
$5/48$.  The hard overlap chart is therefore a different cylinder even though
the deleted divisor tube can be arbitrarily small.
\end{proposition}

\begin{proof}
The mass is the normalized length of the two retained intervals.  Moreover,
\[
 \mathbb E_{\mu_g}t^2
 =\frac{\int_g^1t^2\,dt}{1-g}
 =\frac{1-g^3}{3(1-g)}
 =\frac{1+g+g^2}{3}.
\]
\end{proof}

\subsection{Exact regularization and hard/soft assembly controls}

\begin{proposition}[Rational regularized-sign defect]
\label{prop:V17-rational-regularization}
For the scalar Wilson eigenvalue $\lambda=3$ and regularization $\eta=4$,
\begin{equation}
 \boxed{
 \varepsilon_4(3)=\frac35,
 \qquad
 1-\varepsilon_4(3)^2=\frac{16}{25},
 \qquad
 P_4(3)-P_4(3)^2=\frac4{25},}
 \label{eq:V17-rational-defects}
\end{equation}
while the exact sign error is $2/5$.  The regularized object is therefore
zero-safe but not an overlap projection even in this one-dimensional exact
control.
\end{proposition}

\begin{proof}
The denominator is $\sqrt{3^2+4^2}=5$.  Substitute into
\eqref{eq:V17-regularized-sign}--\eqref{eq:V17-projector-defect}.
\end{proof}

\begin{proposition}[Exact mixed crossing floor from a small soft branch]
\label{prop:V17-mixed-floor-control}
Let a scalar crossing equal $1/2$ on a hard event of probability $9/10$ and
$-1$ on its complement.  Then
\begin{equation}
 \boxed{
 \mathbb EK
 =\frac9{10}\frac12-\frac1{10}
 =\frac7{20}>0.}
 \label{eq:V17-mixed-floor-control}
\end{equation}
This saturates the bound in \eqref{eq:V17-hard-soft-floor} with
$\kappa=1/2$ and $L=1$, while the soft history remains pointwise negative.
\end{proposition}

\begin{proof}
Direct evaluation gives the displayed number.  Equality holds because both
pointwise bounds are attained.
\end{proof}

\begin{proposition}[Exact inertia transport control]
\label{prop:V17-inertia-control}
For
\begin{equation}
 S_0=\operatorname{diag}(1,-1,-1),
 \qquad
 S_1=\operatorname{diag}(1,1,-1),
 \label{eq:V17-inertia-control-matrices}
\end{equation}
one has
\begin{equation}
 \boxed{
 r_1-r_0=1,
 \qquad
 \|S_1-S_0\|_{\mathrm{op}}=2,
 \qquad
 \|S_1-S_0\|_{\mathrm{HS}}^2=4,}
 \label{eq:V17-inertia-control-values}
\end{equation}
so every inequality in \Cref{thm:V17-inertia-transport} is sharp.
\end{proposition}

\begin{proof}
The difference is $\operatorname{diag}(0,2,0)$.
\end{proof}

% =============================================================================
\section{Version V17 theorem, ledger, and next physical acquisition}
\label{sec:V17-terminal}
% =============================================================================

\begin{theorem}[Version V17 Wilson-admissibility and divisor theorem]
\label{thm:V17-terminal}
Versions~V1--V16 are retained.  Version~V17 additionally proves:
\begin{enumerate}[label=\textup{(E17.\arabic*)},leftmargin=5.7em]
\item a uniform Wilson gap is exactly a minimum-gap statement on the actual
      physical measure support;
\item a hard admissibility selector changes the cylinder unless its omitted
      mass is zero;
\item a nontrivial analytic Wilson divisor has smooth measure zero, but any
      occupied neighbourhood destroys the hard gap;
\item divisor-tube probabilities and inverse-gap moments obey the exact
      layer-cake duality \eqref{eq:V17-layer-cake};
\item the zero-safe regularized sign has the exact involution and projection
      defects \eqref{eq:V17-involution-defect}--\eqref{eq:V17-projector-defect};
\item a power-law divisor tube gives the integrated sign rate
      \eqref{eq:V17-integrated-sign-rate};
\item the regularized crossing plus its complete tube radius can certify a
      mixed or record-relative one-letter form, but not historywise Gaussian
      positivity;
\item a transverse simple Wilson zero produces the rank-one sign and overlap
      projector jumps \eqref{eq:V17-sign-jump}--\eqref{eq:V17-projector-jump};
\item any adjacent inertia change has operator defect two and the sharp
      Hilbert--Schmidt lower bounds of
      \Cref{thm:V17-inertia-transport};
\item the hard/soft crossing decomposition has the exact lower floor
      \eqref{eq:V17-hard-soft-floor};
\item the finite admissibility packet terminates in hard, soft, divisor,
      spectral-flow, or unresolved branches;
\item the interval, rational regularization, mixed crossing, and inertia
      controls realize all central distinctions exactly.
\end{enumerate}
No item supplies the active RPESM gauge-history support, its Wilson gap
distribution, divisor mass, line orientation, or direct crossing value.
\end{theorem}

\begin{proof}
Items~\textup{(E17.1)--(E17.2)} are
\Cref{thm:V17-hard-occurrence,cor:V17-hard-conditioning}.  Item
\textup{(E17.3)} is \Cref{lem:V17-analytic-zero-set}; item
\textup{(E17.4)} is \Cref{thm:V17-tube-moment}; items
\textup{(E17.5)--(E17.7)} are
\Cref{thm:V17-regularized-sign,cor:V17-integrated-sign-rate,thm:V17-soft-mixed-crossing}.
Items~\textup{(E17.8)--(E17.9)} are
\Cref{thm:V17-simple-spectral-flow,thm:V17-inertia-transport}; item
\textup{(E17.10)} is \Cref{thm:V17-hard-soft-lower}; item
\textup{(E17.11)} is \Cref{thm:V17-admissibility-alternative}; and the exact
controls are
\Cref{prop:V17-interval-control,prop:V17-conditioning-control,prop:V17-rational-regularization,prop:V17-mixed-floor-control,prop:V17-inertia-control}.
\end{proof}

\begingroup\small
\begin{longtable}{>{\raggedright\arraybackslash}p{0.27\textwidth}
                  >{\raggedright\arraybackslash}p{0.20\textwidth}
                  >{\raggedright\arraybackslash}p{0.45\textwidth}}
\caption{Version V17 proof-status ledger.}\label{tab:V17-ledger}\\
\toprule
\textbf{Row}&\textbf{Status}&\textbf{Exact scope}\\
\midrule
\endfirsthead
\toprule
\textbf{Row}&\textbf{Status}&\textbf{Exact scope}\\
\midrule
\endhead
V1--V16 source, Duhamel, record, Grassmann, overlap, collar, and direct-sign
results
&Imported exactly
&Preserved verbatim; no prior theorem or control is recomputed.\\[0.3em]
Wilson gap on physical support
&\passstatus
&Essential infimum equals the minimum on every full-support positive
cylinder.\\[0.3em]
Hard admissibility conditioning
&\passstatus
&Exact omitted-mass writer and bounded-observable comparison; not a
normalization gauge.\\[0.3em]
Analytic divisor
&\passstatus
&Measure zero when nontrivial, but incompatible with a hard full-support gap
when nonempty.\\[0.3em]
Divisor tube and inverse moments
&\passstatus
&Exact layer-cake identity and two-way quantitative implications.\\[0.3em]
Regularized sign
&\passstatus
&Zero-safe contraction with exact involution/projector defects and a complete
integrated error.\\[0.3em]
Historywise overlap replacement
&\obsstatus
&Forbidden: the regularized effect is not a chiral projection and mixed
positivity is not conditional positivity.\\[0.3em]
Simple spectral flow
&\passstatus
&Rank-one sign/projector jump and explicit cross-boundary jump.\\[0.3em]
Adjacent inertia change
&\obsstatus
&Order-one sign defect; must be exported as a divisor/line sector event.\\[0.3em]
Hard/soft mixed crossing
&\passstatus
&Exact integrated lower floor with the soft negative contribution retained.\\[0.3em]
Selected RPESM admissibility packet
&\stopstatus
&The actual support, divisor mass/tube law, component inertia, and direct
crossing are not numerically populated.\\[0.3em]
Exterior, Duhamel, current/stress, charge, continuum, Einstein
&Not reopened
&All remain downstream of a passing crossing, line, and divisor packet.\\
\bottomrule
\end{longtable}
\endgroup

\begin{assessment}[Physical status after Version V17]
\label{ass:V17-status}
The protected Wilson window in the finite field-shell construction is a valid
conditional chart hypothesis.  The first unclosed physical arrow is now more
precise: one must show how the actual RPESM bosonic cylinder meets that chart
and the Wilson divisor.  If the actual support is hard-admissible, Version~V16
immediately supplies the finite crossing compiler.  If the divisor is null but
the gap is soft, the new tube card may close only the integrated or
record-relative target.  If the divisor is occupied or the inertia changes,
the zero-mode exterior and line packet precedes every normalized sign or
covariance.

The current status is
\begin{equation}
 \boxed{
 \begin{array}{cl}
 \passstatus:&\text{support-gap, tube, regularization, and flow compilers;}\\[0.25em]
 \obsstatus:&\begin{array}{l}\text{soft constraints or regularized signs used}\\[-0.1em]
                         \text{as hard overlap data;}\end{array}\\[0.25em]
 \stopstatus:&\begin{array}{l}\text{first populated RPESM admissibility/divisor}\\[-0.1em]
                          \text{occurrence card.}\end{array}
 \end{array}}
 \label{eq:V17-status}
\end{equation}
This moves the programme upstream of the repeated unpopulated sign matrix
without reopening the finite fermionic, Duhamel, or continuum theorem layers.
\end{assessment}

\begin{openproblem}[Version V18: populate the Wilson admissibility/divisor card]
\label{op:V17-next}
Preserve Versions~V1--V17.  On one predeclared active RPESM cutoff and one
adjacent cutoff:
\begin{enumerate}[label=\textup{(V18.\arabic*)},leftmargin=5.8em]
\item export the actual positive unnormalized bosonic cylinder, its support,
      the complete Wilson kernel family, and the physical boundary record;
\item acquire an outward interval for the divisor mass and a frozen finite
      tube panel $\mu(0<\delta_W<r_j)$, or certify a hard support gap directly;
\item record the positive spectral rank on every occupied invertible component
      and acquire the oriented zero-mode card on every occupied divisor
      component;
\item on a hard branch, run Version~V16 without modification; on a soft branch,
      use one predeclared regularization scale and charge the complete tube
      radius before testing only the intended mixed/record-relative form;
\item evaluate the adjacent inertia and sign-transport rows before dividing by
      any determinant or normalizing a covariance;
\item retain one held-out direct Wilson-sign or zero-mode saturation Read and
      the determinant/Pfaffian orientation packet independently.
\end{enumerate}
If neither a hard gap nor a divisor-tube/zero-mode packet is supplied, the
verdict remains \stopstatus{} at this upstream physical incidence.  Another
sign polynomial, fitted collar, exterior kernel, or Duhamel theorem is not an
advance.
\end{openproblem}

\section*{V17 exact replay and append-only record}
\addcontentsline{toc}{section}{V17 exact replay and append-only record}
The accompanying verifier uses Python standard-library exact rational
arithmetic.  It checks the interval divisor tube and inverse half-moment; the
change of the conditioned second moment; the exact regularized-sign and
projector defects at the $3$--$4$--$5$ control; the rank-one spectral-flow and
inertia bounds; and the hard/soft mixed crossing floor.  Ordinary and
optimized executions produce byte-identical JSON output.  The verifier does
not populate an RPESM gauge-history measure or Wilson matrix.

The complete Version~V16 source preceding its sole final active
\verb|\end{document}| is preserved byte for byte.  Its preterminal SHA--256
digest is recorded in the validation manifest.  A later continuation must
remove only the final active document-closing command, append new material
immediately before it, restore one terminal command, and increment the filename
to end in \texttt{\_V18.tex}.

% =============================================================================
% END APPENDED VERSION V17 MATERIAL
% =============================================================================


% =============================================================================
% START APPENDED VERSION V18 MATERIAL
% =============================================================================

\clearpage
\hypersetup{
 pdftitle={Renewal Einstein--Standard--Model Physical Source Metric, Duhamel Ancestry and Quantum Continuum Programme V18},
 pdfsubject={Overlap-index sector occurrence, connected-support inertia, regular Wilson divisors, and singular sector-wall sources}
}
\pdfinfo{
 /Title (Renewal Einstein--Standard--Model Physical Source Metric, Duhamel Ancestry and Quantum Continuum Programme V18)
 /Subject (Overlap-index sector occurrence, connected-support inertia, regular Wilson divisors, and singular sector-wall sources)
}

\section{Version V18: the overlap-index sector is an occurrence and source row}
\label{sec:V18-direction}

Versions~V1--V17 are retained above.  Version~V17 placed the first unresolved
fermionic arrow at the support of the actual Wilson family and separated hard,
soft, divisor, and spectral-flow branches.  The supplied finite field-shell
contains one additional premise which must now be made explicit: its chiral
block is defined on an index-zero sector.  A sector label is locally constant
off the Wilson divisor, but it is not manufactured by a positive determinant
modulus, a Wilson gap on a separately chosen chart, or a normalized bosonic
law.

The present continuation therefore asks two prior questions.
\begin{equation}
 \boxed{
 \begin{gathered}
 \text{Does the desired overlap-index sector occur in the physical cylinder?}\\
 \text{Does a declared deformation preserve that sector, or move its Wilson-divisor wall?}
 \end{gathered}}
 \label{eq:V18-two-upstream-questions}
\end{equation}
The second question is essential for the source-metric programme.  When a hard
sector wall moves through a positive configuration-space density, the
resulting measure tangent has a singular surface part and a Hellinger cusp.
It is then not represented by an ordinary square-integrable score.  A bulk
action derivative and its Fisher Gram cannot silently absorb that source.

\begin{firewall}[A sector hypothesis is not a sector occurrence]
\label{fire:V18-sector-hypothesis}
Writing ``on the index-zero sector'' specifies the domain of a formula.  It
does not prove that the original unnormalized cylinder is supported there,
that the sector has nonzero mass, that conditioning on it is physically
licensed, or that the sector remains fixed under the commands used to define
the source metric.  Those are separate occurrence, absolute-normalization,
and transport rows.
\end{firewall}

% =============================================================================
\section{Inertia sectors, gauge covariance, and the connected-support test}
\label{sec:V18-inertia-sectors}
% =============================================================================

Let $X$ be a topological space and let
\begin{equation}
 H:X\longrightarrow\operatorname{Herm}(\mathcal E)
 \label{eq:V18-H-family}
\end{equation}
be continuous on a finite-dimensional Hilbert space $\mathcal E$ of dimension
$d$.  Put
\begin{equation}
 \mathscr D=\{x:\det H(x)=0\},
 \qquad X^\times=X\setminus\mathscr D,
 \label{eq:V18-divisor-complement}
\end{equation}
and, on $X^\times$,
\begin{equation}
 \nu(x)=n_+(H(x)),
 \qquad
 P(x)=\one_{(0,\infty)}(H(x))
     =\frac12\bigl(I+\operatorname{sgn}H(x)\bigr).
 \label{eq:V18-inertia-projector}
\end{equation}
The integer $\nu$ is used below because it is independent of the sign
convention chosen for the usual overlap index; on a balanced chiral carrier
the two differ by a fixed affine transformation.

\begin{theorem}[Inertia-sector structure and gauge covariance]
\label{thm:V18-inertia-sectors}
For the family in \eqref{eq:V18-H-family}:
\begin{enumerate}[label=\textup{(\roman*)},leftmargin=3em]
\item each set
      \[
       X_k=\{x\in X^\times:\nu(x)=k\}
      \]
      is open in $X$, and $\nu$ is constant on every connected component of
      $X^\times$;
\item the spectral projection $P$ is norm-continuous on $X^\times$;
\item if a group acts on $X$ and
      \[
       H(gx)=U_g(x)H(x)U_g(x)^*
      \]
      for unitary $U_g(x)$, then
      \begin{equation}
       \boxed{
       \nu(gx)=\nu(x),
       \qquad
       P(gx)=U_g(x)P(x)U_g(x)^*.}
       \label{eq:V18-gauge-covariance}
      \end{equation}
\end{enumerate}
In particular, gauge transformations move a history inside one inertia
sector; they do not create a sector transition.
\end{theorem}

\begin{proof}
At every invertible $H(x_0)$, the smallest absolute eigenvalue is positive.
Norm continuity of $H$ and Weyl's eigenvalue inequality give a neighbourhood
on which no eigenvalue crosses zero.  The positive eigenvalue count is
therefore locally constant, proving the first assertion.  On the same
neighbourhood the Riesz formula
\[
 P(x)=\frac{1}{2\pi i}\int_\Gamma(z-H(x))^{-1}\,dz
\]
with a contour $\Gamma$ enclosing the positive spectrum proves norm
continuity.  Unitary conjugacy preserves the spectrum and functional calculus,
which gives \eqref{eq:V18-gauge-covariance}.
\end{proof}

\begin{theorem}[Two-anchor divisor and spectral-flow witness]
\label{thm:V18-two-anchor-flow}
Let $\gamma:[0,1]\to X$ be continuous, suppose $H(\gamma(0))$ and
$H(\gamma(1))$ are invertible, and put $H_t=H(\gamma(t))$.  If
\begin{equation}
 \nu(\gamma(0))\ne\nu(\gamma(1)),
 \label{eq:V18-different-inertia}
\end{equation}
then $H_t$ is singular for at least one $t\in(0,1)$.

Suppose in addition that $H_t$ is $C^1$, that its zeroes occur at finitely many
parameters $t_j$, and that each zero consists of one simple eigenvalue branch
$\lambda_j(t)$ with $\dot\lambda_j(t_j)\ne0$.  Then
\begin{equation}
 \boxed{
 \nu(\gamma(1))-\nu(\gamma(0))
 =\sum_j\operatorname{sgn}\dot\lambda_j(t_j).}
 \label{eq:V18-spectral-flow-count}
\end{equation}
Consequently at least
$|\nu(\gamma(1))-\nu(\gamma(0))|$ transverse Wilson zeroes occur along the
path.
\end{theorem}

\begin{proof}
If every $H_t$ were invertible, the integer-valued function
$t\mapsto\nu(\gamma(t))$ would be locally constant by
\Cref{thm:V18-inertia-sectors}, hence constant on the interval.  This proves
the first claim.  At a simple transverse zero, all noncrossing eigenvalues keep
their signs, while the branch $\lambda_j$ changes from negative to positive
when $\dot\lambda_j(t_j)>0$ and in the opposite direction when the derivative
is negative.  Thus the positive inertia changes by the displayed sign.
Summation over the finitely many crossings proves
\eqref{eq:V18-spectral-flow-count}.
\end{proof}

\begin{corollary}[Connected gauge--Higgs support forces one hard sector]
\label{cor:V18-connected-gauge-support}
Fix all discrete boundary, spin-structure, and regulator labels.  The finite
configuration factor
\begin{equation}
 \left(
   \frac{\mathrm{SU}(3)\times\mathrm{SU}(2)\times\mathrm{U}(1)}{\mathbb Z_6}
 \right)^{E}
 \times \mathscr K_H^{V}
 \label{eq:V18-connected-field-factor}
\end{equation}
is connected whenever $\mathscr K_H$ is a ball.  Hence, if the physical
cylinder has full support on this factor and $H_W$ is invertible everywhere,
its positive spectral rank is one constant integer.  If two histories in that
support have different inertia, a Wilson divisor is unavoidable.
\end{corollary}

\begin{proof}
Each of $\mathrm{SU}(3)$, $\mathrm{SU}(2)$, and $\mathrm{U}(1)$ is connected.  A quotient of a
connected space by the finite central subgroup $\mathbb Z_6$ is connected;
finite products preserve connectedness; and a ball is convex.  Apply
\Cref{thm:V18-inertia-sectors,thm:V18-two-anchor-flow}.
\end{proof}

\begin{remark}[Disconnected physical records]
\label{rem:V18-discrete-sectors}
A physical boundary condition, spin structure, orientation, or retained
terminal record may split the complete history space into disconnected
components.  The preceding corollary is then applied on each occurring
component.  The discrete component label is not erased: it is precisely the
record on which sectorwise line and crossing cards must be conditioned.
\end{remark}

% =============================================================================
\section{Sector conditioning and the impossibility of a hidden fixed-rank source}
\label{sec:V18-sector-conditioning}
% =============================================================================

Let $\mu$ be a probability on $X$ and suppose $\mu(\mathscr D)=0$.  Put
\begin{equation}
 p_k=\mu(X_k),
 \qquad
 \mu_k=p_k^{-1}\one_{X_k}\mu
 \quad(p_k>0).
 \label{eq:V18-sector-laws}
\end{equation}

\begin{theorem}[Exact sector decomposition and conditioning cost]
\label{thm:V18-sector-decomposition}
The physical law has the exact decomposition
\begin{equation}
 \boxed{
 \mu=\sum_{k=0}^{d}p_k\mu_k,}
 \label{eq:V18-law-sector-sum}
\end{equation}
and every bounded operator-valued writer $F$ obeys
\begin{equation}
 \boxed{
 \int F\,d\mu
 =\sum_{k:p_k>0}p_k\int F\,d\mu_k.}
 \label{eq:V18-writer-sector-sum}
\end{equation}
For the total-variation convention
$d_{\rm TV}(\alpha,\beta)=\sup_A|\alpha(A)-\beta(A)|$,
\begin{equation}
 \boxed{d_{\rm TV}(\mu_k,\mu)=1-p_k.}
 \label{eq:V18-sector-TV}
\end{equation}
Therefore a restriction to one overlap-index sector is a different physical
experiment whenever its mass is smaller than one.
\end{theorem}

\begin{proof}
The disjoint inertia strata partition $X^\times$ and the divisor is null, so
\eqref{eq:V18-law-sector-sum} and \eqref{eq:V18-writer-sector-sum} are the
ordinary disintegration over the finite-valued map $\nu$.  The density of
$\mu_k$ relative to $\mu$ is $p_k^{-1}\one_{X_k}$.  Hence
\[
 \frac12\int\left|p_k^{-1}\one_{X_k}-1\right|\,d\mu
 =\frac12\bigl[(1-p_k)+(1-p_k)\bigr]
 =1-p_k,
\]
which equals the displayed total-variation distance.
\end{proof}

\begin{theorem}[No unrecorded fixed-rank chiral source across inertia strata]
\label{thm:V18-no-fixed-rank-source}
Suppose $p_kp_\ell>0$ for two distinct integers $k\ne\ell$.  There is no
integer $r$ and no measurable family of isometries
\begin{equation}
 J_x:\mathbb C^r\longrightarrow\mathcal E
 \label{eq:V18-fixed-rank-trivialization}
\end{equation}
whose range equals $\operatorname{Ran}P(x)$ for $\mu$-almost every history in
both $X_k$ and $X_\ell$.

Consequently a historywise overlap, Grassmann, or reflected-hopping packet on
such a law must retain the inertia-sector label and a variable-rank measurable
field, or else enlarge to an ambient carrier while retaining the physical
support projector.  A single fixed-rank source coordinate system without that
record is ill typed.
\end{theorem}

\begin{proof}
An isometry from $\mathbb C^r$ has rank $r$.  If its range equals
$\operatorname{Ran}P(x)$ on $X_k$, then $r=k$; equality on $X_\ell$ gives
$r=\ell$, a contradiction.  The final statement is the same rank obstruction
written in the two permissible coordinate conventions.
\end{proof}

\begin{firewall}[Positive sector mixing is not a conditional crossing theorem]
\label{fire:V18-sector-mixture}
The integral of sectorwise crossing matrices may be positive even when one
occupied sector contains a negative crossing.  Such an integral is a mixed or
record-coarsened target.  It does not supply the historywise hopping
factorization required before conditional exterior second quantization.
\end{firewall}

% =============================================================================
\section{The norm-dominant full-support branch is homotopically fixed}
\label{sec:V18-dominant-branch}
% =============================================================================

The easiest way to prove a Wilson window on an unrestricted gauge cylinder is
a dominant diagonal term.  The following result records what that route does
and does not buy.

\begin{theorem}[Global hopping-dominance certificate and fixed spectral asymmetry]
\label{thm:V18-dominant-hard-gap}
Let $\Gamma=\Gamma^*=\Gamma^{-1}$ and suppose
\begin{equation}
 H_W(x)=\Gamma\bigl(cI+K(x)\bigr)
 \label{eq:V18-Wilson-diagonal-hopping}
\end{equation}
is Hermitian for every $x\in X$.  Assume
\begin{equation}
 \sup_{x\in X}\|K(x)\|_{\mathrm{op}}\le\kappa<|c|.
 \label{eq:V18-hopping-bound}
\end{equation}
Then
\begin{equation}
 \boxed{|H_W(x)|\succeq(|c|-\kappa)I
 \quad\text{for every }x\in X.}
 \label{eq:V18-global-dominant-gap}
\end{equation}
The homotopy
\begin{equation}
 H_{W,s}(x)=\Gamma\bigl(cI+sK(x)\bigr),
 \qquad 0\le s\le1,
 \label{eq:V18-dominant-homotopy}
\end{equation}
remains gapped, and therefore
\begin{equation}
 \boxed{
 \Tr\operatorname{sgn}H_W(x)
 =\operatorname{sgn}(c)\Tr\Gamma.}
 \label{eq:V18-fixed-asymmetry}
\end{equation}
In particular, when $\Tr\Gamma=0$, this certificate lands in the balanced
zero-spectral-asymmetry sector.
\end{theorem}

\begin{proof}
Left multiplication by the unitary $\Gamma$ preserves singular values, so
\[
 s_{\min}(H_{W,s}(x))
 =s_{\min}(cI+sK(x))
 \ge |c|-s\|K(x)\|
 \ge |c|-\kappa.
\]
This proves the gap and the gapped homotopy.  The trace of the sign is the
difference between positive and negative inertia and is integer-valued and
continuous along a gapped homotopy.  At $s=0$,
$\operatorname{sgn}(c\Gamma)=\operatorname{sgn}(c)\Gamma$, proving
\eqref{eq:V18-fixed-asymmetry}.
\end{proof}

\begin{corollary}[Executable local hopping bound]
\label{cor:V18-local-hopping-bound}
If
\begin{equation}
 K(x)=\sum_{j=1}^{N}\bigl(A_j(x)+B_j(x)\bigr),
 \qquad
 \|A_j(x)\|\le a_j,
 \quad
 \|B_j(x)\|\le b_j,
 \label{eq:V18-local-hopping-sum}
\end{equation}
then \eqref{eq:V18-hopping-bound} holds with
\begin{equation}
 \boxed{\kappa=\sum_{j=1}^{N}(a_j+b_j).}
 \label{eq:V18-kappa-sum}
\end{equation}
Thus one finite row of local link-operator norms gives a full-support hard-gap
certificate whenever $|c|>\kappa$.
\end{corollary}

\begin{proof}
Use the triangle inequality in operator norm and apply
\Cref{thm:V18-dominant-hard-gap}.
\end{proof}

\begin{assessment}[What the norm-dominant branch cannot decide]
\label{ass:V18-dominant-scope}
The certificate is useful for a massive index-zero control.  Because it is
gapped-homotopic to the constant diagonal operator, it cannot exhibit a
spectral-flow transition.  It also says nothing about the sign of the
reflection-adapted crossing: that sign still depends on the physical
reflection, Berezin orientation, Hodge factor, and cross-boundary block tested
in Versions~V13--V16.
\end{assessment}

% =============================================================================
\section{Regular Wilson divisors have a computable linear tube law}
\label{sec:V18-coarea}
% =============================================================================

Let $X$ now be a smooth $n$-dimensional Riemannian manifold and let
\begin{equation}
 d\mu=\rho\,d\operatorname{vol},
 \qquad \rho>0,
 \label{eq:V18-smooth-density}
\end{equation}
where $\rho$ is continuous.  Let $U\Subset X$ and suppose $H\in C^2(U)$.
Assume there is a simple real eigenvalue branch $\lambda\in C^2(U)$ with
spectral projection $E_\lambda$, such that every other eigenvalue has absolute
value at least $\gamma>0$ on $U$, and
\begin{equation}
 \mathscr D_0=\{x\in U:\lambda(x)=0\}
 \label{eq:V18-regular-divisor}
\end{equation}
is nonempty, compact, and regular:
\begin{equation}
 |\nabla\lambda|>0\quad\text{on }\mathscr D_0.
 \label{eq:V18-transverse-spatial-zero}
\end{equation}
After shrinking $U$, choose $r_0>0$ such that
$\{|\lambda|<r_0\}\Subset U$ and the smallest absolute eigenvalue there is
$|\lambda|$.

\begin{theorem}[Coarea density and exact regular-divisor tube law]
\label{thm:V18-coarea-tube}
For $|t|<r_0$, define
\begin{equation}
 h(t)=\int_{\{\lambda=t\}}
       \frac{\rho}{|\nabla\lambda|}\,d\sigma_t.
 \label{eq:V18-coarea-density}
\end{equation}
Then $h$ is continuous near zero, $h(0)>0$, and for every
$0<r<r_0$,
\begin{equation}
 \boxed{
 \mu\{x\in U:0<\delta_W(x)<r\}
 =\int_{-r}^{r}h(t)\,dt.}
 \label{eq:V18-exact-tube-coarea}
\end{equation}
Consequently,
\begin{equation}
 \boxed{
 \lim_{r\downarrow0}
 \frac{\mu\{0<\delta_W<r\}}{2r}
 =h(0)
 =\int_{\mathscr D_0}
   \frac{\rho}{|\nabla\lambda|}\,d\sigma>0.}
 \label{eq:V18-linear-tube-coefficient}
\end{equation}
A regular simple Wilson divisor therefore has generic tube exponent one; the
coefficient is a physical surface-density Read, not an abstract unspecified
constant.
\end{theorem}

\begin{proof}
The regular-value theorem gives a tubular neighbourhood of $\mathscr D_0$ in
which the level sets $\{\lambda=t\}$ vary smoothly.  Integration in the normal
coordinate, or equivalently the coarea formula, gives
\[
 \int_{\{|\lambda|<r\}}\rho\,d\operatorname{vol}
 =\int_{-r}^{r}
   \left(
     \int_{\{\lambda=t\}}
     \frac{\rho}{|\nabla\lambda|}\,d\sigma_t
   \right)dt.
\]
The divisor itself has zero volume, so removing it does not change the
integral.  Smooth tubular transport and continuity of $\rho$ and
$|\nabla\lambda|^{-1}$ imply continuity of $h$.  Positivity of $\rho$ and
nonemptiness of the compact divisor give $h(0)>0$.  Divide by $2r$ and use
continuity at zero.
\end{proof}

\begin{corollary}[Sharp inverse-gap threshold at a regular divisor]
\label{cor:V18-inverse-gap-threshold}
Under the hypotheses of \Cref{thm:V18-coarea-tube},
\begin{equation}
 \boxed{
 \int_{\{|\lambda|<r_0\}}
 \delta_W^{-q}\,d\mu
 \begin{cases}
 <\infty,&0<q<1,\\
 =\infty,&q\ge1.
 \end{cases}}
 \label{eq:V18-inverse-gap-threshold}
\end{equation}
Thus the first inverse gap moment already diverges on a regular codimension-one
Wilson wall with positive density.
\end{corollary}

\begin{proof}
By coarea the integral equals
\[
 \int_{-r_0}^{r_0}|t|^{-q}h(t)\,dt.
\]
Continuity and positivity of $h$ near zero make this integral finite exactly
when $q<1$.
\end{proof}

\begin{corollary}[Regularized sign and projector defect measure the divisor]
\label{cor:V18-regularization-calibrates-divisor}
Let
\begin{equation}
 \varepsilon_\eta(H)=H(H^2+\eta^2I)^{-1/2},
 \qquad
 P_\eta=\frac12(I+\varepsilon_\eta(H)).
 \label{eq:V18-regularized-sign-again}
\end{equation}
Define the transverse local readings
\begin{align}
 \mathscr E_\eta
 &=\int_{\{|\lambda|<r_0\}}
   \Tr\!
   \left[
    E_\lambda
    \left|\operatorname{sgn}H-\varepsilon_\eta(H)\right|
   \right]d\mu,
 \label{eq:V18-sign-error-read}\\
 \mathscr P_\eta
 &=\int_{\{|\lambda|<r_0\}}
   \Tr\!
   \left[E_\lambda(P_\eta-P_\eta^2)\right]d\mu.
 \label{eq:V18-projector-defect-read}
\end{align}
Then
\begin{equation}
 \boxed{
 \lim_{\eta\downarrow0}\frac{\mathscr E_\eta}{\eta}=2h(0),
 \qquad
 \lim_{\eta\downarrow0}\frac{\mathscr P_\eta}{\eta}
 =\frac{\pi}{4}h(0).}
 \label{eq:V18-regularization-asymptotics}
\end{equation}
If $0<h_-\le h(t)\le h_+$ for $|t|\le r_0$, the finite-scale enclosures are
\begin{equation}
 \boxed{
 2h_-\Phi(\eta,r_0)
 \le\mathscr E_\eta
 \le2h_+\Phi(\eta,r_0),}
 \label{eq:V18-sign-finite-bounds}
\end{equation}
where
\begin{equation}
 \Phi(\eta,r)=r+\eta-\sqrt{r^2+\eta^2},
 \label{eq:V18-Phi}
\end{equation}
and
\begin{equation}
 \boxed{
 \frac{h_-\eta}{2}\arctan\frac{r_0}{\eta}
 \le\mathscr P_\eta
 \le
 \frac{h_+\eta}{2}\arctan\frac{r_0}{\eta}.}
 \label{eq:V18-projector-finite-bounds}
\end{equation}
Thus a pair of zero-safe regularization readings can calibrate the regular
surface density without pretending that $P_\eta$ is an exact overlap
projection.
\end{corollary}

\begin{proof}
On the simple eigenline,
\[
 \Tr E_\lambda
 \left|\operatorname{sgn}H-\varepsilon_\eta(H)\right|
 =1-\frac{|\lambda|}{\sqrt{\lambda^2+\eta^2}},
\]
and
\[
 \Tr E_\lambda(P_\eta-P_\eta^2)
 =\frac{\eta^2}{4(\lambda^2+\eta^2)}.
\]
Apply \eqref{eq:V18-exact-tube-coarea}.  After the substitution $t=\eta u$,
the first quotient tends by dominated convergence to
\[
 h(0)\int_{-\infty}^{\infty}
 \left(1-\frac{|u|}{\sqrt{1+u^2}}\right)du
 =2h(0),
\]
while the second tends to
\[
 \frac{h(0)}4\int_{-\infty}^{\infty}\frac{du}{1+u^2}
 =\frac\pi4h(0).
\]
The finite bounds follow by replacing $h(t)$ by $h_-$ and $h_+$ and using
\[
 \int_0^r\left(1-\frac{t}{\sqrt{t^2+\eta^2}}\right)dt
 =r+\eta-\sqrt{r^2+\eta^2},
\]
\[
 \int_{-r}^{r}\frac{\eta^2}{4(t^2+\eta^2)}dt
 =\frac\eta2\arctan\frac r\eta.
\]
\end{proof}

\begin{firewall}[Nonregular divisors retain their own exponent]
\label{fire:V18-nonregular-divisor}
The exponent one and the limits in
\eqref{eq:V18-regularization-asymptotics} require one simple transverse
zero.  A multiple eigenvalue, a symmetry-enforced higher-codimension zero, a
critical point of the crossing eigenvalue, or an intersection of divisor
strata is not assigned this law.  It requires its own stratified tube or
zero-mode card and remains \stopstatus{} until that card is supplied.
\end{firewall}

% =============================================================================
\section{Moving a hard index sector creates a singular source}
\label{sec:V18-moving-sector}
% =============================================================================

The preceding coarea coefficient becomes a source term when the Wilson family
itself is deformed.  Let $s\mapsto\lambda_s\in C^2(U)$ and
$s\mapsto w_s\in C^1(U;(0,\infty))$ be smooth near $s=0$.  Put
\begin{equation}
 A_s=\{x\in U:\lambda_s(x)>0\},
 \qquad
 \Lambda_s=\one_{A_s}w_s\,d\operatorname{vol},
 \qquad
 Z_s=\Lambda_s(U),
 \label{eq:V18-moving-sector-law}
\end{equation}
and assume $D_0=\{\lambda_0=0\}$ is compact and regular.

\begin{theorem}[Exact sector-wall derivative]
\label{thm:V18-sector-wall-derivative}
For every $C^1$ scalar writer $F_s$,
\begin{equation}
 \boxed{
 \left.\frac d{ds}\right|_{0}
 \int_{A_s}F_sw_s\,d\operatorname{vol}
 =\int_{A_0}(\dot F_0w_0+F_0\dot w_0)\,d\operatorname{vol}
 +\int_{D_0}
   F_0w_0\frac{\dot\lambda_0}{{|\nabla\lambda_0|}}\,d\sigma.}
 \label{eq:V18-shape-derivative}
\end{equation}
If $w_s=e^{-S_s}a$ with fixed positive density factor $a$, the normalized
sector expectation obeys
\begin{equation}
 \boxed{
 \begin{aligned}
 \left.\frac d{ds}\right|_0\mathbb E_sF_s
 ={}&\mathbb E_0\dot F_0
   -\operatorname{Cov}_0(F_0,\dot S_0)\\
 &+\frac1{Z_0}\int_{D_0}
   (F_0-\mathbb E_0F_0)
   w_0\frac{\dot\lambda_0}{|\nabla\lambda_0|}\,d\sigma.
 \end{aligned}}
 \label{eq:V18-normalized-wall-response}
\end{equation}
The final integral is the configuration-space Wilson-sector wall source.  It
vanishes for all $F_0$ only when the signed wall functional itself vanishes.
\end{theorem}

\begin{proof}
The regular-value theorem provides normal coordinates $(y,t)$ with
$t=\lambda_0$ near $D_0$.  Differentiating the moving upper half-space in
these coordinates gives the Reynolds transport formula.  The normal velocity
of the boundary, measured in the outward normal of $A_0$, is
$\dot\lambda_0/|\nabla\lambda_0|$: differentiating
$\lambda_s(x_s)=0$ gives
$x'_0\cdot\nabla\lambda_0=-\dot\lambda_0$, while the outward normal of
$\{\lambda_0>0\}$ is $-\nabla\lambda_0/|\nabla\lambda_0|$.  This proves
\eqref{eq:V18-shape-derivative}.

Apply the formula first to $F_s$ and then to $F_s=1$, use
$\dot w_0=-\dot S_0w_0$, and differentiate the quotient by $Z_s$.  The bulk
terms give the covariance and the two boundary terms combine into the centred
integrand in \eqref{eq:V18-normalized-wall-response}.
\end{proof}

\begin{theorem}[Hellinger cusp and failure of an ordinary Fisher score]
\label{thm:V18-Hellinger-cusp}
Let
\begin{equation}
 d\mu_s=Z_s^{-1}\one_{A_s}w_s\,d\operatorname{vol},
 \label{eq:V18-normalized-moving-sector}
\end{equation}
and define the squared Hellinger distance by
\begin{equation}
 \mathsf H^2(\mu_s,\mu_0)
 =\int_U
   \left(\sqrt{\frac{d\mu_s}{d\operatorname{vol}}}
        -\sqrt{\frac{d\mu_0}{d\operatorname{vol}}}\right)^2
   d\operatorname{vol}.
 \label{eq:V18-Hellinger-definition}
\end{equation}
Then
\begin{equation}
 \boxed{
 \lim_{s\to0}
 \frac{\mathsf H^2(\mu_s,\mu_0)}{|s|}
 =\frac1{Z_0}\int_{D_0}
   w_0\frac{|\dot\lambda_0|}{|\nabla\lambda_0|}\,d\sigma
 =:\mathfrak c_{\rm wall}.}
 \label{eq:V18-Hellinger-cusp}
\end{equation}
If $\mathfrak c_{\rm wall}>0$, the family $(\mu_s)$ is not differentiable in
quadratic mean at $s=0$.  In particular, there is no
$\ell\in L^2(\mu_0)$ satisfying
\begin{equation}
 \sqrt{d\mu_s}
 =\sqrt{d\mu_0}
  \left(1+\frac{s}{2}\ell\right)
  +o_{L^2}(|s|).
 \label{eq:V18-QMD}
\end{equation}
Thus the hard moving-sector command has no finite ordinary Fisher score.
\end{theorem}

\begin{proof}
On the overlap $A_s\cap A_0$, smoothness of $w_s$ and $Z_s$ makes the square
of the difference of the two square-root densities $O(s^2)$ in $L^1$.
On the symmetric difference $A_s\triangle A_0$, exactly one density is zero,
so the Hellinger contribution is the corresponding normalized mass.  Tubular
coordinates and the boundary velocity computed in
\Cref{thm:V18-sector-wall-derivative} give
\[
 \lim_{s\to0}\frac1{|s|}
 \int_{A_s\triangle A_0}w_0\,d\operatorname{vol}
 =\int_{D_0}w_0
   \frac{|\dot\lambda_0|}{|\nabla\lambda_0|}\,d\sigma.
\]
Replacing $w_0/Z_0$ by $w_s/Z_s$ on the gained portion changes the result only
by $o(|s|)$.  This proves \eqref{eq:V18-Hellinger-cusp}.  If
\eqref{eq:V18-QMD} held, the squared Hellinger distance would be $O(s^2)$,
contradicting the positive linear limit.
\end{proof}

\begin{corollary}[The source metric must choose one of four honest branches]
\label{cor:V18-source-metric-branches}
For a hard index-sector family under a physical command, an ordinary finite
source-action/Fisher metric is licensed only after one of the following is
proved:
\begin{enumerate}[label=\textup{(S18.\arabic*)},leftmargin=5.4em]
\item the sector is fixed and $\dot\lambda_0=0$ on every active divisor face;
\item the physical density vanishes on the moving wall strongly enough to make
      $\mathfrak c_{\rm wall}=0$ and a separate quadratic-mean estimate passes;
\item a soft selector is part of the actual action, with its complete
      regularization debit retained;
\item the source is enlarged from an $L^2$ bulk score to the measure-valued
      bulk-plus-wall derivative in \eqref{eq:V18-shape-derivative}.
\end{enumerate}
A source Gram formed only from the bulk action score omits the wall source on
every branch with $\mathfrak c_{\rm wall}>0$.
\end{corollary}

\begin{proof}
The first branch removes the moving-support term.  The second is the only way
a moving hard selector can avoid the positive coefficient in
\eqref{eq:V18-Hellinger-cusp} without becoming stationary.  The third changes
the physical experiment to one with common support.  The fourth retains the
singular derivative rather than forcing it into $L^2(\mu_0)$.  The no-score
statement follows directly from \Cref{thm:V18-Hellinger-cusp}.
\end{proof}

\begin{assessment}[Relation to the earlier material-wall calculation]
\label{ass:V18-two-walls}
The wall in \Cref{thm:V18-sector-wall-derivative} lies in bosonic
configuration space and is defined by a Wilson eigenvalue.  It is distinct
from the spatial Dirichlet-sector wall studied in Version~V3.  Both produce
boundary terms, but they live on different carriers and cannot be identified
or paid by one source without a same-history incidence map.
\end{assessment}

% =============================================================================
\section{A terminating sector-occurrence and wall-source card}
\label{sec:V18-card}
% =============================================================================

\begin{definition}[Overlap-index sector occurrence card]
\label{def:V18-sector-card}
At one predeclared cutoff, the target-minimal upstream packet is
\begin{equation}
 \boxed{
 \begin{aligned}
 \mathfrak I_W=\bigl(&X,\Lambda^+,H_W,\mathscr D_W,
 \{p_k\}_{k=0}^{d},\text{ sector selector and decoder},\\
 &\text{two anchor inertias and one physical connecting path},\\
 &\text{hard gap or regular-divisor charts }(\lambda,E_\lambda,
   \nabla\lambda,\rho),\\
 &\text{command velocities }\dot\lambda_u,
   \text{ bulk action scores},\\
 &\text{zero-mode orientation, determinant/Pfaffian line, and one adjacent card}
 \bigr).
 \end{aligned}}
 \label{eq:V18-sector-card}
\end{equation}
The sector selector and every deformation direction are frozen before the
sector masses, tube coefficients, or crossing signs are read.
\end{definition}

\begin{theorem}[Ordered finite sector and wall-source alternative]
\label{thm:V18-sector-alternative}
A well-typed packet in \eqref{eq:V18-sector-card} has exactly one of the
following mathematical outcomes.
\begin{enumerate}[label=\textup{(V18.\arabic*)},leftmargin=5.5em]
\item \emph{Global hard one-sector branch.}  The actual support has a positive
      Wilson gap and one inertia.  If that inertia is the declared sector,
      Versions~V13--V16 apply on the original cylinder.
\item \emph{Physically selected hard sector.}  The desired sector has
      $0<p_k<1$, its selector occurs, and the conditioned cylinder and
      normalization are retained.  The hard crossing compiler applies only
      to this new sector-conditioned experiment.
\item \emph{Regular soft divisor branch.}  The divisor is null but has a
      regular simple face with $h(0)>0$.  The tube exponent, inverse-moment
      threshold, and regularization asymptotics are those of
      \Cref{thm:V18-coarea-tube,cor:V18-regularization-calibrates-divisor}.
      Historywise hard overlap is unavailable.
\item \emph{Moving-wall source branch.}  A physical command has
      $\mathfrak c_{\rm wall}>0$.  The ordinary Fisher/source metric is
      obstructed and the singular wall source must be frozen, softened, or
      retained measure-valuedly before any Duhamel ancestry calculation.
\item \emph{Divisor/stratified branch.}  The divisor has positive mass, a
      multiple zero, a nontransverse face, or an inertia-changing adjacent
      card.  Export the zero-mode exterior, line, and stratified tube packet.
\item \emph{Unresolved.}  The sector mass, support, gap, divisor regularity,
      or command wall velocity lacks a strict outward certificate.  Tighten
      this same card.
\end{enumerate}
An ill-typed selector, path, orientation, or common-cylinder relation is an
input error and is not a seventh physical outcome.
\end{theorem}

\begin{proof}
The support and hard-gap split is \Cref{thm:V17-hard-occurrence}.  Sector
conditioning and its exact cost are
\Cref{thm:V18-sector-decomposition}.  The regular divisor branch is
\Cref{thm:V18-coarea-tube,cor:V18-inverse-gap-threshold,cor:V18-regularization-calibrates-divisor}.
The moving-wall branch is
\Cref{thm:V18-sector-wall-derivative,thm:V18-Hellinger-cusp}.  Positive divisor
mass, nonregular strata, and inertia change are the zero-safe alternatives of
Version~V17 and \Cref{fire:V18-nonregular-divisor}.  The cases are ordered by
the first applicable physical row and exhaust every well-typed packet.
\end{proof}

\begin{firewall}[No return to the Duhamel card before sector differentiation]
\label{fire:V18-no-Duhamel-before-sector}
The V5--V8 source metric and Duhamel matrices presuppose a physically defined
linear command and its source action.  If that command moves a hard overlap
sector with $\mathfrak c_{\rm wall}>0$, the first derivative is not an
$L^2$ score.  A finite bulk Gram computed after deleting the divisor therefore
cannot be promoted to the physical source metric, regardless of how accurately
its later Duhamel residual is evaluated.
\end{firewall}

% =============================================================================
\section{Exact controls: sector cost, regular wall, and a trivial hard branch}
\label{sec:V18-controls}
% =============================================================================

\subsection{One regular index wall}

Let $X=[-1,1]$ with Lebesgue measure, and put
\begin{equation}
 H(t)=\operatorname{diag}(t,2).
 \label{eq:V18-interval-H}
\end{equation}
The normalized ambient law is uniform on $[-1,1]$.

\begin{proposition}[Exact sector and coarea control]
\label{prop:V18-interval-sector-control}
For \eqref{eq:V18-interval-H},
\begin{equation}
 \boxed{
 p_1=p_2=\frac12,
 \qquad
 d_{\rm TV}(\mu_1,\mu)
 =d_{\rm TV}(\mu_2,\mu)=\frac12.}
 \label{eq:V18-sector-control-masses}
\end{equation}
For $0<r\le1$,
\begin{equation}
 \boxed{\mu\{0<\delta_W<r\}=r,\qquad h(0)=\frac12.}
 \label{eq:V18-control-tube}
\end{equation}
Moreover,
\begin{equation}
 \boxed{
 \int\delta_W^{-q}\,d\mu
 =\frac1{1-q}\quad(0<q<1),}
 \label{eq:V18-control-inverse-moment}
\end{equation}
while the integral diverges for $q\ge1$.
Conditioning on the positive-rank-two sector changes the mean of the writer
$t$ from zero to $1/2$.
\end{proposition}

\begin{proof}
The two inertia sectors are $[-1,0)$ and $(0,1]$, each of ambient probability
$1/2$.  The divisor is one point and the crossing eigenvalue is $\lambda=t$;
the ambient density is $1/2$, so $h(0)=1/2$.  The tube has length $2r$ and
therefore probability $r$.  Finally,
\[
 \int\delta_W^{-q}\,d\mu
 =\int_0^1t^{-q}\,dt,
\]
which gives the displayed threshold.  The conditional mean on $(0,1]$ is
$1/2$.
\end{proof}

\begin{proposition}[Exact regularization readings]
\label{prop:V18-interval-regularization}
For the crossing eigenvalue in \eqref{eq:V18-interval-H},
\begin{equation}
 \boxed{
 \mathscr E_\eta
 =1+\eta-\sqrt{1+\eta^2},}
 \label{eq:V18-control-sign-error}
\end{equation}
\begin{equation}
 \boxed{
 \mathscr P_\eta
 =\frac\eta4\arctan\frac1\eta.}
 \label{eq:V18-control-projector-defect}
\end{equation}
Hence
\begin{equation}
 \boxed{
 \frac{\mathscr E_\eta}{\eta}\longrightarrow1=2h(0),
 \qquad
 \frac{\mathscr P_\eta}{\eta}\longrightarrow\frac\pi8
 =\frac\pi4h(0).}
 \label{eq:V18-control-regularization-limits}
\end{equation}
At the rational Pythagorean scale $\eta=3/4$,
\begin{equation}
 \boxed{\mathscr E_{3/4}=\frac12.}
 \label{eq:V18-control-rational-sign-error}
\end{equation}
\end{proposition}

\begin{proof}
Integrate the scalar formulas in
\Cref{cor:V18-regularization-calibrates-divisor} against the uniform density.
The first integral is
\[
 \int_0^1\left(1-\frac{t}{\sqrt{t^2+\eta^2}}\right)dt
 =1+\eta-\sqrt{1+\eta^2},
\]
and the second is
\[
 \frac14\int_0^1\frac{\eta^2}{t^2+\eta^2}\,dt
 =\frac\eta4\arctan\frac1\eta.
\]
The limits follow, and $\sqrt{1+(3/4)^2}=5/4$ gives the rational value.
\end{proof}

\subsection{A moving sector with an exact Hellinger cusp}

Put
\begin{equation}
 \lambda_s(t)=t-s,
 \qquad
 A_s=(s,1)\subset[-1,1],
 \qquad
 d\mu_s=(1-s)^{-1}\one_{(s,1)}(t)\,dt
 \quad(0\le s<1).
 \label{eq:V18-moving-interval-sector}
\end{equation}

\begin{proposition}[Exact shape response and failure of quadratic mean]
\label{prop:V18-moving-control}
For \eqref{eq:V18-moving-interval-sector},
\begin{equation}
 \boxed{Z_s=1-s,
 \qquad
 \mathbb E_s t=\frac{1+s}{2},}
 \label{eq:V18-moving-control-means}
\end{equation}
and the wall term in \eqref{eq:V18-normalized-wall-response} gives
\begin{equation}
 \boxed{
 \left.\frac d{ds}\right|_0\mathbb E_st=\frac12.}
 \label{eq:V18-moving-control-derivative}
\end{equation}
The squared Hellinger distance is exactly
\begin{equation}
 \boxed{
 \mathsf H^2(\mu_s,\mu_0)
 =2\bigl(1-\sqrt{1-s}\bigr),}
 \label{eq:V18-moving-control-Hellinger}
\end{equation}
so
\begin{equation}
 \boxed{
 \lim_{s\downarrow0}
 \frac{\mathsf H^2(\mu_s,\mu_0)}s=1.}
 \label{eq:V18-moving-control-cusp}
\end{equation}
No square-integrable score exists for this hard-sector displacement.
\end{proposition}

\begin{proof}
The first two identities are elementary integration.  At $s=0$, the boundary
point is $t=0$, $\dot\lambda=-1$, $|\nabla\lambda|=1$, and
$F(0)-\mathbb E_0F=-1/2$.  The product in the wall term is therefore $1/2$.
For the Hellinger distance, the lost interval $(0,s)$ contributes $s$, while
on $(s,1)$ the square-root density difference contributes
$2-s-2\sqrt{1-s}$.  Their sum is
\eqref{eq:V18-moving-control-Hellinger}.  The limit is one, so
\Cref{thm:V18-Hellinger-cusp} rules out quadratic-mean differentiability.
\end{proof}

\subsection{A globally hard but homotopically fixed control}

Let
\begin{equation}
 \Gamma=\operatorname{diag}(1,-1),
 \qquad
 H_s(t)=3\Gamma+stI,
 \qquad
 |t|\le1,
 \quad 0\le s\le1.
 \label{eq:V18-dominant-control}
\end{equation}

\begin{proposition}[Exact norm-dominant index-zero control]
\label{prop:V18-dominant-control}
Every matrix in \eqref{eq:V18-dominant-control} has gap at least two and
positive inertia one.  Moreover,
\begin{equation}
 \boxed{
 \operatorname{sgn}H_s(t)=\Gamma,
 \qquad
 \Tr\operatorname{sgn}H_s(t)=0.}
 \label{eq:V18-dominant-control-sign}
\end{equation}
This is a full-support hard branch with no spectral flow.
\end{proposition}

\begin{proof}
The eigenvalues are $3+st$ and $-3+st$.  For $|t|\le1$ and $0\le s\le1$,
the first is at least two and the second is at most minus two.  The remaining
claims follow immediately.
\end{proof}

% =============================================================================
\section{Version V18 theorem, ledger, and next physical acquisition}
\label{sec:V18-terminal}
% =============================================================================

\begin{theorem}[Version V18 sector-occurrence and wall-source theorem]
\label{thm:V18-terminal}
Versions~V1--V17 are retained.  Version~V18 additionally proves:
\begin{enumerate}[label=\textup{(E18.\arabic*)},leftmargin=5.7em]
\item Wilson inertia is locally constant off the divisor and gauge invariant;
\item two different anchor inertias on one path force a divisor, and transverse
      crossings obey the exact spectral-flow count
      \eqref{eq:V18-spectral-flow-count};
\item the fixed-sector gauge--Higgs configuration factor is connected, so a
      full-support hard Wilson family has one inertia there;
\item a null-divisor law decomposes exactly into sector-conditioned cylinders,
      with total-variation cost $1-p_k$ for selecting sector $k$;
\item two occupied inertia ranks cannot be represented by one unrecorded
      fixed-rank chiral source;
\item a global norm-dominant Wilson window is gapped-homotopic to its diagonal
      reference and has fixed spectral asymmetry;
\item a regular simple Wilson divisor has the exact coarea tube law
      \eqref{eq:V18-exact-tube-coarea} and positive linear coefficient
      \eqref{eq:V18-linear-tube-coefficient};
\item the local inverse-gap threshold is exactly $q=1$;
\item zero-safe sign and projector regularizations recover the divisor surface
      density through \eqref{eq:V18-regularization-asymptotics};
\item motion of a hard sector has the exact bulk-plus-wall derivative
      \eqref{eq:V18-normalized-wall-response};
\item nonzero wall velocity produces the Hellinger cusp
      \eqref{eq:V18-Hellinger-cusp} and obstructs every ordinary finite Fisher
      score;
\item the ordered sector-occurrence card terminates in global hard, selected
      hard, regular soft, moving-wall, divisor/stratified, or unresolved
      branches;
\item the interval, regularization, moving-sector, and norm-dominant controls
      realize the central alternatives exactly.
\end{enumerate}
No item populates the actual RPESM index-sector mass, Wilson anchor inertias,
regular-divisor geometry, command wall velocity, or reflected crossing.
\end{theorem}

\begin{proof}
Items~\textup{(E18.1)--(E18.3)} are
\Cref{thm:V18-inertia-sectors,thm:V18-two-anchor-flow,cor:V18-connected-gauge-support}.
Items~\textup{(E18.4)--(E18.5)} are
\Cref{thm:V18-sector-decomposition,thm:V18-no-fixed-rank-source}.  Item
\textup{(E18.6)} is \Cref{thm:V18-dominant-hard-gap}.  Items
\textup{(E18.7)--(E18.9)} are
\Cref{thm:V18-coarea-tube,cor:V18-inverse-gap-threshold,cor:V18-regularization-calibrates-divisor}.
Items~\textup{(E18.10)--(E18.11)} are
\Cref{thm:V18-sector-wall-derivative,thm:V18-Hellinger-cusp}; item
\textup{(E18.12)} is \Cref{thm:V18-sector-alternative}; and the controls are
\Cref{prop:V18-interval-sector-control,prop:V18-interval-regularization,prop:V18-moving-control,prop:V18-dominant-control}.
\end{proof}

\begingroup\small
\begin{longtable}{>{\raggedright\arraybackslash}p{0.27\textwidth}
                  >{\raggedright\arraybackslash}p{0.20\textwidth}
                  >{\raggedright\arraybackslash}p{0.45\textwidth}}
\caption{Version V18 proof-status ledger.}\label{tab:V18-ledger}\\
\toprule
\textbf{Row}&\textbf{Status}&\textbf{Exact scope}\\
\midrule
\endfirsthead
\toprule
\textbf{Row}&\textbf{Status}&\textbf{Exact scope}\\
\midrule
\endhead
V1--V17 source, record, Grassmann, overlap, crossing, sign, and divisor results
&Imported exactly
&Preserved verbatim; no earlier theorem or control is recomputed.\\[0.3em]
Inertia-sector and gauge covariance
&\passstatus
&Locally constant off the divisor; unitary gauge orbits remain in one sector.\\[0.3em]
Two-anchor path test
&\passstatus
&Different endpoint inertias force a divisor; simple crossings carry the exact
spectral-flow count.\\[0.3em]
Index-sector conditioning
&\passstatus
&Exact mixture and total-variation cost; a proper sector is a new cylinder.\\[0.3em]
Fixed-rank source across sectors
&\obsstatus
&Impossible without a sector record or ambient supported formulation.\\[0.3em]
Norm-dominant hard branch
&\passstatus
&Full-support gap and fixed spectral asymmetry, but no crossing-sign theorem.\\[0.3em]
Regular divisor tube
&\passstatus
&Exact coarea density, linear tube law, and sharp inverse-moment threshold.\\[0.3em]
Regularization calibration
&\passstatus
&Sign error and projection defect recover the regular divisor surface density.\\[0.3em]
Moving hard sector
&\obsstatus{} for an ordinary score
&Exact singular wall source and linear Hellinger cusp when wall velocity is
nonzero.\\[0.3em]
Bulk-only source metric
&\obsstatus
&Cannot represent a moving hard-sector command with positive wall stock.\\[0.3em]
Selected RPESM sector packet
&\stopstatus
&Sector mass, selector occurrence, anchor inertias, wall geometry, and command
velocity are not numerically populated.\\[0.3em]
Exterior, Duhamel, current/stress, charge, continuum, Einstein
&Not reopened
&All remain downstream of a well-typed crossing, line, divisor, and source
metric packet.\\
\bottomrule
\end{longtable}
\endgroup

\begin{assessment}[Physical status after Version V18]
\label{ass:V18-status}
The finite field-shell's index-zero clause is now recognized as an upstream
sector-occurrence statement.  On an unrestricted connected gauge--Higgs
support, a hard Wilson gap forces one global inertia; a different occupied
inertia necessarily introduces a divisor.  If the desired sector is selected
by conditioning, its mass and normalization belong to the physical cylinder.
If a source command moves its divisor wall, the derivative contains a singular
surface source and does not possess an ordinary finite Fisher score.

The current status is
\begin{equation}
 \boxed{
 \begin{array}{cl}
 \passstatus:&\begin{array}{l}
   \text{sector, connected-support, coarea, regularization, and}\\[-0.1em]
   \text{shape-derivative compilers;}
  \end{array}\\[0.35em]
 \obsstatus:&\begin{array}{l}
   \text{unrecorded sector restriction or a bulk-only source metric}\\[-0.1em]
   \text{for a moving hard Wilson wall;}
  \end{array}\\[0.35em]
 \stopstatus:&\begin{array}{l}
   \text{first populated RPESM index-sector occurrence and}\\[-0.1em]
   \text{sector-wall source card.}
  \end{array}
 \end{array}}
 \label{eq:V18-status}
\end{equation}
This is earlier than the direct Wilson-sign crossing.  Returning immediately
to another sign approximation would skip the sector and source domain of that
matrix.
\end{assessment}

\begin{openproblem}[Version V19: populate one index-sector and wall-source card]
\label{op:V18-next}
Preserve Versions~V1--V18.  On one active RPESM cutoff and one adjacent
cutoff, proceed in this order.
\begin{enumerate}[label=\textup{(V19.\arabic*)},leftmargin=5.8em]
\item identify the actual bosonic support and determine whether the
      ``index-zero sector'' is the whole support, a physically occurring
      selector, or only a formula-domain hypothesis;
\item evaluate the sector masses and two predeclared anchor inertias on every
      connected occurring component; retain a physical connecting path when
      the inertias differ;
\item on a hard one-sector branch, give the support gap and run Version~V16;
      on a selected hard branch, include the selector normalization and one
      held-out sector-conditioned writer;
\item on every null regular divisor face, acquire the crossing eigenvalue,
      spectral projection, density, gradient floor, and coarea coefficient;
      use two predeclared regularization scales as a held-out consistency test;
\item for each source command entering the physical source-action metric,
      acquire $\dot\lambda_u$ on the divisor and decide whether the wall stock
      $\mathfrak c_{\rm wall}(u)$ vanishes;
\item if a wall stock is positive, choose explicitly between a fixed-sector
      command, an occurring soft selector, or a measure-valued wall source;
      do not form a bulk Fisher/Duhamel card first;
\item retain the zero-mode exterior, determinant/Pfaffian orientation, direct
      crossing, and adjacent inertia rows independently.
\end{enumerate}
If the sector selector and its command incidence are absent, the verdict
remains \stopstatus{} at this upstream source domain.  No additional collar,
polynomial sign, exterior, Duhamel, or continuum theorem is licensed.
\end{openproblem}

\section*{V18 exact replay and append-only record}
\addcontentsline{toc}{section}{V18 exact replay and append-only record}
The accompanying verifier uses Python standard-library rational and
high-precision decimal arithmetic.  It checks the two inertia sectors of the
interval Wilson family, their exact masses and conditioning cost, the linear
divisor tube and inverse-moment threshold, the rational regularized-sign read,
the exact moving-sector expectation and Hellinger cusp, the two-anchor inertia
change, and the globally hard homotopically fixed control.  Ordinary and
optimized executions produce byte-identical JSON output.  The verifier does
not populate an RPESM gauge-history Wilson matrix or sector selector.

The complete Version~V17 source preceding its sole final active
\verb|\end{document}| is preserved byte for byte.  Its preterminal SHA--256
digest is recorded in the validation manifest.  A later continuation must
remove only the final active document-closing command, append new material
immediately before it, restore one terminal command, and increment the filename
to end in \texttt{\_V19.tex}.

% =============================================================================
% END APPENDED VERSION V18 MATERIAL
% =============================================================================

% =============================================================================
% START APPENDED VERSION V19 MATERIAL
% =============================================================================

\clearpage
\hypersetup{
 pdftitle={Renewal Einstein--Standard--Model Physical Source Metric, Duhamel Ancestry and Quantum Continuum Programme V19},
 pdfsubject={Control-to-active line reweighting, weighted Wilson-divisor tubes, and the Fisher threshold for moving overlap sectors}
}
\pdfinfo{
 /Title (Renewal Einstein--Standard--Model Physical Source Metric, Duhamel Ancestry and Quantum Continuum Programme V19)
 /Subject (Control-to-active line reweighting, weighted Wilson-divisor tubes, and the Fisher threshold for moving overlap sectors)
}

\section{Version V19: the control cylinder and the active line-weighted cylinder are different experiments}
\label{sec:V19-direction}

Versions~V1--V18 are retained above.  Version~V18 showed that a moving hard
Wilson-inertia sector has a configuration-space surface source and, when the
physical density is nonzero at the moving wall, a linear Hellinger cusp rather
than an ordinary square-integrable score.  The supplied projective
Einstein--Standard--Model regulator contains one still earlier distinction.  Its
positive population law is built from separately chosen nonzero massive
control sections, while the active chiral determinant, real Pfaffian, their
divisors, and their orientations are retained as independent line data.
Population of all finite writers under the control law does not decide which
positive modulus, if any, defines the source metric used by the active quantum
claim.

The present continuation therefore separates three objects:
\begin{equation}
 \boxed{
 \begin{gathered}
 \text{the nonvanishing control cylinder used to populate the finite packet},\\
 \text{an optional active phase-quenched positive line modulus},\\
 \text{the complex or real line-valued amplitude and its zero-mode saturation.}
 \end{gathered}}
 \label{eq:V19-three-line-objects}
\end{equation}
They may be related by an exact Radon--Nikodym writer, but they are not
interchangeable.  This distinction resolves the first fork left by V18.  A
line modulus which vanishes at the Wilson wall can remove the surface term and,
if it vanishes sufficiently rapidly, restore quadratic-mean differentiability.
A nonzero control section does neither.  Conversely, a vanishing line modulus
does not create a hard Wilson gap: every punctured low-gap tube still has
positive mass whenever the two laws are equivalent off the divisor.

\begin{firewall}[A population measure is not automatically the active quantum measure]
\label{fire:V19-control-not-active}
A positive control law may be used to carry and evaluate line-valued writers,
zero-safe exterior coefficients, and same-history crossing data.  It becomes
the positive modulus of the active fermion experiment only after an occurring
line-incidence or Radon--Nikodym row identifies those two uses.  Replacing the
control section by the active determinant or Pfaffian modulus changes the
normalized law, sector masses, source metric, and cutoff transport unless the
resulting density ratio is constant.
\end{firewall}

% =============================================================================
\section{Exact control-to-active Radon--Nikodym calculus}
\label{sec:V19-RN}
% =============================================================================

Let $(X,\mathscr B,\lambda)$ be one finite-cutoff bosonic history space.  Let
$S_B$ be the assembled bosonic action.  On one common line trivialization let
$\sigma^+$ and $p^+$ be the nonvanishing complex and real control sections, and
let $\sigma$ and $p$ be the active complex and real sections.  Define
\begin{align}
 d\Lambda^+
 &=e^{-S_B}\|\sigma^+\|^2|p^+|^2\,d\lambda,
 \label{eq:V19-control-cylinder}\\
 W_{\rm line}
 &=\frac{\|\sigma\|^2|p|^2}
         {\|\sigma^+\|^2|p^+|^2},
 \label{eq:V19-line-ratio}\\
 d\Lambda^{\rm abs}
 &=W_{\rm line}\,d\Lambda^+.
 \label{eq:V19-active-absolute-cylinder}
\end{align}
The ratio is intrinsic under simultaneous nonzero changes of local line frames.
It records only positive modulus.  Complex determinant phase, real Pfaffian
orientation, and zero-mode saturation remain separate rows.

Assume
\begin{equation}
 0<Z_+:=\Lambda^+(X)<\infty,
 \qquad
 0<Z_{\rm abs}:=\Lambda^{\rm abs}(X)<\infty,
 \label{eq:V19-two-normalizers}
\end{equation}
and put $\mu^+=Z_+^{-1}\Lambda^+$ and
$\mu^{\rm abs}=Z_{\rm abs}^{-1}\Lambda^{\rm abs}$.

\begin{theorem}[Exact line reweighting, sector masses, and metric distance]
\label{thm:V19-RN-exact}
Under \eqref{eq:V19-two-normalizers},
\begin{equation}
 \boxed{
 \frac{d\mu^{\rm abs}}{d\mu^+}
 =\frac{W_{\rm line}}{\mathbb E_{\mu^+}W_{\rm line}}.}
 \label{eq:V19-normalized-RN}
\end{equation}
Consequently, for every integrable writer $F$ and every measurable sector
$A$,
\begin{align}
 \boxed{
 \mathbb E_{\mu^{\rm abs}}F
 =\frac{\mathbb E_{\mu^+}(W_{\rm line}F)}
        {\mathbb E_{\mu^+}W_{\rm line}},}
 \label{eq:V19-reweighted-expectation}\\
 \boxed{
 \mu^{\rm abs}(A)
 =\frac{\mathbb E_{\mu^+}
        (W_{\rm line}\one_A)}
        {\mathbb E_{\mu^+}W_{\rm line}}.}
 \label{eq:V19-reweighted-sector-mass}
\end{align}
With total variation normalized by
$d_{\rm TV}(\mu,\nu)=\frac12\int|d\mu-d\nu|$, and squared Hellinger distance
$\mathsf H^2(\mu,\nu)=\int(\sqrt{d\mu}-\sqrt{d\nu})^2$, one has
\begin{align}
 \boxed{
 d_{\rm TV}(\mu^{\rm abs},\mu^+)
 =\frac12\mathbb E_{\mu^+}
 \left|
 \frac{W_{\rm line}}{\mathbb E_{\mu^+}W_{\rm line}}-1
 \right|,}
 \label{eq:V19-TV-exact}\\
 \boxed{
 \mathsf H^2(\mu^{\rm abs},\mu^+)
 =2\left(
 1-\mathbb E_{\mu^+}
 \sqrt{\frac{W_{\rm line}}
 {\mathbb E_{\mu^+}W_{\rm line}}}
 \right).}
 \label{eq:V19-Hellinger-exact}
\end{align}
These quantities vanish simultaneously exactly when $W_{\rm line}$ is
constant $\mu^+$-almost everywhere.
\end{theorem}

\begin{proof}
Equation \eqref{eq:V19-active-absolute-cylinder} and normalization give
\eqref{eq:V19-normalized-RN}.  Integrating $F$ and $\one_A$ gives
\eqref{eq:V19-reweighted-expectation}--\eqref{eq:V19-reweighted-sector-mass}.
The total-variation formula is the $L^1(\mu^+)$ norm of the displayed density.
For probability measures, expansion of the square roots gives
$\mathsf H^2=2-2\int\sqrt{d\mu^{\rm abs}d\mu^+}$, proving
\eqref{eq:V19-Hellinger-exact}.  Equality in either distance forces the
Radon--Nikodym derivative to be one almost everywhere, and the converse is
immediate.
\end{proof}

\begin{proposition}[Support and essential-gap invariance under equivalent line reweighting]
\label{prop:V19-support-equivalence}
Suppose $X$ is a compact metric space, $\mu^+$ has full support, and
$W_{\rm line}$ is continuous.  Then
\begin{equation}
 \boxed{
 \operatorname{supp}\mu^{\rm abs}
 =\overline{\{x:W_{\rm line}(x)>0\}}.}
 \label{eq:V19-active-support}
\end{equation}
If $W_{\rm line}>0$ $\mu^+$-almost everywhere, the two measures are equivalent.
For every measurable nonnegative gap writer $\delta$,
\begin{equation}
 \boxed{
 \operatorname*{ess\,inf}_{\mu^{\rm abs}}\delta
 =\operatorname*{ess\,inf}_{\mu^+}\delta.}
 \label{eq:V19-ess-gap-invariance}
\end{equation}
In particular, a line weight which vanishes only on a $\mu^+$-null divisor can
improve tube and inverse-moment exponents, but it cannot manufacture a hard
positive gap.
\end{proposition}

\begin{proof}
A point with $W_{\rm line}>0$ has a neighbourhood on which the continuous
weight is bounded below, so every such neighbourhood has positive active mass.
A point outside the closure has a neighbourhood on which the active density is
zero.  This proves \eqref{eq:V19-active-support}.  If the weight is positive
almost everywhere, a set is null for one normalized measure exactly when it is
null for the other.  Essential infima depend only on null sets, proving
\eqref{eq:V19-ess-gap-invariance}.
\end{proof}

\begin{countertheorem}[Control population does not determine active sector stock]
\label{cth:V19-control-no-active-stock}
There are equivalent control and active laws with identical support and
identical essential Wilson gap, but different inertia-sector masses, source
means, total variation, and Hellinger distance.  Thus the active sector card
cannot be populated from control-law occurrence alone.
\end{countertheorem}

\begin{proof}
Take $X=[0,1]$, $d\mu^+=dx$, and $W_{\rm line}=x^2$.  Then
$d\mu^{\rm abs}=3x^2dx$.  The laws are equivalent and both have essential
infimum zero for the gap $\delta(x)=x$.  For $A=[0,1/2]$ their masses are
$1/2$ and $1/8$.  The exact distances are calculated in
\Cref{sec:V19-controls} below and are strictly positive.
\end{proof}

\begin{firewall}[Positive modulus does not orient the line]
\label{fire:V19-modulus-no-phase}
Replacing $\sigma$ by $e^{i\theta}\sigma$, or $p$ by $-p$, leaves
$W_{\rm line}$, $\mu^{\rm abs}$, every positive source Gram, and every sector
mass unchanged.  It changes the complex phase or real orientation row.  No
positive reweighting theorem below supplies determinant-line holonomy or a
Pfaffian sign.
\end{firewall}

% =============================================================================
\section{Weighted Wilson-divisor tubes and the shifted inverse-gap threshold}
\label{sec:V19-weighted-tubes}
% =============================================================================

Let $M$ be a smooth finite-dimensional bosonic history manifold.  Let
$\lambda:M\to\mathbb R$ be a smooth simple Wilson crossing eigenvalue and
\begin{equation}
 D=\{\lambda=0\},
 \qquad
 |\nabla\lambda|>0\quad\text{on }D.
 \label{eq:V19-regular-wall}
\end{equation}
Fix a relatively compact tube $U$ on which $\lambda$ is the unique small
Wilson eigenvalue.  Suppose the control density is $a\,d\operatorname{vol}$
and the active line ratio has the form
\begin{equation}
 W_{\rm line}=|\lambda|^\beta b,
 \qquad
 \beta\ge0,
 \qquad
 a,b>0\text{ continuous on }\overline U.
 \label{eq:V19-wall-order}
\end{equation}
Put
\begin{equation}
 h_b(t)
 =\int_{\{\lambda=t\}\cap U}
   \frac{ab}{|\nabla\lambda|}\,d\sigma_t.
 \label{eq:V19-weighted-coarea-density}
\end{equation}
After shrinking $U$, $h_b$ is continuous at zero and $h_b(0)>0$.

\begin{theorem}[Exact weighted tube law]
\label{thm:V19-weighted-tube}
For $0<r<r_0$,
\begin{equation}
 \boxed{
 \Lambda^{\rm abs}
 \bigl(U\cap\{0<|\lambda|<r\}\bigr)
 =\int_{-r}^{r}|t|^\beta h_b(t)\,dt.}
 \label{eq:V19-weighted-tube-exact}
\end{equation}
Consequently,
\begin{equation}
 \boxed{
 \lim_{r\downarrow0}
 \frac{\Lambda^{\rm abs}
 (U\cap\{0<|\lambda|<r\})}{r^{\beta+1}}
 =\frac{2h_b(0)}{\beta+1}.}
 \label{eq:V19-weighted-tube-asymptotic}
\end{equation}
On the one-sided sector $\{\lambda>0\}$ the factor two is omitted.  After
normalization by $Z_{\rm abs}$ the same formulas hold with $h_b$ replaced by
$Z_{\rm abs}^{-1}h_b$.
\end{theorem}

\begin{proof}
Apply the coarea formula to the integrand
$\one_{0<|\lambda|<r}|\lambda|^\beta ab$.  This gives
\eqref{eq:V19-weighted-tube-exact}.  Substitute $t=ru$ and use continuity of
$h_b$ together with dominated convergence:
\[
 r^{-\beta-1}\int_{-r}^{r}|t|^\beta h_b(t)\,dt
 =\int_{-1}^{1}|u|^\beta h_b(ru)\,du
 \longrightarrow\frac{2h_b(0)}{\beta+1}.
\]
The one-sided and normalized statements are immediate.
\end{proof}

\begin{corollary}[Sharp inverse-gap threshold after line reweighting]
\label{cor:V19-inverse-gap-threshold}
Under the hypotheses of \Cref{thm:V19-weighted-tube}, for $q>0$,
\begin{equation}
 \boxed{
 \int_U |\lambda|^{-q}\,d\mu^{\rm abs}<\infty
 \iff q<\beta+1.}
 \label{eq:V19-inverse-gap-threshold}
\end{equation}
Thus an order-$\beta$ line zero raises the regular-divisor inverse-gap threshold
from one to $\beta+1$, without changing the essential gap.
\end{corollary}

\begin{proof}
By coarea, the singular part is comparable above and below to
$\int_0^{r_0}t^{\beta-q}\,dt$, since $h_b$ has positive finite bounds near
zero.  This integral is finite exactly when $\beta-q>-1$.
\end{proof}

\begin{corollary}[Line-divisor order and the Wilson-wall exponent]
\label{cor:V19-line-order}
Suppose local nonvanishing control frames have been fixed and the active line
sections admit
\begin{equation}
 \sigma=\lambda^{m_\sigma}\widehat\sigma,
 \qquad
 p=\lambda^{m_p}\widehat p,
 \label{eq:V19-line-orders}
\end{equation}
where $m_\sigma,m_p\in\mathbb N_0$ and
$\widehat\sigma,\widehat p$ are nonzero on $D$.  Then
\begin{equation}
 \boxed{
 \beta_{\rm line}=2m_\sigma+2m_p.}
 \label{eq:V19-beta-line}
\end{equation}
In particular, one simple zero of either squared line modulus gives
$\beta_{\rm line}=2$ and makes all inverse moments of order $q<3$ finite near a
regular Wilson wall.  If neither line vanishes on the Wilson divisor, then
$\beta_{\rm line}=0$ and the original threshold $q<1$ remains.
\end{corollary}

\begin{proof}
Taking squared norms in \eqref{eq:V19-line-orders} and dividing by the
nonzero control norms gives
$W_{\rm line}=|\lambda|^{2m_\sigma+2m_p}b$ with $b>0$ continuous.  Apply
\Cref{cor:V19-inverse-gap-threshold}.
\end{proof}

\begin{assessment}[The missing incidence is not automatic]
\label{ass:V19-wall-line-incidence}
A Wilson zero is a zero of $H_W$.  An active determinant or Pfaffian zero is a
zero of a separately assembled fermionic line section.  Neither implies the
other.  The integers $m_\sigma,m_p$ in
\eqref{eq:V19-line-orders} are therefore physical same-history incidence data.
A simple zero of an internal finite-Dirac constituent, a control section, or a
spectator line cannot be substituted for a zero of the complete active line on
the Wilson wall.
\end{assessment}

\begin{firewall}[Zero-mode saturation is not a probability reweighting]
\label{fire:V19-saturation-not-law}
A zero-mode insertion may saturate the line-valued Grassmann amplitude and
replace a vanishing determinant by a nonzero reduced exterior coefficient.
That changes the target writer and its line degree.  It changes the positive
sampling law only when the saturated modulus itself is independently adopted
and normalized as the physical cylinder.  The net wall order must therefore be
computed after the exact target and positive-law roles have been declared.
\end{firewall}

% =============================================================================
\section{Moving weighted sectors: surface term, integrable score, and Fisher threshold}
\label{sec:V19-moving-weighted-sector}
% =============================================================================

Let $s\mapsto\lambda_s$ and $s\mapsto a_s$ be smooth near $s=0$, with
$a_s>0$, and put
\begin{equation}
 d\Lambda_s
 = (\lambda_s)_+^\beta a_s\,d\operatorname{vol},
 \qquad
 Z_s=\Lambda_s(M),
 \qquad
 \mu_s=Z_s^{-1}\Lambda_s.
 \label{eq:V19-moving-weighted-law}
\end{equation}
Assume $D_0=\{\lambda_0=0\}$ is regular.  Dots denote derivatives at zero.

\begin{theorem}[The line zero converts the surface derivative into an interior singular score]
\label{thm:V19-weighted-shape-derivative}
Let $F_s$ be a smooth bounded writer.
\begin{enumerate}[label=\textup{(\roman*)},leftmargin=3em]
\item If $\beta=0$, the derivative contains the surface source
      \[
       \int_{D_0}F_0a_0
       \frac{\dot\lambda}{|\nabla\lambda_0|}\,d\sigma,
      \]
      as in Version~V18.
\item If $\beta>0$, there is no surface delta term and
      \begin{equation}
       \boxed{
       \left.\frac d{ds}\right|_0
       \int F_s\,d\Lambda_s
       =\int_{\{\lambda_0>0\}}
       \left(
       \dot F+F_0\dot{\log a}
       +\beta F_0\frac{\dot\lambda}{\lambda_0}
       \right)d\Lambda_0.}
       \label{eq:V19-weighted-shape-unnormalized}
      \end{equation}
      The singular writer $\dot\lambda/\lambda_0$ is integrable under
      $\Lambda_0$ for every $\beta>0$.
\item With
      \begin{equation}
       g=\dot{\log a}+\beta\frac{\dot\lambda}{\lambda_0},
       \label{eq:V19-L1-score}
      \end{equation}
      the normalized derivative is
      \begin{equation}
       \boxed{
       \left.\frac d{ds}\right|_0\mathbb E_{\mu_s}F_s
       =\mathbb E_{\mu_0}\dot F
       +\operatorname{Cov}_{\mu_0}(F_0,g).}
       \label{eq:V19-normalized-response}
      \end{equation}
\end{enumerate}
Thus every positive wall order removes the measure-valued first derivative,
but it does not yet guarantee a finite Fisher metric.
\end{theorem}

\begin{proof}
For $\beta=0$, differentiation of the moving Heaviside factor gives the
surface distribution already evaluated in Version~V18.  For $\beta>0$, the
locally integrable distribution $x\mapsto x_+^\beta$ has derivative
$\beta x_+^{\beta-1}$.  Differentiating
$F_s(\lambda_s)_+^\beta a_s$ gives
\eqref{eq:V19-weighted-shape-unnormalized}.  In a coarea coordinate
$t=\lambda_0$, the singular term has local absolute integral bounded by a
constant times
$\int_0^{r_0}t^{\beta-1}\,dt$, which is finite exactly for $\beta>0$.
Apply the same formula to $F_s=1$ and differentiate $Z_s^{-1}$ to obtain the
covariance identity \eqref{eq:V19-normalized-response}.
\end{proof}

We now distinguish $L^1$ differentiability from differentiability in quadratic
mean.  Use the squared Hellinger convention of
\eqref{eq:V19-Hellinger-exact}.

\begin{theorem}[Sharp quadratic-mean threshold at a transverse moving wall]
\label{thm:V19-DQM-threshold}
Assume there is a relatively open patch $D_*\subseteq D_0$ of positive surface
measure on which $a_0$ is bounded above and below and
\begin{equation}
 |\dot\lambda|\ge v_*>0.
 \label{eq:V19-transverse-velocity}
\end{equation}
Then the family \eqref{eq:V19-moving-weighted-law} has the following sharp
local alternatives.
\begin{enumerate}[label=\textup{(\roman*)},leftmargin=3em]
\item If $0\le\beta<1$, there is $c>0$ such that
      \begin{equation}
       \boxed{
       \mathsf H^2(\mu_s,\mu_0)
       \ge c|s|^{\beta+1}}
       \label{eq:V19-H-lower-subcritical}
      \end{equation}
      for one sign of sufficiently small $s$.  Hence
      $\mathsf H^2(\mu_s,\mu_0)/s^2\to\infty$ and the family is not
      differentiable in quadratic mean.
\item If $\beta=1$, there are $c,r>0$ such that
      \begin{equation}
       \boxed{
       \mathsf H^2(\mu_s,\mu_0)
       \ge c s^2\log\frac r{|s|}}
       \label{eq:V19-H-lower-borderline}
      \end{equation}
      for one sign of sufficiently small $s$.  Again there is no quadratic-mean
      derivative.
\item If $\beta>1$, the family is differentiable in quadratic mean and its
      score is
      \begin{equation}
       \boxed{
       \ell
       =g-\mathbb E_{\mu_0}g,
       \qquad
       g=\dot{\log a}
       +\beta\frac{\dot\lambda}{\lambda_0}.}
       \label{eq:V19-DQM-score}
      \end{equation}
      In particular,
      \begin{equation}
       \boxed{
       \sqrt{d\mu_s}
       =\sqrt{d\mu_0}
       \left(1+\frac s2\ell\right)
       +o_{L^2}(|s|),
       \qquad
       \mathcal I=\mathbb E_{\mu_0}\ell^2<\infty.}
       \label{eq:V19-DQM-expansion}
      \end{equation}
\end{enumerate}
The threshold $\beta>1$ is sharp whenever the wall velocity is nonzero on a
positive surface patch.
\end{theorem}

\begin{proof}
Choose tubular coordinates $(t,y)$ with $t=\lambda_0\ge0$ on the occupied side
and $y\in D_*$.  After restricting the patch, the Jacobian and $a_s$ have
positive finite bounds, and transversality implies that for one sign of $s$ a
strip of width comparable to $|s|$ is removed or added.

For $0\le\beta<1$, the squared root on the strip present in only one of the two
laws contributes at least a constant multiple of
$\int_0^{c|s|}t^\beta\,dt=c'|s|^{\beta+1}$.  Normalization changes the square
roots only by a smooth scalar factor and cannot cancel a positive strip
contribution.  This proves \eqref{eq:V19-H-lower-subcritical}.

For $\beta=1$, restrict further to $2v_*|s|\le t\le r$.  Up to uniformly bounded
positive factors, the two unnormalized square roots contain $\sqrt t$ and
$\sqrt{t-vs}$.  The identity
\[
 \left(\sqrt t-\sqrt{t-vs}\right)^2
 =\frac{v^2s^2}{(\sqrt t+\sqrt{t-vs})^2}
\]
gives a lower bound $c s^2/t$.  Integration from a constant multiple of
$|s|$ to $r$ proves \eqref{eq:V19-H-lower-borderline}.  The scalar
normalization correction is $O(s)$ in the square root and contributes only
$O(s^2)$.

For $\beta>1$, the unnormalized square root is
$q_s=(\lambda_s)_+^{\beta/2}a_s^{1/2}$.  Away from the wall it has derivative
\[
 \dot q_0
 =\frac12q_0
 \left(\dot{\log a}+\beta\frac{\dot\lambda}{\lambda_0}\right).
\]
Its squared singular part is locally bounded by a constant multiple of
$t^{\beta-2}$, which is integrable exactly when $\beta>1$.  The moving strip
has squared $L^2$ mass $O(|s|^{\beta+1})=o(s^2)$.  Dominated convergence on the
remaining tube and ordinary differentiation away from the tube therefore give
$L^2$ differentiability of $q_s$.  Differentiating the scalar normalizer
centres the logarithmic derivative and yields
\eqref{eq:V19-DQM-score}--\eqref{eq:V19-DQM-expansion}.  The two lower-bound
arguments show sharpness.
\end{proof}

\begin{corollary}[Squared line zeros give a clean integer alternative]
\label{cor:V19-integer-line-DQM}
Under \Cref{cor:V19-line-order}, suppose the wall moves transversely.
\begin{enumerate}[label=\textup{(\roman*)},leftmargin=3em]
\item If $m_\sigma=m_p=0$, then $\beta_{\rm line}=0$ and the Version~V18
      Hellinger cusp remains.
\item If $m_\sigma+m_p\ge1$, then $\beta_{\rm line}\ge2$ and the active
      positive line-modulus law is differentiable in quadratic mean.  Its
      finite score contains the singular but square-integrable term
      \[
       2(m_\sigma+m_p)\frac{\dot\lambda}{\lambda}.
      \]
\end{enumerate}
A simple squared line zero therefore does not delete the wall response; it
converts the surface source into an ordinary finite-energy score.
\end{corollary}

\begin{proof}
Use \eqref{eq:V19-beta-line} and apply
\Cref{thm:V19-DQM-threshold}.
\end{proof}

\begin{corollary}[Finite command-bank source metric on the damped-wall branch]
\label{cor:V19-wall-source-metric}
Let $u,v$ range over a finite physical command space.  Suppose the common
active positive law has wall exponent $\beta>1$ for every command and define
\begin{equation}
 g_u
 =D_u\log a+\beta\frac{D_u\lambda}{\lambda}.
 \label{eq:V19-command-log-writer}
\end{equation}
Then the physical source metric contributed by this probability law is
\begin{equation}
 \boxed{
 \mathsf M_{uv}^{\rm wall}
 =\operatorname{Cov}_{\mu^{\rm abs}}(g_u,g_v).}
 \label{eq:V19-wall-Fisher-metric}
\end{equation}
It is positive semidefinite and finite.  If
$a=e^{-S_B}\|\widehat\sigma\|^2|\widehat p|^2$ in a local line factorization,
then $D_u\log a$ contains the bosonic action and the nonvanishing line-amplitude
responses on the same history.  Separate control-law, determinant, and
Pfaffian diagonal metrics do not determine the mixed covariance in
\eqref{eq:V19-wall-Fisher-metric}.
\end{corollary}

\begin{proof}
Apply the multivariate quadratic-mean derivative of
\Cref{thm:V19-DQM-threshold} in each command direction.  The Gram of the
centred scores is positive semidefinite, and finiteness follows from
$\beta>1$.  Expansion of the logarithm gives the stated same-history terms.
\end{proof}

\begin{assessment}[What line damping does and does not close]
\label{ass:V19-line-damping-scope}
A branch with $\beta>1$ licenses an ordinary finite source metric for motion of
the active positive line-modulus law.  It does not give a historywise hard
Wilson gap, because \eqref{eq:V19-ess-gap-invariance} still applies.  The
historywise overlap sign remains undefined on the divisor and discontinuous
across changing inertia.  Therefore:
\begin{itemize}[leftmargin=2em]
\item a selected hard-sector historywise crossing still requires the sector
      selector and V13--V16;
\item a divisor-null active law with $\beta>1$ may enter a mixed or
      record-relative source/Fisher compiler with the exact line-weighted
      metric;
\item a zero-mode-saturated line-valued target remains in the exterior/divisor
      packet until its own positive scalarization is supplied.
\end{itemize}
\end{assessment}

% =============================================================================
\section{Control-law score versus active-law score}
\label{sec:V19-score-reweight}
% =============================================================================

The preceding wall theorem can be combined with ordinary positive
reweighting away from the divisor.  Let $u$ be a physical command and suppose
on a fixed common support that
\begin{equation}
 d\Lambda_u^+=w_u^+\,d\lambda,
 \qquad
 d\Lambda_u^{\rm abs}=W_u\,d\Lambda_u^+,
 \label{eq:V19-command-reweight}
\end{equation}
with positive differentiable densities.  Put
\begin{equation}
 h_u^+=D_u\log w_u^+,
 \qquad
 h_u^{\rm line}=D_u\log W_u.
 \label{eq:V19-two-log-derivatives}
\end{equation}

\begin{theorem}[Exact active score and nonadditive Fisher assembly]
\label{thm:V19-active-score}
On every common-support quadratic-mean branch, the active normalized score is
\begin{equation}
 \boxed{
 \ell_u^{\rm abs}
 =h_u^++h_u^{\rm line}
 -\mathbb E_{\mu^{\rm abs}}
  (h_u^++h_u^{\rm line}).}
 \label{eq:V19-active-score}
\end{equation}
For a finite command bank, its metric is
\begin{equation}
 \boxed{
 \mathsf M_{uv}^{\rm abs}
 =\operatorname{Cov}_{\mu^{\rm abs}}
 (h_u^++h_u^{\rm line},h_v^++h_v^{\rm line}).}
 \label{eq:V19-active-Fisher}
\end{equation}
In particular, it contains the two mixed control/line blocks.  The control
metric and the line metric, even if both are known, do not determine
$\mathsf M^{\rm abs}$ without those mixed same-history rows.
\end{theorem}

\begin{proof}
Differentiate the logarithm of the unnormalized product density in
\eqref{eq:V19-command-reweight}.  Differentiation of its normalizer subtracts
the active expectation, giving \eqref{eq:V19-active-score}.  Taking the Gram
of the centred scores gives \eqref{eq:V19-active-Fisher}; expansion displays
the two cross covariances.
\end{proof}

\begin{countertheorem}[The two diagonal score metrics do not determine the active metric]
\label{cth:V19-diagonal-metrics-no-active}
There are pairs of centred scalar score writers $(B,X)$ with identical
variances
\[
 \mathbb E B^2=\mathbb E X^2=1
\]
but with
\[
 \operatorname{Var}(B+X)=0
 \quad\text{or}\quad
 \operatorname{Var}(B+X)=4.
\]
Thus separately evaluated bosonic-control and fermion-line source metrics can
cancel or reinforce exactly.
\end{countertheorem}

\begin{proof}
Take one centred unit-variance writer $B$.  The choices $X=-B$ and $X=B$ have
the same two diagonal variances and give the two stated assembled variances.
\end{proof}

\begin{firewall}[Importance sampling is not source-metric inheritance]
\label{fire:V19-importance-no-metric}
The control law may estimate active expectations through
\eqref{eq:V19-reweighted-expectation}.  That computational identity does not
make the control Fisher metric the active Fisher metric.  The latter is
\eqref{eq:V19-active-Fisher}, or the weighted-wall specialization
\eqref{eq:V19-wall-Fisher-metric}, and must be populated before leverage,
reserve, Duhamel ancestry, or current/stress costs are interpreted.
\end{firewall}

% =============================================================================
\section{Exact controls and replay values}
\label{sec:V19-controls}
% =============================================================================

\subsection{A nonconstant line ratio changes the law without creating a hard gap}

Take the control law $d\mu^+=dx$ on $[0,1]$ and the line ratio
$W_{\rm line}=x^2$.  Then
\begin{equation}
 d\mu^{\rm abs}=3x^2\,dx.
 \label{eq:V19-beta2-law}
\end{equation}
For $A=[0,1/2]$,
\begin{equation}
 \boxed{
 \mu^+(A)=\frac12,
 \qquad
 \mu^{\rm abs}(A)=\frac18.}
 \label{eq:V19-sector-mass-control}
\end{equation}
The exact distances are
\begin{equation}
 \boxed{
 d_{\rm TV}(\mu^{\rm abs},\mu^+)
 =\frac{2}{3\sqrt3},
 \qquad
 \mathsf H^2(\mu^{\rm abs},\mu^+)=2-\sqrt3.}
 \label{eq:V19-control-active-distances}
\end{equation}
Both laws are equivalent and both have essential infimum zero for the gap
$\delta(x)=x$.

\begin{proof}
The sector mass is $\int_0^{1/2}3x^2dx=1/8$.  The sign change of
$3x^2-1$ occurs at $x=1/\sqrt3$, and direct integration gives the total
variation.  The Hellinger affinity is
$\int_0^1\sqrt3x\,dx=\sqrt3/2$, yielding the displayed square distance.
\end{proof}

\subsection{A simple squared line zero restores a finite Fisher metric}

For $0\le s<1$, define
\begin{equation}
 p_s(x)
 =\frac{3(x-s)^2}{(1-s)^3}\one_{(s,1)}(x).
 \label{eq:V19-moving-beta2-density}
\end{equation}
This is the normalized one-dimensional model of a transverse wall with
$\beta=2$.  Its score at $s=0$ is
\begin{equation}
 \boxed{
 \ell(x)=3-\frac2x,}
 \label{eq:V19-beta2-score}
\end{equation}
and
\begin{equation}
 \boxed{
 \mathbb E_{p_0}\ell=0,
 \qquad
 \mathcal I=\mathbb E_{p_0}\ell^2=3.}
 \label{eq:V19-beta2-Fisher}
\end{equation}
For the writer $F(x)=x$,
\begin{equation}
 \boxed{
 \mathbb E_{p_s}x=\frac34+\frac s4,
 \qquad
 \operatorname{Cov}_{p_0}(x,\ell)=\frac14.}
 \label{eq:V19-beta2-response}
\end{equation}
The exact Hellinger distance is
\begin{equation}
 \boxed{
 \mathsf H^2(p_s,p_0)
 =2\left[
 1-\frac{1-\frac32s+\frac12s^3}{(1-s)^{3/2}}
 \right],}
 \label{eq:V19-beta2-Hellinger}
\end{equation}
so
\begin{equation}
 \boxed{
 \lim_{s\downarrow0}\frac{\mathsf H^2(p_s,p_0)}{s^2}
 =\frac34=\frac{\mathcal I}{4}.}
 \label{eq:V19-beta2-Hellinger-limit}
\end{equation}
At $s=1/4$,
\begin{equation}
 \boxed{
 \mathsf H^2(p_{1/4},p_0)
 =2-\frac{9\sqrt3}{8}.}
 \label{eq:V19-beta2-quarter}
\end{equation}

\begin{proof}
Differentiate
$\log p_s=\log3+2\log(x-s)-3\log(1-s)$ on the common interior.  At zero this
gives \eqref{eq:V19-beta2-score}.  Integrating against $3x^2dx$ gives
$\mathbb E(1/x)=3/2$, $\mathbb E(1/x^2)=3$, and hence
\eqref{eq:V19-beta2-Fisher}.  The change of variable
$x=s+(1-s)y$ gives \eqref{eq:V19-beta2-response}.  Finally,
\[
 \int_s^1\sqrt{p_s(x)p_0(x)}\,dx
 =\frac{3}{(1-s)^{3/2}}
   \int_s^1x(x-s)\,dx
 =\frac{1-\frac32s+\frac12s^3}{(1-s)^{3/2}}.
\]
This proves \eqref{eq:V19-beta2-Hellinger}; expansion at zero and substitution
$s=1/4$ give the last two identities.
\end{proof}

\subsection{The same positive law carries inequivalent phase and Pfaffian rows}

On the same interval, the complex sections
\begin{equation}
 \sigma_\theta(x)=e^{i\theta}x
 \label{eq:V19-phase-family}
\end{equation}
and the real sections
\begin{equation}
 p_+(x)=x,
 \qquad
 p_-(x)=-x
 \label{eq:V19-pfaffian-sign-family}
\end{equation}
all induce the same positive modulus $x^2$.  Their phase or real orientation
is different.  This is an exact finite replay of
\Cref{fire:V19-modulus-no-phase}.

% =============================================================================
\section{The first executable control-to-active Wilson-wall card}
\label{sec:V19-card}
% =============================================================================

\begin{definition}[Control/active line-incidence card]
\label{def:V19-line-incidence-card}
At one predeclared finite cutoff and on one common bosonic history space,
freeze the packet
\begin{equation}
 \boxed{
 \begin{aligned}
 \mathfrak L_W=\bigl(
 &\Lambda^+,\sigma^+,p^+,\sigma,p,
 W_{\rm line},Z_+,Z_{\rm abs},\\
 &H_W,\lambda,D,\{p_k^+,p_k^{\rm abs}\},
 h_b(0),m_\sigma,m_p,\\
 &\{D_u\lambda\}_u,
 \mathsf M^{\rm abs},
 \text{line phase/Pfaffian orientation},\\
 &\text{one held-out reweighted writer and one adjacent card}
 \bigr).
 \end{aligned}}
 \label{eq:V19-line-incidence-card}
\end{equation}
Here $\lambda$ is the selected crossing eigenvalue on every regular divisor
face; a nonregular or multiple face carries its full zero-mode exterior packet
instead of a fitted scalar order.
\end{definition}

\begin{theorem}[Finite decision semantics]
\label{thm:V19-decision}
After common-history, positivity, line-frame, and normalization checks, the
card of \Cref{def:V19-line-incidence-card} has the following ordered outcomes.
\begin{enumerate}[label=\textup{(V19.\arabic*)},leftmargin=5.2em]
\item If $Z_{\rm abs}=0$ or $Z_{\rm abs}=\infty$, return the active line
      normalization/divisor obstruction.  Do not form a normalized active
      source metric.
\item If $W_{\rm line}$ is not an occurring same-history writer, retain only
      the control law and line-valued target.  The active probability claim is
      \stopstatus{}.
\item If the wall is stationary for every admitted command,
      $D_u\lambda|_D=0$, use the ordinary fixed-sector bulk score and continue
      to the source-metric card.
\item If the wall moves and $\beta=0$, return the Version~V18 surface-source and
      Hellinger-cusp obstruction to an ordinary Fisher/Duhamel entrance.
\item If the wall moves and $0<\beta\le1$, the first derivative is the
      integrable interior formula \eqref{eq:V19-normalized-response}, but no
      finite Fisher metric exists.  This is a typed $L^1$-but-not-$L^2$ source
      obstruction.
\item If the wall moves and $\beta>1$, populate
      \eqref{eq:V19-wall-Fisher-metric}.  This is a \passstatus{} for the
      probability-source metric, not for a historywise hard overlap crossing.
\item Independently, a nonzero line phase, Pfaffian-orientation, zero-mode, or
      held-out reweighting residual is retained at its own line or occurrence
      gate.
\end{enumerate}
An interval which does not separate the relevant threshold returns
\stopstatus{} on the same finite card.  It is not repaired by choosing a
preferred line modulus after the source metric has been inspected.
\end{theorem}

\begin{proof}
The first two items are the domain of
\Cref{thm:V19-RN-exact}.  Stationary-wall reduction is immediate from
\eqref{eq:V19-L1-score}.  The next three alternatives are exactly
\Cref{thm:V19-DQM-threshold}.  The final item is the independent line and
held-out incidence firewall.  These cases are ordered by the earliest failed
physical row and exhaust the regular-face branches.
\end{proof}

\begin{theorem}[Adjacent-cutoff transport of the active line law]
\label{thm:V19-cutoff-transport}
Let $\pi:X_{n+1}\to X_n$ be the declared bosonic cutoff map.  Transport the
fine laws to $X_n$ and write
\begin{equation}
 \nu_n^+=\pi_*\mu_{n+1}^+,
 \qquad
 \nu_n^{\rm abs}=\pi_*\mu_{n+1}^{\rm abs}.
 \label{eq:V19-pushed-fine-laws}
\end{equation}
Let
\begin{equation}
 r_n
 =\frac{W_n}{\mathbb E_{\mu_n^+}W_n}
 =\frac{d\mu_n^{\rm abs}}{d\mu_n^+},
 \qquad
 \widetilde r_{n+1}
 =\frac{d\nu_n^{\rm abs}}{d\nu_n^+}.
 \label{eq:V19-normalized-weight-cutoff}
\end{equation}
Assume
\begin{equation}
 \|\mu_n^+-\nu_n^+\|_{\rm TV}
 \le\delta_n^+,
 \qquad
 \int_{X_n}|r_n-\widetilde r_{n+1}|\,d\nu_n^+
 \le\delta_n^{\rm line}.
 \label{eq:V19-cutoff-assumptions}
\end{equation}
Then
\begin{equation}
 \boxed{
 \|\mu_n^{\rm abs}-\nu_n^{\rm abs}\|_{\rm TV}
 \le
 \|r_n\|_\infty\delta_n^+
 +\frac12\delta_n^{\rm line}.}
 \label{eq:V19-active-TV-transport}
\end{equation}
More generally, without an $L^\infty$ bound the same conclusion holds with the
first term replaced by the exact weighted variation
$\frac12\int r_n\,d|\mu_n^+-\nu_n^+|$.
Summability of these complete defects transports every bounded active writer
and every fixed active inertia-sector mass.  It does not transport the phase or
Pfaffian orientation, which remain separate line-cocycle rows.
\end{theorem}

\begin{proof}
Using
$d\mu_n^{\rm abs}=r_n\,d\mu_n^+$ and
$d\nu_n^{\rm abs}=\widetilde r_{n+1}\,d\nu_n^+$, add and subtract
$r_n\,d\nu_n^+$.  Total variation and the triangle inequality give the exact
weighted term plus one half of the displayed $L^1(\nu_n^+)$ density
difference.  The $L^\infty$ bound gives
\eqref{eq:V19-active-TV-transport}.  Summability gives a Cauchy sequence on
every fixed old screen, and total-variation duality transports bounded
writers and sector indicators.
\end{proof}

\begin{firewall}[The line ratio must be transported with its denominator]
\label{fire:V19-ratio-transport}
Transport of $\|\sigma\|^2|p|^2$ alone does not transport
$W_{\rm line}$.  The control denominator, its normalization, and their
same-history incidence must travel on the same line trivialization.  Near a
control zero the ratio is not a valid coordinate; such a card leaves the
nonvanishing-control branch and enters the full divisor/line-overlap packet.
\end{firewall}

% =============================================================================
\section{Version V19 theorem, ledger, and next acquisition}
\label{sec:V19-terminal}
% =============================================================================

\begin{theorem}[Version V19 control-to-active wall theorem]
\label{thm:V19-terminal}
Versions~V1--V18 are retained.  Version~V19 additionally proves:
\begin{enumerate}[label=\textup{(E19.\arabic*)},leftmargin=5.5em]
\item the nonzero control cylinder, active positive line-modulus cylinder, and
      line-valued saturated amplitude are three distinct physical objects;
\item their normalized positive laws are related by the exact
      Radon--Nikodym, sector-mass, total-variation, and Hellinger formulas
      \eqref{eq:V19-normalized-RN}--\eqref{eq:V19-Hellinger-exact};
\item equivalent line reweighting preserves the essential Wilson gap, so a
      divisor zero can soften moments without creating hard admissibility;
\item a transverse line weight $|\lambda|^\beta$ changes the regular divisor
      tube exponent from one to $\beta+1$ and the inverse-gap threshold from
      $q<1$ to $q<\beta+1$;
\item squared determinant/Pfaffian line zeros have even wall order
      $2m_\sigma+2m_p$;
\item every positive wall order removes the surface delta from the first
      response and replaces it by the integrable writer
      $\beta\dot\lambda/\lambda$;
\item the sharp Fisher threshold is $\beta>1$: below it the squared Hellinger
      distance decays more slowly than $s^2$, at $\beta=1$ it has a
      logarithmic divergence, and above it the score
      \eqref{eq:V19-DQM-score} is square integrable;
\item on the passing branch the source-action metric is the complete
      same-history covariance \eqref{eq:V19-wall-Fisher-metric}, including
      control/line cross blocks;
\item the exact $\beta=2$ control has Fisher information three and satisfies
      the explicit Hellinger formula \eqref{eq:V19-beta2-Hellinger};
\item adjacent-cutoff transport requires the control law and normalized line
      ratio together, while phase and real orientation remain independent.
\end{enumerate}
No item supplies the actual RPESM line ratio on a Wilson divisor, proves that
the active determinant or Pfaffian vanishes there, selects a hard index sector,
or evaluates the physical reflected crossing.
\end{theorem}

\begin{proof}
Items~\textup{(E19.1)--(E19.3)} are
\Cref{thm:V19-RN-exact,prop:V19-support-equivalence}.  Items
\textup{(E19.4)--(E19.5)} are
\Cref{thm:V19-weighted-tube,cor:V19-inverse-gap-threshold,cor:V19-line-order}.
Items~\textup{(E19.6)--(E19.8)} are
\Cref{thm:V19-weighted-shape-derivative,thm:V19-DQM-threshold,cor:V19-wall-source-metric,thm:V19-active-score}.
Item~\textup{(E19.9)} is \Cref{sec:V19-controls}, and item
\textup{(E19.10)} is \Cref{thm:V19-cutoff-transport}.  The final limitations
name physical rows not populated by any of those theorems.
\end{proof}

\begingroup\small
\begin{longtable}{>{\raggedright\arraybackslash}p{0.27\textwidth}
                  >{\raggedright\arraybackslash}p{0.20\textwidth}
                  >{\raggedright\arraybackslash}p{0.45\textwidth}}
\caption{Version V19 proof-status ledger.}\label{tab:V19-ledger}\\
\toprule
\textbf{Row}&\textbf{Status}&\textbf{Exact conclusion}\\
\midrule
\endfirsthead
\toprule
\textbf{Row}&\textbf{Status}&\textbf{Exact conclusion}\\
\midrule
\endhead
Control versus active positive law
&\passstatus
&Exact density ratio, expectation, sector mass, total variation, and Hellinger
formulas.\\[0.3em]
Essential Wilson gap under nonzero reweighting
&\passstatus
&Equivalent laws have identical null sets and identical essential gap; line
damping is not hard admissibility.\\[0.3em]
Regular weighted divisor tube
&\passstatus
&Tube exponent $\beta+1$ and inverse-gap threshold $q<\beta+1$.\\[0.3em]
Line/Wilson incidence order
&Finite compiler; physical value open
&Squared line sections give $\beta=2m_\sigma+2m_p$, but the active RPESM order
on the Wilson wall is not supplied.\\[0.3em]
Weighted moving-wall first derivative
&\passstatus
&For $\beta>0$ the surface delta disappears and the response uses the integrable
writer $\beta\dot\lambda/\lambda$.\\[0.3em]
Quadratic-mean threshold
&\passstatus
&Finite Fisher metric exactly on the transverse branch $\beta>1$; explicit
subcritical and logarithmic Hellinger lower bounds otherwise.\\[0.3em]
Active source metric
&Finite formula; physical value open
&Complete covariance of bosonic/control and line responses, including both
mixed blocks.\\[0.3em]
Line phase and Pfaffian orientation
&Independent
&Unchanged by positive line reweighting; no positive metric chooses them.\\[0.3em]
Zero-mode saturation
&Retained downstream
&May change the line-valued target order but changes the probability law only
through a separately normalized positive cylinder.\\[0.3em]
Adjacent active-law transport
&\passstatus
&Exact control-law plus normalized-ratio defect; phase/orientation transport is
separate.\\[0.3em]
Selected RPESM wall/line card
&\stopstatus
&No supplied finite card gives $W_{\rm line}$, its wall order, sector masses,
wall velocity, and held-out active reweighting on one selected history.\\
\bottomrule
\end{longtable}
\endgroup

\begin{assessment}[Physical status after Version V19]
\label{ass:V19-status}
The projective regulator proves finite population under its positive control
cylinder and separately carries determinant/Pfaffian line data.  Version~V19
shows exactly what must be added before that control population is interpreted
as the active source-metric experiment.  The physical verdict is
\begin{equation}
 \boxed{
 \begin{array}{ll}
 \mathsf{PASS}:&
 \begin{gathered}
 \text{control/active reweighting, weighted tube,}\\
 \text{and Fisher-threshold compiler}
 \end{gathered}\\[0.35em]
 \mathsf{TYPED\ OBSTRUCTION}:&
 \begin{gathered}
 \text{using the control Fisher metric for an}\\
 \text{unverified active line law}
 \end{gathered}\\[0.35em]
 \mathsf{STOP}:&
 \begin{gathered}
 \text{at the first common-history Wilson-wall/}\\
 \text{active-line incidence Read}
 \end{gathered}
 \end{array}}
 \label{eq:V19-status}
\end{equation}
The result goes upstream of V13--V18 rather than adding another sign or
crossing theorem.  A later reflected-hopping, Duhamel, current/stress, or
continuum certificate cannot choose which normalized line law was physically
used.
\end{assessment}

\begin{openproblem}[Version V20: populate one control/active wall-incidence card]
\label{op:V19-next}
Preserve Versions~V1--V19.  Do not open another Wilson-sign approximation,
collar fit, exterior hierarchy, or Duhamel tail.  On one selected RPESM cutoff:
\begin{enumerate}[label=\textup{(V20.\arabic*)},leftmargin=5.2em]
\item freeze the positive control cylinder and the complete active fermion and
      real-line sections on one common bosonic history;
\item acquire $W_{\rm line}$ directly, its normalization, one held-out
      reweighted writer, and the exact active inertia-sector masses;
\item locate every regular Wilson-divisor face met by the support and determine
      whether the active line is nonzero there or has an integer transverse
      order $(m_\sigma,m_p)$;
\item acquire the command velocities $D_u\lambda$ on those faces and run the
      Version~V19 $\beta=0$, $0<\beta\le1$, or $\beta>1$ decision;
\item on the $\beta>1$ branch, populate the complete active source metric,
      including the mixed bosonic/line blocks, before interpreting reserve or
      leverage;
\item retain the complex phase, Pfaffian orientation, zero-mode saturation,
      hard-sector selector, and direct reflected crossing as separate rows;
\item repeat the normalized line ratio and one sector/writer check at the
      adjacent cutoff.
\end{enumerate}
If the active line ratio is not an occurring writer or its wall order remains
unresolved, stop at this card.  The nonzero control population theorem remains
valid at its stated scope, but it does not supply the active quantum source
metric.
\end{openproblem}

\section*{V19 exact replay and append-only record}
\addcontentsline{toc}{section}{V19 exact replay and append-only record}
The accompanying verifier uses Python standard-library rational and
high-precision decimal arithmetic.  It checks the control-to-active sector
mass, total-variation and Hellinger formulas; the $\beta=2$ normalized moving
wall, its score, Fisher information, response derivative, and exact Hellinger
value at $s=1/4$; the inverse-moment threshold on monomial wall weights; and the
phase/Pfaffian nonidentifiability of a positive modulus.  Ordinary and
optimized executions produce byte-identical JSON output.  The verifier does
not populate an RPESM line ratio or Wilson wall.

The complete Version~V18 source preceding its sole final active
\verb|\end{document}| is preserved byte for byte.  Its preterminal SHA--256
digest is recorded in the validation manifest.  A later continuation must
remove only the final active document-closing command, append new material
immediately before it, restore one terminal command, and increment the filename
to end in \texttt{\_V20.tex}.

% =============================================================================
% END APPENDED VERSION V19 MATERIAL
% =============================================================================


% =============================================================================
% BEGIN APPENDED VERSION V20 MATERIAL
% =============================================================================

\clearpage
\hypersetup{
 pdftitle={Renewal Einstein--Standard--Model Physical Source Metric, Duhamel Ancestry and Quantum Continuum Programme V20},
 pdfsubject={Reference-gauge elimination, intrinsic active line cylinder, and line-atlas incidence}
}
\pdfinfo{
 /Title (Renewal Einstein--Standard--Model Physical Source Metric, Duhamel Ancestry and Quantum Continuum Programme V20)
 /Subject (Reference-gauge elimination, intrinsic active line cylinder, and line-atlas incidence)
}

\section{Version V20: reference-gauge elimination and the intrinsic active line cylinder}
\label{sec:V20-reference-gauge}
% =============================================================================

Version~V19 correctly separated the positive control cylinder from the active
positive determinant/Pfaffian modulus and from the line-valued phase packet.
Its proposed next acquisition nevertheless treated the ratio
$W_{\rm line}$ and the control/line score blocks as if they were individually
intrinsic.  That is stronger than the populated regulator supports.  A
nonvanishing control section is a reference presentation of a positive
cylinder.  Multiplying that reference by an arbitrary positive function and
dividing the likelihood ratio by the same function leaves the active cylinder
unchanged.

The first upstream target is therefore not an isolated ratio but the complete
positive product
\begin{equation}
 \boxed{
 \dd\Lambda^{\rm abs}
 =e^{-S_B}\norm{\sigma}^{2}|p|^{2}\,\dd\lambda.}
 \label{eq:V20-active-cylinder}
\end{equation}
A control cylinder and a line ratio are admissible coordinates only when their
same-history product equals \eqref{eq:V20-active-cylinder}, including every
support component.  The active source metric is then differentiated from the
product.  The separate control, line, and mixed blocks depend on the chosen
reference gauge, although their assembled sum does not.

\begin{firewall}[What Version V20 changes]
\label{fire:V20-scope-correction}
No theorem of Version~V19 is deleted.  Its Radon--Nikodym, weighted-wall, and
Fisher-threshold formulas remain valid after one positive reference has been
physically frozen.  Version~V20 changes the acquisition order: first populate
the intrinsic active positive cylinder, or certify one common-history
reference/ratio factorization of it; only then use the corresponding ratio as a
sampling coordinate.  Another Wilson-sign, collar, exterior, or Duhamel
estimate would skip this earlier incidence row.
\end{firewall}

% =============================================================================
\section{The positive-reference gauge and its exact quotient}
\label{sec:V20-positive-reference}
% =============================================================================

Let $(X,\lambda)$ be the finite-cutoff bosonic history space with its declared
base measure.  Throughout Version~V20, for a finite signed measure $\nu$ write
\begin{equation}
 \norm{\nu}_{\rm var}:=|\nu|(X),
 \label{eq:V20-variation-norm}
\end{equation}
the full variation norm.  For two probability measures this is twice the
half-$L^1$ total-variation distance used in some earlier probability formulas.
Let $a,c\in L^1_+(\lambda)$, with $c>0$ almost everywhere, and suppose
\begin{equation}
 a=cw,
 \qquad w\ge0.
 \label{eq:V20-reference-factorization}
\end{equation}
Write
\begin{equation}
 \dd\Lambda_a=a\,\dd\lambda,
 \qquad
 \dd\Lambda_c=c\,\dd\lambda,
 \qquad
 Z_a=\Lambda_a(X),
 \qquad
 Z_c=\Lambda_c(X).
 \label{eq:V20-reference-measures}
\end{equation}

\begin{theorem}[Positive-reference gauge quotient]
\label{thm:V20-reference-gauge}
Assume $0<Z_a,Z_c<\infty$.  For every measurable $h>0$ such that
$hc\in L^1(\lambda)$, put
\begin{equation}
 c^{h}=hc,
 \qquad
 w^{h}=h^{-1}w.
 \label{eq:V20-gauge-action}
\end{equation}
Then
\begin{equation}
 \boxed{c^{h}w^{h}=cw=a.}
 \label{eq:V20-product-invariance}
\end{equation}
Conversely, if $c'>0$ and $c'w'=a$, then the unique positive gauge carrying
$(c,w)$ to $(c',w')$ is $h=c'/c$ almost everywhere.  Thus the orbits of the
positive-reference gauge are in one-to-one correspondence with intrinsic
active positive measures $\Lambda_a$.

For the normalized laws
\begin{equation}
 \dd\mu_c=Z_c^{-1}\dd\Lambda_c,
 \qquad
 \dd\mu_a=Z_a^{-1}\dd\Lambda_a,
 \label{eq:V20-normalized-laws}
\end{equation}
one has
\begin{equation}
 \boxed{
 \mathbb E_{\mu_c}w=\frac{Z_a}{Z_c},
 \qquad
 \dd\mu_a
 =\frac{w}{\mathbb E_{\mu_c}w}\,\dd\mu_c.}
 \label{eq:V20-reference-RN}
\end{equation}
The right side is unchanged by \eqref{eq:V20-gauge-action} after the reference
law is changed from $\mu_c$ to $\mu_{c^h}$.
\end{theorem}

\begin{proof}
Equation \eqref{eq:V20-product-invariance} is immediate.  If two positive
references factor the same $a$, then $h=c'/c$ is defined almost everywhere and
$c'=hc$; division of $c'w'=cw$ by $c'$ gives $w'=w/h$.  This also proves
uniqueness.  Finally,
\[
 \mathbb E_{\mu_c}w
 =Z_c^{-1}\int cw\,\dd\lambda
 =Z_a/Z_c,
\]
and substitution proves \eqref{eq:V20-reference-RN}.  Replacing $(c,w)$ by
$(hc,w/h)$ leaves the numerator measure $cw\,\dd\lambda$ unchanged.
\end{proof}

\begin{corollary}[The ratio alone is not an intrinsic writer]
\label{cor:V20-ratio-not-intrinsic}
Unless the control density $c$ is itself a physically occurring and calibrated
row, the function $w$ in \eqref{eq:V20-reference-factorization} is not an
intrinsic observable of the active cylinder.  Its complete invariant content
is the product measure $w\Lambda_c=\Lambda_a$.  In particular, a statement
about the size, variance, or cutoff transport of $w$ alone is reference-gauge
dependent.
\end{corollary}

\begin{proof}
Choose any nonconstant admissible $h$ in
\eqref{eq:V20-gauge-action}.  The active measure is unchanged while $w$ is
replaced by $w/h$.  Hence no quantity which changes under that replacement is
an invariant of $\Lambda_a$ alone.
\end{proof}

% =============================================================================
\section{Reference-gauge cancellation in the physical source metric}
\label{sec:V20-score-gauge}
% =============================================================================

Let the active and reference densities depend differentiably on a finite
command parameter.  Work on a positive branch on which the logarithmic
derivatives below belong to $L^2(\mu_a)$.  For a command $u$, put
\begin{equation}
 b_u=D_u\log c,
 \qquad
 \ell_u=D_u\log w,
 \qquad
 s_u=b_u+\ell_u.
 \label{eq:V20-score-split}
\end{equation}
The normalized active score is
\begin{equation}
 \widehat s_u
 =s_u-\mathbb E_{\mu_a}s_u.
 \label{eq:V20-active-score}
\end{equation}

\begin{theorem}[Gauge invariance of the assembled active score]
\label{thm:V20-score-gauge}
Let the reference gauge also depend differentiably on the command and write
\begin{equation}
 \phi_u=D_u\log h.
 \label{eq:V20-gauge-score}
\end{equation}
Under \eqref{eq:V20-gauge-action},
\begin{equation}
 \boxed{
 b_u^{h}=b_u+\phi_u,
 \qquad
 \ell_u^{h}=\ell_u-\phi_u,
 \qquad
 s_u^{h}=s_u.}
 \label{eq:V20-score-transform}
\end{equation}
Consequently the physical active source metric
\begin{equation}
 \boxed{
 \Msrc^{\rm abs}(u,v)
 =\operatorname{Cov}_{\mu_a}(s_u,s_v)}
 \label{eq:V20-intrinsic-active-metric}
\end{equation}
is reference-gauge invariant.

By contrast, the four blocks
\begin{equation}
 \operatorname{Cov}(b_u,b_v),\quad
 \operatorname{Cov}(b_u,\ell_v),\quad
 \operatorname{Cov}(\ell_u,b_v),\quad
 \operatorname{Cov}(\ell_u,\ell_v)
 \label{eq:V20-gauge-dependent-blocks}
\end{equation}
are not individually invariant.  Only their assembled sum
\begin{equation}
 \boxed{
 \Msrc^{\rm abs}
 =\Msrc^{bb}+\Msrc^{b\ell}
  +\Msrc^{\ell b}+\Msrc^{\ell\ell}}
 \label{eq:V20-block-sum}
\end{equation}
is physical without an independently frozen control gauge.
\end{theorem}

\begin{proof}
Taking the logarithmic derivative of $c^h=hc$ and $w^h=w/h$ gives
\eqref{eq:V20-score-transform}.  Therefore the centered sum
\eqref{eq:V20-active-score} and its covariance are unchanged.  Expanding the
covariance of $b+\ell$ gives \eqref{eq:V20-block-sum}.  Under the gauge, the
centered control column acquires the centered column of $\phi$, while the line
column loses the same column.  Unless that column vanishes, the individual
terms in \eqref{eq:V20-gauge-dependent-blocks} change.
\end{proof}

\begin{countertheorem}[The separate control and line metrics can be arbitrarily large]
\label{cth:V20-separate-metrics-arbitrary}
Let $X=\{-1,+1\}$ with the uniform active law.  For any $k>0$ and
$|t|<k^{-1}$, put
\begin{equation}
 a_t(x)=1,
 \qquad
 c_t(x)=1+tkx,
 \qquad
 w_t(x)=(1+tkx)^{-1}.
 \label{eq:V20-two-point-gauge}
\end{equation}
Then $c_tw_t=a_t$ for every $t$, so the active experiment is constant.  At
$t=0$,
\begin{equation}
 D_t\log c_t=kx,
 \qquad
 D_t\log w_t=-kx.
 \label{eq:V20-two-point-scores}
\end{equation}
Hence
\begin{equation}
 \boxed{
 \Msrc^{bb}=k^2,
 \qquad
 \Msrc^{\ell\ell}=k^2,
 \qquad
 \Msrc^{b\ell}=\Msrc^{\ell b}=-k^2,
 \qquad
 \Msrc^{\rm abs}=0.}
 \label{eq:V20-two-point-metrics}
\end{equation}
The separate diagonal source metrics may therefore be made arbitrarily large
while the intrinsic active source metric remains exactly zero.
\end{countertheorem}

\begin{proof}
The product identity is exact.  Under the uniform law, $x$ is centered with
variance one.  Substitution of \eqref{eq:V20-two-point-scores} into the four
covariances gives \eqref{eq:V20-two-point-metrics}.
\end{proof}

\begin{assessment}[Correction to the Version V19 metric mandate]
\label{ass:V20-metric-correction}
Version~V19 was right that the mixed control/line terms cannot be discarded
once a control gauge has been physically selected.  Version~V20 sharpens the
statement: those blocks are a decomposition coordinate, not four independent
physical source rows.  The required invariant acquisition is either the direct
active score \eqref{eq:V20-active-score}, or all four blocks on one frozen
control gauge together with their exact cancellation law
\eqref{eq:V20-block-sum}.
\end{assessment}

% =============================================================================
\section{Intrinsic active measure and the singular-reference branch}
\label{sec:V20-singular-reference}
% =============================================================================

A globally positive control density need not dominate an active cylinder if the
submitted control vanishes on an active component.  This failure is invisible
if one merely writes a formal ratio on the control support.

\begin{theorem}[Exact singular-reference residual]
\label{thm:V20-singular-reference}
Let $\Lambda_a$ and $\Lambda_c$ be finite positive measures on the same
history space.  Write the Lebesgue decomposition
\begin{equation}
 \boxed{
 \Lambda_a=w\Lambda_c+\Lambda_a^{\perp},
 \qquad
 \Lambda_a^{\perp}\perp\Lambda_c.}
 \label{eq:V20-Lebesgue-line}
\end{equation}
Then for every measurable $v\ge0$,
\begin{equation}
 \boxed{
 \norm{\Lambda_a-v\Lambda_c}_{\rm var}
 =\int|w-v|\,\dd\Lambda_c
  +\Lambda_a^{\perp}(X).}
 \label{eq:V20-TV-residual-split}
\end{equation}
Consequently
\begin{equation}
 \boxed{
 \inf_{v\ge0}
 \norm{\Lambda_a-v\Lambda_c}_{\rm var}
 =\Lambda_a^{\perp}(X),}
 \label{eq:V20-min-reference-residual}
\end{equation}
and the minimum is attained by the Radon--Nikodym derivative $v=w$.

Thus a ratio representation of the complete active cylinder exists exactly
when
\begin{equation}
 \boxed{\Lambda_a^{\perp}(X)=0.}
 \label{eq:V20-reference-complete}
\end{equation}
The singular mass is the minimum active positive stock missing from every
control-ratio representation based on $\Lambda_c$.
\end{theorem}

\begin{proof}
The signed measure $(w-v)\Lambda_c$ is supported on a set disjoint from
$\Lambda_a^{\perp}$.  Total variation is additive on mutually singular signed
measures, which gives \eqref{eq:V20-TV-residual-split}.  The first term is
nonnegative and vanishes for $v=w$, proving
\eqref{eq:V20-min-reference-residual} and
\eqref{eq:V20-reference-complete}.
\end{proof}

\begin{corollary}[Sharp target interval from unresolved singular line mass]
\label{cor:V20-singular-payoff}
Let $0<Z_a=\Lambda_a(X)<\infty$ and let $F$ be a bounded real writer.  Then
\begin{equation}
 \boxed{
 \mathbb E_{\mu_a}F
 =\frac{1}{Z_a}
 \left[
  \int Fw\,\dd\Lambda_c
  +\int F\,\dd\Lambda_a^{\perp}
 \right].}
 \label{eq:V20-singular-payoff-exact}
\end{equation}
If only
\begin{equation}
 m\le F\le M
 \quad\text{on }\supp\Lambda_a^{\perp},
 \qquad
 s_\perp=\Lambda_a^{\perp}(X)
 \label{eq:V20-singular-payoff-bounds}
\end{equation}
is known, the sharp compatible interval is
\begin{equation}
 \boxed{
 \frac{\int Fw\,\dd\Lambda_c+ms_\perp}{Z_a}
 \le
 \mathbb E_{\mu_a}F
 \le
 \frac{\int Fw\,\dd\Lambda_c+Ms_\perp}{Z_a}.}
 \label{eq:V20-singular-payoff-interval}
\end{equation}
\end{corollary}

\begin{proof}
Integrate $F$ against \eqref{eq:V20-Lebesgue-line} and divide by $Z_a$.
The bounds follow from \eqref{eq:V20-singular-payoff-bounds}.  They are sharp
over all singular measures of mass $s_\perp$ compatible with the displayed
information, by concentrating that mass where the lower or upper bound is
attained, or by the usual approximating argument for essential bounds.
\end{proof}

\begin{definition}[Gauge-invariant line-occurrence residual]
\label{def:V20-line-occurrence}
For a submitted control measure $\Lambda_c$ and ratio writer $v$, define
\begin{equation}
 \boxed{
 \mathbb C_{\rm line}^{\rm occ}(v\mid\Lambda_c,\Lambda_a)
 :=\norm{\Lambda_a-v\Lambda_c}_{\rm var}.}
 \label{eq:V20-line-occurrence-residual}
\end{equation}
If $\Lambda_c$ is replaced by $h\Lambda_c$ and $v$ by $v/h$, the residual is
unchanged.  Its minimum over all ratios is the singular mass in
\eqref{eq:V20-min-reference-residual}.
\end{definition}

\begin{firewall}[No zero-over-zero likelihood coordinate]
\label{fire:V20-zero-over-zero}
On a control zero, the quotient of active and control line norms is not a
writer.  If the active measure also vanishes there, the value of the formal
ratio is irrelevant and may be chosen only as a reference convention.  If the
active measure is positive there, the missing stock is the singular term in
\eqref{eq:V20-Lebesgue-line} and cannot be repaired by assigning an infinite or
regularized ratio.
\end{firewall}

% =============================================================================
\section{Global line geometry: section, density, phase, and orientation}
\label{sec:V20-line-geometry}
% =============================================================================

The positive cylinder uses Hermitian norms and is globally meaningful even
when the determinant or Pfaffian line has nontrivial topology.  A global
nonzero control \emph{section} is stronger: it trivializes the corresponding
line and supplies phase or orientation information which a positive density
does not contain.

\begin{theorem}[Nonzero controls, line triviality, and relative phase]
\label{thm:V20-line-triviality}
Let $\mathscr L\to X$ be a complex Hermitian line bundle over a paracompact
base.
\begin{enumerate}[label=\textup{(L20.\arabic*)},leftmargin=5.3em]
\item A continuous nowhere-zero global section $s_+$ exists if and only if
      $\mathscr L$ is trivial.
\item If $s$ and $s_+$ are sections of the same line and $s_+$ is nowhere zero,
      there is a unique complex function $r$ such that
      \begin{equation}
       s=rs_+,
       \qquad
       \boxed{\norm{s}^{2}/\norm{s_+}^{2}=|r|^2.}
       \label{eq:V20-relative-line-ratio}
      \end{equation}
      The phase of $r$ is a relative line phase.
\item If the active and control sections live in line bundles not equipped with
      a specified unitary isomorphism, their norm ratio is still a scalar
      function where the denominator is nonzero, but no relative complex phase
      is defined.
\end{enumerate}
The analogous real statement holds for a Pfaffian line.  A nowhere-zero real
section both trivializes and orients that line.  Squared norms alone do not
select the orientation.
\end{theorem}

\begin{proof}
A nowhere-zero section defines the bundle isomorphism
$(x,z)\mapsto z s_+(x)$ from $X\times\C$ to $\mathscr L$; a trivialization sends
the constant section one to a nowhere-zero section.  In a one-dimensional
fibre, every vector is a unique scalar multiple of a nonzero vector, which
defines $r$ and gives \eqref{eq:V20-relative-line-ratio} by the Hermitian norm.
Without a fibrewise line isomorphism, vectors in the two line fibres cannot be
divided, although their real norm functions can be compared.  The same proof
over $\R$ gives the real-line statement; choosing the sign of a nonzero real
frame is precisely choosing an orientation.
\end{proof}

\begin{proposition}[A positive control density needs no global control section]
\label{prop:V20-local-control-density}
Let $\mathscr L\to X$ be any Hermitian line bundle over a compact paracompact
base.  Choose a finite trivializing cover $(U_i)$, nowhere-zero local sections
$s_i$ on $U_i$, and a subordinate partition of unity $(\chi_i)$.  Then
\begin{equation}
 \boxed{
 c_{\rm loc}(x)
 =\sum_i\chi_i(x)\norm{s_i(x)}^2>0}
 \label{eq:V20-local-positive-density}
\end{equation}
is a global continuous positive scalar density.  It can be used as a sampling
reference in \Cref{thm:V20-reference-gauge} even when $\mathscr L$ is
nontrivial.  It does not define a global phase section, a Berry connection, or
a Pfaffian orientation.
\end{proposition}

\begin{proof}
Each term in \eqref{eq:V20-local-positive-density} extends by zero outside the
support of $\chi_i$.  At every point at least one subordinate partition
function is positive, and its local section is nonzero, so the sum is strictly
positive.  The construction uses only norms; it supplies no compatible choice
of local phases on overlaps and therefore no global section or orientation.
\end{proof}

\begin{corollary}[What a global RPESM control section would certify]
\label{cor:V20-control-section-certificate}
A globally nowhere-zero determinant control and a globally nowhere-zero real
Pfaffian control certify more than normalizability of the control law: on the
selected base they trivialize the corresponding line bundles and orient the
real line.  If the regulator supplies only local controls, the positive control
cylinder may still be assembled from their norm densities, but the phase,
transition cocycle, and real orientation must remain as a separate line-atlas
packet.
\end{corollary}

\begin{firewall}[Reflection positivity is not supplied by a partition of unity]
\label{fire:V20-local-density-not-RP}
\Cref{prop:V20-local-control-density} constructs a positive sampling density.
It does not prove that this density is the modulus of one reflection-positive
fermionic control section or that it obeys the required cutoff line transport.
Those are additional same-history and line-cocycle rows.
\end{firewall}

% =============================================================================
\section{A direct active Hellinger source, without a control decomposition}
\label{sec:V20-active-Hellinger}
% =============================================================================

On the branch where Version~V19 proves quadratic-mean differentiability, the
physical source metric can be acquired directly from neighbouring active
positive cylinders.  No control/line split is required.

\begin{theorem}[Intrinsic active square-root tangent]
\label{thm:V20-active-Hellinger}
Let
\begin{equation}
 \dd\mu_\theta
 =Z_\theta^{-1}a_\theta\,\dd\lambda,
 \qquad
 \psi_\theta=\sqrt{a_\theta/Z_\theta}
 \label{eq:V20-active-root}
\end{equation}
be a differentiable family in $L^2(\lambda)$, and suppose
$h_u=D_u\log a_\theta|_{\theta=0}\in L^2(\mu_0)$.  Then
\begin{equation}
 \boxed{
 2D_u\psi_0
 =\left(h_u-\mathbb E_{\mu_0}h_u\right)\psi_0.}
 \label{eq:V20-root-tangent}
\end{equation}
Consequently
\begin{equation}
 \boxed{
 4\ip{D_u\psi_0}{D_v\psi_0}_{L^2(\lambda)}
 =\operatorname{Cov}_{\mu_0}(h_u,h_v)
 =\Msrc^{\rm abs}(u,v).}
 \label{eq:V20-Hellinger-Fisher}
\end{equation}
For a physical shrinking family, the finite secant Gram
\begin{equation}
 \mathsf H_{\varepsilon}(u,v)
 =\left\langle
   \frac{2(\psi_{\varepsilon u}-\psi_0)}{\varepsilon},
   \frac{2(\psi_{\varepsilon v}-\psi_0)}{\varepsilon}
  \right\rangle
 \label{eq:V20-Hellinger-secant}
\end{equation}
converges to \eqref{eq:V20-Hellinger-Fisher}.  It depends only on the active
positive cylinders and is invariant under every reference gauge.
\end{theorem}

\begin{proof}
Differentiate $\psi_\theta=a_\theta^{1/2}Z_\theta^{-1/2}$.  Since
$D_uZ_0/Z_0=\mathbb E_{\mu_0}h_u$, one obtains
\eqref{eq:V20-root-tangent}.  Taking inner products and using
$|\psi_0|^2\dd\lambda=\dd\mu_0$ gives
\eqref{eq:V20-Hellinger-Fisher}.  The secant convergence is the definition of
the $L^2$ derivative.  Since $\psi_\theta$ is constructed from $a_\theta$
only, reference-gauge invariance is automatic.
\end{proof}

\begin{corollary}[Line-weighted Wilson wall]
\label{cor:V20-wall-Hellinger}
On a regular moving Wilson wall with active transverse line order $\beta>1$,
Version~V19 supplies the quadratic-mean hypothesis of
\Cref{thm:V20-active-Hellinger}.  The direct active secant therefore recovers
the complete finite Fisher/source metric, including every cancellation which
would appear as a control/line mixed block in a chosen reference gauge.  For
$\beta\le1$, the failure of the secant to remain bounded is the already proved
wall-source obstruction, not a reason to choose a different control gauge.
\end{corollary}

% =============================================================================
\section{Gauge-invariant cutoff transport of the active cylinder}
\label{sec:V20-active-transport}
% =============================================================================

Let $\pi_n:X_{n+1}\to X_n$ be the physical old-record map and let
$\Lambda_n^{\rm abs}$ be the intrinsic active positive measure at cutoff $n$.
Define
\begin{equation}
 \boxed{
 \delta_n^{\rm abs}
 =\norm{
   \Lambda_n^{\rm abs}
   -(\pi_n)_*\Lambda_{n+1}^{\rm abs}
  }_{\rm var}.}
 \label{eq:V20-active-transport-defect}
\end{equation}

\begin{theorem}[Intrinsic active-measure transport]
\label{thm:V20-active-transport}
For $N>m$, let $\pi_{N/m}$ be the composed cutoff map.  Then
\begin{equation}
 \boxed{
 \norm{
  \Lambda_m^{\rm abs}
  -(\pi_{N/m})_*\Lambda_N^{\rm abs}
 }_{\rm var}
 \le
 \sum_{n=m}^{N-1}\delta_n^{\rm abs}.}
 \label{eq:V20-active-transport-telescope}
\end{equation}
If the right tail is summable, the active measures have one unique
fixed-screen total-variation limit.  If in addition their masses satisfy
$Z_n^{\rm abs}\ge z_*>0$, then the normalized active laws obey
\begin{equation}
 \boxed{
 \norm{
  \mu_m^{\rm abs}
  -(\pi_{N/m})_*\mu_N^{\rm abs}
 }_{\rm var}
 \le
 \frac{2}{z_*}
 \sum_{n=m}^{N-1}\delta_n^{\rm abs}.}
 \label{eq:V20-normalized-active-transport}
\end{equation}
Whenever control coordinates are used and the product incidence is exact,
\begin{equation}
 \Lambda_n^{\rm abs}=W_n\Lambda_n^+,
 \label{eq:V20-product-coordinate}
\end{equation}
the defects \eqref{eq:V20-active-transport-defect} are unchanged by independent
positive-reference gauges at every cutoff.
\end{theorem}

\begin{proof}
Pushforward is a contraction in total variation.  Insert the adjacent measures
one at a time and telescope to obtain
\eqref{eq:V20-active-transport-telescope}.  For two positive finite measures
$\alpha,\beta$ with masses $A,B\ge z_*$ and
$\delta=\norm{\alpha-\beta}_{\rm var}$,
\[
 \norm{\alpha/A-\beta/B}_{\rm var}
 \le \delta/A+|A-B|/A
 \le 2\delta/z_*.
\]
Apply this after transport to obtain
\eqref{eq:V20-normalized-active-transport}.  Finally, the gauge transformation
changes $(\Lambda_n^+,W_n)$ to $(h_n\Lambda_n^+,W_n/h_n)$, so their product and
therefore the direct defect remain unchanged.
\end{proof}

\begin{assessment}[Revised transport order]
\label{ass:V20-transport-order}
Version~V19's separate control-law and normalized-ratio rectangle remains a
valid sufficient coordinate estimate when those two objects are physically
required.  For transport of the active positive law itself, the minimal target
is the direct product defect \eqref{eq:V20-active-transport-defect}.  Separate
control and ratio convergence is not charged unless the control sampler or the
ratio writer is independently part of the target.
\end{assessment}

% =============================================================================
\section{Exact finite controls and executable acquisition}
\label{sec:V20-controls}
% =============================================================================

\subsection{Two positive-reference gauges of one active cylinder}

On a two-point base, take
\begin{equation}
 a=(1,1),
 \qquad
 c=(2,1/2),
 \qquad
 w=(1/2,2).
 \label{eq:V20-exact-gauge-one}
\end{equation}
Then $cw=a$.  With $h=(3,4)$,
\begin{equation}
 c^h=(6,2),
 \qquad
 w^h=(1/6,1/2),
 \label{eq:V20-exact-gauge-two}
\end{equation}
and again $c^hw^h=a$.  Moreover,
\begin{equation}
 Z_c=5/2,
 \qquad
 Z_a=2,
 \qquad
 \boxed{\mathbb E_{\mu_c}w=4/5.}
 \label{eq:V20-exact-gauge-normalization}
\end{equation}
The normalized active law is uniform in both gauges.

\subsection{An exact singular-reference packet}

Take active and control weights
\begin{equation}
 a=(1,2,3),
 \qquad
 c=(1,0,1).
 \label{eq:V20-exact-singular-data}
\end{equation}
The exact Radon--Nikodym derivative on the control support is $(1,3)$ and the
singular active mass is two.  Hence
\begin{equation}
 \boxed{
 \inf_{v\ge0}\norm{a-vc}_{\ell^1}=2.}
 \label{eq:V20-exact-singular-residual}
\end{equation}
A mutated ratio giving represented weights $(1,0,2)$ has residual three.

Suppose a target writer has values one and four on the two dominated atoms and
is only known to lie in $[-2,5]$ on the singular atom.  Since $Z_a=6$, the
sharp normalized interval is
\begin{equation}
 \boxed{
 \frac32
 \le\mathbb E_{\mu_a}F
 \le\frac{23}{6}.}
 \label{eq:V20-exact-singular-interval}
\end{equation}

\subsection{Finite executable card}

\begin{acquisition}[Intrinsic active-line cylinder card]
\label{acq:V20-active-line-card}
At one selected RPESM cutoff, acquire in this order:
\begin{enumerate}[label=\textup{(A20.\arabic*)},leftmargin=5.5em]
\item the common bosonic base measure, complete active determinant and Pfaffian
      norm writers, and the direct unnormalized active density $a$;
\item either a direct active-cylinder Read, or one physically occurring control
      density $c$ and ratio $w$ together with an outward interval for
      $\mathbb C_{\rm line}^{\rm occ}$ in
      \eqref{eq:V20-line-occurrence-residual};
\item the singular active mass relative to the proposed control and one
      held-out bounded active writer;
\item the determinant-line atlas, transition functions, active phase section,
      and real Pfaffian orientation; a global control section is used only after
      its nonvanishing and line incidence are certified;
\item on a differentiable branch, the direct active Hellinger secant
      \eqref{eq:V20-Hellinger-secant}, or all reference-gauge score blocks with
      the held-out cancellation identity \eqref{eq:V20-block-sum};
\item the Wilson wall order and velocity rows of Version~V19 only after the
      active cylinder and source metric pass;
\item one adjacent-cutoff direct active product and its line-atlas transport.
\end{enumerate}
\end{acquisition}

\begin{theorem}[Terminating Version V20 card semantics]
\label{thm:V20-card-semantics}
After the common-history and positivity checks pass, the card of
\Cref{acq:V20-active-line-card} has exactly one of the following outcomes.
\begin{enumerate}[label=\textup{(C20.\arabic*)},leftmargin=5.5em]
\item \passstatus: the direct active positive cylinder is finite and nonzero,
      and either it is directly acquired or the submitted reference product has
      zero occurrence residual and zero singular mass.
\item \obsstatus: a strictly positive lower bound is certified for the minimum
      residual \eqref{eq:V20-min-reference-residual}, for a held-out product
      mismatch, or for a line-atlas/orientation inconsistency.
\item \stopstatus: all domain checks pass, but a required interval meets zero,
      the active normalization is not bounded away from zero, or the line atlas
      is incomplete.
\end{enumerate}
A nonpositive submitted density, an ill-typed line comparison, or a ratio used
outside the support of its control is an input rejection rather than a fourth
mathematical outcome.
\end{theorem}

\begin{proof}
The first two outcomes are the zero and positive branches of
\Cref{thm:V20-singular-reference,def:V20-line-occurrence} together with the line
incidence theorem \Cref{thm:V20-line-triviality}.  Every remaining valid
outward interval fails to establish a strict branch and is therefore the stated
stop.  The final cases violate the domains of the positive-measure or line
comparison theorems.
\end{proof}

% =============================================================================
\section{Version V20 theorem, ledger, and revised next acquisition}
\label{sec:V20-terminal}
% =============================================================================

\begin{theorem}[Version V20 intrinsic active-cylinder theorem]
\label{thm:V20-terminal}
Versions~V1--V19 are retained.  Version~V20 additionally proves:
\begin{enumerate}[label=\textup{(E20.\arabic*)},leftmargin=5.5em]
\item positive control cylinders and line ratios form a reference-gauge orbit
      whose exact quotient is the intrinsic active positive measure;
\item the normalized active law is invariant under that gauge and satisfies the
      exact density formula \eqref{eq:V20-reference-RN};
\item the assembled active score and Fisher/source metric are invariant, while
      the separate control, line, and mixed covariance blocks are not;
\item the separate diagonal metrics may be arbitrarily large even when the
      active source metric is zero;
\item the minimum variation-norm mismatch over all ratios equals the active
      mass singular to the proposed control measure;
\item unresolved singular mass gives the sharp bounded-writer interval
      \eqref{eq:V20-singular-payoff-interval};
\item a nowhere-zero global control section is equivalent to line triviality,
      while a global positive sampling density exists from local frames without
      supplying phase or orientation;
\item on a differentiable branch the direct active square-root secant converges
      to the complete physical source metric without a control decomposition;
\item direct active-measure total-variation defects give the target-minimal
      adjacent and cofinal transport theorem;
\item the finite card terminates with a passing active cylinder, a singular or
      incidence witness, or an unresolved interval on the same acquisition.
\end{enumerate}
No item supplies the active RPESM determinant/Pfaffian norm cylinder, its line
atlas, its Wilson-wall incidence, or its reflected crossing at one selected
cutoff.
\end{theorem}

\begin{proof}
Items~\textup{(E20.1)--(E20.4)} are
\Cref{thm:V20-reference-gauge,thm:V20-score-gauge,cth:V20-separate-metrics-arbitrary}.
Items~\textup{(E20.5)--(E20.6)} are
\Cref{thm:V20-singular-reference,cor:V20-singular-payoff}.  Item
\textup{(E20.7)} is
\Cref{thm:V20-line-triviality,prop:V20-local-control-density}.  Items
\textup{(E20.8)--(E20.9)} are
\Cref{thm:V20-active-Hellinger,thm:V20-active-transport}, and item
\textup{(E20.10)} is \Cref{thm:V20-card-semantics}.  The final limitations name
physical rows not evaluated by these structural identities.
\end{proof}

\begingroup\small
\begin{longtable}{>{\raggedright\arraybackslash}p{0.27\textwidth}
                  >{\raggedright\arraybackslash}p{0.20\textwidth}
                  >{\raggedright\arraybackslash}p{0.45\textwidth}}
\caption{Version V20 proof-status ledger.}\label{tab:V20-ledger}\\
\toprule
\textbf{Row}&\textbf{Status}&\textbf{Exact conclusion}\\
\midrule
\endfirsthead
\toprule
\textbf{Row}&\textbf{Status}&\textbf{Exact conclusion}\\
\midrule
\endhead
Positive-reference gauge
&\passstatus
&Every control/ratio factorization of one active density differs by one positive
multiplicative gauge; the active product is the quotient invariant.\\[0.3em]
Control and line score blocks
&Gauge dependent
&Their four-block sum is the physical active metric; the separate blocks can be
arbitrarily large under a pure reference change.\\[0.3em]
Intrinsic active source metric
&\passstatus
&Direct square-root/Hellinger tangents recover it without choosing a control
factorization.\\[0.3em]
Singular control support
&Exact finite alternative
&The minimum ratio mismatch is precisely the active singular mass; bounded
payoffs have the sharp interval \eqref{eq:V20-singular-payoff-interval}.\\[0.3em]
Global control section
&Topological occurrence row
&A nonzero section trivializes the line; local norm densities can sample a
nontrivial line but do not supply phase or Pfaffian orientation.\\[0.3em]
Line phase and real orientation
&Independent
&Require a common line atlas and orientation packet; positive norm data do not
select them.\\[0.3em]
Active cutoff transport
&\passstatus
&The direct positive-product defect is gauge invariant and gives a summable
fixed-screen variation-norm theorem.\\[0.3em]
Version V19 wall/Fisher compiler
&Retained downstream
&Run only after the intrinsic active cylinder and, on the differentiable branch,
its direct source metric have been acquired.\\[0.3em]
Selected RPESM active line cylinder
&\stopstatus
&No supplied finite card gives the direct active norm density, a zero product
residual, singular mass, global/local line atlas, and held-out active writer on
one selected cutoff.\\
\bottomrule
\end{longtable}
\endgroup

\begin{assessment}[Physical status after Version V20]
\label{ass:V20-status}
The nonvanishing control population remains a valid existential regulator
statement.  The new result is that its factorization into a control law and a
line ratio is not itself the physical source metric.  The physical verdict is
\begin{equation}
 \boxed{
 \begin{array}{ll}
 \mathsf{PASS}:&
 \begin{gathered}
 \text{reference-gauge quotient, intrinsic active score,}\\
 \text{singular-support and line-atlas compiler}
 \end{gathered}\\[0.35em]
 \mathsf{TYPED\ OBSTRUCTION}:&
 \begin{gathered}
 \text{pricing a reference-gauge block as an intrinsic source,}\\
 \text{or dropping active mass outside the control support}
 \end{gathered}\\[0.35em]
 \mathsf{STOP}:&
 \begin{gathered}
 \text{at the first direct active determinant/Pfaffian}\\
 \text{positive cylinder and line-atlas incidence card}
 \end{gathered}
 \end{array}}
 \label{eq:V20-status}
\end{equation}
This is upstream of the Version~V19 wall calculation and of all Wilson-sign,
reflected-hopping, Duhamel, current/stress, and continuum rows.
\end{assessment}

\begin{openproblem}[Version V21: populate the intrinsic active line cylinder]
\label{op:V20-next}
Preserve Versions~V1--V20.  On one selected RPESM cutoff, proceed in this order.
\begin{enumerate}[label=\textup{(V21.\arabic*)},leftmargin=5.2em]
\item acquire the complete active positive density
      $e^{-S_B}\norm{\sigma}^2|p|^2$ and certify that its partition lies in a
      strictly positive finite interval;
\item if a control sampler is used, acquire its density and the same-history
      line ratio together, evaluate
      $\mathbb C_{\rm line}^{\rm occ}$, and retain any singular active mass;
\item provide the determinant and Pfaffian line atlas, transition cocycles,
      active complex phase, real orientation, and the provenance of every
      claimed global nonzero control section;
\item retain one bounded active writer outside the ratio fit and one direct
      active-cylinder normalization check;
\item acquire the active neighbouring-cylinder Hellinger panel.  On a
      $\beta>1$ wall branch use it as the physical source metric; on a
      $\beta\le1$ branch return the Version~V19 wall-source obstruction;
\item only then evaluate the Wilson inertia masses, wall order and velocity,
      and the V13--V16 reflected crossing on that active law;
\item repeat the direct active product and line-atlas packet at one adjacent
      cutoff.
\end{enumerate}
If the active product is not populated, do not continue with another ratio,
gap, sign, collar, exterior, or Duhamel theorem.  The next result must be the
literal active positive cylinder or a strict singular/incidence witness against
the proposed control representation.
\end{openproblem}

\section*{V20 exact replay and append-only record}
\addcontentsline{toc}{section}{V20 exact replay and append-only record}
The accompanying standard-library verifier checks two distinct positive
references of one active cylinder, their exact product invariance and
normalization, a pure reference deformation with control and line metrics
$49,49,-49$ but zero active metric, the exact singular-reference residual two,
the larger mutated residual three, the sharp payoff interval
$[3/2,23/6]$, and the normalized active-measure transport bound.  Ordinary and
optimized executions produce byte-identical JSON output.  The replay is a
finite theorem certificate and does not populate an RPESM line cylinder.

The complete Version~V19 source preceding its sole final active
\verb|\end{document}| is preserved byte for byte.  Its preterminal SHA--256
digest is recorded in the validation manifest.  A later continuation must
remove only the final active document-closing command, append new material
immediately before it, restore one terminal command, and increment the filename
to end in \texttt{\_V21.tex}.

% =============================================================================
% END APPENDED VERSION V20 MATERIAL
% =============================================================================


% =============================================================================
% BEGIN APPENDED VERSION V21 MATERIAL
% =============================================================================

\clearpage
\hypersetup{
 pdftitle={Renewal Einstein--Standard--Model Physical Source Metric, Duhamel Ancestry and Quantum Continuum Programme V21},
 pdfsubject={Protected active determinant cylinder and phase-resolved line-source geometry}
}
\pdfinfo{
 /Title (Renewal Einstein--Standard--Model Physical Source Metric, Duhamel Ancestry and Quantum Continuum Programme V21)
 /Subject (Protected active determinant cylinder and phase-resolved line-source geometry)
}

\section{Version V21: protected active determinant cylinder and phase-resolved line-source geometry}
\label{sec:V21-active-line}
% =============================================================================

Version~V20 correctly required the intrinsic positive product before any
control-law ratio, Wilson-wall, or Duhamel calculation.  The source audit now
closes that stop on one literal finite branch.  On the protected compact
one-generation field shell, the same chiral coefficient which defines the
Berezin section also defines the positive determinant cylinder
\begin{equation}
 \boxed{
 \sigma=\operatorname{Det}D_\chi,
 \qquad
 \dd\Lambda_F=e^{-S_B}\norm{\sigma}^{2}\,\dd\lambda.}
 \label{eq:V21-det-literal-cylinder}
\end{equation}
The protected chiral singular-value floor makes this density everywhere
positive on the selected chart.  Thus no auxiliary massive control section is
needed to define its finite active norm-square experiment.

This landing is deliberately branch specific.  The full unbounded RPESM
population uses separately retained complex and real line families and a
nonvanishing massive control cylinder.  The compact determinant branch below
does not populate its real Pfaffian orientation, its active mass outside the
protected chart, or its numerical gauge-history crossing matrix.

\subsection{Finite partition and support theorem}

Let \(X\) be a compact field chart, let \(\lambda\) be a finite positive Borel
measure of full support, and put \(V_\lambda=\lambda(X)\).  Let \(E,F\to X\)
be Hermitian bundles of the same rank \(r\), and let
\(D:E\to F\) be continuous.  Its determinant section is
\begin{equation}
 \sigma_D=\operatorname{Det}D
 \in\Gamma\!\left(\det F\otimes(\det E)^*\right).
 \label{eq:V21-det-section}
\end{equation}

\begin{theorem}[Protected active determinant cylinder]
\label{thm:V21-det-cylinder}
Assume
\begin{equation}
 \tau I\preceq |D(x)|\preceq MI
 \quad(x\in X),
 \qquad 0<\tau\le M<\infty,
 \label{eq:V21-det-window}
\end{equation}
and
\begin{equation}
 S_-\le S_B(x)\le S_+
 \quad(x\in X).
 \label{eq:V21-det-action-window}
\end{equation}
For
\begin{equation}
 a_F=e^{-S_B}\norm{\sigma_D}^{2},
 \qquad
 Z_F=\int_Xa_F\,\dd\lambda,
 \qquad
 \dd\mu_F=Z_F^{-1}a_F\,\dd\lambda,
 \label{eq:V21-det-active-law}
\end{equation}
one has
\begin{equation}
 \boxed{
 V_\lambda e^{-S_+}\tau^{2r}
 \le Z_F\le
 V_\lambda e^{-S_-}M^{2r}.}
 \label{eq:V21-det-partition-bounds}
\end{equation}
In particular \(0<Z_F<\infty\), \(a_F\) is continuous and strictly positive,
and
\begin{equation}
 \boxed{\supp\mu_F=X.}
 \label{eq:V21-det-full-support}
\end{equation}
The determinant section is nowhere zero and therefore trivializes
\(\det F\otimes(\det E)^*\) on this chart.
\end{theorem}

\begin{proof}
If \(s_1(x),\ldots,s_r(x)\) are the singular values of \(D(x)\), the induced
line metric gives
\[
 \norm{\sigma_D(x)}^2
 =\det(D(x)^*D(x))
 =\prod_{j=1}^rs_j(x)^2.
\]
Hence \(\tau^{2r}\le\norm{\sigma_D}^2\le M^{2r}\).  Multiplication by the
action bounds and integration prove \eqref{eq:V21-det-partition-bounds}.  A
strictly positive continuous density over a full-support measure has support
all of \(X\).  Finally a complex line bundle carrying a continuous nowhere-zero
section is trivial.
\end{proof}

Define the bosonic law on the same chart by
\begin{equation}
 Z_B=\int_Xe^{-S_B}\,\dd\lambda,
 \qquad
 \dd\nu_B=Z_B^{-1}e^{-S_B}\,\dd\lambda.
 \label{eq:V21-det-bosonic-law}
\end{equation}

\begin{theorem}[Bounded equivalence and independence of the Wilson window]
\label{thm:V21-det-equivalence}
Put \(q=\norm{\sigma_D}^2\), \(q_-=\tau^{2r}\), and \(q_+=M^{2r}\).  Then
\begin{equation}
 \boxed{
 \frac{\dd\mu_F}{\dd\nu_B}
 =\frac{q}{\mathbb E_{\nu_B}q},
 \qquad
 \frac{q_-}{q_+}
 \le\frac{\dd\mu_F}{\dd\nu_B}
 \le\frac{q_+}{q_-}.}
 \label{eq:V21-det-RN}
\end{equation}
For every Borel set \(A\),
\begin{equation}
 \frac{q_-}{q_+}\nu_B(A)
 \le\mu_F(A)
 \le\frac{q_+}{q_-}\nu_B(A).
 \label{eq:V21-det-event-bounds}
\end{equation}
The sharp universal total-variation estimate is
\begin{equation}
 \boxed{
 d_{\rm TV}(\mu_F,\nu_B)
 \le
 \frac{\sqrt{q_+}-\sqrt{q_-}}
      {\sqrt{q_+}+\sqrt{q_-}}.}
 \label{eq:V21-det-TV}
\end{equation}
The two measures have the same null sets.  Consequently, for every nonnegative
measurable Wilson-gap writer \(\delta_W\),
\begin{equation}
 \boxed{
 \operatorname*{ess\,inf}_{\mu_F}\delta_W
 =\operatorname*{ess\,inf}_{\nu_B}\delta_W.}
 \label{eq:V21-det-gap-independence}
\end{equation}
A determinant weight bounded above and below cannot manufacture a hard Wilson
window absent from the bosonic law.
\end{theorem}

\begin{proof}
The density identity follows from
\(Z_F=Z_B\mathbb E_{\nu_B}q\); the pointwise and eventwise bounds follow from
\(q_-\le q\le q_+\).  To prove the sharp total-variation estimate, write
\(m=q_-\), \(M=q_+\), and \(\bar q=\mathbb E q\).  Convexity of
\(x\mapsto|x-\bar q|\) on \([m,M]\) gives
\[
 \mathbb E|q-\bar q|
 \le\frac{2(M-\bar q)(\bar q-m)}{M-m}.
\]
Therefore
\[
 d_{\rm TV}(\mu_F,\nu_B)
 \le\frac{(M-\bar q)(\bar q-m)}{\bar q(M-m)}.
\]
The right side is maximized at \(\bar q=\sqrt{mM}\), where it equals
\((\sqrt M-\sqrt m)/(\sqrt M+\sqrt m)\); a two-point endpoint law attains the
bound.  The Radon--Nikodym derivative and its reciprocal are strictly positive,
so the null sets, and hence all essential infima, agree.
\end{proof}

\subsection{Gauge descent and the intrinsic determinant score}

\begin{theorem}[Positive-density descent versus section descent]
\label{thm:V21-det-gauge}
Suppose a compact gauge group acts on \(X\), the bosonic action and base measure
are invariant, and
\begin{equation}
 D(gx)=U_F(g,x)D(x)U_E(g,x)^{-1}
 \label{eq:V21-det-gauge-covariance}
\end{equation}
with unitary source and target representations.  Then
\begin{equation}
 \boxed{
 \norm{\sigma_D(gx)}^2=\norm{\sigma_D(x)}^2,}
 \label{eq:V21-det-norm-descent}
\end{equation}
so the positive cylinder descends to the gauge quotient.  The section itself
obeys
\begin{equation}
 \boxed{
 \sigma_D(gx)
 =\det U_F(g,x)\det U_E(g,x)^{-1}\sigma_D(x).}
 \label{eq:V21-det-cocycle}
\end{equation}
It descends as an ordinary section only when this determinant cocycle is
trivial.  A physical antiunitary reflection preserves the norm density whenever
it carries \(D\) to the reflected coefficient and preserves \(S_B,\lambda\).
\end{theorem}

\begin{proof}
Equation \eqref{eq:V21-det-gauge-covariance} makes \(D(gx)^*D(gx)\) unitarily
conjugate to \(D(x)^*D(x)\), proving the norm identity.  Taking top exterior
powers gives \eqref{eq:V21-det-cocycle}.  Antiunitary conjugation likewise
preserves the determinant norm.
\end{proof}

Let \(\theta\mapsto(S_{B,\theta},D_\theta)\) be differentiable on one fixed
protected chart and use predeclared unitary bundle connections to compare the
neighbouring source and target fibres.

\begin{theorem}[Covariant determinant writer and direct Hellinger metric]
\label{thm:V21-det-score}
At \(\theta=0\), the unnormalized logarithmic writer in command direction
\(c\) is
\begin{equation}
 \boxed{
 h_c=-D_cS_B
 +2\operatorname{Re}\Tr\!\left(D^{-1}\nabla_cD\right).}
 \label{eq:V21-det-log-writer}
\end{equation}
The normalized score and physical source metric are
\begin{equation}
 \boxed{
 \ell_c=h_c-\mathbb E_{\mu_F}h_c,
 \qquad
 \mathsf G^{\rm abs}(c,d)
 =\operatorname{Cov}_{\mu_F}(h_c,h_d).}
 \label{eq:V21-det-source-metric}
\end{equation}
If \(\psi_\theta=(a_{F,\theta}/Z_{F,\theta})^{1/2}\), then
\begin{equation}
 \boxed{
 2D_c\psi_0=\ell_c\psi_0,
 \qquad
 4\ip{D_c\psi_0}{D_d\psi_0}=\mathsf G^{\rm abs}(c,d).}
 \label{eq:V21-det-root-metric}
\end{equation}
For two active determinant cylinders on the same base chart,
\begin{equation}
 \boxed{
 \mathsf H^2(\mu_0,\mu_1)
 =2-
 \frac{2}{\sqrt{Z_0Z_1}}
 \int_Xe^{-(S_0+S_1)/2}
 \norm{\sigma_0}\norm{\sigma_1}\,\dd\lambda.}
 \label{eq:V21-det-Hellinger}
\end{equation}
Consequently
\begin{equation}
 \boxed{
 \lim_{t\to0}
 \frac{\mathsf H^2(\mu_{tc},\mu_0)}{t^2}
 =\frac14\mathsf G^{\rm abs}(c,c).}
 \label{eq:V21-det-Hellinger-limit}
\end{equation}
All formulas are invariant under unitary frame changes and under the positive
reference gauge of Version~V20.
\end{theorem}

\begin{proof}
The line metric gives \(\norm{\sigma_D}^2=\det(D^*D)\).  Jacobi's identity and
\((D^*D)^{-1}D^*=D^{-1}\) yield
\[
 D_c\log\norm{\sigma_D}^2
 =2\operatorname{Re}\Tr(D^{-1}\nabla_cD).
\]
Unitary connection terms have purely imaginary determinant trace and disappear
under the real part.  Normalizing the density subtracts the mean, which proves
the score.  Differentiation of the square-root density gives
\eqref{eq:V21-det-root-metric}.  Expanding the squared distance between the two
normalized roots gives \eqref{eq:V21-det-Hellinger}; the tangent limit follows
from \eqref{eq:V21-det-root-metric}.
\end{proof}

\begin{firewall}[Fixed support and constant rank]
\label{fire:V21-det-fixed-support}
The determinant score theorem is a fixed-chart, constant-rank statement.  A
command moving a hard Wilson-index wall or crossing a determinant divisor is
routed to the Version~V18--V19 wall and zero-mode packets.  The bulk logarithmic
writer is not used to erase a support derivative.
\end{firewall}

\subsection{Literal finite-shell landing and the remaining real line}

\begin{theorem}[Protected finite field-shell active-law landing]
\label{thm:V21-det-field-shell-landing}
On the inherited compact field-shell branch with the same-operator chiral
coefficient, chiral floor, gauge-invariant bosonic action, and anomaly-free
one-generation determinant character, the density
\eqref{eq:V21-det-literal-cylinder} satisfies the hypotheses of
\Cref{thm:V21-det-cylinder,thm:V21-det-gauge,thm:V21-det-score}.  Therefore:
\begin{enumerate}[label=\textup{(\roman*)},leftmargin=2.8em]
\item the direct active determinant cylinder has a finite strictly positive
      partition and full support on the protected chart;
\item the determinant section and its positive norm descend to the protected
      gauge quotient;
\item the direct neighbouring-cylinder Hellinger panel is the physical
      determinant source metric;
\item the separately assumed Wilson window is a genuine hard active-law window
      on that chart, not an effect of determinant reweighting.
\end{enumerate}
The theorem does not populate the full unbounded RPESM active line cylinder or
an independently real Pfaffian line.
\end{theorem}

\begin{proof}
The inherited chiral floor supplies \(\tau>0\); continuity on the compact chart
supplies finite \(M,S_-,S_+\) and base volume.  The same-operator Berezin
section is exactly \(\operatorname{Det}D_\chi\), and the inherited field
cylinder is its norm-square density.  The determinant character of the complete
anomaly-free generation is one, so \eqref{eq:V21-det-cocycle} is trivial.  The
four conclusions now follow from the cited theorems and from
\eqref{eq:V21-det-gap-independence}.
\end{proof}

\begin{assessment}[The first Version V20 stop is discharged on one literal branch]
\label{ass:V21-det-landing}
The protected compact determinant field shell is now a populated intrinsic
positive-cylinder and source-metric branch.  The remaining line question is
not another reference-law ratio.  It is the complex phase/Berry packet of the
same determinant section, together with the independently real Pfaffian line
when a Majorana branch is claimed.  Those rows are treated next.
\end{assessment}

\subsection{Phase-resolved line-source geometry before quantum positive-law promotion}
\label{sec:V21-phase-resolved}

The preceding landing closes the intrinsic positive norm-square cylinder on
the protected determinant branch.  That positive law is necessary but not yet
sufficient for the complete quantum fermion source.  The determinant and
Pfaffian data retained by the finite regulator are line-valued and carry a
complex phase and a real orientation in addition to their positive norms.  The
general complex-quotient theorem imported from the Grand Tensor says that a
positive-law shadow of a complex measure exists only on the constant-phase
branch, and that the canonical positive modulus law does not replace the phase
packet.  The present subsection does not reprove that theorem.  It resolves its
infinitesimal and source-metric consequences for the SM--Einstein ancestry
programme.

The new order is
\begin{equation}
 \boxed{
 \begin{gathered}
 \text{positive norm-square cylinder and physical line transport}\\
 \Downarrow\\
 \text{modulus score plus centred phase velocity}\\
 \Downarrow\\
 \text{projective line-amplitude metric, Berry curvature, and phase short}\\
 \Downarrow\\
 \text{positive-shadow, phase-follower, new-phase-source, or partition-zero branch}\\
 \Downarrow\\
 \text{joint line/crossing gate, and only then Duhamel ancestry}.
 \end{gathered}}
 \label{eq:V21-order}
\end{equation}
A positive modulus cylinder is therefore called the \emph{phase-quenched
norm-square cylinder} below.  It is the complete active positive law only after
the phase and orientation rows have passed their own incidence tests.

\subsection{Local phase packet and normalized complex functional}

Let $\Theta\subset\R^d$ be a parameter neighbourhood of the origin, let
$(\Omega,\lambda)$ be a finite or standard probability carrier, and suppose
\begin{equation}
 d\mu_\theta=p_\theta\,d\lambda,
 \qquad p_\theta>0,
 \qquad \int p_\theta\,d\lambda=1.
 \label{eq:V21-positive-law}
\end{equation}
Let $u_\theta:\Omega\to\mathrm U(1)$ be the unit phase of the physically
transported determinant/Pfaffian line packet in one local unitary
trivialization.  Put
\begin{equation}
 m_\theta=\mathbb E_{\mu_\theta}u_\theta,
 \qquad
 d\zeta_\theta=Z_{{\rm abs},\theta}u_\theta\,d\mu_\theta.
 \label{eq:V21-complex-measure}
\end{equation}
Thus the unnormalized complex partition amplitude is
\begin{equation}
 Z_\theta=Z_{{\rm abs},\theta}m_\theta.
 \label{eq:V21-partition-factor}
\end{equation}
Whenever $m_\theta\ne0$, define the normalized complex functional
\begin{equation}
 \mathcal Q_\theta(F)
 =\frac{\mathbb E_{\mu_\theta}(u_\theta F)}{m_\theta}.
 \label{eq:V21-quantum-functional}
\end{equation}
It is a normalized complex functional, not a positive expectation in general.

Assume first that $u_0$ has constant phase.  Multiplication by one global phase
sets $u_0=1$ without changing \eqref{eq:V21-quantum-functional}.  For a command
$c\in\R^d$ define the positive score and phase velocity
\begin{equation}
 a_c=D_c\log p_\theta\big|_0,
 \qquad
 b_c=-\mathrm i\,\overline{u_0}D_cu_\theta\big|_0\in L^2(\mu_0;\R).
 \label{eq:V21-score-phase}
\end{equation}
Normalization gives $\mathbb E a_c=0$.  Write
\begin{equation}
 \widetilde b_c=b_c-\mathbb E b_c.
 \label{eq:V21-centred-phase}
\end{equation}
The phase velocity depends on the physical line transport.  A parameter-only
global phase changes $b_c$ by a constant and therefore leaves
$\widetilde b_c$ unchanged.  A history-dependent change of line frame must be
accompanied by the corresponding physical connection; otherwise the comparison
is ill typed.

\begin{theorem}[Exact complex response derivative]
\label{thm:V21-response-derivative}
Let $F_\theta$ be differentiable and square integrable.  At a constant-phase
base point,
\begin{equation}
 \boxed{
 D_c\mathcal Q_0(F)
 =\mathbb E\,D_cF_0
  +\operatorname{Cov}(F_0,a_c)
  +\mathrm i\operatorname{Cov}(F_0,b_c).}
 \label{eq:V21-response-derivative}
\end{equation}
More generally, at a point with $m\ne0$, put
$s_c=D_c\log p+u^{-1}D_cu$.  Then
\begin{equation}
 \boxed{
 D_c\mathcal Q(F)
 =\mathcal Q(D_cF)
  +\mathcal Q(Fs_c)-\mathcal Q(F)\mathcal Q(s_c).}
 \label{eq:V21-response-general}
\end{equation}
If $\mathfrak v=|m|$, then
\begin{equation}
 \boxed{
 |\mathcal Q(F)|\le\frac{\|F\|_{L^2(\mu)}}{\mathfrak v},}
 \label{eq:V21-Q-condition}
\end{equation}
and
\begin{equation}
 \boxed{
 |D_c\mathcal Q(F)-\mathcal Q(D_cF)|
 \le
 \left(\mathfrak v^{-1}+\mathfrak v^{-2}\right)
 \|F\|_2\|s_c\|_2.}
 \label{eq:V21-response-condition}
\end{equation}
Thus average-phase collapse is an intrinsic conditioning obstruction even when
the phase-quenched positive cylinder is perfectly regular.
\end{theorem}

\begin{proof}
Differentiate the numerator and denominator of
\eqref{eq:V21-quantum-functional}.  At the base point,
$D_cp=p a_c$, $D_cu=\mathrm i b_c$, and
$D_cm=\mathrm i\mathbb E b_c$.  The quotient rule gives
\eqref{eq:V21-response-derivative}.  The same quotient rule before setting
$u=1$ gives \eqref{eq:V21-response-general}.  Since $|u|=1$, Cauchy--Schwarz
gives $|\mathbb E(uF)|\le\|F\|_2$ and
$|\mathbb E(uFs_c)|\le\|F\|_2\|s_c\|_2$.  Apply these bounds separately to the
two covariance terms.
\end{proof}

\subsection{Average-phase Hessian and the missing positive source metric}

Define the phase Gram on the command space by
\begin{equation}
 \boxed{
 \mathsf G^{\rm ph}(c,d)
 =\operatorname{Cov}(b_c,b_d)
 =\mathbb E(\widetilde b_c\widetilde b_d).}
 \label{eq:V21-phase-Gram}
\end{equation}
It is positive semidefinite and invariant under parameter-dependent global
phase changes.

\begin{theorem}[Visibility Hessian equals the phase Gram]
\label{thm:V21-visibility-Hessian}
Assume $p_\theta$ and $u_\theta$ are twice differentiable and $u_0=1$.  Then
for all commands $c,d$,
\begin{equation}
 \boxed{
 -\frac12D_cD_d\log|m_\theta|^2\big|_0
 =\mathsf G^{\rm ph}(c,d).}
 \label{eq:V21-visibility-Hessian}
\end{equation}
The right side is independent of the modulus score, the phase acceleration,
and the chosen global phase gauge.
\end{theorem}

\begin{proof}
Normalization gives $D_cm=\mathrm i\mathbb E b_c$.  In the mixed second
derivative of $m$, the real part of the density second derivative vanishes by
twice differentiating $\int p_\theta=1$.  The density--phase cross terms and
the phase acceleration are imaginary.  The only real phase term is
$-\mathbb E(b_cb_d)$.  Hence
\[
 \operatorname{Re}D_cD_dm=-\mathbb E(b_cb_d).
\]
Since $m_0=1$,
\[
 D_cD_d\log|m|^2
 =2\operatorname{Re}\left(D_cD_dm-D_cmD_dm\right)
 =-2\operatorname{Cov}(b_c,b_d).
\]
This is \eqref{eq:V21-visibility-Hessian}.
\end{proof}

\begin{corollary}[First-order positive-shadow criterion]
\label{cor:V21-positive-shadow-tangent}
At a constant-phase base point, the following are equivalent.
\begin{enumerate}[label=\textup{(\roman*)},leftmargin=2.8em]
\item $\mathsf G^{\rm ph}=0$ on the retained command space;
\item every retained $b_c$ is constant $\mu_0$-almost everywhere;
\item for every real bounded writer $F$, the first derivative
      $D_c\mathcal Q_0(F)$ is real after the positive-law derivative is removed;
\item the normalized complex functional has, in every retained direction, the
      same first-order tangent as a positive probability law.
\end{enumerate}
A positive value of $\mathsf G^{\rm ph}$ is therefore an exact tangent-level
obstruction to promoting the phase-quenched cylinder to the full quantum
source experiment.  Exact positive-law promotion on a neighbourhood remains
the stronger constant-phase condition imported from the complex-quotient
theorem.
\end{corollary}

\begin{proof}
A covariance Gram vanishes exactly when every centred phase velocity vanishes.
Equation \eqref{eq:V21-response-derivative} then has no imaginary phase term,
and the positive law $\mu_\theta$ supplies the same first-order functional.
Conversely, a tangent of positive probability laws is real on every real
writer.  If the imaginary term in
\eqref{eq:V21-response-derivative} vanishes for every bounded real $F$, then
$b_c-\mathbb E b_c$ is orthogonal to every bounded function and hence is zero.
\end{proof}

\subsection{Projective line-amplitude metric and Berry incidence}

The phase-quenched density and the transported unit line section combine into
the normalized line-amplitude vector
\begin{equation}
 \Psi_\theta=p_\theta^{1/2}u_\theta
 \in L^2(\lambda),
 \qquad
 \|\Psi_\theta\|_2=1.
 \label{eq:V21-amplitude-vector}
\end{equation}
This vector is defined only after the physical line transport has identified
the neighbouring fibres.  Its projective class is independent of an overall
phase.

\begin{theorem}[Exact modulus--phase decomposition of the projective metric]
\label{thm:V21-projective-metric}
Let
\begin{equation}
 D_c^{\rm h}\Psi
 =D_c\Psi-\Psi\langle\Psi,D_c\Psi\rangle
 \label{eq:V21-horizontal-derivative}
\end{equation}
be the horizontal derivative at the constant-phase base point.  Then
\begin{equation}
 \boxed{
 D_c^{\rm h}\Psi_0
 =\left(\frac12a_c+\mathrm i\widetilde b_c\right)\Psi_0.}
 \label{eq:V21-horizontal-formula}
\end{equation}
Consequently
\begin{equation}
 \boxed{
 4\operatorname{Re}
 \langle D_c^{\rm h}\Psi_0,D_d^{\rm h}\Psi_0\rangle
 =\mathsf G^{\rm abs}(c,d)+4\mathsf G^{\rm ph}(c,d),}
 \label{eq:V21-projective-real}
\end{equation}
where
\begin{equation}
 \mathsf G^{\rm abs}(c,d)=\operatorname{Cov}(a_c,a_d)
 \label{eq:V21-modulus-metric}
\end{equation}
is the intrinsic positive-cylinder metric of Version~V20.  The Berry
connection and curvature are
\begin{equation}
 \boxed{
 \mathcal A_c=-\mathrm i\langle\Psi,D_c\Psi\rangle
 =\mathbb E b_c,}
 \label{eq:V21-Berry-connection}
\end{equation}
\begin{equation}
 \boxed{
 \Omega(c,d)
 =2\operatorname{Im}
 \langle D_c^{\rm h}\Psi,D_d^{\rm h}\Psi\rangle
 =\operatorname{Cov}(a_c,b_d)-\operatorname{Cov}(a_d,b_c).}
 \label{eq:V21-Berry-curvature}
\end{equation}
Thus the positive Hellinger metric sees only the modulus summand.  The phase
Gram and the mixed Berry curvature are independent physical rows.
\end{theorem}

\begin{proof}
Differentiate \eqref{eq:V21-amplitude-vector}:
\[
 D_c\Psi_0=\left(\frac12a_c+\mathrm i b_c\right)\Psi_0.
\]
Since $\mathbb E a_c=0$,
$\langle\Psi,D_c\Psi\rangle=\mathrm i\mathbb E b_c$, which proves
\eqref{eq:V21-horizontal-formula}.  Expanding the Hilbert inner product gives
\[
 \langle D_c^{\rm h}\Psi,D_d^{\rm h}\Psi\rangle
 =\frac14\operatorname{Cov}(a_c,a_d)
  +\operatorname{Cov}(b_c,b_d)
  +\frac{\mathrm i}{2}
   \left[\operatorname{Cov}(a_c,b_d)
        -\operatorname{Cov}(a_d,b_c)\right].
\]
This proves \eqref{eq:V21-projective-real}.  Differentiating
\eqref{eq:V21-Berry-connection} and using commutation of the local phase mixed
derivatives yields \eqref{eq:V21-Berry-curvature}; equivalently it is twice the
imaginary part of the horizontal Gram.
\end{proof}

\begin{firewall}[A line connection is not a positive-density convention]
\label{fire:V21-line-connection}
The covariance $\mathsf G^{\rm ph}$ is invariant under a global phase change,
but comparison of $u_\theta$ at distinct histories, commands, or cutoffs still
requires the physically retained determinant/Pfaffian line atlas and its
transport.  A history-dependent rephasing without the corresponding connection
changes the local phase velocities and is not a gauge of the populated card.
Likewise the real Pfaffian orientation is a discrete row and is not reconstructed
by \eqref{eq:V21-projective-real}.
\end{firewall}

\subsection{No double charging: the modulus-relative phase short}

Let
\begin{equation}
 A:E_{\rm abs}\longrightarrow L^2_0(\mu_0;\R),
 \qquad
 B:E_{\rm ph}\longrightarrow L^2_0(\mu_0;\R)
 \label{eq:V21-score-syntheses}
\end{equation}
be the syntheses of the retained centred modulus scores and centred phase
velocities.  Put
\begin{equation}
 G_A=A^*A,
 \qquad
 C=A^*B,
 \qquad
 G_B=B^*B.
 \label{eq:V21-complete-score-Gram}
\end{equation}

\begin{theorem}[Source-minimal phase innovation]
\label{thm:V21-phase-short}
The complete block Gram is positive and the modulus-relative phase residual is
\begin{equation}
 \boxed{
 R_{{\rm ph}\mid{\rm abs}}
 =G_B-C^*G_A^\dagger C
 =B^*(I-P_A)B\succeq0,}
 \label{eq:V21-phase-short}
\end{equation}
where $P_A=AG_A^\dagger A^*$ is the orthogonal projection onto $\operatorname{Ran}A$.
Moreover,
\begin{equation}
 \boxed{
 \rank R_{{\rm ph}\mid{\rm abs}}
 =\dim\operatorname{Ran}(A,B)-\dim\operatorname{Ran}A.}
 \label{eq:V21-phase-rank}
\end{equation}
Its polar synthesis is the unique source-minimal new real phase bank.  A
positive phase Gram need not be a new source: it may be a deterministic
follower of the already occurring modulus score.  Only
\eqref{eq:V21-phase-short} may enlarge the source atlas.
\end{theorem}

\begin{proof}
The block matrix is the Gram of $(A,B)$.  The Moore--Penrose projection identity
gives $C=A^*B$ and
\[
 C^*G_A^\dagger C=B^*P_AB.
\]
Subtracting yields \eqref{eq:V21-phase-short}.  The rank identity is the
dimension of the orthogonal component $(I-P_A)\operatorname{Ran}B$.  Polar
decomposition gives the minimum-dimensional isometric synthesis onto that
component.
\end{proof}

\begin{countertheorem}[The positive norm-square cylinder does not determine the quantum source]
\label{cth:V21-modulus-no-phase}
There are line packets with identical phase-quenched law, identical positive
source metric, and identical line modulus at every command, but inequivalent
normalized complex responses and different phase-source Grams.
\end{countertheorem}

\begin{proof}
Take $\Omega=\{-1,+1\}$ with the uniform law and no modulus deformation.  The
constant packet $u_t(x)=e^{\mathrm it}$ has
$\mathsf G^{\rm ph}=0$ and its normalized complex functional equals the
uniform positive expectation.  The packet
\[
 u_t(x)=e^{\mathrm itx}
\]
has the same positive law and modulus, but
\[
 m_t=\cos t,
 \qquad
 \mathsf G^{\rm ph}=1.
\]
For the real writer $F(x)=x$,
\begin{equation}
 \boxed{
 \mathcal Q_t(F)=\mathrm i\tan t.}
 \label{eq:V21-two-point-response}
\end{equation}
No positive probability law can reproduce this value for nonzero small $t$.
\end{proof}

\subsection{Finite direct acquisition without choosing a control gauge}

For unit vectors $\Phi,\Psi$ in one transported line-amplitude Hilbert space,
define the projective chordal distance
\begin{equation}
 d_{\rm pr}([\Phi],[\Psi])
 =\min_{|z|=1}\|\Phi-z\Psi\|_2
 =\sqrt{2(1-|\langle\Phi,\Psi\rangle|)}.
 \label{eq:V21-projective-distance}
\end{equation}

\begin{theorem}[Direct projective secant reconstructs the complete line-source metric]
\label{thm:V21-projective-secant}
For a command $c$, choose the phase $z_{t,c}$ which makes
$\langle\Psi_0,z_{t,c}\Psi_{tc}\rangle$ nonnegative real and put
\begin{equation}
 \Delta_{t,c}
 =\frac{2}{t}(z_{t,c}\Psi_{tc}-\Psi_0).
 \label{eq:V21-aligned-secant}
\end{equation}
Then
\begin{equation}
 \boxed{
 \Delta_{t,c}\longrightarrow
 (a_c+2\mathrm i\widetilde b_c)\Psi_0
 \quad\text{in }L^2,}
 \label{eq:V21-secant-limit}
\end{equation}
so that
\begin{equation}
 \boxed{
 \lim_{t\to0}\operatorname{Re}
 \langle\Delta_{t,c},\Delta_{t,d}\rangle
 =\mathsf G^{\rm abs}(c,d)+4\mathsf G^{\rm ph}(c,d),}
 \label{eq:V21-secant-real}
\end{equation}
\begin{equation}
 \boxed{
 \lim_{t\to0}\frac12\operatorname{Im}
 \langle\Delta_{t,c},\Delta_{t,d}\rangle
 =\Omega(c,d).}
 \label{eq:V21-secant-imag}
\end{equation}
In particular,
\begin{equation}
 \boxed{
 \lim_{t\to0}\frac{4}{t^2}
 d_{\rm pr}([\Psi_{tc}],[\Psi_0])^2
 =\mathsf G^{\rm abs}(c,c)+4\mathsf G^{\rm ph}(c,c).}
 \label{eq:V21-projective-speed}
\end{equation}
This finite source panel is invariant under the Version~V20 positive-reference
gauge and under global line phase.  It is the direct phase-sensitive analogue
of the positive Hellinger secant.
\end{theorem}

\begin{proof}
The derivative of the aligning phase is
$-\mathbb E b_c$, so differentiating the aligned vector gives
$(a_c/2+\mathrm i\widetilde b_c)\Psi_0$.  Multiplication by two proves
\eqref{eq:V21-secant-limit}.  The real and imaginary Gram limits follow by
expansion, and the diagonal identity is the definition of the minimizing
projective distance.
\end{proof}

The phase Gram may also be acquired from scalar average-phase data alone.
For a fixed command $c$ define
\begin{equation}
 \mathcal S_c(t)=-\log|m_{tc}|^2.
 \label{eq:V21-visibility-action}
\end{equation}

\begin{theorem}[Validated visibility secant]
\label{thm:V21-visibility-secant}
Assume $|m_{tc}|>0$ for $|t|\le\varepsilon$ and
\begin{equation}
 M_{4,c}(\varepsilon)
 =\sup_{|t|\le\varepsilon}|\mathcal S_c^{(4)}(t)|<\infty.
 \label{eq:V21-visibility-fourth}
\end{equation}
Then
\begin{equation}
 \boxed{
 Q_{\varepsilon}^{\rm ph}(c)
 :=\frac{\mathcal S_c(\varepsilon)+\mathcal S_c(-\varepsilon)
            -2\mathcal S_c(0)}{2\varepsilon^2}}
 \label{eq:V21-phase-secant}
\end{equation}
satisfies
\begin{equation}
 \boxed{
 \left|Q_{\varepsilon}^{\rm ph}(c)
       -\mathsf G^{\rm ph}(c,c)\right|
 \le\frac{\varepsilon^2}{24}M_{4,c}(\varepsilon).}
 \label{eq:V21-phase-secant-error}
\end{equation}
Mixed entries follow from polarization of the commands $c+d$ and $c-d$.  An
outward interval for $|m|$ which meets zero is a partition-zero or
insufficient-visibility \stopstatus, not a negative phase metric.
\end{theorem}

\begin{proof}
By \Cref{thm:V21-visibility-Hessian},
$\mathcal S_c''(0)=2\mathsf G^{\rm ph}(c,c)$.  The symmetric second-difference
remainder is bounded by
$\varepsilon^2M_{4,c}/12$.  Division by two gives
\eqref{eq:V21-phase-secant-error}.
\end{proof}

\begin{proposition}[Outward normalized-amplitude error]
\label{prop:V21-amplitude-error}
Let
\begin{equation}
 n=\mathbb E_\mu(uF),
 \qquad m=\mathbb E_\mu u,
 \qquad |m|\ge v_*>0,
 \label{eq:V21-num-den}
\end{equation}
and suppose $|\widehat n-n|\le\eta_n$,
$|\widehat m-m|\le\eta_m<v_*$.  Then
\begin{equation}
 \boxed{
 \left|\frac{\widehat n}{\widehat m}-\frac{n}{m}\right|
 \le
 \frac{\eta_n}{v_*-\eta_m}
 +\frac{|n|\eta_m}{v_*(v_*-\eta_m)}.}
 \label{eq:V21-amplitude-error}
\end{equation}
One may replace $|n|$ by $\|F\|_{L^2(\mu)}$.  Thus every normalized complex
writer requires a visibility floor; the unnormalized complex measure remains
the zero-safe output when that floor fails.
\end{proposition}

\begin{proof}
Insert and subtract $n/\widehat m$, use
$|\widehat m|\ge v_*-\eta_m$, and apply the triangle inequality.
\end{proof}

\subsection{Exact finite controls for curvature and phase ancestry}

\begin{proposition}[Two-parameter Berry-curvature control]
\label{prop:V21-Berry-control}
On $\Omega=\{-1,+1\}$ put
\begin{equation}
 \mu_s(x)=\frac{1+sx}{2},
 \qquad
 u_t(x)=e^{\mathrm itx}.
 \label{eq:V21-Berry-model}
\end{equation}
At $(s,t)=(0,0)$,
\begin{equation}
 \boxed{
 \mathsf G^{\rm abs}_{ss}=1,
 \qquad
 \mathsf G^{\rm ph}_{tt}=1,
 \qquad
 \Omega_{st}=1.}
 \label{eq:V21-Berry-values}
\end{equation}
The Berry connection has
$\mathcal A_t=s$ and $\mathcal A_s=0$, so its curvature is exactly one.  The
separate modulus and phase diagonal metrics do not determine this signed mixed
incidence.
\end{proposition}

\begin{proof}
The two tangent writers are $a_s=x$ and $b_t=x$.  Their variances are one.
Equation \eqref{eq:V21-Berry-curvature} gives
$\Omega_{st}=\operatorname{Cov}(x,x)=1$.  Directly,
$\mathcal A_t=\mathbb E_{\mu_s}x=s$.
\end{proof}

\begin{proposition}[Phase follower versus genuinely new phase source]
\label{prop:V21-phase-follower-control}
Let $(x,y)$ be independent uniform signs.  Use the modulus score bank
$A1=x$.
\begin{enumerate}[label=\textup{(\alph*)},leftmargin=2.8em]
\item For the phase velocity $B_0 1=x$,
      \begin{equation}
       \boxed{G_A=G_{B_0}=C=1,
       \qquad R_{{\rm ph}\mid{\rm abs}}=0.}
       \label{eq:V21-phase-follower}
      \end{equation}
\item For the phase velocity $B_1 1=y$,
      \begin{equation}
       \boxed{G_A=G_{B_1}=1,
       \qquad C=0,
       \qquad R_{{\rm ph}\mid{\rm abs}}=1.}
       \label{eq:V21-phase-new}
      \end{equation}
\end{enumerate}
The two phase Grams are identical, but only the second packet costs a new
source direction.
\end{proposition}

\begin{proof}
The statements are the covariances of independent centred signs and the Schur
formula \eqref{eq:V21-phase-short}.
\end{proof}

\subsection{Projective cutoff transport}

After transporting every finite card to one fixed physical line-amplitude
screen, write $[\Psi_n]$ for its projective unit vector and define
\begin{equation}
 \delta_n^{\rm pr}
 =d_{\rm pr}([\Psi_{n+1}],[\Psi_n]).
 \label{eq:V21-projective-defect}
\end{equation}

\begin{theorem}[Cofinal projective line landing]
\label{thm:V21-projective-transport}
If
\begin{equation}
 \sum_{n=1}^\infty\delta_n^{\rm pr}<\infty,
 \label{eq:V21-projective-sum}
\end{equation}
then there are phases $z_n\in\mathrm U(1)$ such that $z_n\Psi_n$ is Cauchy in
the common Hilbert screen.  Its nonzero limit is unique up to one global phase,
and the projective classes have a unique limit.

If, in addition, the aligned command secant Grams of
\Cref{thm:V21-projective-secant} have summable adjacent defects and their
shrinking-command remainders vanish, then the complete real projective metric
$\mathsf G^{\rm abs}+4\mathsf G^{\rm ph}$ and the Berry two-form $\Omega$ have
unique fixed-screen limits.

This theorem does not orient a real Pfaffian line and does not prove
path-independence around a cutoff/command loop.  A held-out line-holonomy
rectangle remains an independent physical row.
\end{theorem}

\begin{proof}
Choose $z_{n+1}$ recursively so that
\[
 \|z_{n+1}\Psi_{n+1}-z_n\Psi_n\|=\delta_n^{\rm pr}.
\]
The summable tail makes the aligned vectors Cauchy.  Their norms are one, so the
limit is nonzero.  Any second aligned limit differs by the limiting relative
phase.  The metric and curvature conclusion follows by the same telescoping
argument applied to the finite matrices, followed by the shrinking-command
limits.
\end{proof}

\subsection{Terminating phase-resolved acquisition card}

\begin{acquisition}[Phase-resolved active line-source card]
\label{acq:V21-card}
On one selected RPESM cutoff, acquire on one common bosonic history cylinder
\begin{equation}
 \boxed{
 \begin{aligned}
 \mathfrak P_{\rm line}=\bigl(
 &\Lambda_{\rm abs},\ \mu_{\rm abs},\
 \mathscr L_{\det},\mathscr L_{\rm Pf},\nabla_{\rm line},\ u,\ m,\\
 &\mathsf G^{\rm abs},\ \mathsf G^{\rm ph},\
 \mathsf C^{{\rm abs},{\rm ph}},\
 R_{{\rm ph}\mid{\rm abs}},\ \Omega,\\
 &\text{one direct projective secant panel, one held-out complex writer,}\\
 &\text{one visibility floor or partition-zero interval, and one adjacent card}
 \bigr).
 \end{aligned}}
 \label{eq:V21-card}
\end{equation}
The determinant phase and real Pfaffian orientation are retained separately.
The line transport is frozen before the phase velocities or secants are read.
\end{acquisition}

\begin{theorem}[Version V21 finite alternative]
\label{thm:V21-alternative}
After the common-history, line-atlas, transport, and outward-radius checks pass,
exactly one of the following branches is returned.
\begin{enumerate}[label=\textup{(V21.\arabic*)},leftmargin=5.2em]
\item \emph{Positive-shadow branch.}
      The phase Gram vanishes on the intended command bank and the held-out
      complex writer passes the constant-phase use class.  The Version~V20
      positive cylinder supplies the complete first-order source metric.
\item \emph{Phase-follower branch.}
      $\mathsf G^{\rm ph}\ne0$ but
      $R_{{\rm ph}\mid{\rm abs}}=0$.  The phase response and Berry curvature are
      retained, but no second source price is charged.
\item \emph{New phase-source branch.}
      $R_{{\rm ph}\mid{\rm abs}}\ne0$.  Its polar range is appended as the
      unique source-minimal phase bank before any Duhamel ancestry short.
\item \emph{Partition-zero or visibility-soft branch.}
      The amplitude interval meets zero or the visibility floor collapses.
      Normalized complex expectations are not promoted; the unnormalized
      complex measure, divisor, and phase-conditioning packet remain canonical.
\item \emph{Line-incidence obstruction.}
      The local phases, transition cocycles, physical connection, real
      orientation, or held-out complex writer do not belong to one common
      line-history packet.
\end{enumerate}
An outward interval which fails to separate zero is \stopstatus.  A failed line
or same-history identity is an input rejection, not a negative quantum mode.
\end{theorem}

\begin{proof}
The visibility theorem and
\Cref{cor:V21-positive-shadow-tangent} give the first and fourth branches.
The orthogonal short of \Cref{thm:V21-phase-short} gives the mutually exclusive
second and third branches.  The final branch records failure of the domain on
which those theorems are stated.  These alternatives exhaust the finite
phase-resolved packet.
\end{proof}

\begin{assessment}[Physical status after Version V21]
\label{ass:V21-status}
The protected determinant cylinder populated above is the canonical
phase-quenched norm-square experiment on that finite branch.  It is not called
the complete active quantum source law unless the constant-phase branch has
been certified.  The new physical verdict is
\begin{equation}
 \boxed{
 \begin{array}{ll}
 \mathsf{PASS}:&
 \begin{gathered}
 \text{protected active determinant cylinder and source metric,}\\
 \text{phase Gram, projective metric, Berry incidence, and phase short}
 \end{gathered}\\[0.45em]
 \mathsf{TYPED\ OBSTRUCTION}:&
 \begin{gathered}
 \text{promoting a modulus law to the quantum source without phase,}\\
 \text{or charging a modulus-generated phase direction twice}
 \end{gathered}\\[0.45em]
 \mathsf{STOP}:&
 \begin{gathered}
 \text{at the first populated determinant/Pfaffian phase and the}\\
 \text{same-history gauge crossing/held-out complex-writer card}
 \end{gathered}
 \end{array}}
 \label{eq:V21-status}
\end{equation}
No numerical RPESM phase visibility, phase Gram, Berry curvature, or
phase-relative source short is supplied by the attached corpus.  The protected
field-shell norm cylinder is populated, but its numerical gauge-history
crossing is not.  Consequently the joint line/crossing gate and the downstream
exterior, Duhamel, current/stress, charge, and Einstein rows remain closed.
\end{assessment}

\begin{openproblem}[Version V22: one joint protected-shell line and crossing card]
\label{op:V21-next}
Preserve Versions~V1--V21.  On one actual history of the populated protected
determinant cylinder and on one adjacent cutoff, acquire one joint packet:
\begin{equation}
 \boxed{
 \left(
 \sigma,\nabla_{\rm line},u,m,
 H_W,T_\pm,R_\Theta,\varepsilon_\Theta,
 \star_0,P_+,W_1
 \right).}
 \label{eq:V21-next-card}
\end{equation}
Proceed in this order.
\begin{enumerate}[label=\textup{(V22.\arabic*)},leftmargin=5.2em]
\item provide the determinant-line frames, transition cocycles, physical line
      transport, and one direct complex phase analyzer;
\item use two predeclared command displacements to populate the projective
      line-amplitude secant or the visibility secant with outward errors;
\item compute \(\mathsf G^{\rm ph}\), the modulus--phase mixed block,
      \(R_{{\rm ph}\mid{\rm abs}}\), and the Berry two-form;
\item evaluate \(T_+\operatorname{sgn}(H_W)T_-\) on the same history by exact
      spectral calculus, the Version~V16 boundary-cyclic panel, or a certified
      polynomial head and tail;
\item form the reflected crossing with the frozen reflection and Berezin sign,
      and run the Version~V9 relation, skew, and negative-part residuals;
\item retain one complex writer and one four-phase crossing probe outside the
      fit; if a real Majorana branch is claimed, populate its Pfaffian line and
      orientation separately;
\item repeat the phase and crossing packet at one adjacent cutoff and retain the
      cutoff/command holonomy rectangle.
\end{enumerate}
The phase and crossing calculations are parallel components of the same
terminal gate.  Neither is used to repair the other, and no exterior or Duhamel
compiler is opened until both physically required components pass.
\end{openproblem}

\section*{V21 exact replay and append-only record}
\addcontentsline{toc}{section}{V21 exact replay and append-only record}
The accompanying verifier first checks the exact two-history protected
determinant cylinder, its partition bounds, bosonic/active Radon--Nikodym
window, frame invariance, Jacobi score, and Hellinger--Fisher limit.  It then
checks the phase-only model $\mathcal Q_t(x)=\mathrm i\tan t$, the visibility
Hessian and projective speed, the global-phase null branch, the two-parameter
Berry curvature, and the phase-follower versus orthogonal new-phase-source
Schur alternatives.  Normal and optimized executions are required to produce
byte-identical JSON output.  The replay is a theorem certificate; it does not
populate the active RPESM phase/crossing card.

The complete Version~V20 source preceding its sole final active
\verb|\end{document}| is preserved byte for byte.  Its preterminal SHA--256
digest is recorded in the validation manifest.  A later continuation must
remove only the final active document-closing command, append new material
immediately before it, restore one terminal command, and increment the filename
to end in \texttt{\_V22.tex}.


% =============================================================================
% END APPENDED VERSION V21 MATERIAL
% =============================================================================



% =============================================================================
% BEGIN APPENDED VERSION V22 MATERIAL
% =============================================================================

\clearpage
\hypersetup{
 pdftitle={Renewal Einstein--Standard--Model Physical Source Metric, Duhamel Ancestry and Quantum Continuum Programme V22},
 pdfsubject={Chiral scalarization, determinant doubling, and phase-weighted reflected crossing}
}
\pdfinfo{
 /Title (Renewal Einstein--Standard--Model Physical Source Metric, Duhamel Ancestry and Quantum Continuum Programme V22)
 /Subject (Chiral scalarization, determinant doubling, and phase-weighted reflected crossing)
}


% The cumulative table of contents now has three-digit section numbers and
% double-digit subsection numbers.  Persist wider number and page boxes through
% the auxiliary file so a clean rebuild remains reproducible without changing
% the preserved Version V21 prefix.
\makeatletter
\immediate\write\@auxout{\string\global\string\def\string\@pnumwidth{3.2em}}
\immediate\write\@auxout{\string\global\string\def\string\@tocrmarg{4.2em}}
\immediate\write\@auxout{\string\global\string\def\string\l@subsection{\string\@dottedtocline{2}{1.5em}{3.2em}}}
\makeatother

\section{Version V22: chiral scalarization, determinant doubling, and phase-weighted reflected crossing}
\label{sec:V22-chiral-crossing}
% =============================================================================

Version~V21 populated the protected positive norm-square cylinder and derived
its phase-resolved half-density geometry.  Its final mandate treated the line
phase and the reflected crossing as parallel components of one gate.  Before
that packet is sampled, however, one further source-incidence distinction is
necessary.  For a single chiral Berezin factor the complex amplitude is linear
in the determinant section, whereas the populated field-shell density is
quadratic in its norm.  The latter is intrinsic and positive, but it is the
amplitude of an independently occurring conjugate double only when that second
factor is actually present.  Otherwise it is a control or norm-square
experiment.

The corrected entrance is therefore
\begin{equation}
 \boxed{
 \begin{gathered}
 \text{determinant line and physical scalarization, or an actual conjugate/real branch}
 \\
 \Downarrow\\
 \text{single-determinant polar law versus norm-square control law}
 \\
 \Downarrow\\
 \text{reflection sewing and the same-history phase--crossing measure}
 \\
 \Downarrow\\
 \text{record-relative vacuum and one-letter Osterwalder--Schrader gate}
 \\
 \Downarrow\\
 \text{only then exterior, Hamiltonian, and Duhamel promotion}.
 \end{gathered}}
 \label{eq:V22-corrected-order}
\end{equation}
This continuation preserves the Version~V21 projective line-amplitude
identities.  It determines which modulus power and which complex weight must be
inserted into them for the declared fermionic target.

\subsection{The Hermitian line metric does not scalarize a chiral amplitude}

Let $\mathscr L\to X$ be a Hermitian complex line and let
$\sigma\in\Gamma(\mathscr L)$ be a determinant section.  Write
\begin{equation}
 q(x)=\norm{\sigma(x)}.
 \label{eq:V22-line-modulus}
\end{equation}

\begin{lemma}[Riesz pairing is anti-linear and produces the norm square]
\label{lem:V22-Riesz-not-scalar}
The Hermitian metric gives a canonical anti-linear map
\begin{equation}
 \jmath_h:\mathscr L_x\longrightarrow\mathscr L_x^*,
 \qquad
 \jmath_h(v)(w)=\ip{v}{w},
 \label{eq:V22-Riesz-map}
\end{equation}
and
\begin{equation}
 \boxed{\jmath_h(\sigma)(\sigma)=q^2.}
 \label{eq:V22-Riesz-square}
\end{equation}
It does not give a nonzero complex-linear scalarization
$\mathscr L_x\to\C$ natural under every unitary automorphism of the line.
\end{lemma}

\begin{proof}
With the convention that the Hermitian product is conjugate-linear in its
first variable, \eqref{eq:V22-Riesz-map} is anti-linear in $v$ and linear in
$w$, and \eqref{eq:V22-Riesz-square} is the definition of the norm.  Suppose a
complex-linear functional $\ell$ were determined by the Hermitian line alone
and invariant under every unitary scalar $e^{\mathrm i\theta}$.  Naturality
would give
\[
 \ell(e^{\mathrm i\theta}v)=\ell(v),
\]
whereas linearity gives
$\ell(e^{\mathrm i\theta}v)=e^{\mathrm i\theta}\ell(v)$.  Choosing a
nontrivial phase forces $\ell(v)=0$ for every $v$.
\end{proof}

\begin{corollary}[Physical scalarization is an additional line-incidence row]
\label{cor:V22-physical-scalarization}
A scalar chiral amplitude requires an independently declared nonzero dual
section $\rho\in\Gamma(\mathscr L^*)$, a physical pairing with a conjugate
line, or a genuine real/Pfaffian structure.  On a unit dual branch put
\begin{equation}
 d=\rho(\sigma)=qu,
 \qquad |u|=1.
 \label{eq:V22-scalar-determinant}
\end{equation}
The complex scalar $d$ is invariant under simultaneous dual frame changes,
whereas the coefficient phase of $\sigma$ in an arbitrary local frame is not.
Using the normalized section $\sigma/q$ itself as an after-the-fact frame sets
$u=1$ but changes the physical connection and supplies no scalarization
certificate.
\end{corollary}

\begin{proof}
A dual section evaluates a line vector linearly.  Under a unitary frame change,
the section and dual frame transform inversely, so their pairing is unchanged.
The last assertion follows because $\sigma/q$ depends on the amplitude being
measured and is not an independently frozen physical line frame.
\end{proof}

\subsection{The modulus-power ladder and the exact single-versus-double distinction}

Assume for the moment that the scalarization in
\Cref{cor:V22-physical-scalarization} has been physically supplied and that
$q>0$ on one protected chart.  For $\alpha\ge0$ define
\begin{equation}
 \dd\Lambda_\alpha
 =e^{-S_B}q^\alpha\,\dd\lambda,
 \qquad
 Z_\alpha=\Lambda_\alpha(X),
 \qquad
 \dd\mu_\alpha=Z_\alpha^{-1}\dd\Lambda_\alpha.
 \label{eq:V22-modulus-power-law}
\end{equation}
The three relevant values are
\begin{equation}
 \boxed{
 \begin{array}{c|c}
 \alpha=0&\text{bosonic cylinder},\\
 \alpha=1&\text{polar positive law of one scalar chiral determinant},\\
 \alpha=2&\text{the populated norm-square cylinder}.
 \end{array}}
 \label{eq:V22-three-powers}
\end{equation}

\begin{theorem}[Exact modulus-power transport]
\label{thm:V22-power-ladder}
For any $\alpha,\gamma\ge0$,
\begin{equation}
 \boxed{
 \frac{\dd\mu_\alpha}{\dd\mu_\gamma}
 =\frac{q^{\alpha-\gamma}}
        {\mathbb E_{\mu_\gamma}q^{\alpha-\gamma}}.}
 \label{eq:V22-power-RN}
\end{equation}
If $q$ is bounded above and below by positive constants, all of these laws are
equivalent and have the same support.  Moreover, for
$\Phi(\alpha)=\log Z_\alpha$,
\begin{equation}
 \boxed{
 \Phi'(\alpha)=\mathbb E_{\mu_\alpha}\log q,
 \qquad
 \Phi''(\alpha)=\operatorname{Var}_{\mu_\alpha}(\log q)\ge0.}
 \label{eq:V22-power-convexity}
\end{equation}
For a differentiable physical command $c$ on a fixed full-rank chart, the
unnormalized logarithmic writer and the physical score metric are
\begin{equation}
 h_{\alpha,c}
 =-D_cS_B
  +\alpha\operatorname{Re}\Tr(D^{-1}\nabla_cD),
 \label{eq:V22-power-score}
\end{equation}
\begin{equation}
 \boxed{
 \mathsf G_\alpha(c,d)
 =\operatorname{Cov}_{\mu_\alpha}
   (h_{\alpha,c},h_{\alpha,d}).}
 \label{eq:V22-power-metric}
\end{equation}
Thus the $\alpha=2$ Hellinger metric does not, in general, equal or determine
the single-determinant $\alpha=1$ metric.
\end{theorem}

\begin{proof}
Equation \eqref{eq:V22-power-RN} is obtained by dividing the two normalized
densities.  Differentiation under the finite or dominated integral gives
\[
 Z_\alpha'=\int e^{-S_B}q^\alpha\log q\,\dd\lambda
\]
and a second differentiation gives the variance identity in
\eqref{eq:V22-power-convexity}.  Since
$D_c\log q=\operatorname{Re}\Tr(D^{-1}\nabla_cD)$, differentiation of
\eqref{eq:V22-modulus-power-law} and normalization give
\eqref{eq:V22-power-score}--\eqref{eq:V22-power-metric}.
\end{proof}

\begin{countertheorem}[The norm-square metric is not the one-determinant metric]
\label{cth:V22-power-metrics-differ}
There is a protected scalar determinant family for which
\begin{equation}
 \boxed{\mathsf G_1=1,\qquad \mathsf G_2=4.}
 \label{eq:V22-power-metric-example}
\end{equation}
\end{countertheorem}

\begin{proof}
Take the uniform two-point carrier $x\in\{-1,+1\}$, set $S_B=0$, and let
$D_\theta(x)=e^{\theta x}$.  At $\theta=0$, every $\mu_\alpha$ is uniform and
the centered score is $\alpha x$.  Its variance is $\alpha^2$.
\end{proof}

\begin{theorem}[Independent conjugate doubling identity]
\label{thm:V22-Grassmann-doubling}
Let $D:E\to F$ be a square finite chiral coefficient.  With one fixed Berezin
ordering,
\begin{equation}
 \int e^{-\bar\psi D\psi}\,
      \mathcal D(\bar\psi,\psi)
 =\det D.
 \label{eq:V22-single-Grassmann}
\end{equation}
For an independently occurring conjugate fermion pair with coefficient $D^*$,
\begin{equation}
 \boxed{
 \int e^{-\bar\psi D\psi-\bar\chi D^*\chi}\,
      \mathcal D(\bar\psi,\psi)\mathcal D(\bar\chi,\chi)
 =\det D\,\det D^*
 =\det(D^*D)
 =\norm{\operatorname{Det}D}^{2}.}
 \label{eq:V22-double-Grassmann}
\end{equation}
Therefore the norm-square cylinder is the literal fermionic modulus of a
doubled vectorlike branch only when the conjugate factor and its common action
actually occur.  The algebraic identity alone does not populate that branch.
\end{theorem}

\begin{proof}
The first identity is the finite Grassmann Gaussian formula.  Independence
factorizes the second integral into two Gaussian determinants.  Since
$\det D^*=\overline{\det D}$ and
$\det(D^*D)=\overline{\det D}\det D$, the remaining equalities follow.
Physical occurrence of the second Grassmann system is not a consequence of an
algebraic product identity.
\end{proof}

\begin{firewall}[Nambu rewriting is not conjugate occurrence]
\label{fire:V22-Nambu-not-double}
The skew Nambu matrix
\begin{equation}
 \mathcal A_D=
 \begin{pmatrix}0&D\\-D^{\mathsf T}&0\end{pmatrix}
 \label{eq:V22-Nambu-matrix}
\end{equation}
has
\begin{equation}
 \operatorname{Pf}\mathcal A_D
 =(-1)^{r(r-1)/2}\det D.
 \label{eq:V22-Nambu-Pfaffian}
\end{equation}
It rewrites one determinant; it does not produce
$\det D\det\overline D$.  Exact phase cancellation requires an independently
integrated conjugate factor or a genuine real/Pfaffian certificate.
\end{firewall}

\subsection{The single chiral complex measure relative to the norm-square control}

The linearly scalarized one-determinant measure is
\begin{equation}
 \dd\zeta_\chi=e^{-S_B}d\,\dd\lambda
 =e^{-S_B}qu\,\dd\lambda.
 \label{eq:V22-single-complex-measure}
\end{equation}
Let
\begin{equation}
 m_1=\mathbb E_{\mu_1}u,
 \qquad
 m_{2,-1}=\mathbb E_{\mu_2}(u/q).
 \label{eq:V22-two-partitions}
\end{equation}

\begin{theorem}[Exact one-determinant reweighting]
\label{thm:V22-single-reweighting}
Whenever the relevant denominator is nonzero, the normalized single-chiral
functional is
\begin{equation}
 \boxed{
 \mathcal Q_\chi(F)
 =\frac{\mathbb E_{\mu_1}(uF)}{\mathbb E_{\mu_1}u}
 =\frac{\mathbb E_{\mu_2}(u q^{-1}F)}
        {\mathbb E_{\mu_2}(u q^{-1})}.}
 \label{eq:V22-single-Q}
\end{equation}
By contrast,
\begin{equation}
 \frac{\mathbb E_{\mu_2}(uF)}{\mathbb E_{\mu_2}u}
 \label{eq:V22-square-phase-functional}
\end{equation}
is the normalized functional of the complex density
$e^{-S_B}q^2u$, not of \eqref{eq:V22-single-complex-measure}.  The two
functionals agree for every bounded $F$ only when $q$ is constant almost
everywhere on the nonzero complex support, apart from degenerate partition-zero
cases.
\end{theorem}

\begin{proof}
Since
\[
 \dd\mu_1=Z_1^{-1}e^{-S_B}q\,\dd\lambda,
 \qquad
 \dd\mu_2=Z_2^{-1}e^{-S_B}q^2\,\dd\lambda,
\]
one has
\[
 \dd\zeta_\chi=Z_1u\,\dd\mu_1
               =Z_2uq^{-1}\,\dd\mu_2.
\]
Division by the total complex mass proves \eqref{eq:V22-single-Q}.  If
\eqref{eq:V22-single-Q} and \eqref{eq:V22-square-phase-functional} agree on
all bounded tests, their normalized complex measures agree.  Their densities
with respect to $e^{-S_B}\lambda$ are proportional, so
$uq=c\,uq^2$.  Since $u$ and $q$ are nonzero, $q=c^{-1}$ almost everywhere.
\end{proof}

\begin{countertheorem}[Phase alone does not correct a norm-square control]
\label{cth:V22-phase-alone-wrong-weight}
Let a two-point base have determinant moduli $q=(1,2)$, phases
$u=(1,-1)$, and let $F=(0,1)$.  Then
\begin{equation}
 \boxed{
 \frac{\mathbb E_{\mu_1}(uF)}{\mathbb E_{\mu_1}u}
 =\frac{\mathbb E_{\mu_2}(uq^{-1}F)}
        {\mathbb E_{\mu_2}(uq^{-1})}
 =2,}
 \label{eq:V22-correct-two-point}
\end{equation}
whereas
\begin{equation}
 \boxed{
 \frac{\mathbb E_{\mu_2}(uF)}{\mathbb E_{\mu_2}u}
 =\frac43.}
 \label{eq:V22-wrong-two-point}
\end{equation}
\end{countertheorem}

\begin{proof}
The two polar laws are
$\mu_1=(1/3,2/3)$ and $\mu_2=(1/5,4/5)$.  Direct substitution gives the
three displayed ratios.
\end{proof}

\begin{assessment}[Scope of the Version V21 phase geometry]
\label{ass:V22-V21-scope}
The projective half-density identities of Version~V21 remain exact for every
explicitly declared vector
$\Psi_\alpha=(\dd\mu_\alpha/\dd\lambda)^{1/2}u$.  For the single chiral
complex measure they must be applied to the $\alpha=1$ polar law, or,
equivalently, to the complete weight $u/q$ on the populated $\alpha=2$
control.  Applying only the unit phase to the $\alpha=2$ law is a different
complex experiment unless the determinant modulus is constant.
\end{assessment}

\subsection{Reflection sewing before positivity}

Let $\Theta:X\to X$ be a measurable involution preserving $\mu_1$.  After a
physical line scalarization and division by the total partition phase, let
$v:X\to\mathrm U(1)$ be the aligned phase.  Define the reflected scalar form
\begin{equation}
 \mathfrak q_v(F,G)
 =\int_Xv(x)\overline{F(\Theta x)}G(x)\,\dd\mu_1(x).
 \label{eq:V22-reflected-scalar-form}
\end{equation}

\begin{theorem}[Exact reflection-sewing criterion]
\label{thm:V22-reflection-sewing}
The form \eqref{eq:V22-reflected-scalar-form} is Hermitian on all bounded
writers if and only if
\begin{equation}
 \boxed{v(x)=\overline{v(\Theta x)}\quad\mu_1\text{-almost everywhere}.}
 \label{eq:V22-sewing-equation}
\end{equation}
Equivalently, the positive sewing residual
\begin{equation}
 \boxed{
 \Delta_{\rm sew}^{\rm line}
 =\int_X\abs{v(x)-\overline{v(\Theta x)}}^2\,\dd\mu_1(x)}
 \label{eq:V22-sewing-residual}
\end{equation}
vanishes.  On a reflection-fixed set, sewing forces $v\in\{+1,-1\}$; it does
not by itself select the positive sign.
\end{theorem}

\begin{proof}
Changing variables $x\mapsto\Theta x$ gives
\[
 \overline{\mathfrak q_v(G,F)}
 =\int_X\overline{v(\Theta x)}
       \overline{F(\Theta x)}G(x)\,\dd\mu_1(x).
\]
Equality with \eqref{eq:V22-reflected-scalar-form} for every bounded $F,G$ is
equivalent to \eqref{eq:V22-sewing-equation}.  The residual is the squared
$L^2$ norm of the difference.  At a fixed point the equation reads
$v=\overline v$ and unit modulus leaves the two real signs.
\end{proof}

\subsection{The same-history phase--crossing measure}

Let $E$ be one finite represented one-letter coefficient space and assume that
the phase-quenched crossing gate has produced a measurable matrix
\begin{equation}
 K(x)=K(x)^*\succeq0
 \quad\text{on }E.
 \label{eq:V22-positive-history-crossing}
\end{equation}
For the single-determinant polar law define the positive and complex
operator-valued measures
\begin{equation}
 \mathsf M(A)=\int_AK\,\dd\mu_1,
 \qquad
 \boldsymbol\Gamma(A)=\int_AuK\,\dd\mu_1.
 \label{eq:V22-operator-measures}
\end{equation}
The second object, not the separate phase marginal and crossing marginal, is
the one-letter component of the chiral complex amplitude.

\begin{theorem}[Exact crossing-visible phase criterion]
\label{thm:V22-operator-phase-criterion}
There exists $z\in\mathrm U(1)$ such that
\begin{equation}
 \overline z\,\boldsymbol\Gamma(A)\succeq0
 \quad\text{for every measurable }A
 \label{eq:V22-positive-operator-measure}
\end{equation}
if and only if
\begin{equation}
 \boxed{u(x)=z\quad\mu_1\text{-almost everywhere on }
        \{x:\Tr K(x)>0\}.}
 \label{eq:V22-crossing-phase-constant}
\end{equation}
On that branch
$\overline z\,\boldsymbol\Gamma=\mathsf M$.
\end{theorem}

\begin{proof}
The reverse implication is immediate.  Conversely, taking the trace in
\eqref{eq:V22-positive-operator-measure} gives a positive scalar measure with
Radon--Nikodym density
$\overline zu(x)\Tr K(x)$ relative to $\mu_1$.  On the set where
$\Tr K>0$, that density must be real and nonnegative.  Its phase has modulus
one, so $\overline zu=1$ almost everywhere there.
\end{proof}

Put
\begin{equation}
 \tau_K=\int_X\Tr K\,\dd\mu_1,
 \qquad
 \zeta_K=\int_Xu\Tr K\,\dd\mu_1.
 \label{eq:V22-trace-carrier}
\end{equation}

\begin{corollary}[Crossing-visible phase cancellation debit]
\label{cor:V22-phase-crossing-debit}
If $\tau_K>0$, then
\begin{equation}
 \boxed{
 \Delta_{\rm ph\times}
 :=\inf_{|z|=1}
   \int_X\abs{u-z}^2\Tr K\,\dd\mu_1
 =2(\tau_K-|\zeta_K|)\ge0.}
 \label{eq:V22-phase-crossing-debit}
\end{equation}
It vanishes exactly on the branch of
\Cref{thm:V22-operator-phase-criterion}.  If the total partition phase has
already fixed $\omega$, the anchored debit is
\begin{equation}
 \boxed{
 \Delta_{\rm ph\times}(\omega)
 =2\bigl(\tau_K-\operatorname{Re}(\overline\omega\zeta_K)\bigr).}
 \label{eq:V22-anchored-debit}
\end{equation}
\end{corollary}

\begin{proof}
Expand
$|u-z|^2=2-2\operatorname{Re}(\overline zu)$ and maximize
$\operatorname{Re}(\overline z\zeta_K)$ over the unit circle.  Equality in the
triangle inequality for $\zeta_K$ is equivalent to constancy of the phase on
the support of the positive measure $(\Tr K)\mu_1$.
\end{proof}

\subsection{Record-relative vacuum and one-letter positivity}

Let $r:X\to\mathcal R$ be a finite physically occurring boundary record, with
cells $X_a=r^{-1}(a)$.  Put
\begin{equation}
 z_a=\int_{X_a}u\,\dd\mu_1,
 \qquad
 \Gamma_a=\int_{X_a}uK\,\dd\mu_1,
 \qquad
 m=\sum_az_a=\mathbb E_{\mu_1}u.
 \label{eq:V22-record-complex-blocks}
\end{equation}
For coefficient fields $(\alpha_a,v_a)\in\C\oplus E$, define
\begin{equation}
 \mathfrak Q_r((\alpha,v),(\beta,w))
 =\frac1m\sum_a
 \left(
  \overline{\alpha_a}z_a\beta_a
  +\ip{v_a}{\Gamma_aw_a}
 \right).
 \label{eq:V22-record-form}
\end{equation}

\begin{theorem}[Finite record-relative complex OS criterion]
\label{thm:V22-record-criterion}
Assume $m\ne0$ and put $\omega=m/|m|$.  The vacuum-plus-one-letter form
\eqref{eq:V22-record-form} is Hermitian and nonnegative if and only if, for
every occurring record cell,
\begin{equation}
 \boxed{
 \overline\omega z_a\in[0,\infty),
 \qquad
 \overline\omega\Gamma_a\in\operatorname{Herm}_+(E).}
 \label{eq:V22-record-positive-blocks}
\end{equation}
For the trivial record only the total matrix
$\overline\omega\boldsymbol\Gamma(X)$ is tested.  If all measurable selectors
are admitted and the vacuum modulus has full support, positivity is equivalent
to
\begin{equation}
 \boxed{u=\omega\quad\mu_1\text{-almost everywhere};}
 \label{eq:V22-full-selector-phase}
\end{equation}
in that case every positive historywise crossing mixes positively.
\end{theorem}

\begin{proof}
The record selectors make \eqref{eq:V22-record-form} the orthogonal direct sum
of the scalar blocks $z_a/m$ and matrix blocks $\Gamma_a/m$.  A direct sum is
Hermitian positive exactly when each block is.  Multiplication by the positive
number $|m|^{-1}$ reduces this to
\eqref{eq:V22-record-positive-blocks}.  With all selectors, the scalar set
function
$A\mapsto\overline\omega\int_Au\,\dd\mu_1$ is a positive measure.  Its
Radon--Nikodym density $\overline\omega u$ is nonnegative and of unit modulus,
so it equals one almost everywhere.  The converse is immediate.
\end{proof}

\begin{countertheorem}[Separate phase and crossing marginals do not determine the chiral OS branch]
\label{cth:V22-marginals-no-joint}
There are two four-history cards with the same positive law, the same phase
multiset, the same nonzero average phase, and the same strictly positive
crossing multiset, but with opposite normalized one-letter signs.
\end{countertheorem}

\begin{proof}
Use the uniform law on four histories and
\begin{equation}
 u=(1,1,1,-1),
 \qquad m=\frac12.
 \label{eq:V22-four-phase}
\end{equation}
For Card A and Card B take, respectively,
\begin{equation}
 K_A=(5,1,1,1),
 \qquad
 K_B=(1,1,1,5).
 \label{eq:V22-four-crossings}
\end{equation}
Both crossing multisets have mean two and are positive at every history.  Yet
\begin{equation}
 \boxed{
 \mathbb E(uK_A)=\frac32,
 \qquad
 \frac{\mathbb E(uK_A)}m=3>0,}
 \label{eq:V22-card-A}
\end{equation}
whereas
\begin{equation}
 \boxed{
 \mathbb E(uK_B)=-\frac12,
 \qquad
 \frac{\mathbb E(uK_B)}m=-1<0.}
 \label{eq:V22-card-B}
\end{equation}
The phase visibility, phase marginal, positive crossing distribution, and
phase-quenched crossing mean are identical.  Only their same-history mixed
incidence differs.
\end{proof}

\begin{assessment}[The phase and crossing rows are not merely parallel]
\label{ass:V22-not-parallel}
The Version~V21 phase Gram and the Version~V9--V16 phase-quenched crossing are
separate marginals of the actual chiral gate.  The physical one-letter target
is the mixed matrix
\begin{equation}
 \boxed{
 P_\chi
 =\frac{\mathbb E_{\mu_1}(uK)}{\mathbb E_{\mu_1}u}
 =\frac{\mathbb E_{\mu_2}(uq^{-1}K)}
        {\mathbb E_{\mu_2}(uq^{-1})}.}
 \label{eq:V22-physical-joint-crossing}
\end{equation}
A phase-only panel and a crossing-only panel cannot replace one direct
same-history phase--crossing Read or an equivalent common-cylinder theorem.
\end{assessment}

\subsection{Validated finite joint certificate}

Suppose a finite record card supplies complex interval centres
$\widehat m,\widehat z_a$ and matrices $\widehat\Gamma_a$, with
\begin{equation}
 |m-\widehat m|\le\eta_m,
 \quad
 |z_a-\widehat z_a|\le\eta_a^0,
 \quad
 \norm{\Gamma_a-\widehat\Gamma_a}_{\rm HS}\le\eta_a^1.
 \label{eq:V22-data-radii}
\end{equation}
Assume
\begin{equation}
 m_*:=|\widehat m|-\eta_m>0,
 \qquad
 \widehat\omega=\widehat m/|\widehat m|.
 \label{eq:V22-partition-floor}
\end{equation}

\begin{theorem}[Outward phase-aligned block radii]
\label{thm:V22-outward-joint}
The true partition phase $\omega=m/|m|$ obeys
\begin{equation}
 \boxed{
 |\omega-\widehat\omega|
 \le\delta_\omega:=\frac{2\eta_m}{m_*}.}
 \label{eq:V22-phase-radius}
\end{equation}
Consequently
\begin{equation}
 \boxed{
 |\overline\omega z_a-
   \overline{\widehat\omega}\widehat z_a|
 \le
 \rho_a^0:=\eta_a^0+\delta_\omega|\widehat z_a|,}
 \label{eq:V22-scalar-radius}
\end{equation}
\begin{equation}
 \boxed{
 \norm{\overline\omega\Gamma_a-
   \overline{\widehat\omega}\widehat\Gamma_a}_{\rm HS}
 \le
 \rho_a^1:=\eta_a^1+
 \delta_\omega\norm{\widehat\Gamma_a}_{\rm HS}.}
 \label{eq:V22-matrix-radius}
\end{equation}
If physical reflection sewing proves exact Hermiticity of the true aligned
blocks, strict inequalities
\begin{equation}
 \operatorname{Re}
 (\overline{\widehat\omega}\widehat z_a)>\rho_a^0,
 \qquad
 \lambda_{\min}
 \left(\operatorname{Herm}
 (\overline{\widehat\omega}\widehat\Gamma_a)\right)
 >\rho_a^1
 \label{eq:V22-strict-pass}
\end{equation}
for every occurring cell give a finite strict pass.  Conversely, if the
Hilbert--Schmidt distance of an estimated aligned block from its positive cone
exceeds its radius, the card has a strict obstruction.  Every remaining valid
interval is unresolved.
\end{theorem}

\begin{proof}
For nonzero complex numbers $z,w$,
\[
 \left|\frac z{|z|}-\frac w{|w|}\right|
 \le\frac{2|z-w|}{\min\{|z|,|w|\}}.
\]
Here $|m|\ge m_*$, proving \eqref{eq:V22-phase-radius}.  Add and subtract the
estimated aligned scalar and matrix to obtain
\eqref{eq:V22-scalar-radius}--\eqref{eq:V22-matrix-radius}.  Under exact
Hermiticity, the minimum eigenvalue is one-Lipschitz in operator norm, which is
bounded by the Hilbert--Schmidt radius.  Distance to a closed convex cone is
one-Lipschitz, proving the obstruction statement.
\end{proof}

\begin{firewall}[Approximate Hermiticity is not a positive pass]
\label{fire:V22-Hermitian-domain}
The positive cone lies in the real subspace of Hermitian matrices and has no
interior in the full complex matrix space.  A small anti-Hermitian error is not
made positive by an eigenvalue margin.  A pass in
\Cref{thm:V22-outward-joint} therefore requires exact physical reflection
Hermiticity or a validated identity forcing the anti-Hermitian part to vanish.
A strict nonzero lower bound on that part is an obstruction; an interval
meeting zero is unresolved.
\end{firewall}

\subsection{Real Pfaffian and independently conjugate branches}

\begin{theorem}[Three distinct scalarization branches]
\label{thm:V22-three-line-branches}
At the first fermionic line gate, the following branches are mathematically and
physically distinct.
\begin{enumerate}[label=\textup{(\roman*)},leftmargin=2.8em]
\item \emph{Single chiral complex branch.}  A physical dual line or scalarization
      supplies $d=qu$; the polar law is $\mu_1$, and the joint gate is
      \eqref{eq:V22-physical-joint-crossing}.
\item \emph{Independently conjugate doubled branch.}  Two physically occurring
      Grassmann factors $D$ and $D^*$ give the positive density $q^2$ by
      \eqref{eq:V22-double-Grassmann}; the complex phase cancels because the
      dual factor actually occurs.
\item \emph{Real/Pfaffian branch.}  A real skew coefficient and an oriented
      Pfaffian line give a sign $s\in\{\pm1\}$.  Its square
      $\operatorname{Pf}(A)^2=\det A$ loses that orientation.  On a full
      selector algebra the signed vacuum and one-letter kernel is positive
      exactly when, after one global orientation choice, $s=+1$ on every
      visible nonzero sector and the phase-quenched crossing is positive.
\end{enumerate}
Neither a Nambu rewrite nor a positive norm-square density selects among these
branches.
\end{theorem}

\begin{proof}
The first two statements are
\Cref{cor:V22-physical-scalarization,thm:V22-Grassmann-doubling}.  For a real
skew matrix the Pfaffian is real after an orientation is chosen and squares to
the determinant.  Apply \Cref{thm:V22-record-criterion} with unit phase
$u=s\in\{\pm1\}$.  Full selector positivity forces the sign to be the same
positive global orientation wherever the scalar or crossing weight is nonzero.
\end{proof}

\subsection{Adjacent-cutoff transport of the joint complex gate}

Let an adjacent fine card be transported to the old record and one-letter
screen.  Denote the old and transported fine partition amplitudes and joint
crossing blocks by $(m_0,\Gamma_0)$ and $(m_1,\Gamma_1)$.  Put
\begin{equation}
 \delta_m=|m_1-m_0|,
 \qquad
 \delta_\Gamma=\norm{\Gamma_1-\Gamma_0}_{\rm op}.
 \label{eq:V22-joint-defects}
\end{equation}

\begin{theorem}[Normalized joint crossing rectangle]
\label{thm:V22-joint-transport}
If
\begin{equation}
 |m_0|,|m_1|\ge m_*>0,
 \label{eq:V22-joint-partition-floor}
\end{equation}
then
\begin{equation}
 \boxed{
 \left\|\frac{\Gamma_1}{m_1}-
          \frac{\Gamma_0}{m_0}\right\|_{\rm op}
 \le
 \frac{\delta_\Gamma}{m_*}
 +\frac{\norm{\Gamma_0}_{\rm op}}{m_*^2}\delta_m.}
 \label{eq:V22-normalized-rectangle}
\end{equation}
The same estimate holds cellwise after descendant aggregation.  Uniform
bounds on the joint blocks, a uniform nonzero partition floor, and summable
adjacent joint defects give a unique fixed-screen limit of the normalized
single-chiral crossing.  If the partition floor fails, the unnormalized
complex measures and their zero/divisor packet remain the canonical outputs;
normalization is not transported through a zero.
\end{theorem}

\begin{proof}
Write
\[
 \frac{\Gamma_1}{m_1}-\frac{\Gamma_0}{m_0}
 =\frac{\Gamma_1-\Gamma_0}{m_1}
  +\Gamma_0\frac{m_0-m_1}{m_0m_1}
\]
and take norms.  The cofinal assertion is the Cauchy criterion after summing
the displayed bound.
\end{proof}

\subsection{Exact finite replay and revised physical acquisition}

The exact verifier accompanying this continuation checks:
\begin{enumerate}[label=\textup{(R\arabic*)},leftmargin=3.4em]
\item a Gaussian-integer determinant with
      $\det(D^*D)=|\det D|^2=8$;
\item the distinct polar laws $\mu_1=(1/3,2/3)$ and
      $\mu_2=(1/5,4/5)$, together with the correct $u/q$ reweighting;
\item the exact source metrics $\mathsf G_1=1$ and $\mathsf G_2=4$;
\item passing and failing reflection-sewing phases;
\item the four-history same-marginal cards
      \eqref{eq:V22-card-A}--\eqref{eq:V22-card-B};
\item the corresponding phase--crossing debits $1$ and $3$;
\item a real Pfaffian sign flip with unchanged determinant modulus;
\item a nonzero adjacent normalized-crossing change lying inside the exact
      rectangle \eqref{eq:V22-normalized-rectangle}.
\end{enumerate}
Every check uses exact integer, rational, or Gaussian-integer arithmetic.
Normal and optimized Python executions must produce byte-identical JSON.

\begin{theorem}[Version V22 chiral line/crossing theorem]
\label{thm:V22-terminal}
Versions~V1--V21 are retained.  Version~V22 additionally proves:
\begin{enumerate}[label=\textup{(V22.\arabic*)},leftmargin=5.3em]
\item the Hermitian determinant-line metric canonically supplies the norm square
      but no nonzero linear scalarization of a single chiral amplitude;
\item the populated norm-square cylinder is the $\alpha=2$ member of the
      modulus-power ladder, whereas the polar law of one scalar determinant is
      the $\alpha=1$ member;
\item an independently occurring conjugate Grassmann factor produces the norm
      square exactly, while a Nambu rewrite produces only one determinant;
\item relative to the norm-square control, the exact single-chiral weight is
      $u/q$, not the unit phase $u$;
\item reflection Hermiticity is the independent sewing equation
      \eqref{eq:V22-sewing-equation};
\item the actual one-letter chiral target is the same-history mixed measure
      $\boldsymbol\Gamma(A)=\int_AuK\,d\mu_1$ and is not determined by separate
      phase and crossing marginals;
\item full record-relative vacuum and one-letter positivity is equivalent to
      the cellwise criterion \eqref{eq:V22-record-positive-blocks};
\item the crossing-visible phase cancellation debit
      \eqref{eq:V22-phase-crossing-debit} vanishes exactly when one global phase
      makes the complete one-letter operator measure positive;
\item outward matrix radii and the adjacent-cutoff rectangle give terminating
      finite and cofinal joint-gate certificates.
\end{enumerate}
No item populates the actual RPESM determinant scalarization, conjugate species,
Pfaffian orientation, complex partition, gauge-history crossing, or joint
phase--crossing matrix.
\end{theorem}

\begin{proof}
Items~\textup{(V22.1)--(V22.4)} are
\Cref{lem:V22-Riesz-not-scalar,thm:V22-power-ladder,thm:V22-Grassmann-doubling,thm:V22-single-reweighting}.
Item~\textup{(V22.5)} is \Cref{thm:V22-reflection-sewing}.  Items
\textup{(V22.6)--(V22.8)} are
\Cref{thm:V22-operator-phase-criterion,thm:V22-record-criterion,cor:V22-phase-crossing-debit,cth:V22-marginals-no-joint}.
Item~\textup{(V22.9)} is
\Cref{thm:V22-outward-joint,thm:V22-joint-transport}.
\end{proof}

\begin{assessment}[Version V22 verdict]
\label{ass:V22-verdict}
The current result is
\begin{equation}
 \boxed{
 \begin{gathered}
 \passstatus\\[-0.15em]
 \text{line scalarization, modulus-power, and independent-doubling audit;}\\
 \text{reflection sewing, joint phase--crossing and record criterion;}\\
 \text{outward finite certificate and adjacent joint transport}
 \\[0.45em]
 \obsstatus\\[-0.15em]
 \text{norm-square control identified with one chiral determinant;}\\
 \text{phase }u\text{ used where }u/q\text{ is required;}\\
 \text{Nambu rewrite used as conjugate occurrence, or joint incidence inferred}
 \\[0.45em]
 \stopstatus\\[-0.15em]
 \text{first physically scalarized or independently doubled RPESM line,}\\
 \text{together with its direct same-history complex crossing card}
 \end{gathered}}
 \label{eq:V22-verdict}
\end{equation}
The protected norm-square field shell remains a literal positive control and
source-metric experiment.  It is promoted to a physical doubled fermion theory
only if the conjugate factor occurs, and it is promoted to a single chiral
complex measure only through the exact $u/q$ incidence or the $\alpha=1$ polar
law.  No exterior or Duhamel gate is opened before this fork is resolved.
\end{assessment}

\begin{acquisition}[Version V23: one amplitude-semantics and joint line--crossing card]
\label{acq:V22-next}
Preserve Versions~V1--V22.  On one protected field-shell history and one
adjacent cutoff, first declare exactly one of the following physical targets:
\begin{enumerate}[label=\textup{(A\arabic*)},leftmargin=3.8em]
\item a single chiral complex determinant with a physical dual-line
      scalarization;
\item an independently integrated conjugate double whose common action proves
      \eqref{eq:V22-double-Grassmann};
\item a real/Pfaffian branch with a transported orientation and direct sign.
\end{enumerate}
On the single chiral branch acquire, without separate-history fitting,
\begin{equation}
 \boxed{
 \left(
 q,\mu_1,u,m,
 \Delta_{\rm sew}^{\rm line},
 K_F,
 \{z_a,\Gamma_a\}_{a\in\mathcal R},
 \text{one held-out mixed phase--crossing Read}
 \right).}
 \label{eq:V22-next-card}
\end{equation}
If the populated norm-square cylinder is used as the sampler, acquire the
complete weight $u/q$ and verify its product incidence.  Test the complex
partition floor, reflection sewing, cellwise scalar and matrix positivity,
relation descent, and the held-out mixed Read.  Repeat the unnormalized complex
partition and crossing measures at one adjacent cutoff.  A real branch must
retain its Pfaffian sign; a doubled branch must retain the actual conjugate
Grassmann occurrence.  Until one branch passes, do not add another Wilson-sign,
phase-Gram, exterior, localizer, Duhamel, or continuum compiler.
\end{acquisition}

\section*{V22 exact replay and append-only record}
\addcontentsline{toc}{section}{V22 exact replay and append-only record}
The accompanying verifier reconstructs every stated finite control in exact
arithmetic and rejects the two deliberate substitutions: phase-only
reweighting of the norm-square control and permutation of the phase/crossing
incidence while all separate marginals are held fixed.  It is a theorem replay,
not a populated RPESM line or crossing measurement.

The complete Version~V21 source preceding its sole final active
\verb|\end{document}| is preserved byte for byte.  Its preterminal SHA--256
digest is recorded in the validation manifest.  A later continuation must
remove only the final active document-closing command, append new material
immediately before it, restore one terminal command, and increment the filename
to end in \texttt{\_V23.tex}.

% =============================================================================
% END APPENDED VERSION V22 MATERIAL
% =============================================================================


% =============================================================================
% BEGIN APPENDED VERSION V23 MATERIAL
% =============================================================================

\clearpage
\hypersetup{
 pdftitle={Renewal Einstein--Standard--Model Physical Source Metric, Duhamel Ancestry and Quantum Continuum Programme V23},
 pdfsubject={Positive joint line--crossing Gram, source ancestry, response dispersion, and route curvature}
}
\pdfinfo{
 /Title (Renewal Einstein--Standard--Model Physical Source Metric, Duhamel Ancestry and Quantum Continuum Programme V23)
 /Subject (Positive joint line--crossing Gram, source ancestry, response dispersion, and route curvature)
}

\section{Version V23: positive joint line--crossing Gram, source ancestry, and route curvature}
\label{sec:V23-positive-joint-line-crossing}
% =============================================================================

Version~V22 identified the physical one-letter chiral target as the
same-history complex block
\(
 \Gamma=\int wK\,\dd\mu
\),
where \(\mu\) is the declared positive sampler, \(w\) is its complete complex
line weight, and \(K\succeq0\) is the phase-quenched reflected crossing.  It
also proved that separate phase and crossing marginals do not determine this
block.  The remaining acquisition should therefore not be split again into a
phase panel and a crossing panel.

The complementary boundary and terminal programmes sharpen this point in two
ways.  First, a signed mixed coefficient can be recovered from one positive
amplitude-difference or block-Gram experiment, without constructing a stronger
stochastic coupling.  Second, source tightness does not imply tightness of its
causal response, and a vanishing static short can acquire a nonzero positive
curvature when the physical source map moves.  The present continuation
imports those lessons only at the level of method.  No Navier--Stokes or
Yang--Mills coefficient is inserted into the fermionic card.

The corrected finite order is
\begin{equation}
 \boxed{
 \begin{gathered}
 \text{one declared positive sampler and complete complex line weight}\
 \Downarrow\\
 \text{one positive joint line--crossing Gram}\
 \Downarrow\\
 \text{source-follower short or orthogonal phase--crossing source}\
 \Downarrow\\
 \text{obstruction-first mixed verdict and response-dispersion audit}\
 \Downarrow\\
 \text{physical route curvature and adjacent-cutoff transport}\
 \Downarrow\\
 \text{only then the conditional exterior and Duhamel gates.}
 \end{gathered}}
 \label{eq:V23-corrected-order}
\end{equation}
The positive sampler may be the single-determinant polar law of Version~V22,
the populated norm-square law together with the complete weight \(u/q\), an
independently doubled positive law, or an oriented real/Pfaffian branch.  The
weight is part of the target and is not refitted after the Gram is read.

\subsection{The positive joint line--crossing Gram}

Let \((X,\mu)\) be one finite or standard-Borel positive history cylinder, let
\(E\) be a finite-dimensional represented one-letter space, and let
\begin{equation}
 K:X\longrightarrow\operatorname{Herm}_+(E),
 \qquad
 w:X\longrightarrow\C
 \label{eq:V23-K-w}
\end{equation}
be measurable.  Assume
\begin{equation}
 \int_X(1+|w|^2)\Tr K\,\dd\mu<\infty.
 \label{eq:V23-integrability}
\end{equation}
For the single-chiral polar law one has \(w=u\) and \(|w|=1\).  If the
norm-square law is used as the sampler, then \(w=u/q\), so
\eqref{eq:V23-integrability} is a genuine inverse-modulus incidence condition.

Put \(\mathcal H=L^2(X,\mu;E)\) and define the two source syntheses
\begin{equation}
 (Av)(x)=K(x)^{1/2}v,
 \qquad
 (Bv)(x)=w(x)K(x)^{1/2}v.
 \label{eq:V23-two-syntheses}
\end{equation}
Their finite matrices are
\begin{equation}
 \boxed{
 M=A^*A=\int_XK\,\dd\mu,
 \quad
 N=B^*B=\int_X|w|^2K\,\dd\mu,
 \quad
 \Gamma=A^*B=\int_XwK\,\dd\mu.}
 \label{eq:V23-MNGamma}
\end{equation}

\begin{theorem}[Positive joint line--crossing Gram and exact source short]
\label{thm:V23-positive-joint-Gram}
The block matrix
\begin{equation}
 \boxed{
 \mathbb G_{w\times}
 =\begin{pmatrix}M&\Gamma\\ \Gamma^*&N\end{pmatrix}
 =\begin{pmatrix}A^*\\B^*\end{pmatrix}(A\;B)
 \succeq0}
 \label{eq:V23-joint-Gram}
\end{equation}
is the source-minimal positive carrier of the phase-quenched and complex
weighted crossing arms.  Positivity implies the canonical support incidences
\begin{equation}
 \Ker M\subseteq\Ker\Gamma^*,
 \qquad
 \Ker N\subseteq\Ker\Gamma.
 \label{eq:V23-support-incidence}
\end{equation}
With \(P_A=AM^\dagger A^*\), the supported Schur complement is
\begin{equation}
 \boxed{
 D_{w\times}
 :=N-\Gamma^*M^\dagger\Gamma
 =B^*(I-P_A)B\succeq0.}
 \label{eq:V23-joint-Schur}
\end{equation}
Moreover,
\begin{equation}
 \boxed{
 \rank D_{w\times}
 =\dim\Ran(A,B)-\dim\Ran A}
 \label{eq:V23-joint-rank}
\end{equation}
is the minimum number of additional linear response-source directions needed
beyond the phase-quenched crossing source.

The following are equivalent:
\begin{enumerate}[label=\textup{(\roman*)},leftmargin=2.8em]
\item \(D_{w\times}=0\);
\item \(\Ran B\subseteq\Ran A\);
\item
      \begin{equation}
       \boxed{B=AC,
       \qquad C=M^\dagger\Gamma.}
       \label{eq:V23-follower-map}
      \end{equation}
\end{enumerate}
On this branch \(C^*MC=N\).  If \(|w|=1\), then \(N=M\), so \(C\) is an
\(M\)-isometry on the supported coefficient quotient.  It need not be a scalar
phase.
\end{theorem}

\begin{proof}
For \(v,z\in E\),
\[
 \left\langle\binom vz,
 \mathbb G_{w\times}\binom vz\right\rangle
 =\|Av+Bz\|_{\mathcal H}^2,
\]
which proves positivity and the Gram identity.  If \(v\in\Ker M\), then
\(Av=0\), hence \(\Gamma^*v=B^*Av=0\); the second support inclusion is
analogous.  The Moore--Penrose formula makes \(P_A\) the orthogonal projection
onto \(\Ran A\).  Therefore
\[
 B^*(I-P_A)B
 =B^*B-B^*AM^\dagger A^*B
 =N-\Gamma^*M^\dagger\Gamma.
\]
Its rank is the dimension of \((I-P_A)\Ran B\), which is exactly the increase
from \(\Ran A\) to \(\Ran(A,B)\).  The Schur complement vanishes precisely
when \((I-P_A)B=0\).  In that case
\(B=P_AB=AM^\dagger A^*B=AC\), and the converse is immediate.  Finally,
\(B^*B=C^*A^*AC=C^*MC=N\).
\end{proof}

\begin{corollary}[Vacuum grade and average-phase visibility]
\label{cor:V23-vacuum-grade}
For the scalar vacuum crossing \(K=1\),
\begin{equation}
 \boxed{
 \mathbb G_{w,0}
 =\begin{pmatrix}1&m\\ \overline m&n_w\end{pmatrix}\succeq0,
 \qquad
 m=\mathbb E_\mu w,
 \quad
 n_w=\mathbb E_\mu|w|^2,}
 \label{eq:V23-vacuum-Gram}
\end{equation}
and
\begin{equation}
 \boxed{
 D_{w,0}=n_w-|m|^2
 =\mathbb E_\mu|w-m|^2.}
 \label{eq:V23-vacuum-Schur}
\end{equation}
For a unit phase this is \(1-|m|^2\).  For an actually doubled positive
branch \(w=1\), it is zero.
\end{corollary}

\begin{proof}
Apply \Cref{thm:V23-positive-joint-Gram} with the one-dimensional constant
source and expand the scalar variance.
\end{proof}

\begin{countertheorem}[Zero source short does not imply a global scalar phase]
\label{cth:V23-zero-short-not-scalar}
Let \(X=\{1,2\}\) with equal weights, let \(E=\C^2\), and put
\begin{equation}
 K(1)=2|e_1\rangle\langle e_1|,
 \qquad
 K(2)=2|e_2\rangle\langle e_2|,
 \qquad
 w(1)=1,
 \quad
 w(2)=-1.
 \label{eq:V23-nonscalar-follower-data}
\end{equation}
Then
\begin{equation}
 \boxed{
 M=N=I_2,
 \qquad
 \Gamma=\diag(1,-1),
 \qquad
 D_{w\times}=0,}
 \label{eq:V23-nonscalar-follower}
\end{equation}
although the phase is not constant.  The phase-weighted source is a follower
through the non-scalar coefficient map \(C=\diag(1,-1)\).  Full selector
positivity or a one-dimensional line scalarization therefore remains a
separate gate.
\end{countertheorem}

\begin{proof}
Direct averaging gives \eqref{eq:V23-nonscalar-follower}; then
\(\Gamma^*M^{-1}\Gamma=I_2=N\).  The two values of \(w\) are different, so no
global scalar phase represents the pointwise weight.
\end{proof}

\subsection{Reconstruction from positive coherent probes}

For \(\theta\in\R\), define the coherent source and its positive Gram
\begin{equation}
 S_\theta=A+e^{i\theta}B,
 \qquad
 Q_\theta=S_\theta^*S_\theta\succeq0.
 \label{eq:V23-coherent-probe}
\end{equation}

\begin{theorem}[Four positive probes reconstruct the complex mixed block]
\label{thm:V23-four-positive-probes}
For every \(\theta\),
\begin{equation}
 \boxed{
 Q_\theta
 =M+N+e^{i\theta}\Gamma+e^{-i\theta}\Gamma^*.}
 \label{eq:V23-Qtheta}
\end{equation}
Hence
\begin{align}
 \boxed{
 \operatorname{Herm}\Gamma
 =\frac{Q_0-Q_\pi}{4},}
 \label{eq:V23-Herm-reconstruction}\\
 \boxed{
 \operatorname{Im}\Gamma
 :=\frac{\Gamma-\Gamma^*}{2i}
 =\frac{Q_{-\pi/2}-Q_{\pi/2}}{4},}
 \label{eq:V23-Im-reconstruction}\\
 \boxed{
 \Gamma
 =\frac{Q_0-Q_\pi}{4}
 +\frac{i}{4}\left(Q_{-\pi/2}-Q_{\pi/2}\right).}
 \label{eq:V23-Gamma-reconstruction}
\end{align}
Together with the two positive arm Grams \(M\) and \(N\), these four
positive matrices reconstruct the complete joint Gram
\eqref{eq:V23-joint-Gram}.

For an independently retained analyzer angle \(\theta_*\), the positive
same-instrument occurrence residual is
\begin{equation}
 \boxed{
 \Delta_{\theta_*}^{\rm coh}
 =\left\|
 Q_{\theta_*}
 -M-N-e^{i\theta_*}\Gamma-e^{-i\theta_*}\Gamma^*
 \right\|_{\rm HS}^2.}
 \label{eq:V23-heldout-coherent-residual}
\end{equation}
It vanishes for one common coherent source experiment.  Algebraically
compatible probes acquired on different history populations do not establish
that occurrence.
\end{theorem}

\begin{proof}
Expand \(S_\theta^*S_\theta\).  The pairs \(\theta=0,\pi\) and
\(\theta=-\pi/2,\pi/2\) isolate the Hermitian and imaginary Hermitian parts,
respectively.  The held-out equation is the same expansion evaluated at an
angle not used in the reconstruction.
\end{proof}

\begin{remark}[No square-root phase is fitted]
The canonical positive square root \(K^{1/2}\) is used only to prove the Gram
representation.  Every acquired matrix in
\Cref{thm:V23-four-positive-probes} is the integral of the nonnegative
quadratic writer
\[
 v\longmapsto
 \left\|K(x)^{1/2}v+e^{i\theta}w(x)K(x)^{1/2}v\right\|^2.
\]
Changing a one-sided factor of \(K\) by a historywise unitary leaves the
matrices unchanged.  The card therefore does not fit a Kraus or hopping-source
phase.
\end{remark}

\subsection{Static followerhood and clean route curvature}

The preceding short is static.  A physical command or cutoff path can move
both source arms and the phase-weighted incidence.  The first nonzero motion
of a previously vanishing short is a positive curvature, not a new arbitrary
source declaration.

Let \(t\mapsto A_t,B_t:E\to\mathcal H\) be twice continuously differentiable
on one physically transported Hilbert carrier.  Assume \(A_t\) has constant
rank near \(t=0\), and put
\begin{equation}
 M_t=A_t^*A_t,
 \qquad
 \Gamma_t=A_t^*B_t,
 \qquad
 P_t=A_tM_t^\dagger A_t^*,
 \qquad
 D_t=B_t^*(I-P_t)B_t.
 \label{eq:V23-moving-short}
\end{equation}

\begin{theorem}[Clean phase--crossing birth curvature]
\label{thm:V23-clean-birth-curvature}
Suppose \(D_0=0\) and define
\begin{equation}
 C_0=M_0^\dagger\Gamma_0,
 \qquad
 Z_0=(I-P_0)(\dot B_0-\dot A_0C_0).
 \label{eq:V23-normal-relative-velocity}
\end{equation}
Then
\begin{equation}
 \boxed{
 D_0=0,
 \qquad
 \dot D_0=0,
 \qquad
 \ddot D_0=2Z_0^*Z_0\succeq0.}
 \label{eq:V23-clean-curvature}
\end{equation}
Furthermore,
\begin{equation}
 \boxed{
 \rank(Z_0^*Z_0)=\rank Z_0}
 \label{eq:V23-curvature-rank}
\end{equation}
is the minimum number of first-order phase--crossing response directions not
carried by the moving phase-quenched source range.
\end{theorem}

\begin{proof}
Let \(R_t=(I-P_t)B_t\), so \(D_t=R_t^*R_t\).  The hypothesis
\(D_0=0\) gives \(R_0=0\), and
\Cref{thm:V23-positive-joint-Gram} gives \(B_0=A_0C_0\).  Differentiating
\((I-P_t)A_t=0\) at zero yields
\[
 \dot P_0A_0=(I-P_0)\dot A_0.
\]
Therefore
\[
 \dot R_0
 =(I-P_0)\dot B_0-\dot P_0B_0
 =(I-P_0)(\dot B_0-\dot A_0C_0)
 =Z_0.
\]
Differentiating \(R_t^*R_t\) gives \(\dot D_0=0\) and
\(\ddot D_0=2\dot R_0^*\dot R_0\).  The rank statement is the ordinary Gram
rank identity.
\end{proof}

\begin{corollary}[Phase velocity after shorting by the crossing source]
\label{cor:V23-phase-velocity-short}
Assume at \(t=0\) that \(w_0=z\in\mathrm U(1)\) is constant and that
\begin{equation}
 \dot w_0(x)=izb(x),
 \qquad b:X\to\R.
 \label{eq:V23-phase-velocity}
\end{equation}
Let \(\mathsf M_b\) denote multiplication by \(b\) on \(\mathcal H\).  Then
\begin{equation}
 \boxed{
 \frac12\ddot D_0
 =A_0^*\mathsf M_b(I-P_0)\mathsf M_bA_0\succeq0.}
 \label{eq:V23-phase-velocity-Gram}
\end{equation}
It vanishes exactly when the phase-velocity source
\(\mathsf M_bA_0\) is already a coefficient follower of \(A_0\).  Otherwise
its rank is the source-minimal new phase response.  The positive phase Gram is
therefore charged only after the complete crossing-source short.
\end{corollary}

\begin{proof}
At zero, \(B_0=zA_0\) and \(C_0=zI\) on the supported coefficient space.
Since \(B_t=\mathsf M_{w_t}A_t\),
\[
 \dot B_0-\dot A_0C_0
 =iz\mathsf M_bA_0.
\]
Insert this into \eqref{eq:V23-clean-curvature}; the scalar phase cancels.
The kernel and rank statements follow from the Gram representation.
\end{proof}

\begin{corollary}[Target-native route-length bound]
\label{cor:V23-route-length}
Let \(R_t=(I-P_t)B_t\) and assume \(R_0=0\).  Then
\begin{equation}
 \boxed{
 \|D_1\|_{\rm op}^{1/2}
 =\|R_1\|_{\rm op}
 \le
 \int_0^1
 \|(I-P_t)\dot B_t-\dot P_tB_t\|_{\rm op}\,\dd t.}
 \label{eq:V23-route-length-bound}
\end{equation}
Thus an integrated bound on the actual transported source pair controls the
endpoint ancestry short.  A failure of a convenient upper bound is
\textsf{UNRESOLVED}; only a positive endpoint short or a positive curvature
Gram is a source-birth certificate.
\end{corollary}

\begin{proof}
The displayed integrand is \(\|\dot R_t\|_{\rm op}\).  Apply the fundamental
theorem of calculus and the triangle inequality, then use
\(D_1=R_1^*R_1\).
\end{proof}

\subsection{Obstruction-first use of the mixed chiral card}

Let
\begin{equation}
 m=\mathbb E_\mu w\ne0,
 \qquad
 P_\chi=\frac{\Gamma}{m}.
 \label{eq:V23-normalized-chiral-crossing}
\end{equation}
Exact reflection sewing and the Berezin convention must independently make
\(P_\chi\) Hermitian before positivity is tested.

\begin{theorem}[A negative mixed chiral direction exposes a bad history set]
\label{thm:V23-obstruction-first}
Let \(v\in E\) and put \(k_v(x)=\langle v,K(x)v\rangle\ge0\).  If
\begin{equation}
 \left\langle v,
 \operatorname{Herm}(P_\chi)v\right\rangle<0,
 \label{eq:V23-negative-mixed-direction}
\end{equation}
then
\begin{equation}
 \boxed{
 \mu\left(
 \left\{x:
 \operatorname{Re}(\overline m w(x))\,k_v(x)<0
 \right\}
 \right)>0.}
 \label{eq:V23-bad-history-set}
\end{equation}
For a finite occurring boundary record, at least one cell has a negative
aligned one-letter contribution on the same vector.  Thus a negative mixed
interval is already terminal for the mixed target and for a claim of
positivity on every visible cell.  It does not exclude the existence of some
other unselected good history.
\end{theorem}

\begin{proof}
Since \(|m|^2>0\),
\[
 \left\langle v,\operatorname{Herm}(P_\chi)v\right\rangle
 =\frac1{|m|^2}
 \int_X\operatorname{Re}(\overline m w(x))k_v(x)\,\dd\mu(x).
\]
A real integrable function with negative integral is negative on a set of
positive measure.  Decomposing the same integral over finitely many record
cells gives the cell statement.
\end{proof}

\begin{assessment}[Positive mixture semantics]
\label{ass:V23-positive-mixture-scope}
The converse direction is one-way.  Historywise aligned positivity implies
mixed positivity by convexity, but a positive mixed chiral card can hide a
negative history.  The exact Card~A of
\eqref{eq:V22-card-A} already gives such a case.  Phase or boundary selection
is not manufactured by the positive average.
\end{assessment}

\subsection{A quantified dispersion route to a positive mixed pass}

The absence of an exact joint block cannot in general be repaired from
marginals.  There is, however, a sharp sufficient branch when both the line
weight and the positive crossing fluctuate little relative to their means.
Put
\begin{equation}
 d_w^2=\mathbb E_\mu|w-m|^2
 =\mathbb E_\mu|w|^2-|m|^2,
 \label{eq:V23-weight-dispersion}
\end{equation}
and define the relative crossing-dispersion coefficient
\begin{equation}
 \boxed{
 \sigma_K^2
 :=\sup_{v:\,\langle v,Mv\rangle>0}
 \frac{\mathbb E_\mu
  |k_v-\mathbb E_\mu k_v|^2}
 {\langle v,Mv\rangle^2}.}
 \label{eq:V23-relative-K-dispersion}
\end{equation}
The value may be infinite.  On a finite coefficient screen it is an ordinary
compact generalized supremum.  Each numerator has the positive two-copy
representation
\begin{equation}
 \mathbb E|k_v-\mathbb E k_v|^2
 =\frac12\mathbb E|k_v(X_1)-k_v(X_2)|^2.
 \label{eq:V23-two-copy-dispersion}
\end{equation}

\begin{theorem}[Relative phase--crossing dispersion bound]
\label{thm:V23-relative-dispersion}
For every \(v\in E\),
\begin{equation}
 \boxed{
 \left\langle v,
 \operatorname{Herm}(P_\chi)v\right\rangle
 \ge
 \left(1-\frac{d_w\sigma_K}{|m|}\right)
 \langle v,Mv\rangle.}
 \label{eq:V23-dispersion-lower}
\end{equation}
Consequently, if exact reflection sewing gives \(P_\chi=P_\chi^*\),
\begin{equation}
 \boxed{|m|>d_w\sigma_K}
 \label{eq:V23-dispersion-pass}
\end{equation}
implies
\begin{equation}
 \boxed{
 P_\chi
 \succeq
 \left(1-\frac{d_w\sigma_K}{|m|}\right)M.}
 \label{eq:V23-dispersion-pass-floor}
\end{equation}
If \(M\succeq\kappa P_M\), the right side gives the explicit supported floor
\(
 \kappa(1-d_w\sigma_K/|m|)
\).

The same theorem applies conditionally on each finite boundary-record cell.
A cellwise visibility floor and cellwise dispersion bounds therefore give a
record-relative pass without claiming pointwise history positivity.
\end{theorem}

\begin{proof}
Since \(M=\mathbb E K\),
\begin{equation}
 \Gamma-mM
 =\mathbb E\bigl[(w-m)(K-M)\bigr].
 \label{eq:V23-covariance-decomposition}
\end{equation}
For \(\overline k_v=\langle v,Mv\rangle\), Cauchy--Schwarz and
\eqref{eq:V23-relative-K-dispersion} give
\[
 \left|
 \langle v,(\Gamma-mM)v\rangle
 \right|
 \le d_w\sigma_K\overline k_v.
\]
Therefore
\[
 \operatorname{Re}
 \langle v,\overline m\Gamma v\rangle
 \ge
 |m|^2\overline k_v-|m|d_w\sigma_K\overline k_v.
\]
Divide by \(|m|^2\).  Exact Hermiticity converts the quadratic-form bound into
the stated Loewner inequality.  Conditional application uses the same proof
with the normalized cell law.
\end{proof}

\begin{remark}[A sufficient marginal route is not joint occurrence]
The strict inequality \eqref{eq:V23-dispersion-pass} can certify the sign of
the mixed matrix when exact Hermiticity is supplied independently.  It does
not reconstruct \(\Gamma\), its source ancestry, or its cutoff transport.  If
the bound fails, the programme returns to the direct positive joint Gram; it
does not fit a phase or crossing coupling.
\end{remark}

\subsection{Source tightness is not joint-response tightness}

\begin{countertheorem}[Fixed phase and crossing marginals can have an oscillating joint response]
\label{cth:V23-marginal-tight-no-joint}
Use the four-history cards of
\eqref{eq:V22-four-phase}--\eqref{eq:V22-four-crossings}, choosing Card~A at
even cutoffs and Card~B at odd cutoffs.  At every cutoff the phase law, the
crossing multiset, the positive crossing mean, the vacuum amplitude, and the
two diagonal source Grams are identical:
\begin{equation}
 \boxed{
 m=\frac12,
 \qquad
 M=N=2.}
 \label{eq:V23-fixed-marginals}
\end{equation}
Nevertheless,
\begin{equation}
 \boxed{
 \Gamma_{2j}=\frac32,
 \qquad
 \Gamma_{2j+1}=-\frac12,}
 \label{eq:V23-oscillating-Gamma}
\end{equation}
so
\begin{equation}
 \boxed{
 P_{\chi,2j}=3,
 \qquad
 P_{\chi,2j+1}=-1.}
 \label{eq:V23-oscillating-normalized}
\end{equation}
The positive block Grams and their source shorts alternate between
\begin{equation}
 \mathbb G_A=
 \begin{pmatrix}2&3/2\\3/2&2\end{pmatrix},
 \qquad
 D_A=\frac78,
 \label{eq:V23-card-A-Gram}
\end{equation}
and
\begin{equation}
 \mathbb G_B=
 \begin{pmatrix}2&-1/2\\-1/2&2\end{pmatrix},
 \qquad
 D_B=\frac{15}{8}.
 \label{eq:V23-card-B-Gram}
\end{equation}
Thus tightness and even exact cutoff constancy of the separate source
marginals do not imply tightness of the physical joint response.
\end{countertheorem}

\begin{proof}
The phase and crossing data are those of the exact V22 counterexample.  Direct
averaging gives \eqref{eq:V23-fixed-marginals}--
\eqref{eq:V23-oscillating-Gamma}.  Division by \(m=1/2\) gives
\eqref{eq:V23-oscillating-normalized}; the two scalar Schur complements are
\(N-|\Gamma|^2/M\).
\end{proof}

\begin{assessment}[Response-tightness correction]
\label{ass:V23-response-tightness}
The diagonal blocks \(M\) and \(N\) are source stocks.  The mixed block
\(\Gamma\) is the source-to-complex-response incidence.  A cofinal theorem
must transport the complete block Gram or an equivalent direct mixed Read.
Separate compactness of the source, line, and crossing marginals leaves a
joint-response escape exactly as in
\Cref{cth:V23-marginal-tight-no-joint}.
\end{assessment}

\subsection{Direct joint-Gram transport}

Transport adjacent cutoff cards to one fixed coefficient screen.  Let
\begin{equation}
 \mathbb G_j=\begin{pmatrix}M_j&\Gamma_j\\\Gamma_j^*&N_j\end{pmatrix},
 \qquad
 D_j=N_j-\Gamma_j^*M_j^\dagger\Gamma_j,
 \qquad j\in\{0,1\}.
 \label{eq:V23-two-joint-cards}
\end{equation}
Assume a common supported projection \(P\) and a floor
\begin{equation}
 M_j=P M_jP\succeq s_*P,
 \qquad s_*>0.
 \label{eq:V23-common-support-floor}
\end{equation}
Put
\begin{equation}
 \delta_M=\|M_1-M_0\|,
 \quad
 \delta_N=\|N_1-N_0\|,
 \quad
 \delta_\Gamma=\|\Gamma_1-\Gamma_0\|,
 \quad
 \gamma=\max_j\|\Gamma_j\|.
 \label{eq:V23-component-defects}
\end{equation}

\begin{theorem}[Supported transport of the joint source short]
\label{thm:V23-joint-short-transport}
Under \eqref{eq:V23-common-support-floor},
\begin{equation}
 \boxed{
 \|D_1-D_0\|_{\rm op}
 \le
 \delta_N
 +\frac{2\gamma}{s_*}\delta_\Gamma
 +\frac{\gamma^2}{s_*^2}\delta_M.}
 \label{eq:V23-short-transport-bound}
\end{equation}
If, in addition, the complex partition amplitudes obey
\begin{equation}
 |m_0|,|m_1|\ge m_*>0,
 \qquad
 \delta_m=|m_1-m_0|,
 \label{eq:V23-partition-transport-data}
\end{equation}
then
\begin{equation}
 \boxed{
 \left\|
 \frac{\Gamma_1}{m_1}-\frac{\Gamma_0}{m_0}
 \right\|_{\rm op}
 \le
 \frac{\delta_\Gamma}{m_*}
 +\frac{\gamma}{m_*^2}\delta_m.}
 \label{eq:V23-normalized-transport-bound}
\end{equation}
A change of the canonical support \(P\), or loss of either floor, is retained
as a support or partition-visibility event rather than inserted into a
pseudoinverse perturbation bound.
\end{theorem}

\begin{proof}
On the common support,
\[
 M_1^\dagger-M_0^\dagger
 =M_1^\dagger(M_0-M_1)M_0^\dagger,
\]
so
\(
 \|M_1^\dagger-M_0^\dagger\|
 \le s_*^{-2}\delta_M
\).
Insert and subtract intermediate terms:
\begin{align*}
 \Gamma_1^*M_1^\dagger\Gamma_1-
 \Gamma_0^*M_0^\dagger\Gamma_0
 ={}&
 (\Gamma_1-\Gamma_0)^*M_1^\dagger\Gamma_1\\
 &+\Gamma_0^*(M_1^\dagger-M_0^\dagger)\Gamma_1\\
 &+\Gamma_0^*M_0^\dagger(\Gamma_1-\Gamma_0).
\end{align*}
Taking norms gives \eqref{eq:V23-short-transport-bound}.  The normalized
bound follows by
\[
 \frac{\Gamma_1}{m_1}-\frac{\Gamma_0}{m_0}
 =\frac{\Gamma_1-\Gamma_0}{m_1}
  +\Gamma_0\frac{m_0-m_1}{m_0m_1}.
\]
\end{proof}

\begin{theorem}[Cofinal joint line--crossing alternative]
\label{thm:V23-cofinal-alternative}
On every fixed physical one-letter screen, first audit the following rows:
\begin{enumerate}[label=\textup{(J\arabic*)},leftmargin=3.2em]
\item the positive block Grams \(\mathbb G_n\) are transported by the declared
      source and line maps and their adjacent operator-norm defects are
      summable;
\item the supported \(M_n\) have one fixed canonical support and floor
      \(s_*>0\);
\item the complex partition amplitudes have a floor \(|m_n|\ge m_*>0\) and
      summable adjacent defects;
\item reflection sewing and Hermiticity are transported on the same card.
\end{enumerate}
If one row fails, the corresponding support, partition, sewing, or joint-Gram
transport packet is returned before normalization.  If all four rows pass,
then \(M_n,N_n,\Gamma_n,D_n\), and the normalized chiral crossings
\(P_{\chi,n}=\Gamma_n/m_n\) converge uniquely on that screen, and exactly one
of the following prioritized branches occurs along a subsequence:
\begin{enumerate}[label=\textup{(C23.\arabic*)},leftmargin=5.8em]
\item a strict negative or skew normalized crossing gives a mixed chiral
      obstruction;
\item the limit short \(D_\infty\) has a positive protected eigenvalue and
      supplies a source-minimal phase--crossing response survivor;
\item \(D_n\to0\), so the complex weighted crossing is an asymptotic linear
      follower of the phase-quenched source in the joint source norm; scalar
      phase alignment remains a separate line criterion.
\end{enumerate}
There is no branch in which separate source marginals are declared sufficient
without the mixed block.
\end{theorem}

\begin{proof}
A failed audit row has the stated typed output by definition of the card.  On
the passing domain, summable block-Gram defects make each matrix block Cauchy.
The common support floor and \Cref{thm:V23-joint-short-transport} make the
Schur complements Cauchy, while the partition floor gives convergence of the
normalized crossings.  Hermiticity and cone membership are closed
finite-dimensional conditions.  A failing strict cone test gives the first
branch.  Otherwise the nonnegative limit \(D_\infty\) either has a positive
eigenvalue on a protected screen or vanishes.  In the latter case
\(\|(I-P_{A_n})B_n\|^2=\|D_n\|\to0\), which is the asserted asymptotic
follower statement.  The exact zero criterion is
\Cref{thm:V23-positive-joint-Gram}.
\end{proof}

\subsection{Exact finite replay}

The standard-library verifier accompanying this continuation checks the
following exact controls.
\begin{enumerate}[label=\textup{(R\arabic*)},leftmargin=3.4em]
\item A scalar two-history phase card with
      \(K=(1,3)\) and \(w=(1,i)\) has
      \begin{equation}
       \boxed{
       M=N=2,
       \quad
       \Gamma=\frac{1+3i}{2},
       \quad
       D_{w\times}=\frac34.}
       \label{eq:V23-complex-control}
      \end{equation}
      Its four positive coherent Grams are
      \begin{equation}
       \boxed{Q_0=5,
       \quad Q_\pi=3,
       \quad Q_{\pi/2}=1,
       \quad Q_{-\pi/2}=7,}
       \label{eq:V23-complex-four-probes}
      \end{equation}
      which reconstruct \(\Gamma\) exactly.
\item The nonconstant phase-follower card
      \eqref{eq:V23-nonscalar-follower} has zero short with a non-scalar
      coefficient map.
\item A rotating source line
      \(A_t=e_1\),
      \(B_t=(\cos t)e_1+(\sin t)e_2\)
      obeys
      \begin{equation}
       \boxed{D_t=\sin^2t,
       \qquad D_0'=0,
       \qquad D_0''=2.}
       \label{eq:V23-rotation-control}
      \end{equation}
\item With weights \((3/4,1/4)\), line values \((1,-1)\), and scalar
      crossings \((2,1)\),
      \begin{equation}
       \boxed{
       m=\frac12,
       \quad M=\frac74,
       \quad \Gamma=\frac54,
       \quad P_\chi=\frac52,}
       \label{eq:V23-dispersion-control}
      \end{equation}
      while
      \begin{equation}
       d_w^2=\frac34,
       \qquad
       \sigma_K^2=\frac3{49},
       \qquad
       1-\frac{d_w\sigma_K}{|m|}=\frac47.
       \label{eq:V23-dispersion-control-bound}
      \end{equation}
      The certified lower bound is therefore \(P_\chi\ge1\).
\item The alternating V22 cards have the constant marginals and oscillating
      joint responses in
      \eqref{eq:V23-fixed-marginals}--\eqref{eq:V23-card-B-Gram}.
\end{enumerate}
All rational and Gaussian-rational quantities are represented exactly.  The
trigonometric curvature check is performed from the exact derivatives at zero.
Ordinary and optimized Python executions must produce byte-identical JSON.

\begin{theorem}[Version V23 positive joint line--crossing theorem]
\label{thm:V23-terminal}
Versions~V1--V22 are retained.  Version~V23 additionally proves:
\begin{enumerate}[label=\textup{(V23.\arabic*)},leftmargin=5.3em]
\item every declared sampler/line/crossing branch has the intrinsic positive
      block Gram \eqref{eq:V23-joint-Gram};
\item its supported Schur complement is the exact source-ancestry residual and
      its rank is the minimum new phase--crossing response dimension;
\item six positive arm/coherent matrices reconstruct the complete complex
      mixed block, with one held-out coherent analyzer testing occurrence;
\item a static follower short has zero first derivative, and its first possible
      reopening is the positive route curvature \eqref{eq:V23-clean-curvature};
\item after a constant phase, that curvature is precisely the phase-velocity
      source shorted by the complete crossing-source range;
\item a negative mixed chiral direction exposes a positive-measure bad history
      set, while a positive mixed card is not a historywise pass;
\item the relative dispersion estimate \eqref{eq:V23-dispersion-lower} gives a
      strict sufficient mixed or cellwise pass without fitting the joint law;
\item separate line and crossing marginal tightness does not imply joint
      response tightness;
\item a summably transported positive joint Gram, stable support floor,
      partition floor, and sewing row give a unique cofinal normalized crossing
      and ancestry short.
\end{enumerate}
No item populates the actual RPESM scalarization, complete weight, coherent
line--crossing analyzer, gauge-history crossing, or route derivative.
\end{theorem}

\begin{proof}
Items~\textup{(V23.1)--(V23.3)} are
\Cref{thm:V23-positive-joint-Gram,thm:V23-four-positive-probes}.  Items
\textup{(V23.4)--(V23.5)} are
\Cref{thm:V23-clean-birth-curvature,cor:V23-phase-velocity-short,cor:V23-route-length}.
Items~\textup{(V23.6)--(V23.7)} are
\Cref{thm:V23-obstruction-first,thm:V23-relative-dispersion}.  Items
\textup{(V23.8)--(V23.9)} are
\Cref{cth:V23-marginal-tight-no-joint,thm:V23-joint-short-transport,thm:V23-cofinal-alternative}.
\end{proof}

\begin{assessment}[Version V23 verdict]
\label{ass:V23-verdict}
The current result is
\begin{equation}
 \boxed{
 \begin{gathered}
 \passstatus\\[-0.15em]
 \text{positive joint line--crossing Gram and source-minimal short;}\\
 \text{positive coherent acquisition, obstruction-first use, dispersion bound;}\\
 \text{clean route curvature and direct cofinal joint transport}
 \\[0.45em]
 \obsstatus\\[-0.15em]
 \text{separate phase/crossing marginals used as a joint response;}\\
 \text{zero static short promoted through a moving route without curvature;}\\
 \text{or a non-scalar follower map promoted to global line phase}
 \\[0.45em]
 \stopstatus\\[-0.15em]
 \text{first physically populated positive joint line--crossing Gram,}\\
 \text{held-out coherent analyzer, and physical route derivative.}
 \end{gathered}}
 \label{eq:V23-verdict}
\end{equation}
The new card closes the acquisition form of Version~V22 without choosing a
stronger process or refitting an amplitude branch.  It also prevents a static
follower verdict from being transported through a moving determinant line,
reflection frame, or crossing source without the corresponding normal route
velocity.
\end{assessment}

\begin{acquisition}[Version V24: execute one positive joint RPESM line--crossing card]
\label{acq:V23-next}
Preserve Versions~V1--V23.  On one protected RPESM history, one finite
boundary record, and one adjacent cutoff, declare exactly one amplitude
semantics from Version~V22 and populate
\begin{equation}
 \boxed{
 \left(
 \mu,w,m,
 M,N,
 Q_0,Q_\pi,Q_{\pi/2},Q_{-\pi/2},
 Q_{\theta_*},
 \Delta_{\rm sew}^{\rm line},
 \text{supports and outward radii}
 \right).}
 \label{eq:V23-next-card}
\end{equation}
Reconstruct \(\Gamma\), test the full positive block Gram, the V22 normalized
record criterion, the source short \(D_{w\times}\), and the held-out coherent
analyzer.  If \(D_{w\times}=0\), acquire one physical command or adjacent-cutoff
secant of the two source arms and evaluate the normal relative velocity in
\eqref{eq:V23-normal-relative-velocity}; do not infer dynamic followerhood from
the static short.  If the direct joint block is unavailable but exact
Hermiticity and the strict dispersion bound
\eqref{eq:V23-dispersion-pass} are available, record only the corresponding
mixed or cellwise positivity pass.

A negative normalized mixed block returns the obstruction of
\Cref{thm:V23-obstruction-first}.  A positive ancestry short appends only its
polar source and is not charged again as an independent determinant phase,
reflection source, and crossing source.  Repeat the complete positive block
Gram and partition amplitude at one adjacent cutoff.  Until this card lands,
do not add another phase, crossing, exterior, localizer, Duhamel, or continuum
compiler.
\end{acquisition}

\section*{V23 exact replay and append-only record}
\addcontentsline{toc}{section}{V23 exact replay and append-only record}
The accompanying verifier reconstructs the joint Gram, coherent probes,
source shorts, clean curvature, relative dispersion bound, and alternating
response counterexample in exact arithmetic.  It is a theorem replay and not a
measurement of the active RPESM determinant line or reflected crossing.

The complete Version~V22 source preceding its sole final active
\verb|\end{document}| is preserved byte for byte.  Its preterminal SHA--256
digest is recorded in the validation manifest.  A later continuation must
remove only the final active document-closing command, append new material
immediately before it, restore one terminal command, and increment the filename
to end in \texttt{\_V24.tex}.

% =============================================================================
% END APPENDED VERSION V23 MATERIAL
% =============================================================================


% =============================================================================
% BEGIN APPENDED VERSION V24 MATERIAL
% =============================================================================

\clearpage
\hypersetup{
 pdftitle={Renewal Einstein--Standard--Model Physical Source Metric, Duhamel Ancestry and Quantum Continuum Programme V24},
 pdfsubject={Boundary-record-adapted line--crossing assembly, four-outcome positive tomography, and router curvature}
}
\pdfinfo{
 /Title (Renewal Einstein--Standard--Model Physical Source Metric, Duhamel Ancestry and Quantum Continuum Programme V24)
 /Subject (Boundary-record-adapted line--crossing assembly, four-outcome positive tomography, and router curvature)
}

\section{Version V24: boundary-record-adapted line--crossing assembly, four-outcome tomography, and router curvature}
\label{sec:V24-record-adapted-line-crossing}
% =============================================================================

Version~V23 placed the phase-quenched crossing source and its complex
line-weighted response in one positive block Gram and identified the supported
Schur complement as the source-minimal response innovation.  That short is
well defined only after one further upstream choice has been fixed: whether
the declared positive-time algebra contains a boundary record whose selectors
allow the follower coefficient to depend on the record cell.  A record-blind
coefficient and a record-adapted coefficient are different physical reuse
classes and need not have the same source short.

The related boundary and terminal programmes point to the same order of
operations.  A positive module-distance writer can recover a signed mixed
coefficient, but the source and its later response must still be assembled on
one common record.  Source tightness does not imply response tightness, and a
moving material or routing map can create a positive curvature even when each
instantaneous source is a follower.  The present continuation imports these
lessons only as structural guidance.  It introduces no Navier--Stokes or
Yang--Mills coefficient into the fermionic card.

The corrected order is
\begin{equation}
 \boxed{
 \begin{gathered}
 \text{physical reflection-fixed boundary record and its selector algebra}
 \\\Downarrow\\
 \text{record-cell positive line--crossing Grams}
 \\\Downarrow\\
 \text{within-cell source short plus record-router dispersion}
 \\\Downarrow\\
 \text{one four-outcome positive arm instrument and held-out analyzer}
 \\\Downarrow\\
 \text{record-refinement and adjacent-cutoff transport}
 \\\Downarrow\\
 \text{only then phase--crossing source birth or follower promotion.}
 \end{gathered}}
 \label{eq:V24-corrected-order}
\end{equation}
A finer record may be used only when its selectors physically occur in the
positive-time algebra and are transported with the card.  A partition chosen
after inspecting the fitted follower maps is not a source reduction.

\subsection{Record-cell Grams and the assembly-before-short identity}

Let \(\mathcal R\) be a finite reflection-fixed boundary record and let
\(X=\bigsqcup_{a\in\mathcal R}X_a\) be its measurable cell decomposition.
The cells may be atoms of a finite physical record or a predeclared finite
coarsening of a standard-Borel boundary record.  Let \(E\) be a finite
coefficient space.  On
\(
 \mathcal H_a=L^2(X_a,\mu;E)
\)
define the two source syntheses
\begin{equation}
 (A_av)(x)=K(x)^{1/2}v,
 \qquad
 (B_av)(x)=w(x)K(x)^{1/2}v,
 \qquad x\in X_a,
 \label{eq:V24-cell-syntheses}
\end{equation}
under the integrability hypothesis of
\eqref{eq:V23-integrability}.  Put
\begin{equation}
 \boxed{
 M_a=A_a^*A_a,
 \qquad
 \Gamma_a=A_a^*B_a,
 \qquad
 N_a=B_a^*B_a,}
 \label{eq:V24-cell-blocks}
\end{equation}
\begin{equation}
 \boxed{
 C_a=M_a^\dagger\Gamma_a,
 \qquad
 D_a=N_a-\Gamma_a^*M_a^\dagger\Gamma_a\succeq0.}
 \label{eq:V24-cell-decoder-short}
\end{equation}
The matrices are unnormalized cell masses; no division by
\(\mu(X_a)\) is needed.  Empty cells are omitted.

The vacuum grade is retained separately.  Put
\begin{equation}
 \boxed{
 p_a=\mu(X_a),
 \qquad
 z_a=\int_{X_a}w\,\dd\mu,
 \qquad
 n_a^{w}=\int_{X_a}|w|^2\,\dd\mu,}
 \label{eq:V24-vacuum-cell-blocks}
\end{equation}
so that
\begin{equation}
 \mathbb G_{a,0}=
 \begin{pmatrix}p_a&z_a\\\overline z_a&n_a^w\end{pmatrix}\succeq0,
 \qquad
 m=\sum_az_a.
 \label{eq:V24-vacuum-cell-Gram}
\end{equation}
These vacuum blocks carry the cell partition amplitudes required by the
record-relative Version~V22 normalization.  They are not inferred from the
one-letter crossing blocks.

Assemble the record-blind source maps on the common coefficient space by
\begin{equation}
 \boxed{
 \begin{aligned}
 A:E&\longrightarrow\bigoplus_{a\in\mathcal R}\mathcal H_a,
 &Av&=(A_av)_{a\in\mathcal R},\\
 B:E&\longrightarrow\bigoplus_{a\in\mathcal R}\mathcal H_a,
 &Bv&=(B_av)_{a\in\mathcal R}.
 \end{aligned}}
 \label{eq:V24-assembled-syntheses}
\end{equation}
The same coefficient vector is used in every cell.  Then
\begin{equation}
 M=\sum_aM_a,
 \qquad
 \Gamma=\sum_a\Gamma_a,
 \qquad
 N=\sum_aN_a,
 \qquad
 C=M^\dagger\Gamma.
 \label{eq:V24-aggregate-blocks}
\end{equation}

\begin{theorem}[Assembly before short and exact record-router Pythagoras]
\label{thm:V24-assembly-before-short}
For every occurring cell,
\begin{equation}
 \mathbb G_a=
 \begin{pmatrix}M_a&\Gamma_a\\\Gamma_a^*&N_a\end{pmatrix}
 \succeq0,
 \qquad
 D_a=B_a^*(I-P_{A_a})B_a,
 \label{eq:V24-cell-positive-Gram}
\end{equation}
and the aggregate block is
\begin{equation}
 \mathbb G=
 \sum_a\mathbb G_a
 =\begin{pmatrix}M&\Gamma\\\Gamma^*&N\end{pmatrix}
 \succeq0.
 \label{eq:V24-aggregate-positive-Gram}
\end{equation}
Its record-blind source short has the exact decomposition
\begin{equation}
 \boxed{
 \begin{aligned}
 D_{\rm blind}
 &:={N-\Gamma^*M^\dagger\Gamma}\\
 &=\sum_{a\in\mathcal R}D_a
 +\underbrace{\sum_{a\in\mathcal R}
   (C_a-C)^*M_a(C_a-C)}_{\displaystyle
   \mathcal D_{\rm rout}(\mathcal R)\succeq0}.
 \end{aligned}}
 \label{eq:V24-router-Pythagoras}
\end{equation}
The global decoder is the supported \(M_a\)-weighted barycentre:
\begin{equation}
 \boxed{
 \sum_aM_a(C_a-C)=0.}
 \label{eq:V24-decoder-normal-equation}
\end{equation}
Equivalently,
\begin{equation}
 \boxed{
 \mathcal D_{\rm rout}(\mathcal R)
 =\sum_a\Gamma_a^*M_a^\dagger\Gamma_a
  -\Gamma^*M^\dagger\Gamma.}
 \label{eq:V24-router-Schur-difference}
\end{equation}
Thus the record-blind short contains two different positive quantities:
within-cell phase--crossing novelty \(\sum_aD_a\), and the price of forcing
one common follower map across physically distinct boundary cells.
\end{theorem}

\begin{proof}
Each cell block is the Gram of \((A_a\;B_a)\), so
\Cref{thm:V23-positive-joint-Gram} gives
\begin{equation}
 B_a=A_aC_a+Z_a,
 \qquad
 Z_a=(I-P_{A_a})B_a,
 \qquad
 A_a^*Z_a=0,
 \qquad
 Z_a^*Z_a=D_a.
 \label{eq:V24-cell-orthogonal-split}
\end{equation}
The aggregate block is the sum of the positive cell blocks.  Its support
incidence gives \(MM^\dagger\Gamma=\Gamma\), and therefore
\[
 \sum_aM_a(C_a-C)
 =\sum_a\Gamma_a-MC
 =\Gamma-MM^\dagger\Gamma
 =0.
\]
The record-blind residual synthesis is
\[
 B-AC
 =\bigoplus_a\bigl(A_a(C_a-C)+Z_a\bigr).
\]
The direct-sum cells are orthogonal, and inside each cell
\(A_a^*Z_a=0\).  Taking the Gram of this residual gives
\eqref{eq:V24-router-Pythagoras}.  Expanding the second term and using
\eqref{eq:V24-decoder-normal-equation} yields
\eqref{eq:V24-router-Schur-difference}.
\end{proof}

\begin{corollary}[Record-adapted and record-blind source prices]
\label{cor:V24-adapted-versus-blind}
Let
\(
 E_{\mathcal R}=\bigoplus_{a\in\mathcal R}E
\)
be the coefficient space in which a boundary selector permits an independent
coefficient in every record cell.  The record-adapted joint Gram is
\begin{equation}
 \mathbb G_{\rm ad}
 =\bigoplus_a\mathbb G_a,
 \qquad
 D_{\rm ad}=\bigoplus_aD_a.
 \label{eq:V24-adapted-Gram-short}
\end{equation}
For the diagonal embedding
\(Jv=(v)_{a\in\mathcal R}\),
\begin{equation}
 \boxed{
 D_{\rm blind}-J^*D_{\rm ad}J
 =\mathcal D_{\rm rout}(\mathcal R)\succeq0.}
 \label{eq:V24-short-diagonal-noncommutation}
\end{equation}
Therefore restriction to record-blind coefficients and source shorting do not
commute.  If the record selectors physically occur and adaptive reuse is the
declared target, only the cell shorts \(D_a\) are new source dimensions; the
router term is deterministic record-dependent follower variation.  If those
selectors do not occur, the record-blind short is the correct physical target.
\end{corollary}

\begin{proof}
The Schur complement of the block-diagonal record-adapted Gram is the direct
sum of the cell shorts.  Compressing it by \(J\) gives \(\sum_aD_a\).
Subtract this from \eqref{eq:V24-router-Pythagoras}.
\end{proof}

\begin{countertheorem}[The aggregate joint Gram does not identify ancestry]
\label{cth:V24-aggregate-no-ancestry}
There are two two-cell scalar cards with the same aggregate joint Gram
\begin{equation}
 \boxed{
 M=N=2,
 \qquad
 \Gamma=0,
 \qquad
 \mathbb G=2I_2,
 \qquad
 D_{\rm blind}=2,}
 \label{eq:V24-same-aggregate}
\end{equation}
but opposite ancestry decompositions.

For Card~F, take
\begin{equation}
 (M_1,N_1,\Gamma_1)=(1,1,1),
 \qquad
 (M_2,N_2,\Gamma_2)=(1,1,-1).
 \label{eq:V24-card-F}
\end{equation}
Then
\begin{equation}
 \boxed{
 D_1=D_2=0,
 \qquad
 \mathcal D_{\rm rout}=2.}
 \label{eq:V24-card-F-decomposition}
\end{equation}
Every cell is an exact follower; the entire aggregate short is record-router
dispersion.

For Card~N, take
\begin{equation}
 (M_1,N_1,\Gamma_1)=(1,1,0),
 \qquad
 (M_2,N_2,\Gamma_2)=(1,1,0).
 \label{eq:V24-card-N}
\end{equation}
Then
\begin{equation}
 \boxed{
 D_1=D_2=1,
 \qquad
 \mathcal D_{\rm rout}=0.}
 \label{eq:V24-card-N-decomposition}
\end{equation}
The complete aggregate short is genuine within-cell novelty.  Hence the
unconditional matrices \((M,N,\Gamma)\) do not determine whether their Schur
residual should enlarge the source atlas.
\end{countertheorem}

\begin{proof}
Every displayed cell matrix
\(
 \begin{psmallmatrix}M_a&\Gamma_a\\\Gamma_a&N_a\end{psmallmatrix}
\)
is positive semidefinite.  Direct calculation gives the common aggregate
\eqref{eq:V24-same-aggregate}.  For Card~F the cell decoders are
\(C_1=1\), \(C_2=-1\), and the global decoder is zero, so the router term is
\(1+1=2\).  For Card~N all decoders vanish and each scalar Schur complement is
one.
\end{proof}

\begin{firewall}[A record is not fitted by the source short]
\label{fire:V24-no-record-fitting}
The record map, its reflection-fixed selector algebra, and the declared reuse
class are frozen before \(C_a\), \(D_a\), or
\(\mathcal D_{\rm rout}\) is evaluated.  Refining a record until the fitted
cell decoders become constant, or coarsening it until a desired positive short
appears, is not a physical source theorem.  The two choices correspond to
different observable algebras.
\end{firewall}

\subsection{Record refinement and decoder-dispersion curvature}

Let \(\mathcal S\) be a finer finite physical boundary record and let
\(\pi:\mathcal S\to\mathcal R\) be a deterministic coarsening.  Write
\(b\mapsto a\) when \(\pi(b)=a\).  At the fine level let
\((M_b,\Gamma_b,N_b,C_b,D_b)\) be defined by
\eqref{eq:V24-cell-blocks}--\eqref{eq:V24-cell-decoder-short}.  Aggregate inside
one coarse cell by
\begin{equation}
 M_a=\sum_{b\mapsto a}M_b,
 \qquad
 \Gamma_a=\sum_{b\mapsto a}\Gamma_b,
 \qquad
 N_a=\sum_{b\mapsto a}N_b.
 \label{eq:V24-coarse-cell-aggregation}
\end{equation}

\begin{theorem}[Exact refinement Pythagoras]
\label{thm:V24-refinement-Pythagoras}
For every occurring coarse cell,
\begin{equation}
 \boxed{
 D_a
 =\sum_{b\mapsto a}D_b
 +\sum_{b\mapsto a}(C_b-C_a)^*M_b(C_b-C_a).}
 \label{eq:V24-refinement-Pythagoras}
\end{equation}
Consequently,
\begin{equation}
 \boxed{
 \sum_{a\in\mathcal R}D_a
 -\sum_{b\in\mathcal S}D_b
 =\sum_{a\in\mathcal R}\sum_{b\mapsto a}
   (C_b-C_a)^*M_b(C_b-C_a)
 \succeq0.}
 \label{eq:V24-refinement-monotonicity}
\end{equation}
A finer physically readable record can only decrease the source innovation
required after optimal cellwise follower shorting.  The difference is the
positive follower-router information hidden inside the coarse fibres.
\end{theorem}

\begin{proof}
Apply \Cref{thm:V24-assembly-before-short} to the fine cells contained in one
coarse cell.  Summing the resulting identities over coarse cells gives
\eqref{eq:V24-refinement-monotonicity}.
\end{proof}

\begin{corollary}[Finite record filtration and exact terminal attribution]
\label{cor:V24-record-filtration}
For a finite increasing sequence of physically readable records
\(
 \mathscr R_0\subseteq\cdots\subseteq\mathscr R_L
\), let
\(
 \mathfrak D_j=\sum_{a\in\mathscr R_j}D_{j,a}
\).
Then
\begin{equation}
 \boxed{
 \mathfrak D_0\succeq\mathfrak D_1\succeq\cdots\succeq\mathfrak D_L\succeq0,}
 \label{eq:V24-filtration-monotone}
\end{equation}
and every decrement is the explicit decoder-dispersion Gram in
\eqref{eq:V24-refinement-monotonicity}.  On the full physically readable
record, the terminal short is the unique within-record phase--crossing source
innovation.  No coarser provisional short is priced in addition.
\end{corollary}

\begin{proof}
Apply \Cref{thm:V24-refinement-Pythagoras} at each refinement step and telescope.
\end{proof}

\begin{theorem}[Clean curvature of a record-blind follower]
\label{thm:V24-router-curvature}
Let the cell cards depend twice differentiably on a physical parameter \(t\)
on one fixed record, with constant canonical source supports and ranks so that
the supported decoders are twice differentiable.  Assume on an interval that
every cell remains
an exact follower,
\begin{equation}
 D_a(t)=0,
 \qquad
 B_a(t)=A_a(t)C_a(t),
 \label{eq:V24-cellwise-follower-path}
\end{equation}
and that at \(t=0\) all cell decoders coincide:
\begin{equation}
 C_a(0)=C_0
 \qquad\text{for every }a.
 \label{eq:V24-common-base-decoder}
\end{equation}
Let \(C(t)=M(t)^\dagger\Gamma(t)\) be the record-blind decoder and put
\begin{equation}
 \dot E_a=\dot C_a(0)-\dot C(0).
 \label{eq:V24-relative-decoder-velocity}
\end{equation}
Then
\begin{equation}
 \boxed{
 D_{\rm blind}(0)=0,
 \qquad
 \dot D_{\rm blind}(0)=0,
 \qquad
 \frac12\ddot D_{\rm blind}(0)
 =\sum_a\dot E_a^*M_a(0)\dot E_a\succeq0.}
 \label{eq:V24-router-curvature}
\end{equation}
The first reopening of a common follower relation is therefore a positive
record-router curvature.  If the record selectors are physical, this is not a
new within-cell source birth.
\end{theorem}

\begin{proof}
Under \eqref{eq:V24-cellwise-follower-path},
\Cref{thm:V24-assembly-before-short} gives
\[
 D_{\rm blind}(t)
 =\sum_aE_a(t)^*M_a(t)E_a(t),
 \qquad
 E_a(t)=C_a(t)-C(t).
\]
The common-base condition makes every \(E_a(0)\) zero.  The value and first
derivative therefore vanish.  In the second derivative, every term containing
an undifferentiated \(E_a(0)\) vanishes, leaving precisely
\(2\sum_a\dot E_a^*M_a(0)\dot E_a\).
\end{proof}

\begin{countertheorem}[A static cellwise follower can have positive blind curvature]
\label{cth:V24-router-curvature-control}
Take two scalar record cells with
\begin{equation}
 M_1=M_2=1,
 \qquad
 C_1(t)=t,
 \qquad
 C_2(t)=-t,
 \qquad
 D_1(t)=D_2(t)=0.
 \label{eq:V24-curvature-control-data}
\end{equation}
Then the global decoder is zero and
\begin{equation}
 \boxed{
 D_{\rm blind}(t)=2t^2,
 \qquad
 D_{\rm blind}''(0)=4.}
 \label{eq:V24-curvature-control-result}
\end{equation}
The source is a perfect follower in every cell at every \(t\), yet the
record-blind response has strictly positive curvature.
\end{countertheorem}

\begin{proof}
Substitute the two decoders into
\eqref{eq:V24-router-Pythagoras}; the cell shorts vanish and the two squared
decoder deviations are both \(t^2\).
\end{proof}

\subsection{A four-outcome positive arm instrument}

Version~V23 reconstructed \(\Gamma\) from two arm Grams and four coherent
Grams.  Once one common two-arm instrument physically occurs, the full block
can be reconstructed more economically from four positive outcomes.  This is
stronger than a trine which reconstructs only the trace and off-diagonal arm
coordinates: the Schur short also needs the arm imbalance \(M-N\).

For a Hermitian block
\begin{equation}
 \mathbb G=
 \begin{pmatrix}M&\Gamma\\\Gamma^*&N\end{pmatrix},
 \label{eq:V24-general-joint-block}
\end{equation}
put
\begin{equation}
 \boxed{
 T=M+N,
 \qquad
 X=\Gamma+\Gamma^*,
 \qquad
 Y=i(\Gamma-\Gamma^*),
 \qquad
 Z=M-N.}
 \label{eq:V24-Pauli-coordinates}
\end{equation}
Then
\begin{equation}
 \mathbb G
 =\frac12
 \left(
 I_2\otimes T+\sigma_x\otimes X+\sigma_y\otimes Y
 +\sigma_z\otimes Z
 \right),
 \label{eq:V24-Pauli-expansion}
\end{equation}
where the Pauli matrices have their standard order.

Define four arm effects
\begin{equation}
 \boxed{
 \begin{aligned}
 E_1&=\frac14I_2+\frac18\sigma_x,
 &E_2&=\frac14I_2+\frac18\sigma_y,\\
 E_3&=\frac14I_2+\frac18\sigma_z,
 &E_4&=\frac14I_2-\frac18(\sigma_x+\sigma_y+\sigma_z).
 \end{aligned}}
 \label{eq:V24-four-effects}
\end{equation}
For a block matrix \(G=[G_{ij}]_{i,j=0}^1\), use the arm contraction
\begin{equation}
 \mathcal R_E(G)=\sum_{i,j=0}^1E_{ji}G_{ij}.
 \label{eq:V24-arm-contraction}
\end{equation}

\begin{theorem}[Dimension-minimal four-outcome reconstruction]
\label{thm:V24-four-outcome-tomography}
The effects in \eqref{eq:V24-four-effects} satisfy
\begin{equation}
 E_j\succ0,
 \qquad
 \sum_{j=1}^4E_j=I_2,
 \label{eq:V24-effect-POVM}
\end{equation}
and are linearly independent in \(\operatorname{Herm}(\C^2)\).  If
\(\mathbb G\succeq0\), the four outcome matrices
\begin{equation}
 R_j=\mathcal R_{E_j}(\mathbb G)
 \label{eq:V24-outcome-matrices}
\end{equation}
are positive semidefinite and obey
\begin{equation}
 \boxed{
 \begin{aligned}
 R_1&=\frac14T+\frac18X,
 &R_2&=\frac14T+\frac18Y,\\
 R_3&=\frac14T+\frac18Z,
 &R_4&=\frac14T-\frac18(X+Y+Z).
 \end{aligned}}
 \label{eq:V24-outcome-forward}
\end{equation}
Conversely, the complete block is reconstructed by
\begin{equation}
 \boxed{
 \begin{gathered}
 T=R_1+R_2+R_3+R_4,\\
 X=8R_1-2T,
 \qquad
 Y=8R_2-2T,
 \qquad
 Z=8R_3-2T,\\
 M=\frac12(T+Z),
 \qquad
 N=\frac12(T-Z),
 \qquad
 \Gamma=\frac12(X-iY).
 \end{gathered}}
 \label{eq:V24-outcome-inverse}
\end{equation}
Four scalar arm outcomes are minimum for universal reconstruction of the full
joint block.  Three outcomes can reconstruct at most a three-dimensional real
arm operator system and cannot, in general, determine the Schur short.
\end{theorem}

\begin{proof}
For \(E_1,E_2,E_3\), the trace is \(1/2\) and the determinant is
\(3/64\).  For \(E_4\), the trace is \(1/2\) and the determinant is
\(1/64\).  Hence all four effects are positive definite.  Their sum is the
identity.  If \(\sum_jc_jE_j=0\), the three Pauli coordinates give
\(c_1=c_2=c_3=c_4\), while the trace coordinate gives
\(\sum_jc_j=0\); hence all coefficients vanish.

For any coefficient vector \(v\), the scalar matrix
\(
 [\langle v,G_{ij}v\rangle]_{ij}
\)
is positive semidefinite.  Therefore
\(
 \langle v,\mathcal R_E(G)v\rangle
\)
is the trace of a product of two positive matrices and is nonnegative.  This
proves positivity of every outcome.  Taking Pauli traces yields
\eqref{eq:V24-outcome-forward}, and direct inversion gives
\eqref{eq:V24-outcome-inverse}.

The real vector space \(\operatorname{Herm}(\C^2)\) has dimension four.  A
measurement with fewer than four scalar effects cannot span it, so some arm
coordinate of the joint block is invisible.  Since the Schur complement
depends on all of \(M,N,\Gamma\), universal reconstruction is impossible with
fewer outcomes.
\end{proof}

\begin{countertheorem}[A positive trine does not determine the source short]
\label{cth:V24-trine-no-short}
Let \(F_k\), \(k=0,1,2\), be the equatorial trine effects
\begin{equation}
 F_k=\frac13
 \left(I_2+\cos\frac{2\pi k}{3}\,\sigma_x
              +\sin\frac{2\pi k}{3}\,\sigma_y\right).
 \label{eq:V24-trine-effects}
\end{equation}
The two positive scalar blocks
\begin{equation}
 \mathbb G_0=
 \begin{pmatrix}1&0\\0&1\end{pmatrix},
 \qquad
 \mathbb G_1=
 \begin{pmatrix}3/2&0\\0&1/2\end{pmatrix}
 \label{eq:V24-trine-counterexample-blocks}
\end{equation}
have identical trine outcomes, because they have the same
\(T=2\), \(X=Y=0\).  Nevertheless their source shorts are
\begin{equation}
 \boxed{D_0=1,
 \qquad D_1=\frac12.}
 \label{eq:V24-trine-different-shorts}
\end{equation}
The missing arm coordinate is \(Z=M-N\).  A three-outcome trine is therefore
minimum for a trace/off-diagonal complex quotient, but not for the full
line--crossing Schur target.
\end{countertheorem}

\begin{proof}
Each equatorial effect has zero \(\sigma_z\) coefficient, so its contraction
depends only on \(T,X,Y\).  Both blocks have the same values of those three
coordinates.  Since \(\Gamma=0\), the short equals \(N\), which is one in the
first card and one half in the second.
\end{proof}

\begin{corollary}[Record-cell informationally complete card]
\label{cor:V24-cell-IC-card}
If the boundary selectors \(\mathbf1_{X_a}\) are physical, apply the four arm
effects conditionally in each cell at both the vacuum and one-letter grades:
\begin{equation}
 R_{a,j}^{(g)}=\mathcal R_{E_j}(\mathbb G_{a,g}),
 \qquad
 a\in\mathcal R,
 \quad 1\le j\le4,
 \quad g\in\{0,1\},
 \label{eq:V24-cell-outcomes}
\end{equation}
where \(\mathbb G_{a,0}\) is \eqref{eq:V24-vacuum-cell-Gram} and
\(\mathbb G_{a,1}=\mathbb G_a\).  Then
\eqref{eq:V24-outcome-inverse} reconstructs every
\((p_a,z_a,n_a^w)\) and \((M_a,N_a,\Gamma_a)\), and
\Cref{thm:V24-assembly-before-short} reconstructs both the within-cell source
short and the record-router dispersion.  Summing the outcome matrices over
cells recovers the unconditional four-outcome cards at both grades.
\end{corollary}

\begin{proof}
The selector makes the direct-integral cells orthogonal.  Apply
\Cref{thm:V24-four-outcome-tomography} in each cell and use linearity of the
arm contraction under cell aggregation.
\end{proof}

\subsection{Held-out arm incidence and outward certification}

Let \(E_*\) be one arm effect fixed before the four informationally complete
outcomes are acquired.  The reconstructed block predicts
\begin{equation}
 R_*^{\rm pred}=\mathcal R_{E_*}(\mathbb G).
 \label{eq:V24-heldout-prediction}
\end{equation}

\begin{definition}[Same-instrument held-out residual]
\label{def:V24-heldout-residual}
For a directly acquired held-out positive matrix \(R_*^{\rm dir}\) at either
the vacuum or one-letter grade, put
\begin{equation}
 \boxed{
 \Delta_*^{\rm arm}
 =\|R_*^{\rm dir}-R_*^{\rm pred}\|_{\rm HS}^2.}
 \label{eq:V24-heldout-residual}
\end{equation}
The effect \(E_*\), route, gain convention, boundary record, and source
coefficient bank are frozen before the four fitting outcomes are inspected.
\end{definition}

\begin{proposition}[Exact held-out occurrence test]
\label{prop:V24-heldout-occurrence}
One common arm instrument on one common history cylinder satisfies
\(
 \Delta_*^{\rm arm}=0
\).
A positive residual is a same-instrument or same-history incidence failure; it
is not repaired by projecting the reconstructed block onto the positive cone.
\end{proposition}

\begin{proof}
Both sides of \eqref{eq:V24-heldout-prediction} are the same linear arm
contraction of the same block Gram.  A nonzero difference therefore falsifies
at least one of the common-instrument, common-history, common-gain, or common
source-identification rows.
\end{proof}

Suppose the four exact outcome matrices are enclosed by submitted matrices
\(\widehat R_j\) with operator radii
\begin{equation}
 \|\widehat R_j-R_j\|_{\rm op}\le\eta_j.
 \label{eq:V24-outcome-radii}
\end{equation}
Define
\begin{equation}
 \begin{aligned}
 \eta_T&=\sum_{j=1}^4\eta_j,\\
 \eta_X&=8\eta_1+2\eta_T,
 &\eta_Y&=8\eta_2+2\eta_T,
 &\eta_Z&=8\eta_3+2\eta_T,\\
 \eta_M&=\eta_N=\frac12(\eta_T+\eta_Z),
 &\eta_\Gamma&=\frac12(\eta_X+\eta_Y),\\
 \eta_{\mathbb G}&=\eta_M+\eta_N+2\eta_\Gamma.
 \end{aligned}
 \label{eq:V24-reconstruction-radii}
\end{equation}

\begin{theorem}[Outward four-outcome block and short certificate]
\label{thm:V24-outward-certificate}
The matrices reconstructed from \(\widehat R_j\) by
\eqref{eq:V24-outcome-inverse} satisfy
\begin{equation}
 \boxed{
 \begin{gathered}
 \|\widehat T-T\|\le\eta_T,
 \quad
 \|\widehat X-X\|\le\eta_X,
 \quad
 \|\widehat Y-Y\|\le\eta_Y,
 \quad
 \|\widehat Z-Z\|\le\eta_Z,\\
 \|\widehat M-M\|\le\eta_M,
 \quad
 \|\widehat N-N\|\le\eta_N,
 \quad
 \|\widehat\Gamma-\Gamma\|\le\eta_\Gamma,\\
 \|\widehat{\mathbb G}-\mathbb G\|\le\eta_{\mathbb G}.
 \end{gathered}}
 \label{eq:V24-outward-block-bounds}
\end{equation}
Consequently:
\begin{enumerate}[label=\textup{(C\arabic*)},leftmargin=3em]
\item if
      \(
       \lambda_{\min}(\widehat{\mathbb G})>\eta_{\mathbb G}
      \), then the exact block is strictly positive;
\item if
      \(
       \lambda_{\min}(\widehat{\mathbb G})<-\eta_{\mathbb G}
      \), then no exact positive joint Gram lies in the submitted uncertainty
      ball;
\item otherwise the positivity verdict is unresolved unless an exact positive
      factorization is supplied.
\end{enumerate}

Assume additionally that the exact and reconstructed \(M\) have one verified
canonical support \(P\),
\begin{equation}
 M\succeq s_*P,
 \qquad
 \widehat M\succeq s_*P,
 \qquad
 \max\{\|\Gamma\|,\|\widehat\Gamma\|\}\le\gamma.
 \label{eq:V24-short-domain}
\end{equation}
Then the reconstructed source short
\(
 \widehat D=\widehat N-\widehat\Gamma^*\widehat M^\dagger\widehat\Gamma
\)
satisfies
\begin{equation}
 \boxed{
 \|\widehat D-D\|_{\rm op}
 \le
 \eta_N+\frac{2\gamma}{s_*}\eta_\Gamma
       +\frac{\gamma^2}{s_*^2}\eta_M.}
 \label{eq:V24-outward-short-radius}
\end{equation}
A support change is a typed support event and is not hidden inside this
pseudoinverse radius.
\end{theorem}

\begin{proof}
All reconstruction formulas are linear, so the first seven bounds follow from
the triangle inequality.  Decomposing the block error into its two diagonal
and two off-diagonal blocks gives the conservative final bound.  Weyl's
inequality gives the first three alternatives.  The short estimate is
\Cref{thm:V23-joint-short-transport} applied to the exact and reconstructed
cards.
\end{proof}

\begin{remark}[Exact semidefinite passes]
A physically important card may have a singular positive joint Gram or a zero
cell short.  Such a boundary pass is certified by exact arithmetic, an exact
Gram factorization, or a validated support-and-Schur certificate.  Merely
having an eigenvalue interval which touches zero is \stopstatus{}, not
\passstatus{}.
\end{remark}

\subsection{Record-aware adjacent-cutoff transport}

Let a fine boundary record \(\mathcal R_{n+1}\) map to the old record
\(\mathcal R_n\) by \(\pi_n\), and let \(J_n:E_n\to E_{n+1}\) be the declared
one-letter coefficient isometry.  For each old cell, arm outcome, and grade define the positive transport defect
\begin{equation}
 \varepsilon_{n,a,j,g}^{\rm arm}
 =\left\|
 R_{n,a,j}^{(g)}
 -\sum_{\pi_n(b)=a}J_n^*R_{n+1,b,j}^{(g)}J_n
 \right\|_{\rm op},
 \qquad g\in\{0,1\}.
 \label{eq:V24-arm-transport-defect}
\end{equation}
The linear inversion
\eqref{eq:V24-outcome-inverse} turns these four positive defects into explicit
radii for the transported \(M_a,N_a,\Gamma_a\), while
\Cref{thm:V24-refinement-Pythagoras} separates exact fine-cell novelty from
coarse router dispersion.

\begin{theorem}[Record-aware cofinal phase--crossing alternative]
\label{thm:V24-cofinal-alternative}
Transport every card to one fixed physical boundary and coefficient screen.
Assume:
\begin{enumerate}[label=\textup{(R\arabic*)},leftmargin=3em]
\item the boundary records and their coarsening maps physically occur, and the
      four conditional arm outcomes have summable total defect
      \begin{equation}
       \sum_n\sum_{a,j,g}\varepsilon_{n,a,j,g}^{\rm arm}<\infty;
       \label{eq:V24-summable-arm-defect}
      \end{equation}
\item every occurring cell source Gram has a transported canonical support and
      one positive supported floor on the retained screen;
\item the vacuum cell blocks, line sewing, global partition amplitude, and
      held-out analyzers have summable transported defects and the partition
      amplitude has a nonzero floor;
\item the declared reuse class---record adapted or record blind---is fixed
      across the sequence.
\end{enumerate}
Then the cell Grams, cell shorts, coarse shorts, and router-dispersion matrices
have unique fixed-screen limits.  After the normalized V22 crossing test, one
of the following prioritized branches occurs:
\begin{enumerate}[label=\textup{(V24.\arabic*)},leftmargin=5em]
\item a cellwise negative, skew, sewing, held-out, or block-positivity witness
      gives a finite physical obstruction;
\item a cell short retains a positive protected eigenvalue and yields a genuine
      within-record phase--crossing source survivor;
\item every cell short vanishes asymptotically, but the router dispersion has a
      positive limit; the line-weighted response is a recordwise follower and
      the surviving cost is failure of one common record-blind decoder;
\item both the cell shorts and the router dispersion vanish, giving an
      asymptotic follower on the declared reuse class.
\end{enumerate}
On the record-adapted branch, item~\textup{(V24.3)} does not append a new
source.  On the record-blind branch it remains part of the physical short.
\end{theorem}

\begin{proof}
The four-outcome inversion is a fixed finite linear map, so
\eqref{eq:V24-summable-arm-defect} makes all transported cell blocks Cauchy.
The supported floors and
\Cref{thm:V23-joint-short-transport} transport every cell Schur complement.
The exact identity \eqref{eq:V24-refinement-Pythagoras} then transports the
coarse shorts and router terms without ambiguity.  Sewing, held-out, and
partition-floor hypotheses pass the normalized V22 criterion.  Closedness of
the finite-dimensional positive cone gives the first branch when a strict
witness survives.  Otherwise the nonnegative cell shorts either retain a
positive protected eigenvalue or vanish.  In the latter case the nonnegative
router term either survives or vanishes, giving the last two branches.  The
source attribution follows from \Cref{cor:V24-adapted-versus-blind}.
\end{proof}

\subsection{Exact finite replay}

The standard-library verifier accompanying Version~V24 checks the following
exact cards.
\begin{enumerate}[label=\textup{(E\arabic*)},leftmargin=3.2em]
\item The two ancestry cards of
      \Cref{cth:V24-aggregate-no-ancestry} have the same aggregate block
      \(2I_2\), while Card~F has cell-short/router split \((0,2)\) and Card~N
      has split \((2,0)\).
\item The nontrivial scalar refinement card
      \begin{equation}
       \begin{aligned}
       &(M_1,\Gamma_1,N_1)=(1,1,2),\\
       &(M_2,\Gamma_2,N_2)=(3,-3,4)
       \end{aligned}
       \label{eq:V24-refinement-control-data}
      \end{equation}
      has
      \begin{equation}
       \boxed{
       D_1=D_2=1,
       \qquad
       \mathcal D_{\rm rout}=3,
       \qquad
       D_{\rm blind}=5.}
       \label{eq:V24-refinement-control-result}
      \end{equation}
\item For the V23 complex scalar card
      \(
       M=N=2
      \),
      \(
       \Gamma=(1+3i)/2
      \), the four informationally complete outcomes are
      \begin{equation}
       \boxed{
       R_1=\frac98,
       \qquad
       R_2=\frac58,
       \qquad
       R_3=1,
       \qquad
       R_4=\frac54.}
       \label{eq:V24-IC-control-outcomes}
      \end{equation}
      They reconstruct \(T=4\), \(X=1\), \(Y=-3\), \(Z=0\), and the original
      complex mixed block exactly.
\item The held-out rank-one effect
      \begin{equation}
       E_*=\frac12
       \left(I_2+\frac35\sigma_x+\frac45\sigma_y\right)
       \label{eq:V24-heldout-effect-control}
      \end{equation}
      predicts
      \begin{equation}
       \boxed{R_*^{\rm pred}=\frac{11}{10}.}
       \label{eq:V24-heldout-control-value}
      \end{equation}
\item The two-cell decoder path in
      \eqref{eq:V24-curvature-control-data} has
      \(D_{\rm blind}(t)=2t^2\) and curvature four.
\item The equatorial trine cards of
      \Cref{cth:V24-trine-no-short} have identical three positive outcomes and
      different source shorts \(1\) and \(1/2\).
\end{enumerate}
All rational and Gaussian-rational quantities are represented exactly.
Ordinary and optimized Python executions must produce byte-identical JSON.

\begin{theorem}[Version V24 boundary-record-adapted line--crossing theorem]
\label{thm:V24-terminal}
Versions~V1--V23 are retained.  Version~V24 additionally proves:
\begin{enumerate}[label=\textup{(V24.\arabic*)},leftmargin=5.3em]
\item the record-blind joint source short splits exactly into within-cell
      phase--crossing novelty and positive record-router dispersion;
\item a physically readable refinement decreases the adapted source short by
      the exact decoder-dispersion Gram hidden in the coarse fibres;
\item the unconditional joint Gram alone cannot determine the ancestry split;
\item a cellwise follower can develop positive record-blind curvature solely
      through differential motion of its follower decoder;
\item one fixed four-outcome positive arm instrument reconstructs the full
      joint block \((M,N,\Gamma)\), including the arm imbalance needed by the
      Schur short;
\item four outcomes are minimum for the complete arm operator system, whereas
      a trine can leave different source shorts indistinguishable;
\item one held-out arm effect tests same-instrument occurrence, and explicit
      outward radii propagate from positive outcomes to the block and its
      supported Schur complement;
\item summable conditional-outcome transport gives a cofinal alternative
      separating a true within-cell survivor from a record-router survivor and
      a complete follower.
\end{enumerate}
No item populates the actual RPESM boundary record, conditional arm outcomes,
line scalarization, gauge-history crossing, or adjacent-cutoff values.
\end{theorem}

\begin{proof}
Items~\textup{(V24.1)--(V24.3)} are
\Cref{thm:V24-assembly-before-short,cor:V24-adapted-versus-blind,cth:V24-aggregate-no-ancestry}.
Item~\textup{(V24.4)} is
\Cref{thm:V24-router-curvature}.  Items~\textup{(V24.5)--(V24.6)} are
\Cref{thm:V24-four-outcome-tomography,cth:V24-trine-no-short}.  Item
\textup{(V24.7)} is
\Cref{prop:V24-heldout-occurrence,thm:V24-outward-certificate}, and item
\textup{(V24.8)} is \Cref{thm:V24-cofinal-alternative}.
\end{proof}

\begin{assessment}[Version V24 verdict]
\label{ass:V24-verdict}
The current result is
\begin{equation}
 \boxed{
 \begin{gathered}
 \passstatus\\[-0.15em]
 \text{record-cell assembly, exact router Pythagoras, refinement monotonicity;}\\
 \text{four-outcome positive tomography, held-out incidence, outward transport}
 \\[0.45em]
 \obsstatus\\[-0.15em]
 \text{an aggregate Schur short priced before the physical record is fixed;}\\
 \text{a record-router residual charged as a second source birth;}\\
 \text{or a three-outcome trace/off-diagonal card used as a full Schur card}
 \\[0.45em]
 \stopstatus\\[-0.15em]
 \text{first physically populated record-indexed RPESM arm instrument,}\\
 \text{held-out analyzer, line sewing, and adjacent boundary-record map.}
 \end{gathered}}
 \label{eq:V24-verdict}
\end{equation}
The continuation goes upstream of Version~V23's source price without reopening
the line, crossing, exterior, Duhamel, or continuum compilers.  It determines
which part of a positive joint short is a true within-record source innovation
and which part is only the cost of suppressing an already readable boundary
router.
\end{assessment}

\begin{acquisition}[Version V25: execute one record-indexed RPESM arm card]
\label{acq:V24-next}
Preserve Versions~V1--V24.  On one protected RPESM history and one adjacent
cutoff, freeze the reflection-fixed boundary record \(r_0\), its physically
available selector algebra, the V22 amplitude semantics, the one-letter source
bank, and the common arm instrument.  Populate
\begin{equation}
 \boxed{
 \begin{gathered}
 \left(
 r_0,
 \{R_{a,j}^{(g)}:a\in\mathcal R,1\le j\le4,g\in\{0,1\}\},
 \right.\\
 \left.
 \{R_{a,*}^{(g)}:a\in\mathcal R,g\in\{0,1\}\},
 m,\Delta_{\rm sew}^{\rm line},
 \text{supports and outward radii}
 \right)
 \end{gathered}.}
 \label{eq:V24-next-card}
\end{equation}
Use \eqref{eq:V24-outcome-inverse} to reconstruct every vacuum cell
\((p_a,z_a,n_a^w)\) and every one-letter cell
\((M_a,N_a,\Gamma_a)\).  Evaluate the cell shorts, the record-router
Pythagoras, the V22 normalized cell criterion, and the held-out residual.  If a
cell short is positive, append only its polar source.  If all cell shorts
vanish but the router term is positive, retain the physical record-dependent
follower map and do not charge a new source on the record-adapted branch.  If
the boundary selectors do not physically occur, use the record-blind card and
state that stronger reuse class explicitly.

Repeat the four positive conditional outcomes, partition amplitude, sewing
row, and record map at one adjacent cutoff.  A fitted record refinement, a
nearest-positive replacement, or a trine lacking the arm imbalance is not an
accepted substitute.  Until this card lands, do not add another phase,
crossing, exterior, localizer, Duhamel, or continuum compiler.
\end{acquisition}

\section*{V24 exact replay and append-only record}
\addcontentsline{toc}{section}{V24 exact replay and append-only record}
The accompanying verifier checks the assembly-before-short identity, the two
indistinguishable aggregate ancestry cards, exact record refinement, decoder
curvature, the rational four-outcome informationally complete instrument, the
held-out analyzer, and the trine insufficiency counterexample.  It is a theorem
replay and not a measurement of the active RPESM boundary record or fermionic
line.

The complete Version~V23 source preceding its sole final active
\verb|\end{document}| is preserved byte for byte.  Its preterminal SHA--256
digest is recorded in the validation manifest.  A later continuation must
remove only the final active document-closing command, append new material
immediately before it, restore one terminal command, and increment the filename
to end in \texttt{\_V25.tex}.

% =============================================================================
% END APPENDED VERSION V24 MATERIAL
% =============================================================================


% =============================================================================
% BEGIN APPENDED VERSION V25 MATERIAL
% =============================================================================

\clearpage
\hypersetup{
 pdftitle={Renewal Einstein--Standard--Model Physical Source Metric, Duhamel Ancestry and Quantum Continuum Programme V25},
 pdfsubject={Primitive-complete line--crossing ancestry, positive square acquisition, and the first numerical-value stop}
}
\pdfinfo{
 /Title (Renewal Einstein--Standard--Model Physical Source Metric, Duhamel Ancestry and Quantum Continuum Programme V25)
 /Subject (Primitive-complete line--crossing ancestry, positive square acquisition, and the first numerical-value stop)
}
\section{Version V25: primitive-complete line--crossing ancestry, positive square acquisition, and the first numerical-value stop}
\label{sec:V25-main}

Versions~V23--V24 correctly showed that determinant-line information and the
reflected crossing must be acquired on one history and that record-dependent
follower maps must be assembled before a source price is assigned.  The
upstream audit of the inherited RPESM population theorem changes one part of
the stopping semantics.  The retained-terminal regulator already places every
predeclared finite source, line, reflected, mixed boson--fermion, boundary, and
held-out writer on one common event law.  Therefore the four-arm tomography of
Version~V24 is an optional positive acquisition coordinate, not an additional
existence postulate, whenever its arm maps are deterministic functions of that
already populated common-history packet.

There is, however, an earlier ancestry short which Version~V24 did not include.
The line--crossing pair must first be shorted by the \emph{complete saturated
prior primitive bank}.  Otherwise an already occurring primitive follower can
be counted as a new phase--crossing source.  The correct finite order is
\begin{equation}
 \boxed{
 \begin{gathered}
 \text{same-history saturated primitive, crossing, and line-weighted arms}\\
 \Downarrow\\
 \text{complete three-arm positive Gram}\\
 \Downarrow\\
 \text{prior-primitive short}\\
 \Downarrow\\
 \text{record-cell crossing short and router Pythagoras}\\
 \Downarrow\\
 \text{semigroup-cyclic short and only then dynamic source price}.
 \end{gathered}}
 \label{eq:V25-order}
\end{equation}
This version proves the complete finite algebra behind the first four arrows.
It also replaces the proposed new arm instrument by deterministic positive
square writers of the already occurring three-arm synthesis.  The remaining
RPESM task is numerical evaluation, supported outward certification, and
adjacent-cutoff transport on one selected physical card.

% -----------------------------------------------------------------------------
\subsection{The primitive-complete three-arm Gram}
\label{subsec:V25-three-arm}
% -----------------------------------------------------------------------------

Let \(\mathcal H\) be one finite reflected source Hilbert space.  Let
\begin{equation}
 C:E_0\longrightarrow\mathcal H,
 \qquad
 A:E\longrightarrow\mathcal H,
 \qquad
 B:E\longrightarrow\mathcal H
 \label{eq:V25-three-syntheses}
\end{equation}
be, respectively, the saturated prior primitive synthesis, the phase-quenched
crossing synthesis, and the line-weighted crossing synthesis selected by one
fixed amplitude semantics.  Define
\begin{equation}
 \begin{gathered}
 S=C^*C,
 \qquad U=C^*A,
 \qquad V=C^*B,\\
 M=A^*A,
 \qquad \Gamma=A^*B,
 \qquad N=B^*B,
 \end{gathered}
 \label{eq:V25-six-blocks}
\end{equation}
and the complete block Gram
\begin{equation}
 \boxed{
 \mathbb G_{C,A,B}
 =
 \begin{pmatrix}
 S&U&V\\
 U^*&M&\Gamma\\
 V^*&\Gamma^*&N
 \end{pmatrix}
 =(C\ \ A\ \ B)^*(C\ \ A\ \ B)
 \succeq0.}
 \label{eq:V25-three-arm-Gram}
\end{equation}
Put
\begin{equation}
 P_C=CS^\dagger C^*,
 \qquad Q_C=I-P_C,
 \qquad A_C=Q_CA,
 \qquad B_C=Q_CB.
 \label{eq:V25-primitive-residual-arms}
\end{equation}
The primitive-shorted arm blocks are
\begin{equation}
 \boxed{
 \begin{aligned}
 M_C&=M-U^*S^\dagger U=A_C^*A_C,\\
 \Gamma_C&=\Gamma-U^*S^\dagger V=A_C^*B_C,\\
 N_C&=N-V^*S^\dagger V=B_C^*B_C.
 \end{aligned}}
 \label{eq:V25-primitive-shorted-blocks}
\end{equation}

\begin{theorem}[Nested primitive--crossing ancestry and exact mass allocation]
\label{thm:V25-nested-short}
Let the syntheses in \eqref{eq:V25-three-syntheses} be physically assembled on
one Hilbert space.  Then
\begin{equation}
 P_{A\mid C}:=A_CM_C^\dagger A_C^*
 \label{eq:V25-residual-A-projection}
\end{equation}
is the orthogonal projection onto \(\Ran A_C\),
\begin{equation}
 \boxed{P_{(C,A)}=P_C+P_{A\mid C}}
 \label{eq:V25-nested-range-projection}
\end{equation}
is the orthogonal projection onto \(\Ran(C\ \ A)\), and the final
line--crossing ancestry residual is
\begin{equation}
 \boxed{
 D_{B\mid(C,A)}
 :=N_C-\Gamma_C^*M_C^\dagger\Gamma_C
 =B^*(I-P_{(C,A)})B
 \succeq0.}
 \label{eq:V25-final-nested-short}
\end{equation}
The complete line-weighted mass has the orthogonal allocation
\begin{equation}
 \boxed{
 N
 =V^*S^\dagger V
 +\Gamma_C^*M_C^\dagger\Gamma_C
 +D_{B\mid(C,A)}.}
 \label{eq:V25-three-level-mass}
\end{equation}
The three terms are, in order, the prior-primitive follower, the crossing
follower after primitive removal, and the only remaining linear source
innovation.
\end{theorem}

\begin{proof}
The Moore--Penrose range formula makes \(P_C\) the orthogonal projection onto
\(\Ran C\).  Hence \(A_C\) and \(B_C\) take values in
\((\Ran C)^\perp\), and \eqref{eq:V25-primitive-shorted-blocks} follows by
expansion.  The same range formula applied to \(A_C\) makes
\(P_{A\mid C}\) the projection onto \(\Ran A_C\).  The two projections are
orthogonal.  Every vector in \(\Ran(C\ \ A)\) is the sum of a vector in
\(\Ran C\) and the residual part of a vector in \(\Ran A\), proving
\eqref{eq:V25-nested-range-projection}.  Therefore
\[
 B^*(I-P_{(C,A)})B
 =B_C^*(I-P_{A\mid C})B_C
 =N_C-\Gamma_C^*M_C^\dagger\Gamma_C,
\]
which proves \eqref{eq:V25-final-nested-short}.  Finally decompose \(B\)
orthogonally into
\[
 P_CB,
 \qquad P_{A\mid C}B_C,
 \qquad (I-P_{(C,A)})B.
\]
Taking their coefficient Grams gives \eqref{eq:V25-three-level-mass}.
\end{proof}

\begin{corollary}[The Version V24 two-arm short is an upper ancestry card]
\label{cor:V25-V24-upper}
Put \(P_A=AM^\dagger A^*\).  The two-arm residual formed before the prior
primitive short is
\begin{equation}
 D_{B\mid A}:=N-\Gamma^*M^\dagger\Gamma
 =B^*(I-P_A)B.
 \label{eq:V25-premature-two-arm-short}
\end{equation}
It obeys
\begin{equation}
 \boxed{
 D_{B\mid A}-D_{B\mid(C,A)}
 =B^*(P_{(C,A)}-P_A)B
 \succeq0.}
 \label{eq:V25-premature-overprice}
\end{equation}
Equality holds exactly when
\((P_{(C,A)}-P_A)B=0\).  Thus a positive Version~V24 short is a valid
upper card for the final source novelty, but it is not itself the final source
price until the complete prior primitive bank has been removed.
\end{corollary}

\begin{proof}
The projections satisfy \(P_A\preceq P_{(C,A)}\) because
\(\Ran A\subseteq\Ran(C\ \ A)\).  Their difference is the orthogonal
projection onto \(\Ran(C\ \ A)\cap(\Ran A)^\perp\).  Subtract the two
projection formulas.  A positive Gram \(B^*QB\) vanishes exactly when
\(QB=0\).
\end{proof}

\begin{firewall}[No replacement of the physical primitive bank]
\label{fire:V25-complete-primitive}
The synthesis \(C\) is the complete saturated prior primitive bank on the
selected physical card.  It is not replaced by the subset most correlated
with \(A\) or \(B\), and it is not enlarged after observing the residual.
Changing \(C\) changes the declared ancestry question.  The final cyclic
Duhamel short remains a later operation and is not identified with
\eqref{eq:V25-final-nested-short}.
\end{firewall}

% -----------------------------------------------------------------------------
\subsection{Record adaptation after the primitive short}
\label{subsec:V25-record}
% -----------------------------------------------------------------------------

Let \(a\in\mathcal R\) range over one physically readable finite boundary
record.  On cell \(a\), let \((C_a,A_a,B_a)\) be the corresponding three
syntheses.  Short the complete primitive bank first and write
\begin{equation}
 \begin{gathered}
 \overline A_a=(I-P_{C_a})A_a,
 \qquad
 \overline B_a=(I-P_{C_a})B_a,\\
 \overline M_a=\overline A_a^*\overline A_a,
 \qquad
 \overline\Gamma_a=\overline A_a^*\overline B_a,
 \qquad
 \overline N_a=\overline B_a^*\overline B_a,\\
 L_a=\overline M_a^\dagger\overline\Gamma_a,
 \qquad
 D_a^{\rm fin}
 =\overline N_a-
  \overline\Gamma_a^*\overline M_a^\dagger\overline\Gamma_a.
 \end{gathered}
 \label{eq:V25-cell-residualized-data}
\end{equation}
Assemble
\begin{equation}
 \overline M=\sum_a\overline M_a,
 \qquad
 \overline\Gamma=\sum_a\overline\Gamma_a,
 \qquad
 \overline N=\sum_a\overline N_a,
 \qquad
 L=\overline M^\dagger\overline\Gamma.
 \label{eq:V25-record-assembled-data}
\end{equation}

\begin{theorem}[Primitive-complete record-router Pythagoras]
\label{thm:V25-record-router}
Assume the boundary selectors are physical and the primitive bank is
record-adapted, so that its global range is the orthogonal direct sum of the
cell primitive ranges.  Then the record-blind final short after primitive
removal is
\begin{equation}
 \boxed{
 \begin{aligned}
 D_{\rm blind}^{\rm fin}
 &:={\overline N}
 -\overline\Gamma^*\overline M^\dagger\overline\Gamma\\
 &=\sum_aD_a^{\rm fin}
 +\underbrace{\sum_a
 (L_a-L)^*\overline M_a(L_a-L)}_
 {\mathcal D_{\rm rout}^{\rm fin}(\mathcal R)\succeq0}.
 \end{aligned}}
 \label{eq:V25-final-router-Pythagoras}
\end{equation}
Moreover the complete unshorted line mass has the allocation
\begin{equation}
 \boxed{
 \begin{aligned}
 \sum_aN_a
 ={}&\sum_a V_a^*S_a^\dagger V_a\\
 &+\sum_a
 \overline\Gamma_a^*\overline M_a^\dagger\overline\Gamma_a
 +\sum_aD_a^{\rm fin},
 \end{aligned}}
 \label{eq:V25-record-adapted-mass}
\end{equation}
where \(S_a=C_a^*C_a\) and \(V_a=C_a^*B_a\).  If the crossing coefficient is
forced to be record blind, the second line in
\eqref{eq:V25-record-adapted-mass} is replaced by
\begin{equation}
 \overline\Gamma^*\overline M^\dagger\overline\Gamma
 +\mathcal D_{\rm rout}^{\rm fin}(\mathcal R)
 +\sum_aD_a^{\rm fin}.
 \label{eq:V25-record-blind-mass}
\end{equation}
Thus prior-primitive follower mass, record-dependent crossing follower mass,
router dispersion, and final novelty are four distinct nonduplicating rows.
\end{theorem}

\begin{proof}
Orthogonal direct-sum primitive shorting gives the cell residuals in
\eqref{eq:V25-cell-residualized-data}.  Expand
\[
 \sum_a(L_a-L)^*\overline M_a(L_a-L).
\]
The cross terms vanish because
\[
 \sum_a\overline M_a(L_a-L)
 =\sum_a\overline\Gamma_a-\overline M L=0
\]
on the supported assembled range.  The remaining terms are
\[
 \sum_a\overline\Gamma_a^*\overline M_a^\dagger
          \overline\Gamma_a
 -\overline\Gamma^*\overline M^\dagger\overline\Gamma.
\]
Adding \(\sum_aD_a^{\rm fin}\) proves
\eqref{eq:V25-final-router-Pythagoras}.  Applying
\eqref{eq:V25-three-level-mass} in every cell and summing proves
\eqref{eq:V25-record-adapted-mass}; substitution of the router identity gives
\eqref{eq:V25-record-blind-mass}.
\end{proof}

\begin{corollary}[Correct source-pricing order on a readable boundary record]
\label{cor:V25-source-price-order}
On the record-adapted branch, only \(\bigoplus_aD_a^{\rm fin}\) may enlarge
the source atlas.  The primitive and crossing follower terms are charged to
their original sources, and
\(\mathcal D_{\rm rout}^{\rm fin}(\mathcal R)\) records the cost of suppressing
a physically readable decoder.  On a declared record-blind reuse class, the
router term remains part of that stronger target's residual but is not
reinterpreted as a new history occurrence.
\end{corollary}

% -----------------------------------------------------------------------------
\subsection{Instrument-free positive acquisition of the complete three-arm Gram}
\label{subsec:V25-positive-square}
% -----------------------------------------------------------------------------

The common-history population theorem licenses deterministic finite writers of
an already occurring source synthesis.  A separate interferometric instrument
is therefore unnecessary for the present finite Gram target.  The complete
block can be reconstructed from positive norm squares alone.

When the three coefficient spaces differ, extend them by zero to the canonical
labelled direct sum
\begin{equation}
 \mathcal E=E_0\oplus E\oplus E
 \label{eq:V25-common-coefficient-universe}
\end{equation}
and retain the corresponding support projections.  In the following theorem
\(X,Y\) denote any two of the zero-extended syntheses \(C,A,B\), acting on
\(\mathcal E\).  Put
\begin{equation}
 G_X=X^*X,
 \qquad
 Q_{XY}^{(+)}=(X+Y)^*(X+Y),
 \qquad
 Q_{XY}^{(i)}=(X+iY)^*(X+iY).
 \label{eq:V25-pair-square-writers}
\end{equation}
Every displayed matrix is positive semidefinite.

\begin{theorem}[Nine positive square panels reconstruct the three-arm Gram]
\label{thm:V25-nine-positive-panels}
For each unordered pair \(\{X,Y\}\subseteq\{C,A,B\}\), define
\begin{equation}
 H_{XY}=\frac12
 \left(Q_{XY}^{(+)}-G_X-G_Y\right),
 \qquad
 J_{XY}=\frac12
 \left(G_X+G_Y-Q_{XY}^{(i)}\right).
 \label{eq:V25-pair-reconstruction}
\end{equation}
Then
\begin{equation}
 \boxed{X^*Y=H_{XY}+iJ_{XY}.}
 \label{eq:V25-cross-reconstruction}
\end{equation}
Consequently the nine positive matrices
\begin{equation}
 \boxed{
 G_C,G_A,G_B,
 Q_{CA}^{(+)},Q_{CA}^{(i)},
 Q_{CB}^{(+)},Q_{CB}^{(i)},
 Q_{AB}^{(+)},Q_{AB}^{(i)}}
 \label{eq:V25-nine-card}
\end{equation}
reconstruct the complete block
\(\mathbb G_{C,A,B}\), including all three complex mixed rows.

If all three arms share one coefficient screen and the acquisition is required
to be linear in an arm effect, nine real linearly independent positive arm
effects are minimum: \(\operatorname{Herm}(\C^3)\) has real dimension nine.
For unequal coefficient spaces, the zero-extension and supported-corner
construction above gives the same finite result without identifying unrelated
coefficient labels.
\end{theorem}

\begin{proof}
Expansion gives
\[
 Q_{XY}^{(+)}=G_X+G_Y+X^*Y+Y^*X
\]
and
\[
 Q_{XY}^{(i)}=G_X+G_Y+iX^*Y-iY^*X.
\]
Thus \(H_{XY}=(X^*Y+Y^*X)/2\) and
\(J_{XY}=(X^*Y-Y^*X)/(2i)\), proving
\eqref{eq:V25-cross-reconstruction}.  Apply the formula to the three pairs.
For the minimum statement, any arm-effect acquisition is a real-linear map on
\(\operatorname{Herm}(\C^3)\).  Injectivity requires at least its real
dimension, nine.  A positive spanning family exists because every Hermitian
basis can be shifted and rescaled into the positive cone.
\end{proof}



\begin{corollary}[The arm instrument is not a new RPESM occurrence row]
\label{cor:V25-no-new-instrument}
Suppose \(C,A,B\) are deterministic finite writer syntheses of one already
populated retained-terminal RPESM history, with their amplitude semantics,
reflection, boundary record, and coefficient supports frozen before the
values are inspected.  Then every matrix in \eqref{eq:V25-nine-card} is a
deterministic positive square writer of that same event.  It requires no new
randomizer, branch law, or arm instrument.

What remains physical is the numerical evaluation of these writers, their
common gain and calibration, the exact support incidences, and their outward
radii.  If the maps defining \(A\) or \(B\) are not deterministic writers of
the populated packet, this corollary does not manufacture them.
\end{corollary}

\begin{proof}
A deterministic complex linear combination of finitely many same-history
Hilbert-valued writers is again a same-history writer.  Its squared norm is a
nonnegative scalar writer, and its coefficient Gram is positive.  The
population premise places every predeclared finite deterministic writer in the
same event grammar.  No stochastic arm selection is used.
\end{proof}

\begin{proposition}[Held-out three-arm coherence test]
\label{prop:V25-heldout-three-arm}
Fix \(\alpha,\beta,\gamma\in\C\) before fitting.  The reconstructed block
predicts the positive matrix
\begin{equation}
 \boxed{
 Q_{\alpha,\beta,\gamma}^{\rm pred}
 =
 (\alpha C+\beta A+\gamma B)^*
 (\alpha C+\beta A+\gamma B),}
 \label{eq:V25-heldout-prediction}
\end{equation}
whose expansion is determined by the six blocks in
\eqref{eq:V25-six-blocks}.  A direct same-event reading gives
\(Q_{\alpha,\beta,\gamma}^{\rm dir}\).  The residual
\begin{equation}
 \boxed{
 \Delta_{\alpha,\beta,\gamma}^{\rm occ}
 =\left\|
 Q_{\alpha,\beta,\gamma}^{\rm dir}
 -Q_{\alpha,\beta,\gamma}^{\rm pred}
 \right\|_{\rm HS}^{2}}
 \label{eq:V25-heldout-residual}
\end{equation}
is an occurrence, common-gain, or coefficient-identification test.  A
positive value is not repaired by projecting the reconstructed three-arm block
to the positive cone.
\end{proposition}

% -----------------------------------------------------------------------------
\subsection{Validated nested-short radius and cutoff transport}
\label{subsec:V25-outward}
% -----------------------------------------------------------------------------

\begin{theorem}[Two-stage supported Schur perturbation bound]
\label{thm:V25-two-stage-radius}
Let exact and submitted six-block packets have common canonical supports and
satisfy
\begin{equation}
 S,\widehat S\succeq s_*P_0,
 \qquad
 M_C,\widehat M_C\succeq m_*P_1,
 \qquad s_*,m_*>0.
 \label{eq:V25-two-floor-assumption}
\end{equation}
Assume
\begin{equation}
 \begin{gathered}
 \max\{\|U\|,\|\widehat U\|\}\le u,
 \qquad
 \max\{\|V\|,\|\widehat V\|\}\le v,\\
 \max\{\|\Gamma_C\|,\|\widehat\Gamma_C\|\}\le\gamma,
 \end{gathered}
 \label{eq:V25-block-norm-windows}
\end{equation}
and write the six block errors as
\(\epsilon_S,\epsilon_U,\epsilon_V,\epsilon_M,\epsilon_\Gamma,
\epsilon_N\).  Define
\begin{equation}
 \boxed{
 \epsilon_{M_C}
 =\epsilon_M+\frac{2u}{s_*}\epsilon_U
  +\frac{u^2}{s_*^2}\epsilon_S,}
 \label{eq:V25-MC-radius}
\end{equation}
\begin{equation}
 \boxed{
 \epsilon_{\Gamma_C}
 =\epsilon_\Gamma
  +\frac{v}{s_*}\epsilon_U
  +\frac{u}{s_*}\epsilon_V
  +\frac{uv}{s_*^2}\epsilon_S,}
 \label{eq:V25-GammaC-radius}
\end{equation}
\begin{equation}
 \boxed{
 \epsilon_{N_C}
 =\epsilon_N+\frac{2v}{s_*}\epsilon_V
  +\frac{v^2}{s_*^2}\epsilon_S.}
 \label{eq:V25-NC-radius}
\end{equation}
Then
\begin{equation}
 \boxed{
 \|\widehat D_{B\mid(C,A)}-D_{B\mid(C,A)}\|
 \le
 \epsilon_D
 :=\epsilon_{N_C}
  +\frac{2\gamma}{m_*}\epsilon_{\Gamma_C}
  +\frac{\gamma^2}{m_*^2}\epsilon_{M_C}.}
 \label{eq:V25-final-radius}
\end{equation}
Every block error can be propagated linearly from the nine positive square
panels.  A change of either canonical support is a typed support event and is
not covered by \eqref{eq:V25-final-radius}.
\end{theorem}

\begin{proof}
On the fixed support,
\[
 \widehat S^\dagger-S^\dagger
 =-\widehat S^\dagger(\widehat S-S)S^\dagger,
 \qquad
 \|S^\dagger\|,\|\widehat S^\dagger\|\le s_*^{-1}.
\]
Expand the three differences in
\eqref{eq:V25-primitive-shorted-blocks}, insert and subtract one factor at a
time, and apply submultiplicativity.  This gives
\eqref{eq:V25-MC-radius}--\eqref{eq:V25-NC-radius}.  Repeat the same supported
pseudoinverse identity for
\[
 D=N_C-\Gamma_C^*M_C^\dagger\Gamma_C
\]
with \(\|M_C^\dagger\|,\|\widehat M_C^\dagger\|\le m_*^{-1}\).
The resulting three-term bound is \eqref{eq:V25-final-radius}.  The
reconstruction formulas of \Cref{thm:V25-nine-positive-panels} are linear, so
positive-panel radii propagate by the triangle inequality.
\end{proof}

\begin{corollary}[Cofinal transport of the primitive-complete source short]
\label{cor:V25-cofinal}
After transport to a fixed physical record and coefficient screen, suppose the
canonical primitive and residual crossing supports stabilize, the floors in
\eqref{eq:V25-two-floor-assumption} remain positive, and the adjacent errors
of all six blocks are summable.  Then the primitive-shorted blocks and
\(D_{B\mid(C,A)}\) are Cauchy in operator norm and have unique limits on that
screen.

If the boundary record is retained, the cell shorts and router term converge
separately when their transported positive square panels are summable.  A
positive limiting cell short is a source survivor; a vanishing cell short with
a positive router limit is a record-dependent follower survivor; and
vanishing of both gives asymptotic linear followerhood before the later cyclic
short.
\end{corollary}

\begin{proof}
Apply \Cref{thm:V25-two-stage-radius} to every adjacent pair and telescope.
The record statement follows from the exact identity
\eqref{eq:V25-final-router-Pythagoras} and continuity of supported
Moore--Penrose inversion under the stated floors.
\end{proof}

% -----------------------------------------------------------------------------
\subsection{Exact two-cell replay and the size of premature source pricing}
\label{subsec:V25-control}
% -----------------------------------------------------------------------------

Let two boundary cells have copies of \(\C^3\) with standard basis
\(e_1,e_2,e_3\).  In both cells set
\begin{equation}
 C_\pm(1)=e_1,
 \qquad
 A_\pm(1)=e_1+e_2,
 \qquad
 B_\pm(1)=e_1\pm i e_2+e_3.
 \label{eq:V25-control-vectors}
\end{equation}

\begin{proposition}[Exact primitive, crossing, router, and novelty allocation]
\label{prop:V25-exact-control}
For each sign,
\begin{equation}
 \boxed{
 S_\pm=1,
 \quad U_\pm=V_\pm=1,
 \quad M_\pm=2,
 \quad \Gamma_\pm=1\pm i,
 \quad N_\pm=3.}
 \label{eq:V25-control-full-blocks}
\end{equation}
The complete three-arm Gram is positive definite with determinant one.  After
the primitive short,
\begin{equation}
 \boxed{
 (M_{C,\pm},\Gamma_{C,\pm},N_{C,\pm})
 =(1,\pm i,2),
 \qquad
 D_\pm^{\rm fin}=1.}
 \label{eq:V25-control-cell-short}
\end{equation}
The premature two-arm short in either cell is two.

For one record-blind coefficient assembled across the two cells,
\begin{equation}
 \boxed{
 M=4,
 \qquad\Gamma=2,
 \qquad N=6,
 \qquad D_{B\mid A}=5.}
 \label{eq:V25-control-premature-aggregate}
\end{equation}
After the complete primitive short,
\begin{equation}
 \boxed{
 \overline M=2,
 \qquad\overline\Gamma=0,
 \qquad\overline N=4,
 \qquad D_{\rm blind}^{\rm fin}=4.}
 \label{eq:V25-control-primitive-aggregate}
\end{equation}
The record-adapted final novelty and router term are
\begin{equation}
 \boxed{
 \sum_{\pm}D_\pm^{\rm fin}=2,
 \qquad
 \mathcal D_{\rm rout}^{\rm fin}=2.}
 \label{eq:V25-control-router}
\end{equation}
Thus the exact line mass has the two equivalent allocations
\begin{equation}
 \boxed{
 6=2_{\rm primitive}+2_{\rm crossing\ follower}+2_{\rm novelty}}
 \label{eq:V25-control-adapted-allocation}
\end{equation}
and
\begin{equation}
 \boxed{
 6=2_{\rm primitive}+4_{\rm record\mbox{-}blind\ residual},
 \qquad
 4=2_{\rm router}+2_{\rm novelty}.}
 \label{eq:V25-control-blind-allocation}
\end{equation}
The Version~V24 short formed before primitive removal overprices the
record-blind final residual by one and overprices the record-adapted new source
by three.
\end{proposition}

\begin{proof}
All identities follow by direct inner products.  The column matrix
\((C_+\ A_+\ B_+)\) is upper triangular with determinant one, so its Gram has
determinant one; the minus card is its complex conjugate.  Subtracting the
\(e_1\) primitive leaves \(e_2\) and \(\pm i e_2+e_3\), which proves
\eqref{eq:V25-control-cell-short}.  Summation cancels the two residual decoders
\(\pm i\), producing \eqref{eq:V25-control-primitive-aggregate}; the router
identity then gives \eqref{eq:V25-control-router}.
\end{proof}

For the plus cell, the nine positive square values are
\begin{equation}
 \boxed{
 \begin{array}{c|ccc}
 &G_C&G_A&G_B\\ \hline
 \text{value}&1&2&3
 \end{array}}
 \label{eq:V25-control-diagonal-squares}
\end{equation}
and
\begin{equation}
 \boxed{
 \begin{array}{c|cc}
 \{X,Y\}&Q_{XY}^{(+)}&Q_{XY}^{(i)}\\ \hline
 \{C,A\}&5&3\\
 \{C,B\}&6&4\\
 \{A,B\}&7&3
 \end{array}.}
 \label{eq:V25-control-pair-squares}
\end{equation}
They reconstruct \(U=V=1\) and \(\Gamma=1+i\) exactly.  The held-out
three-arm writer with coefficients \((1,-1,i)\) has
\begin{equation}
 \boxed{
 \|C_+-A_++iB_+\|^2=6.}
 \label{eq:V25-control-heldout}
\end{equation}
Changing only the reported held-out value by \(1/7\) gives the exact
occurrence residual \(1/49\).

\begin{theorem}[Version V25 primitive-complete line--crossing theorem]
\label{thm:V25-terminal}
Versions~V1--V24 are retained.  Version~V25 additionally proves:
\begin{enumerate}[label=\textup{(V25.\arabic*)},leftmargin=5.3em]
\item the final line--crossing source is the short of the line-weighted arm by
      the joint range of the complete prior primitive and crossing arms;
\item the Version~V24 two-arm residual is a positive upper ancestry card and
      can strictly overprice an already occurring primitive follower;
\item after primitive removal, the record-blind residual splits exactly into
      within-cell novelty and record-router dispersion;
\item the complete primitive, crossing, and line-weighted Gram is reconstructed
      from nine deterministic positive square panels and one held-out
      three-arm writer;
\item on the inherited population-complete RPESM history, those panels require
      no additional arm instrument, although their selected numerical values,
      supports, calibration, and outward radii remain physical acquisitions;
\item a two-stage supported perturbation theorem transports the complete final
      short without hiding primitive or residual-crossing support changes;
\item the exact two-cell card separates primitive follower, crossing follower,
      router dispersion, and final novelty and exhibits strict premature
      overpricing.
\end{enumerate}
No item supplies the numerical three-arm matrices of a selected active RPESM
boundary card or its adjacent-cutoff margins.
\end{theorem}

\begin{proof}
Items~\textup{(V25.1)--(V25.2)} are
\Cref{thm:V25-nested-short,cor:V25-V24-upper}.  Item~\textup{(V25.3)} is
\Cref{thm:V25-record-router}.  Items~\textup{(V25.4)--(V25.5)} are
\Cref{thm:V25-nine-positive-panels,cor:V25-no-new-instrument,prop:V25-heldout-three-arm}.
Item~\textup{(V25.6)} is
\Cref{thm:V25-two-stage-radius,cor:V25-cofinal}, and item~\textup{(V25.7)} is
\Cref{prop:V25-exact-control}.
\end{proof}

\begin{assessment}[Version V25 verdict]
\label{ass:V25-verdict}
The current result is
\begin{equation}
 \boxed{
 \begin{gathered}
 \passstatus\\[-0.15em]
 \text{primitive-complete nested ancestry, record-router Pythagoras,}\\
 \text{deterministic positive square acquisition, outward transport, exact replay}
 \\[0.45em]
 \obsstatus\\[-0.15em]
 \text{a line--crossing short priced before the complete primitive bank;}\\
 \text{a record-router term charged as a new source;}\\
 \text{or a fitted arm instrument demanded although the joint writers already occur}
 \\[0.45em]
 \stopstatus\\[-0.15em]
 \text{first selected RPESM numerical three-arm Gram, line sewing,}\\
 \text{supported floors, held-out positive square, and adjacent-cutoff values.}
 \end{gathered}}
 \label{eq:V25-verdict}
\end{equation}
The V24 arm instrument remains a valid optional positive implementation, but
its occurrence is no longer treated as the first missing physical row.  The
first unresolved task is the actual matrix evaluation of an already populated
common-history writer packet.
\end{assessment}

\begin{acquisition}[Version V26: evaluate one primitive-complete RPESM three-arm card]
\label{acq:V25-next}
Preserve Versions~V1--V25.  Choose exactly one V22 amplitude branch, one
protected retained-terminal RPESM history, one physical boundary record, and
one adjacent cutoff.  Freeze the complete saturated primitive synthesis
\(C_a\), phase-quenched crossing arm \(A_a\), line-weighted crossing arm
\(B_a\), line sewing, reflection and Berezin signs, coefficient supports, and
physical gain calibration before reading any value.

Populate in every occurring record cell
\begin{equation}
 \boxed{
 \left(
 G_{C,a},G_{A,a},G_{B,a},
 Q_{CA,a}^{(+)},Q_{CA,a}^{(i)},
 Q_{CB,a}^{(+)},Q_{CB,a}^{(i)},
 Q_{AB,a}^{(+)},Q_{AB,a}^{(i)},
 Q_{*,a}^{\rm dir}
 \right),}
 \label{eq:V25-next-card}
\end{equation}
with outward radii and canonical support floors.  Reconstruct the complete
six-block packet, test the held-out three-arm square, short the complete
primitive bank, and only then evaluate the cell line--crossing shorts and the
record-router term.  Append only the polar range of the positive final cell
short.  A primitive follower, crossing follower, or readable router is not a
second occurrence.

Repeat the complete positive packet at one adjacent cutoff and transport the
primitive support, residual crossing support, line sewing, partition amplitude,
and final short together.  A separate stochastic arm device is required only
if that device, rather than the numerical Gram, is itself the declared target.
Until this numerical card lands, do not add another phase, crossing, exterior,
localizer, Duhamel, or continuum compiler.
\end{acquisition}

\section*{V25 exact replay and append-only record}
\addcontentsline{toc}{section}{V25 exact replay and append-only record}
The accompanying verifier checks the complete three-arm Gram, the nested
primitive and crossing shorts, the strict overpricing by the premature V24
short, the record-router decomposition, the nine positive square
reconstruction, the held-out three-arm writer, and the exact source-mass
allocation.  It is a theorem replay and not a numerical reading of an active
RPESM history.

The complete Version~V24 source preceding its sole final active
\verb|\end{document}| is preserved byte for byte.  Its preterminal SHA--256
digest is recorded in the validation manifest.  A later continuation must
remove only the final active document-closing command, append new material
immediately before it, restore one terminal command, and increment the filename
to end in \texttt{\_V26.tex}.

% =============================================================================
% END APPENDED VERSION V25 MATERIAL
% =============================================================================


% =============================================================================
% BEGIN APPENDED VERSION V26 MATERIAL
% =============================================================================

\clearpage
\clearpage
\hypersetup{
  pdftitle={Renewal Einstein--Standard--Model Physical Source Metric, Duhamel Ancestry and Quantum Continuum Programme V26},
  pdfsubject={Selection-free boundary-record disintegration, primitive-complete refinement telescope, and one-anchor cofinal landing}
}
\pdfinfo{
 /Title (Renewal Einstein--Standard--Model Physical Source Metric, Duhamel Ancestry and Quantum Continuum Programme V26)
 /Subject (Selection-free boundary-record disintegration, primitive-complete refinement telescope, and one-anchor cofinal landing)
}

\section{Version V26: selection-free boundary-record disintegration, a primitive-complete refinement telescope, and one-anchor cofinal landing}
\label{sec:V26-selection-free}

\subsection{Why the singleton-card mandate must be weakened}
\label{subsec:V26-context}

Version~V25 correctly moved the line-weighted crossing behind the complete
primitive bank.  Its proposed next acquisition nevertheless asked for one
selected retained-terminal history.  That is stronger than the physical target
unless a singleton history selector belongs to the readable boundary record.
The retained-terminal construction supplies literal provenance and one common
history for every mixed writer, but provenance and readable conditioning are
not the same operation.  The source price must be computed at the finest record
which the physical positive-time algebra actually carries.

The correct target is therefore an unnormalized positive Gram measure on the
physical boundary record.  Singleton history values may be used only when the
corresponding point selectors occur.  This change also removes an unnecessary
rare-cell instability: unnormalized cell blocks transport linearly, whereas
normalizing each cell before transport introduces inverse cell probabilities
which are not part of the source geometry.

The new order is
\begin{equation}
 \boxed{
 \begin{gathered}
 \text{one common retained-terminal history law and physical boundary record}
 \\[-0.1em]\Downarrow\\[-0.1em]
 \text{unnormalized primitive-complete Gram measure on that record}
 \\[-0.1em]\Downarrow\\[-0.1em]
 \text{conditional Schur short at the readable resolution}
 \\[-0.1em]\Downarrow\\[-0.1em]
 \text{record-refinement router telescope and final within-record novelty}
 \\[-0.1em]\Downarrow\\[-0.1em]
 \text{one finite anchor plus the inherited summable path debit.}
 \end{gathered}}
 \label{eq:V26-order}
\end{equation}
No new phase, crossing, exterior, localizer, or Duhamel source is introduced.

% -----------------------------------------------------------------------------
\subsection{The primitive-complete Gram measure}
\label{subsec:V26-Gram-measure}
% -----------------------------------------------------------------------------

Let $(\Omega,\mathscr F,\mu)$ be one finite positive physical history law and
let $\mathscr R\subseteq\mathscr F$ be the sigma-algebra generated by the
physical reflection-fixed boundary record.  The coefficient carriers below are
fixed finite-dimensional Hilbert spaces.  At each history let
\begin{equation}
 \mathbb G(\omega)
 =\begin{pmatrix}
    H(\omega)&W(\omega)\\
    W(\omega)^*&N(\omega)
  \end{pmatrix}\succeq0,
 \label{eq:V26-history-Gram}
\end{equation}
where
\begin{equation}
 H=\begin{pmatrix}S&U\\U^*&M\end{pmatrix},
 \qquad
 W=\binom{V}{\Gamma}.
 \label{eq:V26-V25-blocking}
\end{equation}
Thus $H$ is the joint prior-primitive-plus-crossing block of Version~V25,
$W$ is its mixed block with the line-weighted crossing arm, and $N$ is the
line-weighted diagonal block.  Assume the matrix entries are integrable.

For $E\in\mathscr R$, define the unnormalized operator-valued measure
\begin{equation}
 \boxed{
 \mathbf G(E)
 :=\int_E\mathbb G(\omega)\,\mu(\dd\omega)
 =\begin{pmatrix}H_E&W_E\\W_E^*&N_E\end{pmatrix}\succeq0.}
 \label{eq:V26-Gram-measure}
\end{equation}
The decoder and final cell short are
\begin{equation}
 \boxed{
 L_E:=H_E^\dagger W_E,
 \qquad
 D_E:=N_E-W_E^*H_E^\dagger W_E.}
 \label{eq:V26-cell-decoder-short}
\end{equation}
When $\mu(E)=0$, all four blocks and both displayed quantities are defined to
be zero.

\begin{lemma}[Automatic support incidence for every record cell]
\label{lem:V26-cell-support}
For every $E\in\mathscr R$,
\begin{equation}
 \boxed{
 (I-H_EH_E^\dagger)W_E=0,
 \qquad
 H_EL_E=W_E,
 \qquad
 D_E\succeq0.}
 \label{eq:V26-cell-support}
\end{equation}
Moreover,
\begin{equation}
 \boxed{
 W_E^*H_E^\dagger W_E=L_E^*H_EL_E.}
 \label{eq:V26-decoder-energy}
\end{equation}
\end{lemma}

\begin{proof}
The integral in \eqref{eq:V26-Gram-measure} is positive.  If $x\in\Ker H_E$,
positivity of the quadratic polynomial obtained from $(tx,y)$ forces
$x^*W_Ey=0$ for every $y$.  Hence $\Ran W_E\subseteq\Ran H_E$, proving the
support incidence and $H_EL_E=W_E$.  The supported Moore--Penrose Schur
criterion gives $D_E\succeq0$.  Finally,
$H_E^\dagger H_EH_E^\dagger=H_E^\dagger$, which gives
\eqref{eq:V26-decoder-energy}.
\end{proof}

\begin{remark}[Unnormalized cells are the native transport objects]
\label{rem:V26-unnormalized}
If $p_E=\mu(E)>0$ and
$\widehat H_E=p_E^{-1}H_E$, with analogous notation for $W,N$, then
\begin{equation}
 L_E=\widehat H_E^\dagger\widehat W_E,
 \qquad
 D_E=p_E\left(
  \widehat N_E-\widehat W_E^*\widehat H_E^\dagger\widehat W_E
 \right).
 \label{eq:V26-conditional-scaling}
\end{equation}
The decoder is unchanged by cell normalization, while the short carries the
physical cell mass.  Transporting $(H_E,W_E,N_E)$ therefore avoids every
spurious factor $p_E^{-1}$.
\end{remark}

% -----------------------------------------------------------------------------
\subsection{Exact refinement Pythagoras}
\label{subsec:V26-refinement}
% -----------------------------------------------------------------------------

\begin{theorem}[Primitive-complete record-refinement Pythagoras]
\label{thm:V26-refinement-Pythagoras}
Let $\mathcal Q$ be a finite measurable partition refining a finite measurable
partition $\mathcal P$.  For $E\in\mathcal P$, let
$\mathcal Q(E)=\{F\in\mathcal Q:F\subseteq E\}$.  Then
\begin{equation}
 \boxed{
 D_E
 =\sum_{F\in\mathcal Q(E)}D_F
  +\sum_{F\in\mathcal Q(E)}
   (L_F-L_E)^*H_F(L_F-L_E).}
 \label{eq:V26-refinement-Pythagoras}
\end{equation}
Consequently, with
\begin{equation}
 \mathfrak D(\mathcal P):=\sum_{E\in\mathcal P}D_E,
 \label{eq:V26-partition-short}
\end{equation}
one has
\begin{equation}
 \boxed{
 \mathfrak D(\mathcal P)
 =\mathfrak D(\mathcal Q)
  +\mathcal R(\mathcal Q/\mathcal P),}
 \label{eq:V26-partition-telescope-step}
\end{equation}
where
\begin{equation}
 \boxed{
 \mathcal R(\mathcal Q/\mathcal P)
 :=\sum_{E\in\mathcal P}
   \sum_{F\in\mathcal Q(E)}
   (L_F-L_E)^*H_F(L_F-L_E)\succeq0.}
 \label{eq:V26-router-increment}
\end{equation}
Thus readable refinement can only decrease the source novelty.  The decrement
is exactly the primitive-complete router dispersion revealed by the finer
record.
\end{theorem}

\begin{proof}
For every $F$, Lemma~\ref{lem:V26-cell-support} gives
$W_F=H_FL_F$ and
$D_F=N_F-L_F^*H_FL_F$.  Additivity of the operator-valued measure gives
\begin{equation}
 H_E=\sum_FH_F,
 \qquad W_E=\sum_FW_F=\sum_FH_FL_F,
 \qquad N_E=\sum_FN_F,
 \label{eq:V26-additive-blocks}
\end{equation}
where the sums range over $\mathcal Q(E)$.  Since $H_EL_E=W_E$, expansion of
the router term yields
\begin{align*}
 &\sum_F(L_F-L_E)^*H_F(L_F-L_E)\\
 &\quad=\sum_FL_F^*H_FL_F-L_E^*H_EL_E.
\end{align*}
Adding $\sum_FD_F$ cancels the first sum and leaves
$N_E-L_E^*H_EL_E=D_E$.  Summation over $E$ proves the remaining statements.
\end{proof}

\begin{corollary}[A zero coarse short is refinement robust]
\label{cor:V26-zero-coarse}
If $D_E=0$ for one coarse cell $E$, then every refinement in
\eqref{eq:V26-refinement-Pythagoras} satisfies
\begin{equation}
 \boxed{
 D_F=0,
 \qquad
 H_F^{1/2}(L_F-L_E)=0
 \quad(F\subseteq E).}
 \label{eq:V26-zero-coarse-conclusion}
\end{equation}
In particular, a zero mixed primitive-complete short cannot hide either a fine
within-record novelty or a decoder-routing residual.
\end{corollary}

\begin{proof}
Every term on the right of
\eqref{eq:V26-refinement-Pythagoras} is positive semidefinite.  A finite sum of
positive matrices is zero only when every summand is zero.
\end{proof}

\begin{countertheorem}[A positive coarse short has no unique ancestry]
\label{cth:V26-positive-coarse-ambiguous}
A positive coarse short may consist entirely of decoder routing, entirely of
within-record novelty, or any positive combination of the two.  Therefore it
is only an upper ancestry card until the physical record resolution or a
router bound has been supplied.
\end{countertheorem}

\begin{proof}
The two scalar cards of Version~V24 already give the first two alternatives
with identical coarse diagonal and mixed blocks.  Direct sums and positive
rescalings produce every convex combination.  Equation
\eqref{eq:V26-refinement-Pythagoras} identifies the two contributions exactly.
\end{proof}

% -----------------------------------------------------------------------------
\subsection{Conditional sigma-algebras and the martingale telescope}
\label{subsec:V26-conditional}
% -----------------------------------------------------------------------------

For a sub-sigma-algebra $\mathscr A\subseteq\mathscr F$, set
\begin{equation}
 H_{\mathscr A}=\mathbb E[H\mid\mathscr A],
 \quad
 W_{\mathscr A}=\mathbb E[W\mid\mathscr A],
 \quad
 N_{\mathscr A}=\mathbb E[N\mid\mathscr A],
 \label{eq:V26-conditional-blocks}
\end{equation}
and define
\begin{equation}
 L_{\mathscr A}=H_{\mathscr A}^\dagger W_{\mathscr A},
 \qquad
 D_{\mathscr A}
 =N_{\mathscr A}
  -W_{\mathscr A}^*H_{\mathscr A}^\dagger W_{\mathscr A}.
 \label{eq:V26-conditional-short}
\end{equation}
Conditional expectation preserves positivity of the block Gram, so these
quantities satisfy the support identities of
Lemma~\ref{lem:V26-cell-support} almost surely.

\begin{theorem}[Conditional primitive-complete Pythagoras]
\label{thm:V26-conditional-Pythagoras}
If $\mathscr A\subseteq\mathscr B\subseteq\mathscr F$, then almost surely
\begin{equation}
 \boxed{
 \begin{aligned}
 D_{\mathscr A}
 ={}&\mathbb E[D_{\mathscr B}\mid\mathscr A]\\
 &+\mathbb E\!\left[
 (L_{\mathscr B}-L_{\mathscr A})^*
 H_{\mathscr B}
 (L_{\mathscr B}-L_{\mathscr A})
 \middle|\mathscr A\right].
 \end{aligned}}
 \label{eq:V26-conditional-Pythagoras}
\end{equation}
Both terms are positive semidefinite.  Hence
\begin{equation}
 \boxed{
 \mathbb E D_{\mathscr A}\succeq
 \mathbb E D_{\mathscr B}.}
 \label{eq:V26-conditional-monotonicity}
\end{equation}
\end{theorem}

\begin{proof}
By the tower property,
$H_{\mathscr A}=\mathbb E[H_{\mathscr B}\mid\mathscr A]$, and similarly for
$W,N$.  Moreover
$H_{\mathscr B}L_{\mathscr B}=W_{\mathscr B}$ and
$H_{\mathscr A}L_{\mathscr A}=W_{\mathscr A}$.  Expand the second term in
\eqref{eq:V26-conditional-Pythagoras}.  The two conditional cross terms are
both $L_{\mathscr A}^*H_{\mathscr A}L_{\mathscr A}$, while the final quadratic
term restores one copy.  Adding
\begin{equation*}
 \mathbb E[D_{\mathscr B}\mid\mathscr A]
 =N_{\mathscr A}
  -\mathbb E[L_{\mathscr B}^*H_{\mathscr B}L_{\mathscr B}
             \mid\mathscr A]
\end{equation*}
leaves precisely $D_{\mathscr A}$.  Taking expectations proves
\eqref{eq:V26-conditional-monotonicity}.
\end{proof}

\begin{corollary}[Finite filtration telescope]
\label{cor:V26-filtration-telescope}
For
$\mathscr A_0\subseteq\cdots\subseteq\mathscr A_m$, put
\begin{equation}
 \mathcal R_k
 :=\mathbb E\left[
 (L_{\mathscr A_{k+1}}-L_{\mathscr A_k})^*
 H_{\mathscr A_{k+1}}
 (L_{\mathscr A_{k+1}}-L_{\mathscr A_k})
 \right].
 \label{eq:V26-filtration-router}
\end{equation}
Then
\begin{equation}
 \boxed{
 \mathbb E D_{\mathscr A_0}
 =\mathbb E D_{\mathscr A_m}
  +\sum_{k=0}^{m-1}\mathcal R_k.}
 \label{eq:V26-filtration-telescope}
\end{equation}
No source mass is lost: every decrement is recorded as a positive decoder
innovation.
\end{corollary}

\begin{proof}
Take expectations in
\eqref{eq:V26-conditional-Pythagoras} for every adjacent pair and telescope.
\end{proof}

\begin{theorem}[Infinite readable-record landing]
\label{thm:V26-infinite-record-landing}
Let $\mathscr A_n\uparrow\mathscr A_\infty$.  Assume there is one fixed
projection $P$ and $h_*>0$ such that
\begin{equation}
 H_{\mathscr A_n}\succeq h_*P,
 \qquad
 H_{\mathscr A_n}=PH_{\mathscr A_n}P,
 \qquad
 W_{\mathscr A_n}=PW_{\mathscr A_n}
 \label{eq:V26-record-support-floor}
\end{equation}
almost surely for every $n$, and assume $H,W,N$ are square integrable.  Then
\begin{equation}
 \boxed{
 \mathbb E D_{\mathscr A_0}
 =\mathbb E D_{\mathscr A_\infty}
  +\sum_{n=0}^{\infty}\mathcal R_n,}
 \label{eq:V26-infinite-telescope}
\end{equation}
with convergence in operator norm of the finite-dimensional matrices.  The
series of positive router increments converges, and
$D_{\mathscr A_n}\to D_{\mathscr A_\infty}$ in $L^1$.
\end{theorem}

\begin{proof}
Finite-dimensional martingale convergence gives $L^2$ convergence of the
three conditional block fields.  On the supported cone
$H\succeq h_*P$, inversion is Lipschitz:
\begin{equation}
 \norm{H_2^\dagger-H_1^\dagger}
 \le h_*^{-2}\norm{H_2-H_1}.
 \label{eq:V26-inverse-Lipschitz}
\end{equation}
Hence $L_{\mathscr A_n}$ and $D_{\mathscr A_n}$ converge in $L^1$ to the
corresponding $\mathscr A_\infty$ quantities.  Apply
Corollary~\ref{cor:V26-filtration-telescope} and pass to the limit.  The partial
sums of the positive router series are increasing and bounded above by
$\mathbb E D_{\mathscr A_0}$, so they converge in finite dimension.
\end{proof}

\begin{firewall}[A provenance history is not automatically a conditioning port]
\label{fire:V26-provenance-not-selector}
The full retained-terminal provenance may be used as
$\mathscr A_\infty$ in \eqref{eq:V26-infinite-telescope} only when the target
admits its point selectors or an equivalent physically occurring conditional
Read.  Otherwise the source price stops at the actual boundary sigma-algebra.
Mathematical disintegration on a larger latent record remains a diagnostic
allocation, not an executable record-adapted source policy.
\end{firewall}

% -----------------------------------------------------------------------------
\subsection{Router-debited birth certificates}
\label{subsec:V26-router-debit}
% -----------------------------------------------------------------------------

\begin{corollary}[Strict birth after a validated router debit]
\label{cor:V26-router-debited-birth}
In the setting of Theorem~\ref{thm:V26-conditional-Pythagoras}, let $P_Q$ be a
protected target projection.  Suppose on one coarse cell
\begin{equation}
 P_QD_{\mathscr A}P_Q\succeq\delta P_Q
 \label{eq:V26-coarse-margin}
\end{equation}
and the conditional router term has the validated upper bound
\begin{equation}
 P_Q\mathbb E\!\left[
 (L_{\mathscr B}-L_{\mathscr A})^*H_{\mathscr B}
 (L_{\mathscr B}-L_{\mathscr A})
 \middle|\mathscr A\right]P_Q
 \preceq\eta P_Q.
 \label{eq:V26-router-upper}
\end{equation}
If $\delta>\eta$, then
\begin{equation}
 \boxed{
 P_Q\mathbb E[D_{\mathscr B}\mid\mathscr A]P_Q
 \succeq(\delta-\eta)P_Q.}
 \label{eq:V26-fine-birth-lower}
\end{equation}
For every nonzero $v\in\Ran P_Q$, the set of finer records on which
$v^*D_{\mathscr B}v>0$ has positive conditional measure.
\end{corollary}

\begin{proof}
Subtract \eqref{eq:V26-router-upper} from
\eqref{eq:V26-coarse-margin} in
\eqref{eq:V26-conditional-Pythagoras}.  If the last assertion failed, the
conditional expectation of the nonnegative scalar
$v^*D_{\mathscr B}v$ would vanish.
\end{proof}

\begin{corollary}[Exact pure-router and mixed branches]
\label{cor:V26-pure-router}
On a coarse cell, the following branches are exact:
\begin{enumerate}[label=\textup{(R\arabic*)},leftmargin=3.5em]
\item if the router term equals $D_{\mathscr A}$, then
      $D_{\mathscr B}=0$ almost surely and the coarse residual is pure routing;
\item if the router term vanishes, then
      $D_{\mathscr A}=\mathbb E[D_{\mathscr B}\mid\mathscr A]$ and the coarse
      residual contains no decoder dispersion;
\item otherwise the two positive terms must be retained separately.
\end{enumerate}
\end{corollary}

\begin{proof}
This is the positive decomposition
\eqref{eq:V26-conditional-Pythagoras}.  Equality of one positive summand with
the total forces the other to vanish.
\end{proof}

% -----------------------------------------------------------------------------
\subsection{The retained-terminal debit transports the unnormalized square packet}
\label{subsec:V26-path-transport}
% -----------------------------------------------------------------------------

The retained-terminal theorem may be rerun after adjoining the finitely many
normalized scalar coordinates needed to reconstruct the Version~V25 positive
square packet in every cell of one finite boundary partition.  Each cell
indicator is multiplied into the writer before the expectation is taken.  The
objects transported are therefore the unnormalized cell blocks of
\eqref{eq:V26-Gram-measure}, not conditional matrices divided by a rare cell
mass.

Fix a Hermitian coordinate frame for all nine positive square matrices in one
record packet.  Let $\mathbf q_n$ denote the vector of its normalized scalar
coordinates after transport to a common old card, and let
\begin{equation}
 \eta_n^{\rm sq}:=\norm{\mathbf q_{n+1\downarrow n}-\mathbf q_n}_{\infty}.
 \label{eq:V26-square-coordinate-defect}
\end{equation}
Let $\mathcal T$ be the fixed linear reconstruction from these coordinates to
the complete record-direct-sum blocks $(H_n,W_n,N_n)$, and define
\begin{equation}
 \kappa_{\mathcal T}
 :=\norm{\mathcal T}_{\ell^\infty\to\max(\mathrm{op},\mathrm{op},\mathrm{op})}.
 \label{eq:V26-reconstruction-constant}
\end{equation}
This is a finite, explicitly computable constant fixed by the coefficient
basis and the nine-square polarization formulas.

\begin{theorem}[Summable positive-square transport from the populated path]
\label{thm:V26-square-transport}
Augment the bounded retained-terminal grammar by the normalized scalar
coordinates defining $\mathbf q_n$.  The cofinal selection can be chosen so
that
\begin{equation}
 \boxed{
 \sum_{n\ge0}\eta_n^{\rm sq}
 \le D_{\rm cond}\le\frac7{512}.}
 \label{eq:V26-square-debit}
\end{equation}
Consequently
\begin{equation}
 \boxed{
 \max\{\norm{\Delta H_n},\norm{\Delta W_n},\norm{\Delta N_n}\}
 \le\kappa_{\mathcal T}\eta_n^{\rm sq}.}
 \label{eq:V26-block-transport}
\end{equation}
The same statement holds cellwise after descendant cells are added before
comparison.
\end{theorem}

\begin{proof}
There are only finitely many coefficient coordinates in the chosen record
packet.  Scale each corresponding scalar square writer to the interval
$[0,1]$ and include it, together with its cell indicator, in the finite bounded
path grammar before applying the retained-terminal diagonal selection.  The
selected endpoint and compact-layer omissions therefore bound every coordinate
difference by the same path-plus-layer debit.  Their sums are
$3/256+1/512=7/512$.  Applying the fixed linear reconstruction map gives
\eqref{eq:V26-block-transport}.  Descendant aggregation is linear and is
performed before the comparison, so it introduces no inverse cell mass.
\end{proof}

\begin{remark}[What has and has not been populated]
\label{rem:V26-population-status}
The theorem populates a summable transport budget for the bounded positive
square packet on a newly selected cofinal retained-terminal sequence.  It does
not print the base numerical matrices, their supported floors, or the physical
line-sewing sign.  Those remain one finite-anchor acquisition.
\end{remark}

% -----------------------------------------------------------------------------
\subsection{One finite anchor controls the cofinal final short}
\label{subsec:V26-one-anchor}
% -----------------------------------------------------------------------------

Let
\begin{equation}
 \mathscr S(H,W,N):=N-W^*H^\dagger W.
 \label{eq:V26-short-map}
\end{equation}

\begin{lemma}[Supported Lipschitz bound for the assembled short]
\label{lem:V26-short-Lipschitz}
Let two transported packets have one common support $P$ and satisfy
\begin{equation}
 H_i\succeq h_*P,
 \qquad
 H_i=PH_iP,
 \qquad
 W_i=PW_i,
 \qquad
 \norm{W_i}\le w_*.
 \label{eq:V26-short-domain}
\end{equation}
Then
\begin{equation}
 \boxed{
 \begin{aligned}
 \norm{\mathscr S(H_2,W_2,N_2)-\mathscr S(H_1,W_1,N_1)}
 \le{}&\norm{\Delta N}\\
 &+\frac{2w_*}{h_*}\norm{\Delta W}
 +\frac{w_*^2}{h_*^2}\norm{\Delta H}.
 \end{aligned}}
 \label{eq:V26-short-Lipschitz}
\end{equation}
\end{lemma}

\begin{proof}
On the fixed support,
\begin{equation}
 H_2^\dagger-H_1^\dagger
 =-H_2^\dagger(H_2-H_1)H_1^\dagger.
 \label{eq:V26-pseudoinverse-difference}
\end{equation}
Add and subtract
$W_2^*H_2^\dagger W_1$ and
$W_2^*H_1^\dagger W_1$, then use
$\norm{H_i^\dagger}\le h_*^{-1}$ and
$\norm{W_i}\le w_*$.
\end{proof}

Define
\begin{equation}
 \Lambda_*(h_*,w_*)
 :=\kappa_{\mathcal T}
 \left(1+\frac{2w_*}{h_*}+\frac{w_*^2}{h_*^2}\right)
 \label{eq:V26-short-transport-constant}
\end{equation}
and the remaining positive-square tail
\begin{equation}
 \mathfrak e_{n_0}^{\rm sq}
 :=\sum_{n\ge n_0}\eta_n^{\rm sq}.
 \label{eq:V26-square-tail}
\end{equation}

\begin{theorem}[One-anchor cofinal source certificate]
\label{thm:V26-one-anchor}
Assume the record map, coefficient identification, line sewing, and the support
conditions \eqref{eq:V26-short-domain} hold from $n_0$ onward.  Let $D_n$ be
the direct sum of the final primitive-complete cell shorts at cutoff $n$,
transported to the $n_0$ screen.  Then $D_n$ converges in operator norm to a
unique $D_\infty$ and
\begin{equation}
 \boxed{
 \norm{D_\infty-D_{n_0}}
 \le
 \Lambda_*(h_*,w_*)\mathfrak e_{n_0}^{\rm sq}.}
 \label{eq:V26-one-anchor-radius}
\end{equation}
In particular:
\begin{enumerate}[label=\textup{(A\arabic*)},leftmargin=3.5em]
\item if, on a protected coefficient screen $P_Q$,
      \begin{equation}
       P_QD_{n_0}P_Q
       \succeq
       \left(\delta+
       \Lambda_*\mathfrak e_{n_0}^{\rm sq}\right)P_Q
       \label{eq:V26-anchor-pass}
      \end{equation}
      for some $\delta>0$, then
      $P_QD_\infty P_Q\succeq\delta P_Q$;
\item if
      \begin{equation}
       \norm{D_{n_0}}+
       \Lambda_*\mathfrak e_{n_0}^{\rm sq}
       \le\varepsilon,
       \label{eq:V26-anchor-follower}
      \end{equation}
      then $\norm{D_\infty}\le\varepsilon$;
\item a change of record incidence, line sewing, or canonical support is a
      typed transport event and is not absorbed into
      \eqref{eq:V26-one-anchor-radius}.
\end{enumerate}
The cell novelty and router contributions satisfy the same conclusion
separately when their exact record-refinement packet is transported.
\end{theorem}

\begin{proof}
Combine Theorem~\ref{thm:V26-square-transport} with
Lemma~\ref{lem:V26-short-Lipschitz} and telescope.  The summable bound makes
$(D_n)$ Cauchy.  The two margin statements are Weyl's inequality and the
triangle inequality.  The final sentence follows by applying the same
argument to the exact positive decomposition
\eqref{eq:V26-refinement-Pythagoras}; no cancellation between novelty and
router transport is used.
\end{proof}

\begin{corollary}[The adjacent numerical copy is no longer the first missing row]
\label{cor:V26-adjacent-discharged}
For the bounded positive square packet, the retained-terminal construction and
Theorem~\ref{thm:V26-square-transport} already provide a summable adjacent
transport route after the packet is inserted into the selected grammar.  The
first remaining numerical acquisition is one finite anchor containing:
\begin{equation}
 \boxed{
 \begin{gathered}
 \text{the physical boundary record and unnormalized square values,}\\
 h_*,\ w_*,\ \text{line sewing, and one held-out square}.
 \end{gathered}}
 \label{eq:V26-reduced-anchor}
\end{equation}
An additional adjacent numerical card is required only for a different route,
an unbounded writer, or an independent validation of the selected cofinal
presentation.
\end{corollary}

\begin{proof}
Insert the finite normalized square-coordinate bank and its cell indicators
into the bounded retained-terminal grammar before the diagonal cofinal
selection.  Theorem~\ref{thm:V26-square-transport} then supplies the complete
summable adjacent debit for those coordinates, and
Theorem~\ref{thm:V26-one-anchor} converts one supported anchor into the stated
cofinal interval.  The listed exceptions are precisely the hypotheses not
covered by that bounded common-route argument.
\end{proof}

% -----------------------------------------------------------------------------
\subsection{Exact matrix replay}
\label{subsec:V26-replay}
% -----------------------------------------------------------------------------

Let $\Omega=\{1,2,3,4\}$ with equal weights.  On a two-dimensional coefficient
space take $H_i=I_2$ and history decoders
\begin{equation}
 \begin{aligned}
 L_1&=\diag(0,0),&L_2&=\diag(2,0),\\
 L_3&=\diag(2,2),&L_4&=\diag(4,2),
 \end{aligned}
 \label{eq:V26-control-decoders}
\end{equation}
with historywise novelty
\begin{equation}
 D_1=D_2=\diag(0,1),
 \qquad
 D_3=D_4=\diag(1,0).
 \label{eq:V26-control-novelty}
\end{equation}
Set
\begin{equation}
 W_i=H_iL_i,
 \qquad
 N_i=L_i^*H_iL_i+D_i.
 \label{eq:V26-control-blocks}
\end{equation}
Every history block is positive and has short $D_i$.

Let $\mathscr A_0$ be the trivial record,
$\mathscr A_1$ the two-cell record
$\{\{1,2\},\{3,4\}\}$, and
$\mathscr A_2$ the full history record.  With the physical weights included,
V26 obtains
\begin{equation}
 \boxed{
 \mathbb E D_{\mathscr A_0}
 =\diag\left(\frac52,\frac32\right),}
 \label{eq:V26-control-level0}
\end{equation}
\begin{equation}
 \boxed{
 \mathbb E D_{\mathscr A_1}
 =\diag\left(\frac32,\frac12\right),}
 \label{eq:V26-control-level1}
\end{equation}
and
\begin{equation}
 \boxed{
 \mathbb E D_{\mathscr A_2}
 =\diag\left(\frac12,\frac12\right).}
 \label{eq:V26-control-level2}
\end{equation}
The two router increments are
\begin{equation}
 \boxed{
 \mathcal R_0=I_2,
 \qquad
 \mathcal R_1=\diag(1,0),}
 \label{eq:V26-control-router-increments}
\end{equation}
so the exact telescope is
\begin{equation}
 \boxed{
 \diag\left(\frac52,\frac32\right)
 =\diag\left(\frac12,\frac12\right)
  +I_2+\diag(1,0).}
 \label{eq:V26-control-telescope}
\end{equation}

The first coordinate therefore contains one half unit of final historywise
novelty and two full units of router dispersion.  The second contains one half
unit of novelty and one unit of router dispersion.  Pricing the trivial-record
short as an atomic source would overcount both coordinates.

If all $D_i$ are changed to zero while the decoders are retained, the coarse
short becomes
\begin{equation}
 \boxed{\diag(2,1)}
 \label{eq:V26-pure-router-control}
\end{equation}
and is entirely routing.  If all decoders are made equal and all $D_i$ are
zero, every coarse and fine short vanishes, illustrating
Corollary~\ref{cor:V26-zero-coarse}.

\begin{theorem}[Version V26 selection-free primitive-complete theorem]
\label{thm:V26-terminal}
Versions~V1--V25 are retained.  Version~V26 additionally proves:
\begin{enumerate}[label=\textup{(V26.\arabic*)},leftmargin=5.5em]
\item the complete primitive-plus-crossing versus line-weighted Gram defines an
      unnormalized positive operator-valued measure on the actual boundary
      record;
\item every physical record refinement has the exact primitive-complete
      novelty-plus-router Pythagoras
      \eqref{eq:V26-refinement-Pythagoras};
\item a zero coarse short is robust under every finer readable conditioning,
      whereas a positive coarse short is only an upper ancestry card;
\item conditional sigma-algebras give the exact martingale telescope
      \eqref{eq:V26-filtration-telescope}, with an infinite landing under one
      supported floor;
\item a router upper bound converts a coarse margin into a strict fine novelty
      certificate;
\item the populated retained-terminal path transports the finite normalized
      positive-square coordinate packet with total debit at most $7/512$ after
      those writers are inserted into the selected grammar;
\item one finite anchor and the explicit supported Lipschitz constant
      \eqref{eq:V26-short-transport-constant} decide a cofinal positive floor or
      an outward follower interval;
\item the exact matrix replay separates two successive router increments from
      the final historywise novelty.
\end{enumerate}
No item supplies the base numerical RPESM square values or licenses
conditioning on a singleton provenance history which is not physically
readable.
\end{theorem}

\begin{proof}
Items~\textup{(V26.1)--(V26.3)} are
Lemma~\ref{lem:V26-cell-support},
Theorem~\ref{thm:V26-refinement-Pythagoras}, and
Corollary~\ref{cor:V26-zero-coarse}.  Items~\textup{(V26.4)--(V26.5)} are
Theorems~\ref{thm:V26-conditional-Pythagoras}--
\ref{thm:V26-infinite-record-landing} and
Corollary~\ref{cor:V26-router-debited-birth}.  Items~\textup{(V26.6)--(V26.7)}
are Theorems~\ref{thm:V26-square-transport} and
\ref{thm:V26-one-anchor}.  Item~\textup{(V26.8)} is the direct calculation
\eqref{eq:V26-control-level0}--\eqref{eq:V26-control-telescope}.
\end{proof}

\begin{assessment}[Version V26 verdict]
\label{ass:V26-verdict}
The current result is
\begin{equation}
 \boxed{
 \begin{gathered}
 \passstatus\\[-0.15em]
 \text{selection-free primitive-complete Gram measure, exact record telescope,}\\
 \text{router-debited birth, inherited summable transport, and one-anchor landing}
 \\[0.45em]
 \obsstatus\\[-0.15em]
 \text{a singleton provenance card used without a physical selector;}\\
 \text{a normalized rare-cell packet transported instead of its unnormalized measure;}\\
 \text{or a positive coarse short priced before decoder dispersion is removed}
 \\[0.45em]
 \stopstatus\\[-0.15em]
 \text{one finite boundary-record RPESM anchor: unnormalized square values,}\\
 \text{supported floors, physical line sewing, and one held-out square.}
 \end{gathered}}
 \label{eq:V26-verdict}
\end{equation}
The adjacent bounded-panel route is no longer a separate first acquisition: it
is supplied, with summable debit, by the retained-terminal selection after the
panel is inserted into the bounded grammar.
\end{assessment}

\begin{acquisition}[Version V27: one boundary-record anchor, not one latent history]
\label{acq:V26-next}
Preserve Versions~V1--V26.  Choose one V22 amplitude semantics, one physically
readable finite reflection-fixed boundary record, and one retained-terminal
card sufficiently far along the cofinal selection that the remaining square
transport debit is explicit.

In every occurring record cell populate the unnormalized nine-square packet
of Version~V25 and one held-out three-arm square, together with the cell vacuum
amplitude, line sewing, canonical top-block support, a supported floor $h_*$,
and a bound $w_*$ for the mixed block.  Reconstruct the complete
primitive-plus-crossing block $H_E$, line mixed block $W_E$, and line diagonal
block $N_E$.  Run the exact cell short and the record-refinement telescope.
Append only the polar range of the final within-record short.  Retain every
positive decoder-router term as a follower-routing row, not as another source.

Compare the anchor margin with the explicit tail in
\eqref{eq:V26-one-anchor-radius}.  A strict lower interval gives a cofinal
source survivor.  A small upper interval gives an asymptotic follower.  A
support, sewing, or record transition is returned before the Moore--Penrose
short is continued.  A singleton history evaluation is admitted only when its
selector is an actual boundary Read.

Do not open the conditional exterior, locality, current/stress, charge,
Duhamel, or Einstein gates until this one anchor has been evaluated.
\end{acquisition}

\section*{V26 exact replay and append-only record}
\addcontentsline{toc}{section}{V26 exact replay and append-only record}
The accompanying exact verifier checks positivity of every history block, the
three record levels, both router increments, the complete matrix telescope, the
pure-router card, and the zero-coarse robustness control.  It also checks a
synthetic one-anchor radius inequality.  These are theorem replays and are not
numerical measurements of the selected RPESM regulator.

The complete Version~V25 source preceding its sole final active
\verb|\end{document}| is preserved byte for byte.  Its preterminal SHA--256
digest is recorded in the validation manifest.  A later continuation must
remove only the final active document-closing command, append new material
immediately before it, restore one terminal command, and increment the filename
to end in \texttt{\_V27.tex}.

% =============================================================================
% END APPENDED VERSION V26 MATERIAL
% =============================================================================


% =============================================================================
% BEGIN APPENDED VERSION V27 MATERIAL
% =============================================================================

\clearpage
\clearpage
\hypersetup{
  pdftitle={Renewal Einstein--Standard--Model Physical Source Metric, Duhamel Ancestry and Quantum Continuum Programme V27},
  pdfsubject={Source-total gain calibration, bounded shape cylinder, and response-tight cofinal transport}
}
\pdfinfo{
 /Title (Renewal Einstein--Standard--Model Physical Source Metric, Duhamel Ancestry and Quantum Continuum Programme V27)
 /Subject (Source-total gain calibration, bounded shape cylinder, and response-tight cofinal transport)
}

\section{Version V27: source-total gain calibration, a bounded shape cylinder, and response-tight cofinal transport}
\label{sec:V27-source-total-gain}

\subsection{Why the one-anchor statement still hid an absolute-scale premise}
\label{subsec:V27-context}

Version~V26 correctly replaced singleton-history conditioning by an
unnormalized Gram measure on the physical boundary record.  Its path argument,
however, inserted \emph{normalized} bounded square coordinates into the
retained-terminal grammar and then reconstructed \emph{unnormalized} Gram
blocks through one fixed linear map.  This is valid only when the inverse
normalization is itself fixed in physical units, or when its cutoff motion is
carried as an additional calibrated row.  Scaling each cutoff's writer to
$[0,1]$ by a different bound does not supply that row.

The missing quantity is not another phase or crossing matrix.  It is the
absolute source-total gain of the already assembled primitive--crossing--line
packet and the source-weighted mass of every omitted path layer.  A small
probability tail can carry order-one source action.  Consequently the
retained-terminal probability debit by itself does not transport the
unnormalized source Gram.

The corrected order is
\begin{equation}
 \boxed{
 \begin{gathered}
 \text{physical primitive--crossing--line square writers}
 \\[-0.1em]\Downarrow\\[-0.1em]
 \text{canonical source-total positive cylinder and bounded shape densities}
 \\[-0.1em]\Downarrow\\[-0.1em]
 \text{absolute gain and relative record-cell source masses}
 \\[-0.1em]\Downarrow\\[-0.1em]
 \text{primitive-complete conditional short and router telescope}
 \\[-0.1em]\Downarrow\\[-0.1em]
 \text{source-weighted tail transport and one-anchor cofinal interval.}
 \end{gathered}}
 \label{eq:V27-order}
\end{equation}
This section qualifies the transport clause of Version~V26.  It leaves every
finite record and Schur identity there unchanged.

% -----------------------------------------------------------------------------
\subsection{The canonical source-total cylinder}
\label{subsec:V27-source-total-cylinder}
% -----------------------------------------------------------------------------

Let $(\Omega,\mathscr F,\mu)$ be one positive physical history law.  After the
physical coefficient metrics have been frozen and whitened, let
\begin{equation}
 \mathcal X(\omega)=(C(\omega)\ \ A(\omega)\ \ B(\omega))
 \label{eq:V27-three-arm-synthesis}
\end{equation}
be the complete prior-primitive, phase-quenched crossing, and line-weighted
crossing synthesis of Version~V25.  Its pointwise positive Gram is
\begin{equation}
 \mathbb G(\omega)
 :=\mathcal X(\omega)^*\mathcal X(\omega)
 =\begin{pmatrix}
   S&U&V\\
   U^*&M&\Gamma\\
   V^*&\Gamma^*&N
  \end{pmatrix}\succeq0.
 \label{eq:V27-pointwise-three-arm-Gram}
\end{equation}
Assume its entries are integrable.  Define the source-total writer
\begin{equation}
 \boxed{
 t(\omega):=\Tr\mathbb G(\omega)
 =\norm{C(\omega)}_{\HS}^{2}
  +\norm{A(\omega)}_{\HS}^{2}
  +\norm{B(\omega)}_{\HS}^{2}.}
 \label{eq:V27-source-total-writer}
\end{equation}
The associated finite positive measure is
\begin{equation}
 \boxed{\nu(E):=\int_Et\,\dd\mu,\qquad E\in\mathscr F.}
 \label{eq:V27-source-total-measure}
\end{equation}
On $\{t>0\}$ put
\begin{equation}
 \widehat{\mathbb G}(\omega):=t(\omega)^{-1}\mathbb G(\omega),
 \label{eq:V27-pointwise-shape}
\end{equation}
and put $\widehat{\mathbb G}=0$ on $\{t=0\}$.

\begin{theorem}[Canonical source-total domination and bounded shape]
\label{thm:V27-source-total-domination}
The source-total measure has the following exact properties.
\begin{enumerate}[label=\textup{(S\arabic*)},leftmargin=3.5em]
\item $\mathbb G=0$ on $\{t=0\}$ and
      \begin{equation}
       \boxed{
       \int_E\mathbb G\,\dd\mu
       =\int_E\widehat{\mathbb G}\,\dd\nu.}
       \label{eq:V27-Gram-shape-disintegration}
      \end{equation}
\item On $\{t>0\}$,
      \begin{equation}
       \boxed{
       0\preceq\widehat{\mathbb G}\preceq I,
       \qquad \Tr\widehat{\mathbb G}=1.}
       \label{eq:V27-shape-density-bound}
      \end{equation}
\item Let $X,Y\in\{C,A,B\}$ be zero-extended to the common labelled
      coefficient sum and define
      \begin{equation}
       Q_{XY}^{(+)}=(X+Y)^*(X+Y),
       \qquad
       Q_{XY}^{(i)}=(X+iY)^*(X+iY).
       \label{eq:V27-pair-square-writers}
      \end{equation}
      Then the diagonal and pair-square measures are absolutely continuous
      with respect to $\nu$, and their density matrices obey
      \begin{equation}
       \boxed{
       \begin{aligned}
       0&\preceq t^{-1}X^*X\preceq I,\\
       0&\preceq t^{-1}Q_{XY}^{(+)},\ t^{-1}Q_{XY}^{(i)}\preceq2I.
       \end{aligned}}
       \label{eq:V27-nine-square-bounds}
      \end{equation}
\item The measure $\nu$ is the sum of the three diagonal arm-energy measures.
      Hence any positive scalar measure dominating all three diagonal arm
      measures also dominates $\nu$.
\end{enumerate}
Thus the source-total cylinder is a canonical common dominating measure for
all nine positive squares, and their shapes form a cutoff-independent bounded
writer bank.
\end{theorem}

\begin{proof}
A positive finite matrix with zero trace is zero, proving the assertion on
$\{t=0\}$.  Equation~\eqref{eq:V27-Gram-shape-disintegration} is then the
definition of $\nu$ and $\widehat{\mathbb G}$.  A positive trace-one matrix has
all eigenvalues in $[0,1]$, proving
\eqref{eq:V27-shape-density-bound}.

For the pair writers, the Hilbert--Schmidt parallelogram estimate gives
\[
 \Tr Q_{XY}^{(+)}=\norm{X+Y}_{\HS}^{2}
 \le2\norm X_{\HS}^{2}+2\norm Y_{\HS}^{2}\le2t,
\]
and the same estimate holds for $X+iY$.  A positive matrix is bounded above by
its trace times the identity, which proves
\eqref{eq:V27-nine-square-bounds}.  Finally,
\[
 \nu(E)=\int_E\Tr(C^*C)\,\dd\mu
       +\int_E\Tr(A^*A)\,\dd\mu
       +\int_E\Tr(B^*B)\,\dd\mu,
\]
so domination of the three summands implies domination of their sum.
\end{proof}

\begin{remark}[Metric covariance]
\label{rem:V27-metric-covariance}
The trace in \eqref{eq:V27-source-total-writer} is taken after whitening by the
frozen physical coefficient metrics.  It is invariant under metric-unitary
changes of source coordinates.  A nonunitary reparametrization must transport
the coefficient metric with it; otherwise it changes the physical gain rather
than merely changing notation.
\end{remark}

% -----------------------------------------------------------------------------
\subsection{Cell mass, shape, and the homogeneous Schur short}
\label{subsec:V27-scale-shape-short}
% -----------------------------------------------------------------------------

For a physical boundary-record cell $E$, put
\begin{equation}
 \mathbf G_E:=\int_E\mathbb G\,\dd\mu
 =\begin{pmatrix}H_E&W_E\\W_E^*&N_E\end{pmatrix},
 \qquad
 \tau_E:=\Tr\mathbf G_E=\nu(E).
 \label{eq:V27-cell-Gram-mass}
\end{equation}
When $\tau_E>0$, define the trace-one cell shape
\begin{equation}
 \widehat{\mathbf G}_E:=\tau_E^{-1}\mathbf G_E
 =\begin{pmatrix}\widehat H_E&\widehat W_E\\
                  \widehat W_E^*&\widehat N_E\end{pmatrix}.
 \label{eq:V27-cell-shape}
\end{equation}
Let
\begin{equation}
 L_E:=H_E^\dagger W_E,
 \qquad
 D_E:=N_E-W_E^*H_E^\dagger W_E,
 \label{eq:V27-cell-raw-short}
\end{equation}
and define $\widehat L_E,\widehat D_E$ from the hatted blocks in the same way.

\begin{theorem}[Exact gain--shape factorization]
\label{thm:V27-gain-shape-short}
For every cell with $\tau_E>0$,
\begin{equation}
 \boxed{
 L_E=\widehat L_E,
 \qquad
 D_E=\tau_E\widehat D_E.}
 \label{eq:V27-short-homogeneity}
\end{equation}
Moreover,
\begin{equation}
 \boxed{
 0\preceq\widehat D_E\preceq\widehat N_E\preceq I.}
 \label{eq:V27-normalized-short-bound}
\end{equation}
The canonical supports of $H_E$ and $\widehat H_E$ are identical.
\end{theorem}

\begin{proof}
For $\tau_E>0$,
$(\tau_E\widehat H_E)^\dagger=\tau_E^{-1}\widehat H_E^\dagger$.
Substitution into \eqref{eq:V27-cell-raw-short} gives
\[
 H_E^\dagger W_E
 =\tau_E^{-1}\widehat H_E^\dagger(\tau_E\widehat W_E)
 =\widehat L_E
\]
and
\[
 D_E
 =\tau_E\widehat N_E
  -(\tau_E\widehat W_E)^*
    (\tau_E^{-1}\widehat H_E^\dagger)
    (\tau_E\widehat W_E)
 =\tau_E\widehat D_E.
\]
The supported Schur criterion gives $\widehat D_E\succeq0$, and the defining
subtraction gives $\widehat D_E\preceq\widehat N_E$.  Since
$\widehat{\mathbf G}_E$ is positive with trace one,
$\widehat N_E\preceq I$.  Positive scalar multiplication does not change a
support projection.
\end{proof}

\begin{corollary}[Absolute versus projective source verdict]
\label{cor:V27-absolute-projective}
A strict normalized shape floor
$\widehat D_E\succeq\delta P$ gives the absolute source floor
\begin{equation}
 \boxed{D_E\succeq\tau_E\delta P.}
 \label{eq:V27-absolute-floor}
\end{equation}
Consequently normalized ancestry geometry alone proves a noncollapsed physical
source only together with a positive lower bound on $\tau_E$, or when the
declared target metric carries exactly the same scale and that relative
normalization has been frozen in advance.
\end{corollary}

\begin{proof}
Multiply the normalized inequality by $\tau_E$ and use
\eqref{eq:V27-short-homogeneity}.
\end{proof}

% -----------------------------------------------------------------------------
\subsection{The source-weighted record, not the history-weighted record}
\label{subsec:V27-source-weighted-record}
% -----------------------------------------------------------------------------

Let a coarse record cell $E$ be the disjoint union of finer readable cells
$F$.  For $\tau_F>0$, let $\widehat{\mathbf G}_F$ and $\widehat D_F$ be the
corresponding shapes and shorts.  Put
\begin{equation}
 p_F^{\rm src}:=\frac{\tau_F}{\tau_E}.
 \label{eq:V27-source-cell-weight}
\end{equation}
Cells with $\tau_F=0$ are omitted from the sums below.

\begin{theorem}[Source-weighted refinement Pythagoras]
\label{thm:V27-source-weighted-Pythagoras}
The coarse shape is the source-mass barycentre
\begin{equation}
 \boxed{
 \widehat{\mathbf G}_E
 =\sum_{F\subseteq E}p_F^{\rm src}\widehat{\mathbf G}_F.}
 \label{eq:V27-shape-barycentre}
\end{equation}
Its normalized short satisfies
\begin{equation}
 \boxed{
 \widehat D_E
 =\sum_{F\subseteq E}p_F^{\rm src}\widehat D_F
 +\sum_{F\subseteq E}p_F^{\rm src}
  (L_F-L_E)^*\widehat H_F(L_F-L_E).}
 \label{eq:V27-source-weighted-router}
\end{equation}
Thus record refinement is weighted by relative source mass $\tau_F/\tau_E$,
not by the ordinary history probabilities $\mu(F)/\mu(E)$.
\end{theorem}

\begin{proof}
Additivity of the unnormalized Gram measure gives
$\mathbf G_E=\sum_F\mathbf G_F$ and
$\tau_E=\sum_F\tau_F$, which proves
\eqref{eq:V27-shape-barycentre}.  Divide the exact unnormalized refinement
Pythagoras of Version~V26 by $\tau_E$ and use
$H_F=\tau_F\widehat H_F$, $D_F=\tau_F\widehat D_F$, and the decoder invariance
of Theorem~\ref{thm:V27-gain-shape-short}.
\end{proof}

\begin{firewall}[Cellwise normalization erases a physical router weight]
\label{fire:V27-no-cellwise-normalization}
Normalizing every cell separately retains $\widehat{\mathbf G}_F$ but erases
the source weights $p_F^{\rm src}$.  Those weights are needed both for the
coarse decoder and for the router term in
\eqref{eq:V27-source-weighted-router}.  They cannot be restored from the
ordinary record probabilities unless a same-history source-total incidence
law has independently been supplied.
\end{firewall}

% -----------------------------------------------------------------------------
\subsection{Global gain and relative cell gains are different rows}
\label{subsec:V27-gain-gauge}
% -----------------------------------------------------------------------------

\begin{theorem}[Common-gain gauge and absolute source scale]
\label{thm:V27-common-gain}
For $g>0$, replace all three arms by
\begin{equation}
 (C,A,B)\longmapsto\sqrt g\,(C,A,B).
 \label{eq:V27-common-gain-action}
\end{equation}
Then
\begin{equation}
 \boxed{
 \mathbb G\mapsto g\mathbb G,
 \quad
 \nu\mapsto g\nu,
 \quad
 \widehat{\mathbb G}\mapsto\widehat{\mathbb G},
 \quad
 L_E\mapsto L_E,
 \quad
 D_E\mapsto gD_E.}
 \label{eq:V27-common-gain-effects}
\end{equation}
Every normalized nine-square shape and every homogeneous held-out consistency
ratio is unchanged.  Therefore one absolute gain anchor is necessary to turn
the projective square packet into an absolute source-action packet.
\end{theorem}

\begin{proof}
Every Gram and square writer is quadratic in the arms.  The source-total writer
and measure therefore scale by $g$, while division by $t$ removes the factor.
The decoder and short statements are Theorem~\ref{thm:V27-gain-shape-short}.
Any homogeneous degree-zero consistency ratio is unchanged by common scalar
multiplication.
\end{proof}

The common gain above is a calibration coordinate.  A separate gain for every
record cell is not a harmless gauge because it changes the physical
source-weighted mixture.

\begin{countertheorem}[Relative cell gains change a pure router]
\label{cth:V27-cell-gain-router}
Consider two scalar readable cells with exact follower cards
\begin{equation}
 (M_1,\Gamma_1,N_1)=(1,1,1),
 \qquad
 (M_2,\Gamma_2,N_2)=(1,-1,1).
 \label{eq:V27-two-cell-follower-shapes}
\end{equation}
After multiplying the two cell Grams by gains $g_1,g_2>0$, both cell shorts
remain zero, but the record-blind short is
\begin{equation}
 \boxed{
 D_{\rm blind}(g_1,g_2)
 =g_1+g_2-\frac{(g_1-g_2)^2}{g_1+g_2}
 =\frac{4g_1g_2}{g_1+g_2}.}
 \label{eq:V27-relative-gain-router}
\end{equation}
Hence separately normalized cell packets do not determine even a pure-router
coarse short.
\end{countertheorem}

\begin{proof}
The scaled aggregate blocks are
$M=N=g_1+g_2$ and $\Gamma=g_1-g_2$.  Their scalar Schur complement is the
first expression in \eqref{eq:V27-relative-gain-router}; simplification gives
the second.  Each cell short is
$g_j-g_j^2/g_j=0$.
\end{proof}

% -----------------------------------------------------------------------------
\subsection{Why probability-tail control is not source-response tightness}
\label{subsec:V27-weighted-tail}
% -----------------------------------------------------------------------------

For any event $E$, positivity and the source-total definition give the operator
bound
\begin{equation}
 \boxed{
 0\preceq\mathbf G(E)\preceq\nu(E)I,
 \qquad
 \norm{\mathbf G(E)}_{\mathrm{op}}\le\nu(E).}
 \label{eq:V27-Gram-tail-bound}
\end{equation}
Thus the omitted source response is controlled by source-total mass, not by
probability alone.

\begin{countertheorem}[A summable probability debit can carry infinite source debit]
\label{cth:V27-probability-not-source-tail}
For $n\ge0$, let $E_n$ have probability
\begin{equation}
 d_n=\frac7{1024}\,2^{-n},
 \qquad
 \sum_{n\ge0}d_n=\frac7{512}.
 \label{eq:V27-probability-debits}
\end{equation}
Let the complete Gram vanish off $E_n$ and have source-total writer
\begin{equation}
 t_n=d_n^{-1}
 \quad\text{on }E_n.
 \label{eq:V27-tail-amplification}
\end{equation}
Then
\begin{equation}
 \boxed{\nu_n(E_n)=1\quad\text{for every }n,}
 \label{eq:V27-source-debits-one}
\end{equation}
so
$\sum_n\nu_n(E_n)=\infty$ although the complete probability debit equals
$7/512$.
\end{countertheorem}

\begin{proof}
By definition,
$\nu_n(E_n)=\int_{E_n}t_n\,\dd\mu_n=d_nt_n=1$.
\end{proof}

\begin{theorem}[Three sufficient source-tail upgrades]
\label{thm:V27-source-tail-upgrades}
Let $E_n$ be the path or compact-layer omission at step $n$, let
$d_n=\mu_n(E_n)$, and let $t_n$ be the source-total writer.
Each of the following independently upgrades probability control to source
control.
\begin{enumerate}[label=\textup{(T\arabic*)},leftmargin=3.5em]
\item If $t_n\le T_*$ almost surely for all $n$, then
      \begin{equation}
       \boxed{\nu_n(E_n)\le T_*d_n.}
       \label{eq:V27-uniform-envelope-tail}
      \end{equation}
\item If $p>1$ and $\norm{t_n}_{L^p(\mu_n)}\le C_p$ uniformly, then
      \begin{equation}
       \boxed{
       \nu_n(E_n)\le C_p d_n^{1-1/p}.}
       \label{eq:V27-Lp-tail}
      \end{equation}
\item For every $R>0$,
      \begin{equation}
       \boxed{
       \nu_n(E_n)
       \le Rd_n+\int_{\{t_n>R\}}t_n\,\dd\mu_n.}
       \label{eq:V27-UI-tail}
      \end{equation}
      Hence schedules $R_n\uparrow\infty$ satisfying
      \begin{equation}
       \sum_n\left[
        R_nd_n+\int_{\{t_n>R_n\}}t_n\,\dd\mu_n
       \right]<\infty
       \label{eq:V27-source-tail-schedule}
      \end{equation}
      give summable source-total loss.
\end{enumerate}
\end{theorem}

\begin{proof}
The first statement is immediate.  The second is H\"older's inequality:
\[
 \int_{E_n}t_n\,\dd\mu_n
 \le\norm{t_n}_{L^p}\,\mu_n(E_n)^{1-1/p}.
\]
For the third, split $E_n$ into $\{t_n\le R\}$ and $\{t_n>R\}$.
Summing \eqref{eq:V27-UI-tail} with $R=R_n$ gives the last claim.
\end{proof}

\begin{assessment}[Exact qualification of the V26 transport statement]
\label{ass:V27-V26-qualification}
The block-transport conclusion of Theorem~\ref{thm:V26-square-transport} is
valid on either of the following branches:
\begin{enumerate}[label=\textup{(Q\arabic*)},leftmargin=3.5em]
\item all square writers are normalized by one cutoff-independent physical
      scale, so the inverse reconstruction map is genuinely fixed;
\item the cutoff-dependent scale or source-total measure is transported with a
      summable debit satisfying Theorem~\ref{thm:V27-source-tail-upgrades}.
\end{enumerate}
Per-cutoff normalization to $[0,1]$ without one of these rows transports only
the projective shape packet.  It does not imply transport of the unnormalized
Gram or its absolute Schur short.
\end{assessment}

% -----------------------------------------------------------------------------
\subsection{Exact scale--shape transport}
\label{subsec:V27-scale-shape-transport}
% -----------------------------------------------------------------------------

Let $\mathbf G_n^\oplus$ be the direct sum of all unnormalized cell Grams after
transport to one fixed boundary screen, and put
\begin{equation}
 \tau_n:=\Tr\mathbf G_n^\oplus,
 \qquad
 \widehat{\mathbf G}_n^\oplus:=\tau_n^{-1}\mathbf G_n^\oplus
 \quad(\tau_n>0).
 \label{eq:V27-global-scale-shape}
\end{equation}
This \emph{global} normalization retains the relative source masses of all
record cells.  Let $\widehat D_n$ be the Schur short of
$\widehat{\mathbf G}_n^\oplus$.  Equivalently, if
$\widehat D_{n,E}^{\rm cell}$ denotes the trace-one short of cell $E$, then
\begin{equation}
 \boxed{
 \widehat D_n
 =\bigoplus_E\frac{\tau_{n,E}}{\tau_n}
   \widehat D_{n,E}^{\rm cell},
 \qquad
 D_n=\tau_n\widehat D_n.}
 \label{eq:V27-global-short-scaling}
\end{equation}

Suppose the normalized positive-square coordinates have adjacent defect
$\epsilon_n^{\rm sh}$ and reconstruct normalized blocks with
\begin{equation}
 \max\{\norm{\Delta\widehat H_n},
       \norm{\Delta\widehat W_n},
       \norm{\Delta\widehat N_n}\}
 \le\kappa_{\mathcal T}\epsilon_n^{\rm sh}.
 \label{eq:V27-shape-block-defect}
\end{equation}
Assume one common normalized support and
\begin{equation}
 \widehat H_n\succeq\widehat h_*P,
 \qquad
 \norm{\widehat W_n}\le\widehat w_*.
 \label{eq:V27-normalized-short-domain}
\end{equation}
Define
\begin{equation}
 \widehat\Lambda_*
 :=\kappa_{\mathcal T}
 \left(
 1+\frac{2\widehat w_*}{\widehat h_*}
  +\frac{\widehat w_*^2}{\widehat h_*^2}
 \right).
 \label{eq:V27-normalized-Lipschitz}
\end{equation}

\begin{theorem}[Scale--shape adjacent identity and cofinal landing]
\label{thm:V27-scale-shape-landing}
Under \eqref{eq:V27-shape-block-defect}--
\eqref{eq:V27-normalized-short-domain},
\begin{equation}
 \boxed{
 \norm{\widehat D_{n+1}-\widehat D_n}
 \le\widehat\Lambda_*\epsilon_n^{\rm sh}.}
 \label{eq:V27-shape-short-transport}
\end{equation}
The unnormalized shorts satisfy the exact estimate
\begin{equation}
 \boxed{
 \norm{D_{n+1}-D_n}
 \le
 \tau_{n+1}\widehat\Lambda_*\epsilon_n^{\rm sh}
 +|\tau_{n+1}-\tau_n|.}
 \label{eq:V27-absolute-short-transport}
\end{equation}
If
\begin{equation}
 \boxed{
 \mathfrak e_{n_0}^{\rm abs}
 :=\sum_{n\ge n_0}
 \left[
  \tau_{n+1}\widehat\Lambda_*\epsilon_n^{\rm sh}
  +|\tau_{n+1}-\tau_n|
 \right]<\infty,}
 \label{eq:V27-absolute-tail}
\end{equation}
then $D_n$ has a unique operator-norm limit $D_\infty$ and
\begin{equation}
 \boxed{
 \norm{D_\infty-D_{n_0}}
 \le\mathfrak e_{n_0}^{\rm abs}.}
 \label{eq:V27-corrected-one-anchor-radius}
\end{equation}
Consequently:
\begin{enumerate}[label=\textup{(A\arabic*)},leftmargin=3.5em]
\item if $P_QD_{n_0}P_Q\succeq
      (\delta+\mathfrak e_{n_0}^{\rm abs})P_Q$, then
      $P_QD_\infty P_Q\succeq\delta P_Q$;
\item if $\norm{D_{n_0}}+\mathfrak e_{n_0}^{\rm abs}\le\varepsilon$, then
      $\norm{D_\infty}\le\varepsilon$.
\end{enumerate}
\end{theorem}

\begin{proof}
The normalized supported Lipschitz estimate is
Lemma~\ref{lem:V26-short-Lipschitz} applied to the hatted blocks, which gives
\eqref{eq:V27-shape-short-transport}.  By
Theorem~\ref{thm:V27-gain-shape-short},
$D_n=\tau_n\widehat D_n$, so
\[
 D_{n+1}-D_n
 =\tau_{n+1}(\widehat D_{n+1}-\widehat D_n)
  +(\tau_{n+1}-\tau_n)\widehat D_n.
\]
Equation~\eqref{eq:V27-normalized-short-bound} gives
$\norm{\widehat D_n}\le1$.  This proves
\eqref{eq:V27-absolute-short-transport}.  Summation proves that $(D_n)$ is
Cauchy and gives \eqref{eq:V27-corrected-one-anchor-radius}.  The two interval
claims follow from the triangle inequality and Weyl monotonicity.
\end{proof}

\begin{corollary}[The exact scale alternatives]
\label{cor:V27-scale-alternatives}
Assume the normalized shape shorts converge to $\widehat D_\infty$.
\begin{enumerate}[label=\textup{(G\arabic*)},leftmargin=3.5em]
\item If $\tau_n\to\tau_\infty\in(0,\infty)$, then
      $D_n\to\tau_\infty\widehat D_\infty$.
\item If $\tau_n\to0$, then $D_n\to0$ regardless of a positive normalized
      shape floor.
\item If $(\tau_n)$ is unbounded or has no Cauchy subsequence compatible with
      the selected route, the output is an absolute source-gain escape rather
      than a cofinal source birth or follower verdict.
\end{enumerate}
\end{corollary}

\begin{proof}
The first two statements follow from
$D_n=\tau_n\widehat D_n$ and
$\norm{\widehat D_n}\le1$.  The third is the failure of the absolute
calibration component required by
\eqref{eq:V27-absolute-short-transport}; projective shape convergence cannot
supply it.
\end{proof}

\begin{countertheorem}[One normalized anchor does not determine absolute ancestry]
\label{cth:V27-shape-no-absolute}
Let the normalized complete Gram be the constant matrix
\begin{equation}
 \widehat{\mathbf G}
 =\begin{pmatrix}
  1/3&0&1/6\\
  0&1/3&0\\
  1/6&0&1/3
 \end{pmatrix}\succ0.
 \label{eq:V27-control-shape}
\end{equation}
Treat its first two coordinates as the complete old block and its third as the
line arm.  Then
\begin{equation}
 \boxed{\widehat D=\frac14.}
 \label{eq:V27-control-shape-short}
\end{equation}
The three sequences
\begin{equation}
 \tau_n=1,
 \qquad
 \tau_n=\frac1{n+1},
 \qquad
 \tau_n=n+1
 \label{eq:V27-three-gain-sequences}
\end{equation}
have identical normalized nine-square packets and zero shape-transport defect,
but their absolute shorts respectively stay at $1/4$, converge to zero, and
diverge.  Even agreement of all three sequences at one chosen anchor can be
forced before their later gains separate.
\end{countertheorem}

\begin{proof}
The old block is $\widehat H=\diag(1/3,1/3)$, the mixed column is
$\widehat W=(1/6,0)^{\mathsf T}$, and $\widehat N=1/3$.  Hence
\[
 \widehat D=\frac13-\frac{(1/6)^2}{1/3}=\frac14.
\]
All normalized positive squares depend only on
$\widehat{\mathbf G}$, while the absolute short is
$D_n=\tau_n/4$.  The final statement follows by taking gain sequences which
coincide through the selected anchor and then follow the three displayed
tails.
\end{proof}

% -----------------------------------------------------------------------------
\subsection{Exact rational replay}
\label{subsec:V27-replay}
% -----------------------------------------------------------------------------

For the shape \eqref{eq:V27-control-shape}, the three diagonal arm squares are
\begin{equation}
 G_C=G_A=G_B=\frac13.
 \label{eq:V27-replay-diagonals}
\end{equation}
The six pair squares of Version~V25 are
\begin{equation}
 \boxed{
 \begin{aligned}
 Q_{CA}^{(+)}&=Q_{CA}^{(i)}=\frac23,\\
 Q_{CB}^{(+)}&=1,
 &Q_{CB}^{(i)}&=\frac23,\\
 Q_{AB}^{(+)}&=Q_{AB}^{(i)}=\frac23.
 \end{aligned}}
 \label{eq:V27-replay-nine-squares}
\end{equation}
The held-out square for coefficients $(1,1,i)$ is exactly one.

At absolute gain $\tau=12$, the complete Gram is
\begin{equation}
 \mathbf G_{12}
 =\begin{pmatrix}
 4&0&2\\
 0&4&0\\
 2&0&4
 \end{pmatrix},
 \qquad
 \boxed{D_{12}=3.}
 \label{eq:V27-replay-gain12}
\end{equation}
At gain $\tau=15$, the shape is unchanged while
\begin{equation}
 \boxed{D_{15}=\frac{15}{4},
 \qquad
 |D_{15}-D_{12}|=\frac34.}
 \label{eq:V27-replay-gain15}
\end{equation}
Thus a zero shape defect coexists with a nonzero absolute source defect equal
to the gain increment times the normalized short.

For the two-cell exact-follower control of
Countertheorem~\ref{cth:V27-cell-gain-router}, the choices
$(g_1,g_2)=(1,1)$ and $(1,3)$ give
\begin{equation}
 \boxed{
 D_{\rm blind}(1,1)=2,
 \qquad
 D_{\rm blind}(1,3)=3.}
 \label{eq:V27-replay-cell-gains}
\end{equation}
Every cell short remains zero in both cards.  The change is entirely a
relative-gain router effect.

\begin{theorem}[Version V27 source-total calibration theorem]
\label{thm:V27-terminal}
Versions~V1--V26 are retained, with the transport clause of Version~V26 read on
the fixed absolute-scale branch or with the source-total correction below.
Version~V27 additionally proves:
\begin{enumerate}[label=\textup{(V27.\arabic*)},leftmargin=5.5em]
\item the complete three-arm Gram has a canonical source-total positive measure
      whose nine square densities are uniformly bounded;
\item every record-cell Gram and final short factor exactly into absolute
      source mass times a trace-one shape, while the decoder is scale free;
\item record refinement is a Pythagoras theorem under source-mass weights rather
      than ordinary history probabilities;
\item one global gain is an independent absolute calibration row, whereas
      uncalibrated cellwise gains alter the router itself;
\item a summable retained-terminal probability debit does not imply summable
      source response, with exact uniform-envelope, $L^p$, and
      uniform-integrability upgrades;
\item normalized shape transport and absolute gain transport combine through
      the exact adjacent estimate
      \eqref{eq:V27-absolute-short-transport};
\item one anchor controls the cofinal absolute short only through the corrected
      source-weighted tail \eqref{eq:V27-absolute-tail};
\item identical normalized square packets admit absolute survivor, collapse,
      and gain-escape branches.
\end{enumerate}
\end{theorem}

\begin{proof}
Items~\textup{(V27.1)--(V27.3)} are
Theorems~\ref{thm:V27-source-total-domination}--
\ref{thm:V27-source-weighted-Pythagoras}.  Item~\textup{(V27.4)} is
Theorem~\ref{thm:V27-common-gain} and
Countertheorem~\ref{cth:V27-cell-gain-router}.  Item~\textup{(V27.5)} is
Countertheorem~\ref{cth:V27-probability-not-source-tail} and
Theorem~\ref{thm:V27-source-tail-upgrades}.  Items~\textup{(V27.6)--(V27.8)}
are Theorem~\ref{thm:V27-scale-shape-landing},
Corollary~\ref{cor:V27-scale-alternatives}, and
Countertheorem~\ref{cth:V27-shape-no-absolute}.
\end{proof}

\begin{assessment}[Version V27 verdict]
\label{ass:V27-verdict}
The current result is
\begin{equation}
 \boxed{
 \begin{gathered}
 \passstatus\\[-0.15em]
 \text{canonical source-total cylinder, bounded shape packet,}\\
 \text{source-weighted record Pythagoras, gain--shape short,}\\
 \text{and corrected one-anchor transport}
 \\[0.45em]
 \obsstatus\\[-0.15em]
 \text{per-cutoff or per-cell normalization used as absolute calibration;}\\
 \text{a probability path debit used as an unnormalized source debit;}\\
 \text{or a positive normalized short promoted through source-total collapse}
 \\[0.45em]
 \stopstatus\\[-0.15em]
 \text{one RPESM source-total square cylinder, its absolute gain cocycle,}\\
 \text{source-weighted compact-layer tail, physical line sewing,}\\
 \text{and one bounded shape/held-out square packet.}
 \end{gathered}}
 \label{eq:V27-verdict}
\end{equation}
The first missing row is therefore an absolute source-response calibration and
tightness row, not another phase, crossing, record, or Schur theorem.
\end{assessment}

\begin{acquisition}[Version V28: source-total gain and response-tight anchor]
\label{acq:V27-next}
Preserve Versions~V1--V27.  On one physically readable finite boundary record
and one selected amplitude semantics, acquire the positive source-total writer
\begin{equation}
 t=\norm C_{\HS}^2+\norm A_{\HS}^2+\norm B_{\HS}^2
 \label{eq:V27-next-total-writer}
\end{equation}
on the same retained-terminal history as the nine positive squares.  Populate
its unnormalized cell masses $\tau_E$, one absolute gain anchor, and either:
\begin{enumerate}[label=\textup{(V28.\arabic*)},leftmargin=5.5em]
\item a uniform pointwise source-total envelope;
\item a source-total $L^p$, $p>1$, bound together with a sufficiently fast
      selected path debit;
\item or a direct source-weighted tail schedule
      \eqref{eq:V27-source-tail-schedule}.
\end{enumerate}
Use $\nu=t\mu$ to acquire the bounded trace-one shape densities and one held-out
three-arm square.  Reconstruct the unnormalized Gram from source mass times
shape, run the complete primitive and record-router shorts, and compare the
finite margin only with the corrected absolute tail
\eqref{eq:V27-absolute-tail}.  Append only the polar range of the final
within-record short.

A common unknown gain must be fixed by one physical absolute square or
partition calibration.  If cells are acquired in separate runs, their relative
gains require a same-instrument calibration before aggregation.  Failure of
source-total uniform integrability is returned as a source-response escape and
is not repaired by choosing a smaller probability tail.

Do not open the conditional exterior, locality, current/stress, charge,
Duhamel, or Einstein gates until this source-total anchor has landed.
\end{acquisition}

\section*{V27 exact replay and append-only record}
\addcontentsline{toc}{section}{V27 exact replay and append-only record}
The accompanying exact verifier checks the trace-one shape, all nine positive
square values, the held-out square, Schur homogeneity at two absolute gains,
the three scale alternatives, the cellwise-gain router, and the summable
probability/non-summable source-tail counterexample.  These are theorem replays
and are not numerical measurements of the selected RPESM regulator.

The complete Version~V26 source preceding its sole final active
\verb|\end{document}| is preserved byte for byte.  Its preterminal SHA--256
digest is recorded in the validation manifest.  A later continuation must
remove only the final active document-closing command, append new material
immediately before it, restore one terminal command, and increment the filename
to end in \texttt{\_V28.tex}.

% =============================================================================
% END APPENDED VERSION V27 MATERIAL
% =============================================================================


% =============================================================================
% BEGIN APPENDED VERSION V28 MATERIAL
% =============================================================================

\clearpage
\clearpage
\hypersetup{
  pdftitle={Renewal Einstein--Standard--Model Physical Source Metric, Duhamel Ancestry and Quantum Continuum Programme V28},
  pdfsubject={Physical command incidence, metric-native source total, reserve-safe tails, and direct cofinal transport}
}
\pdfinfo{
 /Title (Renewal Einstein--Standard--Model Physical Source Metric, Duhamel Ancestry and Quantum Continuum Programme V28)
 /Subject (Physical command incidence, metric-native source total, reserve-safe tails, and direct cofinal transport)
}

\section{Version V28: physical command incidence before source-total gain, metric-native response, and reserve-safe transport}
\label{sec:V28-command-incidence}

\subsection{The whitening premise hidden in the Version V27 acquisition}
\label{subsec:V28-context}

Version~V27 correctly separated an absolute source-total mass from a bounded
trace-one shape and proved that a probability tail need not control the
corresponding source-response tail.  Its construction began, however, only
\emph{after the physical coefficient metrics had been frozen and whitened}.
That phrase contains the first still-unpopulated physical row.

The three arms
\[
 C,\qquad A,\qquad B
\]
are role-labelled copies used to expose prior-primitive, phase-quenched
crossing, and line-weighted crossing incidence.  Their direct-sum Gram is the
correct positive acquisition object.  It does not by itself say which linear
combinations are generated by one physical command, how those commands are
priced, or whether a zero-cost command has a nonzero response.  Consequently
\[
 \Tr(C\;A\;B)^*(C\;A\;B)
\]
is a canonical \emph{role-envelope mass} after a role metric has been chosen,
but it is not automatically the physical command-source action.

The present version inserts the missing map
\begin{equation}
 \boxed{
 \begin{gathered}
 \text{physical command space and cost}
 \\[-0.2em]\mathop{\Downarrow}^{J}\\[-0.2em]
 \text{primitive--crossing--line role space}
 \\[-0.2em]\mathop{\Downarrow}^{\mathcal X}\\[-0.2em]
 \text{reflected source carrier}
 \end{gathered}}
 \label{eq:V28-command-incidence-chain}
\end{equation}
It imports the null-cost and generalized-reserve theorem of
\Cref{thm:metric-reserve-matrix}; it does not reprove the earlier single-bank
result.  The new work is the assembly of that theorem with the three-arm role
packet, the separation of range visibility from incidence gain, and the
resulting correction to the Version~V27 source-total tail.

The corrected order is
\begin{equation}
 \boxed{
 \begin{gathered}
 \text{physical command occurrence and source-action metric}
 \\[-0.1em]\Downarrow\\[-0.1em]
 \text{command-to-role incidence and null-cost test}
 \\[-0.1em]\Downarrow\\[-0.1em]
 \text{metric-whitened physical command Gram and reserve}
 \\[-0.1em]\Downarrow\\[-0.1em]
 \text{role-envelope domination plus incidence-gain control}
 \\[-0.1em]\Downarrow\\[-0.1em]
 \text{source-weighted tail and direct cofinal command transport}.
 \end{gathered}}
 \label{eq:V28-order}
\end{equation}

% -----------------------------------------------------------------------------
\subsection{The physical command-incidence packet}
\label{subsec:V28-command-packet}
% -----------------------------------------------------------------------------

Let
\begin{equation}
 \mathcal E_{\rm role}
 =\mathcal E_C\oplus\mathcal E_A\oplus\mathcal E_B,
 \qquad
 \mathcal X=(C\ \ A\ \ B):\mathcal E_{\rm role}\to\mathcal H,
 \label{eq:V28-role-synthesis}
\end{equation}
and let
\begin{equation}
 \mathbb G=\mathcal X^*\mathcal X\succeq0
 \label{eq:V28-role-Gram}
\end{equation}
be the complete role Gram of Versions~V25--V27.  Let
\(\mathcal U\) be the physically occurring command coefficient space,
\begin{equation}
 J:\mathcal U\longrightarrow\mathcal E_{\rm role}
 \label{eq:V28-incidence-map}
\end{equation}
its same-history role-incidence map, and
\begin{equation}
 \mathsf M_{\rm cmd}=\mathsf M_{\rm cmd}^*\succeq0
 \label{eq:V28-command-metric}
\end{equation}
the physical source-action metric.  The actual source synthesis is
\begin{equation}
 \boxed{Y:=\mathcal XJ.}
 \label{eq:V28-physical-synthesis}
\end{equation}
Put
\begin{equation}
 P_{\mathsf M}:=\mathsf M_{\rm cmd}\mathsf M_{\rm cmd}^\dagger
 \label{eq:V28-command-support}
\end{equation}
and define the null-cost incidence residual
\begin{equation}
 \boxed{
 \Delta_{0}^{\rm cmd}
 :=\norm{Y(I-P_{\mathsf M})}_{\HS}^{2}.}
 \label{eq:V28-null-cost-residual}
\end{equation}

\begin{theorem}[Physical command null-cost and metric-native reserve]
\label{thm:V28-command-reserve}
The following are equivalent:
\begin{equation}
 \boxed{
 \Delta_{0}^{\rm cmd}=0
 \iff
 \Ker\mathsf M_{\rm cmd}\subseteq\Ker Y.}
 \label{eq:V28-null-cost-equivalence}
\end{equation}
On this branch define the whitened physical command synthesis and its Gram by
\begin{equation}
 \boxed{
 Z:=Y\mathsf M_{\rm cmd}^{\dagger/2},
 \qquad
 Q_{\rm cmd}:=Z^*Z
 =\mathsf M_{\rm cmd}^{\dagger/2}
  J^*\mathbb GJ
  \mathsf M_{\rm cmd}^{\dagger/2}.}
 \label{eq:V28-command-Gram}
\end{equation}
Then:
\begin{enumerate}[label=\textup{(C\arabic*)},leftmargin=3.5em]
\item the nonzero spectrum of \(Q_{\rm cmd}\) is the generalized reserve
      spectrum of
      \begin{equation}
       J^*\mathbb GJ\,u=\gamma\mathsf M_{\rm cmd}u
       \quad\text{on }\mathcal U/\Ker\mathsf M_{\rm cmd};
       \label{eq:V28-generalized-command-reserve}
      \end{equation}
\item its total physical command action is
      \begin{equation}
       \boxed{
       \tau_{\rm cmd}:=\Tr Q_{\rm cmd}=\norm Z_{\HS}^{2};}
       \label{eq:V28-command-total}
      \end{equation}
\item its supported lower and upper reserves are
      \begin{equation}
       \boxed{
       g_-^{\rm cmd}:=\lambda_{\min}^{+}(Q_{\rm cmd}),
       \qquad
       g_+^{\rm cmd}:=\norm{Q_{\rm cmd}}_{\mathrm{op}}.}
       \label{eq:V28-command-reserve-window}
      \end{equation}
\end{enumerate}
The role trace \(\Tr\mathbb G\), the command total
\(\tau_{\rm cmd}\), and the reserve floor \(g_-^{\rm cmd}\) are three
logically distinct quantities.
\end{theorem}

\begin{proof}
The orthogonal projection onto \((\Ker\mathsf M_{\rm cmd})^\perp\) is
\(P_{\mathsf M}\).  Hence
\(Y(I-P_{\mathsf M})=0\) exactly when \(Y\) vanishes on
\(\Ker\mathsf M_{\rm cmd}\), proving
\eqref{eq:V28-null-cost-equivalence}.  On this branch
\Cref{thm:metric-reserve-matrix}, applied to the single physically assembled
synthesis \(Y=\mathcal XJ\), gives the generalized spectrum and ambient
unit-action interpretation.  The trace identity follows from
\(Q_{\rm cmd}=Z^*Z\).  The final quantities are definitions on the canonical
support of that positive Gram.
\end{proof}

\begin{countertheorem}[Pseudowhitening does not repair a null-cost source]
\label{cth:V28-null-cost-pseudoinverse}
Take
\begin{equation}
 \mathsf M_{\rm cmd}=\begin{pmatrix}1&0\\0&0\end{pmatrix},
 \qquad
 Y=\begin{pmatrix}0&1\end{pmatrix}.
 \label{eq:V28-null-cost-control}
\end{equation}
Then
\begin{equation}
 \boxed{
 \Delta_0^{\rm cmd}=1,
 \qquad
 \norm{Y\mathsf M_{\rm cmd}^{\dagger/2}}_{\HS}^{2}=0.}
 \label{eq:V28-null-cost-control-values}
\end{equation}
Thus forming the Moore--Penrose-whitened Gram before testing
\eqref{eq:V28-null-cost-equivalence} can erase a live zero-cost response and
return a false zero source total.
\end{countertheorem}

\begin{proof}
Here \(I-P_{\mathsf M}=\diag(0,1)\), so
\(Y(I-P_{\mathsf M})=Y\) has squared norm one.  On the other hand
\(\mathsf M_{\rm cmd}^{\dagger/2}=\diag(1,0)\), and the displayed product is
zero.
\end{proof}

% -----------------------------------------------------------------------------
\subsection{Role coordinates, command coordinates, and physical covariance}
\label{subsec:V28-coordinate-covariance}
% -----------------------------------------------------------------------------

\begin{theorem}[Two-sided role/command coordinate covariance]
\label{thm:V28-two-sided-covariance}
Let
\begin{equation}
 S:\widetilde{\mathcal E}_{\rm role}\to\mathcal E_{\rm role},
 \qquad
 R:\widetilde{\mathcal U}\to\mathcal U
 \label{eq:V28-coordinate-isomorphisms}
\end{equation}
be invertible coordinate maps.  Set
\begin{equation}
 \widetilde{\mathcal X}=\mathcal XS,
 \qquad
 \widetilde J=S^{-1}JR,
 \qquad
 \widetilde{\mathsf M}_{\rm cmd}=R^*\mathsf M_{\rm cmd}R.
 \label{eq:V28-coordinate-transform}
\end{equation}
Then
\begin{equation}
 \widetilde{\mathcal X}\widetilde J=YR.
 \label{eq:V28-physical-transform}
\end{equation}
The null-cost branch is preserved, and the metric-normalized command Grams are
unitarily equivalent on their supports.  In particular
\begin{equation}
 \boxed{
 \widetilde\tau_{\rm cmd}=\tau_{\rm cmd},
 \qquad
 \widetilde g_\pm^{\rm cmd}=g_\pm^{\rm cmd}.}
 \label{eq:V28-coordinate-invariants}
\end{equation}
By contrast, \(\Tr\widetilde{\mathbb G}=\Tr(S^*\mathbb GS)\) need not equal
\(\Tr\mathbb G\) unless the role-coordinate change is metric unitary.
\end{theorem}

\begin{proof}
Equation~\eqref{eq:V28-physical-transform} is direct substitution.  The pairs
\((Y,\mathsf M_{\rm cmd})\) and
\((YR,R^*\mathsf M_{\rm cmd}R)\) are congruent physical source/action pairs.
The invariance statement is therefore item~\textup{(R4)} of
\Cref{thm:metric-reserve-matrix}; the trace is the sum of the generalized
reserve eigenvalues with multiplicity.  The final observation follows by
choosing a nonunitary \(S\).
\end{proof}

\begin{countertheorem}[Raw Hilbert--Schmidt mass is a coordinate gauge]
\label{cth:V28-raw-coordinate-gauge}
In one dimension take \(Y=3\) and physical cost \(\mathsf M_{\rm cmd}=4\).
The metric-native total is
\begin{equation}
 \boxed{\tau_{\rm cmd}=\frac94.}
 \label{eq:V28-one-dimensional-total}
\end{equation}
Under the coordinate rescaling \(u=5\widetilde u\), the synthesis and metric
become \(\widetilde Y=15\) and
\(\widetilde{\mathsf M}_{\rm cmd}=100\).  The metric-native total remains
\(9/4\), whereas the raw squared coefficient norm changes from \(9\) to
\(225\).
\end{countertheorem}

\begin{proof}
Both physical totals are the response square divided by the command cost:
\(3^2/4=15^2/100=9/4\).
\end{proof}

% -----------------------------------------------------------------------------
\subsection{Range visibility and incidence gain inside the role envelope}
\label{subsec:V28-role-envelope}
% -----------------------------------------------------------------------------

On the null-cost branch put
\begin{equation}
 \mathcal I:=J\mathsf M_{\rm cmd}^{\dagger/2}.
 \label{eq:V28-whitened-incidence}
\end{equation}
Let
\begin{equation}
 \mathcal I=V_{\rm inc}R_{\rm inc}
 \label{eq:V28-incidence-polar}
\end{equation}
be its polar decomposition, where
\(R_{\rm inc}=(\mathcal I^*\mathcal I)^{1/2}\) and
\(P_{\rm inc}:=V_{\rm inc}V_{\rm inc}^*\) is the projection onto the role
subspace reached by unit-cost physical commands.

Define
\begin{equation}
 \tau_{\rm role}:=\Tr\mathbb G,
 \qquad
 \tau_{\rm vis}:=\Tr(V_{\rm inc}^*\mathbb GV_{\rm inc}),
 \label{eq:V28-role-visible-masses}
\end{equation}
and
\begin{equation}
 \tau_{\perp}:=
 \Tr\bigl((I-P_{\rm inc})\mathbb G(I-P_{\rm inc})\bigr).
 \label{eq:V28-role-perp-mass}
\end{equation}

\begin{theorem}[Role-envelope Pythagoras and incidence-gain sandwich]
\label{thm:V28-role-envelope-Pythagoras}
The role envelope splits exactly as
\begin{equation}
 \boxed{\tau_{\rm role}=\tau_{\rm vis}+\tau_{\perp}.}
 \label{eq:V28-role-envelope-split}
\end{equation}
Moreover,
\begin{equation}
 \boxed{
 \tau_{\rm cmd}
 =\Tr\!
 \left(
 R_{\rm inc}V_{\rm inc}^*\mathbb GV_{\rm inc}R_{\rm inc}
 \right).}
 \label{eq:V28-command-gain-formula}
\end{equation}
If the nonzero singular values of \(\mathcal I\) lie in
\([r_-,r_+]\), then
\begin{equation}
 \boxed{
 r_-^2\tau_{\rm vis}
 \le\tau_{\rm cmd}
 \le r_+^2\tau_{\rm vis}.}
 \label{eq:V28-incidence-gain-sandwich}
\end{equation}
Finally,
\begin{equation}
 \boxed{
 \tau_{\perp}=0
 \iff
 \mathbb G=P_{\rm inc}\mathbb GP_{\rm inc}.}
 \label{eq:V28-role-support-test}
\end{equation}
Thus the Version~V27 role trace carries two independent debits: role response
outside the actual command incidence range and gain distortion inside that
range.
\end{theorem}

\begin{proof}
Let \(Q=I-P_{\rm inc}\).  Since \(P_{\rm inc}Q=0\), cyclicity of trace gives
\[
 \Tr\mathbb G
 =\Tr(P_{\rm inc}\mathbb GP_{\rm inc})
  +\Tr(Q\mathbb GQ).
\]
The first trace equals
\(\Tr(V_{\rm inc}^*\mathbb GV_{\rm inc})\), proving
\eqref{eq:V28-role-envelope-split}.  Since
\(Z=\mathcal X\mathcal I\), polar substitution gives
\[
 Z^*Z
 =R_{\rm inc}V_{\rm inc}^*\mathbb GV_{\rm inc}R_{\rm inc},
\]
which proves \eqref{eq:V28-command-gain-formula}.  On the initial support of
\(V_{\rm inc}\),
\(r_-^2I\preceq R_{\rm inc}^2\preceq r_+^2I\).  Taking the trace against the
positive matrix \(V_{\rm inc}^*\mathbb GV_{\rm inc}\) proves the sandwich.

If \(\tau_\perp=0\), then
\[
 0=\Tr(Q\mathbb GQ)
  =\norm{\mathbb G^{1/2}Q}_{\HS}^{2}.
\]
Hence \(\mathbb GQ=Q\mathbb G=0\), and
\(\mathbb G=P_{\rm inc}\mathbb GP_{\rm inc}\).  The converse is immediate.
\end{proof}

\begin{corollary}[Metric-isometric command incidence]
\label{cor:V28-isometric-incidence}
If \(\mathcal I\) is a partial isometry on the supported command quotient,
then \(R_{\rm inc}\) is the initial support projection and
\begin{equation}
 \boxed{
 \tau_{\rm role}
 =\tau_{\rm cmd}+\tau_{\perp}.}
 \label{eq:V28-isometric-Pythagoras}
\end{equation}
Thus \(\tau_{\rm role}\) is an exact upper physical source card in this branch,
with overestimate equal to the command-invisible role mass.
\end{corollary}

\begin{proof}
For a partial isometry the nonzero incidence gains equal one, so
\(\tau_{\rm cmd}=\tau_{\rm vis}\) in
\eqref{eq:V28-incidence-gain-sandwich}.  Substitute into
\eqref{eq:V28-role-envelope-split}.
\end{proof}

\begin{countertheorem}[The same role total permits physical reinforcement or cancellation]
\label{cth:V28-role-total-no-physical-total}
Let \(\mathcal E_{\rm role}=\R^2\), \(\mathcal U=\R\),
\begin{equation}
 J=\binom11,
 \qquad
 \mathsf M_{\rm cmd}=2,
 \label{eq:V28-aligned-incidence}
\end{equation}
and compare
\begin{equation}
 \mathcal X_+=\begin{pmatrix}1&1\end{pmatrix},
 \qquad
 \mathcal X_-=\begin{pmatrix}1&-1\end{pmatrix}.
 \label{eq:V28-plus-minus-syntheses}
\end{equation}
Both role Grams have trace two and both diagonal arm squares equal one.  Yet
\begin{equation}
 \boxed{
 \tau_{\rm cmd}^{+}=2,
 \qquad
 \tau_{\rm cmd}^{-}=0.}
 \label{eq:V28-plus-minus-command-totals}
\end{equation}
For the plus card \(\tau_\perp=0\); for the minus card
\(\tau_\perp=2\).  Thus the role-total gain alone cannot decide whether the
one physical command is reinforced or exactly cancelled by the role assembly.
\end{countertheorem}

\begin{proof}
The whitened incidence is \(J/\sqrt2\), an isometry onto the span of
\((1,1)^{\mathsf T}\).  The physical responses are
\(\mathcal X_+J=2\) and \(\mathcal X_-J=0\).  Division by the command cost two
gives the displayed totals.  The orthogonal role masses then follow from
\eqref{eq:V28-isometric-Pythagoras}.
\end{proof}

\begin{countertheorem}[A tight role envelope can have divergent command gain]
\label{cth:V28-incidence-gain-escape}
Take one role source with \(\mathbb G_n=1\), one command with cost one, and
incidence \(J_n=n\).  Then
\begin{equation}
 \boxed{
 \tau_{{\rm role},n}=1,
 \qquad
 \tau_{{\rm cmd},n}=n^2.}
 \label{eq:V28-gain-escape}
\end{equation}
Thus source-envelope tightness does not imply tightness of the physically
commanded response unless the incidence gain is controlled.
\end{countertheorem}

\begin{proof}
The physical synthesis is \(Y_n=n\), so its unit-cost action is \(n^2\).
\end{proof}

% -----------------------------------------------------------------------------
\subsection{A finite metric-native command acquisition}
\label{subsec:V28-command-acquisition}
% -----------------------------------------------------------------------------

\begin{theorem}[Unit-action command frame and positive-square reconstruction]
\label{thm:V28-unit-command-frame}
Assume the null-cost branch and let
\(r=\rank\mathsf M_{\rm cmd}\).  Choose a command frame
\(u_1,\ldots,u_r\) satisfying
\begin{equation}
 \ip{u_j}{\mathsf M_{\rm cmd}u_k}=\delta_{jk},
 \qquad
 \Span\{u_j\}=\supp(\mathsf M_{\rm cmd})\mathcal U.
 \label{eq:V28-unit-command-frame}
\end{equation}
Then
\begin{equation}
 \boxed{
 \sum_{j=1}^{r}u_ju_j^*=\mathsf M_{\rm cmd}^\dagger,}
 \label{eq:V28-frame-resolution}
\end{equation}
and the physical command Gram in this frame is
\begin{equation}
 Q_{jk}=\ip{Yu_j}{Yu_k}.
 \label{eq:V28-frame-Gram}
\end{equation}
In particular,
\begin{equation}
 \boxed{
 \tau_{\rm cmd}=\sum_{j=1}^{r}\norm{Yu_j}^{2},}
 \label{eq:V28-unit-frame-total}
\end{equation}
and the eigenvalues of \((Q_{jk})\) are the physical reserve spectrum.

Every matrix entry is reconstructed from positive command-square writers.  Put
\begin{equation}
 q_j=\norm{Yu_j}^{2},
 \qquad
 q_{jk}^{(+)}=\norm{Y(u_j+u_k)}^{2},
 \qquad
 q_{jk}^{(i)}=\norm{Y(u_j+i u_k)}^{2}.
 \label{eq:V28-command-square-writers}
\end{equation}
For \(j<k\),
\begin{equation}
 \boxed{
 Q_{jk}
 =\frac{q_{jk}^{(+)}-q_j-q_k}{2}
 +\frac{i}{2}\bigl(q_j+q_k-q_{jk}^{(i)}\bigr).}
 \label{eq:V28-command-square-polarization}
\end{equation}
One command \(u_*\), frozen before reconstruction, gives the held-out
incidence residual
\begin{equation}
 \boxed{
 \Delta_*^{\rm cmd}
 =\left|
 q_*^{\rm dir}-u_*^*Q_{\rm cmd}u_*
 \right|^2.}
 \label{eq:V28-command-heldout}
\end{equation}
\end{theorem}

\begin{proof}
The map
\(\mathsf M_{\rm cmd}^{1/2}U\), where \(Ue_j=u_j\), is unitary from
\(\C^r\) onto \(\supp(\mathsf M_{\rm cmd})\mathcal U\).  Hence
\[
 UU^*
 =\mathsf M_{\rm cmd}^{\dagger/2}
   P_{\mathsf M}
   \mathsf M_{\rm cmd}^{\dagger/2}
 =\mathsf M_{\rm cmd}^\dagger,
\]
which proves \eqref{eq:V28-frame-resolution}.  Therefore
\[
 \sum_j\norm{Yu_j}^{2}
 =\Tr(Y^*YUU^*)
 =\Tr(Y^*Y\mathsf M_{\rm cmd}^\dagger)
 =\Tr Q_{\rm cmd}.
\]
The frame Gram is unitarily equivalent to
\eqref{eq:V28-command-Gram} on its support.  Expanding the two positive squares
in \eqref{eq:V28-command-square-writers} gives the real and imaginary parts in
\eqref{eq:V28-command-square-polarization}.  The held-out formula is the
definition of same-event incidence consistency.
\end{proof}

\begin{remark}[Target-minimal alternative to the role envelope]
\label{rem:V28-target-minimal-command-card}
The role Gram remains necessary when the target is the full
primitive--crossing--line ancestry decomposition.  If the only intended target
is the physical command reserve, the unit-action frame of
\Cref{thm:V28-unit-command-frame} is smaller: it reads the already assembled
physical command synthesis directly and does not reconstruct unused role
combinations.
\end{remark}

% -----------------------------------------------------------------------------
\subsection{Metric uncertainty and the difference between total mass and reserve}
\label{subsec:V28-metric-reserve}
% -----------------------------------------------------------------------------

\begin{theorem}[Metric-comparison interval]
\label{thm:V28-metric-comparison}
Let \(\mathsf M_0,\mathsf M_1\succeq0\) have one common support, let the
null-cost condition hold for one fixed physical synthesis \(Y\), and assume
\begin{equation}
 c_-\mathsf M_0
 \preceq\mathsf M_1
 \preceq c_+\mathsf M_0,
 \qquad
 0<c_-\le c_+<\infty.
 \label{eq:V28-metric-comparison-hypothesis}
\end{equation}
Then every ordered generalized reserve eigenvalue satisfies
\begin{equation}
 \boxed{
 c_+^{-1}\gamma_k^{(0)}
 \le\gamma_k^{(1)}
 \le c_-^{-1}\gamma_k^{(0)},}
 \label{eq:V28-reserve-comparison}
\end{equation}
and the total physical actions satisfy
\begin{equation}
 \boxed{
 c_+^{-1}\tau_{\rm cmd}^{(0)}
 \le\tau_{\rm cmd}^{(1)}
 \le c_-^{-1}\tau_{\rm cmd}^{(0)}.}
 \label{eq:V28-total-comparison}
\end{equation}
Thus uncertainty in the source-action metric must be propagated before an
absolute gain or reserve interval is reported.
\end{theorem}

\begin{proof}
On the common support the metrics are positive definite.  For every nonzero
command \(u\),
\[
 \frac{u^*Y^*Yu}{c_+u^*\mathsf M_0u}
 \le
 \frac{u^*Y^*Yu}{u^*\mathsf M_1u}
 \le
 \frac{u^*Y^*Yu}{c_-u^*\mathsf M_0u}.
\]
The generalized Courant--Fischer principle gives
\eqref{eq:V28-reserve-comparison}.  Equivalently,
\(c_+^{-1}\mathsf M_0^{-1}\preceq\mathsf M_1^{-1}
\preceq c_-^{-1}\mathsf M_0^{-1}\).  Taking the trace against
\(Y^*Y\succeq0\) proves \eqref{eq:V28-total-comparison}.
\end{proof}

\begin{theorem}[Trace, rank, and reserve conditioning]
\label{thm:V28-trace-reserve-conditioning}
Let \(Q_{\rm cmd}\succeq0\) have rank \(r>0\), total
\(\tau_{\rm cmd}\), lower reserve \(g_-\), upper reserve \(g_+\), and
condition number \(\chi=g_+/g_-\).  Then
\begin{equation}
 \boxed{
 g_-\le\frac{\tau_{\rm cmd}}r\le g_+,}
 \label{eq:V28-trace-average-reserve}
\end{equation}
and
\begin{equation}
 \boxed{
 g_-\ge\frac{\tau_{\rm cmd}}{r\chi}.}
 \label{eq:V28-conditioned-reserve-lower}
\end{equation}
Consequently a source-total lower bound becomes a reserve lower bound only
with a rank and conditioning control, or with a direct protected-screen
minimum.
\end{theorem}

\begin{proof}
The trace is the sum of the \(r\) positive eigenvalues, each lying in
\([g_-,g_+]\), which proves \eqref{eq:V28-trace-average-reserve}.  Since
\(g_+=\chi g_-\),
\(\tau_{\rm cmd}\le rg_+=r\chi g_-\), proving the second inequality.
\end{proof}

\begin{countertheorem}[A constant physical total can have collapsing reserve]
\label{cth:V28-total-no-reserve}
For \(n\ge2\), let
\begin{equation}
 Q_n=\diag\left(1-\frac1n,\frac1n\right).
 \label{eq:V28-reserve-collapse-control}
\end{equation}
Then
\begin{equation}
 \boxed{
 \Tr Q_n=1,
 \qquad
 \rank Q_n=2,
 \qquad
 \lambda_{\min}(Q_n)=\frac1n\longrightarrow0.}
 \label{eq:V28-reserve-collapse-values}
\end{equation}
Thus even the metric-native source total is a tail and gain quantity, not a
substitute for the supported reserve floor required by the Duhamel compiler.
\end{countertheorem}

\begin{proof}
The claims are read directly from the diagonal entries.
\end{proof}

% -----------------------------------------------------------------------------
\subsection{Response-tight tails require incidence-gain control}
\label{subsec:V28-response-tail}
% -----------------------------------------------------------------------------

Return to a measurable history family.  At each history let
\begin{equation}
 t_{\rm role}=\Tr\mathbb G,
 \qquad
 \kappa_{\rm inc}
 =\norm{J\mathsf M_{\rm cmd}^{\dagger/2}}_{\mathrm{op}},
 \label{eq:V28-pointwise-role-gain}
\end{equation}
and let
\begin{equation}
 t_{\rm cmd}
 =\Tr\!
 \left(
 \mathsf M_{\rm cmd}^{\dagger/2}
 J^*\mathbb GJ
 \mathsf M_{\rm cmd}^{\dagger/2}
 \right).
 \label{eq:V28-pointwise-command-total}
\end{equation}
Define \(d\nu_{\rm role}=t_{\rm role}\,d\mu\) and
\(d\nu_{\rm cmd}=t_{\rm cmd}\,d\mu\).

\begin{theorem}[Envelope-to-command tail transfer]
\label{thm:V28-envelope-command-tail}
Pointwise,
\begin{equation}
 \boxed{
 t_{\rm cmd}\le\kappa_{\rm inc}^{2}t_{\rm role}.}
 \label{eq:V28-pointwise-tail-bound}
\end{equation}
Consequently, for every measurable omitted layer \(E\),
\begin{equation}
 \boxed{
 \nu_{\rm cmd}(E)
 \le\int_E\kappa_{\rm inc}^{2}\,d\nu_{\rm role}.}
 \label{eq:V28-command-tail-bound}
\end{equation}
In particular:
\begin{enumerate}[label=\textup{(T\arabic*)},leftmargin=3.5em]
\item if \(\kappa_{\rm inc}\le\kappa_*\), then
      \begin{equation}
       \nu_{\rm cmd}(E)
       \le\kappa_*^2\nu_{\rm role}(E);
       \label{eq:V28-uniform-incidence-tail}
      \end{equation}
\item if \(\kappa_{\rm inc}^2\in L^p(\nu_{\rm role})\), \(p>1\), then
      \begin{equation}
       \nu_{\rm cmd}(E)
       \le
       \norm{\kappa_{\rm inc}^2}_{L^p(\nu_{\rm role})}
       \nu_{\rm role}(E)^{1-1/p}.
       \label{eq:V28-Lp-incidence-tail}
      \end{equation}
\end{enumerate}
Thus the Version~V27 role-envelope uniform-integrability row promotes to the
physical command response only after a bounded or uniformly integrable
incidence-gain row has been supplied.
\end{theorem}

\begin{proof}
Put \(R=J\mathsf M_{\rm cmd}^{\dagger/2}\).  Then
\[
 t_{\rm cmd}=\norm{\mathcal XR}_{\HS}^{2}
 \le\norm R_{\mathrm{op}}^{2}\norm{\mathcal X}_{\HS}^{2}
 =\kappa_{\rm inc}^{2}t_{\rm role},
\]
which proves the first two claims.  The remaining estimates are respectively
the essential-supremum bound and Holder's inequality with respect to
\(\nu_{\rm role}\).
\end{proof}

\begin{countertheorem}[Summable envelope probability and source mass can coexist with response escape]
\label{cth:V28-envelope-no-command-tail}
Let omitted layers \(E_n\) satisfy
\begin{equation}
 d_n:=\mu(E_n)=\frac7{1024}2^{-n},
 \qquad
 \sum_{n\ge0}d_n=\frac7{512}.
 \label{eq:V28-tail-probabilities}
\end{equation}
Take \(t_{{\rm role},n}=1\) on \(E_n\) and
\(\kappa_{{\rm inc},n}^{2}=d_n^{-1}\).  Then
\begin{equation}
 \boxed{
 \nu_{{\rm role},n}(E_n)=d_n,
 \qquad
 \nu_{{\rm cmd},n}(E_n)=1.}
 \label{eq:V28-tail-gain-escape}
\end{equation}
Hence both the probability tail and role-envelope tail are summable, while the
physical command-response tail is not.
\end{countertheorem}

\begin{proof}
The role mass equals \(d_n\).  Equality holds in
\eqref{eq:V28-pointwise-tail-bound} for a scalar synthesis aligned with the
incidence, so the command mass is \(d_nd_n^{-1}=1\).
\end{proof}

% -----------------------------------------------------------------------------
\subsection{Direct cofinal transport on the physical command carrier}
\label{subsec:V28-command-transport}
% -----------------------------------------------------------------------------

For each cutoff let
\begin{equation}
 Z_n
 =\mathcal X_nJ_n\mathsf M_{{\rm cmd},n}^{\dagger/2}
 :\mathcal U_n^{\rm s}\to\mathcal H_n
 \label{eq:V28-whitened-command-synthesis-n}
\end{equation}
be the physically whitened synthesis on its supported command quotient, and put
\begin{equation}
 Q_n=Z_n^*Z_n,
 \qquad
 \tau_n=\norm{Z_n}_{\HS}^{2}.
 \label{eq:V28-command-Gram-n}
\end{equation}
Let
\begin{equation}
 R_n:\mathcal U_n^{\rm s}\to\mathcal U_{n+1}^{\rm s},
 \qquad
 V_n:\mathcal H_n\to\mathcal H_{n+1}
 \label{eq:V28-command-output-isometries}
\end{equation}
be the physically declared command and output isometries.  Define
\begin{equation}
 \delta_n^{\rm cmd}
 :=\norm{Z_{n+1}R_n-V_nZ_n}_{\HS},
 \label{eq:V28-command-transport-defect}
\end{equation}
and the new-command mass
\begin{equation}
 \beta_n^{\rm cmd}
 :=\norm{Z_{n+1}(I-R_nR_n^*)}_{\HS}^{2}.
 \label{eq:V28-new-command-mass}
\end{equation}

\begin{theorem}[Direct physical command transport]
\label{thm:V28-direct-command-transport}
For every adjacent pair,
\begin{equation}
 \boxed{
 \norm{R_n^*Q_{n+1}R_n-Q_n}_{\mathrm{op}}
 \le
 \left(
orm{Z_{n+1}R_n}_{\HS}+\norm{Z_n}_{\HS}\right)
 \delta_n^{\rm cmd}.}
 \label{eq:V28-command-Gram-transport}
\end{equation}
The absolute command totals satisfy
\begin{equation}
 \boxed{
 \begin{aligned}
 |\tau_{n+1}-\tau_n|
 \le{}&
 \left(
orm{Z_{n+1}R_n}_{\HS}+\norm{Z_n}_{\HS}\right)
 \delta_n^{\rm cmd}\
 &+\beta_n^{\rm cmd}.
 \end{aligned}}
 \label{eq:V28-command-total-transport}
\end{equation}
Consequently, if the Hilbert--Schmidt norms are uniformly bounded and
\begin{equation}
 \sum_n\delta_n^{\rm cmd}<\infty,
 \qquad
 \sum_n\beta_n^{\rm cmd}<\infty,
 \label{eq:V28-command-transport-summability}
\end{equation}
then every fixed transported command screen has a unique Gram limit and the
absolute totals are Cauchy.  A reserve floor on a fixed screen survives when
the remaining tail in \eqref{eq:V28-command-Gram-transport} is smaller than
that floor.  The masses \(\beta_n^{\rm cmd}\) are new physical command
supports and are not hidden inside an old-screen metric perturbation.
\end{theorem}

\begin{proof}
Put \(A_n=Z_{n+1}R_n\) and \(B_n=V_nZ_n\).  Since \(V_n\) is an isometry,
\(B_n^*B_n=Q_n\), and
\[
 A_n^*A_n-B_n^*B_n
 =A_n^*(A_n-B_n)+(A_n-B_n)^*B_n.
\]
Taking operator norms and bounding the operator norms by the Hilbert--Schmidt
norms proves \eqref{eq:V28-command-Gram-transport}.

The command domain at level \(n+1\) is the orthogonal sum of
\(\Ran R_n\) and its complement.  Therefore
\[
 \tau_{n+1}
 =\norm{Z_{n+1}R_n}_{\HS}^{2}
  +\beta_n^{\rm cmd}.
\]
The scalar identity
\[
 \left|\norm A_{\HS}^{2}-\norm B_{\HS}^{2}\right|
 \le(\norm A_{\HS}+\norm B_{\HS})\norm{A-B}_{\HS}
\]
gives \eqref{eq:V28-command-total-transport}.  Summation proves the Cauchy
claims, while Weyl's inequality transports a strict screen reserve.  The last
statement follows because \(I-R_nR_n^*\) is the canonical new-command
projection.
\end{proof}

\begin{remark}[Direct transport versus separate role, incidence, and metric transport]
\label{rem:V28-direct-versus-factor-transport}
Equation~\eqref{eq:V28-command-transport-defect} transports the physically
whitened synthesis itself.  It is therefore invariant under simultaneous
role-coordinate, command-coordinate, and metric changes of
\Cref{thm:V28-two-sided-covariance}.  Separate transport of
\(\mathbb G_n\), \(J_n\), and \(\mathsf M_{{\rm cmd},n}\) is a sufficient
factorized implementation, but the direct defect is the target-native row.
\end{remark}

% -----------------------------------------------------------------------------
\subsection{Exact replay and terminal classification}
\label{subsec:V28-replay}
% -----------------------------------------------------------------------------

For the metric-native unit-action control
\begin{equation}
 \mathsf M_{\rm cmd}=\diag(4,9),
 \qquad
 Y=\diag(2,6),
 \label{eq:V28-replay-unit-frame}
\end{equation}
the frame \(u_1=(1/2,0)^{\mathsf T}\),
\(u_2=(0,1/3)^{\mathsf T}\) gives
\begin{equation}
 \boxed{
 Q_{\rm cmd}=\diag(1,4),
 \qquad
 \tau_{\rm cmd}=5,
 \qquad
 (g_-^{\rm cmd},g_+^{\rm cmd})=(1,4).}
 \label{eq:V28-replay-command-Gram}
\end{equation}
Changing the metric to
\(\diag(8,9/2)\) while retaining the same physical response gives
\begin{equation}
 \boxed{\tau_{\rm cmd}=\frac{17}{2},}
 \label{eq:V28-replay-metric-change}
\end{equation}
inside the comparison interval \([5/2,10]\).

The complex command Gram
\begin{equation}
 Q=
 \begin{pmatrix}
 2&1+i\\
 1-i&3
 \end{pmatrix}
 \label{eq:V28-replay-complex-Gram}
\end{equation}
has determinant four.  Its positive square values are
\begin{equation}
 q_1=2,
 \qquad
 q_2=3,
 \qquad
 q_{12}^{(+)}=7,
 \qquad
 q_{12}^{(i)}=3,
 \label{eq:V28-replay-positive-squares}
\end{equation}
and \eqref{eq:V28-command-square-polarization} reconstructs
\(Q_{12}=1+i\).  The held-out command \(u_*=(1,2)^{\mathsf T}\) has exact
square
\begin{equation}
 \boxed{u_*^*Qu_*=18.}
 \label{eq:V28-replay-heldout}
\end{equation}

For the direct transport control, let the old whitened response be \(Z_0=2\),
let the transported old response be \(Z_1R=5/2\), and let the new command
carry squared action \(1/9\).  Then
\begin{equation}
 \boxed{
 \left|\left(\frac52\right)^2-2^2\right|
 =\left(\frac52+2\right)\left|\frac52-2\right|
 =\frac94,}
 \label{eq:V28-replay-old-transport}
\end{equation}
and
\begin{equation}
 \boxed{
 |\tau_1-\tau_0|
 =\frac94+\frac19
 =\frac{85}{36}.}
 \label{eq:V28-replay-total-transport}
\end{equation}
Both inequalities of \Cref{thm:V28-direct-command-transport} are attained.

\begin{theorem}[Version V28 physical command-incidence theorem]
\label{thm:V28-terminal}
Versions~V1--V27 are retained, with the Version~V27 source-total trace read as
a role-envelope domination mass until the physical command-incidence packet
below has passed.  Version~V28 additionally proves:
\begin{enumerate}[label=\textup{(V28.\arabic*)},leftmargin=5.8em]
\item the actual physical source is \(Y=\mathcal XJ\), and its source-action
      metric must pass the null-cost test before pseudoinverse whitening;
\item the metric-native command Gram and reserve are invariant under simultaneous
      role and command reparametrization;
\item the role-envelope trace splits into command-visible and command-invisible
      masses, while the physical total carries an independent incidence-gain
      sandwich;
\item one unit-action command frame reconstructs the physical reserve through
      positive square writers and a held-out command;
\item source-action metric uncertainty has an exact generalized-eigenvalue and
      total-action interval;
\item a total source mass does not imply a reserve floor without rank and
      conditioning control;
\item role-envelope tail control promotes to physical command-response control
      only through a bounded or uniformly integrable incidence gain;
\item the physically whitened synthesis has a direct adjacent-cutoff transport
      theorem which separates old-screen motion from genuinely new command
      support.
\end{enumerate}
\end{theorem}

\begin{proof}
Items~\textup{(V28.1)--(V28.2)} are
Theorems~\ref{thm:V28-command-reserve} and
\ref{thm:V28-two-sided-covariance}.  Item~\textup{(V28.3)} is
Theorem~\ref{thm:V28-role-envelope-Pythagoras}.  Item~\textup{(V28.4)} is
Theorem~\ref{thm:V28-unit-command-frame}.  Items~\textup{(V28.5)--(V28.6)}
are Theorems~\ref{thm:V28-metric-comparison} and
\ref{thm:V28-trace-reserve-conditioning}, together with
Countertheorem~\ref{cth:V28-total-no-reserve}.  Items~\textup{(V28.7)--(V28.8)}
are Theorems~\ref{thm:V28-envelope-command-tail} and
\ref{thm:V28-direct-command-transport}.
\end{proof}

\begin{assessment}[Version V28 verdict]
\label{ass:V28-verdict}
The current result is
\begin{equation}
 \boxed{
 \begin{gathered}
 \passstatus\\[-0.15em]
 \text{physical command-incidence calculus, metric-native reserve,}\\
 \text{role-envelope Pythagoras, gain-aware tails, unit-command}\\
 \text{tomography, and direct command transport}
 \\[0.45em]
 \obsstatus\\[-0.15em]
 \text{pseudowhitening before the null-cost test; treating the role trace}\\
 \text{as physical source action without command incidence; or using}\\
 \text{source-total mass as a reserve floor without conditioning}
 \\[0.45em]
 \stopstatus\\[-0.15em]
 \text{one selected RPESM physical command-to-role incidence map,}\\
 \text{its source-action metric and null-cost residual, unit-command}\\
 \text{squares, incidence-gain tail, line sewing, and one adjacent card.}
 \end{gathered}}
 \label{eq:V28-verdict}
\end{equation}
The first missing quantity is no longer the role-space source-total writer by
itself.  It is the same-history map which says how one physically calibrated
command populates the primitive, crossing, and line arms.
\end{assessment}

\begin{acquisition}[Version V29: populated physical command-incidence and reserve card]
\label{acq:V28-next}
Preserve Versions~V1--V28.  On one physically readable RPESM boundary record,
one selected amplitude semantics, and one finite command screen, populate
\begin{equation}
 \boxed{
 \left(
 \mathcal U,\mathsf M_{\rm cmd},J,
 \mathbb G,
 \Delta_0^{\rm cmd},
 Q_{\rm cmd},
 g_-^{\rm cmd},g_+^{\rm cmd},
 \tau_{\rm cmd}
 \right).}
 \label{eq:V28-next-packet}
\end{equation}
The source-action metric must come from the already declared active Hellinger,
direct-action, or equivalent calibrated command experiment.  The map \(J\)
must be the actual same-history incidence of those commands with the saturated
primitive, phase-quenched crossing, and line-weighted crossing arms; it may not
be chosen from the observed Gram.

Acquire a physical unit-action frame and the positive squares of
\Cref{thm:V28-unit-command-frame}, including one held-out command.  Retain the
Version~V27 role-total writer as a bounded acquisition envelope, but transport
its tail to the physical command measure only after evaluating either a
uniform incidence-gain bound or the direct source-weighted tail of
\Cref{thm:V28-envelope-command-tail}.  At one adjacent cutoff acquire the
direct whitened-synthesis defect \eqref{eq:V28-command-transport-defect} and
the new-command mass \eqref{eq:V28-new-command-mass}.

A positive \(\tau_{\rm cmd}\) is not promoted without a protected reserve
margin.  A positive role-complement mass is retained as a command-incidence
obstruction or as an independently occurring command family; it is not
silently appended to the physical source atlas.  Only after this packet passes
may the primitive-complete line--crossing short be interpreted in physical
source units and passed to the semigroup-cyclic Duhamel ancestry layer.

Do not open the conditional exterior, locality, current/stress, charge, or
Einstein gates before this command-incidence card lands.
\end{acquisition}

\section*{V28 exact replay and append-only record}
\addcontentsline{toc}{section}{V28 exact replay and append-only record}
The accompanying exact verifier checks the metric-covariant one-dimensional
control, the null-cost pseudoinverse obstruction, the aligned/cancelled
role-envelope pair, incidence-gain escape, the unit-action command frame,
metric-comparison bounds, constant-total reserve collapse, complex positive-
square reconstruction, the held-out command, direct old-screen transport, new
command birth, and gain-amplified tail.  These are theorem replays and are not
numerical measurements of the selected RPESM regulator.

The complete Version~V27 source preceding its sole final active
\verb|\end{document}| is preserved byte for byte.  Its preterminal SHA--256
digest is recorded in the validation manifest.  A later continuation must
remove only the final active document-closing command, append new material
immediately before it, restore one terminal command, and increment the filename
to end in \texttt{\_V29.tex}.

% =============================================================================
% END APPENDED VERSION V28 MATERIAL
% =============================================================================

% =============================================================================
% BEGIN APPENDED VERSION V29 MATERIAL
% =============================================================================

\clearpage
\clearpage
\hypersetup{
  pdftitle={Renewal Einstein--Standard--Model Physical Source Metric, Duhamel Ancestry and Quantum Continuum Programme V29},
  pdfsubject={Shrinking neighbouring-cylinder Hellinger secants, validated command-response incidence, and tangent-safe reserve}
}
\pdfinfo{
 /Title (Renewal Einstein--Standard--Model Physical Source Metric, Duhamel Ancestry and Quantum Continuum Programme V29)
 /Subject (Shrinking neighbouring-cylinder Hellinger secants, validated command-response incidence, and tangent-safe reserve)
}

\section{Version V29: shrinking neighbouring-cylinder command geometry before tangent incidence}
\label{sec:V29-neighbouring-cylinder}

\subsection{The tangent assumption remaining in the Version V28 packet}
\label{subsec:V29-context}

Version~V28 correctly inserted the physical command-to-role incidence map and
its source-action metric before role-space whitening.  Its packet was still a
\emph{tangent-level} packet: the linear map \(J\) and the positive form
\(\mathsf M_{\rm cmd}\) were assumed to have been produced by one physically
occurring command experiment.  Version~V20 proves the intrinsic Hellinger
formula once an \(L^2\)-differentiable active cylinder has been supplied, but a
finite list of neighbouring endpoint laws does not itself prove that such a
tangent exists or identify it uniquely.

The first unresolved arrow is therefore
\begin{equation}
 \boxed{
 \begin{gathered}
 \text{one common physical command family of positive cylinders}\
 \Downarrow\\[-0.2em]
 \text{finite Hellinger and direct source-response secants}\
 \Downarrow\\[-0.2em]
 \text{shrinking-family or validated-remainder tangent}\
 \Downarrow\\[-0.2em]
 \text{metric-native command reserve and role incidence}\
 \Downarrow\\[-0.2em]
 \text{Version V28 source-total and cofinal transport packet}.
 \end{gathered}}
 \label{eq:V29-upstream-order}
\end{equation}
A fixed finite endpoint panel is a secant experiment.  It becomes a tangent
experiment only through a physical shrinking family, a direct derivative Read,
or a proof-carrying continuation class with a quantitative remainder.  This
version supplies the finite positive compiler, the exact nonidentifiability
counterexample, and the stable-support passage to the Version~V28 reserve.

% -----------------------------------------------------------------------------
\subsection{One common neighbouring-cylinder and source-response family}
\label{subsec:V29-family}
% -----------------------------------------------------------------------------

Let \(\mathcal U\) be a finite-dimensional real command space with a frozen
Euclidean coordinate frame \(e_1,\ldots,e_r\); all operator norms on the
command coefficients in this version use that declared reference norm.  Let \(\mathcal O\subset\mathcal U\) be a
symmetric neighbourhood of zero.  On one common measurable history carrier
\((\Omega,m)\), suppose the physical command setting \(\theta\in\mathcal O\)
produces a positive finite cylinder
\begin{equation}
 \dd\Lambda_\theta=a_\theta\,\dd m,
 \qquad
 Z_\theta=\int a_\theta\,\dd m\in(0,\infty),
 \qquad
 \dd\mu_\theta=Z_\theta^{-1}\dd\Lambda_\theta,
 \label{eq:V29-positive-family}
\end{equation}
and write its canonical positive square root as
\begin{equation}
 \psi_\theta=\sqrt{a_\theta/Z_\theta}\in L^2(m),
 \qquad
 \norm{\psi_\theta}_2=1.
 \label{eq:V29-normalized-root}
\end{equation}
All roots in this section are compared after the physically declared transport
of their history, determinant/Pfaffian-line, and boundary-record fibres.

Let \(\mathcal K\) be the transported reflected source carrier.  In addition to
the positive law, suppose the same command setting produces a direct physical
source-amplitude Read
\begin{equation}
 \Phi_\theta\in\mathcal K.
 \label{eq:V29-direct-source-amplitude}
\end{equation}
When role attribution is intended, also retain a physically written role
coordinate
\begin{equation}
 \rho_\theta\in\mathcal E_{\rm role}
 \label{eq:V29-role-coordinate}
\end{equation}
and the base role synthesis
\begin{equation}
 \mathcal X_0:\mathcal E_{\rm role}\longrightarrow\mathcal K
 \label{eq:V29-base-role-synthesis}
\end{equation}
from Versions~V25--V28.  The direct amplitude \(\Phi_\theta\) is the
target-native object.  The factorization through \(\rho_\theta\) is an
additional same-history attribution statement.

For \(0<\varepsilon\) with \(\pm\varepsilon e_j\in\mathcal O\), define the
central normalized-root, direct-response, and role secants by
\begin{equation}
 \begin{aligned}
 S_\varepsilon e_j
 &:=\frac{\psi_{\varepsilon e_j}-\psi_{-\varepsilon e_j}}{2\varepsilon},\\
 R_\varepsilon e_j
 &:=\frac{\Phi_{\varepsilon e_j}-\Phi_{-\varepsilon e_j}}{2\varepsilon},\\
 J_\varepsilon e_j
 &:=\frac{\rho_{\varepsilon e_j}-\rho_{-\varepsilon e_j}}{2\varepsilon}.
 \end{aligned}
 \label{eq:V29-central-secants}
\end{equation}
They are extended linearly from the frozen command frame.  The finite positive
command and response Grams are
\begin{equation}
 \boxed{
 \mathsf M_\varepsilon:=4S_\varepsilon^*S_\varepsilon\succeq0,
 \qquad
 \mathsf G_\varepsilon:=R_\varepsilon^*R_\varepsilon\succeq0.}
 \label{eq:V29-secant-Grams}
\end{equation}
The factor four makes \(\mathsf M_\varepsilon\) converge to the ordinary
Hellinger--Fisher metric when the physical tangent exists.

% -----------------------------------------------------------------------------
\subsection{Positive module-distance reconstruction}
\label{subsec:V29-positive-reconstruction}
% -----------------------------------------------------------------------------

For probability laws dominated by \(m\), use the squared Hellinger convention
\begin{equation}
 \mathsf H^2(\mu,\nu)
 :=\norm{\sqrt{\dd\mu/\dd m}-\sqrt{\dd\nu/\dd m}}_2^2.
 \label{eq:V29-Hellinger-convention}
\end{equation}
For the basis command \(e_j\), put
\begin{equation}
 d_j^\psi
 :=\norm{\psi_{\varepsilon e_j}-\psi_{-\varepsilon e_j}}_2^2,
 \qquad
 d_{jk}^\psi
 :=\norm{
 (\psi_{\varepsilon e_j}-\psi_{-\varepsilon e_j})
 -(\psi_{\varepsilon e_k}-\psi_{-\varepsilon e_k})}_2^2.
 \label{eq:V29-law-module-distances}
\end{equation}
Likewise, define the positive direct-response distances
\begin{equation}
 d_j^\Phi
 :=\norm{\Phi_{\varepsilon e_j}-\Phi_{-\varepsilon e_j}}_{\mathcal K}^2,
 \qquad
 d_{jk}^\Phi
 :=\norm{
 (\Phi_{\varepsilon e_j}-\Phi_{-\varepsilon e_j})
 -(\Phi_{\varepsilon e_k}-\Phi_{-\varepsilon e_k})}_{\mathcal K}^2.
 \label{eq:V29-response-module-distances}
\end{equation}
When \(\mathcal K\) is complex, all command Grams in this version use its
underlying real Hilbert product \(\operatorname{Re}\langle\cdot,\cdot\rangle\).

\begin{theorem}[Positive module-distance reconstruction of the finite command card]
\label{thm:V29-module-reconstruction}
The complete finite secant Grams are reconstructed from the nonnegative
quantities in \eqref{eq:V29-law-module-distances}--
\eqref{eq:V29-response-module-distances} by
\begin{equation}
 \boxed{
 (\mathsf M_\varepsilon)_{jj}=\frac{d_j^\psi}{\varepsilon^2},
 \qquad
 (\mathsf M_\varepsilon)_{jk}
 =\frac12
 \left(
 (\mathsf M_\varepsilon)_{jj}
 +(\mathsf M_\varepsilon)_{kk}
 -\frac{d_{jk}^\psi}{\varepsilon^2}
 \right),}
 \label{eq:V29-M-from-distances}
\end{equation}
and
\begin{equation}
 \boxed{
 (\mathsf G_\varepsilon)_{jj}=\frac{d_j^\Phi}{4\varepsilon^2},
 \qquad
 (\mathsf G_\varepsilon)_{jk}
 =\frac12
 \left(
 (\mathsf G_\varepsilon)_{jj}
 +(\mathsf G_\varepsilon)_{kk}
 -\frac{d_{jk}^\Phi}{4\varepsilon^2}
 \right).}
 \label{eq:V29-G-from-distances}
\end{equation}
Equivalently, with Hellinger affinities
\begin{equation}
 \mathcal A_{jk}^{\sigma\tau}
 :=\ip{\psi_{\sigma\varepsilon e_j}}
          {\psi_{\tau\varepsilon e_k}},
 \qquad \sigma,\tau\in\{+1,-1\},
 \label{eq:V29-affinities}
\end{equation}
one has
\begin{equation}
 \boxed{
 (\mathsf M_\varepsilon)_{jk}
 =\frac{1}{\varepsilon^2}
 \left(
 \mathcal A_{jk}^{++}-\mathcal A_{jk}^{+-}
 -\mathcal A_{jk}^{-+}+\mathcal A_{jk}^{--}
 \right).}
 \label{eq:V29-M-from-affinities}
\end{equation}
Thus a finite signed cross coefficient is derived from positive endpoint or
module-distance writers; it is not introduced as an independently oriented
source row.
\end{theorem}

\begin{proof}
Let
\(
 \Delta_j^\psi
 =\psi_{\varepsilon e_j}-\psi_{-\varepsilon e_j}
\).
By \eqref{eq:V29-central-secants},
\(
 \mathsf M_{\varepsilon,jk}
 =\varepsilon^{-2}
   \operatorname{Re}\langle\Delta_j^\psi,\Delta_k^\psi\rangle
\).
The real polarization identity
\[
 2\operatorname{Re}\langle x,y\rangle
 =\norm x^2+\norm y^2-\norm{x-y}^2
\]
gives \eqref{eq:V29-M-from-distances}.  The same argument with
\(
 \Delta_j^\Phi
 =\Phi_{\varepsilon e_j}-\Phi_{-\varepsilon e_j}
\)
and the factor \((2\varepsilon)^{-1}\) gives
\eqref{eq:V29-G-from-distances}.  Expanding
\(
 \langle\Delta_j^\psi,\Delta_k^\psi\rangle
\)
proves \eqref{eq:V29-M-from-affinities}.  Positivity follows independently
from the Gram representations in \eqref{eq:V29-secant-Grams}.
\end{proof}

\begin{remark}[No stochastic interferometer is required]
\label{rem:V29-no-instrument-inflation}
The distances in \eqref{eq:V29-law-module-distances} and
\eqref{eq:V29-response-module-distances} are deterministic positive squares on
one common transported module.  A physical multi-arm instrument is required
only when the instrument itself is part of the target semantics.  It is not
charged merely to recover the finite command and response Grams.
\end{remark}

% -----------------------------------------------------------------------------
\subsection{Finite secant null-cost and reserve}
\label{subsec:V29-secant-reserve}
% -----------------------------------------------------------------------------

Put
\begin{equation}
 P_\varepsilon
 :=\mathsf M_\varepsilon\mathsf M_\varepsilon^\dagger
 \label{eq:V29-secant-support}
\end{equation}
and define the finite null-cost residual
\begin{equation}
 \boxed{
 \Delta_{0,\varepsilon}^{\rm sec}
 :=\Tr\!
 \left[(I-P_\varepsilon)\mathsf G_\varepsilon\right].}
 \label{eq:V29-secant-null-cost}
\end{equation}

\begin{theorem}[Finite neighbouring-cylinder reserve]
\label{thm:V29-finite-reserve}
The following are equivalent:
\begin{equation}
 \boxed{
 \Delta_{0,\varepsilon}^{\rm sec}=0
 \iff
 R_\varepsilon(I-P_\varepsilon)=0
 \iff
 \Ker\mathsf M_\varepsilon\subseteq\Ker R_\varepsilon.}
 \label{eq:V29-finite-null-equivalence}
\end{equation}
On this branch the finite secant reserve is
\begin{equation}
 \boxed{
 Q_\varepsilon
 :=\mathsf M_\varepsilon^{\dagger/2}
   \mathsf G_\varepsilon
   \mathsf M_\varepsilon^{\dagger/2}
 \succeq0.}
 \label{eq:V29-finite-reserve}
\end{equation}
Its nonzero eigenvalues are the generalized response-per-Hellinger-action
values of the actual finite neighbouring-cylinder experiment.  The trace,
minimum supported eigenvalue, and maximum eigenvalue give the finite secant
analogues of \(\tau_{\rm cmd}\), \(g_-^{\rm cmd}\), and \(g_+^{\rm cmd}\).
\end{theorem}

\begin{proof}
Since \(\mathsf G_\varepsilon=R_\varepsilon^*R_\varepsilon\),
\[
 \Tr[(I-P_\varepsilon)\mathsf G_\varepsilon]
 =\norm{R_\varepsilon(I-P_\varepsilon)}_{\rm HS}^2.
\]
The first equivalence follows.  The orthogonal complement of
\(\Ran P_\varepsilon\) is \(\Ker\mathsf M_\varepsilon\), proving the second.
The displayed reserve is a positive congruence.  Its Rayleigh quotient on the
supported command space is exactly
\(
 u^*\mathsf G_\varepsilon u/(u^*\mathsf M_\varepsilon u)
\), which proves the generalized-eigenvalue statement.
\end{proof}

\begin{firewall}[A finite secant reserve is not yet a tangent reserve]
\label{fire:V29-secant-not-tangent}
The matrix \(Q_\varepsilon\) is an exact property of the declared finite
neighbouring-cylinder experiment.  It is not promoted to the Version~V28
tangent reserve from one fixed value of \(\varepsilon\).  The promotion
requires either the validated remainder theorem below or convergence along one
physical shrinking family with stable support.
\end{firewall}

% -----------------------------------------------------------------------------
\subsection{No finite endpoint bank determines a physical tangent}
\label{subsec:V29-no-finite-tangent}
% -----------------------------------------------------------------------------

\begin{countertheorem}[Any fixed finite endpoint bank can conceal a different Fisher tangent]
\label{cth:V29-finite-endpoints-no-tangent}
Fix a finite collection of nonzero radii
\begin{equation}
 0<\varepsilon_1<\cdots<\varepsilon_m\le R.
 \label{eq:V29-fixed-radii}
\end{equation}
There are two smooth strictly positive Bernoulli families
\(\mu_t^{(0)}\) and \(\mu_t^{(1)}\), defined for \(|t|\le R\), such that
\begin{equation}
 \boxed{
 \mu_0^{(0)}=\mu_0^{(1)},
 \qquad
 \mu_{\pm\varepsilon_j}^{(0)}
 =\mu_{\pm\varepsilon_j}^{(1)}
 \quad(1\le j\le m),}
 \label{eq:V29-same-finite-laws}
\end{equation}
but their Hellinger--Fisher informations at zero are different.  One may take
\begin{equation}
 p_t^{(0)}=\frac12,
 \qquad
 p_t^{(1)}=\frac12+c\,g(t),
 \qquad
 g(t)=t\prod_{j=1}^m(t^2-\varepsilon_j^2),
 \label{eq:V29-Bernoulli-families}
\end{equation}
with any sufficiently small nonzero \(c\).  Then
\begin{equation}
 \boxed{
 \mathcal I^{(0)}(0)=0,
 \qquad
 \mathcal I^{(1)}(0)
 =4c^2\prod_{j=1}^m\varepsilon_j^4>0.}
 \label{eq:V29-different-Fisher}
\end{equation}
Consequently no fixed finite collection of sampled endpoint laws or
Hellinger distances identifies the physical tangent without an independent
continuation or shrinking-family premise.
\end{countertheorem}

\begin{proof}
The polynomial \(g\) vanishes at zero and at every
\(\pm\varepsilon_j\), which proves \eqref{eq:V29-same-finite-laws}.  On the
compact interval \([-R,R]\), choose \(c\ne0\) so small that
\(|cg(t)|\le1/4\); then \(p_t^{(1)}\in[1/4,3/4]\), so the family is strictly
positive and smooth.  At zero,
\[
 (p^{(1)})'(0)
 =c(-1)^m\prod_{j=1}^m\varepsilon_j^2\ne0.
\]
The Fisher information of a Bernoulli family is
\(
 (p')^2/[p(1-p)]
\).  Since \(p_0=1/2\), this gives
\eqref{eq:V29-different-Fisher}.  The constant family has zero derivative and
zero information.
\end{proof}

\begin{corollary}[A zero central secant may hide a nonzero source]
\label{cor:V29-zero-secant-hidden-tangent}
For the single-radius case \(m=1\), both endpoint laws at
\(\pm\varepsilon_1\) in \Cref{cth:V29-finite-endpoints-no-tangent} coincide,
so the exact central secant Gram at that radius is zero.  Nevertheless the
physical tangent Fisher information is strictly positive.  Failure to see a
source on one finite symmetric panel is therefore \(\stopstatus\), not an exact
follower verdict.
\end{corollary}

% -----------------------------------------------------------------------------
\subsection{Validated central secants and held-out command linearity}
\label{subsec:V29-validated-secants}
% -----------------------------------------------------------------------------

Suppose now that the command family is physically declared on a neighbourhood,
with the fibre transport frozen independently of the observed response.  Put
\begin{equation}
 S e_j:=D_{e_j}\psi_0,
 \qquad
 Y e_j:=D_{e_j}\Phi_0,
 \qquad
 J e_j:=D_{e_j}\rho_0.
 \label{eq:V29-tangent-columns}
\end{equation}
Then the intended tangent metric and direct response Gram are
\begin{equation}
 \mathsf M:=4S^*S,
 \qquad
 \mathsf G:=Y^*Y.
 \label{eq:V29-tangent-Grams}
\end{equation}

For every basis direction define the third-derivative majorants
\begin{equation}
 \begin{aligned}
 K_j^\psi(\varepsilon)
 &:=\sup_{|t|\le\varepsilon}
 \norm{\frac{\dd^3}{\dd t^3}\psi_{te_j}}_2,\\
 K_j^\Phi(\varepsilon)
 &:=\sup_{|t|\le\varepsilon}
 \norm{\frac{\dd^3}{\dd t^3}\Phi_{te_j}}_{\mathcal K},\\
 K_j^\rho(\varepsilon)
 &:=\sup_{|t|\le\varepsilon}
 \norm{\frac{\dd^3}{\dd t^3}\rho_{te_j}}_{\mathcal E_{\rm role}}.
 \end{aligned}
 \label{eq:V29-third-derivative-majorants}
\end{equation}
Set
\begin{equation}
 \eta_\psi=\frac{\varepsilon^2}{6}
 \left(\sum_j(K_j^\psi)^2\right)^{1/2},
 \quad
 \eta_\Phi=\frac{\varepsilon^2}{6}
 \left(\sum_j(K_j^\Phi)^2\right)^{1/2},
 \quad
 \eta_\rho=\frac{\varepsilon^2}{6}
 \left(\sum_j(K_j^\rho)^2\right)^{1/2}.
 \label{eq:V29-secant-radii}
\end{equation}

\begin{theorem}[Validated tangent and same-history incidence]
\label{thm:V29-validated-tangent}
Assume the three transported families are three times continuously
Fr\'echet differentiable on the displayed command neighbourhood.  Then
\begin{equation}
 \boxed{
 \norm{S_\varepsilon-S}_{\rm op}\le\eta_\psi,
 \qquad
 \norm{R_\varepsilon-Y}_{\rm op}\le\eta_\Phi,
 \qquad
 \norm{J_\varepsilon-J}_{\rm op}\le\eta_\rho.}
 \label{eq:V29-secant-column-bounds}
\end{equation}
Consequently
\begin{equation}
 \boxed{
 \norm{\mathsf M_\varepsilon-\mathsf M}_{\rm op}
 \le4(2\norm{S_\varepsilon}_{\rm op}+\eta_\psi)\eta_\psi,}
 \label{eq:V29-M-error}
\end{equation}
\begin{equation}
 \boxed{
 \norm{\mathsf G_\varepsilon-\mathsf G}_{\rm op}
 \le(2\norm{R_\varepsilon}_{\rm op}+\eta_\Phi)\eta_\Phi.}
 \label{eq:V29-G-error}
\end{equation}
Moreover the tangent role-incidence defect obeys
\begin{equation}
 \boxed{
 \norm{Y-\mathcal X_0J}_{\rm op}
 \le
 \norm{R_\varepsilon-\mathcal X_0J_\varepsilon}_{\rm op}
 +\eta_\Phi+\norm{\mathcal X_0}_{\rm op}\eta_\rho.}
 \label{eq:V29-incidence-error}
\end{equation}
Thus a finite secant card plus its declared third-derivative radii gives an
outward tangent metric, direct response, and role-incidence certificate.
\end{theorem}

\begin{proof}
For a Banach-valued \(C^3\) function \(f\), the central-difference remainder is
\begin{equation}
 \left\|
 \frac{f(\varepsilon)-f(-\varepsilon)}{2\varepsilon}-f'(0)
 \right\|
 \le\frac{\varepsilon^2}{6}
 \sup_{|t|\le\varepsilon}\norm{f'''(t)}.
 \label{eq:V29-central-remainder}
\end{equation}
Apply this to every basis column and use the Hilbert--Schmidt column bound to
obtain \eqref{eq:V29-secant-column-bounds}.  If
\(E=S_\varepsilon-S\), then
\[
 S_\varepsilon^*S_\varepsilon-S^*S
 =S_\varepsilon^*E+E^*S,
\]
while \(\norm S\le\norm{S_\varepsilon}+\norm E\).  This proves
\eqref{eq:V29-M-error}; the response bound is identical without the factor
four.  Finally add and subtract
\(R_\varepsilon\) and \(\mathcal X_0J_\varepsilon\) to obtain
\eqref{eq:V29-incidence-error}.
\end{proof}

For a command
\(u_*=\sum_jc_je_j\) fixed before the basis columns are acquired, define the
direct central secants
\begin{equation}
 s_\varepsilon^{\rm dir}(u_*)
 =\frac{\psi_{\varepsilon u_*}-\psi_{-\varepsilon u_*}}{2\varepsilon},
 \qquad
 r_\varepsilon^{\rm dir}(u_*)
 =\frac{\Phi_{\varepsilon u_*}-\Phi_{-\varepsilon u_*}}{2\varepsilon}.
 \label{eq:V29-heldout-secants}
\end{equation}
The held-out positive residual is
\begin{equation}
 \boxed{
 \Delta_{\varepsilon,u_*}^{\rm lin}
 :=4\norm{s_\varepsilon^{\rm dir}(u_*)-S_\varepsilon u_*}_2^2
  +\norm{r_\varepsilon^{\rm dir}(u_*)-R_\varepsilon u_*}_{\mathcal K}^2.}
 \label{eq:V29-heldout-linearity}
\end{equation}

\begin{corollary}[Validated held-out command]
\label{cor:V29-heldout-command}
Let \(K_{u_*}^\psi(\varepsilon)\) and
\(K_{u_*}^\Phi(\varepsilon)\) be the analogues of
\eqref{eq:V29-third-derivative-majorants} on the line \(tu_*\).  Then
\begin{equation}
 \begin{aligned}
 \norm{s_\varepsilon^{\rm dir}(u_*)-S_\varepsilon u_*}
 &\le\frac{\varepsilon^2}{6}
 \left(
 K_{u_*}^\psi+\sum_j|c_j|K_j^\psi
 \right),\\
 \norm{r_\varepsilon^{\rm dir}(u_*)-R_\varepsilon u_*}
 &\le\frac{\varepsilon^2}{6}
 \left(
 K_{u_*}^\Phi+\sum_j|c_j|K_j^\Phi
 \right).
 \end{aligned}
 \label{eq:V29-heldout-bound}
\end{equation}
A measured residual larger than the corresponding outward bound is a
same-family, command-superposition, or fibre-transport obstruction.  A small
residual validates only the declared finite command screen and continuation
class; it does not identify an unrestricted nonlinear response law.
\end{corollary}

\begin{proof}
The derivative at zero is linear in the command.  Apply
\eqref{eq:V29-central-remainder} once to the held-out line and once to every
basis line, then use the triangle inequality.
\end{proof}

% -----------------------------------------------------------------------------
\subsection{Normalized Hellinger geometry versus absolute cylinder geometry}
\label{subsec:V29-scale-shape}
% -----------------------------------------------------------------------------

Define the unnormalized positive amplitude
\begin{equation}
 \varphi_\theta:=\sqrt{a_\theta}=\sqrt{Z_\theta}\,\psi_\theta.
 \label{eq:V29-unnormalized-root}
\end{equation}

\begin{theorem}[Exact finite scale--shape distance]
\label{thm:V29-scale-shape-distance}
For every two command settings \(\theta,\vartheta\),
\begin{equation}
 \boxed{
 \norm{\varphi_\theta-\varphi_\vartheta}_2^2
 =\left(\sqrt{Z_\theta}-\sqrt{Z_\vartheta}\right)^2
 +\sqrt{Z_\theta Z_\vartheta}\,
  \mathsf H^2(\mu_\theta,\mu_\vartheta).}
 \label{eq:V29-scale-shape-distance}
\end{equation}
If the family is differentiable at zero and
\begin{equation}
 h_u=D_u\log a_0,
 \qquad
 b_u=D_u\log Z_0=\mathbb E_{\mu_0}h_u,
 \label{eq:V29-score-and-mass-slope}
\end{equation}
then
\begin{equation}
 \boxed{
 4\ip{D_u\varphi_0}{D_v\varphi_0}
 =Z_0\left(
 \mathsf M(u,v)+b_ub_v
 \right),}
 \label{eq:V29-absolute-root-metric}
\end{equation}
where
\(\mathsf M(u,v)=\operatorname{Cov}_{\mu_0}(h_u,h_v)\) is the normalized
Hellinger--Fisher metric.
\end{theorem}

\begin{proof}
Expand the left side using
\(\varphi_\theta=\sqrt{Z_\theta}\psi_\theta\) and
\(\norm{\psi_\theta}=\norm{\psi_\vartheta}=1\).  The identity
\(
 \mathsf H^2=2-2\langle\psi_\theta,\psi_\vartheta\rangle
\)
gives \eqref{eq:V29-scale-shape-distance}.  Differentiating the normalized and
unnormalized roots gives
\[
 2D_u\psi_0=(h_u-b_u)\psi_0,
 \qquad
 2D_u\varphi_0=h_u\varphi_0.
\]
Taking inner products and using
\(|\varphi_0|^2\dd m=Z_0\dd\mu_0\) proves
\eqref{eq:V29-absolute-root-metric}.
\end{proof}

\begin{countertheorem}[Normalized laws do not determine absolute source scale]
\label{cth:V29-normalized-no-absolute-scale}
On a one-point history space, let \(a_t=e^{2t}\).  Then the normalized law is
constant and its Hellinger--Fisher metric is zero, while
\begin{equation}
 \boxed{
 4\norm{D_t\sqrt{a_t}|_{t=0}}^2=4.}
 \label{eq:V29-singleton-scale}
\end{equation}
Thus a normalized neighbouring-law panel cannot by itself populate an
unnormalized amplitude action, partition slope, or direct-action Hessian.
\end{countertheorem}

\begin{firewall}[Three source-action semantics]
\label{fire:V29-three-metrics}
The normalized Hellinger metric, the unnormalized square-root-amplitude metric
\eqref{eq:V29-absolute-root-metric}, and an independently declared direct-action
Hessian are distinct physical targets.  Version~V28 may use any one of them
only after the programme states which experiment prices the command and
populates its missing scale or second-action rows.  Agreement of normalized
laws does not identify the other two.
\end{firewall}

% -----------------------------------------------------------------------------
\subsection{Stable-support passage to the Version V28 tangent reserve}
\label{subsec:V29-reserve-limit}
% -----------------------------------------------------------------------------

Let
\begin{equation}
 Q
 :=\mathsf M^{\dagger/2}\mathsf G\mathsf M^{\dagger/2}
 \label{eq:V29-tangent-reserve}
\end{equation}
be the intended Version~V28 direct tangent reserve.  The next theorem gives a
quantitative finite-secants-to-tangent bridge.

\begin{theorem}[Stable-support reserve perturbation]
\label{thm:V29-reserve-perturbation}
Assume the finite and tangent null-cost conditions hold, and that
\(\mathsf M_\varepsilon\) and \(\mathsf M\) have the same canonical support
\(P\) with
\begin{equation}
 \mathsf M_\varepsilon\succeq m_*P,
 \qquad
 \mathsf M\succeq m_*P,
 \qquad
 m_*>0.
 \label{eq:V29-common-floor}
\end{equation}
Put
\begin{equation}
 \eta_M=\norm{\mathsf M_\varepsilon-\mathsf M}_{\rm op},
 \qquad
 \eta_G=\norm{\mathsf G_\varepsilon-\mathsf G}_{\rm op},
 \qquad
 g=\norm{\mathsf G}_{\rm op}.
 \label{eq:V29-Gram-radii}
\end{equation}
Then
\begin{equation}
 \boxed{
 \norm{Q_\varepsilon-Q}_{\rm op}
 \le\frac{\eta_G}{m_*}+\frac{g\eta_M}{m_*^2}.}
 \label{eq:V29-reserve-error}
\end{equation}
For any protected command screen \(P_Q\preceq P\), if
\begin{equation}
 \lambda_{\min}
 \left(P_QQ_\varepsilon P_Q\big|_{\Ran P_Q}\right)
 >\frac{\eta_G}{m_*}+\frac{g\eta_M}{m_*^2},
 \label{eq:V29-reserve-pass}
\end{equation}
then the tangent reserve has the strict floor
\begin{equation}
 \boxed{
 P_QQP_Q
 \succeq
 \left[
 \lambda_{\min}
 \left(P_QQ_\varepsilon P_Q\big|_{\Ran P_Q}\right)
 -\frac{\eta_G}{m_*}-\frac{g\eta_M}{m_*^2}
 \right]P_Q.}
 \label{eq:V29-reserve-screen-floor}
\end{equation}
If the metric supports differ, the event is a command-rank transition and
\eqref{eq:V29-reserve-error} is not applied.
\end{theorem}

\begin{proof}
Restrict to \(\Ran P\).  For positive operators \(A,B\succeq m_*I\), the
resolvent representation
\[
 A^{-1/2}=\frac{1}{\pi}\int_0^\infty
 t^{-1/2}(A+tI)^{-1}\,\dd t
\]
and the resolvent identity give
\begin{equation}
 \norm{A^{-1/2}-B^{-1/2}}_{\rm op}
 \le\frac{1}{2m_*^{3/2}}\norm{A-B}_{\rm op}.
 \label{eq:V29-inverse-root-Lipschitz}
\end{equation}
Write
\(H_\varepsilon=\mathsf M_\varepsilon^{-1/2}\) and
\(H=\mathsf M^{-1/2}\) on this support.  Then
\[
 Q_\varepsilon-Q
 =H_\varepsilon(\mathsf G_\varepsilon-\mathsf G)H_\varepsilon
 +(H_\varepsilon-H)\mathsf G H_\varepsilon
 +H\mathsf G(H_\varepsilon-H).
\]
Use
\(\norm{H_\varepsilon},\norm H\le m_*^{-1/2}\) and
\eqref{eq:V29-inverse-root-Lipschitz}.  The first term is at most
\(\eta_G/m_*\), and the last two sum to at most
\(g\eta_M/m_*^2\).  Compression to \(P_Q\) and Weyl's inequality give
\eqref{eq:V29-reserve-screen-floor}.
\end{proof}

\begin{corollary}[Shrinking-family landing]
\label{cor:V29-shrinking-landing}
Let \(\varepsilon_n\downarrow0\) be one physically occurring shrinking command
family.  Suppose:
\begin{enumerate}[label=\textup{(S\arabic*)},leftmargin=3em]
\item the third-derivative radii of \Cref{thm:V29-validated-tangent} tend to
      zero;
\item the held-out residuals \eqref{eq:V29-heldout-linearity} lie inside their
      declared remainder intervals;
\item the finite null-cost rows \eqref{eq:V29-secant-null-cost} pass;
\item one canonical command support and one positive metric floor persist;
\item the direct role-incidence radii in \eqref{eq:V29-incidence-error} tend to
      zero when role attribution is intended.
\end{enumerate}
Then
\begin{equation}
 \boxed{
 \mathsf M_{\varepsilon_n}\to\mathsf M,
 \qquad
 \mathsf G_{\varepsilon_n}\to\mathsf G,
 \qquad
 Q_{\varepsilon_n}\to Q}
 \label{eq:V29-shrinking-convergence}
\end{equation}
in operator norm.  The tangent metric and reserve are therefore populated by
the neighbouring-cylinder experiment rather than postulated.  When the
incidence row also closes,
\begin{equation}
 \boxed{Y=\mathcal X_0J.}
 \label{eq:V29-tangent-incidence-landing}
\end{equation}
\end{corollary}

\begin{proof}
The first convergence pair is
\Cref{thm:V29-validated-tangent}; the reserve convergence is
\Cref{thm:V29-reserve-perturbation}.  The incidence identity follows by taking
the limit in \eqref{eq:V29-incidence-error}.
\end{proof}

% -----------------------------------------------------------------------------
\subsection{Exact controls and terminal classification}
\label{subsec:V29-controls}
% -----------------------------------------------------------------------------

For the one-radius Bernoulli obstruction, take
\begin{equation}
 p_t=\frac12+\frac18t(t^2-1),
 \qquad |t|\le1.
 \label{eq:V29-replay-Bernoulli}
\end{equation}
Then \(p_{-1}=p_0=p_1=1/2\), so the central endpoint secant is zero, while
\begin{equation}
 \boxed{p'_0=-\frac18,\qquad \mathcal I(0)=\frac{1}{16}.}
 \label{eq:V29-replay-hidden-Fisher}
\end{equation}
The family stays in \([3/8,5/8]\), so no support singularity is involved.

For the exact finite scale--shape identity, use normalized roots
\begin{equation}
 \psi_0=\left(\frac35,\frac45\right),
 \qquad
 \psi_1=\left(\frac45,\frac35\right),
 \qquad
 Z_0=4,
 \qquad
 Z_1=9.
 \label{eq:V29-replay-roots}
\end{equation}
Then
\begin{equation}
 \boxed{
 \mathsf H^2(\mu_0,\mu_1)=\frac{2}{25},
 \qquad
 \norm{\sqrt{Z_1}\psi_1-\sqrt{Z_0}\psi_0}^2=\frac{37}{25},}
 \label{eq:V29-replay-scale-shape}
\end{equation}
and
\begin{equation}
 \frac{37}{25}
 =\left(3-2\right)^2+6\,\frac{2}{25}.
 \label{eq:V29-replay-scale-shape-sum}
\end{equation}

For a rank-two finite secant card, take four-point positive roots
\begin{equation}
 \begin{aligned}
 \psi_{+1}&=(4/5,3/5,0,0),
 &\psi_{-1}&=(3/5,4/5,0,0),\\
 \psi_{+2}&=(0,0,4/5,3/5),
 &\psi_{-2}&=(0,0,3/5,4/5),
 \end{aligned}
 \label{eq:V29-replay-four-roots}
\end{equation}
with \(\varepsilon=1\), and direct response secants
\begin{equation}
 R_\varepsilon=\diag(1/5,2/5).
 \label{eq:V29-replay-response}
\end{equation}
The positive module-distance reconstruction gives
\begin{equation}
 \boxed{
 \mathsf M_\varepsilon=\frac{2}{25}I_2,
 \qquad
 \mathsf G_\varepsilon=\diag(1/25,4/25),}
 \label{eq:V29-replay-Grams}
\end{equation}
and therefore
\begin{equation}
 \boxed{
 Q_\varepsilon=\diag(1/2,2),
 \qquad
 \tau_\varepsilon=\frac52,
 \qquad
 (g_-^\varepsilon,g_+^\varepsilon)=(1/2,2).}
 \label{eq:V29-replay-reserve}
\end{equation}

The central-difference constant is attained by the scalar cubic
\begin{equation}
 f(t)=f_0+tv+t^3w,
 \label{eq:V29-replay-cubic}
\end{equation}
for which
\begin{equation}
 \boxed{
 \frac{f(\varepsilon)-f(-\varepsilon)}{2\varepsilon}-f'(0)
 =\varepsilon^2w
 =\frac{\varepsilon^2}{6}f'''(0).}
 \label{eq:V29-replay-sharp-remainder}
\end{equation}

\begin{theorem}[Version V29 neighbouring-cylinder command theorem]
\label{thm:V29-terminal}
Versions~V1--V28 are retained.  Version~V29 additionally proves:
\begin{enumerate}[label=\textup{(V29.\arabic*)},leftmargin=5.9em]
\item one common family of positive cylinders and direct source amplitudes gives
      finite positive command and response Grams through module distances;
\item the finite secant null-cost test precedes its metric-native reserve;
\item no fixed finite endpoint bank identifies a physical tangent, even for
      smooth strictly positive Bernoulli families;
\item a proof-carrying third-derivative bound gives outward tangent, response,
      role-incidence, and held-out command intervals;
\item normalized Hellinger geometry and absolute cylinder-amplitude geometry
      satisfy an exact scale--shape identity but remain different source-action
      semantics;
\item a common supported floor gives an explicit finite-secants-to-tangent
      reserve radius;
\item one shrinking physical family with stable support lands the Version~V28
      metric and reserve without postulating either.
\end{enumerate}
\end{theorem}

\begin{proof}
Items~\textup{(V29.1)--(V29.2)} are
Theorems~\ref{thm:V29-module-reconstruction} and
\ref{thm:V29-finite-reserve}.  Item~\textup{(V29.3)} is
Countertheorem~\ref{cth:V29-finite-endpoints-no-tangent}.  Item~\textup{(V29.4)}
is Theorem~\ref{thm:V29-validated-tangent} and
Corollary~\ref{cor:V29-heldout-command}.  Item~\textup{(V29.5)} is
Theorem~\ref{thm:V29-scale-shape-distance}.  Items~\textup{(V29.6)--(V29.7)}
are Theorem~\ref{thm:V29-reserve-perturbation} and
Corollary~\ref{cor:V29-shrinking-landing}.
\end{proof}

\begin{assessment}[Version V29 verdict]
\label{ass:V29-verdict}
The current result is
\begin{equation}
 \boxed{
 \begin{gathered}
 \passstatus\\[-0.15em]
 \text{positive neighbouring-cylinder module geometry, finite secant reserve,}\\
 \text{validated tangent and incidence radii, exact scale--shape separation,}\\
 \text{and stable-support passage to the physical command reserve}
 \\[0.45em]
 \obsstatus\\[-0.15em]
 \text{promoting finitely many endpoint laws to a tangent; using a zero finite}\\
 \text{secant as exact followerhood; or identifying normalized Hellinger action}\\
 \text{with an unnormalized amplitude or direct-action metric}
 \\[0.45em]
 \stopstatus\\[-0.15em]
 \text{one shrinking RPESM command-cylinder and direct source-response family,}\\
 \text{with physical fibre transport, remainder or derivative control, stable}\\
 \text{support, line sewing, and one held-out command.}
 \end{gathered}}
 \label{eq:V29-verdict}
\end{equation}
The first missing row is no longer an abstract command metric.  It is the
neighbouring physical deformation experiment from which the metric, source
response, and optional role incidence are jointly differentiated.
\end{assessment}

\begin{acquisition}[Version V30: shrinking active command-cylinder and direct response card]
\label{acq:V29-next}
Preserve Versions~V1--V29.  On one physically readable RPESM boundary record,
one amplitude semantics, and one finite real command screen, acquire a common
family at
\begin{equation}
 \boxed{
 \theta=0,
 \quad
 \theta=\pm\varepsilon_ne_j,
 \quad
 \theta=\pm\varepsilon_nu_*,
 \qquad
 \varepsilon_n\downarrow0,}
 \label{eq:V29-next-settings}
\end{equation}
including:
\begin{enumerate}[label=\textup{(A\arabic*)},leftmargin=3em]
\item the unnormalized positive masses \(Z_\theta\) and normalized roots or
      Hellinger affinities;
\item the direct transported source amplitudes \(\Phi_\theta\);
\item, when role attribution is claimed, the transported role coordinates
      \(\rho_\theta\) and the base synthesis \(\mathcal X_0\);
\item either direct derivative writers or validated third-derivative remainder
      bounds on the entire shrinking panel;
\item the finite null-cost residual, canonical support, supported floor, and
      one held-out command-superposition residual;
\item the determinant/Pfaffian line sewing and every support or index-wall event
      encountered by the deformation.
\end{enumerate}
Reconstruct \(\mathsf M_{\varepsilon_n}\),
\(\mathsf G_{\varepsilon_n}\), and \(Q_{\varepsilon_n}\) before fitting a
tangent.  A support transition is exported to the divisor/index-wall packet.
A bounded but nonconvergent secant family is a tangent obstruction.  On a
stable shrinking branch, use \Cref{thm:V29-reserve-perturbation} to populate the
Version~V28 physical command reserve, and only then interpret the
primitive-complete line--crossing short in those units.

Do not open another phase, crossing, exterior, localizer, Duhamel, locality,
current/stress, charge, or Einstein compiler before this neighbouring-cylinder
card lands.
\end{acquisition}

\section*{V29 exact replay and append-only record}
\addcontentsline{toc}{section}{V29 exact replay and append-only record}
The accompanying exact verifier checks the finite-anchor no-tangent family, the
zero-secant/nonzero-Fisher witness, the finite scale--shape identity, the
singleton absolute-scale obstruction, the positive module-distance
reconstruction, the finite metric-native reserve, the sharp central-difference
constant, a same-history role-incidence residual, and the stable-support reserve
bound.  These are theorem replays and are not measurements of the selected
RPESM regulator.

The complete Version~V28 source preceding its sole final active
\verb|\end{document}| is preserved byte for byte.  Its preterminal SHA--256
digest is recorded in the validation manifest.  A later continuation must
remove only the final active document-closing command, append new material
immediately before it, restore one terminal command, and increment the filename
to end in \texttt{\_V30.tex}.

% =============================================================================
% END APPENDED VERSION V29 MATERIAL
% =============================================================================


% =============================================================================
% BEGIN APPENDED VERSION V30 MATERIAL
% =============================================================================

\hypersetup{
 pdftitle={Renewal Einstein--Standard--Model Physical Source Metric, Duhamel Ancestry and Quantum Continuum Programme V30},
 pdfsubject={Relation-covariant command tangents, moving word supports, and common-action plaquettes}
}

\part*{Version V30 continuation: relation-covariant command tangents and common-action plaquettes}
\addcontentsline{toc}{part}{Version V30 continuation: relation-covariant command tangents and common-action plaquettes}

\section{Upstream correction: quotient transport before differentiation}
\label{sec:V30-upstream-correction}

Versions~V1--V29 are retained verbatim.  Version~V29 reconstructed finite
Hellinger and direct-response secants and proved that a shrinking physical
family, rather than finitely many endpoint laws, is needed to identify a
command tangent.  The acquisition proposed there still presupposed that the
coefficient fibres at the neighbouring commands had already been placed on one
physical relation quotient.  That presupposition is not automatic when the
returned-word support, chiral support, boundary source support, or another
physical relation carrier moves with the command.

The first unresolved arrow is therefore refined to
\begin{equation}
 \boxed{
 \begin{gathered}
 \text{one common unnormalized multiparameter command cylinder}\\
 \Downarrow\\
 \text{physical word relations and their support transport}\\
 \Downarrow\\
 \text{relation-covariant neighbouring-command secants}\\
 \Downarrow\\
 \text{Version~V29 tangent metric and Version~V28 command reserve}.
 \end{gathered}}
 \label{eq:V30-corrected-order}
\end{equation}
A raw ambient derivative may otherwise assign positive action to a coefficient
which is a zero physical word at the base command.  Conversely, separately
populated one-parameter command branches need not be the restrictions of one
common multiparameter action.  These are loading and incidence defects, not new
source births.

\section{Moving relation supports and the Kato normal connection}
\label{sec:V30-moving-support}

Let \(E\) and \(\mathcal K\) be finite-dimensional Hilbert spaces.  Let
\begin{equation}
 W_t:E\longrightarrow\mathcal K,
 \qquad
 H_t=W_t^*W_t,
 \qquad
 P_t=H_tH_t^\dagger
 \label{eq:V30-word-support}
\end{equation}
be a \(C^1\) family on an interval about zero.  Assume that
\(\rank H_t\) is constant there.  Thus \(P_t\) is the orthogonal projection
onto \((\Ker W_t)^\perp\) and is \(C^1\).  Let
\begin{equation}
 X_t:E\longrightarrow\mathcal Y
 \label{eq:V30-response-synthesis}
\end{equation}
be a \(C^1\) response synthesis satisfying the physical relation descent
\begin{equation}
 \boxed{X_t=X_tP_t.}
 \label{eq:V30-endpoint-descent}
\end{equation}

\begin{theorem}[Kato support transport and exact relation-covariant derivative]
\label{thm:V30-Kato-support}
Put
\begin{equation}
 K_t=[\dot P_t,P_t]
 =\dot P_tP_t-P_t\dot P_t.
 \label{eq:V30-Kato-generator}
\end{equation}
Then \(K_t^*=-K_t\).  The initial-value problem
\begin{equation}
 \dot U_t=K_tU_t,
 \qquad
 U_0=I
 \label{eq:V30-Kato-ODE}
\end{equation}
has a unitary solution satisfying
\begin{equation}
 \boxed{P_tU_t=U_tP_0.}
 \label{eq:V30-support-intertwining}
\end{equation}
At \(t=0\), the corresponding covariant derivative of the response is
\begin{equation}
 \boxed{
 \nabla^{\rm K}X
 :=\left.\frac{\dd}{\dd t}\right|_{0}X_tU_t
 =\dot X_0+X_0K_0
 =\dot X_0P_0.}
 \label{eq:V30-covariant-derivative}
\end{equation}
In particular,
\begin{equation}
 \boxed{\nabla^{\rm K}X(I-P_0)=0.}
 \label{eq:V30-covariant-descent}
\end{equation}
Thus the Kato-aligned tangent descends through every base-point word relation.
\end{theorem}

\begin{proof}
Differentiating \(P_t^2=P_t\) gives
\begin{equation}
 P_t\dot P_tP_t=0,
 \qquad
 \dot P_t=P_t\dot P_t(I-P_t)+(I-P_t)\dot P_tP_t.
 \label{eq:V30-projector-derivative}
\end{equation}
Hence \([\dot P_t,P_t]^*=-[\dot P_t,P_t]\).  The solution of
\eqref{eq:V30-Kato-ODE} is therefore unitary.  A direct computation using
\begin{equation}
 [K_t,P_t]=\dot P_t
 \label{eq:V30-Kato-commutator}
\end{equation}
shows that
\begin{equation}
 \frac{\dd}{\dd t}(P_tU_t-U_tP_0)
 =K_t(P_tU_t-U_tP_0).
\end{equation}
The difference vanishes at zero, so it vanishes identically, proving
\eqref{eq:V30-support-intertwining}.

Differentiate \eqref{eq:V30-endpoint-descent} at zero:
\begin{equation}
 \dot X_0=\dot X_0P_0+X_0\dot P_0.
 \label{eq:V30-X-derivative-split-first}
\end{equation}
Since \(X_0=X_0P_0\) and \(P_0\dot P_0P_0=0\), one has
\begin{equation}
 X_0K_0
 =X_0\dot P_0P_0-X_0P_0\dot P_0
 =-X_0\dot P_0.
\end{equation}
Substitution into the derivative of \(X_tU_t\) proves
\eqref{eq:V30-covariant-derivative}.  The last assertion follows because
\(P_0(I-P_0)=0\).
\end{proof}

\begin{theorem}[Raw-tangent Pythagoras and the exact moving-relation debit]
\label{thm:V30-raw-Pythagoras}
Under the hypotheses of \Cref{thm:V30-Kato-support},
\begin{equation}
 \boxed{
 \dot X_0
 =\nabla^{\rm K}X
 +X_0\dot P_0(I-P_0).}
 \label{eq:V30-raw-covariant-split}
\end{equation}
The two summands have orthogonal initial spaces.  Consequently
\begin{equation}
 \boxed{
 \norm{\dot X_0}_{\mathrm{HS}}^{2}
 =\norm{\nabla^{\rm K}X}_{\mathrm{HS}}^{2}
 +\norm{X_0\dot P_0(I-P_0)}_{\mathrm{HS}}^{2}.}
 \label{eq:V30-raw-Pythagoras}
\end{equation}
More precisely, the raw tangent Gram has the block decomposition
\begin{equation}
 \boxed{
 \begin{aligned}
 \dot X_0^*\dot X_0
 ={}&P_0\dot X_0^*\dot X_0P_0\\
 &+(I-P_0)\dot P_0X_0^*X_0\dot P_0(I-P_0)\\
 &+P_0\dot X_0^*X_0\dot P_0(I-P_0)
 +(I-P_0)\dot P_0X_0^*\dot X_0P_0.
 \end{aligned}}
 \label{eq:V30-raw-Gram-blocks}
\end{equation}
The lower-right positive block
\begin{equation}
 \boxed{
 \mathcal L_{\rm rel}
 :=(I-P_0)\dot P_0X_0^*X_0\dot P_0(I-P_0)\succeq0}
 \label{eq:V30-relation-motion-leakage}
\end{equation}
is exactly the action assigned by fixed ambient coordinates to the motion of
base-point zero words.  It is removed by relation-covariant transport and is
not a new physical source.
\end{theorem}

\begin{proof}
Equation~\eqref{eq:V30-X-derivative-split-first} and
\(X_0\dot P_0=X_0\dot P_0(I-P_0)\) give
\eqref{eq:V30-raw-covariant-split}.  The first summand vanishes on
\(\Ran(I-P_0)\), while the second vanishes on \(\Ran P_0\).  Orthogonality of
the domain decomposition gives \eqref{eq:V30-raw-Pythagoras}.  Expanding the
Gram gives \eqref{eq:V30-raw-Gram-blocks}; its lower-right block is the Gram of
\(X_0\dot P_0(I-P_0)\), proving the final statement.
\end{proof}

\begin{corollary}[Physical internal connection remains an independent row]
\label{cor:V30-internal-connection}
Let \(V_t\) be any other unitary support transport satisfying
\begin{equation}
 P_tV_t=V_tP_0,
 \qquad
 V_0=I.
 \label{eq:V30-physical-support-transport}
\end{equation}
Then
\begin{equation}
 G_t:=U_t^*V_t
 \label{eq:V30-internal-gauge}
\end{equation}
is unitary and commutes with \(P_0\).  Conversely every such supported unitary
\(G_t\) produces a support transport \(V_t=U_tG_t\).  The physical derivative is
\begin{equation}
 \boxed{
 \nabla^{\rm phys}X
 =\nabla^{\rm K}X+X_0\dot G_0.}
 \label{eq:V30-physical-versus-Kato}
\end{equation}
Thus the moving relation support fixes the off-diagonal normal connection, but
it does not determine the supported phase, flavour, reflection, or line
connection \(G_t\).
\end{corollary}

\begin{proof}
Both \(U_t\) and \(V_t\) are unitary.  The intertwining equations imply
\(P_0G_t=G_tP_0\).  The converse is immediate.  Differentiating
\(X_tU_tG_t\) at zero gives \eqref{eq:V30-physical-versus-Kato}.
\end{proof}

\section{Exact finite support alignment}
\label{sec:V30-finite-alignment}

\begin{theorem}[Polar chord between nearby relation supports]
\label{thm:V30-polar-chord}
Let \(P,Q\) be orthogonal projections of the same finite rank and assume
\begin{equation}
 \delta:=\norm{P-Q}_{\mathrm{op}}<1.
 \label{eq:V30-angle-condition}
\end{equation}
On \(\Ran P\), the positive operator \(PQP\) obeys
\begin{equation}
 PQP\succeq(1-\delta^2)P.
 \label{eq:V30-PQP-floor}
\end{equation}
Define
\begin{equation}
 \boxed{
 U_{Q\leftarrow P}
 :=QP(PQP)^{\dagger/2}.}
 \label{eq:V30-polar-chord}
\end{equation}
Then
\begin{equation}
 \boxed{
 U_{Q\leftarrow P}^*U_{Q\leftarrow P}=P,
 \qquad
 U_{Q\leftarrow P}U_{Q\leftarrow P}^*=Q.}
 \label{eq:V30-polar-chord-isometry}
\end{equation}
It is the polar partial isometry of \(QP\) and is the unique such isometry with
positive old-support overlap
\begin{equation}
 P U_{Q\leftarrow P}P=(PQP)^{1/2}.
 \label{eq:V30-positive-overlap}
\end{equation}
\end{theorem}

\begin{proof}
For \(x\in\Ran P\),
\begin{equation}
 \norm{Qx}^2
 =\norm{x}^2-\norm{(I-Q)x}^2
 \ge(1-\delta^2)\norm{x}^2,
\end{equation}
which proves \eqref{eq:V30-PQP-floor}.  Hence \(QP\) is injective from
\(\Ran P\) to \(\Ran Q\).  Since the ranks agree, it is bijective.  The polar
formula gives
\begin{equation}
 U^*U
 =(PQP)^{\dagger/2}PQP(PQP)^{\dagger/2}=P.
\end{equation}
Its range is all of \(\Ran Q\), so \(UU^*=Q\).  Finally,
\(PUP=(PQP)^{1/2}\).  Uniqueness is the uniqueness of the polar partial
isometry when the initial and final supports are fixed and the overlap is
positive.
\end{proof}

\begin{theorem}[Exactly relation-descended neighbouring-command secant]
\label{thm:V30-aligned-secant}
Let \(P_0,P_+,P_-\) have the same rank and satisfy
\(\norm{P_\pm-P_0}<1\).  Put
\begin{equation}
 U_\pm=U_{P_\pm\leftarrow P_0}.
 \label{eq:V30-two-chords}
\end{equation}
Suppose \(X_\pm=X_\pm P_\pm\).  Define the chordally aligned central secant
\begin{equation}
 \boxed{
 \mathscr S_\varepsilon^{\rm K}
 =\frac{X_+U_+-X_-U_-}{2\varepsilon}.}
 \label{eq:V30-aligned-central-secant}
\end{equation}
Then
\begin{equation}
 \boxed{
 \mathscr S_\varepsilon^{\rm K}
 =\mathscr S_\varepsilon^{\rm K}P_0.}
 \label{eq:V30-finite-exact-descent}
\end{equation}
Thus its Gram
\begin{equation}
 G_\varepsilon^{\rm K}
 =(\mathscr S_\varepsilon^{\rm K})^*
  \mathscr S_\varepsilon^{\rm K}
 \label{eq:V30-aligned-Gram}
\end{equation}
has zero base-relation residual exactly:
\begin{equation}
 \boxed{
 \Tr\bigl((I-P_0)G_\varepsilon^{\rm K}\bigr)=0.}
 \label{eq:V30-aligned-zero-relation}
\end{equation}
If \(P_t\) and \(X_t\) are \(C^3\), then for
\(F(t)=X_tU_{P_t\leftarrow P_0}\),
\begin{equation}
 \boxed{
 \norm{
 \mathscr S_\varepsilon^{\rm K}-\nabla^{\rm K}X
 }
 \le
 \frac{\varepsilon^2}{6}
 \sup_{|t|\le\varepsilon}\norm{F^{(3)}(t)}.}
 \label{eq:V30-aligned-remainder}
\end{equation}
\end{theorem}

\begin{proof}
The polar chord satisfies \(U_\pm=U_\pm P_0\), so each term in
\eqref{eq:V30-aligned-central-secant} vanishes on \(\Ker P_0\).  This proves
\eqref{eq:V30-finite-exact-descent} and hence
\eqref{eq:V30-aligned-zero-relation}.  The first derivative of the polar chord
at zero is \(\dot P_0P_0\), obtained by differentiating
\eqref{eq:V30-polar-chord}; therefore
\begin{equation}
 F'(0)=\dot X_0P_0+X_0\dot P_0P_0=\dot X_0P_0=\nabla^{\rm K}X.
\end{equation}
The Banach-valued central-difference remainder gives
\eqref{eq:V30-aligned-remainder}.
\end{proof}

\begin{countertheorem}[A raw secant can manufacture action on a zero word]
\label{cth:V30-raw-secant-spurious}
Let
\begin{equation}
 W_t=X_t=\begin{pmatrix}1&t\end{pmatrix}:
 \R^2\longrightarrow\R.
 \label{eq:V30-moving-line-example}
\end{equation}
Then
\begin{equation}
 P_t
 =\frac1{1+t^2}
 \begin{pmatrix}1&t\\t&t^2\end{pmatrix},
 \qquad
 P_0=\begin{pmatrix}1&0\\0&0\end{pmatrix}.
 \label{eq:V30-moving-line-projector}
\end{equation}
The base coefficient \(e_2\) is a zero physical word, but for every
\(\varepsilon>0\),
\begin{equation}
 \boxed{
 \frac{X_\varepsilon-X_{-\varepsilon}}{2\varepsilon}
 =\begin{pmatrix}0&1\end{pmatrix},}
 \label{eq:V30-raw-spurious-secant}
\end{equation}
so the raw secant assigns unit action to \(e_2\).

By contrast, the polar chords satisfy
\begin{equation}
 U_\pm e_1
 =\frac{(1,\pm\varepsilon)^{\mathsf T}}
        {\sqrt{1+\varepsilon^2}},
 \label{eq:V30-example-chords}
\end{equation}
and therefore
\begin{equation}
 X_\pm U_\pm e_1=\sqrt{1+\varepsilon^2}.
\end{equation}
The relation-covariant central secant is exactly zero.  At the tangent level,
\begin{equation}
 \boxed{
 \norm{\dot X_0}_{\mathrm{HS}}^2=1
 =0+1
 =\norm{\nabla^{\rm K}X}^2
 +\norm{X_0\dot P_0(I-P_0)}^2.}
 \label{eq:V30-example-Pythagoras}
\end{equation}
Thus the positive raw null-word action is entirely relation-support motion.
\end{countertheorem}

\begin{proof}
All formulas follow by direct multiplication.  In particular,
\(\dot P_0=\begin{psmallmatrix}0&1\\1&0\end{psmallmatrix}\),
\(K_0=\begin{psmallmatrix}0&-1\\1&0\end{psmallmatrix}\), and
\(\dot X_0+X_0K_0=0\).
\end{proof}

\section{Common-action plaquettes before mixed command Grams}
\label{sec:V30-common-action}

A collection of one-parameter neighbouring cylinders may be individually
positive and smoothly parameterized without being the restriction of one
common multiparameter physical action.  The missing incidence is a finite
command plaquette.

Let \(\Theta_h\subset h\mathbb{Z}^d\) be a finite rectangular command grid.  For every
oriented positive coordinate edge, let
\begin{equation}
 r_i(\theta,\omega)>0
 \label{eq:V30-edge-ratio}
\end{equation}
be the submitted unnormalized likelihood ratio for the command increment
\(\theta\to\theta+he_i\).  Every ratio is required to be a same-history writer
on one common base event space.

\begin{definition}[Command plaquette residual]
\label{def:V30-plaquette}
For every elementary \((i,j)\)-plaquette define
\begin{equation}
 \boxed{
 \begin{aligned}
 \Pi_{ij,h}(\theta)
 :={}&r_i(\theta)r_j(\theta+he_i)\\
 &-r_j(\theta)r_i(\theta+he_j),
 \end{aligned}}
 \label{eq:V30-multiplicative-plaquette}
\end{equation}
and, relative to any full-support finite dominating measure \(m\),
\begin{equation}
 \boxed{
 \Delta_{ij,h}^{\rm act}(\theta)
 =\int|\Pi_{ij,h}(\theta,\omega)|^2\,m(\dd\omega).}
 \label{eq:V30-plaquette-residual}
\end{equation}
Since the ratios are positive, the equivalent logarithmic curvature is
\begin{equation}
 \boxed{
 \begin{aligned}
 \Omega_{ij,h}(\theta)
 :={}&\log r_i(\theta)+\log r_j(\theta+he_i)\\
 &-\log r_j(\theta)-\log r_i(\theta+he_j).
 \end{aligned}}
 \label{eq:V30-log-plaquette}
\end{equation}
\end{definition}

\begin{theorem}[Finite common-action reconstruction]
\label{thm:V30-common-action-grid}
Fix one positive anchor density \(a_0\).  The following are equivalent on a
finite simply connected rectangular command grid:
\begin{enumerate}[label=\textup{(\roman*)},leftmargin=2.8em]
\item there is a family of positive densities \(a_\theta\) satisfying
      \begin{equation}
       a_{\theta+he_i}=r_i(\theta)a_\theta
       \label{eq:V30-edge-density-relation}
      \end{equation}
      on every edge and having the prescribed anchor;
\item every elementary plaquette residual
      \(\Delta_{ij,h}^{\rm act}(\theta)\) vanishes.
\end{enumerate}
On this branch the family is unique and is reconstructed by the product of edge
ratios along any path from the anchor.  Hence all normalized Hellinger panels,
absolute partition rows, direct source-response panels, and mixed command Grams
are restrictions of one common unnormalized deformation cylinder.
\end{theorem}

\begin{proof}
If \eqref{eq:V30-edge-density-relation} holds, multiplying along the two paths
around one elementary square gives \(\Pi_{ij,h}=0\).  Conversely, define
\(a_\theta\) by multiplying the ratios along a lattice path from zero to
\(\theta\).  Any two monotone paths in a rectangle are related by finitely many
interchanges of adjacent coordinate steps.  Vanishing of the corresponding
plaquette makes each interchange preserve the product.  Backward edges use the
inverse positive ratio, so the same conclusion holds for arbitrary paths.
Thus the definition is path independent and satisfies every edge relation.
Uniqueness follows by propagation from the anchor.
\end{proof}

\begin{theorem}[Infinitesimal action-curl limit]
\label{thm:V30-curl-limit}
Suppose the edge ratios have the locally uniform expansion
\begin{equation}
 \log r_{i,h}(\theta,\omega)
 =h b_i(\theta,\omega)+h^2c_i(\theta,\omega)+o(h^2),
 \label{eq:V30-edge-log-expansion}
\end{equation}
where \(b_i\) is continuously differentiable and \(c_i\) is continuous in
\(\theta\).  Then
\begin{equation}
 \boxed{
 \frac{\Omega_{ij,h}(\theta)}{h^2}
 \longrightarrow
 \partial_i b_j(\theta)-\partial_j b_i(\theta)}
 \label{eq:V30-curl-limit}
\end{equation}
pointwise and in every norm in which the displayed expansion is uniform.  A
common differentiable unnormalized density has
\begin{equation}
 b_i=\partial_i\log a_\theta
 \label{eq:V30-score-potential}
\end{equation}
and therefore zero curl.  Conversely, on a simply connected command domain, a
\(C^1\) zero-curl field \((b_i)\) integrates, history by history, to a common
log-density up to one command-independent positive factor.
\end{theorem}

\begin{proof}
Insert \eqref{eq:V30-edge-log-expansion} into
\eqref{eq:V30-log-plaquette}.  Taylor expansion gives
\begin{equation}
 \Omega_{ij,h}
 =h^2(\partial_i b_j-\partial_j b_i)+o(h^2).
\end{equation}
The common-density implication is equality of mixed partial derivatives.  The
converse is the Poincare lemma applied pointwise in the history variable; the
integration constant is the positive anchor density.
\end{proof}

\begin{countertheorem}[Positive command edges do not imply one common cylinder]
\label{cth:V30-positive-edges-no-square}
On a one-point history space, take a two-command square with positive edge
ratios
\begin{equation}
 r_1(0,0)=2,
 \qquad
 r_2(0,0)=2,
 \qquad
 r_2(1,0)=3,
 \qquad
 r_1(0,1)=2.
 \label{eq:V30-incompatible-ratios}
\end{equation}
The first path predicts the terminal mass \(6a_0\), while the second predicts
\(4a_0\).  The plaquette residual is
\begin{equation}
 \boxed{\Delta_{12,1}^{\rm act}=|6-4|^2=4.}
 \label{eq:V30-incompatible-residual}
\end{equation}
Every one-parameter edge is strictly positive and every normalized law is the
same singleton law, but no common unnormalized command cylinder exists.  If the
last ratio is changed from \(2\) to \(3\), both paths give \(6a_0\) and the
common square is reconstructed exactly.
\end{countertheorem}

\begin{proof}
This is the two-path product test in
\eqref{eq:V30-multiplicative-plaquette}.  The normalized singleton law erases
all absolute mass information, so it cannot detect the failure.
\end{proof}

\section{Relation-covariant V29 command cards}
\label{sec:V30-covariant-command-card}

\begin{theorem}[Finite relation-covariant neighbouring-cylinder compiler]
\label{thm:V30-covariant-compiler}
Fix a finite command screen and one shrinking radius \(\varepsilon\).  Suppose:
\begin{enumerate}[label=\textup{(C\arabic*)},leftmargin=3.4em]
\item the positive command ratios on the required finite command grid pass all
      common-action plaquettes;
\item the returned-word Grams \(H_\theta\) have one stable rank and their
      canonical support projections \(P_\theta\) satisfy
      \(\norm{P_\theta-P_0}<1\);
\item every direct root or response amplitude \(X_\theta\) passes the endpoint
      relation test \(X_\theta=X_\theta P_\theta\);
\item the supported physical internal transport \(G_\theta\), when required by
      the target, is independently supplied; otherwise only the Kato-normal
      quotient tangent is claimed;
\item the determinant/Pfaffian line sewing and every orientation row pass on
      the same command grid.
\end{enumerate}
Then replacing every raw neighbouring response in Version~V29 by
\begin{equation}
 \widetilde X_\theta
 =X_\theta U_{P_\theta\leftarrow P_0}G_\theta
 \label{eq:V30-aligned-response}
\end{equation}
produces a finite command and response Gram which:
\begin{enumerate}[label=\textup{(\roman*)},leftmargin=2.8em]
\item descends exactly through the base returned-word quotient;
\item is reconstructed by the same positive module-distance formulas as in
      Version~V29;
\item converges, under the Version~V29 remainder bounds, to the physically
      transported tangent Gram;
\item separates raw relation motion, supported connection motion, and genuine
      quotient response.
\end{enumerate}
If item~\textup{(C1)} fails, the output is a common-action obstruction.  If
item~\textup{(C2)} fails, the output is a support or divisor event.  If
item~\textup{(C3)} fails, the output is a word/score occurrence obstruction.
None is a source birth.
\end{theorem}

\begin{proof}
Common-action reconstruction is \Cref{thm:V30-common-action-grid}.  The polar
chords of \Cref{thm:V30-polar-chord} identify all endpoint supports with the
base support.  Endpoint descent then gives exact finite descent by
\Cref{thm:V30-aligned-secant}.  Multiplication by the supported unitary
\(G_\theta\) preserves that relation.  The positive-distance reconstruction
and the shrinking-family limit are the Version~V29 arguments applied to the
aligned amplitudes.  The decomposition into the three stated motions is
\Cref{thm:V30-raw-Pythagoras,cor:V30-internal-connection}.
\end{proof}

\begin{corollary}[Exact finite relation residual hierarchy]
\label{cor:V30-residual-hierarchy}
For a raw endpoint response define
\begin{equation}
 \Delta_\theta^{\rm end}
 =\norm{X_\theta(I-P_\theta)}_{\mathrm{HS}}^2.
 \label{eq:V30-endpoint-relation-residual}
\end{equation}
For a raw central secant \(\mathscr S_\varepsilon^{\rm raw}\), define
\begin{equation}
 \Delta_\varepsilon^{\rm raw}
 =\Tr\left[(I-P_0)
 (\mathscr S_\varepsilon^{\rm raw})^*
 \mathscr S_\varepsilon^{\rm raw}\right].
 \label{eq:V30-raw-secant-residual}
\end{equation}
For the aligned secant define \(\Delta_\varepsilon^{\rm cov}\) analogously.
Then:
\begin{enumerate}[label=\textup{(R\arabic*)},leftmargin=3em]
\item a positive endpoint residual is a direct word/response occurrence
      failure;
\item zero endpoint residuals and positive
      \(\Delta_\varepsilon^{\rm raw}\) with
      \(\Delta_\varepsilon^{\rm cov}=0\) identify pure relation-support motion;
\item a positive aligned residual identifies a failed support transport or a
      response which does not descend through the transported relation;
\item only after all three rows pass is the Version~V29 null-cost and reserve
      test applied.
\end{enumerate}
\end{corollary}

\begin{proof}
Items~\textup{(R1)--(R3)} are the definitions together with
\Cref{thm:V30-aligned-secant,thm:V30-raw-Pythagoras}.  Item~\textup{(R4)} is the
required order of the quotient and metric operations.
\end{proof}

\section{Command--cutoff rectangles at a shrinking radius}
\label{sec:V30-command-cutoff-rectangle}

The shrinking command scale and the regulator cutoff form a second rectangle.
A small endpoint transport error need not remain small after division by the
command radius.

\begin{theorem}[Derivative-scaled endpoint transport]
\label{thm:V30-scaled-transport}
Let \(F_{n,+}(\varepsilon)\) and \(F_{n,-}(\varepsilon)\) be the completely
aligned endpoint amplitudes at cutoff \(n\), including the physical supported
connection.  Let \(J_n\) and \(V_n\) transport the base command and output
carriers to cutoff \(n+1\).  Assume
\begin{equation}
 \norm{F_{n+1,\pm}(\varepsilon)J_n
       -V_nF_{n,\pm}(\varepsilon)}
 \le\delta_{n,\pm}(\varepsilon).
 \label{eq:V30-endpoint-route-errors}
\end{equation}
For the central secants
\begin{equation}
 S_{n,\varepsilon}
 =\frac{F_{n,+}(\varepsilon)-F_{n,-}(\varepsilon)}{2\varepsilon},
 \label{eq:V30-cutoff-secants}
\end{equation}
one has the exact bound
\begin{equation}
 \boxed{
 \norm{S_{n+1,\varepsilon}J_n-V_nS_{n,\varepsilon}}
 \le
 \frac{
 \delta_{n,+}(\varepsilon)+
 \delta_{n,-}(\varepsilon)}{2\varepsilon}.}
 \label{eq:V30-scaled-route-bound}
\end{equation}
Thus a cofinal shrinking family requires derivative-scaled route control.  In
particular, endpoint defects merely tending to zero are insufficient.
\end{theorem}

\begin{proof}
Subtract the two endpoint transport identities, divide by
\(2\varepsilon\), and use the triangle inequality.
\end{proof}

\begin{countertheorem}[Vanishing endpoint error with order-one tangent mismatch]
\label{cth:V30-vanishing-endpoint-not-tangent}
Let \(F_{0,\pm}(\varepsilon)=0\) and
\(F_{1,\pm}(\varepsilon)=\pm\varepsilon\) on one-dimensional carriers, with
identity cutoff maps.  Then
\begin{equation}
 \delta_\pm(\varepsilon)=\varepsilon\longrightarrow0,
 \label{eq:V30-endpoint-error-vanishes}
\end{equation}
but
\begin{equation}
 \boxed{S_{1,\varepsilon}-S_{0,\varepsilon}=1}
 \label{eq:V30-tangent-mismatch-one}
\end{equation}
for every \(\varepsilon>0\).  Endpoint transport convergence therefore does
not imply tangent transport convergence.
\end{countertheorem}

\begin{proof}
The endpoint differences have absolute value \(\varepsilon\), whereas the
central secants are respectively zero and one.
\end{proof}

\begin{corollary}[Cofinal relation-covariant tangent landing]
\label{cor:V30-cofinal-landing}
Choose \(\varepsilon_n\downarrow0\).  Suppose:
\begin{enumerate}[label=\textup{(L\arabic*)},leftmargin=3.2em]
\item the common-action plaquette, endpoint descent, support-angle, line-sewing,
      and held-out command rows pass at every selected card;
\item the aligned central-difference remainder is \(r_n\to0\);
\item the derivative-scaled adjacent defects satisfy
      \begin{equation}
       \sum_n
       \frac{\delta_{n,+}(\varepsilon_n)
             +\delta_{n,-}(\varepsilon_n)}{\varepsilon_n}<\infty;
       \label{eq:V30-summable-scaled-route}
      \end{equation}
\item changing the command radius from \(\varepsilon_n\) to
      \(\varepsilon_{n+1}\) on the common fine card has a summable aligned
      secant defect;
\item the Version~V29 command metric support and its positive floor are stable.
\end{enumerate}
Then the relation-covariant command and response tangents have unique
fixed-screen limits.  Their metric-native reserve converges by the
Version~V29 perturbation theorem.  A positive limiting reserve is a physical
command-response survivor; a zero limiting reserve is an asymptotic follower
only after the complete V25 primitive short and the later semigroup-cyclic
short are applied in order.
\end{corollary}

\begin{proof}
Theorem~\ref{thm:V30-scaled-transport} and item~\textup{(L4)} make the diagonal
aligned secants Cauchy.  Item~\textup{(L2)} identifies their limits with the
physical tangents.  The common supported metric floor permits the
Moore--Penrose reserve perturbation estimate of Version~V29.  The final
ancestry qualification is the nested-short order already established in
Versions~V5--V6 and V25.
\end{proof}

\section{Exact replay and revised finite acquisition}
\label{sec:V30-replay}

At the rational radius
\begin{equation}
 \varepsilon=\frac34,
 \label{eq:V30-replay-epsilon}
\end{equation}
the moving-line projections of
\Cref{cth:V30-raw-secant-spurious} are
\begin{equation}
 \boxed{
 P_+=\frac1{25}
 \begin{pmatrix}16&12\\12&9\end{pmatrix},
 \qquad
 P_-=\frac1{25}
 \begin{pmatrix}16&-12\\-12&9\end{pmatrix}.}
 \label{eq:V30-replay-projectors}
\end{equation}
The polar chord columns are
\begin{equation}
 \boxed{
 U_+e_1=\frac15\binom43,
 \qquad
 U_-e_1=\frac15\binom4{-3}.}
 \label{eq:V30-replay-chords}
\end{equation}
Both aligned endpoint amplitudes equal \(5/4\), so their secant vanishes
exactly, while the raw null-word secant is one.

The incompatible command square
\eqref{eq:V30-incompatible-ratios} has path products six and four and squared
plaquette residual four.  Replacing the last ratio by three makes both path
products six and reconstructs a unique common terminal density.

The derivative-scaled route control is replayed with endpoint discrepancy
\(\varepsilon\) and tangent discrepancy one.  The verification script checks
all projector, chord, Pythagoras, plaquette, and scaled-route identities using
exact rational arithmetic.

\begin{theorem}[Version V30 relation-covariant command theorem]
\label{thm:V30-terminal}
Versions~V1--V29 are retained.  Version~V30 additionally proves:
\begin{enumerate}[label=\textup{(V30.\arabic*)},leftmargin=6.1em]
\item a constant-rank physical word support has a canonical Kato normal
      transport;
\item the raw response tangent splits orthogonally into a quotient-covariant
      tangent and an exact moving-relation debit;
\item a finite polar chord aligns nearby supports and makes every neighbouring
      secant descend exactly through the base relation quotient;
\item a raw secant can assign positive action entirely to a moving zero word;
\item finite command plaquettes are necessary and sufficient for one common
      positive action cylinder on a simply connected command grid;
\item the plaquette converges to the command-score curl, while a positive edge
      bank alone need not possess a common action;
\item endpoint cutoff errors are amplified by the reciprocal shrinking command
      radius, so derivative-scaled rectangle control is necessary for cofinal
      tangent transport;
\item after these rows pass, the Version~V29 metric-native reserve lands on the
      physically transported relation quotient.
\end{enumerate}
\end{theorem}

\begin{proof}
Items~\textup{(V30.1)--(V30.4)} are
\Cref{thm:V30-Kato-support,thm:V30-raw-Pythagoras,thm:V30-polar-chord,thm:V30-aligned-secant,cth:V30-raw-secant-spurious}.
Items~\textup{(V30.5)--(V30.6)} are
\Cref{thm:V30-common-action-grid,thm:V30-curl-limit,cth:V30-positive-edges-no-square}.
Items~\textup{(V30.7)--(V30.8)} are
\Cref{thm:V30-scaled-transport,cth:V30-vanishing-endpoint-not-tangent,cor:V30-cofinal-landing}.
\end{proof}

\begin{assessment}[Version V30 verdict]
\label{ass:V30-verdict}
The current result is
\begin{equation}
 \boxed{
 \begin{gathered}
 \passstatus\\[-0.15em]
 \text{Kato relation transport, raw/covariant tangent Pythagoras, finite polar}\\
 \text{support alignment, common-action plaquettes, score-curl limit, and}\\
 \text{derivative-scaled command--cutoff transport}
 \\[0.45em]
 \obsstatus\\[-0.15em]
 \text{differentiating in a fixed ambient coefficient frame while word relations}\\
 \text{move; inferring a common action from separate command rays; inferring a}\\
 \text{physical supported connection from support projectors; or transporting}\\
 \text{shrinking secants with unscaled endpoint errors}
 \\[0.45em]
 \stopstatus\\[-0.15em]
 \text{one shrinking RPESM command plaquette with returned-word support Grams,}\\
 \text{direct source responses, physical supported transport, line sewing,}\\
 \text{held-out command and word Reads, and derivative-scaled cutoff rectangles.}
 \end{gathered}}
 \label{eq:V30-verdict}
\end{equation}
The first missing row is no longer merely a shrinking list of neighbouring
laws.  It is a shrinking, common-action, relation-covariant command square on
the physical returned-word quotient.
\end{assessment}

\begin{acquisition}[Version V31: common-action relation-covariant RPESM command square]
\label{acq:V30-next}
Preserve Versions~V1--V30.  On one physically readable RPESM boundary record,
one amplitude semantics, and a predeclared two- or higher-dimensional command
screen, acquire a shrinking family of elementary command squares.  Each square
must contain:
\begin{enumerate}[label=\textup{(A\arabic*)},leftmargin=3em]
\item the unnormalized positive corner cylinders and all same-history edge
      likelihood ratios;
\item every multiplicative command plaquette and one held-out larger rectangle;
\item the complete returned-word Gram at every corner, its canonical support,
      rank, and supported floor;
\item the direct root and source-response amplitudes with their endpoint
      relation residuals;
\item the Kato polar chords and, independently, the supported physical
      reflection/line/flavour transport needed by the target;
\item the aligned Hellinger and direct-response module distances, one held-out
      command superposition, and one held-out word;
\item one adjacent cutoff square with endpoint route errors reported in units of
      the shrinking command radius;
\item every determinant/Pfaffian sewing, divisor, index-wall, orientation, or
      support transition encountered by the square.
\end{enumerate}
Only after the plaquette, relation, support, physical-connection, and scaled
rectangle rows pass may the card populate the Version~V29 tangent metric and
Version~V28 reserve.  The complete primitive short of Version~V25 and the
semigroup-cyclic Duhamel short remain later operations.

Do not open another phase, crossing, exterior, localizer, Duhamel, locality,
current/stress, charge, or Einstein compiler before this command square is
populated.
\end{acquisition}

\section*{V30 exact replay and append-only record}
\addcontentsline{toc}{section}{V30 exact replay and append-only record}
The accompanying verifier checks the rational moving-relation projectors, polar
chords, raw/covariant Pythagoras, exact disappearance of the false null-word
secant, compatible and incompatible command plaquettes, common-density path
reconstruction, and the reciprocal-radius amplification of endpoint transport
errors.  These are theorem replays and are not measurements of the selected
RPESM regulator.

The complete Version~V29 source preceding its sole final active
\verb|\end{document}| is preserved byte for byte.  Its preterminal SHA--256
digest is recorded in the validation manifest.  A later continuation must
remove only the final active document-closing command, append new material
immediately before it, restore one terminal command, and increment the filename
to end in \texttt{\_V31.tex}.

% =============================================================================
% END APPENDED VERSION V30 MATERIAL
% =============================================================================

\end{document}
